% Vertex Operator Algebras in Lean -- The Brief: the statements of the release edition of the book, with their
% Lean names and verbatim Lean, without proofs.  Numbers, pages and cross-references are those of book.pdf.
% Build from this folder (twice, for the table of contents):
%   lualatex brief.tex && lualatex brief.tex
%
\documentclass[11pt,openany]{book}
\usepackage{fontspec}
\usepackage{amsmath,amsthm,amssymb,mathtools}
\usepackage[math-style=TeX,bold-style=TeX]{unicode-math}
\setmainfont{STIX Two Text}
\setmathfont{STIX Two Math}
% Monospace font for Lean source: JuliaMono (SIL OFL, distributed with TeX Live),
% chosen because it covers the Unicode used in Lean/Mathlib (subscripts, script letters, ⟦⟧, ⨆, …).
\setmonofont{JuliaMono}[Extension=.ttf, UprightFont=*-Regular, BoldFont=*-Bold, Scale=0.80]
\usepackage{fvextra}
\usepackage[most]{tcolorbox}
\usepackage{xcolor}
\usepackage{microtype}
\usepackage[a4paper,margin=2.6cm]{geometry}
\usepackage{enumitem}
\usepackage{booktabs,longtable,array}
\usepackage[colorlinks=true,linkcolor=blue!55!black,citecolor=green!35!black,urlcolor=blue!65!black,
  bookmarksnumbered=true,bookmarksopen=true,bookmarksopenlevel=1]{hyperref}
\usepackage{cleveref}

\definecolor{leanbg}{RGB}{247,247,250}
\definecolor{leanframe}{RGB}{140,140,170}
\definecolor{notebg}{RGB}{250,248,240}

% Lean source: \begin{leancode} ... \end{leancode}  (verbatim, Unicode-safe, breaks long lines)
\DefineVerbatimEnvironment{leancode}{Verbatim}{breaklines=true,breakanywhere=true,
  fontsize=\small,frame=leftline,framerule=1.2pt,rulecolor=\color{leanframe},
  framesep=6pt,xleftmargin=4pt}
% Non-Lean verbatim text (shell commands, tool output, config files): \begin{plaincode}...\end{plaincode}
\DefineVerbatimEnvironment{plaincode}{Verbatim}{breaklines=true,breakanywhere=true,
  fontsize=\small,frame=leftline,framerule=0.6pt,rulecolor=\color{black!40},framesep=6pt,xleftmargin=4pt}
% inline Lean identifier: \lean{VOA.Moonshine.moonshine}
\newcommand{\lean}[1]{\texttt{\detokenize{#1}}}
\newcommand{\leanpath}[1]{\texttt{\small #1}}

% Explanatory boxes
\newtcolorbox{dictionary}[1][]{enhanced,breakable,colback=notebg,colframe=black!45,
  title={From the mathematics to the Lean statement},fonttitle=\bfseries\small,boxrule=0.4pt,arc=2pt,#1}
\newtcolorbox{warning}[1][]{enhanced,breakable,colback=red!3,colframe=red!50!black,
  title={Trap / dead end},fonttitle=\bfseries\small,boxrule=0.4pt,arc=2pt,#1}
\newtcolorbox{status}[1][]{enhanced,breakable,colback=green!3,colframe=green!35!black,
  title={What the machine verified},fonttitle=\bfseries\small,boxrule=0.4pt,arc=2pt,#1}
\newtcolorbox{physics}[1][]{enhanced,breakable,colback=blue!3,colframe=blue!45!black,
  title={Physical meaning},fonttitle=\bfseries\small,boxrule=0.4pt,arc=2pt,#1}

\theoremstyle{definition}
\newtheorem{definition}{Definition}[chapter]
\newtheorem{example}[definition]{Example}
\newtheorem{remark}[definition]{Remark}
\theoremstyle{plain}
\newtheorem{theorem}[definition]{Theorem}
\newtheorem{proposition}[definition]{Proposition}
\newtheorem{lemma}[definition]{Lemma}
\newtheorem{corollary}[definition]{Corollary}

\newcommand{\C}{{\mathbb{C}}}
\newcommand{\Z}{{\mathbb{Z}}}
\newcommand{\Q}{{\mathbb{Q}}}
\newcommand{\N}{{\mathbb{N}}}
\newcommand{\Vnat}{{V^{\natural}}}
\newcommand{\Leech}{{\Lambda}}
\newcommand{\End}{\operatorname{End}}
\newcommand{\Aut}{\operatorname{Aut}}
\newcommand{\Res}{\operatorname{Res}}
\newcommand{\wt}{\operatorname{wt}}
\newcommand{\tr}{\operatorname{tr}}

% Edition
\newcommand{\bookrelease}{release v1.0}
\newcommand{\bookcommit}{(recorded at release)}
\newcommand{\bookdate}{2026-11-03}

% Brief statement environments carry the book's own number: \begin{btheorem}{5.4}[title] ... \end{btheorem}
\newtheoremstyle{briefplain}{6pt}{3pt}{\itshape}{}{\bfseries}{.}{0.5em}{\thmname{#1}~\thmnote{#3}}
\newtheoremstyle{briefdef}{6pt}{3pt}{\normalfont}{}{\bfseries}{.}{0.5em}{\thmname{#1}~\thmnote{#3}}
\makeatletter
\newcommand{\brieftheoremkind}[3]{% #1 env, #2 printed name, #3 style
  \theoremstyle{#3}\newtheorem*{#1@inner}{#2}%
  \newenvironment{#1}[1]{\def\briefnum{##1}\@ifnextchar[{\csname #1@opt\endcsname}{\begin{#1@inner}[\briefnum]}}{\end{#1@inner}}%
  \expandafter\def\csname #1@opt\endcsname[##1]{\begin{#1@inner}[{\briefnum\ (##1)}]}}
\makeatother
\brieftheoremkind{bdefinition}{Definition}{briefdef}
\brieftheoremkind{btheorem}{Theorem}{briefplain}
\brieftheoremkind{bproposition}{Proposition}{briefplain}
\brieftheoremkind{blemma}{Lemma}{briefplain}
\brieftheoremkind{bcorollary}{Corollary}{briefplain}
\title{\Huge Vertex Operator Algebras in Lean\\ \vspace{0.4em}\LARGE The Brief\\ \vspace{0.6em}
  \large Statements, Lean names and verbatim Lean, without proofs}
\author{Fedor Popov and Leonardo Rastelli\\ \vspace{0.3em}\small with Claude (Anthropic)}
\date{Companion to the release edition of the book; documents commit \texttt{\bookcommit}\\ \small(built \bookdate)}
\begin{document}
\frontmatter
\maketitle
\chapter*{How to read the brief}
This is a companion to the release edition of \emph{Vertex Operator Algebras in Lean}
(\texttt{book.pdf}). For the
headline results of every chapter, and for every definition (in any chapter) that their statements
use, it prints only: the mathematical statement, as in the book; under \textbf{Lean.}, the fully
qualified names of the Lean declarations that formalise it; the Lean file(s); every declaration of
the book's quotation, verbatim; and the hypotheses and status recorded in the chapter's Referee's
checklist. Chapter, theorem and page numbers, and all cross-references, are the \emph{book's}, so
every item can be looked up there with its proof, its fidelity discussion (``From the mathematics to
the Lean statement'') and its physical meaning. Chapter~1 of the book (how to read Lean) has no
statements and is omitted.

Every quoted Lean excerpt occurs verbatim (up to whitespace and comments) in the library at the
commit above. The qualified names are computed from the \texttt{namespace} blocks of the cited files.

A clean \texttt{\#print axioms} does not make a theorem unconditional: explicit hypotheses are
listed in the status line. The library has no project axiom (\texttt{VOA.pbw} is a theorem)
(book Chapter~3).
\tableofcontents
\mainmatter
\pagestyle{plain}% chapter titles are too long for running heads
\part{Foundations}
\setcounter{chapter}{1}
\chapter{Formal calculus and locality}
\begingroup
\providecommand{\loneW}{\mathcal{W}}
\providecommand{\loneV}{V}
\providecommand{\loneiota}[2]{\iota_{#1,#2}}
\providecommand{\lonedelta}{\delta(z_1,z_2)}
\providecommand{\lonecoef}[1]{\big[#1\big]}
\providecommand{\lonebinom}[2]{\binom{#1}{#2}}
\providecommand{\lonewarn}[1]{the trap box on p.~\pageref{#1}}
\noindent{\small 10 of the chapter's 56 numbered statements: its headline results and every definition they use.}\par\medskip
\begin{bdefinition}{2.1}[Formal distribution]
Let $\loneW$ be a set. A \emph{formal distribution} in one variable $z$ with coefficients in
$\loneW$ is a function $f:\Z\to\loneW$, written as the doubly infinite formal series
\begin{equation*}
  f(z) = \sum_{n\in\Z} f_n\,z^n,\qquad f_n := f(n)\in\loneW .
\end{equation*}
The set of formal distributions is denoted $\loneW[[z^{\pm1}]]$. When $\loneW$ is an abelian group,
respectively a module over a semiring $R$, $\loneW[[z^{\pm1}]]$ is an abelian group, respectively an
$R$-module, under coefficientwise addition and scalar multiplication. Two formal distributions are
equal iff all their coefficients agree. No product of two formal distributions is defined.
\end{bdefinition}
{\raggedright\noindent\textbf{Lean.} \lean{VOA.}\allowbreak\lean{FormalDist},\allowbreak\ \lean{VOA.}\allowbreak\lean{FormalDist.}\allowbreak\lean{instAddCommGroup},\allowbreak\ \lean{VOA.}\allowbreak\lean{FormalDist.}\allowbreak\lean{instModule},\allowbreak\ \lean{VOA.}\allowbreak\lean{FormalDist.}\allowbreak\lean{ext}\par}
{\raggedright\noindent File \leanpath{VOA/\allowbreak Foundations/\allowbreak FormalSeries.lean}.\par}
\begin{leancode}
def FormalDist (𝒲 : Type*) : Type _ := ℤ → 𝒲

@[inherit_doc] notation:max 𝒲 "⟦z^±⟧" => FormalDist 𝒲

instance instAddCommGroup [AddCommGroup 𝒲] : AddCommGroup (FormalDist 𝒲) :=
  inferInstanceAs (AddCommGroup (ℤ → 𝒲))

instance instModule {R : Type*} [Semiring R] [AddCommGroup 𝒲] [Module R 𝒲] :
    Module R (FormalDist 𝒲) := inferInstanceAs (Module R (ℤ → 𝒲))

@[ext]
theorem ext {f g : FormalDist 𝒲} (h : ∀ n, f n = g n) : f = g := funext h
\end{leancode}
\noindent\textbf{Hypotheses / status.} Definition\par
\medskip
\begin{bdefinition}{2.7}[Modes, residue, derivative]
Let $f\in\loneW[[z^{\pm1}]]$.
\begin{enumerate}[label=(\roman*)]
\item The \emph{$m$-th mode} of $f$ ($m\in\Z$) is $f_{(m)} := f_{-m-1}$, so that
  $f(z) = \sum_{m\in\Z}f_{(m)}\,z^{-m-1}$.
\item If $\loneW$ is a complex vector space, the \emph{residue} is the linear functional
  $\Res_z f := f_{-1}$, the coefficient of $z^{-1}$.
\item The \emph{formal derivative} is the linear map $\partial_z$ with
  $(\partial_zf)_n := (n+1)f_{n+1}$, i.e.\ $\partial_z\sum_n f_nz^n = \sum_n nf_nz^{n-1}$.
\end{enumerate}
\end{bdefinition}
{\raggedright\noindent\textbf{Lean.} \lean{VOA.}\allowbreak\lean{mode},\allowbreak\ \lean{VOA.}\allowbreak\lean{res},\allowbreak\ \lean{VOA.}\allowbreak\lean{deriv}\par}
{\raggedright\noindent File \leanpath{VOA/\allowbreak Foundations/\allowbreak FormalCalculus.lean}.\par}
\begin{leancode}
def mode (f : FormalDist 𝒲) (m : ℤ) : 𝒲 := f (-m - 1)

def res : FormalDist 𝒲 →ₗ[ℂ] 𝒲 where
  toFun f := f (-1)

def deriv : FormalDist 𝒲 →ₗ[ℂ] FormalDist 𝒲 where
  toFun f := fun n => (n + 1 : ℤ) • f (n + 1)
\end{leancode}
\noindent\textbf{Hypotheses / status.} Definition; \textbf{key convention}\par
\medskip
\begin{bproposition}{2.8}[Modes are residues; residues of derivatives vanish]
For every $f\in\loneW[[z^{\pm1}]]$ and $m\in\Z$,
\begin{equation*}
  f_{(m)} = \Res_z\big(z^mf(z)\big),\qquad \Res_z\big(\partial_zf\big) = 0 .
\end{equation*}
\end{bproposition}
{\raggedright\noindent\textbf{Lean.} \lean{VOA.}\allowbreak\lean{mode_eq_res},\allowbreak\ \lean{VOA.}\allowbreak\lean{res_deriv}\par}
{\raggedright\noindent File \leanpath{VOA/\allowbreak Foundations/\allowbreak FormalCalculus.lean}.\par}
\begin{leancode}
theorem mode_eq_res (f : FormalDist 𝒲) (m : ℤ) : mode f m = res (mulByMonomial m f) := by

@[simp] theorem res_deriv (f : FormalDist 𝒲) : res (deriv f) = 0 := by
\end{leancode}
\noindent\textbf{Hypotheses / status.} None; unconditional\par
\medskip
\begin{btheorem}{2.12}[Integration by parts; Gannon (5.1.3b)]
If $g\in\C[z^{\pm1}]$ is a Laurent polynomial and $f\in\loneW[[z^{\pm1}]]$, then
\begin{equation*}
  \Res_z\big(g\cdot\partial_zf\big) = -\Res_z\big((\partial_zg)\cdot f\big).
\end{equation*}
\end{btheorem}
{\raggedright\noindent\textbf{Lean.} \lean{VOA.}\allowbreak\lean{res_lpSMul_deriv}\par}
{\raggedright\noindent File \leanpath{VOA/\allowbreak Foundations/\allowbreak FormalCalculus.lean}.\par}
\begin{leancode}
theorem res_lpSMul_deriv (g : FormalDist ℂ) (hg : (Function.support g).Finite)
    (f : FormalDist 𝒲) : res (lpSMul g (deriv f)) = - res (lpSMul (deriv g) f) := by
\end{leancode}
\noindent\textbf{Hypotheses / status.} $g$ Laurent polynomial\par
\medskip
\begin{btheorem}{2.27}[Structure theorem; Gannon Prop.~5.1.2, Kac Thm.~2.3]
Let $\loneW$ be a complex vector space, $f\in\loneW[[z_1^{\pm1},z_2^{\pm1}]]$ and $N\in\N$. Then
\begin{equation*}
  (z_1-z_2)^Nf = 0\quad\Longleftrightarrow\quad
  \exists\,c^0,\dots,c^{N-1}\in\loneW[[z^{\pm1}]]:\ f(z_1,z_2) = \sum_{j=0}^{N-1}c^j(z_2)\,\partial_{z_2}^j\lonedelta .
\end{equation*}
Moreover, in the forward direction the coefficients can be taken to be
$c^j = \frac1{j!}\Res_{z_1}(z_1-z_2)^jf$ (Corollary~2.28).
\end{btheorem}
{\raggedright\noindent\textbf{Lean.} \lean{VOA.}\allowbreak\lean{exists_deltaSum_of_mulBySubPow_eq_zero},\allowbreak\ \lean{VOA.}\allowbreak\lean{local_iff_deltaSum}\par}
{\raggedright\noindent File \leanpath{VOA/\allowbreak Foundations/\allowbreak Locality.lean}.\par}
\begin{leancode}
theorem exists_deltaSum_of_mulBySubPow_eq_zero (N : ℕ) (f : FormalDist₂ 𝒲)
    (hf : mulBySubPow N f = 0) :
    ∃ c : Fin N → FormalDist 𝒲, f = ∑ j : Fin N, deltaDeriv (j : ℕ) (c j) := by

theorem local_iff_deltaSum (f : FormalDist₂ 𝒲) (N : ℕ) :
    mulBySubPow N f = 0 ↔
      ∃ c : Fin N → FormalDist 𝒲, f = ∑ j : Fin N, deltaDeriv (j : ℕ) (c j) := by
\end{leancode}
\noindent\textbf{Hypotheses / status.} None; unconditional\par
\medskip
\begin{bdefinition}{2.30}[The two expansions of $(z_1-z_2)^n$]
For $n\in\Z$ define two elements of $\C[[z_1^{\pm1},z_2^{\pm1}]]$:
\begin{align*}
  \loneiota{z_1}{z_2}(z_1-z_2)^n &:= \sum_{k\ge0}(-1)^k\binom nk\,z_1^{n-k}z_2^{k}
  &&\text{(nonnegative powers of $z_2$; ``$|z_1|>|z_2|$'')},\\
  \loneiota{z_2}{z_1}(z_1-z_2)^n &:= \sum_{k\ge0}(-1)^{n+k}\binom nk\,z_1^{k}z_2^{n-k}
  &&\text{(nonnegative powers of $z_1$; ``$|z_2|>|z_1|$'')}.
\end{align*}
Also let $\lonedelta = \sum_kz_1^kz_2^{-k-1}$ be the delta function, with coefficient $[m+n+1 = 0]$.
\end{bdefinition}
{\raggedright\noindent\textbf{Lean.} \lean{VOA.}\allowbreak\lean{expandPos},\allowbreak\ \lean{VOA.}\allowbreak\lean{expandNeg},\allowbreak\ \lean{VOA.}\allowbreak\lean{deltaForm}\par}
{\raggedright\noindent File \leanpath{VOA/\allowbreak Foundations/\allowbreak Expansion.lean}.\par}
\begin{leancode}
noncomputable def expandPos (n : ℤ) : FormalDist₂ ℂ :=
  fun a b => if 0 ≤ b ∧ a = n - b then negOnePow b * gBinom n b.toNat else 0

noncomputable def expandNeg (n : ℤ) : FormalDist₂ ℂ :=
  fun a b => if 0 ≤ a ∧ b = n - a then negOnePow (n + a) * gBinom n a.toNat else 0

def deltaForm : FormalDist₂ ℂ := fun m n => if m + n + 1 = 0 then 1 else 0
\end{leancode}
\noindent\textbf{Hypotheses / status.} Definition\par
\medskip
\begin{btheorem}{2.33}[The delta-function jump; Gannon §5.1.2, FLM §8.10]
\begin{equation*}
  \loneiota{z_1}{z_2}\frac1{z_1-z_2}-\loneiota{z_2}{z_1}\frac1{z_1-z_2} = \lonedelta,
  \qquad (z_1-z_2)\,\lonedelta = 0 .
\end{equation*}
\end{btheorem}
{\raggedright\noindent\textbf{Lean.} \lean{VOA.}\allowbreak\lean{expand_jump},\allowbreak\ \lean{VOA.}\allowbreak\lean{mulBySub_deltaForm}\par}
{\raggedright\noindent File \leanpath{VOA/\allowbreak Foundations/\allowbreak Expansion.lean}.\par}
\begin{leancode}
theorem expand_jump : expandPos (-1) - expandNeg (-1) = deltaForm := by

theorem mulBySub_deltaForm : mulBySub (deltaForm : FormalDist₂ ℂ) = 0 := by
\end{leancode}
\noindent\textbf{Hypotheses / status.} None; unconditional\par
\medskip
\begin{btheorem}{2.36}[Multiplicativity of the expansions]
For all $m,n\in\Z$,
\begin{equation*}
  \iota_{12}(z_1-z_2)^m\cdot\iota_{12}(z_1-z_2)^n = \iota_{12}(z_1-z_2)^{m+n},\qquad
  \iota_{21}(z_1-z_2)^m\cdot\iota_{21}(z_1-z_2)^n = \iota_{21}(z_1-z_2)^{m+n}.
\end{equation*}
\end{btheorem}
{\raggedright\noindent\textbf{Lean.} \lean{VOA.}\allowbreak\lean{expandPos_mul},\allowbreak\ \lean{VOA.}\allowbreak\lean{expandNeg_mul}\par}
{\raggedright\noindent File \leanpath{VOA/\allowbreak Foundations/\allowbreak IotaZExpansion.lean}.\par}
\begin{leancode}
theorem expandPos_mul (m n : ℤ) : convolPos m n = expandPos (m + n) := by

theorem expandNeg_mul (m n : ℤ) : convolNeg m n = expandNeg (m + n) := by
\end{leancode}
\noindent\textbf{Hypotheses / status.} None; \lean{convolPos} is the genuine product (paper)\par
\medskip
\begin{btheorem}{2.38}[All-order delta jump]
For every $k\in\N$,
\begin{equation*}
  \iota_{12}(z_1-z_2)^{-k-1}-\iota_{21}(z_1-z_2)^{-k-1} = \frac1{k!}\,\partial_{z_2}^k\lonedelta .
\end{equation*}
\end{btheorem}
{\raggedright\noindent\textbf{Lean.} \lean{VOA.}\allowbreak\lean{mulBySub_negPowJump},\allowbreak\ \lean{VOA.}\allowbreak\lean{negPowJump_residue_zero},\allowbreak\ \lean{VOA.}\allowbreak\lean{deltaDeriv_succ_residue_zero},\allowbreak\ \lean{VOA.}\allowbreak\lean{negPowJump_eq_deltaDeriv}\par}
{\raggedright\noindent File \leanpath{VOA/\allowbreak Foundations/\allowbreak DeltaJump.lean}.\par}
\begin{leancode}
theorem mulBySub_negPowJump (k : ℕ) :
    mulBySub (negPowJump (k + 1)) = negPowJump k := by

theorem negPowJump_residue_zero (k : ℕ) (n : ℤ) :
    negPowJump (k + 1) (-1) n = 0 := by

theorem deltaDeriv_succ_residue_zero (k : ℕ) (c : FormalDist ℂ) (n : ℤ) :
    deltaDeriv (k + 1) c (-1) n = 0 := by

theorem negPowJump_eq_deltaDeriv (k : ℕ) :
    negPowJump k = (((k.factorial : ℂ))⁻¹) • deltaDeriv k deltaCoeff := by
\end{leancode}
\noindent\textbf{Hypotheses / status.} None; unconditional\par
\medskip
\begin{btheorem}{2.62}[Outer-to-inner transport]
For all orientations $o_1,o_2$ and $N\in\N$,
\begin{equation*}
  (z_{o_1}-z_{o_2})^N\,K_{o_1o_2} = (x_1-x_2)^N\,K_{o_1o_2}.
\end{equation*}
\end{btheorem}
{\raggedright\noindent\textbf{Lean.} \lean{VOA.}\allowbreak\lean{FormalDist₄Delta.}\allowbreak\lean{mulBySub₁₂_comm_mulBySub₃₄},\allowbreak\ \lean{VOA.}\allowbreak\lean{FormalDist₄Delta.}\allowbreak\lean{mulBySub₃₄_doubleDelta},\allowbreak\ \lean{VOA.}\allowbreak\lean{FormalDist₄Delta.}\allowbreak\lean{mulBySub₃₄Pow_doubleDelta}\par}
{\raggedright\noindent File \leanpath{VOA/\allowbreak Foundations/\allowbreak FormalDist4Delta.lean}.\par}
\begin{leancode}
theorem mulBySub₁₂_comm_mulBySub₃₄ :
    Commute
      (FormalDist₄.mulBySub₁₂ : FormalDist₄ (FormalDist ℂ) →ₗ[ℂ] _)
      (FormalDist₄.mulBySub₃₄ : FormalDist₄ (FormalDist ℂ) →ₗ[ℂ] _) := by

theorem mulBySub₃₄_doubleDelta (o₁ o₂ : Orientation) :
    FormalDist₄.mulBySub₃₄ (doubleDelta o₁ o₂) =
      FormalDist₄.mulBySub₁₂ (doubleDelta o₁ o₂) := by

theorem mulBySub₃₄Pow_doubleDelta (o₁ o₂ : Orientation) (N : ℕ) :
    FormalDist₄.mulBySub₃₄Pow N (doubleDelta o₁ o₂) =
      FormalDist₄.mulBySub₁₂Pow N (doubleDelta o₁ o₂) := by
\end{leancode}
\noindent\textbf{Hypotheses / status.} None; unconditional\par
\medskip
\endgroup
\setcounter{chapter}{2}
\chapter{Lie-theoretic inputs and the trust base}
\begingroup
\providecommand{\ltwog}{\mathfrak{g}}
\providecommand{\ltwoghat}{\widehat{\mathfrak{g}}}
\providecommand{\ltwop}{\mathfrak{p}}
\providecommand{\ltwon}{\mathfrak{n}}
\providecommand{\ltwovac}{|0\rangle}
\providecommand{\ltwoW}{\mathfrak{W}}
\providecommand{\ltwoVir}{\mathfrak{Vir}}
\providecommand{\ltwoUe}{\mathcal{U}}
\providecommand{\ltwoInd}{\operatorname{Ind}}
\providecommand{\ltwoJ}{\mathcal{J}}
\providecommand{\ltwoV}{\mathcal{V}}
\noindent{\small 7 of the chapter's 25 numbered statements: its headline results and every definition they use.}\par\medskip
\begin{btheorem}{3.9}[Poincar\'e--Birkhoff--Witt]
Let $\ltwog$ be a Lie algebra over a field $F$ (of any dimension), and let $(e_i)_{i\in I}$ be a
basis of $\ltwog$ indexed by a totally ordered set $I$. Then the ordered monomials
\begin{equation*}
 \iota(e_{i_1})\iota(e_{i_2})\cdots\iota(e_{i_r}),\qquad r\ge0,\quad i_1\le i_2\le\dots\le i_r,
\end{equation*}
form a basis of $U(\ltwog)$ (the empty monomial being $1$). Equivalently, the associated graded
algebra of $U(\ltwog)$ for its natural filtration is the symmetric algebra $S(\ltwog)$.
\end{btheorem}
{\raggedright\noindent\textbf{Lean.} \lean{VOA.}\allowbreak\lean{PBWPrinciple},\allowbreak\ \lean{VOA.}\allowbreak\lean{pbw}\par}
\noindent \leanpath{VOA/Foundations/AffineLie.lean}; the same declaration is now a theorem:
\begin{leancode}
structure PBWPrinciple : Prop where
  /-- PBW nonvanishing for every affine vacuum module. -/
  affineVacuum : ∀ {𝔤 : Type u} [LieRing 𝔤] [LieAlgebra ℂ 𝔤]
      (κ : LinearMap.BilinForm ℂ 𝔤) (hκ : κ.IsInvariant) (k : ℂ),
    PBW κ hκ k
  /-- The canonical bottom embeds in every module induced from a Lie subalgebra. -/
  inducedBottom : ∀ {g U : Type u} [LieRing g] [LieAlgebra ℂ g]
      [AddCommGroup U] [Module ℂ U] (p : LieSubalgebra ℂ g)
      (rho : p →ₗ⁅ℂ⁆ Module.End ℂ U),
    Function.Injective (PBWInduction.bottom p rho)

theorem pbw : PBWPrinciple :=
  ⟨fun κ hκ k => AffineLie.vacuumVector_ne_zero_of_induction κ hκ k,
    fun p rho => PBWTheorem.bottom_injective p rho⟩
\end{leancode}
\medskip
\begin{bproposition}{3.11}[The Lean induced module is the induced module]
Let $\ltwop\subseteq\ltwog$ be a Lie subalgebra, $\rho$ a representation of $\ltwop$ on $U$, and
$\varphi : U(\ltwop)\to U(\ltwog)$ the algebra map induced by the inclusion. Make $U(\ltwog)$ a right
$U(\ltwop)$-module through $\varphi$ and $U$ a left $U(\ltwop)$-module through $\rho$. Let
$N\subseteq U(\ltwog)\otimes_\C U$ be the span of $\{y\iota(x)\otimes u - y\otimes\rho(x)u : y\in U(\ltwog),
x\in\ltwop,u\in U\}$. Then $N$ equals the span $N'$ of
$\{y\varphi(a)\otimes u - y\otimes a\cdot u : y\in U(\ltwog), a\in U(\ltwop), u\in U\}$, so
\begin{equation*}
 \lean{Induced p rho} = (U(\ltwog)\otimes_\C U)/N = U(\ltwog)\otimes_{U(\ltwop)}U.
\end{equation*}
\end{bproposition}
{\raggedright\noindent\textbf{Lean.} \lean{VOA.}\allowbreak\lean{PBWInduction.}\allowbreak\lean{Cover},\allowbreak\ \lean{VOA.}\allowbreak\lean{PBWInduction.}\allowbreak\lean{balancingSet},\allowbreak\ \lean{VOA.}\allowbreak\lean{PBWInduction.}\allowbreak\lean{balancing},\allowbreak\ \lean{VOA.}\allowbreak\lean{PBWInduction.}\allowbreak\lean{Induced},\allowbreak\ \lean{VOA.}\allowbreak\lean{PBWInduction.}\allowbreak\lean{bottom}\par}
\noindent File \leanpath{VOA/Foundations/AffineLie.lean}.
\begin{leancode}
universe u

namespace PBWInduction

variable {g U : Type u} [LieRing g] [LieAlgebra ℂ g]
variable [AddCommGroup U] [Module ℂ U]

abbrev Cover (g U : Type u) [LieRing g] [LieAlgebra ℂ g]
    [AddCommGroup U] [Module ℂ U] :=
  UniversalEnvelopingAlgebra ℂ g ⊗[ℂ] U

def balancingSet (p : LieSubalgebra ℂ g) (rho : p →ₗ⁅ℂ⁆ Module.End ℂ U) :
    Set (Cover g U) :=
  {t | ∃ (y : UniversalEnvelopingAlgebra ℂ g) (x : p) (u : U),
    t = (y * UniversalEnvelopingAlgebra.ι ℂ (x : g)) ⊗ₜ[ℂ] u - y ⊗ₜ[ℂ] (rho x u)}

noncomputable def balancing (p : LieSubalgebra ℂ g)
    (rho : p →ₗ⁅ℂ⁆ Module.End ℂ U) : Submodule ℂ (Cover g U) :=
  Submodule.span ℂ (balancingSet p rho)

abbrev Induced (p : LieSubalgebra ℂ g) (rho : p →ₗ⁅ℂ⁆ Module.End ℂ U) : Type u :=
  Cover g U ⧸ balancing p rho

noncomputable def bottom (p : LieSubalgebra ℂ g)
    (rho : p →ₗ⁅ℂ⁆ Module.End ℂ U) : U →ₗ[ℂ] Induced p rho :=
  (balancing p rho).mkQ.comp
    (TensorProduct.mk ℂ (UniversalEnvelopingAlgebra ℂ g) U 1)

end PBWInduction
\end{leancode}
\noindent\textbf{Hypotheses / status.} Definitions\par
\medskip
\begin{btheorem}{3.15}[Relative PBW representation; Humphreys \S17.4, Lemma A, relative form]
Let $\ltwog$ be a Lie algebra over a field $K$ (of any dimension), $\ltwop\subseteq\ltwog$ a Lie subalgebra,
$\mathfrak c$ a vector-space complement of $\ltwop$ in $\ltwog$, $(y_a)_{a\in A}$ a basis of $\mathfrak c$ indexed by
a totally ordered set $A$, and $\rho$ a representation of $\ltwop$ on a $K$-vector space $U$. Let $M$ be the
vector space of finitely supported functions from finite multisets of $A$ to $U$; write $m\otimes u$ for
the function with value $u$ at the multiset $m$ and $0$ elsewhere, $|m|$ for the number of elements of
$m$ counted with multiplicity, and $a\le m$ for ``$a\le b$ for every letter $b$ of $m$''. Then there is a
representation $\sigma$ of $\ltwog$ on $M$ such that, for all $u\in U$, $x\in\ltwop$ and all multisets $m$:
\begin{align*}
 \text{(E1)}\quad& \sigma(y_a)(m\otimes u) = (a\,m)\otimes u &&\text{if } a\le m,\\
 \text{(E2)}\quad& \sigma(y_a)\bigl((b\,m)\otimes u\bigr) = \sigma(y_b)\sigma(y_a)(m\otimes u) + \sigma([y_a,y_b])(m\otimes u)
   &&\text{if } b<a,\ b\le m,\\
 \text{(E3)}\quad& \sigma(x)(\varnothing\otimes u) = \varnothing\otimes\rho(x)u,\\
 \text{(E4)}\quad& \sigma(x)\bigl((b\,m)\otimes u\bigr) = \sigma(y_b)\sigma(x)(m\otimes u) + \sigma([x,y_b])(m\otimes u)
   &&\text{if } b\le m,
\end{align*}
where $a\,m$ denotes the multiset $m$ with one more copy of $a$. In particular $u\mapsto\varnothing\otimes u$ is
an injective linear map $U\to M$ which intertwines $\rho(x)$ with $\sigma(x)$ for every $x\in\ltwop$.
\end{btheorem}
{\raggedright\noindent\textbf{Lean.} \lean{VOA.}\allowbreak\lean{PBWTheorem.}\allowbreak\lean{Car},\allowbreak\ \lean{VOA.}\allowbreak\lean{PBWTheorem.}\allowbreak\lean{act_y_low},\allowbreak\ \lean{VOA.}\allowbreak\lean{PBWTheorem.}\allowbreak\lean{act_p_zero},\allowbreak\ \lean{VOA.}\allowbreak\lean{PBWTheorem.}\allowbreak\lean{act_y_cons},\allowbreak\ \lean{VOA.}\allowbreak\lean{PBWTheorem.}\allowbreak\lean{act_p_cons},\allowbreak\ \lean{VOA.}\allowbreak\lean{PBWTheorem.}\allowbreak\lean{defect_eq_zero},\allowbreak\ \lean{VOA.}\allowbreak\lean{PBWTheorem.}\allowbreak\lean{rep},\allowbreak\ \lean{VOA.}\allowbreak\lean{PBWTheorem.}\allowbreak\lean{rep_apply},\allowbreak\ \lean{VOA.}\allowbreak\lean{PBWTheorem.}\allowbreak\lean{rep_p_zero}\par}
{\raggedright\noindent File \leanpath{VOA/\allowbreak Foundations/\allowbreak PBWTheorem/\allowbreak Construction.lean}.\par}
\begin{leancode}
variable {K : Type*} [Field K]
variable {g : Type*} [LieRing g] [LieAlgebra K g]
variable {U : Type*} [AddCommGroup U] [Module K U]
variable {A : Type*} [LinearOrder A]

abbrev Car (A U : Type*) [AddCommGroup U] [Module K U] := Multiset A →₀ U
\end{leancode}
\begin{leancode}
theorem act_y_low {a : A} {m : Multiset A} (h : ∀ b ∈ m, a ≤ b) (u : U) :
    act p hc bc ρ (bc a : g) (single m u) = single (a ::ₘ m) u := by

theorem act_p_zero (x : p) (u : U) :
    act p hc bc ρ (x : g) (single 0 u) = single 0 (ρ x u) := by

theorem act_y_cons {a b : A} {m : Multiset A} (hba : b < a) (hb : ∀ c ∈ m, b ≤ c) (u : U) :
    act p hc bc ρ (bc a : g) (single (b ::ₘ m) u) =
      act p hc bc ρ (bc b : g) (act p hc bc ρ (bc a : g) (single m u)) +
        act p hc bc ρ ⁅(bc a : g), (bc b : g)⁆ (single m u) := by

theorem act_p_cons (x : p) {b : A} {m : Multiset A} (hb : ∀ c ∈ m, b ≤ c) (u : U) :
    act p hc bc ρ (x : g) (single (b ::ₘ m) u) =
      act p hc bc ρ (bc b : g) (act p hc bc ρ (x : g) (single m u)) +
        act p hc bc ρ ⁅(x : g), (bc b : g)⁆ (single m u) := by
\end{leancode}
{\raggedright\noindent File \leanpath{VOA/\allowbreak Foundations/\allowbreak PBWTheorem/\allowbreak LieRep.lean}.\par}
\begin{leancode}
variable (p : LieSubalgebra K g) {c : Submodule K g} (hc : IsCompl p.toSubmodule c)
variable (bc : Basis A K c) (ρ : p →ₗ⁅K⁆ Module.End K U)

theorem defect_eq_zero (z z' : g) (v : Car (K := K) A U) : defect p hc bc ρ z z' v = 0 := by

noncomputable def rep : g →ₗ⁅K⁆ Module.End K (Car (K := K) A U) :=

theorem rep_apply (z : g) (v : Car (K := K) A U) : rep p hc bc ρ z v = act p hc bc ρ z v := rfl

theorem rep_p_zero (x : p) (u : U) :
    rep p hc bc ρ (x : g) (single 0 u) = single 0 (ρ x u) :=
  act_p_zero p hc bc ρ x u
\end{leancode}
\noindent\textbf{Hypotheses / status.} Unconditional, any field\par
\medskip
\begin{btheorem}{3.17}[Bottom injectivity, relative PBW]
Let $\ltwog$ be a Lie algebra over a field $K$, $\ltwop$ a Lie subalgebra and $\rho$ a representation of $\ltwop$
on $U$. Then the map $U\to(U(\ltwog)\otimes_K U)/N$, $u\mapsto[1\otimes u]$, is injective, where $N$ is the span
of the balancing relations $y\,\iota(x)\otimes u - y\otimes\rho(x)u$. By Proposition~3.11 this is
the map $U\to U(\ltwog)\otimes_{U(\ltwop)}U$, and over $K = \C$ it is field (B) of \lean{PBWPrinciple}.
More generally: if there are a representation $\sigma$ of $\ltwog$ on some $M$ and an injective linear
map $j : U\to M$ with $\sigma(x)j(u) = j(\rho(x)u)$ for $x\in\ltwop$, then the bottom map is injective.
\end{btheorem}
{\raggedright\noindent\textbf{Lean.} \lean{VOA.}\allowbreak\lean{PBWTheorem.}\allowbreak\lean{balancingSet},\allowbreak\ \lean{VOA.}\allowbreak\lean{PBWTheorem.}\allowbreak\lean{balancing},\allowbreak\ \lean{VOA.}\allowbreak\lean{PBWTheorem.}\allowbreak\lean{bottom},\allowbreak\ \lean{VOA.}\allowbreak\lean{PBWTheorem.}\allowbreak\lean{bottom_injective_of_embedding},\allowbreak\ \lean{VOA.}\allowbreak\lean{PBWTheorem.}\allowbreak\lean{bottom_injective}\par}
{\raggedright\noindent File \leanpath{VOA/\allowbreak Foundations/\allowbreak PBWTheorem/\allowbreak Induced.lean}.\par}
\begin{leancode}
def balancingSet (p : LieSubalgebra K g) (rho : p →ₗ⁅K⁆ Module.End K U) :
    Set (UniversalEnvelopingAlgebra K g ⊗[K] U) :=
  {t | ∃ (y : UniversalEnvelopingAlgebra K g) (x : p) (u : U),
    t = (y * UniversalEnvelopingAlgebra.ι K (x : g)) ⊗ₜ[K] u - y ⊗ₜ[K] (rho x u)}

noncomputable def balancing (p : LieSubalgebra K g) (rho : p →ₗ⁅K⁆ Module.End K U) :
    Submodule K (UniversalEnvelopingAlgebra K g ⊗[K] U) :=
  Submodule.span K (balancingSet p rho)

noncomputable def bottom (p : LieSubalgebra K g) (rho : p →ₗ⁅K⁆ Module.End K U) :
    U →ₗ[K] (UniversalEnvelopingAlgebra K g ⊗[K] U) ⧸ balancing p rho :=
  (balancing p rho).mkQ.comp (TensorProduct.mk K (UniversalEnvelopingAlgebra K g) U 1)

theorem bottom_injective_of_embedding (p : LieSubalgebra K g) (rho : p →ₗ⁅K⁆ Module.End K U)
    {M : Type*} [AddCommGroup M] [Module K M] (σ : g →ₗ⁅K⁆ Module.End K M) (j : U →ₗ[K] M)
    (hj : Function.Injective j) (hσ : ∀ (x : p) (u : U), σ (x : g) (j u) = j (rho x u)) :
    Function.Injective (bottom p rho) := by

theorem bottom_injective (p : LieSubalgebra K g) (rho : p →ₗ⁅K⁆ Module.End K U) :
    Function.Injective (bottom p rho) := by
\end{leancode}
\noindent\textbf{Hypotheses / status.} Unconditional\par
\medskip
\begin{btheorem}{3.18}[Nonvanishing of the affine vacuum, from bottom injectivity]
For every complex Lie algebra $\ltwog$ with symmetric invariant form $\kappa$ and every $k\in\C$, the vacuum
vector of $V^k(\ltwog)$ is nonzero. The proof uses the subalgebra
$\hat{\ltwop} = \ltwog[t]\oplus\C C\oplus\C\ell_0$ of $\ltwoghat$ and its character $\chi_k(f + aC + p\ell_0) = ka$.
\end{btheorem}
{\raggedright\noindent\textbf{Lean.} \lean{VOA.}\allowbreak\lean{AffineLie.}\allowbreak\lean{nonnegSub},\allowbreak\ \lean{VOA.}\allowbreak\lean{AffineLie.}\allowbreak\lean{vacuumChar},\allowbreak\ \lean{VOA.}\allowbreak\lean{AffineLie.}\allowbreak\lean{tmul_one_mem_balancing_of_mem_vacuumIdeal},\allowbreak\ \lean{VOA.}\allowbreak\lean{AffineLie.}\allowbreak\lean{vacuumVector_ne_zero_of_induction}\par}
{\raggedright\noindent File \leanpath{VOA/\allowbreak Foundations/\allowbreak AffineLie.lean}.\par}
\begin{leancode}
noncomputable def nonnegSub : LieSubalgebra ℂ (AffineLie 𝔤 κ hκ) where
  carrier := {X | X.carrier.1 ∈ supported 𝔤 ℂ (Set.Ici (0 : ℤ))}
  ...

noncomputable def vacuumChar (k : ℂ) : nonnegSub κ hκ →ₗ⁅ℂ⁆ Module.End ℂ ℂ where
  toFun X := (k * (X : AffineLie 𝔤 κ hκ).carrier.2.1) • LinearMap.id
  ...

theorem tmul_one_mem_balancing_of_mem_vacuumIdeal (k : ℂ) {a : UEnv 𝔤 κ hκ}
    (ha : a ∈ vacuumIdeal κ hκ k) :
    a ⊗ₜ[ℂ] (1 : ℂ) ∈ PBWTheorem.balancing (nonnegSub κ hκ) (vacuumChar κ hκ k) := by

theorem vacuumVector_ne_zero_of_induction (k : ℂ) : vacuumVector κ hκ k ≠ 0 := by
\end{leancode}
\noindent\textbf{Hypotheses / status.} Unconditional\par
\medskip
\begin{bdefinition}{3.20}[Ordered creation word; PBW basis hypothesis]
Let $(b_i)_{i\in I}$ be a basis of $\ltwog$ indexed by a totally ordered set $I$, and order $I\times\N$
lexicographically. For a finitely supported exponent function $\mu : I\times\N\to\N$, the
\emph{ordered creation word} lists the pairs $(i,n)$ in the support in increasing order, each repeated
$\mu(i,n)$ times; the corresponding \emph{ordered creation monomial} is
\begin{equation*}
 v_\mu = \prod_{(i,n)\ \text{increasing}}(b_i\otimes t^{-n-1})^{\mu(i,n)}\ltwovac\ \in V^k(\ltwog),
\end{equation*}
with the smallest pair leftmost. \emph{$\mathrm{PBWBasis}(\kappa,b,k)$} is the statement that
$(v_\mu)_\mu$ is a basis of $V^k(\ltwog)$.
\end{bdefinition}
{\raggedright\noindent\textbf{Lean.} \lean{VOA.}\allowbreak\lean{pbwWord},\allowbreak\ \lean{VOA.}\allowbreak\lean{PBWBasis},\allowbreak\ \lean{VOA.}\allowbreak\lean{PBWBasis.}\allowbreak\lean{linearIndependent},\allowbreak\ \lean{VOA.}\allowbreak\lean{PBWBasis.}\allowbreak\lean{vacuumVector_ne_zero},\allowbreak\ \lean{VOA.}\allowbreak\lean{pbwBasis_iff_of_subsingleton}\par}
{\raggedright\noindent File \leanpath{VOA/\allowbreak Foundations/\allowbreak PBWBasis.lean}.\par}
\begin{leancode}
noncomputable def pbwWord {𝔤 : Type u} [LieRing 𝔤] [LieAlgebra ℂ 𝔤] {ι : Type*} [LinearOrder ι]
    (b : Module.Basis ι ℂ 𝔤) (μ : (ι ×ₗ ℕ) →₀ ℕ) : List (𝔤 × ℕ) :=
  (μ.support.sort (· ≤ ·)).flatMap fun p => List.replicate (μ p) (b (ofLex p).1, (ofLex p).2)

def PBWBasis {𝔤 : Type u} [LieRing 𝔤] [LieAlgebra ℂ 𝔤]
    (κ : LinearMap.BilinForm ℂ 𝔤) (hκ : κ.IsInvariant)
    {ι : Type*} [LinearOrder ι] (b : Module.Basis ι ℂ 𝔤) (k : ℂ) : Prop :=
  ∃ B : Module.Basis ((ι ×ₗ ℕ) →₀ ℕ) ℂ (AffineLie.vacuumSpace κ hκ k),
    ∀ μ, B μ = AffineLie.negWord κ hκ k (pbwWord b μ)
\end{leancode}
\begin{leancode}
theorem linearIndependent (h : PBWBasis κ hκ b k) :
    LinearIndependent ℂ fun μ : (ι ×ₗ ℕ) →₀ ℕ => AffineLie.negWord κ hκ k (pbwWord b μ) := by

theorem vacuumVector_ne_zero (h : PBWBasis κ hκ b k) :
    AffineLie.vacuumVector κ hκ k ≠ 0 := by

theorem pbwBasis_iff_of_subsingleton :
    PBWBasis κ hκ (Module.Basis.empty 𝔤 : Module.Basis (Fin 0) ℂ 𝔤) k
      ↔ AffineLie.vacuumVector κ hκ k ≠ 0 :=
\end{leancode}
\medskip
\begin{btheorem}{3.21}[PBW basis of the affine vacuum module]
For every complex Lie algebra $\ltwog$ (of any dimension), every symmetric invariant bilinear form $\kappa$,
every basis $(b_i)_{i\in I}$ of $\ltwog$ indexed by a totally ordered set, and every $k\in\C$ (including the
critical level), $\mathrm{PBWBasis}(\kappa,b,k)$ holds: the ordered creation monomials $v_\mu$ of
Definition~3.20 form a basis of $V^k(\ltwog)$. More precisely, the linear map
$\Phi : V^k(\ltwog)\to S(\ltwon_-)$, $[a]\mapsto a\cdot1$, given by the relative PBW representation of
$\ltwoghat$ on $S(\ltwon_-)\otimes\C_{\chi_k}\cong S(\ltwon_-)$, is a linear isomorphism.
\end{btheorem}
{\raggedright\noindent\textbf{Lean.} \lean{VOA.}\allowbreak\lean{VacuumBasis.}\allowbreak\lean{yMode},\allowbreak\ \lean{VOA.}\allowbreak\lean{VacuumBasis.}\allowbreak\lean{negComp},\allowbreak\ \lean{VOA.}\allowbreak\lean{VacuumBasis.}\allowbreak\lean{negComp_isCompl},\allowbreak\ \lean{VOA.}\allowbreak\lean{VacuumBasis.}\allowbreak\lean{Φ},\allowbreak\ \lean{VOA.}\allowbreak\lean{VacuumBasis.}\allowbreak\lean{Φ_rho},\allowbreak\ \lean{VOA.}\allowbreak\lean{VacuumBasis.}\allowbreak\lean{ΦEquiv},\allowbreak\ \lean{VOA.}\allowbreak\lean{pbwBasis_holds}\par}
{\raggedright\noindent File \leanpath{VOA/\allowbreak Foundations/\allowbreak PBWTheorem/\allowbreak VacuumBasis.lean}.\par}
\begin{leancode}
noncomputable def yMode (a : ι ×ₗ ℕ) : AffineLie 𝔤 κ hκ :=
  ofLoop κ hκ (b (ofLex a).1) (-((ofLex a).2 : ℤ) - 1)

noncomputable def negComp : Submodule ℂ (AffineLie 𝔤 κ hκ) :=
  Submodule.span ℂ (Set.range (yMode κ hκ b))

theorem negComp_isCompl : IsCompl (nonnegSub κ hκ).toSubmodule (negComp κ hκ b) := by

noncomputable def Φ : vacuumSpace κ hκ k →ₗ[ℂ] (Multiset (ι ×ₗ ℕ) →₀ ℂ) :=

theorem Φ_rho (z : AffineLie 𝔤 κ hκ) (v : vacuumSpace κ hκ k) :
    Φ κ hκ b k (vacuumRho κ hκ k z v) = ACT z (Φ κ hκ b k v) := by

noncomputable def ΦEquiv : vacuumSpace κ hκ k ≃ₗ[ℂ] (Multiset (ι ×ₗ ℕ) →₀ ℂ) :=
  LinearEquiv.ofLinear (Φ κ hκ b k) (Ψ κ hκ b k)
    (Finsupp.lhom_ext fun m u => Φ_Ψ_single κ hκ b k m u)
    (LinearMap.ext_on_range (negSpan_eq_top κ hκ k) fun L => Ψ_Φ_negWord κ hκ b k L)

theorem pbwBasis_holds {𝔤 : Type u} [LieRing 𝔤] [LieAlgebra ℂ 𝔤]
    (κ : LinearMap.BilinForm ℂ 𝔤) (hκ : κ.IsInvariant)
    {ι : Type*} [LinearOrder ι] (b : Module.Basis ι ℂ 𝔤) (k : ℂ) : PBWBasis κ hκ b k := by
\end{leancode}
\noindent\textbf{Hypotheses / status.} Unconditional\par
\medskip
\endgroup
\setcounter{chapter}{3}
\chapter[Vertex algebra axioms, Borcherds identity, reconstruction]{Vertex algebra axioms, the Borcherds identity, and reconstruction}
\begingroup
\providecommand{\lthreevac}{\mathbf{1}}
\providecommand{\lthreeV}{\mathcal{V}}
\providecommand{\lthreeY}{Y}
\providecommand{\lthreeBL}{B_{\mathrm L}}
\providecommand{\lthreeBR}{B_{\mathrm R}}
\providecommand{\lthreeD}{\partial}
\providecommand{\lthreeno}[1]{{:}#1{:}}
\noindent{\small 6 of the chapter's 51 numbered statements: its headline results and every definition they use.}\par\medskip
\begin{btheorem}{4.14}[Goddard uniqueness, three forms]
Let $V$ be a vertex algebra in the library's sense, $a\in V$ and $A(z)\in\End(V)[[z^{\pm1}]]$ with
$A(z)\lthreevac = Y(a,z)\lthreevac$ coefficientwise.
\begin{enumerate}[label=(\roman*)]
\item If $A$ is local with every $Y(b,z)$, then $A = Y(a,\cdot)$.
\item It suffices that for every $b$ some $(z_1-z_2)^N$ kills $[A(z_1),Y(b,z_2)]\lthreevac$ (the
  commutator applied to the vacuum only).
\end{enumerate}
(iii) (\emph{Abstract test fields.}) Let $V$ be any vector space with a vector $\lthreevac$, $A$ a
distribution with $A(z)\lthreevac = 0$, and $(G_b)_{b\in V}$ distributions with $G_b(w)\lthreevac\in V[[w]]$
of constant term $b$. If $A$ is local with every $G_b$, then $A = 0$.
\end{btheorem}
{\raggedright\noindent\textbf{Lean.} \lean{VOA.}\allowbreak\lean{PreVertexAlgebra.}\allowbreak\lean{goddard_uniqueness},\allowbreak\ \lean{VOA.}\allowbreak\lean{goddard_zero},\allowbreak\ \lean{VOA.}\allowbreak\lean{PreVertexAlgebra.}\allowbreak\lean{goddard_uniqueness_vac}\par}
\begin{leancode}
theorem goddard_uniqueness (V : PreVertexAlgebra 𝒱) (A : FormalDist (End ℂ 𝒱)) (a : 𝒱)
    (hAloc : ∀ b : 𝒱, Local A (V.Y b))
    (hAvac : ∀ m : ℤ, A m V.vacuum = (V.Y a) m V.vacuum) :
    A = V.Y a

theorem goddard_zero (vac : 𝒱) (A : FormalDist (End ℂ 𝒱)) (G : 𝒱 → FormalDist (End ℂ 𝒱))
    (hloc : ∀ b, Local A (G b))
    (hcre0 : ∀ b, (G b) 0 vac = b)
    (hcreneg : ∀ b, ∀ q : ℤ, q < 0 → (G b) q vac = 0)
    (hAvac : ∀ p : ℤ, A p vac = 0)
    (m : ℤ) (b : 𝒱) : A m b = 0
\end{leancode}
{\raggedright\noindent File \leanpath{VOA/\allowbreak Core/\allowbreak Goddard.lean} (\lean{VOA.PreVertexAlgebra.goddard_uniqueness},
\lean{VOA.goddard_zero}).\par}
\begin{leancode}
theorem goddard_uniqueness_vac (V : PreVertexAlgebra 𝒱) (A : FormalDist (End ℂ 𝒱)) (a : 𝒱)
    (hAloc : ∀ b : 𝒱, ∃ N : ℕ, mulBySubPow N (applyVec V.vacuum (bracketDist A (V.Y b))) = 0)
    (hAvac : ∀ m : ℤ, A m V.vacuum = (V.Y a) m V.vacuum) :
    A = V.Y a
\end{leancode}
{\raggedright\noindent File \leanpath{VOA/\allowbreak Core/\allowbreak WeakAssoc.lean} (\lean{VOA.PreVertexAlgebra.goddard_uniqueness_vac}).\par}
\noindent\textbf{Hypotheses / status.} VA level / abstract; unconditional\par
\medskip
\begin{btheorem}{4.16}[Dong's lemma, all products]
Let $a,b,c$ be pairwise local distributions on $V$. Then for every $k\ge0$ the field product
$a_{(k)}b := \Res_{z_1}(z_1-z_2)^k[a(z_1),b(z_2)]$ is local with $c$. If moreover $a,b$ are fields,
then ${:}ab{:}$ is local with $c$, and hence so are all negative products
$a_{(-m-1)}b = \frac1{m!}{:}(\partial^ma)b{:}$.
\end{btheorem}
{\raggedright\noindent\textbf{Lean.} \lean{VOA.}\allowbreak\lean{jacobi₃},\allowbreak\ \lean{VOA.}\allowbreak\lean{dong_core},\allowbreak\ \lean{VOA.}\allowbreak\lean{dong_nonneg},\allowbreak\ \lean{VOA.}\allowbreak\lean{dong_neg}\par}
\begin{leancode}
theorem jacobi₃ (a b c : FormalDist (End ℂ V)) :
    outerC₁₂ (bracketDist a b) c
      = outerA a (bracketDist b c) - outerB b (bracketDist a c)

theorem dong_core (a b c : FormalDist (End ℂ V)) {r : ℕ}
    (hab : mulBySubPow r (bracketDist a b) = 0)
    (hac : mulBySubPow r (bracketDist a c) = 0)
    (hbc : mulBySubPow r (bracketDist b c) = 0)
    (k M : ℕ) (hM : 2 * r ≤ M) :
    mulBySub₂₃Pow (M + r) (mulBySub₁₂Pow k (outerC₁₂ (bracketDist a b) c)) = 0

theorem dong_nonneg {a b c : FormalDist (End ℂ V)}
    (hab : Local a b) (hac : Local a c) (hbc : Local b c) (k : ℕ) :
    Local (cCoeff k (bracketDist a b)) c

theorem dong_neg {a b c : FormalDist (End ℂ V)} (ha : IsField a) (hb : IsField b)
    (hab : Local a b) (hac : Local a c) (hbc : Local b c) :
    Local (nprod ha hb) c
\end{leancode}
{\raggedright\noindent File \leanpath{VOA/\allowbreak Core/\allowbreak Dong.lean}.\par}
\noindent\textbf{Hypotheses / status.} Pairwise locality (+ field); unconditional\par
\medskip
\begin{btheorem}{4.20}[The $n$-th product theorem, $n\ge0$]
In a VOA, for all $u,v$ and $k\in\N$:
\begin{equation*}
Y(u_{(k)}v,z) = \Res_{z_1}(z_1-z)^k[Y(u,z_1),Y(v,z)] .
\end{equation*}
\end{btheorem}
{\raggedright\noindent\textbf{Lean.} \lean{VOA.}\allowbreak\lean{VertexOperatorAlgebra.}\allowbreak\lean{cCoeff_vacuum_eq},\allowbreak\ \lean{VOA.}\allowbreak\lean{VertexOperatorAlgebra.}\allowbreak\lean{nth_product},\allowbreak\ \lean{VOA.}\allowbreak\lean{VertexOperatorAlgebra.}\allowbreak\lean{assoc_nonneg}\par}
\begin{leancode}
theorem cCoeff_vacuum_eq (W : VertexOperatorAlgebra 𝒱) (u v : 𝒱) (k : ℕ) (m : ℤ) :
    cCoeff k (bracketDist (W.Y u) (W.Y v)) m W.vacuum = (W.Y (W.smul u k v)) m W.vacuum

theorem nth_product (W : VertexOperatorAlgebra 𝒱) (u v : 𝒱) (k : ℕ) :
    cCoeff k (bracketDist (W.Y u) (W.Y v)) = W.Y (W.smul u k v)

theorem assoc_nonneg (W : VertexOperatorAlgebra 𝒱) (u v : 𝒱) (k : ℕ) :
    W.Y (W.smul u k v) = cCoeff k (bracketDist (W.Y u) (W.Y v))
\end{leancode}
{\raggedright\noindent File \leanpath{VOA/\allowbreak Core/\allowbreak Reconstruction.lean}.\par}
\noindent\textbf{Hypotheses / status.} VOA; unconditional\par
\medskip
\begin{btheorem}{4.21}[Commutator formula]
In a VOA, for all $u,v$ and $a,b\in\Z$,
$[u_{(a)},v_{(b)}] = \sum_{j\ge0}\binom aj(u_{(j)}v)_{(a+b-j)}$ (a finite sum).
\end{btheorem}
{\raggedright\noindent\textbf{Lean.} \lean{VOA.}\allowbreak\lean{VertexOperatorAlgebra.}\allowbreak\lean{mode_commutator},\allowbreak\ \lean{VOA.}\allowbreak\lean{VertexOperatorAlgebra.}\allowbreak\lean{commBound}\par}
\begin{leancode}
theorem mode_commutator (W : VertexOperatorAlgebra 𝒱) (u v : 𝒱) (a b : ℤ) :
    ∃ N : ℕ, ⁅mode (W.Y u) a, mode (W.Y v) b⁆
      = ∑ j : Fin N, gBinom a (j : ℕ) • mode (W.Y (W.smul u (j : ℕ) v)) (a + b - (j : ℕ))

theorem commBound (W : VertexOperatorAlgebra 𝒱) (u v : 𝒱) (m n : ℤ) :
    ∃ N : ℕ, (∀ k : ℕ, N ≤ k → W.smul u k v = 0)
      ∧ ⁅mode (W.Y u) m, mode (W.Y v) n⁆
        = ∑ j : Fin N, gBinom m (j : ℕ) • mode (W.Y (W.smul u (j : ℕ) v)) (m + n - (j : ℕ))
\end{leancode}
{\raggedright\noindent File \leanpath{VOA/\allowbreak Core/\allowbreak Reconstruction.lean} (\lean{mode_commutator}) and
\leanpath{VOA/\allowbreak Core/\allowbreak Borcherds.lean} (\lean{commBound}). The theorem
\lean{jacobi_commutator} of \leanpath{VOA/\allowbreak Core/\allowbreak BorcherdsIdentity.lean} is
\lean{mode_commutator} under another name.\par}
\noindent\textbf{Hypotheses / status.} VOA; $\exists N$ made exhaustive by \lean{commBound}\par
\medskip
\begin{btheorem}{4.26}[The Borcherds identity]
In a VOA, for all $u,v,x\in V$ and all $\ell,m,n\in\Z$, $\lthreeBL(\ell,m,n) = \lthreeBR(\ell,m,n)$, i.e.
\begin{equation}
\sum_{i\ge0}(-1)^i\binom\ell i\Big(u_{(\ell+m-i)}v_{(n+i)}x - (-1)^\ell v_{(\ell+n-i)}u_{(m+i)}x\Big)
 = \sum_{i\ge0}\binom mi(u_{(\ell+i)}v)_{(m+n-i)}x .
\end{equation}
\end{btheorem}
{\raggedright\noindent\textbf{Lean.} \lean{VOA.}\allowbreak\lean{VertexOperatorAlgebra.}\allowbreak\lean{borcherds_base},\allowbreak\ \lean{VOA.}\allowbreak\lean{VertexOperatorAlgebra.}\allowbreak\lean{borcherds_boundary},\allowbreak\ \lean{VOA.}\allowbreak\lean{VertexOperatorAlgebra.}\allowbreak\lean{borcherds},\allowbreak\ \lean{VOA.}\allowbreak\lean{VertexOperatorAlgebra.}\allowbreak\lean{jacobi}\par}
\begin{leancode}
theorem borcherds_base (W : VertexOperatorAlgebra 𝒱) (u v x : 𝒱) (ℓ n : ℤ) :
    BL W u v x ℓ 0 n = BR W u v x ℓ 0 n

theorem borcherds_boundary (W : VertexOperatorAlgebra 𝒱) (u v x : 𝒱) (m n : ℤ) :
    BL W u v x 0 m n = BR W u v x 0 m n

theorem borcherds (W : VertexOperatorAlgebra 𝒱) (u v x : 𝒱) (ℓ m n : ℤ) :
    BL W u v x ℓ m n = BR W u v x ℓ m n

theorem jacobi (W : VertexOperatorAlgebra 𝒱) (u v x : 𝒱) (ℓ m n : ℤ) :
    (∑ᶠ i : ℤ, (negOnePow i * cbz ℓ i) •
        (W.smul u (ℓ + m - i) (W.smul v (n + i) x)
          - negOnePow ℓ • W.smul v (ℓ + n - i) (W.smul u (m + i) x)))
      = ∑ᶠ i : ℤ, cbz m i • W.smul (W.smul u (ℓ + i) v) (m + n - i) x
\end{leancode}
{\raggedright\noindent File \leanpath{VOA/\allowbreak Core/\allowbreak Borcherds.lean}.\par}
\noindent\textbf{Hypotheses / status.} VOA; unconditional, axiom-clean\par
\medskip
\begin{btheorem}{4.41}[Operator-first reconstruction]
Let $V$ be a vector space, $\lthreevac\in V$, $T\in\End V$ with $T\lthreevac = 0$, and $(s_i,f_i)_{i\in I}$
states and distributions such that each $f_i$ is a field, the $f_i$ are pairwise local and
$T$-covariant, $f_i(z)\lthreevac\in V[[z]]$ with constant term $s_i$, and the $s_i$ span $V$. Then there is
a linear $Y:V\to\End(V)[[z^{\pm1}]]$ with $Y(s_i,\cdot) = f_i$, all $Y(u,\cdot)$ fields, pairwise
local, $T$-covariant, with $Y(u,z)\lthreevac\in u+zV[[z]]$.
\end{btheorem}
{\raggedright\noindent\textbf{Lean.} \lean{VOA.}\allowbreak\lean{Reconstruction.}\allowbreak\lean{reconstruct}\par}
\begin{leancode}
theorem reconstruct {ι : Type*} (vac : V) (st : ι → V) (fld : ι → FormalDist (End ℂ V))
    (T : End ℂ V)
    (hspan : Submodule.span ℂ (Set.range st) = ⊤)
    (hfield : ∀ i, IsField (fld i))
    (hlocal : ∀ i j, Local (fld i) (fld j))
    (hcre0 : ∀ i, (fld i) 0 vac = st i)
    (hcreneg : ∀ i, ∀ q : ℤ, q < 0 → (fld i) q vac = 0)
    (hTvac : T vac = 0)
    (hTcov : ∀ i, DCovariant T (fld i)) :
    ∃ YA : V →ₗ[ℂ] FormalDist (End ℂ V),
      (∀ i, YA (st i) = fld i) ∧
      (∀ u : V, IsField (YA u)) ∧
      (∀ u v : V, Local (YA u) (YA v)) ∧
      (∀ u : V, (YA u) 0 vac = u) ∧
      (∀ u : V, ∀ q : ℤ, q < 0 → (YA u) q vac = 0) ∧
      (∀ u : V, DCovariant T (YA u))
\end{leancode}
{\raggedright\noindent File \leanpath{VOA/\allowbreak Core/\allowbreak ReconstructionExistence.lean}
(\lean{VOA.Reconstruction.reconstruct}).\par}
\noindent\textbf{Hypotheses / status.} All hypotheses explicit; no VA2/VA1 in conclusion\par
\medskip
\endgroup
\setcounter{chapter}{4}
\chapter{Structure theory and invariant forms}
\begingroup
\providecommand{\lfourvac}{\mathbf{1}}
\providecommand{\lfourV}{V}
\providecommand{\lfourm}{\mathfrak{m}}
\providecommand{\lfourrad}{\operatorname{rad}}
\providecommand{\lfourQP}{\mathcal{QP}}
\providecommand{\lfourvc}{\varepsilon}
\providecommand{\lfourH}{\mathrm{H}}
\noindent{\small 7 of the chapter's 55 numbered statements: its headline results and every definition they use.}\par\medskip
\begin{btheorem}{5.2}[Zero modes are derivations; Gannon (5.2.4)]
For all $u,v,w\in\lfourV$ and $n\in\Z$:
$u_{(0)}(v_{(n)}w) - v_{(n)}(u_{(0)}w) = (u_{(0)}v)_{(n)}w$.
\end{btheorem}
{\raggedright\noindent\textbf{Lean.} \lean{VOA.}\allowbreak\lean{VertexOperatorAlgebra.}\allowbreak\lean{zeroMode_isDerivation},\allowbreak\ \lean{VOA.}\allowbreak\lean{VertexOperatorAlgebra.}\allowbreak\lean{zeroMode_leibniz}\par}
{\raggedright\noindent File \leanpath{VOA/\allowbreak Core/\allowbreak LieAlgebraV1.lean}.\par}
\begin{leancode}
theorem zeroMode_isDerivation (W : VertexOperatorAlgebra 𝒱) (u v w : 𝒱) (n : ℤ) :
    W.smul u 0 (W.smul v n w) - W.smul v n (W.smul u 0 w) = W.smul (W.smul u 0 v) n w := by
\end{leancode}
{\raggedright\noindent File \leanpath{VOA/\allowbreak Core/\allowbreak Automorphism.lean} restates it as a Leibniz rule:\par}
\begin{leancode}
theorem zeroMode_leibniz (W : VertexOperatorAlgebra 𝒱) (u a b : 𝒱) (n : ℤ) :
    W.smul u 0 (W.smul a n b) = W.smul (W.smul u 0 a) n b + W.smul a n (W.smul u 0 b) := by
\end{leancode}
\noindent\textbf{Hypotheses / status.} None; unconditional\par
\medskip
\begin{btheorem}{5.17}[Ideals are two-sided]
If $I$ is an ideal, then $u_{(n)}v\in I$ for all $u\in I$, $v\in\lfourV$, $n\in\Z$.
\end{btheorem}
{\raggedright\noindent\textbf{Lean.} \lean{VOA.}\allowbreak\lean{QuotVOA.}\allowbreak\lean{AbsorbsLeft},\allowbreak\ \lean{VOA.}\allowbreak\lean{QuotVOA.}\allowbreak\lean{absorbsLeft_of_isIdeal}\par}
{\raggedright\noindent File \leanpath{VOA/\allowbreak Core/\allowbreak QuotientVOA.lean}.\par}
\begin{leancode}
def AbsorbsLeft (W : VertexOperatorAlgebra 𝒱) (I : Submodule ℂ 𝒱) : Prop :=
  ∀ (u v : 𝒱) (n : ℤ), u ∈ I → W.smul u n v ∈ I

theorem absorbsLeft_of_isIdeal (hI : W.IsIdeal I) : AbsorbsLeft W I := by
\end{leancode}
\noindent\textbf{Hypotheses / status.} None\par
\medskip
\begin{btheorem}{5.18}[Ideals are graded]
Every ideal $I$ satisfies $I = \bigoplus_n(I\cap\lfourV_n)$.
\end{btheorem}
{\raggedright\noindent\textbf{Lean.} \lean{VOA.}\allowbreak\lean{QuotVOA.}\allowbreak\lean{IsGradedIdeal},\allowbreak\ \lean{VOA.}\allowbreak\lean{QuotVOA.}\allowbreak\lean{isGradedIdeal_of_isIdeal}\par}
{\raggedright\noindent File \leanpath{VOA/\allowbreak Core/\allowbreak QuotientVOA.lean}.\par}
\begin{leancode}
def IsGradedIdeal (W : VertexOperatorAlgebra 𝒱) (I : Submodule ℂ 𝒱) : Prop :=
  I = ⨆ n : ℤ, I ⊓ W.grading n

theorem isGradedIdeal_of_isIdeal (hI : W.IsIdeal I) : IsGradedIdeal W I := by
\end{leancode}
\noindent\textbf{Hypotheses / status.} None\par
\medskip
\begin{bdefinition}{5.26}[Generated subalgebra]
For $S\subseteq\lfourV$, $\langle S\rangle$ is the intersection of all vertex operator subalgebras
(in the sense just recalled) containing $S$.
\end{bdefinition}
{\raggedright\noindent\textbf{Lean.} \lean{VOA.}\allowbreak\lean{VOASubalgebra.}\allowbreak\lean{generated},\allowbreak\ \lean{VOA.}\allowbreak\lean{VOASubalgebra.}\allowbreak\lean{generated_le},\allowbreak\ \lean{VOA.}\allowbreak\lean{VOASubalgebra.}\allowbreak\lean{mem_generated},\allowbreak\ \lean{VOA.}\allowbreak\lean{VOASubalgebra.}\allowbreak\lean{generatedVOA}\par}
{\raggedright\noindent File \leanpath{VOA/\allowbreak Core/\allowbreak VOASubalgebraGenerated.lean}.\par}
\begin{leancode}
def generated (S : Set 𝒱) : VOASubalgebra W :=
  sInf {A : VOASubalgebra W | S ⊆ A.carrier}

theorem generated_le {S : Set 𝒱} (A : VOASubalgebra W) (hA : S ⊆ A.carrier) :
    (generated W S).carrier ≤ A.carrier :=

theorem mem_generated {S : Set 𝒱} {motive : 𝒱 → Prop}
    (base : ∀ x ∈ S, motive x)
    (zero : motive 0)
    (add : ∀ x y : 𝒱, motive x → motive y → motive (x + y))
    (smul : ∀ (c : ℂ) (x : 𝒱), motive x → motive (c • x))
    (vac : motive W.vacuum)
    (conf : motive W.conformalVector)
    (prod : ∀ u : 𝒱, motive u → ∀ (n : ℤ) (v : 𝒱), motive v → motive (W.smul u n v))
    (proj : ∀ (n : ℤ) (v : 𝒱), motive v → motive (piHom W n v))
    {x : 𝒱} (hx : x ∈ (generated W S).carrier) : motive x := by

noncomputable def generatedVOA (S : Set 𝒱) :
    VertexOperatorAlgebra (generated W S).carrier :=
  (generated W S).toVOA
\end{leancode}
\noindent\textbf{Hypotheses / status.} \textbf{always contains $\omega$}\par
\medskip
\begin{btheorem}{5.50}[Orthogonality, radical, uniqueness, nondegeneracy, evaluation on $\lfourV_1$]
Let $B$ be an invariant form. (i) $(\lfourV_m|\lfourV_n) = 0$ for $m\ne n$. Now let $B$ be fully invariant.
(ii) Its left radical is an ideal. (iii) If $B_1,B_2$ are fully invariant and normalised, $B_1 = B_2$.
(iv) If $\lfourV$ is simple and $B$ is fully invariant and normalised, its radical is $0$. (v) If $B$ is
fully invariant and normalised, $u_{(1)}v = -(u|v)\lfourvac$ for $u,v\in\lfourV_1$.
\end{btheorem}
{\raggedright\noindent\textbf{Lean.} \lean{VOA.}\allowbreak\lean{VertexOperatorAlgebra.}\allowbreak\lean{InvariantForm.}\allowbreak\lean{orthogonal_of_ne},\allowbreak\ \lean{VOA.}\allowbreak\lean{VertexOperatorAlgebra.}\allowbreak\lean{InvariantForm.}\allowbreak\lean{radical_isIdeal},\allowbreak\ \lean{VOA.}\allowbreak\lean{VertexOperatorAlgebra.}\allowbreak\lean{InvariantForm.}\allowbreak\lean{invariantForm_unique},\allowbreak\ \lean{VOA.}\allowbreak\lean{VertexOperatorAlgebra.}\allowbreak\lean{InvariantForm.}\allowbreak\lean{invariantForm_nondegenerate},\allowbreak\ \lean{VOA.}\allowbreak\lean{VertexOperatorAlgebra.}\allowbreak\lean{InvariantForm.}\allowbreak\lean{eval_oneMode_grading1}\par}
{\raggedright\noindent File \leanpath{VOA/\allowbreak Core/\allowbreak BilinearForm.lean}.\par}
\begin{leancode}
theorem orthogonal_of_ne (B : InvariantForm W) {m n : ℤ} (hmn : m ≠ n) {a b : 𝒱}
    (ha : a ∈ W.grading m) (hb : b ∈ W.grading n) : B.form a b = 0 := by

theorem radical_isIdeal {B : InvariantForm W} (hB : B.IsFullyInvariant) :
    W.IsIdeal (leftRadical B) := by

theorem invariantForm_unique {B₁ B₂ : InvariantForm W}
    (h₁ : B₁.IsFullyInvariant) (h₂ : B₂.IsFullyInvariant)
    (hn₁ : B₁.Normalized) (hn₂ : B₂.Normalized) : B₁.form = B₂.form := by

theorem invariantForm_nondegenerate {B : InvariantForm W} (hsimple : W.IsSimple)
    (hB : B.IsFullyInvariant) (hn : B.Normalized) : leftRadical B = ⊥ := by

theorem eval_oneMode_grading1 {B : InvariantForm W} (hB : B.IsFullyInvariant)
    (hn : B.Normalized) {u v : 𝒱} (hu : u ∈ W.grading 1) (hv : v ∈ W.grading 1) :
    W.smul u 1 v = -B.form u v • W.vacuum := by
\end{leancode}
\noindent\textbf{Hypotheses / status.} None\par
\medskip
\begin{btheorem}{5.56}[FHL, Li: existence and uniqueness of the invariant form]
Let $\lfourV$ be a VOA (of CFT type) with $\lfourvac\ne0$ and $L_1\lfourV_1 = 0$. There is a fully invariant form
with $(\lfourvac|\lfourvac) = 1$; every fully invariant normalised form equals it; it is nondegenerate if
$\lfourV$ is simple; and $u_{(1)}v = -(u|v)\lfourvac$ on $\lfourV_1$.
\end{btheorem}
{\raggedright\noindent\textbf{Lean.} \lean{VOA.}\allowbreak\lean{canonicalInvariantForm},\allowbreak\ \lean{VOA.}\allowbreak\lean{canonicalInvariantForm_isFullyInvariant},\allowbreak\ \lean{VOA.}\allowbreak\lean{canonicalInvariantForm_normalized},\allowbreak\ \lean{VOA.}\allowbreak\lean{invariantForm_unique_FHL},\allowbreak\ \lean{VOA.}\allowbreak\lean{invariantForm_nondegenerate_FHL},\allowbreak\ \lean{VOA.}\allowbreak\lean{eval_oneMode_grading1_FHL}\par}
{\raggedright\noindent File \leanpath{VOA/\allowbreak Core/\allowbreak InvariantFormFHL2.lean}.\par}
\begin{leancode}
noncomputable def canonicalInvariantForm (hne : W.vacuum ≠ 0)
    (hqp : ∀ v, v ∈ W.grading 1 → W.L 1 v = 0) : InvariantForm W :=
  canFormFHL hne (canDiagonal_unconditional hne hqp)

theorem canonicalInvariantForm_isFullyInvariant (hne : W.vacuum ≠ 0)
    (hqp : ∀ v, v ∈ W.grading 1 → W.L 1 v = 0) :
    (canonicalInvariantForm hne hqp).IsFullyInvariant :=

theorem canonicalInvariantForm_normalized (hne : W.vacuum ≠ 0)
    (hqp : ∀ v, v ∈ W.grading 1 → W.L 1 v = 0) :
    (canonicalInvariantForm hne hqp).Normalized :=

theorem invariantForm_unique_FHL (hne : W.vacuum ≠ 0)
    (hqp : ∀ v, v ∈ W.grading 1 → W.L 1 v = 0) {B : InvariantForm W}
    (hB : B.IsFullyInvariant) (hn : B.Normalized) :
    B.form = (canonicalInvariantForm hne hqp).form :=

theorem invariantForm_nondegenerate_FHL (hne : W.vacuum ≠ 0)
    (hqp : ∀ v, v ∈ W.grading 1 → W.L 1 v = 0) (hsimple : W.IsSimple) :
    InvariantForm.leftRadical (canonicalInvariantForm hne hqp) = ⊥ :=

theorem eval_oneMode_grading1_FHL (hne : W.vacuum ≠ 0)
    (hqp : ∀ v, v ∈ W.grading 1 → W.L 1 v = 0) {u v : V}
    (hu : u ∈ W.grading 1) (hv : v ∈ W.grading 1) :
    W.smul u 1 v = -(canonicalInvariantForm hne hqp).form u v • W.vacuum :=
\end{leancode}
\noindent\textbf{Hypotheses / status.} $\lfourvac\ne0$, $L_1\lfourV_1 = 0$; symmetry \textbf{not proved}\par
\medskip
\begin{btheorem}{5.57}[Li, Cor.~3.2]
Let $B$ be a fully invariant normalised form on $\lfourV$. Then
\begin{enumerate}
\item $\lfourrad_LB = \lfourm(\lfourV) = \lfourrad_RB$;
\item $\lfourV$ is simple if and only if $\lfourrad_LB = 0$;
\item for homogeneous $x\in\lfourV_n$: $x\in\lfourm(\lfourV)\iff(x|y) = 0$ for all $y\in\lfourV_n$; equivalently, the
weight-$n$ part of $\lfourm(\lfourV)$ is the radical of the restriction of $B$ to $\lfourV_n\times\lfourV_n$.
\end{enumerate}
\end{btheorem}
{\raggedright\noindent\textbf{Lean.} \lean{VOA.}\allowbreak\lean{LiRadical.}\allowbreak\lean{leftRadical_eq_maxGradedIdeal},\allowbreak\ \lean{VOA.}\allowbreak\lean{LiRadical.}\allowbreak\lean{leftRadical_eq_rightRadical},\allowbreak\ \lean{VOA.}\allowbreak\lean{LiRadical.}\allowbreak\lean{isSimple_iff_leftRadical_eq_bot},\allowbreak\ \lean{VOA.}\allowbreak\lean{LiRadical.}\allowbreak\lean{mem_maxGradedIdeal_iff_grading},\allowbreak\ \lean{VOA.}\allowbreak\lean{LiRadical.}\allowbreak\lean{gradedForm},\allowbreak\ \lean{VOA.}\allowbreak\lean{LiRadical.}\allowbreak\lean{ker_gradedForm}\par}
{\raggedright\noindent File \leanpath{VOA/\allowbreak Core/\allowbreak InvariantFormRadical.lean}.\par}
\begin{leancode}
theorem leftRadical_eq_maxGradedIdeal {B : InvariantForm W} (hB : B.IsFullyInvariant)
    (hn : B.Normalized) :
    InvariantForm.leftRadical B = QuotVOA.maxGradedIdeal W :=

theorem leftRadical_eq_rightRadical {B : InvariantForm W} (hB : B.IsFullyInvariant)
    (hn : B.Normalized) :
    InvariantForm.leftRadical B = SingVec.rightRadical B := by

theorem isSimple_iff_leftRadical_eq_bot {B : InvariantForm W} (hB : B.IsFullyInvariant)
    (hn : B.Normalized) :
    W.IsSimple ↔ InvariantForm.leftRadical B = ⊥ := by

theorem mem_maxGradedIdeal_iff_grading {B : InvariantForm W} (hB : B.IsFullyInvariant)
    (hn : B.Normalized) {n : ℤ} {x : 𝒱} (hx : x ∈ W.grading n) :
    x ∈ QuotVOA.maxGradedIdeal W ↔ ∀ y ∈ W.grading n, B.form x y = 0 := by

def gradedForm (n : ℤ) : LinearMap.BilinForm ℂ (W.grading n) :=
  B.form.compl₁₂ (W.grading n).subtype (W.grading n).subtype

theorem ker_gradedForm {B : InvariantForm W} (hB : B.IsFullyInvariant) (hn : B.Normalized)
    (n : ℤ) :
    LinearMap.ker (gradedForm B n) = (QuotVOA.maxGradedIdeal W).comap (W.grading n).subtype := by
\end{leancode}
\noindent\textbf{Hypotheses / status.} $B$ fully inv., normalised\par
\medskip
\endgroup
\part{Examples, deformations and rigidity}
\setcounter{chapter}{5}
\chapter{Heisenberg and lattice vertex algebras in the library}
\begingroup
\providecommand{\lfivevac}{\mathbf{1}}
\providecommand{\lfiveF}{\mathcal{F}}
\providecommand{\lfiveR}{\mathbb{R}}
\providecommand{\lfiveO}{O_2}
\providecommand{\lfiveA}{A_2}
\providecommand{\lfivesl}{\mathfrak{sl}_2}
\noindent{\small 6 of the chapter's 34 numbered statements: its headline results and every definition they use.}\par\medskip
\begin{btheorem}{6.7}[The axioms VA1--VA5 of the library's \texttt{PreVertexAlgebra}]
For every $k\in\C$: $Y(\lfivevac,z)=\mathrm{id}$; $Y(a_{-1}\lfivevac,z)=\alpha(z)$; for every $u$,
$u_{(-1)}\lfivevac=u$ and $u_{(n)}\lfivevac=0$ for $n\ge0$; any two $Y(u,z)$, $Y(v,z)$ are mutually local;
for all $u,v$ there is $N$ with $u_{(n)}v=0$ for $n\ge N$; and if $u\in\lfiveF_{(p)}$, $v\in\lfiveF_{(q)}$
then $u_{(n)}v\in\lfiveF_{(p+q-n-1)}$.
\end{btheorem}
{\raggedright\noindent\textbf{Lean.} \lean{VOA.}\allowbreak\lean{Y_vacuum},\allowbreak\ \lean{VOA.}\allowbreak\lean{Y_current},\allowbreak\ \lean{VOA.}\allowbreak\lean{Y_creation},\allowbreak\ \lean{VOA.}\allowbreak\lean{Y_local},\allowbreak\ \lean{VOA.}\allowbreak\lean{va5_regularity},\allowbreak\ \lean{VOA.}\allowbreak\lean{va1_grading_mode}\par}
{\raggedright\noindent Files \leanpath{VOA/\allowbreak Examples/\allowbreak HeisenbergVertex.lean},
\leanpath{VOA/\allowbreak Examples/\allowbreak HeisenbergCreation.lean}, \leanpath{VOA/\allowbreak Examples/\allowbreak HeisenbergStructure.lean},
\leanpath{VOA/\allowbreak Examples/\allowbreak HeisenbergGrading.lean}.\par}
\begin{leancode}
theorem Y_vacuum (k : ℂ) : Y k vacuum = idField

theorem Y_current (k : ℂ) : Y k (osc k (-1) vacuum) = current k

theorem Y_creation (k : ℂ) (u : Fock) :
    mode (Y k u) (-1) vacuum = u ∧ ∀ n : ℤ, 0 ≤ n → mode (Y k u) n vacuum = 0

theorem Y_local (k : ℂ) (u v : Fock) : Local (Y k u) (Y k v)

theorem va5_regularity (k : ℂ) (u v : Fock) :
    ∃ N : ℤ, ∀ n : ℤ, N ≤ n → mode (Y k u) n v = 0

theorem va1_grading_mode (kk : ℂ) {k ℓ : ℤ} (u : Fock) (hu : u ∈ gradingZ k)
    (v : Fock) (hv : v ∈ gradingZ ℓ) (n : ℤ) :
    mode (Y kk u) n v ∈ gradingZ (k + ℓ - n - 1)
\end{leancode}
\noindent\textbf{Hypotheses / status.} Unconditional (any $k$)\par
\medskip
\begin{btheorem}{6.10}[The Heisenberg vertex operator algebra at level $k$]
For every $k\in\C\setminus\{0\}$, $(\lfiveF,Y,\lfivevac,\omega)$ with the grading
$(\lfiveF_{(n)})_{n\in\Z}$ is a vertex operator algebra of CFT type with central charge $1$.
\end{btheorem}
{\raggedright\noindent\textbf{Lean.} \lean{VOA.}\allowbreak\lean{heisVA},\allowbreak\ \lean{VOA.}\allowbreak\lean{heisenbergVOA},\allowbreak\ \lean{VOA.}\allowbreak\lean{heisenberg_central_charge}\par}
{\raggedright\noindent File \leanpath{VOA/\allowbreak Examples/\allowbreak HeisenbergVOA.lean}.\par}
\begin{leancode}
noncomputable def heisVA (k : ℂ) (_hk : k ≠ 0) : PreVertexAlgebra Fock where
  grading := gradingZ
  isInternal := gradingZ_isInternal
  finrank := gradingZ_finrank
  Y := Y k
  vacuum := vacuum
  vacuum_mem := mem_gradingZ_of_gradingW vacuum_mem_gradingW
  grading_mode := fun u hu v hv n => va1_grading_mode k u hu v hv n
  vacuum_Y := by rw [Y_vacuum, idField_eq_one]
  creation := fun u => Y_creation k u
  locality := fun u v => Y_local k u v
  regularity := fun u v => va5_regularity k u v

noncomputable def heisenbergVOA (k : ℂ) (hk : k ≠ 0) : VertexOperatorAlgebra Fock where
  toPreVertexAlgebra := heisVA k hk
  conformalVector := ω k
  conformal_mem := omega_mem_gradingZ_two k
  centralCharge := 1
  virasoro := fun m n => conf_virasoro k hk m n
  L0_grading := fun v hv => conf_L0_grading k hk v hv
  translation := fun v => conf_translation k hk v
  vacuum_grade := gradingZ_zero_eq_span_vacuum
  positivity := fun n hn => gradingZ_positivity n hn

theorem heisenberg_central_charge (k : ℂ) (hk : k ≠ 0) :
    (heisenbergVOA k hk).centralCharge = 1 := rfl
\end{leancode}
\noindent\textbf{Hypotheses / status.} $k\ne0$\par
\medskip
\begin{btheorem}{6.18}[The free boson at general central charge]
For every $\lambda\in\C$, the vertex algebra $\lfiveF$ of Theorem~6.10 at level $1$, with
conformal vector $\omega^\lambda$, is a vertex operator algebra of CFT type with central charge
$1-12\lambda^2$; its vertex algebra, its grading, $L_0$ and $L_{-1}$ are those of the undeformed algebra.
Consequently every complex number is the central charge of a vertex operator algebra structure on $\lfiveF$.
\end{btheorem}
{\raggedright\noindent\textbf{Lean.} \lean{VOA.}\allowbreak\lean{FBgc.}\allowbreak\lean{freeBosonVOA},\allowbreak\ \lean{VOA.}\allowbreak\lean{FBgc.}\allowbreak\lean{freeBosonVOA_centralCharge},\allowbreak\ \lean{VOA.}\allowbreak\lean{FBgc.}\allowbreak\lean{freeBosonVOA_toPreVertexAlgebra},\allowbreak\ \lean{VOA.}\allowbreak\lean{FBgc.}\allowbreak\lean{freeBosonVOA_L_zero},\allowbreak\ \lean{VOA.}\allowbreak\lean{FBgc.}\allowbreak\lean{exists_centralCharge},\allowbreak\ \lean{VOA.}\allowbreak\lean{FBgc.}\allowbreak\lean{centralCharge_surjective}\par}
{\raggedright\noindent File \leanpath{VOA/\allowbreak Examples/\allowbreak FreeBosonGeneralC.lean}.\par}
\begin{leancode}
noncomputable def freeBosonVOA (lam : ℂ) : VertexOperatorAlgebra Fock where
  toPreVertexAlgebra := heisVA 1 one_ne_zero
  conformalVector := omegaLD lam
  conformal_mem := omegaLD_mem_gradingZ_two lam
  centralCharge := 1 - 12 * lam ^ 2
  virasoro := fun m n => confLD_virasoro lam m n
  L0_grading := fun v hv => confLD_L0_grading lam v hv
  translation := fun v => confLD_translation lam v
  vacuum_grade := gradingZ_zero_eq_span_vacuum
  positivity := fun n hn => gradingZ_positivity n hn

theorem freeBosonVOA_centralCharge (lam : ℂ) :
    (freeBosonVOA lam).centralCharge = 1 - 12 * lam ^ 2 := rfl

theorem freeBosonVOA_toPreVertexAlgebra (lam : ℂ) :
    (freeBosonVOA lam).toPreVertexAlgebra = heisVA 1 one_ne_zero := rfl

theorem freeBosonVOA_L_zero (lam : ℂ) :
    (freeBosonVOA lam).L 0 = (heisenbergVOA 1 one_ne_zero).L 0

theorem exists_centralCharge (c : ℂ) :
    ∃ lam : ℂ, (freeBosonVOA lam).centralCharge = c

theorem centralCharge_surjective :
    Function.Surjective (fun lam : ℂ => (freeBosonVOA lam).centralCharge) :=
  fun c => exists_centralCharge c
\end{leancode}
\noindent\textbf{Hypotheses / status.} Unconditional\par
\medskip
\begin{btheorem}{6.20}[The grading is the expected one]
For any integral lattice $L$, $V_L=\bigoplus_n(V_L)_n$ (no hypothesis). If $L$ is even and positive-definite,
then for every $\alpha\in L$ the charged vacuum $e^\alpha=1\otimes e^\alpha$ lies in the weight space
$(V_L)_{\langle\alpha,\alpha\rangle/2}$ of the vertex operator algebra $V_L$; more generally
$p\otimes e^\alpha$ with $p$ of oscillator level $w$ has weight $w+\frac12\langle\alpha,\alpha\rangle$.
\end{btheorem}
{\raggedright\noindent\textbf{Lean.} \lean{VOA.}\allowbreak\lean{LatGen.}\allowbreak\lean{bigradα_le_gradingLatGen}\par}
{\raggedright\noindent File \leanpath{VOA/\allowbreak Examples/\allowbreak LatticeGenGrading.lean}.\par}
\begin{leancode}
theorem bigradα_le_gradingLatGen (w : ℤ) (α : L.carrier) :
    bigradα L (w, α) ≤ gradingLatGen L (w + halfNorm L α)
\end{leancode}
\noindent\textbf{Hypotheses / status.} Unconditional\par
\medskip
\begin{btheorem}{6.21}[Lattice vertex operator algebras]
Let $L$ be an even positive-definite integral lattice of rank $r\ge0$. Then $V_L$, with the FLM vertex
operators, vacuum $1\otimes e^0$, conformal vector $\omega=\frac12\sum_{a=1}^rh_a(-1)^2\lfivevac\otimes e^0$ and the
grading of Theorem~6.20, is a vertex operator algebra of CFT type with central charge $r$.
\end{btheorem}
{\raggedright\noindent\textbf{Lean.} \lean{VOA.}\allowbreak\lean{LatGen.}\allowbreak\lean{latticeVOAGen},\allowbreak\ \lean{VOA.}\allowbreak\lean{LatGen.}\allowbreak\lean{latticeVOAGen_centralCharge}\par}
{\raggedright\noindent File \leanpath{VOA/\allowbreak Examples/\allowbreak LatticeGenInstance.lean}.\par}
\begin{leancode}
noncomputable def latticeVOAGen (hEven : L.Even) (hPos : L.PosDef) :
    VertexOperatorAlgebra (LatModuleGen L) where
  grading := gradingLatGen L
  isInternal := gradingLatGen_isInternal_of_posDef L hEven hPos
  finrank := gradingLatGen_finrank_of_posDef L hEven hPos
  Y := YfullGen L hEven hPos
  vacuum := latVacuumGen L
  vacuum_mem := by
    rw [gradingLatGen_zero_eq_span_vacuum L hEven hPos]
    exact Submodule.mem_span_singleton_self (latVacuumGen L)
  grading_mode := fun u hu v hv j => YfullGen_grading_mode L hEven hPos u hu v hv j
  vacuum_Y := YfullGen_vacuum L hEven hPos
  creation := YfullGen_creation L hEven hPos
  locality := YfullGen_local L hEven hPos
  regularity := YfullGen_regularity L hEven hPos
  conformalVector := omegaLatGen L
  conformal_mem := omegaLatGen_mem L
  centralCharge := (rank L : ℂ)
  virasoro := YfullGen_virasoro L hEven hPos
  L0_grading := fun v hv => YfullGen_L0_grading L hEven hPos v hv
  translation := YfullGen_translation_omega L hEven hPos
  vacuum_grade := gradingLatGen_zero_eq_span_vacuum L hEven hPos
  positivity := fun n hn => gradingLatGen_positivity L hEven hPos n hn

theorem latticeVOAGen_centralCharge (hEven : L.Even) (hPos : L.PosDef) :
    (latticeVOAGen L hEven hPos).centralCharge = (rank L : ℂ) := rfl
\end{leancode}
\noindent\textbf{Hypotheses / status.} E, P\par
\medskip
\begin{btheorem}{6.27}[Charge additivity]
Let $L$ be even and positive-definite. The boson currents shift charge by $0$, $Y(e^\alpha,z)$ shifts charge by
$\alpha$, and for every state $u$ in the charge-$\alpha$ sector, $Y(u,z)$ shifts charge by $\alpha$.
\end{btheorem}
{\raggedright\noindent\textbf{Lean.} \lean{VOA.}\allowbreak\lean{LatGen.}\allowbreak\lean{alphaLatGen_shiftsCharge},\allowbreak\ \lean{VOA.}\allowbreak\lean{LatGen.}\allowbreak\lean{Yeαvec_shiftsCharge},\allowbreak\ \lean{VOA.}\allowbreak\lean{LatGen.}\allowbreak\lean{YfullGen_shiftsCharge}\par}
{\raggedright\noindent File \leanpath{VOA/\allowbreak Examples/\allowbreak LatticeCharge.lean}.\par}
\begin{leancode}
theorem alphaLatGen_shiftsCharge (a : Fin (rank L)) :
    ShiftsCharge L (alphaLatGen L a) 0

theorem Yeαvec_shiftsCharge (α : L.carrier) :
    ShiftsCharge L (Yeαvec L α) α

theorem YfullGen_shiftsCharge (hEven : L.Even) (hPos : L.PosDef) (α : L.carrier)
    {u : LatModuleGen L} (hu : u ∈ chargeSec L α) :
    ShiftsCharge L (YfullGen L hEven hPos u) α
\end{leancode}
\noindent\textbf{Hypotheses / status.} E, P\par
\medskip
\endgroup
\setcounter{chapter}{6}
\chapter[The Virasoro VOA: Verma modules, Kac determinant, screenings]{The Virasoro vertex operator algebra: Verma modules, the Kac determinant, screenings}
\begingroup
\providecommand{\lsixvac}{\mathbf{1}}
\providecommand{\lsixVir}{\mathfrak{Vir}}
\providecommand{\lsixU}{U(\mathfrak{Vir})}
\providecommand{\lsixV}{V(c,0)}
\providecommand{\lsixL}{L(c,0)}
\providecommand{\lsixM}{M(c,h)}
\providecommand{\lsixK}{K}
\providecommand{\lsixhrs}{h_{r,s}}
\providecommand{\lsixct}{c(t)}
\providecommand{\lsixpp}{\mathrm{p}}
\providecommand{\lsixpt}{\mathrm{p}_{2}}
\providecommand{\lsixG}{G}
\providecommand{\lsixform}[2]{\langle #1,#2\rangle}
\noindent{\small 6 of the chapter's 45 numbered statements: its headline results and every definition they use.}\par\medskip
\begin{btheorem}{7.8}[PBW basis of $\lsixV$, at every $c$]
For every $c\in\C$ and $n\in\N$, the monomials $L_w\lsixvac$ with $w$ sorted of energy $n$ form a
basis of $\lsixV_n$. Hence $\dim\lsixV_n = \lsixpt(n)$, and the sorted monomials of all energies form a
basis of $\lsixV$. In particular $\dim\lsixV_n = 1,0,1,1,2,2,4,4,7$ for $n = 0,\dots,8$ and
$\dim\lsixV_{16} = 55$.
\end{btheorem}
{\raggedright\noindent\textbf{Lean.} \lean{VOA.}\allowbreak\lean{VirModule.}\allowbreak\lean{gradedBasis},\allowbreak\ \lean{VOA.}\allowbreak\lean{VirModule.}\allowbreak\lean{finrank_gradingVir},\allowbreak\ \lean{VOA.}\allowbreak\lean{VirModule.}\allowbreak\lean{sortedWord_linearIndependent},\allowbreak\ \lean{VOA.}\allowbreak\lean{VirModule.}\allowbreak\lean{sortedWords},\allowbreak\ \lean{VOA.}\allowbreak\lean{VirModule.}\allowbreak\lean{mem_sortedWords},\allowbreak\ \lean{VOA.}\allowbreak\lean{VirModule.}\allowbreak\lean{finrank_gradingVir_eq_card},\allowbreak\ \lean{VOA.}\allowbreak\lean{VirModule.}\allowbreak\lean{card_sortedWords_le_eight},\allowbreak\ \lean{VOA.}\allowbreak\lean{VirModule.}\allowbreak\lean{finrank_gradingVir_sixteen}\par}
{\raggedright\noindent File \leanpath{VOA/\allowbreak Examples/\allowbreak VirasoroPartitionBasis.lean}.\par}
\begin{leancode}
noncomputable def gradedBasis (c : ℂ) (n : ℕ) : Module.Basis (SortedOfEnergy n) ℂ
    (gradingVir c (n : ℤ)) :=
  Module.Basis.mk (gradedWord_linearIndependent c n) (gradedWord_span c n)

theorem finrank_gradingVir (c : ℂ) (n : ℕ) :
    Module.finrank ℂ (gradingVir c (n : ℤ)) = Fintype.card (SortedOfEnergy n) :=

theorem sortedWord_linearIndependent (c : ℂ) :
    LinearIndependent ℂ (fun w : {w : List ℕ // SortedWord w} => negWordVir c w.1) := by
\end{leancode}
{\raggedright\noindent File \leanpath{VOA/\allowbreak Examples/\allowbreak VirasoroTopSymbol.lean}.\par}
\begin{leancode}
def sortedWords (n : ℕ) : Finset (List ℕ) := sortedAux n n n

theorem mem_sortedWords (n : ℕ) (w : List ℕ) :
    w ∈ sortedWords n ↔ (SortedWord w ∧ energyVir w = n) := by

theorem finrank_gradingVir_eq_card (c : ℂ) (n : ℕ) :
    Module.finrank ℂ (gradingVir c (n : ℤ)) = (sortedWords n).card := by

theorem card_sortedWords_le_eight :
    (sortedWords 0).card = 1 ∧ (sortedWords 1).card = 0 ∧ (sortedWords 2).card = 1 ∧
      (sortedWords 3).card = 1 ∧ (sortedWords 4).card = 2 ∧ (sortedWords 5).card = 2 ∧
      (sortedWords 6).card = 4 ∧ (sortedWords 7).card = 4 ∧ (sortedWords 8).card = 7 :=

theorem finrank_gradingVir_sixteen (c : ℂ) : Module.finrank ℂ (gradingVir c 16) = 55 := by
\end{leancode}
\noindent\textbf{Hypotheses / status.} None; unconditional\par
\medskip
\begin{btheorem}{7.18}[Graded dimension of $\lsixL$ is a Gram rank]
Let $G_n(c)$ be the matrix of $\lsixform\cdot\cdot_c$ on $\lsixV_n$ in the partition basis. Then
$\dim\lsixL_n = \operatorname{rank}G_n(c)$ for every $c$ and $n$. In particular $\dim\lsixL_n =
\lsixpt(n)$ if $\det G_n(c)\ne0$, and $\dim\lsixL_n<\lsixpt(n)$ if $\det G_n(c) = 0$.
\end{btheorem}
{\raggedright\noindent\textbf{Lean.} \lean{VOA.}\allowbreak\lean{VirSimpleGram.}\allowbreak\lean{finrank_virasoroSimple_grading_partition},\allowbreak\ \lean{VOA.}\allowbreak\lean{VirSimpleGram.}\allowbreak\lean{finrank_virasoroSimple_lt_of_det_eq_zero},\allowbreak\ \lean{VOA.}\allowbreak\lean{GramMat.}\allowbreak\lean{canGram}\par}
{\raggedright\noindent File \leanpath{VOA/\allowbreak Examples/\allowbreak VirasoroSimpleGram.lean}.\par}
\begin{leancode}
theorem finrank_virasoroSimple_grading_partition (n : ℕ) :
    finrank ℂ ((virasoroSimple c).grading (n : ℤ))
      = (GramMat.canGram (vacVec_ne_zero c) (virasoroVOA_qp c) (n : ℤ)
          (gradedBasis c n)).rank :=

theorem finrank_virasoroSimple_lt_of_det_eq_zero {ι : Type*} [Fintype ι] [DecidableEq ι] (n : ℤ)
    (e : Basis ι ℂ ((virasoroVOA c).grading n))
    (h : (GramMat.canGram (vacVec_ne_zero c) (virasoroVOA_qp c) n e).det = 0) :
    finrank ℂ ((virasoroSimple c).grading n) < finrank ℂ ((virasoroVOA c).grading n) :=
\end{leancode}
{\raggedright\noindent File \leanpath{VOA/\allowbreak Core/\allowbreak GradedFormMatrix.lean}.\par}
\begin{leancode}
noncomputable def canGram : Matrix ι ι ℂ := gramMatrix (canonicalInvariantForm hne hqp) n e
\end{leancode}
\noindent\textbf{Hypotheses / status.} None; unconditional\par
\medskip
\begin{btheorem}{7.30}[Singular vectors versus the Gram determinant]
Let $n>0$.
\begin{enumerate}
\item If $M(c,h)$ has a singular vector at level $n$, then $\det G_n(c,h) = 0$; conversely if
  $\det G_n(c,h)\ne0$ there is none.
\item If $\det G_n(c,h) = 0$ and $\det G_k(c,h)\ne0$ for all $k<n$, then $M(c,h)$ has a singular vector at
  level $n$.
\item A singular vector $v$ at level $n$ generates an injective image of $M(c,h+n)$; consequently
  $(X-h)^{\lsixpp(m)}$ divides $\det G_{n+m}(c,X)$ as a polynomial in the highest weight $X$, for every
  $m\ge0$.
\item Singular vectors at a given level are unique up to scalar.
\end{enumerate}
\end{btheorem}
{\raggedright\noindent\textbf{Lean.} \lean{VOA.}\allowbreak\lean{InverseDS.}\allowbreak\lean{VirasoroVerma.}\allowbreak\lean{vermaGramMatrix_nonsingular_no_singular},\allowbreak\ \lean{VOA.}\allowbreak\lean{InverseDS.}\allowbreak\lean{VirasoroVerma.}\allowbreak\lean{verma_singular_of_first_gram_root},\allowbreak\ \lean{VOA.}\allowbreak\lean{InverseDS.}\allowbreak\lean{VirasoroVerma.}\allowbreak\lean{vermaSingularMap},\allowbreak\ \lean{VOA.}\allowbreak\lean{InverseDS.}\allowbreak\lean{VirasoroVerma.}\allowbreak\lean{vermaSingularMap_injective},\allowbreak\ \lean{VOA.}\allowbreak\lean{InverseDS.}\allowbreak\lean{VirasoroVerma.}\allowbreak\lean{vermaWeightGramPolynomial},\allowbreak\ \lean{VOA.}\allowbreak\lean{InverseDS.}\allowbreak\lean{VirasoroVerma.}\allowbreak\lean{vermaWeightDeterminant_dvd_of_singular},\allowbreak\ \lean{VOA.}\allowbreak\lean{InverseDS.}\allowbreak\lean{VirasoroVerma.}\allowbreak\lean{verma_singular_eq_smul}\par}
{\raggedright\noindent File \leanpath{VOA/\allowbreak Examples/\allowbreak VirasoroVermaGramPolynomial.lean}.\par}
\begin{leancode}
theorem vermaGramMatrix_nonsingular_no_singular (c h : ℂ) (n : ℕ)
    [DecidableEq (VermaPartitions n)] (hn : 0 < n)
    (hd : (vermaGramMatrix c h n).det ≠ 0)
    (v : Verma c h) (hv : v ∈ sortedSpan c h n)
    (h1 : vermaMode c h 1 v = 0) (h2 : vermaMode c h 2 v = 0) : v = 0 := by

theorem verma_singular_of_first_gram_root (c h : ℂ) (n : ℕ)
    [∀ k : ℕ, DecidableEq (VermaPartitions k)]
    (hroot : (vermaGramMatrix c h n).det = 0)
    (hlower : ∀ k < n, (vermaGramMatrix c h k).det ≠ 0) :
    ∃ v : Verma c h, v ≠ 0 ∧ v ∈ sortedSpan c h n ∧
      ∀ m : ℤ, 0 < m → vermaMode c h m v = 0 := by
\end{leancode}
{\raggedright\noindent File \leanpath{VOA/\allowbreak Examples/\allowbreak VirasoroVermaSingularEmbedding.lean}.\par}
\begin{leancode}
noncomputable def vermaSingularMap (c h : ℂ) (n : ℕ) (v : Verma c h)
    (hv : v ∈ sortedSpan c h n)
    (h1 : vermaMode c h 1 v=0) (h2 : vermaMode c h 2 v=0) :
    Verma c (h+(n:ℂ)) →ₗ[ℂ] Verma c h :=
  universalMap (vermaRep c h) v (verma_singular_isHighest c h n v hv h1 h2)

theorem vermaSingularMap_injective (c h : ℂ) (n : ℕ) (v : Verma c h)
    (hv : v ∈ sortedSpan c h n) (hne : v≠0)
    (h1 : vermaMode c h 1 v=0) (h2 : vermaMode c h 2 v=0) :
    Function.Injective (vermaSingularMap c h n v hv h1 h2) :=

noncomputable def vermaWeightGramPolynomial (c : ℂ) (n : ℕ) :
    Matrix (VermaPartitions n) (VermaPartitions n) (Polynomial ℂ) :=
  vermaAlgebraGramMatrix (Polynomial.C c) Polynomial.X n

theorem vermaWeightDeterminant_dvd_of_singular (c h : ℂ) (n m : ℕ)
    [DecidableEq (VermaPartitions (n+m))] (hn : 0<n)
    (v : sortedSpan c h n) (hv : v≠0)
    (h1 : vermaMode c h 1 (v : Verma c h)=0)
    (h2 : vermaMode c h 2 (v : Verma c h)=0) :
    (Polynomial.X-Polynomial.C h)^(Fintype.card (VermaPartitions m)) ∣
      (vermaWeightGramPolynomial c (n+m)).det := by
\end{leancode}
{\raggedright\noindent File \leanpath{VOA/\allowbreak Examples/\allowbreak VirasoroSingularUniqueness.lean}.\par}
\begin{leancode}
theorem verma_singular_eq_smul (c h : ℂ) (n : ℕ) (v w : Verma c h)
    (hv : v ∈ sortedSpan c h n) (hw : w ∈ sortedSpan c h n) (hwne : w ≠ 0)
    (hv1 : vermaMode c h 1 v = 0) (hv2 : vermaMode c h 2 v = 0)
    (hw1 : vermaMode c h 1 w = 0) (hw2 : vermaMode c h 2 w = 0) :
    v = (vermaTopCoordinate c h n v / vermaTopCoordinate c h n w) • w := by
\end{leancode}
\noindent\textbf{Hypotheses / status.} Singular vector; unconditional as implication\par
\medskip
\begin{btheorem}{7.34}[Kac formula from singular vectors (conditional form)]
Let $c\in\C$, $n\in\N$ and $t\in\C$ with $\operatorname{Im}t\ne0$. Suppose that for every label $(r,s)$ with
$rs\le n$ the Verma module $M(c,h_{r,s}(t))$ has a singular vector at level $rs$. Then
\begin{equation*}
 \det G_n(c,X) = K_n\cdot\mathcal K_n(t;X).
\end{equation*}
\end{btheorem}
{\raggedright\noindent\textbf{Lean.} \lean{VOA.}\allowbreak\lean{InverseDS.}\allowbreak\lean{VirasoroVerma.}\allowbreak\lean{fusionKacWeight_injective_positive},\allowbreak\ \lean{VOA.}\allowbreak\lean{InverseDS.}\allowbreak\lean{VirasoroVerma.}\allowbreak\lean{vermaWeightDeterminant_eq_kac_of_singular}\par}
{\raggedright\noindent File \leanpath{VOA/\allowbreak Examples/\allowbreak VirasoroKacProduct.lean}.\par}
\begin{leancode}
theorem fusionKacWeight_injective_positive (t : ℂ) (ht : t.im ≠ 0)
    (r s r' s' : ℕ) (hr : 0 < r) (hr' : 0 < r')
    (he : fusionKacWeight r s t = fusionKacWeight r' s' t) : r = r' ∧ s = s' := by

theorem vermaWeightDeterminant_eq_kac_of_singular (c t : ℂ) (n : ℕ) (ht : t.im ≠ 0)
    (hs : ∀ i ∈ kacLabels n,
      ∃ v : sortedSpan c (fusionKacWeight (i.1.val+1) (i.2.val+1) t)
        ((i.1.val+1)*(i.2.val+1)), v ≠ 0 ∧
        vermaMode c (fusionKacWeight (i.1.val+1) (i.2.val+1) t) 1 v = 0 ∧
        vermaMode c (fusionKacWeight (i.1.val+1) (i.2.val+1) t) 2 v = 0) :
    (vermaWeightGramPolynomial c n).det =
      Polynomial.C (∏ w : VermaPartitions n, vermaDiagonalContraction w.val) *
        kacWeightProduct t n := by
\end{leancode}
\noindent\textbf{Hypotheses / status.} $\operatorname{Im}t\ne0$, \lean{hs}; conditional (discharged)\par
\medskip
\begin{btheorem}{7.42}[Kac roots and singular vectors at all labels]
For all $r,s\ge1$:
\begin{enumerate}
\item for every integer $m\ge1$, $\det G_{rs}(c(m^2),h_{r,s}(m^2)) = 0$;
\item hence for every $t\ne0$, $\det G_{rs}(c(t),h_{r,s}(t)) = 0$;
\item for every nonreal $t$, $M(c(t),h_{r,s}(t))$ has a singular vector at level $rs$.
\end{enumerate}
\end{btheorem}
{\raggedright\noindent\textbf{Lean.} \lean{VOA.}\allowbreak\lean{VirasoroScreening.}\allowbreak\lean{integer_screening_kac_root},\allowbreak\ \lean{VOA.}\allowbreak\lean{VirasoroScreening.}\allowbreak\lean{infinite_integer_screening_roots},\allowbreak\ \lean{VOA.}\allowbreak\lean{VirasoroScreening.}\allowbreak\lean{screening_kac_root},\allowbreak\ \lean{VOA.}\allowbreak\lean{VirasoroScreening.}\allowbreak\lean{screening_nonreal_kac_singular}\par}
{\raggedright\noindent File \leanpath{VOA/\allowbreak Examples/\allowbreak VirasoroScreeningKacRoots.lean}.\par}
\begin{leancode}
theorem integer_screening_kac_root (r s m : ℕ) (hr : 0 < r) (hs : 0 < s)
    (hm : 0 < m) :
    (vermaGramMatrix (vermaBSACharge ((m:ℂ)^2))
      (fusionKacWeight r s ((m:ℂ)^2)) (r*s)).det = 0 := by

theorem infinite_integer_screening_roots (r s : ℕ) (hr : 0 < r) (hs : 0 < s) :
    ∃ S : Set ℂ, S.Infinite ∧ (∀ t ∈ S, t ≠ 0) ∧
      ∀ t ∈ S, (vermaGramMatrix (vermaBSACharge t)
        (fusionKacWeight r s t) (r*s)).det = 0 := by

theorem screening_kac_root (r s : ℕ) (hr : 0 < r) (hs : 0 < s)
    (t : ℂ) (ht : t ≠ 0) :
    (vermaGramMatrix (vermaBSACharge t) (fusionKacWeight r s t) (r*s)).det = 0 := by

theorem screening_nonreal_kac_singular (t : ℂ) (ht : t.im ≠ 0)
    (r s : ℕ) (hr : 0 < r) (hs : 0 < s) :
    ∃ v : sortedSpan (vermaBSACharge t) (fusionKacWeight r s t) (r*s),
      v ≠ 0 ∧ vermaMode (vermaBSACharge t) (fusionKacWeight r s t) 1 v = 0 ∧
        vermaMode (vermaBSACharge t) (fusionKacWeight r s t) 2 v = 0 := by
\end{leancode}
\noindent\textbf{Hypotheses / status.} $r,s>0$; unconditional\par
\medskip
\begin{btheorem}{7.43}[Kac determinant formula]
For every $t\in\C\setminus\{0\}$, every $n\in\N$ and every $h\in\C$,
\begin{equation*}
 \det G_n\bigl(c(t),h\bigr) = K_n\prod_{\substack{r,s\ge1\\ rs\le n}}\bigl(h-h_{r,s}(t)\bigr)^{\lsixpp(n-rs)} .
\end{equation*}
Since every $c\in\C$ is some $c(t)$, this is the Kac determinant for all Virasoro Verma modules.
\end{btheorem}
\noindent\emph{(The chapter quotes no Lean declaration for this statement; see book p.~340.)}\par
\noindent\textbf{Hypotheses / status.} $t\ne0$; unconditional\ (check not in this release)\par
\medskip
\endgroup
\setcounter{chapter}{7}
\chapter{First-order deformations and the graded first-order rigidity of Virasoro minimal models}
\begingroup
\providecommand{\lsevenvac}{\mathbf{1}}
\providecommand{\lsevenV}{\mathcal{V}}
\providecommand{\lsevenR}{R}
\providecommand{\lsevenE}{\mathcal{E}}
\providecommand{\lsevenK}{\mathsf{K}}
\providecommand{\lsevenN}{\mathsf{N}}
\providecommand{\lsevenLt}{\widetilde{L}}
\providecommand{\lsevencpq}{c_{p,q}}
\noindent{\small 7 of the chapter's 19 numbered statements: its headline results and every definition they use.}\par\medskip
\begin{btheorem}{8.4}[Coboundaries are cocycles]
For every linear $\varphi$ with $\varphi(\lsevenvac) = 0$, the cochain $\delta\varphi$ satisfies
(D1)--(D4). In particular every coboundary is the $Y_1$ of a first-order deformation.
\end{btheorem}
{\raggedright\noindent\textbf{Lean.} \lean{VOA.}\allowbreak\lean{Deform.}\allowbreak\lean{cocycles},\allowbreak\ \lean{VOA.}\allowbreak\lean{Deform.}\allowbreak\lean{coboundaries_le_cocycles},\allowbreak\ \lean{VOA.}\allowbreak\lean{NonzeroDeform.}\allowbreak\lean{deltaCochain_mem_cocycles}\par}
\begin{leancode}
def cocycles (W : VertexOperatorAlgebra 𝒱) : Submodule ℂ (Cochain 𝒱) where
  carrier := {f | Truncated f ∧ f W.vacuum = 0 ∧ CreationAx W f ∧ JacobiAx W f}

theorem coboundaries_le_cocycles (W : VertexOperatorAlgebra 𝒱) :
    coboundaries W ≤ cocycles W := by

theorem deltaCochain_mem_cocycles (W : VertexOperatorAlgebra 𝒱) {φ : 𝒱 →ₗ[ℂ] 𝒱}
    (hφ0 : φ W.vacuum = 0) : deltaCochain W φ ∈ Deform.cocycles W := by
\end{leancode}
\noindent Files \leanpath{VOA/\allowbreak Core/\allowbreak DeformationGraded.lean}, \leanpath{VOA/\allowbreak Examples/\allowbreak NonzeroDeformation.lean}.
\noindent\textbf{Hypotheses / status.} None; unconditional\par
\medskip
\begin{btheorem}{8.14}[The gauge line]
Let $W$ be any VOA with $\lsevenV_2 = \C\omega$ and let $f$ be a coboundary. Then there is $\mu\in\C$
with $f(\omega)_{(1)}\omega = 2\mu\,\omega$ and $f(\omega)_{(3)}\omega = \mu c\,\lsevenvac$. Consequently, for a
trivial graded deformation the scalars of (8.6) satisfy $\alpha = \mu c$,
$\beta = 2\mu$, hence
\begin{equation}
 \delta := 2\alpha - \beta c = 0 .
\end{equation}
\end{btheorem}
{\raggedright\noindent\textbf{Lean.} \lean{VOA.}\allowbreak\lean{VirCDef.}\allowbreak\lean{cmode_conformal_of_isCoboundary}\par}
\begin{leancode}
theorem cmode_conformal_of_isCoboundary (W : VertexOperatorAlgebra 𝒱)
    (hV2 : W.grading 2 = ℂ ∙ W.conformalVector) {f : Cochain 𝒱} (hf : IsCoboundary W f) :
    ∃ μ : ℂ,
      cmode f W.conformalVector 1 W.conformalVector = (2 * μ) • W.conformalVector ∧
      cmode f W.conformalVector 3 W.conformalVector = (μ * W.centralCharge) • W.vacuum := by
\end{leancode}
\noindent File \leanpath{VOA/\allowbreak Examples/\allowbreak VirasoroCDeformation.lean}.
\noindent\textbf{Hypotheses / status.} $\lsevenV_2 = \C\omega$; unconditional\par
\medskip
\begin{btheorem}{8.16}[Reconstruction from a closed lift]
Let $W$ be simple and $D$ a first-order deformation. Suppose $S\subseteq\lsevenV\times\lsevenV$ is mode-closed,
grade-closed, contains $(\lsevenvac,0)$, does not contain $(0,\lsevenvac)$, and its reduction is surjective.
Then $D$ is trivial.
\end{btheorem}
{\raggedright\noindent\textbf{Lean.} \lean{VOA.}\allowbreak\lean{Deform.}\allowbreak\lean{Reconstruction.}\allowbreak\lean{isTrivial_of_closed_lift}\par}
\begin{leancode}
theorem isTrivial_of_closed_lift (S : Submodule ℂ (V × V))
    (hsimple : W.IsSimple) (hclosed : ModeClosed D S)
    (hsurj : Function.Surjective (reduction S))
    (hsep : (0, W.vacuum) ∉ S) (hvac : (W.vacuum, 0) ∈ S)
    (hgrade : GradeClosed W S) : D.IsTrivial := by
\end{leancode}
\noindent File \leanpath{VOA/\allowbreak Core/\allowbreak DeformationReconstruction.lean}.
\noindent\textbf{Hypotheses / status.} $W$ simple + 5 lift properties\par
\medskip
\begin{btheorem}{8.17}[Fixed-charge reconstruction]
Let $D$ be a graded first-order deformation of $L(c,0)$, $c\in\C$ arbitrary, with scalars $\alpha,\beta$
as in (8.6). If $2\alpha-\beta c = 0$ then $D$ is trivial.
\end{btheorem}
{\raggedright\noindent\textbf{Lean.} \lean{VOA.}\allowbreak\lean{VirDeformModes.}\allowbreak\lean{isTrivial_of_zero_central_variation}\par}
\begin{leancode}
variable (c : ℂ) (D : FirstOrderDeformation (virasoroSimple c))

theorem isTrivial_of_zero_central_variation (hD : D.IsGraded) (α β : ℂ)
    (ha : D.smul₁ (virasoroSimple c).conformalVector 3
      (virasoroSimple c).conformalVector = α • (virasoroSimple c).vacuum)
    (hb : D.smul₁ (virasoroSimple c).conformalVector 1
      (virasoroSimple c).conformalVector = β • (virasoroSimple c).conformalVector)
    (hz : 2*α-β*c = 0) : D.IsTrivial := by
\end{leancode}
\noindent File \leanpath{VOA/\allowbreak Examples/\allowbreak VirasoroFixedReconstruction.lean}.
\noindent\textbf{Hypotheses / status.} Graded; unconditional\par
\medskip
\begin{bproposition}{8.18}[Lifting criterion]
For symmetric $B_0$ on finite-dimensional $E$ and any $s\in E$: $s$ lifts to first order if and only if
$B_1(s,v) = 0$ for all $v$ in the radical of $B_0$. Hence, writing $F_1 = \mathrm{rad}\,B_0$ and
$F_2 = \{s\in F_1 : B_1(s,F_1) = 0\}$ (the first two Jantzen layers), $F_2$ is exactly the set of
relations that lift.
\end{bproposition}
{\raggedright\noindent\textbf{Lean.} \lean{VOA.}\allowbreak\lean{DefForm.}\allowbreak\lean{radBil},\allowbreak\ \lean{VOA.}\allowbreak\lean{DefForm.}\allowbreak\lean{lift_iff_pairing_zero},\allowbreak\ \lean{VOA.}\allowbreak\lean{DefForm.}\allowbreak\lean{F1},\allowbreak\ \lean{VOA.}\allowbreak\lean{DefForm.}\allowbreak\lean{F2},\allowbreak\ \lean{VOA.}\allowbreak\lean{RigidBridge.}\allowbreak\lean{F2_lt_F1_of_self_pairing_ne_zero}\par}
\begin{leancode}
def radBil (B : E →ₗ[K] E →ₗ[K] K) : Submodule K E := LinearMap.ker B

theorem lift_iff_pairing_zero [FiniteDimensional K E] {B₀ B₁ : E →ₗ[K] E →ₗ[K] K}
    (hB₀ : ∀ x y : E, B₀ x y = B₀ y x) (s : E) :
    (∃ s₁ : E, ∀ v : E, B₁ s v + B₀ s₁ v = 0) ↔ ∀ v ∈ radBil B₀, B₁ s v = 0 := by

def F1 (B₀ : E →ₗ[K] E →ₗ[K] K) : Submodule K E := radBil B₀

def F2 (B₀ B₁ : E →ₗ[K] E →ₗ[K] K) : Submodule K E where
  carrier := {s | s ∈ F1 B₀ ∧ ∀ v ∈ F1 B₀, B₁ s v = 0}
\end{leancode}
\begin{leancode}
theorem F2_lt_F1_of_self_pairing_ne_zero (hvac : W.vacuum ≠ 0)
    (D : FirstOrderDeformation W) (hD : D.IsGraded) {k : ℤ} {s : W.grading k}
    (hs : s ∈ DefForm.radBil (shapForm W hvac k))
    (hself : shapForm₁ D hD hvac k s s ≠ 0) :
    DefForm.F2 (shapForm W hvac k) (shapForm₁ D hD hvac k) < DefForm.F1 (shapForm W hvac k) :=
\end{leancode}
\noindent Files \leanpath{VOA/\allowbreak Core/\allowbreak DeformedForm.lean}, \leanpath{VOA/\allowbreak Core/\allowbreak RigidityBridge.lean}.
\noindent\textbf{Hypotheses / status.} $B_0$ symmetric, fin.-dim.; unconditional\par
\medskip
\begin{btheorem}{8.20}[Relation obstruction]
Let $c\in\C$. Suppose $s := \sum_ir_iL_{w_i}\lsevenvac = 0$ in $L(c,0)$ for a finite family of words $w_i$ and
coefficients $r_i$, and suppose the norm polynomial satisfies $P'_{w,r}(c)\ne0$. Then every graded
first-order deformation of $L(c,0)$ has central variation $\delta = 0$, and therefore
$L(c,0)$ is graded rigid.
\end{btheorem}
{\raggedright\noindent\textbf{Lean.} \lean{VOA.}\allowbreak\lean{VirDeformModes.}\allowbreak\lean{virWord_pairing_obstruction},\allowbreak\ \lean{VOA.}\allowbreak\lean{VirDeformModes.}\allowbreak\lean{NormalOrdering.}\allowbreak\lean{simple_normalized_word_tangent},\allowbreak\ \lean{VOA.}\allowbreak\lean{VirDeformModes.}\allowbreak\lean{NormalOrdering.}\allowbreak\lean{graded_rigidity_of_relation}\par}
\begin{leancode}
theorem virWord_pairing_obstruction {W : VertexOperatorAlgebra V}
    (D : FirstOrderDeformation W) (β : ℂ) (B : InvariantForm W)
    {ι : Type*} [Fintype ι] (w : ι → List ℤ) (r : ι → ℂ)
    (hr : ∑ i, r i • virWord W.L (w i) W.vacuum = 0) :
    B.form W.vacuum
      ((∑ i, r i • virWord (normalizedL D β) ((w i).reverse.map Neg.neg))
        (∑ j, r j • virWord (normalizedL D β) (w j) (W.vacuum,0))).2 = 0 := by

theorem simple_normalized_word_tangent (c : ℂ)
    (D : FirstOrderDeformation (VirSimple.virasoroSimple c)) (hD : D.IsGraded)
    (α β : ℂ)
    (ha : D.smul₁ (VirSimple.virasoroSimple c).conformalVector 3
      (VirSimple.virasoroSimple c).conformalVector = α • (VirSimple.virasoroSimple c).vacuum)
    (hb : D.smul₁ (VirSimple.virasoroSimple c).conformalVector 1
      (VirSimple.virasoroSimple c).conformalVector = β • (VirSimple.virasoroSimple c).conformalVector)
    (w : List ℤ) :
    (VirQuotForm.Bsimp c).form (VirSimple.virasoroSimple c).vacuum
      (virWord (normalizedL D β) w ((VirSimple.virasoroSimple c).vacuum,0)).2 =
        (wordPolynomial w).derivative.eval c * (2*α-β*c) := by

theorem graded_rigidity_of_relation (c : ℂ) (w : ι → List ℤ) (r : ι → ℂ)
    (hr : ∑ i, r i • virWord (VirSimple.virasoroSimple c).L (w i)
      (VirSimple.virasoroSimple c).vacuum = 0)
    (hd : (relationPolynomial w r).derivative.eval c ≠ 0) :
    IsRigidGraded (VirSimple.virasoroSimple c) := by
\end{leancode}
\noindent Files \leanpath{VOA/\allowbreak Examples/\allowbreak VirasoroWordPairing.lean}, \leanpath{VOA/\allowbreak Examples/\allowbreak VirasoroUniformObstruction.lean}.
\noindent\textbf{Hypotheses / status.} Relation in $L(c,0)$, $P'(c)\ne0$; unconditional\par
\medskip
\begin{btheorem}{8.22}[Graded first-order rigidity of the Virasoro minimal models]
Let $p,q\ge2$ be coprime integers and $\lsevencpq = 1-6(p-q)^2/(pq)$. Every graded first-order deformation of
the simple Virasoro vertex operator algebra $L(\lsevencpq,0)$ is trivial: there is a grade-preserving linear map
$\varphi$ with $\varphi(\lsevenvac) = 0$ such that $u^1_{(n)}v = (\varphi u)_{(n)}v - \varphi(u_{(n)}v) + u_{(n)}\varphi v$ for all
$u,v,n$.
\end{btheorem}
{\raggedright\noindent\textbf{Lean.} \lean{VOA.}\allowbreak\lean{VirModule.}\allowbreak\lean{minimal_graded_rigidity}\par}
\begin{leancode}
theorem minimal_graded_rigidity (p q : ℕ) (hp : 2 ≤ p) (hq : 2 ≤ q)
    (hcop : p.Coprime q) :
    IsRigidGraded (virasoroSimple (minimalCentralCharge p q)) := by
  apply minimal_rigidity_of_infinite_kac_roots p q hp hq hcop
  intro r s hr hs _
  exact infinite_integer_screening_roots r s hr hs
\end{leancode}
\noindent File \leanpath{VOA/\allowbreak Examples/\allowbreak VirasoroMinimalRigidityComplete.lean}.
\noindent\textbf{Hypotheses / status.} $p,q\ge2$ coprime; \textbf{unconditional}, no \lean{pbw}; non-vacuity machine-checked ($V(c,0)$ not graded rigid)\par
\medskip
\endgroup
\setcounter{chapter}{8}
\chapter[Affine \texorpdfstring{$\mathfrak{sl}_2$}{sl(2)}: construction, Wakimoto, Verma, Kac]{Affine \texorpdfstring{$\mathfrak{sl}_2$}{sl(2)}: construction, Verma modules, Wakimoto realisation, screening, non-simplicity}
\begingroup
\providecommand{\leightsl}{\mathfrak{sl}_2}
\providecommand{\leighthat}{\widehat{\mathfrak{sl}}_2}
\providecommand{\leightvac}{|0\rangle}
\providecommand{\leightV}{\mathcal{V}}
\providecommand{\leightU}{\mathcal{U}}
\providecommand{\leightad}{\operatorname{ad}}
\providecommand{\leightM}{\mathcal{M}}
\providecommand{\leightFock}{\mathcal{F}}
\noindent{\small 7 of the chapter's 43 numbered statements: its headline results and every definition they use.}\par\medskip
\begin{bdefinition}{9.1}[$\leightsl$, its structure constants and normalised Killing form]
$\leightsl$ is the complex Lie algebra with basis $e,f,h$ and brackets $[h,e]=2e$, $[h,f]=-2f$,
$[e,f]=h$. Writing $[x_a,x_b] = \sum_c f^{ab}{}_c\,x_c$ for $(x_0,x_1,x_2) = (e,f,h)$, the nonzero
structure constants are $f^{he}{}_e = 2$, $f^{eh}{}_e=-2$, $f^{hf}{}_f=-2$, $f^{fh}{}_f=2$,
$f^{ef}{}_h=1$, $f^{fe}{}_h=-1$. The \emph{normalised Killing form} is the symmetric invariant form with
$\kappa(e,f)=\kappa(f,e)=1$, $\kappa(h,h)=2$ and all other basis pairings zero; it equals the trace
form $\tr(xy)$ in the defining representation and $\frac14$ of the Killing form
$\tr(\leightad x\,\leightad y)$.
\end{bdefinition}
{\raggedright\noindent\textbf{Lean.} \lean{VOA.}\allowbreak\lean{AffineVOA.}\allowbreak\lean{colE},\allowbreak\ \lean{VOA.}\allowbreak\lean{AffineVOA.}\allowbreak\lean{colF},\allowbreak\ \lean{VOA.}\allowbreak\lean{AffineVOA.}\allowbreak\lean{colH},\allowbreak\ \lean{VOA.}\allowbreak\lean{AffineVOA.}\allowbreak\lean{fabc},\allowbreak\ \lean{VOA.}\allowbreak\lean{AffineVOA.}\allowbreak\lean{kappa},\allowbreak\ \lean{VOA.}\allowbreak\lean{AffineVOA.}\allowbreak\lean{Sl2},\allowbreak\ \lean{VOA.}\allowbreak\lean{AffineVOA.}\allowbreak\lean{basis},\allowbreak\ \lean{VOA.}\allowbreak\lean{AffineVOA.}\allowbreak\lean{brComp},\allowbreak\ \lean{VOA.}\allowbreak\lean{AffineVOA.}\allowbreak\lean{AffSl2},\allowbreak\ \lean{VOA.}\allowbreak\lean{AffineVOA.}\allowbreak\lean{WakModule},\allowbreak\ \lean{VOA.}\allowbreak\lean{AffineVOA.}\allowbreak\lean{J}\par}
{\raggedright\noindent File \leanpath{VOA/\allowbreak Examples/\allowbreak AffineSl2.lean}.\par}
\begin{leancode}
abbrev colE : Fin 3 := 0
abbrev colF : Fin 3 := 1
abbrev colH : Fin 3 := 2

def fabc : Fin 3 → Fin 3 → Fin 3 → ℂ
  | 2, 0, 0 => 2    -- ⁅h,e⁆ = 2e
  | 0, 2, 0 => -2   -- ⁅e,h⁆ = -2e
  | 2, 1, 1 => -2   -- ⁅h,f⁆ = -2f
  | 1, 2, 1 => 2    -- ⁅f,h⁆ = 2f
  | 0, 1, 2 => 1    -- ⁅e,f⁆ = h
  | 1, 0, 2 => -1   -- ⁅f,e⁆ = -h
  | _, _, _ => 0

def kappa : Fin 3 → Fin 3 → ℂ
  | 0, 1 => 1
  | 1, 0 => 1
  | 2, 2 => 2
  | _, _ => 0

abbrev Sl2 : Type := Fin 3 → ℂ

noncomputable def basis (a : Fin 3) : Sl2 := Pi.single a 1

def brComp (x y : Sl2) (c : Fin 3) : ℂ := ∑ a, ∑ b, x a * y b * fabc a b c
\end{leancode}
{\raggedright\noindent File \leanpath{VOA/\allowbreak Examples/\allowbreak AffineSl2.lean}.\par}
\begin{leancode}
noncomputable abbrev AffSl2 : Type := AffineLie Sl2 kappaForm kappaForm_isInvariant

noncomputable abbrev WakModule (k : ℂ) : Type :=
  vacuumSpace kappaForm kappaForm_isInvariant k

noncomputable def J (k : ℂ) (a : Fin 3) (m : ℤ) : Module.End ℂ (WakModule k) :=
  vacuumRho kappaForm kappaForm_isInvariant k (ofLoop kappaForm kappaForm_isInvariant (basis a) m)
\end{leancode}
\medskip
\begin{btheorem}{9.2}[The affine bracket on $V^k(\leightsl)$]
For every $k\in\C$, all colours $a,b$ and all $m,n\in\Z$, as operators on $V^k(\leightsl)$,
\begin{equation*}
  J^a_mJ^b_n - J^b_nJ^a_m = \sum_c f^{ab}{}_c\,J^c_{m+n} + \delta_{m+n,0}\,m\,k\,\kappa(a,b)\,\mathrm{id}.
\end{equation*}
\end{btheorem}
{\raggedright\noindent\textbf{Lean.} \lean{VOA.}\allowbreak\lean{AffineVOA.}\allowbreak\lean{current_comm}\par}
{\raggedright\noindent File \leanpath{VOA/\allowbreak Examples/\allowbreak AffineSl2.lean}.\par}
\begin{leancode}
theorem current_comm (k : ℂ) (a b : Fin 3) (m n : ℤ) :
    J k a m * J k b n - J k b n * J k a m
      = (∑ c, fabc a b c • J k c (m + n))
        + (if m + n = 0 then ((m : ℂ) * k * kappa a b) • LinearMap.id else 0) := by
\end{leancode}
\noindent\textbf{Hypotheses / status.} None\par
\medskip
\begin{btheorem}{9.9}[Translation covariance of the currents]
Let $k\ne-2$. Then $[L_{-1},J^c_\mu] = -\mu\,J^c_{\mu-1}$ for every colour $c$ and $\mu\in\Z$. Equivalently,
$[L_{-1},J^c(z)] = \partial_zJ^c(z)$.
\end{btheorem}
{\raggedright\noindent\textbf{Lean.} \lean{VOA.}\allowbreak\lean{AffineVOA.}\allowbreak\lean{Ltrans},\allowbreak\ \lean{VOA.}\allowbreak\lean{AffineVOA.}\allowbreak\lean{SugawaraCov.}\allowbreak\lean{Lcomm_current},\allowbreak\ \lean{VOA.}\allowbreak\lean{AffineVOA.}\allowbreak\lean{SugawaraCov.}\allowbreak\lean{Jfield_DCovariant},\allowbreak\ \lean{VOA.}\allowbreak\lean{Reconstruction.}\allowbreak\lean{DCovariant}\par}
{\raggedright\noindent File \leanpath{VOA/\allowbreak Examples/\allowbreak AffineSl2Reconstruction.lean}.\par}
\begin{leancode}
noncomputable def Ltrans (k : ℂ) : End ℂ (WakModule k) := Lsug k (-1)
\end{leancode}
{\raggedright\noindent File \leanpath{VOA/\allowbreak Examples/\allowbreak AffineSl2CurrentCovStruct.lean}.\par}
\begin{leancode}
theorem Lcomm_current (k : ℂ) (hk : k + 2 ≠ 0) (c : Fin 3) (μ : ℤ) (x : WakModule k) :
    (Ltrans k * J k c μ - J k c μ * Ltrans k) x = ((-μ : ℤ) : ℂ) • J k c (μ - 1) x := by

theorem Jfield_DCovariant (k : ℂ) (hk : k + 2 ≠ 0) (a : Fin 3) :
    Reconstruction.DCovariant (Ltrans k) (Jfield k a) := by
\end{leancode}
{\raggedright\noindent File \leanpath{VOA/\allowbreak Core/\allowbreak ReconstructionExistence.lean}.\par}
\begin{leancode}
def DCovariant (T : End ℂ V) (a : FormalDist (End ℂ V)) : Prop :=
  ∀ m : ℤ, T * a m - a m * T = ((m + 1 : ℤ)) • a (m + 1)
\end{leancode}
\noindent\textbf{Hypotheses / status.} $k\ne-2$ (necessary)\par
\medskip
\begin{btheorem}{9.10}[The Sugawara $L_0$ is the energy]
Let $k\ne-2$. Then $L_0 = -\ell_0$ on $V^k(\leightsl)$; equivalently $L_0v = nv$ for $v\in V^k_n$.
\end{btheorem}
{\raggedright\noindent\textbf{Lean.} \lean{VOA.}\allowbreak\lean{AffineVOA.}\allowbreak\lean{SugawaraCov.}\allowbreak\lean{Lsug0_comm_current},\allowbreak\ \lean{VOA.}\allowbreak\lean{AffineVOA.}\allowbreak\lean{SugawaraCov.}\allowbreak\lean{Lsug0_eq_negEll0OpA},\allowbreak\ \lean{VOA.}\allowbreak\lean{AffineVOA.}\allowbreak\lean{SugawaraCov.}\allowbreak\lean{hYA_L0}\par}
{\raggedright\noindent File \leanpath{VOA/\allowbreak Examples/\allowbreak AffineSl2L0Cov.lean}.\par}
\begin{leancode}
theorem Lsug0_comm_current (k : ℂ) (hk : k + 2 ≠ 0) (c : Fin 3) (μ : ℤ) (x : WakModule k) :
    (Lsug k 0 * J k c μ - J k c μ * Lsug k 0) x = ((-μ : ℤ) : ℂ) • J k c μ x := by

theorem Lsug0_eq_negEll0OpA (k : ℂ) (hk : k + 2 ≠ 0) : Lsug k 0 = -ell0OpA k := by

theorem hYA_L0 (k : ℂ) (hk : k + 2 ≠ 0) {n : ℤ} (v : WakModule k) (hv : v ∈ gradingZA k n) :
    mode (YA k hk (omegaA k)) 1 v = (n : ℂ) • v := by
\end{leancode}
\noindent\textbf{Hypotheses / status.} $k\ne-2$\par
\medskip
\begin{btheorem}{9.25}[PBW for $V^k(\leightsl)$, by Wakimoto triangularity]
For every $k\ne-2$ the ordered creation monomials form a basis of $V^k(\leightsl)$. Consequently the
universal $C_2$-symbol map $\C[e,f,h]\to R_{V^k(\leightsl)} = V^k/C_2(V^k)$ is injective.
\end{btheorem}
{\raggedright\noindent\textbf{Lean.} \lean{VOA.}\allowbreak\lean{AffineSl2WakimotoLiLeading.}\allowbreak\lean{affineSl2PBWBasis},\allowbreak\ \lean{VOA.}\allowbreak\lean{AffineSl2WakimotoLiLeading.}\allowbreak\lean{universalSymbolAlg_injective}\par}
{\raggedright\noindent File \leanpath{VOA/\allowbreak Examples/\allowbreak AffineSl2WakimotoLiLeading.lean}.\par}
\begin{leancode}
theorem affineSl2PBWBasis (k : ℂ) (hk : k + 2 ≠ 0) :
    PBWBasis AffineVOA.kappaForm AffineVOA.kappaForm_isInvariant
      (Pi.basisFun ℂ (Fin 3)) k :=
  pbwBasis_of_wakimotoOrderedCreationState_linearIndependent k hk
    (wakimotoOrderedCreationState_linearIndependent k)

theorem universalSymbolAlg_injective (k : ℂ) (hk : k + 2 ≠ 0) :
    Function.Injective (Sl2Pres.symbolAlg k hk) :=
  AffineSl2UniversalC2PBW.symbolAlg_injective_of_PBWBasis k hk
    (affineSl2PBWBasis k hk)
\end{leancode}
\noindent\textbf{Hypotheses / status.} $k\ne-2$\par
\medskip
\begin{btheorem}{9.28}[The level is forced; the singular vector]
(i) Every solution of the singularity system has $b=kg$, $c=-(k+2)g$, $d=4g$, $a=-(k-2)(k+2)g$ and
$(3k+4)g=0$; a nonzero solution exists iff $3k+4 = 0$. (ii) At $k=-4/3$ the vector
\begin{equation*}
s = 20\,e_{-3}\leightvac - 12\,e_{-2}h_{-1}\leightvac - 6\,h_{-2}e_{-1}\leightvac + 36\,e_{-1}e_{-1}f_{-1}\leightvac
+ 9\,e_{-1}h_{-1}h_{-1}\leightvac
\end{equation*}
has weight $3$, satisfies $h_0s=2s$, and is killed by every current mode $J^a_m$, $m\ge1$, and by $e_0$.
(iii) $s\ne0$.
\end{btheorem}
{\raggedright\noindent\textbf{Lean.} \lean{VOA.}\allowbreak\lean{Sl2Sing.}\allowbreak\lean{SingSys_param},\allowbreak\ \lean{VOA.}\allowbreak\lean{Sl2Sing.}\allowbreak\lean{exists_nonzero_solution_iff},\allowbreak\ \lean{VOA.}\allowbreak\lean{Sl2Sing.}\allowbreak\lean{sl2Singular},\allowbreak\ \lean{VOA.}\allowbreak\lean{Sl2Sing.}\allowbreak\lean{sl2Singular_E0},\allowbreak\ \lean{VOA.}\allowbreak\lean{Sl2Sing.}\allowbreak\lean{sl2Singular_H0},\allowbreak\ \lean{VOA.}\allowbreak\lean{Sl2Sing.}\allowbreak\lean{sl2Singular_mem},\allowbreak\ \lean{VOA.}\allowbreak\lean{Sl2Sing.}\allowbreak\lean{sl2Singular_pos},\allowbreak\ \lean{VOA.}\allowbreak\lean{Wak.}\allowbreak\lean{Gam3_valW_wSing},\allowbreak\ \lean{VOA.}\allowbreak\lean{Wak.}\allowbreak\lean{sl2Singular_ne_zero}\par}
{\raggedright\noindent File \leanpath{VOA/\allowbreak Examples/\allowbreak AffineSl2Singular.lean}.\par}
\begin{leancode}
theorem SingSys_param {k a b c d g : ℂ} (h : SingSys k a b c d g) :
    b = k * g ∧ c = -(k + 2) * g ∧ d = 4 * g ∧ a = -((k - 2) * (k + 2)) * g
      ∧ (3 * k + 4) * g = 0 := by

theorem exists_nonzero_solution_iff (k : ℂ) :
    (∃ a b c d g : ℂ, ¬(a = 0 ∧ b = 0 ∧ c = 0 ∧ d = 0 ∧ g = 0) ∧ SingSys k a b c d g)
      ↔ 3 * k + 4 = 0 := by

noncomputable def sl2Singular : WakModule (-4 / 3 : ℂ) := ansatz (-4 / 3) 20 (-12) (-6) 36 9

theorem sl2Singular_E0 : J (-4 / 3 : ℂ) colE 0 sl2Singular = 0 :=

theorem sl2Singular_H0 : J (-4 / 3 : ℂ) colH 0 sl2Singular = (2 : ℂ) • sl2Singular :=

theorem sl2Singular_mem : sl2Singular ∈ gradingZA (-4 / 3 : ℂ) 3 := ansatz_mem _ _ _ _ _ _

theorem sl2Singular_pos (a : Fin 3) {m : ℤ} (hm : 1 ≤ m) :
    J (-4 / 3 : ℂ) a m sl2Singular = 0 :=
\end{leancode}
{\raggedright\noindent File \leanpath{VOA/\allowbreak Examples/\allowbreak WakimotoFree.lean}.\par}
\begin{leancode}
theorem Gam3_valW_wSing (k : ℂ) (hk : k + 2 ≠ 0) :
    Gam 3 (valW k hk wSing) = (-8 : ℂ) • (1 : WakFock) := by

theorem sl2Singular_ne_zero : Sl2Sing.sl2Singular ≠ 0 := by
\end{leancode}
\noindent\textbf{Hypotheses / status.} None\par
\medskip
\begin{btheorem}{9.34}[PBW for Verma modules, at every $(k,\lambda)$]
For every $k,\lambda\in\C$: $v_\lambda\ne0$; the \emph{occupation words} --- one fixed product of lower
modes for each finite multiset of lower modes --- form a basis of $M(k,\lambda)$; and for every nonzero
$v\in M(k,\lambda)$ the vectors (occupation word)$\cdot v$ are linearly independent.
\end{btheorem}
{\raggedright\noindent\textbf{Lean.} \lean{VOA.}\allowbreak\lean{AffineVOA.}\allowbreak\lean{HighestModule.}\allowbreak\lean{highestVector_ne_zero},\allowbreak\ \lean{VOA.}\allowbreak\lean{AffineVOA.}\allowbreak\lean{HighestModule.}\allowbreak\lean{occupationWords_independent},\allowbreak\ \lean{VOA.}\allowbreak\lean{AffineVOA.}\allowbreak\lean{HighestModule.}\allowbreak\lean{occupationBasis},\allowbreak\ \lean{VOA.}\allowbreak\lean{Wak.}\allowbreak\lean{Charged.}\allowbreak\lean{lowerLetter},\allowbreak\ \lean{VOA.}\allowbreak\lean{Wak.}\allowbreak\lean{Charged.}\allowbreak\lean{occupationWord},\allowbreak\ \lean{VOA.}\allowbreak\lean{AffineVOA.}\allowbreak\lean{HighestModule.}\allowbreak\lean{descendants_independent}\par}
{\raggedright\noindent File \leanpath{VOA/\allowbreak Examples/\allowbreak AffineSl2Verma.lean}.\par}
\begin{leancode}
theorem highestVector_ne_zero (k h : ℂ) : highestVector k h≠0 := by

theorem occupationWords_independent (k h : ℂ) :
    LinearIndependent ℂ (fun d : Occupation => currentWord (mode k h)
      (occupationWord d) (highestVector k h)) := by
\end{leancode}
{\raggedright\noindent File \leanpath{VOA/\allowbreak Examples/\allowbreak AffineSl2VermaStraightening.lean}.\par}
\begin{leancode}
def occupationBasis (k h : ℂ) : Module.Basis Occupation ℂ (Verma k h) :=
  Module.Basis.mk (occupationWords_independent k h)
    (by change ⊤≤occupationSpan k h; rw [occupationSpan_eq_top])
\end{leancode}
{\raggedright\noindent File \leanpath{VOA/\allowbreak Examples/\allowbreak AffineSl2HighestWords.lean}.\par}
\begin{leancode}
def lowerLetter (v : Fin 3×ℕ) : Fin 3×ℤ :=
  (v.1,if v.1=1 then -(v.2 : ℤ) else -(v.2 : ℤ)-1)

def occupationWord (d : Occupation) : List (Fin 3×ℤ) := d.toMultiset.toList.map lowerLetter
\end{leancode}
{\raggedright\noindent File \leanpath{VOA/\allowbreak Examples/\allowbreak AffineSl2VermaDescendants.lean}.\par}
\begin{leancode}
theorem descendants_independent (k h : ℂ) (v : Verma k h) (hv : v≠0) :
    LinearIndependent ℂ (fun d : Occupation => currentWord (mode k h) (occupationWord d) v) := by
\end{leancode}
\noindent\textbf{Hypotheses / status.} None\par
\medskip
\endgroup
\setcounter{chapter}{9}
\chapter{Graded first-order rigidity of affine \texorpdfstring{$\mathfrak{sl}_2$}{sl(2)}: Linshaw--Qi and admissible levels}
\begingroup
\providecommand{\lninesl}{\mathfrak{sl}_2}
\providecommand{\lnineg}{\mathfrak{g}}
\providecommand{\lnineV}{\mathcal{V}}
\providecommand{\lnineone}{\mathbf{1}}
\providecommand{\lnineDer}{\operatorname{Der}}
\providecommand{\lninead}{\operatorname{ad}}
\providecommand{\lninehat}{\widehat{\mathfrak{sl}}_2}
\noindent{\small 6 of the chapter's 12 numbered statements: its headline results and every definition they use.}\par\medskip
\begin{btheorem}{10.8}[Whitehead's second lemma for $\lninesl$]
If a complex Lie algebra $\lnineg$ has a basis $(e,h,f)$ forming an $\lninesl$-triple, then every
2-cocycle $\omega\in Z^2(\lnineg,\lnineg)$ is $d_{12}\psi$ for some $\psi$; and $Z(\lnineg)=0$.
\end{btheorem}
{\raggedright\noindent\textbf{Lean.} \lean{VOA.}\allowbreak\lean{Sl2Whitehead.}\allowbreak\lean{whitehead2_of_sl2Basis},\allowbreak\ \lean{VOA.}\allowbreak\lean{VertexOperatorAlgebra.}\allowbreak\lean{weightOneH2Vanishes_of_sl2Basis}\par}
\begin{leancode}
theorem whitehead2_of_sl2Basis (b : Basis (Fin 3) ℂ g) (ht : IsSl2Triple (b 1) (b 0) (b 2))
    (ω : twoCochain ℂ g g) (hω : ω ∈ twoCocycle ℂ g g) :
    ∃ ψ : oneCochain ℂ g g, d₁₂ ℂ g g ψ = ω := by
\end{leancode}
\begin{leancode}
theorem weightOneH2Vanishes_of_sl2Basis (W : VertexOperatorAlgebra 𝒱)
    (b : Basis (Fin 3) ℂ ↥(W.grading 1)) (ht : IsSl2Triple (b 1) (b 0) (b 2)) :
    WeightOneH2Vanishes W :=
  fun ω hω => Sl2Whitehead.whitehead2_of_sl2Basis b ht ω hω
\end{leancode}
\noindent File \leanpath{VOA/\allowbreak Core/\allowbreak WeightOneCohomology.lean}.
\noindent\textbf{Hypotheses / status.} $\lninesl$-triple basis\par
\medskip
\begin{btheorem}{10.10}[split form of Linshaw--Qi Theorem 5.1]
For any VOA $W$ with centreless $V_1$: if (Theorem~3.1) every graded regular deformation is trivial,
(Piece~1) $H^2(V_1,V_1) = 0$, and (Piece~2) every graded deformation in pole form has vanishing pole
on $V_1$, then $W$ is graded rigid.
\end{btheorem}
{\raggedright\noindent\textbf{Lean.} \lean{VOA.}\allowbreak\lean{LQ43.}\allowbreak\lean{isRigidGraded_of_split}\par}
\noindent File \leanpath{VOA/Examples/AffineSl2Rigidity43.lean}.
\begin{leancode}
theorem isRigidGraded_of_split {𝒱 : Type*} [AddCommGroup 𝒱] [Module ℂ 𝒱]
    (W : VertexOperatorAlgebra 𝒱) (hZ : W.CentrelessWeightOne)
    (h31 : RegularImpliesTrivialGraded W) (h1 : W.WeightOneH2Vanishes)
    (h2 : NormalisedPoleVanishes W) : Deform.IsRigidGraded W := by
\end{leancode}
\noindent\textbf{Hypotheses / status.} \lean{hZ,h31,h1,h2}\par
\medskip
\begin{btheorem}{10.11}[Graded first-order rigidity of $L_{-4/3}(\lninesl)$; Linshaw--Qi, Theorem 5.1]
Every graded first-order deformation of $L_{-4/3}(\lninesl)$ is trivial.
\end{btheorem}
{\raggedright\noindent\textbf{Lean.} \lean{VOA.}\allowbreak\lean{LQ43.}\allowbreak\lean{L43_graded_rigidity}\par}
\begin{leancode}
theorem L43_graded_rigidity : Deform.IsRigidGraded L43 :=
  L43_isRigidGraded_of_normalisedPoleVanishes L43_normalisedPoleVanishes
\end{leancode}
\noindent File \leanpath{VOA/\allowbreak Examples/\allowbreak AffineSl2RigidityComplete.lean}.
\noindent\textbf{Hypotheses / status.} None\par
\medskip
\begin{btheorem}{10.12}[relation obstruction]
Let $k\ne0,-2$. Suppose $\sum_i c_i J_{u_i}\lnineone = 0$ and $\sum_j d_jJ_{v_j}\lnineone = 0$ in $L_k(\lninesl)$
(two \emph{relations}), and let $P(C) = \sum_{i,j}c_id_j\,G_{u_i,v_j}(C)$ be their Gram polynomial, where
$G_{u,v}(k) = (J_u\lnineone\,|\,J_v\lnineone)$. If $P'(k)\neq0$, then $L_k(\lninesl)$ is graded rigid.
\end{btheorem}
{\raggedright\noindent\textbf{Lean.} \lean{VOA.}\allowbreak\lean{AffineSl2GeneralRigidity.}\allowbreak\lean{graded_rigidity_of_relations}\par}
\noindent File \leanpath{VOA/Examples/AffineSl2UniformObstruction.lean}.
\begin{leancode}
theorem graded_rigidity_of_relations {ι κ : Type*} [Fintype ι] [Fintype κ]
    (u : ι → List (Fin 3 × ℤ)) (c : ι → ℂ)
    (v : κ → List (Fin 3 × ℤ)) (d : κ → ℂ)
    (hc : ∑ i, c i • currentWord k hk (u i) (simple k hk).vacuum = 0)
    (hd : ∑ j, d j • currentWord k hk (v j) (simple k hk).vacuum = 0)
    (hp : (relationPairPolynomial u c v d).derivative.eval k ≠ 0) :
    IsRigidGraded (simple k hk) := by
\end{leancode}
\noindent\textbf{Hypotheses / status.} Relations, $P'(k)\neq0$\par
\medskip
\begin{blemma}{10.13}[kernel pair]
Let $M(X)$ be a square matrix of polynomials over an integral domain, and $z$ a simple root of
$\det M$: $\det M(z) = 0\neq(\det M)'(z)$. Then there are vectors $r,s$ with $r^{\mathsf T}M(z) = 0$,
$M(z)s = 0$ and $r^{\mathsf T}M'(z)s\neq0$.
\end{blemma}
{\raggedright\noindent\textbf{Lean.} \lean{VOA.}\allowbreak\lean{PolynomialMatrix.}\allowbreak\lean{kernel_pair_of_simple_root},\allowbreak\ \lean{VOA.}\allowbreak\lean{AffineSl2GeneralRigidity.}\allowbreak\lean{graded_rigidity_of_vacuum_simple_root}\par}
{\raggedright\noindent File \leanpath{VOA/\allowbreak Core/\allowbreak PolynomialMatrixSimpleRoot.lean}.\par}
\begin{leancode}
theorem kernel_pair_of_simple_root [NoZeroDivisors R]
    (M : Matrix ι ι R[X]) (z : R)
    (hz : M.det.eval z = 0) (hd : M.det.derivative.eval z ≠ 0) :
    ∃ r c : ι → R,
      (∀ j, ∑ i, r i * (M i j).eval z = 0) ∧
      (∀ i, ∑ j, (M i j).eval z * c j = 0) ∧
      (∑ i, ∑ j, (r i * c j) * (M i j).derivative.eval z) ≠ 0 := by
\end{leancode}
\begin{leancode}
theorem graded_rigidity_of_vacuum_simple_root (hk0 : k≠0) (n : ℕ) (c : ℤ)
    (hz : (vacuumChevalleyPolynomial n c).det.eval k=0)
    (hd : (vacuumChevalleyPolynomial n c).det.derivative.eval k≠0) :
    Deform.IsRigidGraded (simple k hk) := by
\end{leancode}
\noindent\textbf{Hypotheses / status.} Simple root\par
\medskip
\begin{btheorem}{10.15}[graded first-order rigidity at negative admissible levels]
Let $p,q\in\N$ with $p\ge2$, $q\ge1$, $\gcd(p,q)=1$ and $p<2q$, and let $k = p/q-2$. Then every graded
first-order deformation of the simple affine vertex operator algebra $L_k(\lninesl)$ is trivial: for
every $Y_1$ satisfying the order-$t$ vertex-algebra axioms and the weight rule, there is a linear
$\varphi:L_k\to L_k$ preserving each weight space, with $\varphi\lnineone = 0$ and $Y_1 = \delta\varphi$.
\end{btheorem}
{\raggedright\noindent\textbf{Lean.} \lean{VOA.}\allowbreak\lean{AffineSl2GeneralRigidity.}\allowbreak\lean{negative_admissible_graded_rigidity}\par}
\begin{leancode}
theorem negative_admissible_graded_rigidity (p q : ℕ) (hp : 2≤p) (hq : 0<q)
    (hcop : p.Coprime q) (hneg : p<2*q) :
    Deform.IsRigidGraded (simple ((p : ℂ)/q-2) (admissible_level_noncritical p q hp hq)) :=
\end{leancode}
\noindent File \leanpath{VOA/\allowbreak Examples/\allowbreak AffineSl2AdmissibleRigidity.lean}.
\noindent\textbf{Hypotheses / status.} $2\le p<2q$, $q\ge1$, coprime\par
\medskip
\endgroup
\setcounter{chapter}{10}
\chapter{Filtrations, the \texorpdfstring{$C_2$}{C2} algebra and presentations}
\begingroup
\providecommand{\lelevenvac}{\mathbf{1}}
\providecommand{\lelevenR}{R_V}
\providecommand{\lelevengr}{\operatorname{gr}}
\providecommand{\lelevenTR}{T(R)}
\providecommand{\lelevenF}{F}
\providecommand{\lelevenJ}{J}
\noindent{\small 8 of the chapter's 58 numbered statements: its headline results and every definition they use.}\par\medskip
\begin{btheorem}{11.8}[Li's multiplicativity estimates]
Let $x\in F^pV$, $y\in F^qV$. Then
\begin{enumerate}
\item for every $N\in\Z$, $x_{(N)}y\in F^{p+q-N-1}V$ (indices $\le0$ meaning $V$);
\item $x_{(-1)}y\in F^{p+q}V$ and, more generally, $x_{(-d-1)}y\in F^{p+q+d}V$ for $d\in\N$;
\item for $n\in\N$ with $n\le p+q$, $x_{(n)}y\in F^{p+q-n}V$;
\item $Tx\in F^{p+1}V$; and $F^0V = V$.
\end{enumerate}
\end{btheorem}
{\raggedright\noindent\textbf{Lean.} \lean{VOA.}\allowbreak\lean{adjointGfil_smul_mem},\allowbreak\ \lean{VOA.}\allowbreak\lean{liFiltration_smul_negSucc_mem},\allowbreak\ \lean{VOA.}\allowbreak\lean{liFiltration_smul_nonneg_mem},\allowbreak\ \lean{VOA.}\allowbreak\lean{liFiltration_L_neg_one_mem},\allowbreak\ \lean{VOA.}\allowbreak\lean{liFiltration_zero_eq_top}\par}
{\raggedright\noindent File \leanpath{VOA/\allowbreak Filtered/\allowbreak LiMultiplicativity.lean}.\par}
\begin{leancode}
theorem adjointGfil_smul_mem (W : VertexOperatorAlgebra V)
    (hLW : LWGenerated (adjointWeakModule W) ({W.vacuum} : Set V))
    {p q N : ℤ} {x y : V}
    (hx : x ∈ Gfil (adjointWeakModule W) ({W.vacuum} : Set V) p)
    (hy : y ∈ Gfil (adjointWeakModule W) ({W.vacuum} : Set V) q) :
    W.smul x N y ∈
      Gfil (adjointWeakModule W) ({W.vacuum} : Set V) (p + q - N - 1) := by

theorem liFiltration_smul_negSucc_mem (W : VertexOperatorAlgebra V)
    {p q : ℕ} (d : ℕ) {x y : V}
    (hx : x ∈ liFiltration W p) (hy : y ∈ liFiltration W q) :
    W.smul x (-(d : ℤ) - 1) y ∈ liFiltration W (p + q + d) := by

theorem liFiltration_smul_nonneg_mem (W : VertexOperatorAlgebra V)
    {p q n : ℕ} (hn : n ≤ p + q) {x y : V}
    (hx : x ∈ liFiltration W p) (hy : y ∈ liFiltration W q) :
    W.smul x (n : ℤ) y ∈ liFiltration W (p + q - n) := by

theorem liFiltration_L_neg_one_mem (W : VertexOperatorAlgebra V) {p : ℕ} {x : V}
    (hx : x ∈ liFiltration W p) : W.L (-1) x ∈ liFiltration W (p + 1) := by

theorem liFiltration_zero_eq_top (W : VertexOperatorAlgebra V) : liFiltration W 0 = ⊤ := by
\end{leancode}
\noindent\textbf{Hypotheses / status.} Unconditional\par
\medskip
\begin{btheorem}{11.10}[Li: $F^1V = C_2(V)$; $\lelevengr^F_0V = \lelevenR$]
$C_2(V) = F^1V$. Consequently the degree-zero layers $\bigoplus_N\lelevengr^F_{0,N}V$ form the algebra
$\lelevenR = V/C_2(V)$, and the class map $V_N\to\lelevengr^F_{0,N}V$ kills $C_2(V)\cap V_N$.
\end{btheorem}
{\raggedright\noindent\textbf{Lean.} \lean{VOA.}\allowbreak\lean{O₂_le_liFiltration_one},\allowbreak\ \lean{VOA.}\allowbreak\lean{liDegreeZeroClass},\allowbreak\ \lean{VOA.}\allowbreak\lean{liDegreeZeroClass_eq_zero_of_mem_O₂},\allowbreak\ \lean{VOA.}\allowbreak\lean{liFiltration_one_le_O₂},\allowbreak\ \lean{VOA.}\allowbreak\lean{liFiltration_one_eq_O₂}\par}
{\raggedright\noindent File \leanpath{VOA/\allowbreak Filtered/\allowbreak ClassicalFreeness.lean} (the inclusion $C_2\subseteq F^1$).\par}
\begin{leancode}
theorem O₂_le_liFiltration_one (W : VertexOperatorAlgebra V) :
    O₂ W ≤ liFiltration W 1 := by

noncomputable def liDegreeZeroClass (W : VertexOperatorAlgebra V) (N : ℕ)
    (x : W.grading (N : ℤ)) : liAssociatedGraded W :=
  DirectSum.of (fun d : ℕ × ℕ => liWeightLayer W d.1 d.2) (0, N)

theorem liDegreeZeroClass_eq_zero_of_mem_O₂ (W : VertexOperatorAlgebra V) (N : ℕ)
    (x : W.grading (N : ℤ)) (hx : x.1 ∈ O₂ W) :
    liDegreeZeroClass W N x = 0 := by
\end{leancode}
{\raggedright\noindent File \leanpath{VOA/\allowbreak Filtered/\allowbreak LiC2Bridge.lean} (the reverse inclusion and the equality, now
library theorems):\par}
\begin{leancode}
theorem liFiltration_one_le_O₂ (W : VertexOperatorAlgebra V) :
    liFiltration W 1 ≤ O₂ W := by

theorem liFiltration_one_eq_O₂ (W : VertexOperatorAlgebra V) :
    liFiltration W 1 = O₂ W :=
\end{leancode}
\noindent\textbf{Hypotheses / status.} None; equality proved in library (\leanpath{LiC2Bridge.lean})\par
\medskip
\begin{btheorem}{11.16}[Generator words hit every Li layer]
Let $a^1,\dots,a^g\in V$ be homogeneous of weights $w_i$, and suppose
\begin{equation*}
V = \operatorname{span}\{a^1,\dots,a^g,\lelevenvac\} + C_2(V).
\end{equation*}
Let $I$ be any ideal of relations vanishing on the degree-zero classes $[a^i]\in\lelevengr^F_0V$. Then for every
$p\le m$ and every $N$, every element of the layer $\lelevengr^F_{p,N}V$ is in the image of the Li jet map with
$x_i = [a^i]$. Moreover the jet monomial $\prod x_{i_k}^{(r_k)}$ maps to the class of the state
$a^{i_1}_{(-r_1-1)}\cdots a^{i_s}_{(-r_s-1)}\lelevenvac$, exactly.
\end{btheorem}
{\raggedright\noindent\textbf{Lean.} \lean{VOA.}\allowbreak\lean{genWord},\allowbreak\ \lean{VOA.}\allowbreak\lean{pairsLevel},\allowbreak\ \lean{VOA.}\allowbreak\lean{dividedTranslateState},\allowbreak\ \lean{VOA.}\allowbreak\lean{prod_dividedIterate_degreeZeroClass_eq_generatorWord},\allowbreak\ \lean{VOA.}\allowbreak\lean{liPresentedJetMap_hits_homogeneous_class_of_generators}\par}
{\raggedright\noindent File \leanpath{VOA/\allowbreak Filtered/\allowbreak LiGeneratorWords.lean}; generator words
\leanpath{VOA/\allowbreak Modularity/\allowbreak OrderedSpan.lean} (Chapter~15).\par}
\begin{leancode}
def genWord {g : ℕ} (gen : Fin g → 𝒱) (wt : Fin g → ℕ) (l : List (Fin g × ℕ)) :
    List (Letter 𝒱) :=
  l.map (fun q => (gen q.1, wt q.1, q.2))

def pairsLevel {g : ℕ} (l : List (Fin g × ℕ)) : ℕ := (l.map Prod.snd).sum

noncomputable def dividedTranslateState (W : VertexOperatorAlgebra V) (r : ℕ) (a : V) : V :=
  ((r.factorial : ℂ)⁻¹) • ((W.toPreVertexAlgebra.T ^ r) a)

theorem prod_dividedIterate_degreeZeroClass_eq_generatorWord
    (W : VertexOperatorAlgebra V)
    {g : ℕ} (gen : Fin g → V) (wt : Fin g → ℕ)
    (hgen : ∀ i, gen i ∈ W.grading (wt i : ℤ)) (l : List (Fin g × ℕ)) :
    (l.map fun q =>
      FiniteJet.dividedIterate (liAssociatedGradedTranslation W) q.2
        (liDegreeZeroClass W (wt q.1) ⟨gen q.1, hgen q.1⟩)).prod =
      liHomogeneousClass W (pairsLevel l) (wordDeg (genWord gen wt l))
        ⟨wordAct (adjointWeakModule W) (genWord gen wt l) W.vacuum,
          generatorWord_mem_liWeightPiece W gen wt hgen l⟩ := by

theorem liPresentedJetMap_hits_homogeneous_class_of_generators
    (W : VertexOperatorAlgebra V) (m : ℕ)
    (I : Ideal (MvPolynomial (Fin g) ℂ))
    (gen : Fin g → V) (wt : Fin g → ℕ)
    (hgen : ∀ i, gen i ∈ W.grading (wt i : ℤ))
    (hI : I ≤ RingHom.ker
      (MvPolynomial.aeval fun i => liDegreeZeroClass W (wt i) ⟨gen i, hgen i⟩).toRingHom)
    (hdec : (⊤ : Submodule ℂ V) =
      Submodule.span ℂ (Set.range gen ∪ {W.vacuum}) ⊔ O₂ W)
    (p N : ℕ) (hp : p ≤ m) :
    ∀ z : liWeightLayer W p N,
      ∃ y : FiniteJet.PresentedJet (R := ℂ) m I,
\end{leancode}
\noindent\textbf{Hypotheses / status.} \lean{hdec}: \textbf{forces $C_2$-cofinite}\par
\medskip
\begin{btheorem}{11.19}[Current words and Li's filtration for $V^k(\mathfrak{sl}_2)$]
Let $k\ne-2$ and let $C_p\subseteq V^k(\mathfrak{sl}_2)$ be the span of the current words
$J^{a_1}_{(-r_1-1)}\cdots J^{a_s}_{(-r_s-1)}\lelevenvac$ with $r_1+\cdots+r_s\ge p$. Then $F^pV^k(\mathfrak{sl}_2) =
C_p$ for all $p$. For any graded ideal $I$, the Li filtration of the quotient $V^k(\mathfrak{sl}_2)/I$ is the image
of that of $V^k(\mathfrak{sl}_2)$, and equals the image of $C_p$.
\end{btheorem}
{\raggedright\noindent\textbf{Lean.} \lean{VOA.}\allowbreak\lean{AffineVOA.}\allowbreak\lean{CurrentLi.}\allowbreak\lean{currentWordPiece},\allowbreak\ \lean{VOA.}\allowbreak\lean{AffineVOA.}\allowbreak\lean{CurrentLi.}\allowbreak\lean{liFiltration_eq_currentWordPiece},\allowbreak\ \lean{VOA.}\allowbreak\lean{AffineVOA.}\allowbreak\lean{CurrentLi.}\allowbreak\lean{quotientLiFiltration_eq_map},\allowbreak\ \lean{VOA.}\allowbreak\lean{AffineVOA.}\allowbreak\lean{CurrentLi.}\allowbreak\lean{quotientLiFiltration_eq_currentWordPiece}\par}
{\raggedright\noindent File \leanpath{VOA/\allowbreak Filtered/\allowbreak AffineCurrentLi.lean}.\par}
\begin{leancode}
def currentWordPiece (k : ℂ) (p : ℕ) : Submodule ℂ (WakModule k) :=
  Submodule.span ℂ
    {x | ∃ w : List (Fin 3 × ℕ), p ≤ currentWordLevel w ∧ wordStateA k w = x}

theorem liFiltration_eq_currentWordPiece (k : ℂ) (hk : k + 2 ≠ 0) (p : ℕ) :
    liFiltration (affineSl2VOA_closed k hk) p = currentWordPiece k p :=

theorem quotientLiFiltration_eq_map (p : ℕ) :
    liFiltration (QuotVOA.quotientVOA W I hI hgr) p =
      (liFiltration W p).map I.mkQ :=

theorem quotientLiFiltration_eq_currentWordPiece (p : ℕ) :
    liFiltration (QuotVOA.quotientVOA (affineSl2VOA_closed k hk) I hI hgr) p =
      quotientCurrentWordPiece k I p :=
\end{leancode}
\noindent\textbf{Hypotheses / status.} $k\neq-2$; unconditional\par
\medskip
\begin{bdefinition}{11.20}[Generator words and the word filtration]
Let $V$ be a VOA and $a^1,\dots,a^g\in V$ a finite family of states, with weight labels $w_1,\dots,w_g\in\N$. Encode
a word by a finite list $l = [(i_1,r_1),\dots,(i_s,r_s)]$ of pairs $(i,r)\in\{1,\dots,g\}\times\N$. Its
\emph{state} is
\begin{equation*}
  a_l := a^{i_1}_{(-r_1-1)}a^{i_2}_{(-r_2-1)}\cdots a^{i_s}_{(-r_s-1)}\lelevenvac ,
\end{equation*}
and its \emph{level} is $|l| := r_1+\cdots+r_s$ (the total number of derivatives, since
$a_{(-r-1)}b = (T^ra/r!)_{(-1)}b$). For $p\in\N$ the \emph{word piece} is
\begin{equation*}
  \mathcal W_p := \operatorname{span}_\C\{\,a_l : |l|\ge p\,\}\subseteq V .
\end{equation*}
The family \emph{spans Li's filtration} if $F^pV\subseteq \mathcal W_p$ for every $p\in\N$.
\end{bdefinition}
{\raggedright\noindent\textbf{Lean.} \lean{VOA.}\allowbreak\lean{genWordState},\allowbreak\ \lean{VOA.}\allowbreak\lean{genWordPiece},\allowbreak\ \lean{VOA.}\allowbreak\lean{GeneratorWordsSpanLiFiltration}\par}
{\raggedright\noindent File \leanpath{VOA/\allowbreak Filtered/\allowbreak GeneratorWordLiSurjectivity.lean}.\par}
\begin{leancode}
abbrev genWordState (l : List (Fin g × ℕ)) : V :=
  wordAct (adjointWeakModule W) (genWord gen wt l) W.vacuum

def genWordPiece (p : ℕ) : Submodule ℂ V :=
  Submodule.span ℂ {u | ∃ l : List (Fin g × ℕ), p ≤ pairsLevel l ∧ genWordState W gen wt l = u}

def GeneratorWordsSpanLiFiltration : Prop :=
  ∀ p : ℕ, liFiltration W p ≤ genWordPiece W gen wt p
\end{leancode}
\noindent\textbf{Hypotheses / status.} Definition; \lean{wt} unused until \lean{hgen}\par
\medskip
\begin{btheorem}{11.23}[Spanning of Li's filtration gives surjective jet maps]
Let $V$ be a VOA, let $a^1,\dots,a^g$ be homogeneous of weights $w_1,\dots,w_g$, and assume that they span Li's filtration
(Definition~11.20). Let $x_i := [a^i]\in\lelevengr^F_{0,w_i}V$ be their degree-zero classes, and let
$I\subseteq\C[x_1,\dots,x_g]$ be any ideal with $f([a^1],\dots,[a^g]) = 0$ for all $f\in I$. Then:
\begin{enumerate}
\item for all $p,N$, the layer $\lelevengr^F_{p,N}V$ is spanned by the classes of the words $a_l$ with $|l| = p$ and
  $\sum_k(w_{i_k}+r_k) = N$;
\item for all $p,N$, every element of $\lelevengr^F_{p,N}V$ is the image of an element of $J_p(\C[x]/I)$ under the
  order-$p$ Li jet map;
\item for all $p,N$, the bidegree jet map $J_p(\C[x]/I)_{(p,N)}\to\lelevengr^F_{p,N}V$ (order $m=p$, base weights $w$) is
  surjective;
\item the infinite Li jet map $J_\infty(\C[x]/I)\to\lelevengr^FV$ is surjective.
\end{enumerate}
\end{btheorem}
{\raggedright\noindent\textbf{Lean.} \lean{VOA.}\allowbreak\lean{genWordClassSet},\allowbreak\ \lean{VOA.}\allowbreak\lean{genWordClassSet_spans_layer},\allowbreak\ \lean{VOA.}\allowbreak\lean{liPresentedJetMap_hits_homogeneous_class_of_genWordSpan},\allowbreak\ \lean{VOA.}\allowbreak\lean{liPresentedJetBidegreeMap_surjective_of_genWordSpan},\allowbreak\ \lean{VOA.}\allowbreak\lean{liInfinitePresentedJetMap_surjective_of_genWordSpan}\par}
{\raggedright\noindent File \leanpath{VOA/\allowbreak Filtered/\allowbreak GeneratorWordLiSurjectivity.lean}.\par}
\begin{leancode}
def genWordClassSet (p N : ℕ) : Set (liAssociatedGraded W) :=
  {x | ∃ l : List (Fin g × ℕ),
      pairsLevel l = p ∧ wordDeg (genWord gen wt l) = N ∧
      x = liHomogeneousClass W (pairsLevel l) (wordDeg (genWord gen wt l))
        ⟨genWordState W gen wt l, generatorWord_mem_liWeightPiece W gen wt hgen l⟩}

theorem genWordClassSet_spans_layer (hspan : GeneratorWordsSpanLiFiltration W gen wt)
    (p N : ℕ) (z : liWeightLayer W p N) :
    DirectSum.of (fun d : ℕ × ℕ => liWeightLayer W d.1 d.2) (p, N) z ∈
      Submodule.span ℂ (genWordClassSet W gen wt hgen p N) := by

theorem liPresentedJetMap_hits_homogeneous_class_of_genWordSpan
    (hspan : GeneratorWordsSpanLiFiltration W gen wt)
    (I : Ideal (MvPolynomial (Fin g) ℂ))
    (hI : I ≤ RingHom.ker
      (MvPolynomial.aeval fun i => liDegreeZeroClass W (wt i) ⟨gen i, hgen i⟩).toRingHom) :
    ∀ p N (z : liWeightLayer W p N),
      ∃ y : FiniteJet.PresentedJet (R := ℂ) p I,
        liPresentedJetMap W p (fun i => liDegreeZeroClass W (wt i) ⟨gen i, hgen i⟩) I hI y =
          DirectSum.of (fun d : ℕ × ℕ => liWeightLayer W d.1 d.2) (p, N) z := by

theorem liPresentedJetBidegreeMap_surjective_of_genWordSpan
    (hspan : GeneratorWordsSpanLiFiltration W gen wt)
    (I : Ideal (MvPolynomial (Fin g) ℂ))
    (hI : I ≤ RingHom.ker
      (MvPolynomial.aeval fun i => liDegreeZeroClass W (wt i) ⟨gen i, hgen i⟩).toRingHom)
    (p N : ℕ) :
    Function.Surjective
      (liPresentedJetBidegreeMap W p (fun i => liDegreeZeroClass W (wt i) ⟨gen i, hgen i⟩)
        I hI wt p N) :=

theorem liInfinitePresentedJetMap_surjective_of_genWordSpan
    (hspan : GeneratorWordsSpanLiFiltration W gen wt)
    (I : Ideal (MvPolynomial (Fin g) ℂ))
    (hI : I ≤ RingHom.ker
      (MvPolynomial.aeval fun i => liDegreeZeroClass W (wt i) ⟨gen i, hgen i⟩).toRingHom) :
    Function.Surjective
      (liInfinitePresentedJetMap W (fun i => liDegreeZeroClass W (wt i) ⟨gen i, hgen i⟩) I hI) := by
\end{leancode}
\noindent\textbf{Hypotheses / status.} \lean{hgen}, \lean{hspan} (= Li spanning), \lean{hI}; unconditional; any VOA, no $C_2$-cofiniteness; comparison with \lean{hdec} on paper only (Rem.~11.24)\par
\medskip
\begin{btheorem}{11.26}[Word spanning plus the tracked estimate give spanning of Li's filtration]
If the words $a_l$ span $V$ and the family satisfies the tracked estimate, then it spans Li's filtration, i.e.\
$F^pV\subseteq \mathcal W_p$ for every $p$.
\end{btheorem}
{\raggedright\noindent\textbf{Lean.} \lean{VOA.}\allowbreak\lean{generatorWordsSpanLiFiltration_of_tracked}\par}
{\raggedright\noindent File \leanpath{VOA/\allowbreak Filtered/\allowbreak GeneratorCurrentLiSpanning.lean}.\par}
\begin{leancode}
theorem generatorWordsSpanLiFiltration_of_tracked (htrack : GeneratorModeTracked W gen wt)
    (hspan : Submodule.span ℂ (Set.range (genWordState W gen wt)) = ⊤) :
    GeneratorWordsSpanLiFiltration W gen wt := by
\end{leancode}
\noindent\textbf{Hypotheses / status.} \lean{htrack}, \lean{hspan} (here = word spanning, $\mathcal W_0=V$); unconditional\par
\medskip
\begin{bdefinition}{11.38}[OPE data: De Sole--Kac Definition 3.1, free-generator form]
An \emph{OPE datum} on generators $\iota$ consists of weights $\Delta_i\ge1$, parities $p_i\in\Z/2$, and elements
$a^i_{(n)}a^j\in\mathcal T(R)$ for $n\in\N$ (the singular OPE table) such that
$a^i_{(n)}a^j\in\mathrm{filt}(\Delta_i+\Delta_j-n-1)$. In particular $a^i_{(n)}a^j = 0$ for $n\ge\Delta_i+\Delta_j$.
\end{bdefinition}
{\raggedright\noindent\textbf{Lean.} \lean{VOA.}\allowbreak\lean{NonlinearConformal.}\allowbreak\lean{OPEData},\allowbreak\ \lean{VOA.}\allowbreak\lean{NonlinearConformal.}\allowbreak\lean{OPEData.}\allowbreak\lean{ope_eq_zero},\allowbreak\ \lean{VOA.}\allowbreak\lean{NonlinearConformal.}\allowbreak\lean{virasoroOPE}\par}
{\raggedright\noindent File \leanpath{VOA/\allowbreak Core/\allowbreak NonlinearConformal.lean}.\par}
\begin{leancode}
structure OPEData (ι : Type u) where
  wt : ι → ℕ
  wt_pos : ∀ i, 0 < wt i
  par : ι → ZMod 2
  ope : ι → ι → ℕ → TR ι
  ope_wt : ∀ i j (n : ℕ), ope i j n ∈ filt wt ((wt i : ℤ) + (wt j : ℤ) - (n : ℤ) - 1)

theorem ope_eq_zero (i j : ι) {n : ℕ} (h : D.wt i + D.wt j ≤ n) : D.ope i j n = 0 := by

noncomputable def virasoroOPE (c : ℂ) : OPEData Unit where
  wt := fun _ => 2
  wt_pos := fun _ => by norm_num
  par := fun _ => 0
  ope := fun _ _ n =>
    if n = 0 then gv ((), 1)
    else if n = 1 then (2 : ℂ) • gv ((), 0)
    else if n = 3 then (c / 2) • (1 : TR Unit)
    else 0
  ope_wt := by
    intro i j n
    split_ifs with h0 h1 h3
    · subst h0
      exact gv_mem_filt _ (by norm_num [wtIdx])
    · subst h1
      exact Submodule.smul_mem _ _ (gv_mem_filt _ (by norm_num [wtIdx]))
    · subst h3
      exact Submodule.smul_mem _ _ (one_mem_filt _ (by norm_num))
    · exact Submodule.zero_mem _
\end{leancode}
\noindent\textbf{Hypotheses / status.} Definition\par
\medskip
\endgroup
\part{Representation theory}
\setcounter{chapter}{11}
\chapter{Modules, Zhu's algebra, and the rationality hierarchy}
\begingroup
\providecommand{\ltwelvevac}{\mathbf{1}}
\providecommand{\ltwelveY}{\mathcal{Y}}
\providecommand{\ltwelveO}{O}
\providecommand{\ltwelveA}{A}
\providecommand{\ltwelvecirc}{\circ}
\providecommand{\ltwelveHP}{\mathrm{HP}_0}
\noindent{\small 18 of the chapter's 43 numbered statements: its headline results and every definition they use.}\par\medskip
\begin{bdefinition}{12.1}[Weak module; Gannon Def.~5.3.1, FHL]
Let $V$ be a VOA. A \emph{weak $V$-module} is a vector space $M$ together with a linear map
$Y_M:V\to\End(M)[[z^{\pm1}]]$, $Y_M(u,z) = \sum_{n\in\Z}u^M_{(n)}z^{-n-1}$, such that
\begin{enumerate}[label=(M\arabic*)]
\item $Y_M(\ltwelvevac,z) = \mathrm{id}_M$;
\item each $Y_M(u,z)$ is a field: for all $u\in V$ and $x\in M$, $u^M_{(n)}x = 0$ for $n\gg0$;
\item the Borcherds (Jacobi) identity holds: for all $u,v\in V$, $x\in M$ and $\ell,m,n\in\Z$,
\begin{equation*}
 \sum_{i\ge0}(-1)^i\binom\ell i\Big(u^M_{(\ell+m-i)}v^M_{(n+i)}x - (-1)^\ell v^M_{(\ell+n-i)}u^M_{(m+i)}x\Big)
 = \sum_{i\ge0}\binom mi(u_{(\ell+i)}v)^M_{(m+n-i)}x .
\end{equation*}
\end{enumerate}
\end{bdefinition}
{\raggedright\noindent\textbf{Lean.} \lean{VOA.}\allowbreak\lean{WeakModule},\allowbreak\ \lean{VOA.}\allowbreak\lean{WeakModule.}\allowbreak\lean{act}\par}
{\raggedright\noindent File \leanpath{VOA/\allowbreak Modules/\allowbreak WeakModule.lean}.\par}
\begin{leancode}
structure WeakModule (W : VertexOperatorAlgebra 𝒱) (M : Type v)
    [AddCommGroup M] [Module ℂ M] where
  YM : 𝒱 →ₗ[ℂ] (End ℂ M)⟦z^±⟧
  vacuum_YM : YM W.vacuum = 1
  truncation : ∀ (u : 𝒱) (x : M), ∃ N : ℤ, ∀ n : ℤ, N ≤ n → mode (YM u) n x = 0
  jacobi : ∀ (u v : 𝒱) (x : M) (ℓ m n : ℤ),
    (∑ᶠ i : ℤ, (negOnePow i * cbz ℓ i) •
        (mode (YM u) (ℓ + m - i) (mode (YM v) (n + i) x)
          - negOnePow ℓ • mode (YM v) (ℓ + n - i) (mode (YM u) (m + i) x)))
      = ∑ᶠ i : ℤ, cbz m i • mode (YM (W.smul u (ℓ + i) v)) (m + n - i) x

def act (A : WeakModule W M) (u : 𝒱) (n : ℤ) (x : M) : M := mode (A.YM u) n x
\end{leancode}
\noindent\textbf{Hypotheses / status.} Definition; \lean{∑ᶠ} finite\par
\medskip
\begin{bdefinition}{12.4}[Admissible module; Zhu 1996, §1.4; Dong--Li--Mason 1998, Def.~2.2]
An \emph{admissible} $V$-module is a weak module $M$ together with a decomposition
$M = \bigoplus_{k\in\N}M(k)$ such that for every homogeneous $a\in V_w$, every $p\in\Z$ and every
$k$,
\begin{equation*}
 a^M_{(p)}M(k)\subseteq M(k + w - p - 1)
\end{equation*}
(with $M(j) = 0$ for $j<0$). It is \emph{irreducible} if $M\neq0$ and its only weak submodules are
$0$ and $M$; \emph{completely reducible} if every weak submodule has a complementary weak submodule.
\end{bdefinition}
{\raggedright\noindent\textbf{Lean.} \lean{VOA.}\allowbreak\lean{AdmissibleModule},\allowbreak\ \lean{VOA.}\allowbreak\lean{AdmissibleModule.}\allowbreak\lean{IsIrreducible},\allowbreak\ \lean{VOA.}\allowbreak\lean{AdmissibleModule.}\allowbreak\lean{IsCompletelyReducible},\allowbreak\ \lean{VOA.}\allowbreak\lean{adjointAdmissibleModule}\par}
{\raggedright\noindent File \leanpath{VOA/\allowbreak Modules/\allowbreak AdmissibleModule.lean}.\par}
\begin{leancode}
structure AdmissibleModule (W : VertexOperatorAlgebra 𝒱) (M : Type v)
    [AddCommGroup M] [Module ℂ M] extends WeakModule W M where
  lev : ℤ → Submodule ℂ M
  lev_eq_bot_of_neg : ∀ k : ℤ, k < 0 → lev k = ⊥
  isInternal : DirectSum.IsInternal lev
  act_lev : ∀ ⦃w p k : ℤ⦄ ⦃a : 𝒱⦄, a ∈ W.grading w → ∀ ⦃x : M⦄, x ∈ lev k →
    mode (YM a) p x ∈ lev (k + (w - p - 1))

def IsIrreducible (N : AdmissibleModule W M) : Prop :=
  Nontrivial M ∧ ∀ P : WeakModule.Submodule' N.toWeakModule, P.carrier = ⊥ ∨ P.carrier = ⊤

def IsCompletelyReducible (N : AdmissibleModule W M) : Prop :=
  ∀ P : WeakModule.Submodule' N.toWeakModule,
    ∃ Q : WeakModule.Submodule' N.toWeakModule, IsCompl P.carrier Q.carrier

def adjointAdmissibleModule (W : VertexOperatorAlgebra 𝒱) : AdmissibleModule W 𝒱 where
  toWeakModule := adjointWeakModule W
  lev := W.grading
\end{leancode}
\noindent\textbf{Hypotheses / status.} Definitions\par
\medskip
\begin{bdefinition}{12.5}[Ordinary module; irreducible; completely reducible; Gannon Def.~5.3.1(c)]
A \emph{$V$-module} (an \emph{ordinary} module) is a weak module $M$ together with a decomposition
$M = \bigoplus_{\alpha\in\C}M_\alpha$ such that $L^M_0 = \omega^M_{(1)}$ acts on $M_\alpha$ as the scalar $\alpha$, each
$M_\alpha$ is finite-dimensional, and the real parts of the weights $\alpha$ with $M_\alpha\neq0$ are bounded below. It is
\emph{irreducible} if $M\neq0$ and its only weak submodules are $0$ and $M$, and \emph{completely reducible} if every weak
submodule has a complementary weak submodule.
\end{bdefinition}
{\raggedright\noindent\textbf{Lean.} \lean{VOA.}\allowbreak\lean{OrdVModule},\allowbreak\ \lean{VOA.}\allowbreak\lean{OrdVModule.}\allowbreak\lean{IsIrreducible},\allowbreak\ \lean{VOA.}\allowbreak\lean{OrdVModule.}\allowbreak\lean{IsCompletelyReducible}\par}
{\raggedright\noindent File \leanpath{VOA/\allowbreak Modules/\allowbreak WeakModule.lean}.\par}
\begin{leancode}
structure OrdVModule (W : VertexOperatorAlgebra 𝒱) (M : Type v)
    [AddCommGroup M] [Module ℂ M] extends WeakModule W M where
  weight : ℂ → Submodule ℂ M
  isInternal : DirectSum.IsInternal weight
  L0_weight : ∀ (α : ℂ) (x : M), x ∈ weight α → mode (YM W.conformalVector) 1 x = α • x
  finrank : ∀ α, FiniteDimensional ℂ (weight α)
  boundedBelow : ∃ B : ℝ, ∀ α : ℂ, weight α ≠ ⊥ → B ≤ α.re

def IsIrreducible (N : OrdVModule W M) : Prop :=
  Nontrivial M ∧ ∀ P : WeakModule.Submodule' N.toWeakModule, P.carrier = ⊥ ∨ P.carrier = ⊤

def IsCompletelyReducible (N : OrdVModule W M) : Prop :=
  ∀ P : WeakModule.Submodule' N.toWeakModule,
    ∃ Q : WeakModule.Submodule' N.toWeakModule, IsCompl P.carrier Q.carrier
\end{leancode}
\noindent\textbf{Hypotheses / status.} Definitions; \textbf{global} lower bound on weights\par
\medskip
\begin{bdefinition}{12.6}[Adjoint module]
Every VOA $V$ is a module over itself, the \emph{adjoint module}, with $Y_V = Y$ and weight spaces $V_\alpha = V_n$ if $\alpha = n\in\Z$ and
$V_\alpha = 0$ if $\alpha\notin\Z$.
\end{bdefinition}
{\raggedright\noindent\textbf{Lean.} \lean{VOA.}\allowbreak\lean{adjWeight},\allowbreak\ \lean{VOA.}\allowbreak\lean{adjointWeakModule},\allowbreak\ \lean{VOA.}\allowbreak\lean{adjointModule}\par}
{\raggedright\noindent File \leanpath{VOA/\allowbreak Modules/\allowbreak WeakModule.lean}.\par}
\begin{leancode}
def adjWeight (W : VertexOperatorAlgebra 𝒱) (α : ℂ) : Submodule ℂ 𝒱 :=
  ⨆ (n : ℤ) (_ : (n : ℂ) = α), W.grading n

def adjointWeakModule (W : VertexOperatorAlgebra 𝒱) : WeakModule W 𝒱 where
  YM := W.Y
  vacuum_YM := W.vacuum_Y
  truncation := fun u x => W.regularity u x
  jacobi := W.jacobi

def adjointModule (W : VertexOperatorAlgebra 𝒱) : OrdVModule W 𝒱 where
  toWeakModule := adjointWeakModule W
  weight := adjWeight W
\end{leancode}
\noindent\textbf{Hypotheses / status.} Unconditional\par
\medskip
\begin{btheorem}{12.16}[Zhu: zero modes act on the lowest-weight space through $A(V)$]
Let $M$ be an ordinary $V$-module with lowest weight $h$. For all $u,v\in V$ and $x\in M_h$,
$o(u*v)x = o(u)o(v)x$. Moreover $o(O(V))M_h = 0$, so $u\mapsto o(u)|_{M_h}$ factors through $A(V) = V/O(V)$. For homogeneous
$a\in V_M$ ($M\in\N$) and all $b\in V$, all the residues $\sum_i\binom{M+n}{i}a_{(i-2-n)}b$, $n\ge0$, lie in $O(V)$.
\end{btheorem}
{\raggedright\noindent\textbf{Lean.} \lean{VOA.}\allowbreak\lean{oMod_zhuMul_lowest_all},\allowbreak\ \lean{VOA.}\allowbreak\lean{rhoLW},\allowbreak\ \lean{VOA.}\allowbreak\lean{O_le_ker_rhoLW},\allowbreak\ \lean{VOA.}\allowbreak\lean{Pres},\allowbreak\ \lean{VOA.}\allowbreak\lean{Pres_mem_O}\par}
{\raggedright\noindent File \leanpath{VOA/\allowbreak Modules/\allowbreak ZhuModule.lean}.\par}
\begin{leancode}
theorem oMod_zhuMul_lowest_all (N : OrdVModule W Mod) {h : ℂ} (hcw : IsLowestWeight N h)
    (u v : 𝒱) {x : Mod} (hx : x ∈ N.weight h) :
    oMod N (zhuMul W u v) x = oMod N u (oMod N v x)

noncomputable def rhoLW (N : OrdVModule W Mod) (h : ℂ) : 𝒱 →ₗ[ℂ] End ℂ (N.weight h) where
  toFun u := weightBlock N (oMod N u) h

theorem O_le_ker_rhoLW (N : OrdVModule W Mod) {h : ℂ} (hcw : IsLowestWeight N h) :
    O W ≤ LinearMap.ker (rhoLW N h)

noncomputable def Pres (W : VertexOperatorAlgebra 𝒱) (M : ℕ) (a b : 𝒱) (n : ℕ) : 𝒱 :=
  ∑ i ∈ Finset.range (M + n + 1), cbz ((M : ℤ) + n) (i : ℤ) • W.smul a ((i : ℤ) - 2 - n) b

theorem Pres_mem_O (W : VertexOperatorAlgebra 𝒱) {M : ℕ} {a : 𝒱}
    (ha : a ∈ W.grading (M : ℤ)) (b : 𝒱) (n : ℕ) : Pres W M a b n ∈ O W
\end{leancode}
\noindent\textbf{Hypotheses / status.} \lean{IsLowestWeight} (\textbf{$M_h$ may be $0$}); unconditional\par
\medskip
\begin{btheorem}{12.17}[Zhu 1996, Theorem~2.1.1]
For every VOA $V$: $O(V)$ is a two-sided ideal for $*$, the product $*$ is associative modulo $O(V)$, and the
class of $\ltwelvevac$ is a two-sided unit. Hence $A(V)$ is an associative unital $\C$-algebra with product
$[u][v] = [u*v]$.
\end{btheorem}
{\raggedright\noindent\textbf{Lean.} \lean{VOA.}\allowbreak\lean{ZhuRingData},\allowbreak\ \lean{VOA.}\allowbreak\lean{zhuRingData_holds},\allowbreak\ \lean{VOA.}\allowbreak\lean{instRingA},\allowbreak\ \lean{VOA.}\allowbreak\lean{instAlgebraA}\par}
{\raggedright\noindent Files \leanpath{VOA/\allowbreak Modules/\allowbreak ZhuRing.lean} (the predicate) and
\leanpath{VOA/\allowbreak Modules/\allowbreak ZhuRingInstance.lean} (the theorem).\par}
\begin{leancode}
structure ZhuRingData (W : VertexOperatorAlgebra 𝒱) : Prop where
  left_ideal : ∀ {u : 𝒱}, u ∈ O W → ∀ v, zhuMul W u v ∈ O W
  right_ideal : ∀ (u : 𝒱) {v : 𝒱}, v ∈ O W → zhuMul W u v ∈ O W
  assoc : ∀ u v w : 𝒱, zhuMul W (zhuMul W u v) w - zhuMul W u (zhuMul W v w) ∈ O W

theorem zhuRingData_holds (W : VertexOperatorAlgebra 𝒱) : ZhuRingData W :=
  zhuRingData_of_cores W
    (fun _ _ _ _ hu ha b => ZhuHR.right_ideal_core W hu ha b)
    (fun _ _ _ _ hu hv w => associator_core W hu hv w)

@[reducible] noncomputable def instRingA (W : VertexOperatorAlgebra 𝒱) : Ring (A W) :=
  AsRing (zhuRingData_holds W)

@[reducible] noncomputable def instAlgebraA (W : VertexOperatorAlgebra 𝒱) :
    @Algebra ℂ (A W) _ (instRingA W).toSemiring :=
  AsAlgebra (zhuRingData_holds W)
\end{leancode}
\noindent\textbf{Hypotheses / status.} None; unconditional\par
\medskip
\begin{btheorem}{12.23}[$R_V$ is a Poisson algebra; Zhu, Li]
$C_2(V)$ is a two-sided ideal for both products $u_{(-1)}v$ and $u_{(0)}v$, and on $R_V$ the product
$[u][v] = [u_{(-1)}v]$ is commutative, associative and unital, the bracket $\{[u],[v]\} = [u_{(0)}v]$ is a Lie
bracket, and the Leibniz rule $\{a,bc\} = \{a,b\}c + b\{a,c\}$ holds.
\end{btheorem}
{\raggedright\noindent\textbf{Lean.} \lean{VOA.}\allowbreak\lean{poissonBracket_leibniz},\allowbreak\ \lean{VOA.}\allowbreak\lean{A₂CommRing},\allowbreak\ \lean{VOA.}\allowbreak\lean{A₂LieRing}\par}
{\raggedright\noindent Files \leanpath{VOA/\allowbreak Modules/\allowbreak C2Cofinite.lean},
\leanpath{VOA/\allowbreak Modules/\allowbreak PoissonDescent.lean}.\par}
\begin{leancode}
theorem poissonBracket_leibniz (W : VertexOperatorAlgebra 𝒱) (u v w : 𝒱) :
    poissonBracket W u (poissonMul W v w) - poissonMul W v (poissonBracket W u w)
      = poissonMul W (poissonBracket W u v) w := by

noncomputable instance A₂CommRing (W : VertexOperatorAlgebra 𝒱) : CommRing (A₂ W) :=

noncomputable instance A₂LieRing (W : VertexOperatorAlgebra 𝒱) : LieRing (A₂ W) :=
\end{leancode}
\noindent\textbf{Hypotheses / status.} Definitions; unconditional\par
\medskip
\begin{bproposition}{12.31}[\lean{Unitary} is unsatisfiable in dimension $\ge2$]
The library's ``unitary module'' structure on an ordinary module $M$ asks for a $\C$-bilinear form $(\cdot|\cdot)$ on $M$
with $(u_{(n)}x|y) = (x|u_{(n)}y)$ for all $u,n,x,y$ and $(x|x)\neq0$ for every $x\neq0$. If $M$ contains two linearly
independent vectors, no such structure exists.
\end{bproposition}
{\raggedright\noindent\textbf{Lean.} \lean{VOA.}\allowbreak\lean{Unitary}\par}
{\raggedright\noindent File \leanpath{VOA/\allowbreak Modules/\allowbreak Contragredient.lean}; the proposition is the library theorem
\lean{VOA.}\allowbreak\lean{old_}\allowbreak\lean{unitary_}\allowbreak\lean{isEmpty} (\leanpath{VOA/\allowbreak Modules/\allowbreak CoreFixUnitary.lean}).\par}
\begin{leancode}
structure Unitary [StarRing ℂ] (N : OrdVModule W M) where
  form : M →ₗ[ℂ] M →ₗ[ℂ] ℂ
  invariant : ∀ (u : 𝒱) (n : ℤ) (x y : M),
    form (N.toWeakModule.act u n x) y = form x (N.toWeakModule.act u n y)
  posDef : ∀ x : M, x ≠ 0 → form x x ≠ 0
\end{leancode}
\medskip
\begin{bdefinition}{12.32}[The library's intertwining operators and intertwiner space]
Let $M,N,P$ be ordinary $V$-modules. \emph{Mode data} of type $\binom PM\,N$ is a family of bilinear maps
$\ltwelveY_s:M\otimes N\to P$, $s\in\C$ (the coefficient of $z^{-s-1}$ of a would-be field $\ltwelveY(u,z)x$ with arbitrary
complex powers of $z$). Its \emph{intertwining defect} is
\begin{equation*}
 \Delta(a,p,s,u,x) = a^P_{(p)}\ltwelveY_s(u)x - \ltwelveY_s(u)a^N_{(p)}x - \sum_{j\ge0}\binom pj\ltwelveY_{s+p-j}(a^M_{(j)}u)x ,
\end{equation*}
for $a\in V$, $p\in\Z$, $s\in\C$, $u\in M$, $x\in N$. The library's \emph{intertwiner space} $\mathrm{Int}(M,N;P)$ is the subspace of mode data
with $\Delta\equiv0$; an \lean{Intertwiner} is an element of it which in addition is lower truncated ($\ltwelveY_s(u)x = 0$ for
$\mathrm{Re}\,s\gg0$); and $\mathrm{fusionCoeff} = \dim\mathrm{Int}(M,N;P)$.
\end{bdefinition}
{\raggedright\noindent\textbf{Lean.} \lean{VOA.}\allowbreak\lean{IntertwinerModes},\allowbreak\ \lean{VOA.}\allowbreak\lean{intertwinerSummand},\allowbreak\ \lean{VOA.}\allowbreak\lean{intertwiningDefect},\allowbreak\ \lean{VOA.}\allowbreak\lean{IntertwinerSpace},\allowbreak\ \lean{VOA.}\allowbreak\lean{Intertwiner}\par}
{\raggedright\noindent File \leanpath{VOA/\allowbreak Modules/\allowbreak Fusion.lean}.\par}
\begin{leancode}
abbrev IntertwinerModes (M N P : Type u) [AddCommGroup M] [Module ℂ M] [AddCommGroup N]
    [Module ℂ N] [AddCommGroup P] [Module ℂ P] : Type u := ℂ → (M →ₗ[ℂ] N →ₗ[ℂ] P)

noncomputable def intertwinerSummand (A : OrdVModule W M) (Y : IntertwinerModes M N P)
    (a : 𝒱) (p : ℤ) (s : ℂ) (u : M) (x : N) (j : ℕ) : P :=
  cbz p j • Y (s + p - j) (A.toWeakModule.act a (j : ℤ) u) x

noncomputable def intertwiningDefect (A : OrdVModule W M) (B : OrdVModule W N) (C : OrdVModule W P) :
    IntertwinerModes M N P →ₗ[ℂ] (𝒱 → ℤ → ℂ → M → N → P) where
  toFun Y := fun a p s u x =>
    C.toWeakModule.act a p (Y s u x) - Y s u (B.toWeakModule.act a p x)
      - ∑ᶠ j : ℕ, intertwinerSummand A Y a p s u x j

noncomputable def IntertwinerSpace (A : OrdVModule W M) (B : OrdVModule W N) (C : OrdVModule W P) :
    Submodule ℂ (IntertwinerModes M N P) :=
  LinearMap.ker (intertwiningDefect A B C)

structure Intertwiner (A : OrdVModule W M) (B : OrdVModule W N) (C : OrdVModule W P) where
  modes : IntertwinerModes M N P
  truncation : ∀ (u : M) (x : N), ∃ s₀ : ℝ, ∀ s : ℂ, s₀ ≤ s.re → modes s u x = 0
  isIntertwiner : modes ∈ IntertwinerSpace A B C
\end{leancode}
\noindent\textbf{Hypotheses / status.} Definitions; \textbf{weaker than FHL} (no iterate, no $L_{-1}$ axiom)\par
\medskip
\begin{btheorem}{12.33}[\lean{fusionCoeff} is identically zero]
For all ordinary modules $A,B,C$, \lean{fusionCoeff A B C = 0}. More precisely, \lean{IntertwinerSpace A B C} is either
the zero space or infinite-dimensional.
\end{btheorem}
{\raggedright\noindent\textbf{Lean.} \lean{VOA.}\allowbreak\lean{fusionCoeff}\par}
\begin{leancode}
noncomputable def fusionCoeff (A : OrdVModule W M) (B : OrdVModule W N) (C : OrdVModule W P) : ℕ :=
  Module.finrank ℂ (IntertwinerSpace A B C)
\end{leancode}
\noindent\textbf{Hypotheses / status.} \textbf{$\equiv0$} (proved in the library)\par
\medskip
\begin{bdefinition}{12.35}[Rational; Zhu 1996, Dong--Li--Mason 1997--98]
$V$ is \emph{(admissibly) rational} if every admissible $V$-module is completely reducible.
\end{bdefinition}
{\raggedright\noindent\textbf{Lean.} \lean{VOA.}\allowbreak\lean{AdmissiblyRational},\allowbreak\ \lean{VOA.}\allowbreak\lean{IsRational}\par}
{\raggedright\noindent File \leanpath{VOA/\allowbreak Modules/\allowbreak AdmissibleModule.lean}.\par}
\begin{leancode}
def AdmissiblyRational (W : VertexOperatorAlgebra 𝒱) : Prop :=
  ∀ {M : Type u} [AddCommGroup M] [Module ℂ M] (N : AdmissibleModule W M),
    N.IsCompletelyReducible

abbrev IsRational (W : VertexOperatorAlgebra 𝒱) : Prop := AdmissiblyRational W
\end{leancode}
\noindent\textbf{Hypotheses / status.} Definitions\par
\medskip
\begin{btheorem}{12.36}[DLM rationality implies \lean{AdmissiblyRational}, under $C_2$]
Let $V$ be $C_2$-cofinite and rational in the sense of Dong--Li--Mason: every admissible module is the internal direct sum
of graded submodules that are minimal among nonzero graded submodules. Then $V$ is \lean{AdmissiblyRational}.
\end{btheorem}
{\raggedright\noindent\textbf{Lean.} \lean{VOA.}\allowbreak\lean{Classification.}\allowbreak\lean{COne.}\allowbreak\lean{DLMRat.}\allowbreak\lean{GradedMinimal},\allowbreak\ \lean{VOA.}\allowbreak\lean{Classification.}\allowbreak\lean{COne.}\allowbreak\lean{DLMRat.}\allowbreak\lean{DLMRational},\allowbreak\ \lean{VOA.}\allowbreak\lean{Classification.}\allowbreak\lean{COne.}\allowbreak\lean{DLMRat.}\allowbreak\lean{admissiblyRational_of_dlmRational}\par}
{\raggedright\noindent File \leanpath{VOA/\allowbreak Classification/\allowbreak COne/\allowbreak DLMRationality.lean}, namespace
\lean{VOA.Classification.COne.DLMRat}.\par}
\begin{leancode}
def GradedMinimal (N : AdmissibleModule W M) (P : WeakModule.Submodule' N.toWeakModule) : Prop :=
  N.IsLevelGraded P ∧ P.carrier ≠ ⊥ ∧
    ∀ Q : WeakModule.Submodule' N.toWeakModule, N.IsLevelGraded Q → Q.carrier ≤ P.carrier →
      Q.carrier = ⊥ ∨ Q.carrier = P.carrier

def DLMRational (W : VertexOperatorAlgebra 𝒱) : Prop :=
  ∀ {M : Type u} [AddCommGroup M] [Module ℂ M] (N : AdmissibleModule W M),
    ∃ (ι : Type u) (P : ι → WeakModule.Submodule' N.toWeakModule),
      iSupIndep (fun i => (P i).carrier) ∧ ⨆ i, (P i).carrier = ⊤ ∧ ∀ i, GradedMinimal N (P i)

theorem admissiblyRational_of_dlmRational (hC2 : C2Cofinite W) (h : DLMRational W) :
    AdmissiblyRational W := by
\end{leancode}
\noindent\textbf{Hypotheses / status.} \lean{hC2}; proved; converse open\par
\medskip
\begin{bdefinition}{12.37}[Gannon-rational and holomorphic; Gannon, Def.~5.3.2]
$V$ is \emph{Gannon-rational} if (1) every ordinary module is completely reducible, (2) $\Phi(V)$ is finite, and (3) every
irreducible admissible module is ordinary (on the same weak module). $V$ is \emph{holomorphic} if it is Gannon-rational
and $|\Phi(V)| = 1$.
\end{bdefinition}
{\raggedright\noindent\textbf{Lean.} \lean{VOA.}\allowbreak\lean{GannonAdmissibleIsOrdinary},\allowbreak\ \lean{VOA.}\allowbreak\lean{GannonRational},\allowbreak\ \lean{VOA.}\allowbreak\lean{Holomorphic}\par}
{\raggedright\noindent File \leanpath{VOA/\allowbreak Modules/\allowbreak Rationality.lean}.\par}
\begin{leancode}
def GannonAdmissibleIsOrdinary (W : VertexOperatorAlgebra 𝒱) : Prop :=
  ∀ {M : Type u} [AddCommGroup M] [Module ℂ M] (N : AdmissibleModule W M),
    N.IsIrreducible → ∃ O : OrdVModule W M, O.toWeakModule = N.toWeakModule

structure GannonRational (W : VertexOperatorAlgebra 𝒱) : Prop where
  completelyReducible : ∀ {M : Type u} [AddCommGroup M] [Module ℂ M] (N : OrdVModule W M),
    OrdVModule.IsCompletelyReducible N
  finitelyManyIrred : Finite (Φ W)
  admissibleIsOrdinary : GannonAdmissibleIsOrdinary W

def Holomorphic (W : VertexOperatorAlgebra 𝒱) : Prop :=
  GannonRational W ∧ Nat.card (Φ W) = 1
\end{leancode}
\medskip
\begin{bdefinition}{12.38}[Regular; Dong--Li--Mason 1997, Def.~3.1]
$V$ is \emph{regular} if every weak module is completely reducible and every irreducible weak module is ordinary.
$\Phi_{\mathrm{adm}}(V)$ is the set of isomorphism classes of irreducible admissible modules (isomorphism of the underlying
weak modules).
\end{bdefinition}
{\raggedright\noindent\textbf{Lean.} \lean{VOA.}\allowbreak\lean{Regular},\allowbreak\ \lean{VOA.}\allowbreak\lean{IrrAdmModule},\allowbreak\ \lean{VOA.}\allowbreak\lean{Φadm}\par}
{\raggedright\noindent File \leanpath{VOA/\allowbreak Modules/\allowbreak RationalityLattice.lean}.\par}
\begin{leancode}
structure Regular (W : VertexOperatorAlgebra 𝒱) : Prop where
  weakCompletelyReducible : ∀ {M : Type u} [AddCommGroup M] [Module ℂ M] (A : WeakModule W M)
    (P : WeakModule.Submodule' A), ∃ Q : WeakModule.Submodule' A, IsCompl P.carrier Q.carrier
  irreducibleWeakIsOrdinary : ∀ {M : Type u} [AddCommGroup M] [Module ℂ M] (A : WeakModule W M),
    (Nontrivial M ∧ ∀ P : WeakModule.Submodule' A, P.carrier = ⊥ ∨ P.carrier = ⊤) →
    ∃ N : OrdVModule W M, N.toWeakModule = A

structure IrrAdmModule (W : VertexOperatorAlgebra 𝒱) where
  carrier : Type u
  [addCommGroup : AddCommGroup carrier]
  [module : Module ℂ carrier]
  isMod : AdmissibleModule W carrier
  irreducible : isMod.IsIrreducible

def Φadm (W : VertexOperatorAlgebra 𝒱) : Type (u + 1) := Quotient (IrrAdmModule.setoid W)
\end{leancode}
\noindent\textbf{Hypotheses / status.} Definitions\par
\medskip
\begin{btheorem}{12.39}[Zhu 1996, Theorem~2.2.3; Dong--Li--Mason: rational $\Rightarrow A(V)$ semisimple]
Let $V$ be a vertex operator algebra of the library. If $V$ is admissibly rational (Definition~12.35), then
Zhu's algebra $A(V)$, with Zhu's product, is a semisimple ring. Equivalently, the predicate \lean{ZhuAdmissibleSemisimple W}
holds for every $V$.
\end{btheorem}
{\raggedright\noindent\textbf{Lean.} \lean{VOA.}\allowbreak\lean{Huang.}\allowbreak\lean{ZhuSemisimple.}\allowbreak\lean{isSemisimpleRing_A_of_admissiblyRational},\allowbreak\ \lean{VOA.}\allowbreak\lean{Huang.}\allowbreak\lean{ZhuSemisimple.}\allowbreak\lean{zhuAdmissibleSemisimple_of_Pr}\par}
{\raggedright\noindent File \leanpath{VOA/\allowbreak Huang/\allowbreak ZhuSemisimple.lean} (Chapter~21), namespace
\lean{VOA.Huang.ZhuSemisimple}. The section variables are \lean{{W : VertexOperatorAlgebra 𝒱}} and
\lean{(Pr : BottomProjection}\allowbreak\lean{ (parabolicRepOfAMod W (regularAMod W)))}, and each theorem is preceded by
\lean{include Pr in}.\par}
\begin{leancode}
theorem isSemisimpleRing_A_of_admissiblyRational (hrat : AdmissiblyRational W) :
    letI := instRingA W; IsSemisimpleRing (A W) := by

theorem zhuAdmissibleSemisimple_of_Pr : ZhuAdmissibleSemisimple W :=
  fun hrat => isSemisimpleRing_A_of_admissiblyRational Pr hrat
\end{leancode}
\noindent\textbf{Hypotheses / status.} \lean{Pr} (from \lean{VOA.pbw}), \lean{AdmissiblyRational}; proved; used by the Zhu headline\par
\medskip
\begin{btheorem}{12.40}[Abe--Buhl--Dong Prop.~5.7, Cor.~5.8, Thm.~4.5; Dong--Li--Mason Thm.~8.1(b),(c); Zhu]
Let $V$ be a vertex operator algebra of the library (of CFT type, Chapter~1).
\begin{enumerate}
\item \emph{(ABD Prop.~5.7)} If $V$ is $C_2$-cofinite, then every weak $V$-module $M$ carries an admissible-module structure
  (Definition~12.4) whose underlying weak module is $M$ itself, with the same vertex operators.
\item \emph{(ABD Cor.~5.8)} If $V$ is $C_2$-cofinite, then every irreducible weak $V$-module $M$ (that is, $M\neq0$ and its only
  weak submodules are $0$ and $M$) carries an ordinary-module structure (Definition~12.5) with the same
  vertex operators.
\item \emph{(ABD Thm.~4.5, hard direction)} If $V$ is $C_2$-cofinite and admissibly rational, then $V$ is regular.
\item \emph{(Zhu)} If $V$ is admissibly rational, then $A(V)$ is a semisimple ring (Theorem~12.39, restated).
\item \emph{(DLM Thm.~8.1(c))} If $V$ is $C_2$-cofinite, then every irreducible admissible $V$-module carries an
  ordinary-module structure with the same vertex operators. No rationality is assumed.
\item \emph{(DLM Thm.~8.1(b))} If $V$ is $C_2$-cofinite and admissibly rational, then $V$ has only finitely many isomorphism
  classes of irreducible admissible modules: $\Phi_{\mathrm{adm}}(V)$ is finite.
\end{enumerate}
\end{btheorem}
{\raggedright\noindent\textbf{Lean.} \lean{VOA.}\allowbreak\lean{Classification.}\allowbreak\lean{COne.}\allowbreak\lean{ABDWeakAdmissible.}\allowbreak\lean{abdWeakIsAdmissible},\allowbreak\ \lean{VOA.}\allowbreak\lean{Classification.}\allowbreak\lean{COne.}\allowbreak\lean{AdmissibleOrdinary.}\allowbreak\lean{irreducible_admissible_isOrdinary},\allowbreak\ \lean{VOA.}\allowbreak\lean{Classification.}\allowbreak\lean{COne.}\allowbreak\lean{AdmissibleOrdinary.}\allowbreak\lean{dlmSimpleAdmissibleOrdinary_of_c2Cofinite},\allowbreak\ \lean{VOA.}\allowbreak\lean{Classification.}\allowbreak\lean{COne.}\allowbreak\lean{LiteratureInputs.}\allowbreak\lean{abdIrreducibleWeakIsOrdinary_holds},\allowbreak\ \lean{VOA.}\allowbreak\lean{Classification.}\allowbreak\lean{COne.}\allowbreak\lean{LiteratureInputs.}\allowbreak\lean{abdRegularOfC2AndRational_holds},\allowbreak\ \lean{VOA.}\allowbreak\lean{Classification.}\allowbreak\lean{COne.}\allowbreak\lean{LiteratureInputs.}\allowbreak\lean{zhuAdmissibleSemisimple_holds},\allowbreak\ \lean{VOA.}\allowbreak\lean{Classification.}\allowbreak\lean{COne.}\allowbreak\lean{LiteratureInputs.}\allowbreak\lean{dlmFinitelyManyIrredAdmissible_of_c2Cofinite}\par}
{\raggedright\noindent File \leanpath{VOA/\allowbreak Classification/\allowbreak COne/\allowbreak ABDWeakAdmissibleComplete.lean}, namespace
\lean{VOA.Classification.COne.ABDWeakAdmissible}, with \lean{W} an implicit section variable (item~1):\par}
\begin{leancode}
theorem abdWeakIsAdmissible : ABDWeakIsAdmissible W :=
  abdWeakIsAdmissible_of_degreeNormalForm_boundedBelow
    weakDegreeNormalForm weakSpectrumBoundedBelow
\end{leancode}
{\raggedright\noindent File \leanpath{VOA/\allowbreak Classification/\allowbreak COne/\allowbreak AdmissibleOrdinaryC2.lean}, namespace
\lean{VOA.Classification.COne.AdmissibleOrdinary}, with \lean{W} implicit (item~5):\par}
\begin{leancode}
theorem irreducible_admissible_isOrdinary (hC2 : C2Cofinite W)
    (N : AdmissibleModule W M) (hirr : N.IsIrreducible) :
    ∃ O : OrdVModule W M, O.toWeakModule = N.toWeakModule := by

theorem dlmSimpleAdmissibleOrdinary_of_c2Cofinite (hC2 : C2Cofinite W) :
    DLMSimpleAdmissibleOrdinary W := by
\end{leancode}
{\raggedright\noindent File \leanpath{VOA/\allowbreak Classification/\allowbreak COne/\allowbreak LiteratureInputsProved.lean}, namespace
\lean{VOA.Classification.COne.LiteratureInputs} (items 2, 3, 4, 6):\par}
\begin{leancode}
variable {𝒱 : Type u} [AddCommGroup 𝒱] [Module ℂ 𝒱] {W : VertexOperatorAlgebra 𝒱}

theorem abdIrreducibleWeakIsOrdinary_holds : ABDIrreducibleWeakIsOrdinary W :=
  abdIrreducibleWeakIsOrdinary_of_weakIsAdmissible abdWeakIsAdmissible

theorem abdRegularOfC2AndRational_holds : ABDRegularOfC2AndRational W :=
  abdRegularOfC2AndRational_of_components abdWeakIsAdmissible
    abdIrreducibleWeakIsOrdinary_holds

theorem zhuAdmissibleSemisimple_holds : ZhuAdmissibleSemisimple W :=
  VOA.Huang.ZhuSemisimple.zhuAdmissibleSemisimple_of_Pr
    (voaPBWBottomProjection (parabolicRepOfAMod W (regularAMod W)))

theorem dlmFinitelyManyIrredAdmissible_of_c2Cofinite (hC2 : C2Cofinite W) :
    DLMFinitelyManyIrredAdmissible W := fun hAR =>
  c2SemisimpleAdmissibleFiniteness_of_bottomProjection
    (voaPBWBottomProjection (parabolicRepOfAMod W (regularAMod W))) hC2
    (VOA.Huang.ZhuSemisimple.isSemisimpleRing_A_of_admissiblyRational
      (voaPBWBottomProjection (parabolicRepOfAMod W (regularAMod W))) hAR)
\end{leancode}
\noindent\textbf{Hypotheses / status.} None; \textbf{proved}, unconditional\par
\medskip
\begin{bdefinition}{12.43}[Strongly finite, strongly rational; Creutzig--Gannon, Dong--Mason, McRae]
$V$ is \emph{strongly finite} if it is simple, self-dual and $C_2$-cofinite (and of CFT type, which is automatic here);
\emph{strongly rational} if moreover it is admissibly rational.
\end{bdefinition}
{\raggedright\noindent\textbf{Lean.} \lean{VOA.}\allowbreak\lean{StronglyFinite},\allowbreak\ \lean{VOA.}\allowbreak\lean{StronglyRational},\allowbreak\ \lean{VOA.}\allowbreak\lean{stronglyFinite_iff},\allowbreak\ \lean{VOA.}\allowbreak\lean{stronglyRational_of_regular},\allowbreak\ \lean{VOA.}\allowbreak\lean{stronglyRational_of_span_vacuum}\par}
{\raggedright\noindent File \leanpath{VOA/\allowbreak Modules/\allowbreak StronglyRational.lean}.\par}
\begin{leancode}
structure StronglyFinite (W : VertexOperatorAlgebra 𝒱) : Prop where
  simple : W.IsSimple
  selfDual : W.SelfDual
  c2Cofinite : C2Cofinite W

structure StronglyRational (W : VertexOperatorAlgebra 𝒱) : Prop extends StronglyFinite W where
  rational : AdmissiblyRational W

theorem stronglyFinite_iff {W : VertexOperatorAlgebra 𝒱} (hsimple : W.IsSimple)
    (hc2 : C2Cofinite W) :
    StronglyFinite W ↔ ∀ v, v ∈ W.grading 1 → W.L 1 v = 0 :=

theorem stronglyRational_of_regular {W : VertexOperatorAlgebra 𝒱} (h : StronglyFinite W)
    (hreg : Regular W) : StronglyRational W where

theorem stronglyRational_of_span_vacuum {W : VertexOperatorAlgebra 𝒱} (hne : W.vacuum ≠ 0)
    (hspan : Submodule.span ℂ ({W.vacuum} : Set 𝒱) = ⊤) : StronglyRational W where
\end{leancode}
\noindent\textbf{Hypotheses / status.} Definitions; satisfiable\par
\medskip
\begin{btheorem}{12.44}[Strongly rational VOAs: regular, finitely many admissible irreducibles, Gannon-rational]
Let $V$ be strongly rational (Definition~12.43). Then
\begin{enumerate}
\item $V$ is regular (Definition~12.38);
\item $V$ satisfies the predicate \lean{DLMFinitelyManyIrredAdmissible}, and $V$ has only finitely many isomorphism classes
  of irreducible admissible modules;
\item $V$ is Gannon-rational (Definition~12.37).
\end{enumerate}
No named literature hypothesis is assumed.
\end{btheorem}
{\raggedright\noindent\textbf{Lean.} \lean{VOA.}\allowbreak\lean{StronglyRational.}\allowbreak\lean{regular_of_abd},\allowbreak\ \lean{VOA.}\allowbreak\lean{StronglyRational.}\allowbreak\lean{gannonRational},\allowbreak\ \lean{VOA.}\allowbreak\lean{StronglyRational.}\allowbreak\lean{regular},\allowbreak\ \lean{VOA.}\allowbreak\lean{StronglyRational.}\allowbreak\lean{dlmFinitelyManyIrredAdmissible},\allowbreak\ \lean{VOA.}\allowbreak\lean{StronglyRational.}\allowbreak\lean{finite_Φadm},\allowbreak\ \lean{VOA.}\allowbreak\lean{StronglyRational.}\allowbreak\lean{gannonRational'}\par}
{\raggedright\noindent File \leanpath{VOA/\allowbreak Modules/\allowbreak StronglyRational.lean} (the older forms, which carry the
literature predicates as hypotheses):\par}
\begin{leancode}
theorem StronglyRational.regular_of_abd {W : VertexOperatorAlgebra 𝒱} (h : StronglyRational W)
    (hABD : ABDRegularOfC2AndRational W) : Regular W :=
  hABD h.c2Cofinite h.rational

theorem gannonRational (h : StronglyRational W) (hFin : DLMFinitelyManyIrredAdmissible W)
    (hOrd : DLMSimpleAdmissibleOrdinary W) : GannonRational W :=
  gannonRational_of_admissiblyRational h.rational hFin hOrd
\end{leancode}
{\raggedright\noindent File \leanpath{VOA/\allowbreak Classification/\allowbreak COne/\allowbreak LiteratureInputsProved.lean} (the
discharged forms; same section variables as in Theorem~12.40):\par}
\begin{leancode}
theorem _root_.VOA.StronglyRational.regular (h : StronglyRational W) : Regular W :=
  h.regular_of_abd abdRegularOfC2AndRational_holds

theorem _root_.VOA.StronglyRational.dlmFinitelyManyIrredAdmissible (h : StronglyRational W) :
    DLMFinitelyManyIrredAdmissible W :=
  dlmFinitelyManyIrredAdmissible_of_c2Cofinite h.c2Cofinite

theorem _root_.VOA.StronglyRational.finite_Φadm (h : StronglyRational W) : Finite (Φadm W) :=
  h.dlmFinitelyManyIrredAdmissible h.rational

theorem _root_.VOA.StronglyRational.gannonRational' (h : StronglyRational W) :
    GannonRational W :=
  h.gannonRational h.dlmFinitelyManyIrredAdmissible
    (dlmSimpleAdmissibleOrdinary_of_stronglyRational h)
\end{leancode}
\noindent\textbf{Hypotheses / status.} \lean{StronglyRational}; proved, no literature hypothesis\par
\medskip
\endgroup
\setcounter{chapter}{12}
\chapter{Zhu's bijection via the tower \texorpdfstring{$A_n(V)$}{A\_n(V)}}
\begingroup
\providecommand{\lthirteenvac}{\mathbf{1}}
\providecommand{\lthirteenVh}{\widehat{V}}
\providecommand{\lthirteenUE}{\mathcal{U}}
\providecommand{\lthirteenMbar}{\overline{M}}
\providecommand{\lthirteenIrr}{\mathrm{Irr}}
\providecommand{\lthirteenO}{O}
\noindent{\small 7 of the chapter's 33 numbered statements: its headline results and every definition they use.}\par\medskip
\begin{btheorem}{13.4}[Zhu's correspondence, uniqueness form]
Let $V$ be any vertex operator algebra of the library and assume the datum \lean{Pr} of
Definition~13.13 for the regular $A(V)$-module. Then $\mathrm{fwd}:\Phi(V)\to\lthirteenIrr A(V)$
is a bijection. Equivalently: if $P,Q$ are irreducible bundled lowest-weight $V$-modules and $P_h\cong Q_{h'}$ as
$A(V)$-modules, then $P\cong Q$. The datum \lean{Pr} is supplied, for every $V$, by the theorem \lean{VOA.pbw} (through \lean{voaPBWBottomProjection}).
\end{btheorem}
{\raggedright\noindent\textbf{Lean.} \lean{VOA.}\allowbreak\lean{zhuBijective},\allowbreak\ \lean{VOA.}\allowbreak\lean{zhuBijective_of_voaPBW},\allowbreak\ \lean{VOA.}\allowbreak\lean{IrredBottomRigidity},\allowbreak\ \lean{VOA.}\allowbreak\lean{forwardMapIrr_surjective_uncond},\allowbreak\ \lean{VOA.}\allowbreak\lean{forwardMapIrr_bijective_iff_rigidity}\par}
{\raggedright\noindent Files \leanpath{VOA/\allowbreak Modules/\allowbreak ZhuVermaAnArchLUBijection.lean} and
\leanpath{VOA/\allowbreak Modules/\allowbreak ZhuVermaAnArchLUBijectionPBW.lean}.\par}
\begin{leancode}
theorem zhuBijective (Pr : BottomProjection (parabolicRepOfAMod W (regularAMod W))) :
    Function.Bijective (forwardMapIrr W) :=
theorem zhuBijective_of_voaPBW (W : VertexOperatorAlgebra 𝒱) :
    Function.Bijective (forwardMapIrr W) :=
  zhuBijective
    (voaPBWBottomProjection (parabolicRepOfAMod W (regularAMod W)))
\end{leancode}
\begin{leancode}
def IrredBottomRigidity (W : VertexOperatorAlgebra 𝒱) : Prop :=
  ∀ P Q : LWModule W, P.IsIrreducibleLW → Q.IsIrreducibleLW →
    Nonempty (AModEquiv W P.toAMod Q.toAMod) → LWModule.Iso W P Q
theorem forwardMapIrr_surjective_uncond (W : VertexOperatorAlgebra 𝒱) :
    Function.Surjective (forwardMapIrr W) := by
theorem forwardMapIrr_bijective_iff_rigidity :
    Function.Bijective (forwardMapIrr W) ↔ IrredBottomRigidity W :=
  ⟨fun hb => irredBottomRigidity_of_injective hb.1, forwardMapIrr_bijective_of_rigidity⟩
\end{leancode}
\noindent\textbf{Hypotheses / status.} \lean{Pr}; conditional\par
\medskip
\begin{btheorem}{13.6}[Bottoms of irreducible modules are simple; Zhu 1996, Theorem~2.2.1(1)]
For every vertex operator algebra $V$ of the library and every irreducible bundled lowest-weight module $P$, the bottom
$P_h$ is a simple $A(V)$-module. More precisely, every lowest-weight module $P$ has the \emph{bottom-recovery property}: for
every $A(V)$-submodule $S\subseteq P_h$, the $V$-submodule of $P$ generated by $S$ meets $P_h$ in exactly $S$.
\end{btheorem}
{\raggedright\noindent\textbf{Lean.} \lean{VOA.}\allowbreak\lean{BottomRecovers},\allowbreak\ \lean{VOA.}\allowbreak\lean{AllBottomRecovers},\allowbreak\ \lean{VOA.}\allowbreak\lean{Huang.}\allowbreak\lean{LWBottomRecovery.}\allowbreak\lean{allBottomRecovers},\allowbreak\ \lean{VOA.}\allowbreak\lean{Huang.}\allowbreak\lean{LWBottomRecovery.}\allowbreak\lean{isSimpleAMod_toAMod_of_irreducible}\par}
\begin{leancode}
def BottomRecovers : Prop :=
  ∀ S : AModSubmodule (P.toAMod),
    vSubmoduleBottomCarrier S.genVSubmodule ≤ S.carrier

def AllBottomRecovers (W : VertexOperatorAlgebra 𝒱) : Prop :=
  ∀ P : LWModule W, BottomRecovers P
\end{leancode}
\begin{leancode}
theorem allBottomRecovers (W : VertexOperatorAlgebra 𝒱) : AllBottomRecovers W :=
  fun P => bottomRecovers P

theorem isSimpleAMod_toAMod_of_irreducible (hP : P.IsIrreducibleLW) : IsSimpleAMod P.toAMod :=
  isSimpleAMod_of_bottomRecovers P hP (bottomRecovers P)
\end{leancode}
\noindent Files \leanpath{VOA/\allowbreak Modules/\allowbreak ZhuVermaAnArchLeaf3.lean}, \leanpath{VOA/\allowbreak Modules/\allowbreak ZhuVermaAnArchAModScope.lean}, \leanpath{VOA/\allowbreak Huang/\allowbreak LWBottomRecovery.lean}.
\noindent\textbf{Hypotheses / status.} None; unconditional\par
\medskip
\begin{btheorem}{13.8}[Zhu's correspondence as one bijection; Zhu 1996, Theorem~2.2.2 / Gannon, Theorem~5.3.4, for $C_2$-cofinite $V$]
Let $V$ be $C_2$-cofinite. There is a bijection
\begin{equation*}
  e:\Phi(V)\xrightarrow{\ \sim\ }\mathrm{Simple}A(V)/{\cong}
\end{equation*}
such that for every irreducible lowest-weight module $P$ with bottom $P_h$, $e([P])=[P_h]$.
\end{btheorem}
{\raggedright\noindent\textbf{Lean.} \lean{VOA.}\allowbreak\lean{zhuCorrespondence_of_C2Cofinite}\par}
{\raggedright\noindent File \leanpath{VOA/\allowbreak Modularity/\allowbreak ZhuModularityCanonical.lean}.\par}
\begin{leancode}
theorem zhuCorrespondence_of_C2Cofinite (W : VertexOperatorAlgebra 𝒱) (hC2 : C2Cofinite W) :
    ∃ e : Φ W ≃ SimpleAModIsoClass W, ∀ P : IrrLWModule W,
      e (forgetClass W (Quotient.mk (IrrLWModule.setoid W) P)) =
        Quotient.mk (SimpleAMod.setoid W) (IrrLWModule.toSimpleAMod P) := by
\end{leancode}
\noindent\textbf{Hypotheses / status.} $C_2$; conditional; via theorem \lean{VOA.pbw}\par
\medskip
\begin{bdefinition}{13.13}[Bottom projection; DLM's PBW step]
A \emph{bottom projection} for a level-$n$ parabolic representation $P$ on $U$ is a linear map $\pi:M_n(U)\to U$ with
$\pi([1\otimes u]) = u$ for all $u$, and $\pi = 0$ on every level $M_n(U)(k)$, $k\neq n$. The \emph{PBW freeness} of $P$ is
the conjunction of: the bottom map is injective, and the levels are independent.
\end{bdefinition}
{\raggedright\noindent\textbf{Lean.} \lean{VOA.}\allowbreak\lean{VHat.}\allowbreak\lean{DLM3.}\allowbreak\lean{BottomProjection},\allowbreak\ \lean{VOA.}\allowbreak\lean{VHat.}\allowbreak\lean{DLM3.}\allowbreak\lean{PBWFreeness},\allowbreak\ \lean{VOA.}\allowbreak\lean{VHat.}\allowbreak\lean{DLM3.}\allowbreak\lean{pbwFreeness_iff_bottom_injective},\allowbreak\ \lean{VOA.}\allowbreak\lean{VHat.}\allowbreak\lean{DLM3.}\allowbreak\lean{bottomProjectionOfInjective},\allowbreak\ \lean{VOA.}\allowbreak\lean{VHat.}\allowbreak\lean{DLM3.}\allowbreak\lean{nonempty_bottomProjection_iff_bottom_injective}\par}
{\raggedright\noindent Files \leanpath{VOA/\allowbreak Modules/\allowbreak ZhuVermaAnArchDLM3Survival.lean} and
\leanpath{VOA/\allowbreak Modules/\allowbreak ZhuVermaAnArchDLM3PBWSplit.lean}.\par}
\begin{leancode}
structure BottomProjection (P : ParabolicRep W n U) where
  /-- The projection onto the bottom `U` of `M_n(U)`. -/
  proj : IndMod P →ₗ[ℂ] U
  /-- It retracts the bottom map `u ↦ [1 ⊗ u]`. -/
  proj_bottom : ∀ u : U, proj (indBottom P u) = u
  /-- It annihilates every level other than `n` (DLM3: "`U*` annihilates `⨁_{i≠n} M(U)(i)`"). -/
  proj_levCover : ∀ ⦃k : ℤ⦄, k ≠ (n : ℤ) → ∀ ⦃z : IndMod P⦄, z ∈ levCover P k → proj z = 0
structure PBWFreeness (P : ParabolicRep W n U) : Prop where
  /-- **PBW, part 1** — the bottom `U` embeds: `u ↦ [1 ⊗ u]` is injective. -/
  bottom_injective : Function.Injective (indBottom P)
  /-- **Historical PBW part 2, now a theorem** — the levels `M_n(U)(k) = U(V̂)_{k−n}U` are
  independent. `levCover_iSupIndep` supplies this field unconditionally; together with
  `levCover_spanning` it gives the internal `ℤ`-grading of `M_n(U)`. -/
  levels_iSupIndep : iSupIndep (levCover P)
\end{leancode}
\begin{leancode}
theorem pbwFreeness_iff_bottom_injective (P : ParabolicRep W n U) :
    PBWFreeness P ↔ Function.Injective (indBottom P) :=
noncomputable def bottomProjectionOfInjective (Q : ParabolicRep W 0 U)
    (hinj : Function.Injective (indBottom Q)) : BottomProjection Q :=
theorem nonempty_bottomProjection_iff_bottom_injective (Q : ParabolicRep W 0 U) :
    Nonempty (BottomProjection Q) ↔ Function.Injective (indBottom Q) := by
\end{leancode}
\noindent\textbf{Hypotheses / status.} Definitions; $\iff$ bottom injective\par
\medskip
\begin{btheorem}{13.25}[Survival of the bottom; the gate]
Let \lean{Pr} be a bottom projection for the parabolic representation of the regular $A(V)$-module. Then:
\begin{enumerate}
\item (Dong--Li--Mason, Theorem~4.1 at $n=0$.) For every $A(V)$-module $U$ with a bottom projection, $U$ embeds into the
  weak module $\lthirteenMbar_0(U) = M_0(U)/(\text{defects})$ as its bottom.
\item (Dong--Jiang, Proposition~4.6 at $(0,0)$.) $O_{0,0}(V) = O(V)$; hence $A_{0,0}(V) = A(V)$, and for every $A(V)$-module $U$
  the bottom level of $M(U)$ is $M(U)(0)\cong U$.
\end{enumerate}
\end{btheorem}
{\raggedright\noindent\textbf{Lean.} \lean{VOA.}\allowbreak\lean{VHat.}\allowbreak\lean{DLM3.}\allowbreak\lean{IndBar},\allowbreak\ \lean{VOA.}\allowbreak\lean{VHat.}\allowbreak\lean{DLM3.}\allowbreak\lean{indBarBottom_injective_of_bottomProjection},\allowbreak\ \lean{VOA.}\allowbreak\lean{gate_of_pbwDatum},\allowbreak\ \lean{VOA.}\allowbreak\lean{djAnnEqAn_zero_of_pbwDatum},\allowbreak\ \lean{VOA.}\allowbreak\lean{djCarrierDegZeroEquivOfPbwDatum}\par}
{\raggedright\noindent Files \leanpath{VOA/\allowbreak Modules/\allowbreak ZhuVermaAnArchDLM3Survive.lean} and
\leanpath{VOA/\allowbreak Modules/\allowbreak ZhuVermaAnArchDLM3Induced.lean}.\par}
\begin{leancode}
abbrev IndBar (P : ParabolicRep W n U) : Type u :=
  IndMod P ⧸ VHatRep.defectSubmodule (indRho P)
theorem indBarBottom_injective_of_bottomProjection (W : VertexOperatorAlgebra 𝒱) (U : AMod W)
    (Pr : BottomProjection (parabolicRepOfAMod W U)) :
    Function.Injective (indBarBottom (parabolicRepOfAMod W U)) :=
theorem gate_of_pbwDatum (Pr : BottomProjection (parabolicRepOfAMod W (regularAMod W))) :
    OnmFull W 0 0 = O W :=
theorem djAnnEqAn_zero_of_pbwDatum
    (Pr : BottomProjection (parabolicRepOfAMod W (regularAMod W))) : DJAnnEqAn W 0 :=
noncomputable def djCarrierDegZeroEquivOfPbwDatum
    (Pr : BottomProjection (parabolicRepOfAMod W (regularAMod W))) (V : AmMod W 0) :
    DJCarrier.deg W V 0 ≃ₗ[ℂ] V.carrier :=
\end{leancode}
\noindent\textbf{Hypotheses / status.} \lean{Pr}\par
\medskip
\begin{btheorem}{13.31}[$P\cong L(\Omega(P))$; Dong--Jiang, Theorem~4.13]
Every irreducible bundled lowest-weight module $P$ is isomorphic to $L(\Omega(P))$.
\end{btheorem}
{\raggedright\noindent\textbf{Lean.} \lean{VOA.}\allowbreak\lean{pRealizesLWIso}\par}
{\raggedright\noindent File \leanpath{VOA/\allowbreak Modules/\allowbreak ZhuVermaAnArchLURealizesIrr.lean}.\par}
\begin{leancode}
theorem pRealizesLWIso (P : LWModule W)
    (Pr : BottomProjection (parabolicRepOfAMod W (regularAMod W)))
    (hU : Nontrivial (pBotAmMod P).carrier)
    (hL0 : ∀ (n : ℕ) (x : DJCarrier W (pBotAmMod P)), x ∈ DJCarrier.deg W (pBotAmMod P) n →
      djModeZero W (pBotAmMod P) 1 W.conformalVector x = (P.h + n) • x)
    (hP : P.IsIrreducibleLW)
    (hfin : (djLU Pr (pBotAmMod P) hU P.h hL0).FiniteWeights) :
    LWModule.Iso W ((djLU Pr (pBotAmMod P) hU P.h hL0).toLWModule hfin) P :=
\end{leancode}
\noindent\textbf{Hypotheses / status.} \lean{Pr}, irreducible\par
\medskip
\begin{btheorem}{13.32}[$L(\cdot)$ is invariant under isomorphism of bottoms]
If $U_1\cong U_2$ as $A(V)$-modules (both nonzero, with $L_0$ acting on level $n$ by $h_i+n$, $h_1 = h_2$), then
$L(U_1)\cong L(U_2)$.
\end{btheorem}
{\raggedright\noindent\textbf{Lean.} \lean{VOA.}\allowbreak\lean{djFunctorEquiv},\allowbreak\ \lean{VOA.}\allowbreak\lean{djLUFunctorLWIso'},\allowbreak\ \lean{VOA.}\allowbreak\lean{pH_eq}\par}
{\raggedright\noindent Files \leanpath{VOA/\allowbreak Modules/\allowbreak ZhuVermaAnArchLUFunctorial.lean} and
\leanpath{VOA/\allowbreak Modules/\allowbreak ZhuVermaAnArchLUBijection.lean}.\par}
\begin{leancode}
noncomputable def djFunctorEquiv (U₁ U₂ : AmMod W 0) (e : AmModHom U₁ U₂) :
    DJCarrier W U₁ ≃ₗ[ℂ] DJCarrier W U₂ :=
  DirectSum.congrLinearEquiv (fun n => djFunctorPiece U₁ U₂ e n)
theorem djLUFunctorLWIso' (Pr : BottomProjection (parabolicRepOfAMod W (regularAMod W)))
    {U₁ U₂ : AmMod W 0} (e : AmModHom U₁ U₂) (hU₁ : Nontrivial U₁.carrier)
    (hU₂ : Nontrivial U₂.carrier) (h₁ h₂ : ℂ) (heq : h₁ = h₂)
    (hL0₁ : ∀ (n : ℕ) (x : DJCarrier W U₁), x ∈ DJCarrier.deg W U₁ n →
      djModeZero W U₁ 1 W.conformalVector x = (h₁ + n) • x)
    (hL0₂ : ∀ (n : ℕ) (x : DJCarrier W U₂), x ∈ DJCarrier.deg W U₂ n →
      djModeZero W U₂ 1 W.conformalVector x = (h₂ + n) • x)
    (hfin₁ : (djLU Pr U₁ hU₁ h₁ hL0₁).FiniteWeights)
    (hfin₂ : (djLU Pr U₂ hU₂ h₂ hL0₂).FiniteWeights) :
    LWModule.Iso W ((djLU Pr U₁ hU₁ h₁ hL0₁).toLWModule hfin₁)
      ((djLU Pr U₂ hU₂ h₂ hL0₂).toLWModule hfin₂) := by
theorem pH_eq {P Q : LWModule W} (ψ : AModEquiv W P.toAMod Q.toAMod) : P.h = Q.h := by
\end{leancode}
\noindent\textbf{Hypotheses / status.} \lean{Pr}\par
\medskip
\endgroup
\part{Modularity}
\setcounter{chapter}{13}
\chapter{Modular forms and the analytic toolkit}
\begingroup
\providecommand{\lfourteenH}{\mathbb{H}}
\providecommand{\lfourteenSL}{\mathrm{SL}_2(\mathbb{Z})}
\providecommand{\lfourteenGt}{\Gamma_\theta}
\providecommand{\lfourteenE}{\mathcal{E}_{46}}
\providecommand{\lfourteenGS}{\mathcal{G}}
\providecommand{\lfourteenev}{\mathrm{ev}}
\providecommand{\lfourteenIm}{\operatorname{Im}}
\providecommand{\lfourteenRe}{\operatorname{Re}}
\noindent{\small 9 of the chapter's 68 numbered statements: its headline results and every definition they use.}\par\medskip
\begin{btheorem}{14.2}[The derivative does not commute with slashing]
Let $k\in\Z$, $\gamma\in\lfourteenSL$, $z\in\lfourteenH$, and let $f$ be complex-differentiable at $\gamma z$. Then
\begin{equation*}
  (f')|_{k+2}\gamma\,(z) \;=\; \frac{d}{dz}\bigl(f|_k\gamma\bigr)(z) \;+\; \frac{k\,c}{cz+d}\,(f|_k\gamma)(z).
\end{equation*}
\end{btheorem}
{\raggedright\noindent\textbf{Lean.} \lean{VOA.}\allowbreak\lean{slashC_deriv_eq}\par}
{\raggedright\noindent File \leanpath{VOA/\allowbreak Modularity/\allowbreak SerreDerivative.lean}.\par}
\begin{leancode}
theorem slashC_deriv_eq (k : ℤ) (γ : SL(2, ℤ)) (f : ℂ → ℂ) {z : ℂ} (hz : z ∈ imGT 0)
    (hf : DifferentiableAt ℂ f (mobiusC γ z)) :
    slashC (k + 2) γ (_root_.deriv f) z
      = _root_.deriv (slashC k γ f) z
        + (k : ℂ) * (γ 1 0 : ℂ) * (denomC γ z)⁻¹ * slashC k γ f z
\end{leancode}
\noindent\textbf{Hypotheses / status.} Differentiability; unconditional\par
\medskip
\begin{btheorem}{14.4}[Quasimodularity of $G_2$; Mathlib's \lean{G2_slash_action}]
For $\gamma\in\lfourteenSL$ and $z\in\lfourteenH$: $(G_2|_2\gamma)(z) = G_2(z) - \dfrac{2\pi i\,c}{cz+d}$.
\end{btheorem}
{\raggedright\noindent\textbf{Lean.} \lean{VOA.}\allowbreak\lean{slashC_G2C}\par}
{\raggedright\noindent File \leanpath{VOA/\allowbreak Modularity/\allowbreak SerreDerivative.lean}.\par}
\begin{leancode}
theorem slashC_G2C (γ : SL(2, ℤ)) {z : ℂ} (hz : z ∈ imGT 0) :
    slashC 2 γ G2C z
      = G2C z - 2 * (Real.pi : ℂ) * Complex.I * (γ 1 0 : ℂ) * (denomC γ z)⁻¹
\end{leancode}
\noindent\textbf{Hypotheses / status.} Differentiability; unconditional\par
\medskip
\begin{btheorem}{14.8}[Identification with Mathlib's Eisenstein series]
For $j\ge2$ and $z\in\lfourteenH$, $\mathcal G_{2j}(z) = g_j(0)\,E_{2j}(z)$, where $E_{2j}$ is Mathlib's normalised Eisenstein
modular form of weight $2j$. Consequently, for all $\gamma\in\lfourteenSL$, $\mathcal G_{2j}|_{2j}\gamma = \mathcal G_{2j}$ on
$\lfourteenH$. (Note $g_j(0) = -B_{2j}/(2j)!\neq1$: Zhu's $\mathcal G_{2j}$ is \emph{not} $E_{2j}$, but a nonzero constant multiple of it.)
\end{btheorem}
{\raggedright\noindent\textbf{Lean.} \lean{VOA.}\allowbreak\lean{zhuGC_eq_mul_E},\allowbreak\ \lean{VOA.}\allowbreak\lean{slashC_zhuGC}\par}
{\raggedright\noindent File \leanpath{VOA/\allowbreak Modularity/\allowbreak EisensteinAnalytic.lean}.\par}
\begin{leancode}
theorem zhuGC_eq_mul_E {j : ℕ} (hj : 2 ≤ j) {z : ℂ} (hz : z ∈ imGT 0) :
    zhuGC j z = zhuG j 0 * ModularForm.E (show 3 ≤ 2 * j by omega) (⟨z, hz⟩ : UpperHalfPlane)

theorem slashC_zhuGC {j : ℕ} (hj : 2 ≤ j) (γ : SL(2, ℤ)) {z : ℂ} (hz : z ∈ imGT 0) :
    slashC (2 * j : ℤ) γ (zhuGC j) z = zhuGC j z
\end{leancode}
\noindent\textbf{Hypotheses / status.} $j\ge2$; unconditional\par
\medskip
\begin{bdefinition}{14.14}[Three named propositions: homogeneous representatives, monomial and full Serre closure]
(a) \emph{Homogeneous representatives}: every $f\in\lfourteenE$ of weight $k$ has a homogeneous polynomial representative of
weight $k$. (b) \emph{Monomial Serre closure}: for all $a,b$, $\partial_{4a+6b}(E_4^aE_6^b)$ equals on $\lfourteenH$ an element of
$\lfourteenE$ of weight $4a+6b+2$. (c) \emph{Serre closure}: for every $f$ of weight $k$ there is $g\in\lfourteenE$ of weight $k+2$
with $g = \partial_kf$ on $\lfourteenH$.
\end{bdefinition}
{\raggedright\noindent\textbf{Lean.} \lean{VOA.}\allowbreak\lean{E46HasWeightedPolynomialReps},\allowbreak\ \lean{VOA.}\allowbreak\lean{E46SerreMonomialClosed},\allowbreak\ \lean{VOA.}\allowbreak\lean{E46SerreClosed}\par}
{\raggedright\noindent File \leanpath{VOA/\allowbreak Modularity/\allowbreak EisensteinRamanujan.lean}.\par}
\begin{leancode}
def E46HasWeightedPolynomialReps : Prop :=
  ∀ (k : ℤ) (f : E46), E46Weight k f → E46WeightedPolynomialRep k f

def E46SerreMonomialClosed : Prop :=
  ∀ (a b : ℕ), ∃ g : E46,
    E46Weight (4 * (a : ℤ) + 6 * (b : ℤ) + 2) g ∧
      ∀ z ∈ imGT 0,
        e46C g z = serreDerivC (4 * (a : ℤ) + 6 * (b : ℤ))
          (e46C (e46Monomial a b)) z

def E46SerreClosed : Prop :=
  ∀ (k : ℤ) (f : E46), E46Weight k f →
    ∃ g : E46, E46Weight (k + 2) g ∧
      ∀ z ∈ imGT 0, e46C g z = serreDerivC k (e46C f) z
\end{leancode}
\medskip
\begin{btheorem}{14.15}[Homogeneous representatives exist]
Every $f\in\lfourteenE$ of weight $k$ is a linear combination of monomials $E_4^aE_6^b$ with $4a+6b=k$; that is,
\lean{E46HasWeightedPolynomialReps} holds.
\end{btheorem}
{\raggedright\noindent File \leanpath{VOA/\allowbreak Modularity/\allowbreak EisensteinSerreClosure.lean}.\par}
\noindent\emph{(The chapter quotes no Lean declaration for this statement; see book p.~684.)}\par
\medskip
\begin{btheorem}{14.23}[Holomorphy and term-by-term differentiation]
Let $\chi$ be a formal $q$-series with $\|\chi(h+n)\|$ subexponential. Then $z\mapsto\chi^{\C}_h(z)$ is holomorphic on
$\lfourteenH$, and for $z\in\lfourteenH$
\begin{equation*}
  \frac{d}{dz}\chi^{\C}_h(z) = 2\pi i\,(D\chi)^{\C}_h(z),\qquad (D\chi)(e) := e\,\chi(e).
\end{equation*}
In particular the character of a $C_2$-cofinite VOA (with $\mathbf 1\neq0$) is holomorphic on $\lfourteenH$.
\end{btheorem}
{\raggedright\noindent\textbf{Lean.} \lean{VOA.}\allowbreak\lean{qDeriv},\allowbreak\ \lean{VOA.}\allowbreak\lean{differentiableOn_qEvalC},\allowbreak\ \lean{VOA.}\allowbreak\lean{deriv_qEvalC},\allowbreak\ \lean{VOA.}\allowbreak\lean{differentiableOn_char_adjoint}\par}
{\raggedright\noindent File \leanpath{VOA/\allowbreak Modularity/\allowbreak ZhuBlockConditions.lean} (a Chapter~15 file).\par}
\begin{leancode}
noncomputable def qDeriv (χ : FormalQSeries) : FormalQSeries := fun e => e * χ e
\end{leancode}
{\raggedright\noindent File \leanpath{VOA/\allowbreak Modularity/\allowbreak QEvalHolomorphic.lean}.\par}
\begin{leancode}
theorem differentiableOn_qEvalC (χ : FormalQSeries) (h : ℂ)
    (hsub : Subexponential (fun n : ℕ => ‖χ (h + (n : ℂ))‖)) :
    DifferentiableOn ℂ (qEvalC χ h) (imGT 0)

theorem deriv_qEvalC (χ : FormalQSeries) (h : ℂ)
    (hsub : Subexponential (fun n : ℕ => ‖χ (h + (n : ℂ))‖)) {z : ℂ} (hz : z ∈ imGT 0) :
    _root_.deriv (qEvalC χ h) z = 2 * (Real.pi : ℂ) * Complex.I * qEvalC (qDeriv χ) h z

theorem differentiableOn_char_adjoint (W : VertexOperatorAlgebra 𝒱) (hne : W.vacuum ≠ 0)
    (hC2 : C2Cofinite W) :
    DifferentiableOn ℂ
      (qEvalC (gradedDim (adjointModule W)) (0 - W.centralCharge / 24)) (imGT 0)
\end{leancode}
\noindent\textbf{Hypotheses / status.} Growth\par
\medskip
\begin{btheorem}{14.35}[Equivariance of the Serre derivative]
For $k\in\Z$, $\gamma\in\lfourteenSL$, $z\in\lfourteenH$, and $f$ differentiable at $\gamma z$:
$(\partial_kf)|_{k+2}\gamma\,(z) = \partial_k(f|_k\gamma)(z)$. (No modularity of $f$ is required. This is equivariance only:
it does \emph{not} show that $\partial_k$ maps $M_k$ into $M_{k+2}$, which also needs holomorphy at the cusp; that is
Serre closure, Definition~14.14, proved separately in \leanpath{EisensteinSerreClosure.lean}.)
\end{btheorem}
{\raggedright\noindent\textbf{Lean.} \lean{VOA.}\allowbreak\lean{serreDerivC_slashC}\par}
{\raggedright\noindent File \leanpath{VOA/\allowbreak Modularity/\allowbreak SerreDerivative.lean}.\par}
\begin{leancode}
theorem serreDerivC_slashC (k : ℤ) (γ : SL(2, ℤ)) (f : ℂ → ℂ) {z : ℂ} (hz : z ∈ imGT 0)
    (hf : DifferentiableAt ℂ f (mobiusC γ z)) :
    slashC (k + 2) γ (serreDerivC k f) z = serreDerivC k (slashC k γ f) z
\end{leancode}
\noindent\textbf{Hypotheses / status.} Differentiability; unconditional\par
\medskip
\begin{btheorem}{14.41}[The majorant theorem]
Let $\chi_n, a_n, b_n\ge0$ with $a$, $b$ subexponential, let $p_n\to+\infty$, and suppose that for some $N_0$ and all $n\ge N_0$
\begin{equation*}
  p_n\,\chi_n \;\le\; \sum_{m<n}a_{n-m}\,\chi_m + b_n.
\end{equation*}
Then $\chi$ is subexponential.
\end{btheorem}
{\raggedright\noindent\textbf{Lean.} \lean{VOA.}\allowbreak\lean{subexponential_of_majorant_recursion}\par}
{\raggedright\noindent File \leanpath{VOA/\allowbreak Modularity/\allowbreak Majorant.lean}.\par}
\begin{leancode}
theorem subexponential_of_majorant_recursion {χ a b p : ℕ → ℝ}
    (hχ0 : ∀ n, 0 ≤ χ n) (ha0 : ∀ n, 0 ≤ a n) (hb0 : ∀ n, 0 ≤ b n)
    (ha : Subexponential a) (hb : Subexponential b)
    (hp : Tendsto p atTop atTop) (N₀ : ℕ)
    (hrec : ∀ n, N₀ ≤ n → p n * χ n ≤ (∑ m ∈ Finset.range n, a (n - m) * χ m) + b n) :
    Subexponential χ
\end{leancode}
\noindent\textbf{Hypotheses / status.} $p_n\to\infty$\par
\medskip
\begin{bproposition}{14.64}[Check for the template warning]
For every super character system $S$: $\mathtt{CompleteRationalFamily}\,S\,P\iff P$; $\mathtt{Conjecture\_HalfWeight}\,S\,
\mathrm{False}$ holds; and the zero family satisfies \lean{SuperRealisesLawOn} for every $\Gamma$, $J$, $\rho$.
\end{bproposition}
\noindent\emph{(The chapter quotes no Lean declaration for this statement; see book p.~719.)}\par
\medskip
\endgroup
\setcounter{chapter}{14}
\chapter{Zhu's modularity theorem}
\begingroup
\providecommand{\lfifteenH}{\mathbb{H}}
\providecommand{\lfifteenB}{\mathfrak{B}}
\providecommand{\lfifteenvac}{\mathbf{1}}
\providecommand{\lfifteenSL}{SL_2(\Z)}
\providecommand{\lfifteenZ}{\mathcal{Z}}
\providecommand{\lfifteenq}{\mathsf{q}}
\noindent{\small 10 of the chapter's 42 numbered statements: its headline results and every definition they use.}\par\medskip
\begin{btheorem}{15.1}[Zhu 1996, Theorem~5.3.3; Dong--Li--Mason 2000]
Let $V$ be rational and $C_2$-cofinite, with irreducible modules $M_1,\dots,M_r$. Then:
\begin{enumerate}
\item each $\lfifteenZ_{M_i}(v,\tau)$ converges to a holomorphic function on $\lfifteenH$;
\item there is a representation $\rho:\lfifteenSL\to GL_r(\C)$ such that for every $v$ in the square-bracket weight space
  $V_{[n]}$ (Section~15.3)
  \begin{equation*}
  \lfifteenZ_{M_i}(v,\gamma\tau)=(c\tau+d)^n\sum_{j}\rho(\gamma)_{ij}\lfifteenZ_{M_j}(v,\tau);
  \end{equation*}
\item the functions $\lfifteenZ_{M_1}(\cdot,\tau),\dots,\lfifteenZ_{M_r}(\cdot,\tau)$ are linearly independent, so $\rho$ is
  unique.
\end{enumerate}
\end{btheorem}
\noindent\emph{(The chapter quotes no Lean declaration for this statement; see book p.~731.)}\par
\medskip
\begin{btheorem}{15.16}[Convergence of one-point functions; Zhu 1996 \S4.4, Dong--Li--Mason 2000 \S\S6--7]
Let $V$ be $C_2$-cofinite and $M$ an ordinary module with conformal weight $h$. Then for every $v\in V$ the coefficients
of $\lfifteenZ_M(v;q)$ along $h-c/24+\N$ grow subexponentially. Consequently every trace function $S_i(v,\tau)$ is given by
an absolutely convergent series on $\lfifteenH$.
\end{btheorem}
{\raggedright\noindent\textbf{Lean.} \lean{VOA.}\allowbreak\lean{onePointSubexp_of_c2}\par}
{\raggedright\noindent File \leanpath{VOA/\allowbreak Modularity/\allowbreak BlockGrowth.lean}.\par}
\begin{leancode}
theorem onePointSubexp_of_c2 {M : Type u} [AddCommGroup M] [Module ℂ M]
    (N : OrdVModule W M) {hw : ℂ} (hcw : HasConformalWeight N hw)
    (hzero : VacuumGradeZero W) (hC2 : C2Cofinite W) :
    OnePointSubexp N (hw - W.centralCharge / 24) := by
\end{leancode}
\noindent\textbf{Hypotheses / status.} $C_2$, conformal weight\par
\medskip
\begin{btheorem}{15.34}[Zhu's modularity theorem, pinned to a given family]
Let $V$ be a vertex operator algebra of the library such that
\begin{enumerate}
\item[(a)] $V$ is $C_2$-cofinite;
\item[(b)] Zhu's algebra $A(V)$ is semisimple;
\item[(c)] $\Phi(V)$, the set of isomorphism classes of irreducible ordinary modules, is finite;
\item[(d)] the datum \lean{Pr : BottomProjection (parabolicRepOfAMod W (regularAMod W))} is given (Chapter~13).
\end{enumerate}
Let $\iota$ be a finite set with $\#\iota=\#\Phi(V)$ (hypothesis \lean{hcard}: one index per isomorphism class), and $(M_i,h_i)_{i\in\iota}$ irreducible ordinary modules with conformal
weights $h_i$, whose bottoms $(M_i)_{h_i}$ are simple $A(V)$-modules, pairwise non-isomorphic whenever $h_j=h_k$, $j\neq k$.
Then there are modular data $D$ with $D.\mathrm{Mfam}=(M_i)$ and $D.h=(h_i)$ satisfying Zhu's analytic law: there is a
homomorphism $\rho:\lfifteenSL\to\mathrm{Mat}_\iota(\C)$ with $\rho(T)=\mathrm{diag}(e^{2\pi i(h_i-c/24)})$ such that for all
$\gamma\in\lfifteenSL$, $n\in\N$, $v\in V_{[n]}$, $i\in\iota$, $\tau\in\lfifteenH$,
\begin{equation*}
  \tr_{M_i}o(v)q^{L_0-c/24}\big|_{\gamma\tau}=(c\tau+d)^n\sum_{j\in\iota}\rho(\gamma)_{ij}\,\tr_{M_j}o(v)q^{L_0-c/24}\big|_{\tau}.
\end{equation*}
Since (d) is supplied for every $V$ by the theorem \lean{VOA.pbw} (through \lean{voaPBWBottomProjection}), the
conclusion holds under (a)--(c) (this is \lean{zhuModularity_pinned_of_C2Cofinite_voaPBW}).
\end{btheorem}
{\raggedright\noindent\textbf{Lean.} \lean{VOA.}\allowbreak\lean{zhuModularity_pinned_of_C2Cofinite},\allowbreak\ \lean{VOA.}\allowbreak\lean{zhuModularity_pinned_of_C2Cofinite_voaPBW}\par}
{\raggedright\noindent Files \leanpath{VOA/\allowbreak Modularity/\allowbreak GNProp8Closure.lean} and
\leanpath{VOA/\allowbreak Modularity/\allowbreak GNProp8PBW.lean}.\par}
\begin{leancode}
theorem zhuModularity_pinned_of_C2Cofinite [Finite (Φ W)]
    (Mfam : ι → IrrModule W) (h : ι → ℂ)
    (hC2 : C2Cofinite W)
    (hss : letI := instRingA W; IsSemisimpleRing (A W))
    (Pr : BottomProjection (parabolicRepOfAMod W (regularAMod W)))
    (hcard : Fintype.card ι = Nat.card (Φ W))
    (hcw : ∀ i, HasConformalWeight (Mfam i).isMod (h i))
    (hsep : ∀ j k : ι, j ≠ k → h j = h k →
      IsEmpty (AModEquiv W (lowestAMod (hcw j).2) (lowestAMod (hcw k).2))) :
    ∃ D : ModularData W ι, D.Mfam = Mfam ∧ D.h = h ∧ ZhuModularityAnalytic D :=
  zhuModularity_pinned_of_strictSpanning Mfam h hC2 hss Pr
    (strictSpanning_of_C2Cofinite W hC2) hcard hcw hsep

theorem zhuModularity_pinned_of_C2Cofinite_voaPBW [Finite (Φ W)]
    (Mfam : ι → IrrModule W) (h : ι → ℂ)
    (hC2 : C2Cofinite W)
    (hss : letI := instRingA W; IsSemisimpleRing (A W))
    (hcard : Fintype.card ι = Nat.card (Φ W))
    (hcw : ∀ i, HasConformalWeight (Mfam i).isMod (h i))
    (hsep : ∀ j k : ι, j ≠ k → h j = h k →
      IsEmpty (AModEquiv W (lowestAMod (hcw j).2) (lowestAMod (hcw k).2))) :
    ∃ D : ModularData W ι, D.Mfam = Mfam ∧ D.h = h ∧ ZhuModularityAnalytic D :=
  zhuModularity_pinned_of_C2Cofinite Mfam h hC2 hss
    (voaPBWBottomProjection (parabolicRepOfAMod W (regularAMod W)))
\end{leancode}
\medskip
\begin{btheorem}{15.35}[Finitely many irreducible modules; Dong--Li--Mason, Abe--Buhl--Dong]
If $V$ is $C_2$-cofinite, then $\Phi(V)$ is finite.
\end{btheorem}
{\raggedright\noindent\textbf{Lean.} \lean{VOA.}\allowbreak\lean{Huang.}\allowbreak\lean{ZhuFiniteIrr.}\allowbreak\lean{finite_Φ_of_C2Cofinite_pbw}\par}
{\raggedright\noindent File \leanpath{VOA/\allowbreak Classification/\allowbreak COne/\allowbreak McRaeO1FinitePhi.lean}
(Chapter~22), in the namespace \lean{VOA.Huang.ZhuFiniteIrr}, with section variable
\lean{{W : VertexOperatorAlgebra 𝒱}}.\par}
\begin{leancode}
theorem finite_Φ_of_C2Cofinite_pbw (hC2 : C2Cofinite W) : Finite (Φ W) :=
  finite_Φ_of_C2Cofinite' (voaPBWBottomProjection (parabolicRepOfAMod W (regularAMod W))) hC2
\end{leancode}
\noindent\textbf{Hypotheses / status.} $C_2$\par
\medskip
\begin{btheorem}{15.37}[Zhu's modularity theorem for one-point functions, canonical family, semisimple form]
Let $V$ be a vertex operator algebra of the library (so of CFT type) such that
\begin{enumerate}
\item[(a)] $V$ is $C_2$-cofinite, and
\item[(b)] Zhu's algebra $A(V)$ is a semisimple ring.
\end{enumerate}
Then $\Phi(V)$ is finite, and there are modular data $D$ indexed by $\Phi(V)$, whose family is the canonical family
$(M_c)_{c\in\Phi(V)}$ and whose weights are its conformal weights $(h_c)$, satisfying Zhu's analytic law: there is a group
homomorphism $\rho:\lfifteenSL\to\mathrm{Mat}_{\Phi(V)}(\C)$ with $\rho(T)=\mathrm{diag}(e^{2\pi i(h_c-c_V/24)})$ such that for
all $\gamma=\left(\begin{smallmatrix}a&b\\ c&d\end{smallmatrix}\right)\in\lfifteenSL$, $n\in\N$, $v\in V_{[n]}$,
$c'\in\Phi(V)$ and $\tau\in\lfifteenH$,
\begin{equation*}
  \tr_{M_{c'}}o(v)q^{L_0-c_V/24}\big|_{\gamma\tau}=(c\tau+d)^n\sum_{c''\in\Phi(V)}\rho(\gamma)_{c'c''}\,
  \tr_{M_{c''}}o(v)q^{L_0-c_V/24}\big|_{\tau},
\end{equation*}
where $c_V$ is the central charge.
\end{btheorem}
{\raggedright\noindent\textbf{Lean.} \lean{VOA.}\allowbreak\lean{zhuModularity_canonical_of_C2Cofinite_semisimple}\par}
{\raggedright\noindent File \leanpath{VOA/\allowbreak Modularity/\allowbreak ZhuModularityCanonical.lean}.\par}
\begin{leancode}
theorem zhuModularity_canonical_of_C2Cofinite_semisimple (W : VertexOperatorAlgebra 𝒱)
    (hC2 : C2Cofinite W) (hss : letI := instRingA W; IsSemisimpleRing (A W)) :
    letI : Finite (Φ W) := VOA.Huang.ZhuFiniteIrr.finite_Φ_of_C2Cofinite_pbw hC2
    letI : Fintype (Φ W) := Fintype.ofFinite (Φ W)
    letI : DecidableEq (Φ W) := Classical.decEq (Φ W)
    ∃ D : ModularData W (Φ W), D.Mfam = canonicalIrrModule ∧
      D.h = (fun c => conformalWeightOf (canonicalIrrModule c)) ∧ ZhuModularityAnalytic D := by
\end{leancode}
\medskip
\begin{btheorem}{15.38}[Zhu's semisimplicity theorem; Zhu 1996, Theorem~2.2.3; Dong--Li--Mason]
Let $V$ be a vertex operator algebra of the library, and let the datum \lean{Pr} of Chapter~13 be given. If $V$ is
rational in the sense of \lean{AdmissiblyRational} (every admissible $V$-module is completely reducible;
Chapter~12, Definition~12.35), then Zhu's algebra $A(V)$ is a semisimple ring. Since \lean{Pr} is
supplied for every $V$ by the theorem \lean{VOA.pbw}, the implication holds for every $V$; in particular the named predicate
\lean{ZhuAdmissibleSemisimple W} of Chapter~12 holds for every $V$.
\end{btheorem}
{\raggedright\noindent\textbf{Lean.} \lean{VOA.}\allowbreak\lean{Huang.}\allowbreak\lean{ZhuSemisimple.}\allowbreak\lean{isSemisimpleRing_A_of_admissiblyRational},\allowbreak\ \lean{VOA.}\allowbreak\lean{Huang.}\allowbreak\lean{ZhuSemisimple.}\allowbreak\lean{zhuAdmissibleSemisimple_of_Pr}\par}
{\raggedright\noindent File \leanpath{VOA/\allowbreak Huang/\allowbreak ZhuSemisimple.lean} (Chapter~21), namespace
\lean{VOA.Huang.ZhuSemisimple}, with section variables \lean{{𝒱 : Type u} [AddCommGroup 𝒱] [Module ℂ 𝒱] {W : VertexOperatorAlgebra 𝒱}}
and \lean{(Pr : BottomProjection (parabolicRepOfAMod W (regularAMod W)))}; each theorem is preceded by the line
\lean{include Pr in}.\par}
\begin{leancode}
theorem isSemisimpleRing_A_of_admissiblyRational (hrat : AdmissiblyRational W) :
    letI := instRingA W; IsSemisimpleRing (A W) := by

theorem zhuAdmissibleSemisimple_of_Pr : ZhuAdmissibleSemisimple W :=
  fun hrat => isSemisimpleRing_A_of_admissiblyRational Pr hrat
\end{leancode}
\medskip
\begin{btheorem}{15.39}[Zhu's modularity theorem for rational $C_2$-cofinite $V$; the chapter's headline]
Let $V$ be a vertex operator algebra of the library (so of CFT type) such that
\begin{enumerate}
\item[(a)] $V$ is $C_2$-cofinite, and
\item[(r)] $V$ is rational, in the sense of \lean{AdmissiblyRational}.
\end{enumerate}
Then the conclusion of Theorem~15.37 holds: $\Phi(V)$ is finite, and the one-point functions
$\tr_{M_c}o(v)q^{L_0-c_V/24}$ of the canonical irreducible modules $M_c$, $c\in\Phi(V)$, with their conformal weights $h_c$,
transform under $\lfifteenSL$ by a matrix representation $\rho$ with $\rho(T)=\mathrm{diag}(e^{2\pi i(h_c-c_V/24)})$, with
weight $n$ for $v\in V_{[n]}$. In particular this holds for every strongly rational $V$.
\end{btheorem}
{\raggedright\noindent\textbf{Lean.} \lean{VOA.}\allowbreak\lean{zhuModularity_canonical_of_C2Cofinite_admissiblyRational},\allowbreak\ \lean{VOA.}\allowbreak\lean{zhuModularity_canonical_of_stronglyRational}\par}
{\raggedright\noindent File \leanpath{VOA/\allowbreak Modularity/\allowbreak ZhuModularityCanonical.lean}.\par}
\begin{leancode}
theorem zhuModularity_canonical_of_C2Cofinite_admissiblyRational (W : VertexOperatorAlgebra 𝒱)
    (hC2 : C2Cofinite W) (hrat : AdmissiblyRational W) :
    letI : Finite (Φ W) := VOA.Huang.ZhuFiniteIrr.finite_Φ_of_C2Cofinite_pbw hC2
    letI : Fintype (Φ W) := Fintype.ofFinite (Φ W)
    letI : DecidableEq (Φ W) := Classical.decEq (Φ W)
    ∃ D : ModularData W (Φ W), D.Mfam = canonicalIrrModule ∧
      D.h = (fun c => conformalWeightOf (canonicalIrrModule c)) ∧ ZhuModularityAnalytic D :=
  zhuModularity_canonical_of_C2Cofinite_semisimple W hC2
    (VOA.Huang.ZhuSemisimple.isSemisimpleRing_A_of_admissiblyRational
      (voaPBWBottomProjection (parabolicRepOfAMod W (regularAMod W))) hrat)

theorem zhuModularity_canonical_of_stronglyRational (W : VertexOperatorAlgebra 𝒱)
    (hSR : StronglyRational W) :
    letI : Finite (Φ W) := VOA.Huang.ZhuFiniteIrr.finite_Φ_of_C2Cofinite_pbw hSR.c2Cofinite
    letI : Fintype (Φ W) := Fintype.ofFinite (Φ W)
    letI : DecidableEq (Φ W) := Classical.decEq (Φ W)
    ∃ D : ModularData W (Φ W), D.Mfam = canonicalIrrModule ∧
      D.h = (fun c => conformalWeightOf (canonicalIrrModule c)) ∧ ZhuModularityAnalytic D :=
  zhuModularity_canonical_of_C2Cofinite_admissiblyRational W hSR.c2Cofinite hSR.rational
\end{leancode}
\medskip
\begin{btheorem}{15.40}[Zhu's theorem for strongly rational $V$, the older packaging with two (now proved) literature hypotheses]
Let $V$ be strongly rational (simple, self-dual, $C_2$-cofinite, rational; Definition~12.43), and assume
the Dong--Li--Mason finiteness theorem (rational $\Rightarrow$ finitely many irreducible admissible modules) and Zhu's
semisimplicity theorem (rational $\Rightarrow$ $A(V)$ semisimple) for $V$. Then, with $\iota=\Phi(V)$ and the canonical
family of representatives with their conformal weights, the pinned conclusion of Theorem~15.34 holds. The
proof uses the theorem \lean{VOA.pbw}.
\end{btheorem}
{\raggedright\noindent\textbf{Lean.} \lean{VOA.}\allowbreak\lean{canonicalIrrModule},\allowbreak\ \lean{VOA.}\allowbreak\lean{conformalWeightOf},\allowbreak\ \lean{VOA.}\allowbreak\lean{zhuModularity_stronglyRational_voaPBW}\par}
{\raggedright\noindent Files \leanpath{VOA/\allowbreak Modularity/\allowbreak ZhuStronglyRational.lean},
\leanpath{VOA/\allowbreak Modularity/\allowbreak StronglyRationalFamily.lean}, \leanpath{VOA/\allowbreak Modularity/\allowbreak ConformalWeightChoice.lean}.\par}
\begin{leancode}
noncomputable def canonicalIrrModule (c : Φ W) : IrrModule W := Quotient.out c

noncomputable def conformalWeightOf (V : IrrModule W) : ℂ :=
  (OrdVModule.exists_hasConformalWeight_of_irreducible V.isMod V.irreducible).choose

theorem zhuModularity_stronglyRational_voaPBW
    (hSR : StronglyRational W)
    (hFin : DLMFinitelyManyIrredAdmissible W)
    (hZhu : ZhuAdmissibleSemisimple W) :
    letI : Finite (Φ W) := finite_Φ_of_finite_Φadm W (hFin hSR.rational)
    letI : Fintype (Φ W) := Fintype.ofFinite (Φ W)
    letI : DecidableEq (Φ W) := Classical.decEq (Φ W)
    ∃ D : ModularData W (Φ W),
      D.Mfam = canonicalIrrModule ∧
      D.h = (fun c => conformalWeightOf (canonicalIrrModule c)) ∧
      ZhuModularityAnalytic D := by
\end{leancode}
\medskip
\begin{btheorem}{15.43}[Zhu's theorem for $V_{2m}^+$ and for the tetrahedral model, with no hypotheses]
(i) For every $m\ge2$, the conclusion of Theorem~15.37 holds for $V=V_{2m}^+$: $\Phi(V_{2m}^+)$ is finite, and
the one-point functions of its canonical irreducible modules, with their conformal weights, transform under $\lfifteenSL$ by a
matrix representation with diagonal $\rho(T)$. (ii) The same holds for the tetrahedral model.
\end{btheorem}
{\raggedright\noindent\textbf{Lean.} \lean{VOA.}\allowbreak\lean{Classification.}\allowbreak\lean{COne.}\allowbreak\lean{zhuModularity_V2mPlus},\allowbreak\ \lean{VOA.}\allowbreak\lean{Classification.}\allowbreak\lean{COne.}\allowbreak\lean{zhuModularity_cubicModel},\allowbreak\ \lean{VOA.}\allowbreak\lean{Classification.}\allowbreak\lean{COne.}\allowbreak\lean{V2mPlus_c2Cofinite},\allowbreak\ \lean{VOA.}\allowbreak\lean{Classification.}\allowbreak\lean{COne.}\allowbreak\lean{V2mPlus_zhuSemisimple},\allowbreak\ \lean{VOA.}\allowbreak\lean{Classification.}\allowbreak\lean{COne.}\allowbreak\lean{V2mPlus_zhu_finrank},\allowbreak\ \lean{VOA.}\allowbreak\lean{Classification.}\allowbreak\lean{COne.}\allowbreak\lean{ZhuV2mPlusSemisimple},\allowbreak\ \lean{VOA.}\allowbreak\lean{Classification.}\allowbreak\lean{COne.}\allowbreak\lean{OrbA4.}\allowbreak\lean{cubicModel_c2Cofinite_holds},\allowbreak\ \lean{VOA.}\allowbreak\lean{Classification.}\allowbreak\lean{COne.}\allowbreak\lean{OrbTop.}\allowbreak\lean{cubicModel_zhuSemisimple}\par}
{\raggedright\noindent File \leanpath{VOA/\allowbreak Classification/\allowbreak COne/\allowbreak ZhuModularityV2mPlus.lean},
namespace \lean{VOA.Classification.COne}.\par}
\begin{leancode}
theorem zhuModularity_V2mPlus (m : ℕ) (hm : 1 ≤ m) (hm2 : 2 ≤ m) :
    letI : Finite (Φ (V2mPlus m hm)) :=
      VOA.Huang.ZhuFiniteIrr.finite_Φ_of_C2Cofinite_pbw (V2mPlus_c2Cofinite m hm)
    letI : Fintype (Φ (V2mPlus m hm)) := Fintype.ofFinite _
    letI : DecidableEq (Φ (V2mPlus m hm)) := Classical.decEq _
    ∃ D : VOA.ModularData (V2mPlus m hm) (Φ (V2mPlus m hm)), D.Mfam = canonicalIrrModule ∧
      D.h = (fun c => conformalWeightOf (canonicalIrrModule c)) ∧ VOA.ZhuModularityAnalytic D :=
  zhuModularity_canonical_of_C2Cofinite_semisimple (V2mPlus m hm) (V2mPlus_c2Cofinite m hm)
    (V2mPlus_zhuSemisimple m hm hm2)

theorem zhuModularity_cubicModel :
    letI : Finite (Φ OrbRat2.cubicModel) :=
      VOA.Huang.ZhuFiniteIrr.finite_Φ_of_C2Cofinite_pbw OrbA4.cubicModel_c2Cofinite_holds
    letI : Fintype (Φ OrbRat2.cubicModel) := Fintype.ofFinite _
    letI : DecidableEq (Φ OrbRat2.cubicModel) := Classical.decEq _
    ∃ D : VOA.ModularData OrbRat2.cubicModel (Φ OrbRat2.cubicModel),
      D.Mfam = canonicalIrrModule ∧
      D.h = (fun c => conformalWeightOf (canonicalIrrModule c)) ∧ VOA.ZhuModularityAnalytic D :=
  zhuModularity_canonical_of_C2Cofinite_semisimple OrbRat2.cubicModel
    OrbA4.cubicModel_c2Cofinite_holds OrbTop.cubicModel_zhuSemisimple
\end{leancode}
\begin{leancode}
theorem V2mPlus_c2Cofinite (m : ℕ) (hm : 1 ≤ m) : C2Cofinite (V2mPlus m hm) :=
  v2mPlusC2CofiniteFact_holds m hm

theorem V2mPlus_zhuSemisimple : ZhuV2mPlusSemisimple :=
  zhuSemisimple_of_fourExtraCharacters
    (fourExtraCharacters_of_twistedOmegaECharacters V2mPlus_twistedOmegaECharacters)

theorem V2mPlus_zhu_finrank (m : ℕ) (hm : 1 ≤ m) (hm2 : 2 ≤ m) :
    Module.finrank ℂ (VOA.A (V2mPlus m hm)) = m + 7 :=
\end{leancode}
\begin{leancode}
def ZhuV2mPlusSemisimple : Prop :=
  ∀ (m : ℕ) (hm : 1 ≤ m), 2 ≤ m →
    letI := instRingA (V2mPlus m hm)
    IsSemisimpleRing (VOA.A (V2mPlus m hm))
\end{leancode}
\begin{leancode}
theorem cubicModel_c2Cofinite_holds : C2Cofinite cubicModel := by

theorem cubicModel_zhuSemisimple : IsSemisimpleRing (VOA.A cubicModel) :=
  OrbRat2.isSemisimpleRing_of_charTable OrbA4.cubicModel_c2Cofinite_holds
    OrbSgl.cubicZhuDimensionBound_holds OrbRat6.cubicZhuCharacterTable_holds
\end{leancode}
\medskip
\begin{btheorem}{15.44}[Zhu's theorem for the lattice VOAs $V_{\sqrt{2m}\Z}$, and for every nice VOA]
(i) For every $m\ge1$, the conclusion of Theorem~15.37 holds for the rank-one lattice VOA
$V=V_{2m}=V_{\sqrt{2m}\Z}$, with no further hypotheses. (ii) It holds for every \emph{nice} vertex operator algebra $V$, i.e.\
every $V$ that is strongly rational and pseudo-unitary (Gannon--Rayhaun, Definition~2.6).
\end{btheorem}
{\raggedright\noindent\textbf{Lean.} \lean{VOA.}\allowbreak\lean{Classification.}\allowbreak\lean{COne.}\allowbreak\lean{zhuModularity_V2m},\allowbreak\ \lean{VOA.}\allowbreak\lean{Classification.}\allowbreak\lean{COne.}\allowbreak\lean{zhuModularity_of_nice},\allowbreak\ \lean{VOA.}\allowbreak\lean{Classification.}\allowbreak\lean{COne.}\allowbreak\lean{Nice},\allowbreak\ \lean{VOA.}\allowbreak\lean{Classification.}\allowbreak\lean{COne.}\allowbreak\lean{V2m},\allowbreak\ \lean{VOA.}\allowbreak\lean{Classification.}\allowbreak\lean{COne.}\allowbreak\lean{V2m_stronglyRational}\par}
{\raggedright\noindent File \leanpath{VOA/\allowbreak Classification/\allowbreak COne/\allowbreak ZhuModularityV2mPlus.lean},
namespace \lean{VOA.Classification.COne}.\par}
\begin{leancode}
theorem zhuModularity_V2m (m : ℕ) (hm : 1 ≤ m) :
    letI : Finite (Φ (V2m m hm)) :=
      VOA.Huang.ZhuFiniteIrr.finite_Φ_of_C2Cofinite_pbw (V2m_stronglyRational m hm).c2Cofinite
    letI : Fintype (Φ (V2m m hm)) := Fintype.ofFinite _
    letI : DecidableEq (Φ (V2m m hm)) := Classical.decEq _
    ∃ D : VOA.ModularData (V2m m hm) (Φ (V2m m hm)), D.Mfam = canonicalIrrModule ∧
      D.h = (fun c => conformalWeightOf (canonicalIrrModule c)) ∧ VOA.ZhuModularityAnalytic D :=
  zhuModularity_canonical_of_stronglyRational (V2m m hm) (V2m_stronglyRational m hm)

theorem zhuModularity_of_nice {𝒱 : Type*} [AddCommGroup 𝒱] [Module ℂ 𝒱]
    (V : VertexOperatorAlgebra 𝒱) (hV : Nice V) :
    letI : Finite (Φ V) := VOA.Huang.ZhuFiniteIrr.finite_Φ_of_C2Cofinite_pbw hV.1.c2Cofinite
    letI : Fintype (Φ V) := Fintype.ofFinite _
    letI : DecidableEq (Φ V) := Classical.decEq _
    ∃ D : VOA.ModularData V (Φ V), D.Mfam = canonicalIrrModule ∧
      D.h = (fun c => conformalWeightOf (canonicalIrrModule c)) ∧ VOA.ZhuModularityAnalytic D :=
  zhuModularity_canonical_of_stronglyRational V hV.1
\end{leancode}
\begin{leancode}
def Nice (V : VertexOperatorAlgebra 𝒱) : Prop :=
  StronglyRational V ∧ PseudoUnitary V

noncomputable def V2m (m : ℕ) (hm : 1 ≤ m) : VertexOperatorAlgebra (LatModuleGen (latP m)) :=
  latticeVOAGen (latP m) (latP_even m) (latP_posDef m hm)

theorem V2m_stronglyRational (m : ℕ) (hm : 1 ≤ m) : StronglyRational (V2m m hm) :=
  { V2m_stronglyFinite m hm with rational := V2m_admissiblyRational hm }
\end{leancode}
\medskip
\endgroup
\setcounter{chapter}{15}
\chapter{Quasi-lisse vertex algebras and the Arakawa--Kawasetsu theorem}
\begingroup
\providecommand{\lsixteenvac}{\mathbf{1}}
\providecommand{\lsixteenHP}{\mathrm{HP}_0}
\providecommand{\lsixteensl}{\mathfrak{sl}_2}
\providecommand{\lsixteenCas}{\mathrm{Cas}}
\providecommand{\lsixteenNil}{\mathcal{N}}
\providecommand{\lsixteenR}{R}
\providecommand{\lsixteenVA}{V_{\mathcal A}}
\providecommand{\lsixteenOq}{O_q(V)}
\noindent{\small 13 of the chapter's 43 numbered statements: its headline results and every definition they use.}\par\medskip
\begin{btheorem}{16.4}[Etingof--Schedler 2010]
If $R$ is a finitely generated Poisson algebra whose variety $\operatorname{Specm}R$ has finitely many
symplectic leaves, then $\dim\lsixteenHP(R)<\infty$.
\end{btheorem}
\noindent\emph{(The chapter quotes no Lean declaration for this statement; see book p.~797.)}\par
\medskip
\begin{bdefinition}{16.9}[Etingof-finite]
$V$ is \emph{Etingof-finite} if $\dim_\C\lsixteenHP(R_V)<\infty$.
\end{bdefinition}
{\raggedright\noindent\textbf{Lean.} \lean{VOA.}\allowbreak\lean{EtingofFinite}\par}
{\raggedright\noindent File \leanpath{VOA/\allowbreak Modules/\allowbreak PoissonTrace.lean}.\par}
\begin{leancode}
def EtingofFinite (W : VertexOperatorAlgebra 𝒱) : Prop := FiniteDimensional ℂ (HP₀ W)
\end{leancode}
\noindent\textbf{Hypotheses / status.} Definition; \textbf{not} quasi-lisseness (weaker, by Etingof--Schedler)\par
\medskip
\begin{btheorem}{16.16}[Growth of one-point functions from Etingof-finiteness]
If $V$ is Etingof-finite and $M$ is an ordinary module of conformal weight $h$ (that is, $M_h\ne0$ and every nonzero
weight space $M_\alpha$ has $\alpha\in h+\N$, Chapter~12), then for every $v\in V$ the
coefficients of the torus one-point function $Z_M(v;q) = \operatorname{tr}_M o(v)q^{L_0-c/24}$ grow subexponentially:
for each $r>1$ there is $C$ with $|[q^{h-c/24+n}]Z_M(v;q)|\le C r^n$.
\end{btheorem}
{\raggedright\noindent\textbf{Lean.} \lean{VOA.}\allowbreak\lean{onePointSubexp_of_etingofFinite},\allowbreak\ \lean{VOA.}\allowbreak\lean{OnePointSubexp},\allowbreak\ \lean{VOA.}\allowbreak\lean{Subexponential}\par}
{\raggedright\noindent File \leanpath{VOA/\allowbreak Modules/\allowbreak QuasiLisseFinite.lean}.\par}
\begin{leancode}
theorem onePointSubexp_of_etingofFinite {M : Type u} [AddCommGroup M] [Module ℂ M]
    (N : OrdVModule W M) {hw : ℂ} (hcw : HasConformalWeight N hw)
    (hzero : VacuumGradeZero W) (hE : EtingofFinite W) :
    OnePointSubexp N (hw - W.centralCharge / 24) := by
\end{leancode}
{\raggedright\noindent Files \leanpath{VOA/\allowbreak Modularity/\allowbreak AnalyticBlock.lean},
\leanpath{VOA/\allowbreak Modularity/\allowbreak ConvergenceBridge.lean}.\par}
\begin{leancode}
def OnePointSubexp (N : OrdVModule W M) (h : ℂ) : Prop :=
  ∀ v : 𝒱, Subexponential (fun n : ℕ => ‖torusOnePoint N v (h + (n : ℂ))‖)

def Subexponential (a : ℕ → ℝ) : Prop :=
  ∀ r : ℝ, 1 < r → ∃ C : ℝ, ∀ n : ℕ, a n ≤ C * r ^ n
\end{leancode}
\noindent\textbf{Hypotheses / status.} \lean{hcw}, \lean{hE}\par
\medskip
\begin{bdefinition}{16.19}[Formal genus-one block; Zhu Def.~5.1.1, formal form]
A linear map $S:V\to\{\text{formal }q\text{-series}\}$ is a \emph{formal block} if for every homogeneous $a$ of weight
$k_a\ge1$ and every $v\in V$:
(i) $S(a[0]v) = 0$;
(ii) $S(a[-2]v) + \sum_{m\ge0}(2m+1)\,G_{2m+2}\star S(a[2m]v) = 0$;
(iii) $S(L[-2]v) = \theta S(v) + \sum_{m\ge0}G_{2m+2}\star S(L[2m]v)$,
where $\star$ is the product of a power series with a formal $q$-series (\lean{qConv}) and $G_{2j}$ is Zhu's normalised
Eisenstein series (below).
\end{bdefinition}
{\raggedright\noindent\textbf{Lean.} \lean{VOA.}\allowbreak\lean{IsFormalBlock},\allowbreak\ \lean{VOA.}\allowbreak\lean{zhuG}\par}
{\raggedright\noindent File \leanpath{VOA/\allowbreak Modularity/\allowbreak ZhuBlockConditions.lean} (Chapter~15).\par}
\begin{leancode}
structure IsFormalBlock (W : VertexOperatorAlgebra 𝒱) (S : 𝒱 →ₗ[ℂ] FormalQSeries) : Prop where
  vanish : ∀ (ka : ℕ), 1 ≤ ka → ∀ a ∈ W.grading (ka : ℤ), ∀ v : 𝒱, S (vSqZ W a ka 0 v) = 0
  negTwo : ∀ (ka : ℕ), 1 ≤ ka → ∀ a ∈ W.grading (ka : ℤ), ∀ v : 𝒱,
    S (vSqZ W a ka (-2) v)
      + ∑ᶠ m : ℕ, ((2 * (m + 1) - 1 : ℕ) : ℂ)
          • qConv (zhuG (m + 1)) (S (vSqZ W a ka (2 * (m + 1) - 2 : ℤ) v)) = 0
  omega : ∀ v : 𝒱, S (Lsq W (-2) v)
    = qDeriv (S v) + ∑ᶠ m : ℕ, qConv (zhuG (m + 1)) (S (Lsq W (2 * (m + 1) - 2 : ℤ) v))
\end{leancode}
{\raggedright\noindent File \leanpath{VOA/\allowbreak Modularity/\allowbreak ZhuNormalization.lean} (Chapter~14).\par}
\begin{leancode}
noncomputable def zhuG (j : ℕ) : ℕ → ℂ :=
  fun m =>
    if m = 0 then -(((bernoulli (2 * j) : ℚ) : ℂ) / (Nat.factorial (2 * j) : ℂ))
    else (2 / (Nat.factorial (2 * j - 1) : ℂ)) * lambertCoeff (2 * j - 1) m
\end{leancode}
\medskip
\begin{btheorem}{16.20}[AK Lemma 5.5, for primary vectors]
Let $S$ be a formal block supported on $h+\N$ and $a$ square-bracket primary. For each $i$ there are
$\eta_{i,0},\dots,\eta_{i,i}\in\mathcal S$ with $\eta_{i,i} = 1$ and
$S(L[-2]^ia) = \sum_{j=0}^{i}\eta_{i,j}\,\theta^jS(a)$.
\end{btheorem}
{\raggedright\noindent\textbf{Lean.} \lean{VOA.}\allowbreak\lean{huni_of_sqPrimary}\par}
{\raggedright\noindent File \leanpath{VOA/\allowbreak Modularity/\allowbreak AKModularODE.lean}.\par}
\begin{leancode}
theorem huni_of_sqPrimary (W : VertexOperatorAlgebra 𝒱)
    {S : 𝒱 →ₗ[ℂ] FormalQSeries} {hbase : ℂ} (hsupp : ∀ v, S v ∈ QSeriesAt hbase)
    (hS : IsFormalBlock W S) {a : 𝒱} {Δ : ℂ} (ha : IsSqPrimary W a Δ) (i : ℕ) :
    ∃ hc : ℕ → PowerSeries ℂ, (∀ j, hc j ∈ SubexpPS) ∧ hc i = 1 ∧ (∀ j, i < j → hc j = 0) ∧
      S ((⇑(Lsq W (-2)))^[i] a)
        = ∑ j ∈ Finset.range (i + 1), psSMul (hc j) (qDeriv^[j] (S a)) := by
\end{leancode}
\noindent\textbf{Hypotheses / status.} Primary $a$\par
\medskip
\begin{bproposition}{16.24}[The formal equation is almost a growth statement]
Let $\chi = q^h F(q)$ with $F = \sum_{n\ge0}F_nq^n$, $F_0\ne0$, and suppose $F$ converges on the open unit disc $D$ and its
zeros in $D$ have multiplicity at most $m$. Then for $s = m+1$ there are $g_0,\dots,g_{s-1}\in\mathcal S$ with
$\theta^s\chi + \sum_{i<s}g_i\theta^i\chi = 0$.
\end{bproposition}
{\raggedright\noindent\textbf{Lean.} \lean{VOA.}\allowbreak\lean{character_formalODE_of_etingofFinite},\allowbreak\ \lean{VOA.}\allowbreak\lean{character_formalODE_adjoint_of_etingofFinite}\par}
\noindent File \leanpath{VOA/\allowbreak Modularity/\allowbreak AKModularODE.lean}.
\begin{leancode}
theorem character_formalODE_of_etingofFinite (W : VertexOperatorAlgebra 𝒱) (N : OrdVModule W M)
    {hw : ℂ} (hcw : HasConformalWeight N hw) (hzero : VacuumGradeZero W) (hE : EtingofFinite W) :
    ∃ (s : ℕ) (g : ℕ → ℕ → ℂ), 1 ≤ s ∧ (∀ i, Subexponential (fun k => ‖g i k‖)) ∧
      ∀ e : ℂ, (qDeriv^[s] (torusPartition N)) e
        + ∑ i ∈ Finset.range s, qConv (g i) (qDeriv^[i] (torusPartition N)) e = 0 := by

theorem character_formalODE_adjoint_of_etingofFinite (W : VertexOperatorAlgebra 𝒱)
    (hne : W.vacuum ≠ 0) (hzero : VacuumGradeZero W) (hE : EtingofFinite W) :
    ∃ (s : ℕ) (g : ℕ → ℕ → ℂ), 1 ≤ s ∧ (∀ i, Subexponential (fun k => ‖g i k‖)) ∧
      ∀ e : ℂ, (qDeriv^[s] (torusPartition (adjointModule W))) e
        + ∑ i ∈ Finset.range s,
            qConv (g i) (qDeriv^[i] (torusPartition (adjointModule W))) e = 0 :=
\end{leancode}
\noindent\textbf{Hypotheses / status.} \lean{hcw}, \lean{hE}\par
\medskip
\begin{bdefinition}{16.25}[Monic MLDE on $\mathbb H$]
For $f:\mathbb H\to\C$, $f$ satisfies a monic MLDE of order $s\ge1$ and weight $w$ with coefficients $p_0,\dots,p_{s-1}$
if each $p_i$ is a modular form of weight $2(s-i)$ and $\vartheta^s_{(w)}f + \sum_{i<s}p_i\vartheta^i_{(w)}f = 0$ on $\mathbb H$.
\end{bdefinition}
{\raggedright\noindent\textbf{Lean.} \lean{VOA.}\allowbreak\lean{serreIterC},\allowbreak\ \lean{VOA.}\allowbreak\lean{MonicMLDEOn}\par}
{\raggedright\noindent File \leanpath{VOA/\allowbreak Modularity/\allowbreak AKMLDE.lean}.\par}
\begin{leancode}
noncomputable def serreIterC (wt : ℤ) : ℕ → (ℂ → ℂ) → (ℂ → ℂ)
  | 0, f => f
  | i + 1, f => serreDerivC (wt + 2 * (i : ℤ)) (serreIterC wt i f)

def MonicMLDEOn (wt : ℤ) (f : ℂ → ℂ) (s : ℕ) (p : ℕ → E46) : Prop :=
  1 ≤ s ∧
    (∀ i < s, E46Weight (2 * ((s - i : ℕ) : ℤ)) (p i)) ∧
    ∀ z ∈ imGT 0, serreIterC wt s f z
      + ∑ i ∈ Finset.range s, e46C (p i) z * serreIterC wt i f z = 0
\end{leancode}
\noindent\textbf{Hypotheses / status.} Definitions; \lean{AKWeightGrading} inhabited for every VOA (\lean{akWeightGrading}); \lean{E46SerreClosed} proved (\lean{e46SerreClosed})\par
\medskip
\begin{bdefinition}{16.26}[AK weight grading]
Let $\lsixteenVA^{E} = \C[E_4,E_6]\otimes V$ and let its total-weight-$n$ subspace be spanned by $f\otimes v$ with $f$ of
modular weight $p$, $L[0]v = qv$, $p+q=n$. An \emph{AK weight grading} consists of: a direct-sum decomposition of $\C[E_4,E_6]$
into its weight spaces; a direct-sum decomposition of $\lsixteenVA^E$ into total-weight spaces; and the statement that the
relation module $O^E_q(V)$ (\lean{blockRelSpanE W}, the $E_4,E_6$-version of $\lsixteenOq$) is homogeneous.
\end{bdefinition}
{\raggedright\noindent\textbf{Lean.} \lean{VOA.}\allowbreak\lean{e46WeightSpace},\allowbreak\ \lean{VOA.}\allowbreak\lean{akTotalWeightSpace},\allowbreak\ \lean{VOA.}\allowbreak\lean{AKWeightGrading}\par}
{\raggedright\noindent File \leanpath{VOA/\allowbreak Modularity/\allowbreak AKGrading.lean}.\par}
\begin{leancode}
noncomputable def e46WeightSpace (k : ℤ) : Submodule ℂ E46 where
  carrier := {f | E46Weight k f}

noncomputable def akTotalWeightSpace (W : VertexOperatorAlgebra 𝒱) (n : ℤ) :
    Submodule ℂ (𝒱E 𝒱) :=
  Submodule.span ℂ {x | ∃ (p q : ℤ) (f : E46) (v : 𝒱), p + q = n ∧
    f ∈ e46WeightSpace p ∧ v ∈ sqWeightSpace W (q : ℂ) ∧ x = f ⊗ₜ[ℂ] v}

structure AKWeightGrading (W : VertexOperatorAlgebra 𝒱) where
  coeffDecomposition : DirectSum.Decomposition e46WeightSpace
  totalDecomposition : DirectSum.Decomposition (akTotalWeightSpace W)
  relationHomogeneous :
    letI := totalDecomposition
    DirectSum.SetLike.IsHomogeneous (akTotalWeightSpace W)
      ((blockRelSpanE W).restrictScalars ℂ)
\end{leancode}
\noindent\textbf{Hypotheses / status.} Definitions; \lean{AKWeightGrading} inhabited for every VOA (\lean{akWeightGrading}); \lean{E46SerreClosed} proved (\lean{e46SerreClosed})\par
\medskip
\begin{btheorem}{16.29}[AK Theorem 5.1, conditional form]
Let $V$ be Etingof-finite, $a\in V$ square-bracket primary of integer weight $w$, and assume an AK weight grading for $V$
and \lean{E46SerreClosed}. Then there are $s\ge1$ and $p_0,\dots,p_{s-1}$ such that for \emph{every} analytic genus-one
block $F$ the function $F(a)$ satisfies the monic MLDE of Definition~16.25 of weight $w$ with these
coefficients. Both side hypotheses are theorems: \lean{AKWeightGrading W} is \lean{akWeightGrading W}
(Theorem~16.27) and \lean{E46SerreClosed} is \lean{e46SerreClosed};
\lean{IsSqPrimary.}\allowbreak\lean{exists_}\allowbreak\lean{monicMLDEOn_}\allowbreak\lean{all_analyticBlocks_}\allowbreak\lean{of_etingofFinite'}
discharges the second and Theorem~16.30 below discharges both.
\end{btheorem}
{\raggedright\noindent\textbf{Lean.} \lean{VOA.}\allowbreak\lean{IsSqPrimary.}\allowbreak\lean{exists_monicMLDEOn_all_analyticBlocks_of_etingofFinite}\par}
{\raggedright\noindent File \leanpath{VOA/\allowbreak Modularity/\allowbreak AKMLDE.lean}.\par}
\begin{leancode}
theorem IsSqPrimary.exists_monicMLDEOn_all_analyticBlocks_of_etingofFinite
    {W : VertexOperatorAlgebra 𝒱} {a : 𝒱} {wt : ℤ}
    (ha : IsSqPrimary W a (wt : ℂ)) (hgrade : AKWeightGrading W)
    (hclose : E46SerreClosed) (hzero : VacuumGradeZero W) (hE : EtingofFinite W) :
    ∃ (s : ℕ) (p : ℕ → E46), 1 ≤ s ∧
      ∀ (F : 𝒱 → ℂ → ℂ), IsAnalyticBlock W F → MonicMLDEOn wt (F a) s p := by
\end{leancode}
\noindent\textbf{Hypotheses / status.} \lean{hE}, \lean{hgrade}, \lean{hclose}; \lean{hclose} discharged in the primed form\par
\medskip
\begin{btheorem}{16.30}[AK Theorem 5.1 under Etingof-finiteness; the $L_{-4/3}(\lsixteensl)$ instance]
(i) Let $V$ be Etingof-finite and $a\in V$ square-bracket primary of integer weight $w$. Then there are $s\ge1$ and
coefficients $p_i\in\C[E_4,E_6]$, modular of weight $2(s-i)$, such that for every analytic genus-one block $F$ the
function $F(a)$ satisfies $\vartheta^{(s)}F(a)+\sum_{i<s}p_i\vartheta^{(i)}F(a)=0$ on $\mathbb H$ (iterated Serre
derivatives). (ii) For $V=L_{-4/3}(\lsixteensl)$ and $a=\lsixteenvac$ this holds with no hypothesis.
\end{btheorem}
{\raggedright\noindent\textbf{Lean.} \lean{VOA.}\allowbreak\lean{IsSqPrimary.}\allowbreak\lean{exists_monicMLDEOn_all_analyticBlocks_of_etingofFinite_uncond},\allowbreak\ \lean{VOA.}\allowbreak\lean{exists_monicMLDEOn_all_analyticBlocks_simpleSl2QL}\par}
{\raggedright\noindent File \leanpath{VOA/\allowbreak Modularity/\allowbreak AKMLDEUnconditional.lean}.\par}
\begin{leancode}
theorem IsSqPrimary.exists_monicMLDEOn_all_analyticBlocks_of_etingofFinite_uncond
    {W : VertexOperatorAlgebra 𝒱} {a : 𝒱} {wt : ℤ}
    (ha : IsSqPrimary W a (wt : ℂ)) (hE : EtingofFinite W) :
    ∃ (s : ℕ) (p : ℕ → E46), 1 ≤ s ∧
      ∀ (F : 𝒱 → ℂ → ℂ), IsAnalyticBlock W F → MonicMLDEOn wt (F a) s p :=
  ha.exists_monicMLDEOn_all_analyticBlocks_of_etingofFinite' (akWeightGrading W)
    (VOA.QuasiLisseSl2.vacuumGradeZero_of_vacuum_grade W) hE

theorem exists_monicMLDEOn_all_analyticBlocks_simpleSl2QL :
    ∃ (s : ℕ) (p : ℕ → E46), 1 ≤ s ∧
      ∀ F, IsAnalyticBlock simpleSl2QL F →
        MonicMLDEOn 0 (F simpleSl2QL.vacuum) s p :=
  IsSqPrimary.exists_monicMLDEOn_all_analyticBlocks_of_etingofFinite_uncond
    (wt := 0) (by exact_mod_cast vacuum_isSqPrimary simpleSl2QL) VOA.Sl2Pres.etingofFinite_simpleSl2QL'
\end{leancode}
\noindent\textbf{Hypotheses / status.} \lean{hE} only (despite the name); order unbounded\par
\medskip
\begin{btheorem}{16.32}[Zhu modularity under Etingof-finiteness, conditional]
Let $V$ be Etingof-finite and $(M_i)_{i\in\iota}$ a finite family of irreducible modules. If the trace functions are
independent on each $T$-eigenvalue fibre (\lean{hfib}) and the restricted analytic blocks together with the $S$-slashes of
the traces span a space of dimension at most $|\iota|$ (\lean{hrank}), then there is an $SL_2(\Z)$-representation
$\rho$ with $Z_{M_i}(v;A\tau) = (c\tau+d)^n\sum_j\rho(A)_{ij}Z_{M_j}(v;\tau)$ for all $v$ of $L[0]$-weight $n$.
\end{btheorem}
{\raggedright\noindent\textbf{Lean.} \lean{VOA.}\allowbreak\lean{exists_zhuModularity_of_sqAnalyticRank_le_irr_etingofFinite}\par}
{\raggedright\noindent File \leanpath{VOA/\allowbreak Modularity/\allowbreak AKZhuModularity.lean}.\par}
\begin{leancode}
theorem exists_zhuModularity_of_sqAnalyticRank_le_irr_etingofFinite
    (hE : EtingofFinite W)
    (hfib : FibreIndependent Mfam (fun i => conformalWeightOf (Mfam i)))
    (hrank : Module.rank ℂ
        (sqAnalyticSlashAmbient Mfam (fun i => conformalWeightOf (Mfam i)))
      ≤ (Fintype.card ι : Cardinal)) :
    ∃ D : ModularData W ι, ZhuModularityAnalytic D :=
\end{leancode}
\noindent\textbf{Hypotheses / status.} \lean{hE}, \lean{hfib}, \lean{hrank} (discharged only for the trivial VOA; \lean{hrank} expected \textbf{false} at $L_{-4/3}$, paper argument)\par
\medskip
\begin{btheorem}{16.38}[The presentation of $R_{L_{-4/3}}$ is constructed]
(i) For every non-critical $k$, the algebra map $\operatorname{Sym}(\lsixteensl)\to R_{V^k(\lsixteensl)}$, $x\mapsto[x_{-1}\lsixteenvac]$, is
surjective. (ii) In $R_{L_{-4/3}}$, $[e]\cdot\lsixteenCas = 0$, and the whole ideal $\lsixteenCas\cdot\mathfrak m$ maps to $0$. (iii)
These maps form an $\lsixteensl$ symbol presentation of $L_{-4/3}(\lsixteensl)$. (iv) Hence $L_{-4/3}(\lsixteensl)$ is
Etingof-finite.
\end{btheorem}
{\raggedright\noindent\textbf{Lean.} \lean{VOA.}\allowbreak\lean{Sl2Pres.}\allowbreak\lean{symbolAlg},\allowbreak\ \lean{VOA.}\allowbreak\lean{Sl2Pres.}\allowbreak\lean{symbolAlg_surjective},\allowbreak\ \lean{VOA.}\allowbreak\lean{Sl2Pres.}\allowbreak\lean{symbolAlgL},\allowbreak\ \lean{VOA.}\allowbreak\lean{Sl2Pres.}\allowbreak\lean{eL_mul_casL},\allowbreak\ \lean{VOA.}\allowbreak\lean{Sl2Pres.}\allowbreak\lean{cas_vanishes_L},\allowbreak\ \lean{VOA.}\allowbreak\lean{Sl2Pres.}\allowbreak\lean{sl2PresentationSimpleSl2QL},\allowbreak\ \lean{VOA.}\allowbreak\lean{Sl2Pres.}\allowbreak\lean{etingofFinite_simpleSl2QL'}\par}
{\raggedright\noindent File \leanpath{VOA/\allowbreak Examples/\allowbreak AffineSl2Presentation.lean}.\par}
\begin{leancode}
def symbolAlg : SymSl2 →ₐ[ℂ] A₂ (affineSl2VOA_closed k hk) :=
  MvPolynomial.aeval ![eCl k hk, hCl k hk, fCl k hk]

theorem symbolAlg_surjective : Function.Surjective (symbolAlg k hk) := by

def symbolAlgL : SymSl2 →ₐ[ℂ] A₂ simpleSl2QL := MvPolynomial.aeval ![eL, hL, fL]

theorem eL_mul_casL : eL * casL = 0 := by

theorem cas_vanishes_L : ∀ p ∈ casIdeal, symbolAlgL p = 0 := by

def sl2PresentationSimpleSl2QL : Sl2Presentation simpleSl2QL where

theorem etingofFinite_simpleSl2QL' : EtingofFinite simpleSl2QL :=
  sl2PresentationSimpleSl2QL.etingofFinite
\end{leancode}
\noindent\textbf{Hypotheses / status.} \textbf{unconditional}, \lean{pbw}-free\par
\medskip
\begin{btheorem}{16.40}[Non-lisseness of $L_{-4/3}(\lsixteensl)$]
Let $E_r = e_{-1}^r\lsixteenvac$ and $F_r = f_{-1}^r\lsixteenvac$ in $V = V^{-4/3}(\lsixteensl)$, and $\langle\cdot,\cdot\rangle$ the
canonical invariant bilinear form. (i) $\langle F_r, x\rangle = 0$ for all $x\in C_2(V)$. (ii) $\langle F_r,E_r\rangle =
(-1)^r\,r!\,\prod_{j<r}(k-j)\ne0$ at $k=-4/3$. (iii) The class of $E_r$ in $R_{L}$, $L = L_{-4/3}(\lsixteensl)$, is nonzero,
and these classes are linearly independent. (iv) $L$ is not $C_2$-cofinite. (v) $L$ is not $C_2$-cofinite and is
Etingof-finite.
\end{btheorem}
{\raggedright\noindent\textbf{Lean.} \lean{VOA.}\allowbreak\lean{Sl2NonLisse.}\allowbreak\lean{form_fHarmonic_O₂_zero},\allowbreak\ \lean{VOA.}\allowbreak\lean{Sl2NonLisse.}\allowbreak\lean{form_fHarmonic_eHarmonic_ne_zero},\allowbreak\ \lean{VOA.}\allowbreak\lean{Sl2NonLisse.}\allowbreak\lean{eHarmonicL_not_mem_O₂},\allowbreak\ \lean{VOA.}\allowbreak\lean{Sl2NonLisse.}\allowbreak\lean{eHarmonicL_C₂_linearIndependent},\allowbreak\ \lean{VOA.}\allowbreak\lean{Sl2NonLisse.}\allowbreak\lean{not_C2Cofinite_simpleSl2QL},\allowbreak\ \lean{VOA.}\allowbreak\lean{Sl2NonLisse.}\allowbreak\lean{strictEtingofFinite_simpleSl2QL}\par}
{\raggedright\noindent File \leanpath{VOA/\allowbreak Examples/\allowbreak AffineSl2NonLisse.lean}.\par}
\begin{leancode}
theorem form_fHarmonic_O₂_zero (r : ℕ) {x : WakModule kk} (hx : x ∈ O₂ W43) :
    LQ43.formAdm.form (fHarmonic kk r) x = 0 := by

theorem form_fHarmonic_eHarmonic_ne_zero (r : ℕ) :
    LQ43.formAdm.form (fHarmonic kk r) (eHarmonic kk r) ≠ 0 := by

theorem eHarmonicL_not_mem_O₂ (r : ℕ) : eHarmonicL r ∉ O₂ L43 := by

theorem eHarmonicL_C₂_linearIndependent :
    LinearIndependent ℂ (fun r : ℕ => mk₂ L43 (eHarmonicL r)) := by

theorem not_C2Cofinite_simpleSl2QL : ¬ C2Cofinite L43 := by

theorem strictEtingofFinite_simpleSl2QL :
    ¬ C2Cofinite L43 ∧ EtingofFinite L43 :=
  ⟨not_C2Cofinite_simpleSl2QL, Sl2Pres.etingofFinite_simpleSl2QL'⟩
\end{leancode}
\noindent\textbf{Hypotheses / status.} \textbf{unconditional}, \lean{pbw}-free\par
\medskip
\endgroup
\part{Super, logarithmic and unitary theories}
\setcounter{chapter}{16}
\chapter{The super generalization}
\begingroup
\providecommand{\lseventeenvac}{\mathbf{1}}
\providecommand{\lseventeenV}{\mathcal{V}}
\providecommand{\lseventeeneps}{\varepsilon}
\providecommand{\lseventeenpar}[1]{|#1|}
\providecommand{\lseventeenGt}{\Gamma_\theta}
\providecommand{\lseventeenstr}{\operatorname{str}}
\renewcommand{\arraystretch}{1.3}
\renewcommand{\arraystretch}{1.25}
\noindent{\small 9 of the chapter's 58 numbered statements: its headline results and every definition they use.}\par\medskip
\begin{btheorem}{17.4}[No translation covariance]
There is a three-dimensional super vertex algebra in the library's sense (purely even, weights $0$
and $2$) on which the super Jacobi identity fails at $(u,v,x,\ell,m,n)=(x_1,\lseventeenvac,\lseventeenvac,-2,0,-2)$
and on which $Y(Tu,z)\neq\partial_zY(u,z)$.
\end{btheorem}
\noindent\emph{(The chapter quotes no Lean declaration for this statement; see book p.~847.)}\par
\medskip
\begin{btheorem}{17.14}[Super Borcherds / Jacobi identity]
Let $V$ be a super near-VOA, $u\in V_{\bar p}$, $v\in V_{\bar q}$, $x\in V$. For all $\ell,m,n\in\Z$,
\begin{equation*}
\sum_{i\ge0}(-1)^i\binom{\ell}{i}\Bigl(u_{(\ell+m-i)}v_{(n+i)}x
-(-1)^{pq}(-1)^\ell\,v_{(\ell+n-i)}u_{(m+i)}x\Bigr)
=\sum_{i\ge0}\binom mi\,(u_{(\ell+i)}v)_{(m+n-i)}x .
\end{equation*}
\end{btheorem}
{\raggedright\noindent\textbf{Lean.} \lean{VOA.}\allowbreak\lean{SuperNearVOA.}\allowbreak\lean{superBorcherds},\allowbreak\ \lean{VOA.}\allowbreak\lean{SuperNearVOA.}\allowbreak\lean{superJacobi}\par}
{\raggedright\noindent File \leanpath{VOA/\allowbreak Core/\allowbreak SuperBorcherds.lean}.\par}
\begin{leancode}
theorem superBorcherds (V : SuperNearVOA 𝒱) {p q : ZMod 2} {u v : 𝒱} (hu : u ∈ V.parity p)
    (hv : v ∈ V.parity q) (x : 𝒱) (ℓ m n : ℤ) :
    BLε V (koszul p q) u v x ℓ m n = BRs V u v x ℓ m n := by
...
theorem superJacobi (V : SuperNearVOA 𝒱) {p q : ZMod 2} {u v : 𝒱} (hu : u ∈ V.parity p)
    (hv : v ∈ V.parity q) (x : 𝒱) (ℓ m n : ℤ) :
    (∑ᶠ i : ℤ, (negOnePow i * cbz ℓ i) •
        (V.smul u (ℓ + m - i) (V.smul v (n + i) x)
          - (koszul p q * negOnePow ℓ) • V.smul v (ℓ + n - i) (V.smul u (m + i) x)))
      = ∑ᶠ i : ℤ, cbz m i • V.smul (V.smul u (ℓ + i) v) (m + n - i) x := by
\end{leancode}
\medskip
\begin{bproposition}{17.21}[The adjoint super module]
Every super near-VOA is a super weak module over itself, with $Y_M=Y$ and its own parity.
\end{bproposition}
{\raggedright\noindent\textbf{Lean.} \lean{VOA.}\allowbreak\lean{SuperNearVOA.}\allowbreak\lean{adjointSuperWeakModule}\par}
{\raggedright\noindent File \leanpath{VOA/\allowbreak Modules/\allowbreak SuperAdjointModule.lean}.\par}
\begin{leancode}
def SuperNearVOA.adjointSuperWeakModule (V : SuperNearVOA 𝒱) :
    SuperWeakModule V.toSuperVertexAlgebra 𝒱 where
  YM := V.Y
\end{leancode}
\medskip
\begin{btheorem}{17.24}[Twisted super modules with $N\ge2$ are zero]
Let $W$ be any super VOA, $g$ any even automorphism, $N\ge2$, and $T$ any
\lean{SuperTwistedModule W g N M}. Then every $x\in M$ is $0$. In particular $M$ is a
subsingleton, every weight space of $T$ is zero, and every parity-twisted (Ramond) carrier
\lean{W.ParityTwistedModule M} is zero.
\end{btheorem}
\noindent\emph{(The chapter quotes no Lean declaration for this statement; see book p.~861.)}\par
\medskip
\begin{btheorem}{17.27}[$\psi$ is not a field]
$\neg\,\mathrm{IsField}(\psi)$: there is no $B$ with $\psi(n)\lseventeenvac=0$ for all $n<B$.
\end{btheorem}
\noindent\emph{(The chapter quotes no Lean declaration for this statement; see book p.~864.)}\par
\medskip
\begin{btheorem}{17.38}[The super Zhu ring of symplectic fermions is not semisimple]
The super Zhu algebra $A(SF)$ (Definition~17.42) is an associative ring of dimension
at most $4$, $[\chi^a]\neq0$, and the two-sided ideal generated by $[\chi^0]$ is nonzero and
square-zero; hence $A(SF)$ is not a semisimple ring.
\end{btheorem}
{\raggedright\noindent\textbf{Lean.} \lean{VOA.}\allowbreak\lean{SympFerm.}\allowbreak\lean{finrank_ASF_le_four},\allowbreak\ \lean{VOA.}\allowbreak\lean{SympFerm.}\allowbreak\lean{chiClass_ne_zero},\allowbreak\ \lean{VOA.}\allowbreak\lean{SympFerm.}\allowbreak\lean{ASF_not_isSemisimpleRing}\par}
{\raggedright\noindent File \leanpath{VOA/\allowbreak Examples/\allowbreak SympFermSuperZhu.lean}.\par}
\begin{leancode}
theorem finrank_ASF_le_four : Module.finrank ℂ ASF ≤ 4 := by
...
theorem chiClass_ne_zero (a : Fin 2) : sfN.superMk hSF (chiStateSF a) ≠ 0 :=
...
theorem ASF_not_isSemisimpleRing :
    letI := instRingASF
    ¬ IsSemisimpleRing ASF :=
\end{leancode}
\noindent\textbf{Hypotheses / status.} Unconditional\par
\medskip
\begin{bdefinition}{17.42}[Super Zhu algebra]
Let $V$ be a super near-VOA with integer weights. For $u\in V_k$ put
$u\ast v=\sum_{m=0}^{k}\binom km u_{(m-1)}v$, extended linearly; $O(V)=\operatorname{span}\{(L_{-1}u+L_0u)\ast v\}$;
$A(V)=V/O(V)$.
\end{bdefinition}
{\raggedright\noindent\textbf{Lean.} \lean{VOA.}\allowbreak\lean{SuperNearVOA.}\allowbreak\lean{superZhuFamily},\allowbreak\ \lean{VOA.}\allowbreak\lean{SuperNearVOA.}\allowbreak\lean{superO},\allowbreak\ \lean{VOA.}\allowbreak\lean{SuperNearVOA.}\allowbreak\lean{superA}\par}
{\raggedright\noindent File \leanpath{VOA/\allowbreak Modules/\allowbreak SuperZhuRing.lean}.\par}
\begin{leancode}
noncomputable def superZhuFamily (v : 𝒱) (k : ℤ) : (V.gradingZ k) →ₗ[ℂ] 𝒱 where
  toFun uk := ∑ m ∈ Finset.range (k.toNat + 1), cbz k m • V.smul (uk : 𝒱) ((m : ℤ) - 1) v
...
def superO (hV : V.IntegralWeights) : Submodule ℂ 𝒱 :=
  Submodule.span ℂ {x | ∃ u v : 𝒱, x = V.superZhuMul hV (V.L (-1) u + V.L 0 u) v}
...
abbrev superA (hV : V.IntegralWeights) : Type u := 𝒱 ⧸ V.superO hV
\end{leancode}
\medskip
\begin{btheorem}{17.45}[Twisted Zhu finiteness]
If $W$ is integrally graded and super $C_2$-cofinite, $g^N=1$ with $N>0$, and the eigenspace
decomposition of $g$ is internal and compatible with weight and parity, then $A_g(W)$ is
finite-dimensional.
\end{btheorem}
{\raggedright\noindent\textbf{Lean.} \lean{VOA.}\allowbreak\lean{finiteDimensional_twistedA_of_superC2Cofinite}\par}
{\raggedright\noindent File \leanpath{VOA/\allowbreak Modules/\allowbreak SuperTwistedZhuFinite.lean}.\par}
\begin{leancode}
theorem finiteDimensional_twistedA_of_superC2Cofinite
    (W : SuperVertexOperatorAlgebra 𝒱) (hW : W.toSuperNearVOA.IntegralWeights)
    (g : SuperAut W) {N : ℕ} (hN : 0 < N) (_hpow : g ^ N = 1)
    (hE : TwistEigenspaceDecomposition W g N) (hC2 : SuperC2Cofinite W) :
    FiniteDimensional ℂ (W.twistedA hW g N) := by
\end{leancode}
\medskip
\begin{btheorem}{17.58}[Finite square-spin cusp quotient]
Let $W$ be a half-integrally graded super VOA ($W_h=0$ unless $h\in\frac12\Z$), super $C_2$-cofinite, and
$(g,h)$ a commuting pair of involutive even automorphisms. Then the quotient of $W$ by the
square-spin cusp relation span of the $(+,+)$ inner eigensector of $(g,h)$ is finite-dimensional.
\end{btheorem}
{\raggedright\noindent\textbf{Lean.} \lean{VOA.}\allowbreak\lean{SuperNearVOA.}\allowbreak\lean{BoundaryPair.}\allowbreak\lean{finiteDimensional_superSpinSqCuspQuotient_of_superC2},\allowbreak\ \lean{VOA.}\allowbreak\lean{BoundaryPair}\par}
{\raggedright\noindent File \leanpath{VOA/\allowbreak Modularity/\allowbreak SuperSpinDescentSqFinite.lean}.\par}
\begin{leancode}
theorem BoundaryPair.finiteDimensional_superSpinSqCuspQuotient_of_superC2
    (b : BoundaryPair W) (hhalf : IsHalfIntegerGraded W)
    (hspatial : b.spatial ^ 2 = 1) (htemporal : b.temporal ^ 2 = 1)
    (hC2 : SuperC2Cofinite W) :
    FiniteDimensional Complex
      (VType ⧸ superSpinSqCuspRelSpan W hhalf b.spinEigenspace) :=
\end{leancode}
{\raggedright\noindent File \leanpath{VOA/\allowbreak Modularity/\allowbreak SuperSectorSystem.lean}.\par}
\begin{leancode}
structure BoundaryPair (W : SuperVertexOperatorAlgebra VType) where
  spatial : SuperAut W
  temporal : SuperAut W
  commute : Commute spatial temporal
\end{leancode}
\medskip
\endgroup
\setcounter{chapter}{17}
\chapter{Logarithmic vertex algebras: the triplet and doublet algebras}
\begingroup
\providecommand{\leighteenvac}{\mathbf{1}}
\providecommand{\leighteenV}{\mathcal{V}}
\providecommand{\leighteenW}{\mathcal{W}}
\providecommand{\leighteenA}{\mathcal{A}}
\providecommand{\leighteenVL}{V_{L_p}}
\providecommand{\leighteenQ}{Q}
\providecommand{\leighteenO}{O}
\providecommand{\leighteenAtwo}{A_2}
\providecommand{\leighteenbar}[1]{\overline{#1}}
\providecommand{\leighteenSF}{\mathrm{SF}}
\providecommand{\leighteenHS}{T}
\noindent{\small 12 of the chapter's 43 numbered statements: its headline results and every definition they use.}\par\medskip
\begin{btheorem}{18.6}[Central charge; non-degeneracy; consistency with the Abe model]
$W(p)$ has central charge $c_{1,p}$; it contains $\leighteenvac,\omega,E,H,F$; it is nonzero; and at
$p=2$ its central charge equals that of \lean{SympFerm.tripletW2}.
\end{btheorem}
{\raggedright\noindent\textbf{Lean.} \lean{VOA.}\allowbreak\lean{LatGen.}\allowbreak\lean{Wp.}\allowbreak\lean{tripletW_centralCharge},\allowbreak\ \lean{VOA.}\allowbreak\lean{LatGen.}\allowbreak\lean{Wp.}\allowbreak\lean{tripletW_carrier_ne_bot},\allowbreak\ \lean{VOA.}\allowbreak\lean{LatGen.}\allowbreak\lean{Wp.}\allowbreak\lean{omegaTrip_ne_zero},\allowbreak\ \lean{VOA.}\allowbreak\lean{LatGen.}\allowbreak\lean{Wp.}\allowbreak\lean{tripletW_centralCharge_eq_tripletW2}\par}
{\raggedright\noindent File \leanpath{VOA/\allowbreak Examples/\allowbreak TripletWp.lean}.\par}
\begin{leancode}
@[simp] theorem tripletW_centralCharge (p : ℕ) (hp : 1 ≤ p) :
    (tripletWVOA p hp).centralCharge = 1 - 6 * ((p : ℂ) - 1) ^ 2 / (p : ℂ) := rfl

theorem tripletW_carrier_ne_bot (p : ℕ) (hp : 1 ≤ p) :
    (tripletW p hp).carrier ≠ ⊥ := by

theorem omegaTrip_ne_zero (p : ℕ) (hp : 1 ≤ p) : omegaTrip p ≠ 0 := by

theorem tripletW_centralCharge_eq_tripletW2 (hp : 1 ≤ 2) :
    (tripletWVOA 2 hp).centralCharge = SympFerm.tripletW2.centralCharge := by
\end{leancode}
\medskip
\begin{btheorem}{18.7}[$E\ne0$ and $W(2)\ne V_{L_2}$; Lean check outside the library]
At $p=2$, $E\ne0$, and $e^\alpha\notin W(2)$; hence $W(2)$ is a proper subalgebra of $V_{L_2}$.
\end{btheorem}
\noindent\emph{(The chapter quotes no Lean declaration for this statement; see book p.~909.)}\par
\noindent\textbf{Hypotheses / status.} None\par
\medskip
\begin{btheorem}{18.10}[Adamovi\'c--Milas relations]
For every $p\ge1$ and every $n\ge-2p$, in $V_{L_p}$:
\begin{equation*}
F_{(n)}F=0,\qquad E_{(n)}E=0,\qquad E_{(n)}F+F_{(n)}E+2H_{(n)}H=0,
\end{equation*}
and $Q$ is a derivation of all products: $Q(u_{(n)}v)=(Qu)_{(n)}v+u_{(n)}(Qv)$. Consequently, in
$R_{W(p)}$,
\begin{equation*}
\bar E^2=0,\qquad \bar F^2=0,\qquad \bar H^2=-\bar E\bar F,\qquad \bar H^4=0 .
\end{equation*}
\end{btheorem}
{\raggedright\noindent\textbf{Lean.} \lean{VOA.}\allowbreak\lean{LatGen.}\allowbreak\lean{WpC2.}\allowbreak\lean{smul_Ftrip_Ftrip},\allowbreak\ \lean{VOA.}\allowbreak\lean{LatGen.}\allowbreak\lean{WpC2.}\allowbreak\lean{Qwp_leibniz},\allowbreak\ \lean{VOA.}\allowbreak\lean{LatGen.}\allowbreak\lean{WpC2.}\allowbreak\lean{EF_add_FE_add_two_HH},\allowbreak\ \lean{VOA.}\allowbreak\lean{LatGen.}\allowbreak\lean{WpC2.}\allowbreak\lean{smul_Etrip_Etrip},\allowbreak\ \lean{VOA.}\allowbreak\lean{LatGen.}\allowbreak\lean{WpC2.}\allowbreak\lean{genWp},\allowbreak\ \lean{VOA.}\allowbreak\lean{LatGen.}\allowbreak\lean{WpC2.}\allowbreak\lean{mk₂_genE_sq},\allowbreak\ \lean{VOA.}\allowbreak\lean{LatGen.}\allowbreak\lean{WpC2.}\allowbreak\lean{mk₂_genF_sq},\allowbreak\ \lean{VOA.}\allowbreak\lean{LatGen.}\allowbreak\lean{WpC2.}\allowbreak\lean{mk₂_genH_sq},\allowbreak\ \lean{VOA.}\allowbreak\lean{LatGen.}\allowbreak\lean{WpC2.}\allowbreak\lean{mk₂_genH_pow_four}\par}
{\raggedright\noindent File \leanpath{VOA/\allowbreak Examples/\allowbreak TripletWpC2.lean}.\par}
\begin{leancode}
theorem smul_Ftrip_Ftrip (p : ℕ) (hp : 1 ≤ p) {n : ℤ} (hn : -(2 * (p : ℤ)) ≤ n) :
    (tripletLatticeVOA p hp).smul (Ftrip p) n (Ftrip p) = 0 := by

theorem Qwp_leibniz (p : ℕ) (hp : 1 ≤ p) (u : LatModuleGen (latP p)) (n : ℤ)
    (v : LatModuleGen (latP p)) :
    Qwp p hp ((tripletLatticeVOA p hp).smul u n v)
      = (tripletLatticeVOA p hp).smul (Qwp p hp u) n v
        + (tripletLatticeVOA p hp).smul u n (Qwp p hp v) :=

theorem EF_add_FE_add_two_HH (p : ℕ) (hp : 1 ≤ p) {n : ℤ} (hn : -(2 * (p : ℤ)) ≤ n) :
    (tripletLatticeVOA p hp).smul (Etrip p) n (Ftrip p)
      + (tripletLatticeVOA p hp).smul (Ftrip p) n (Etrip p)
      + (2 : ℂ) • (tripletLatticeVOA p hp).smul (Htrip p) n (Htrip p) = 0 := by

theorem smul_Etrip_Etrip (p : ℕ) (hp : 1 ≤ p) {n : ℤ} (hn : -(2 * (p : ℤ)) ≤ n) :
    (tripletLatticeVOA p hp).smul (Etrip p) n (Etrip p) = 0 := by
\end{leancode}
\begin{leancode}
noncomputable def genWp (p : ℕ) (hp : 1 ≤ p) : Fin 4 → (tripletW p hp).carrier
  | 0 => genOmega p hp
  | 1 => genE p hp
  | 2 => genH p hp
  | 3 => genF p hp

theorem mk₂_genE_sq (p : ℕ) (hp : 1 ≤ p) :
    (mk₂ (tripletWVOA p hp) (genE p hp)) ^ 2 = 0 := by

theorem mk₂_genF_sq (p : ℕ) (hp : 1 ≤ p) :
    (mk₂ (tripletWVOA p hp) (genF p hp)) ^ 2 = 0 := by

theorem mk₂_genH_sq (p : ℕ) (hp : 1 ≤ p) :
    (mk₂ (tripletWVOA p hp) (genH p hp)) ^ 2
      = -(mk₂ (tripletWVOA p hp) (genE p hp) * mk₂ (tripletWVOA p hp) (genF p hp)) := by

theorem mk₂_genH_pow_four (p : ℕ) (hp : 1 ≤ p) :
    (mk₂ (tripletWVOA p hp) (genH p hp)) ^ 4 = 0 := by
\end{leancode}
\noindent\textbf{Hypotheses / status.} None\par
\medskip
\begin{btheorem}{18.12}[$C_2$-cofiniteness of $W(p)$, conditional]
For every $p\ge1$ and $m\in\N$: if $R_{W(p)}$ is generated by $\bar\omega,\bar E,\bar H,\bar F$ and
$L_{-2}^m\leighteenvac\in C_2(W(p))$, then $W(p)$ is $C_2$-cofinite and $\dim R_{W(p)}\le16m$.
\end{btheorem}
{\raggedright\noindent\textbf{Lean.} \lean{VOA.}\allowbreak\lean{LatGen.}\allowbreak\lean{WpC2.}\allowbreak\lean{c2Cofinite_tripletWVOA},\allowbreak\ \lean{VOA.}\allowbreak\lean{LatGen.}\allowbreak\lean{WpC2.}\allowbreak\lean{c2Cofinite_tripletWVOA_of_omegaPowMem}\par}
{\raggedright\noindent File \leanpath{VOA/\allowbreak Examples/\allowbreak TripletWpC2.lean}.\par}
\begin{leancode}
theorem c2Cofinite_tripletWVOA (p : ℕ) (hp : 1 ≤ p) (hgen : AdjoinGenWp p hp)
    (homega : OmegaNilpotentWp p hp) : C2Cofinite (tripletWVOA p hp) := by

theorem c2Cofinite_tripletWVOA_of_omegaPowMem (p : ℕ) (hp : 1 ≤ p) (m : ℕ)
    (hgen : AdjoinGenWp p hp) (homega : OmegaPowMem p hp m) :
    C2Cofinite (tripletWVOA p hp)
      ∧ Module.finrank ℂ (A₂ (tripletWVOA p hp)) ≤ 16 * m :=
\end{leancode}
\noindent\textbf{Hypotheses / status.} \lean{AdjoinGenWp}, \lean{OmegaPowMem}\par
\medskip
\begin{btheorem}{18.14}[Generation from the generator OPE table]
Let $T_p$ be the span of the strictly-negative-mode words in $\omega,E,H,F$. If $g_{(j)}h\in T_p$ for
all generators $g,h$ and $0\le j<4p+2$, then $W(p)$ is strongly generated by $\omega,E,H,F$. The
$\omega$-row of the table holds for every $p$, because $E,H,F$ are Virasoro primary. The remaining
mixed pairs reduce to the two ``lowering rows'' $F_{(j)}H$ and $F_{(j)}E$.
\end{btheorem}
{\raggedright\noindent\textbf{Lean.} \lean{VOA.}\allowbreak\lean{LatGen.}\allowbreak\lean{WpC2.}\allowbreak\lean{GenPairClosedWp},\allowbreak\ \lean{VOA.}\allowbreak\lean{LatGen.}\allowbreak\lean{WpC2.}\allowbreak\lean{strongGenWp_of_genPairClosed},\allowbreak\ \lean{VOA.}\allowbreak\lean{LatGen.}\allowbreak\lean{WpC2.}\allowbreak\lean{LoweringPairsWp},\allowbreak\ \lean{VOA.}\allowbreak\lean{LatGen.}\allowbreak\lean{WpC2.}\allowbreak\lean{screenGens_primary},\allowbreak\ \lean{VOA.}\allowbreak\lean{LatGen.}\allowbreak\lean{WpC2.}\allowbreak\lean{c2Cofinite_tripletWVOA_of_lowering_AM}\par}
{\raggedright\noindent File \leanpath{VOA/\allowbreak Examples/\allowbreak TripletWpPairClosure.lean}.\par}
\begin{leancode}
def GenPairClosedWp (p : ℕ) (hp : 1 ≤ p) : Prop :=
  ∀ g ∈ tripGens p, ∀ h ∈ tripGens p, ∀ j : ℤ, 0 ≤ j →
    (tripletLatticeVOA p hp).smul g j h ∈ iterProdSpan (tripletLatticeVOA p hp) (tripGens p)

theorem strongGenWp_of_genPairClosed (p : ℕ) (hp : 1 ≤ p) (h : GenPairClosedWp p hp) :
    StrongGenWp p hp :=

def LoweringPairsWp (p : ℕ) (hp : 1 ≤ p) : Prop :=
  (∀ j : ℤ, 0 ≤ j →
      (tripletLatticeVOA p hp).smul (Ftrip p) j (Htrip p)
        ∈ iterProdSpan (tripletLatticeVOA p hp) (tripGens p))
    ∧
  (∀ j : ℤ, 0 ≤ j →
      (tripletLatticeVOA p hp).smul (Ftrip p) j (Etrip p)
        ∈ iterProdSpan (tripletLatticeVOA p hp) (tripGens p))

theorem screenGens_primary (p : ℕ) (hp : 1 ≤ p) :
    ∀ g ∈ screenGens p, ∀ n : ℤ, 1 ≤ n → (tripletLatticeVOA p hp).L n g = 0 := by

theorem c2Cofinite_tripletWVOA_of_lowering_AM (p : ℕ) (hp : 1 ≤ p)
    (hlower : LoweringPairsWp p hp) (hAM : AMOmegaIdentity p hp (3 * p - 1)) :
    C2Cofinite (tripletWVOA p hp)
      ∧ Module.finrank ℂ (A₂ (tripletWVOA p hp)) ≤ 16 * (3 * p - 1) :=
\end{leancode}
\noindent\textbf{Hypotheses / status.} \lean{LoweringPairsWp}, \lean{AMOmegaIdentity}\par
\medskip
\begin{btheorem}{18.15}[Mixed products at $p=2$]
With $c=\varepsilon(\alpha,-\alpha)$, $d=\varepsilon(-\alpha,\alpha)$, in $V_{L_2}$:
\begin{equation*}
F_{(0)}H=\tfrac{48c}{5}\,L_{-2}F+\tfrac{4c}{5}\,T^2F,\qquad
F_{(1)}E=-64cd\,L_{-2}\omega+12cd\,T^2\omega+10d\,TH,\qquad
F_{(3)}E=-48cd\,\omega .
\end{equation*}
Consequently all mixed generator pairs close, and $W(2)$ is strongly generated by $\omega,E,H,F$.
\end{btheorem}
{\raggedright\noindent\textbf{Lean.} \lean{VOA.}\allowbreak\lean{LatGen.}\allowbreak\lean{WpC2.}\allowbreak\lean{tripletTwoC},\allowbreak\ \lean{VOA.}\allowbreak\lean{LatGen.}\allowbreak\lean{WpC2.}\allowbreak\lean{tripletTwoD},\allowbreak\ \lean{VOA.}\allowbreak\lean{LatGen.}\allowbreak\lean{WpC2.}\allowbreak\lean{ThreeExplicitMixedIdentitiesWp},\allowbreak\ \lean{VOA.}\allowbreak\lean{LatGen.}\allowbreak\lean{WpC2.}\allowbreak\lean{Ftrip_three_Etrip_two_explicit},\allowbreak\ \lean{VOA.}\allowbreak\lean{LatGen.}\allowbreak\lean{WpC2.}\allowbreak\lean{threeCoordinateIdentitiesWp_two},\allowbreak\ \lean{VOA.}\allowbreak\lean{LatGen.}\allowbreak\lean{WpC2.}\allowbreak\lean{mixedPairsWp_two},\allowbreak\ \lean{VOA.}\allowbreak\lean{LatGen.}\allowbreak\lean{Wp.}\allowbreak\lean{strongGenWp_two}\par}
{\raggedright\noindent Files \leanpath{VOA/\allowbreak Examples/\allowbreak TripletWpTwoMixed.lean},
\leanpath{TripletWpTwoMixedClose.lean}, \leanpath{TripletWpTwoIrrational.lean}.\par}
\begin{leancode}
noncomputable abbrev tripletTwoC : ℂ :=
  latticeCocycle (latP 2) (alphaP 2) (-alphaP 2)

noncomputable abbrev tripletTwoD : ℂ :=
  latticeCocycle (latP 2) (-alphaP 2) (alphaP 2)

def ThreeExplicitMixedIdentitiesWp (hp : 1 ≤ 2) : Prop :=
  let W := tripletLatticeVOA 2 hp
  W.smul (Ftrip 2) 0 (Htrip 2) =
      ((48 : ℂ) / 5 * tripletTwoC) • W.smul (omegaTrip 2) (-1) (Ftrip 2) +
        ((4 : ℂ) / 5 * tripletTwoC) • (W.toPreVertexAlgebra.T ^ 2) (Ftrip 2) ∧
  W.smul (Ftrip 2) 1 (Etrip 2) =
      (-64 * (tripletTwoC * tripletTwoD)) • W.smul (omegaTrip 2) (-1) (omegaTrip 2) +
        (12 * (tripletTwoC * tripletTwoD)) • (W.toPreVertexAlgebra.T ^ 2) (omegaTrip 2) +
          (10 * tripletTwoD) • W.toPreVertexAlgebra.T (Htrip 2) ∧
  W.smul (Ftrip 2) 3 (Etrip 2) =
      (-48 * (tripletTwoC * tripletTwoD)) • omegaTrip 2
\end{leancode}
{\raggedright\noindent File \leanpath{VOA/\allowbreak Examples/\allowbreak TripletWpTwoMixedAssembly.lean}.\par}
\begin{leancode}
theorem Ftrip_three_Etrip_two_explicit (hp : 1 ≤ 2) :
    (tripletLatticeVOA 2 hp).smul (Ftrip 2) 3 (Etrip 2) =
      (-48 * (tripletTwoC * tripletTwoD)) • omegaTrip 2 := by

theorem threeCoordinateIdentitiesWp_two (hp : 1 ≤ 2) :
    ThreeCoordinateIdentitiesWp hp :=

theorem mixedPairsWp_two (hp : 1 ≤ 2) : MixedPairsWp 2 hp :=

theorem strongGenWp_two (hp : 1 ≤ 2) : StrongGenWp 2 hp :=
\end{leancode}
\noindent\textbf{Hypotheses / status.} None\par
\medskip
\begin{btheorem}{18.17}[The lattice triplet $W(2)$ is $C_2$-cofinite]
$W(2)=$ \lean{tripletWVOA 2 hp} is $C_2$-cofinite, and $\dim W(2)/C_2(W(2))\le80$.
\end{btheorem}
{\raggedright\noindent\textbf{Lean.} \lean{VOA.}\allowbreak\lean{LatGen.}\allowbreak\lean{WpC2.}\allowbreak\lean{tripletWVOA_two_c2Cofinite_and_finrank_le},\allowbreak\ \lean{VOA.}\allowbreak\lean{LatGen.}\allowbreak\lean{WpC2.}\allowbreak\lean{tripletWVOA_two_c2Cofinite}\par}
{\raggedright\noindent File \leanpath{VOA/\allowbreak Examples/\allowbreak TripletWpTwoClose.lean}.\par}
\begin{leancode}
theorem tripletWVOA_two_c2Cofinite_and_finrank_le (hp : 1 ≤ 2) :
    C2Cofinite (tripletWVOA 2 hp)
      ∧ Module.finrank ℂ (A₂ (tripletWVOA 2 hp)) ≤ 80 :=

theorem tripletWVOA_two_c2Cofinite (hp : 1 ≤ 2) : C2Cofinite (tripletWVOA 2 hp) :=
\end{leancode}
\noindent\textbf{Hypotheses / status.} None (lattice model only)\par
\medskip
\begin{btheorem}{18.18}[The lattice $W(2)$ is not rational]
$W(2)$ has no nonzero state of weight one. The ambient $V_{L_2}$, regarded as an ordinary $W(2)$-module by
restriction, is not completely reducible: the submodule $W(2)$ has no complement. Consequently $W(2)$ is
neither Gannon-rational nor admissibly rational.
\end{btheorem}
{\raggedright\noindent\textbf{Lean.} \lean{VOA.}\allowbreak\lean{LatGen.}\allowbreak\lean{Wp.}\allowbreak\lean{tripletAmbientOrdModuleTwo},\allowbreak\ \lean{VOA.}\allowbreak\lean{LatGen.}\allowbreak\lean{Wp.}\allowbreak\lean{tripletCarrierWeakSubmoduleTwo},\allowbreak\ \lean{VOA.}\allowbreak\lean{LatGen.}\allowbreak\lean{Wp.}\allowbreak\lean{eq_zero_of_mem_tripletW_two_and_grading_one},\allowbreak\ \lean{VOA.}\allowbreak\lean{LatGen.}\allowbreak\lean{Wp.}\allowbreak\lean{Ftrip_three_latChargeGen_alpha},\allowbreak\ \lean{VOA.}\allowbreak\lean{LatGen.}\allowbreak\lean{Wp.}\allowbreak\lean{tripletAmbientOrdModuleTwo_not_completelyReducible},\allowbreak\ \lean{VOA.}\allowbreak\lean{LatGen.}\allowbreak\lean{Wp.}\allowbreak\lean{not_gannonRational_tripletWVOA_two},\allowbreak\ \lean{VOA.}\allowbreak\lean{LatGen.}\allowbreak\lean{Wp.}\allowbreak\lean{not_admissiblyRational_tripletWVOA_two},\allowbreak\ \lean{VOA.}\allowbreak\lean{LatGen.}\allowbreak\lean{Wp.}\allowbreak\lean{tripletWVOA_two_c2Cofinite_and_not_admissiblyRational}\par}
{\raggedright\noindent Files \leanpath{VOA/\allowbreak Examples/\allowbreak TripletWpTwoFelderStart.lean},
\leanpath{TripletWpTwoIrrational.lean}.\par}
\begin{leancode}
noncomputable def tripletAmbientOrdModuleTwo (hp : 1 ≤ 2) :
    OrdVModule (tripletWVOA 2 hp) (LatModuleGen (latP 2)) :=
  (tripletW 2 hp).adjointRestrictOrd

noncomputable def tripletCarrierWeakSubmoduleTwo (hp : 1 ≤ 2) :
    WeakModule.Submodule' (tripletAmbientOrdModuleTwo hp).toWeakModule :=
  (tripletW 2 hp).carrierWeakSubmodule
\end{leancode}
\begin{leancode}
theorem eq_zero_of_mem_tripletW_two_and_grading_one (hp : 1 ≤ 2)
    {x : LatModuleGen (latP 2)} (hxW : x ∈ (tripletW 2 hp).carrier)
    (hx1 : x ∈ (tripletLatticeVOA 2 hp).grading 1) : x = 0 := by

theorem Ftrip_three_latChargeGen_alpha (hp : 1 ≤ 2) :
    (tripletLatticeVOA 2 hp).smul (Ftrip 2) 3
        (latChargeGen (latP 2) (alphaP 2))
      = latticeCocycle (latP 2) (-alphaP 2) (alphaP 2) •
          latChargeGen (latP 2) (-alphaP 2 + alphaP 2) := by

theorem tripletAmbientOrdModuleTwo_not_completelyReducible (hp : 1 ≤ 2) :
    ¬ OrdVModule.IsCompletelyReducible (tripletAmbientOrdModuleTwo hp) := by

theorem not_gannonRational_tripletWVOA_two (hp : 1 ≤ 2) :
    ¬ GannonRational (tripletWVOA 2 hp) := by

theorem not_admissiblyRational_tripletWVOA_two (hp : 1 ≤ 2) :
    ¬ AdmissiblyRational (tripletWVOA 2 hp) := by

theorem tripletWVOA_two_c2Cofinite_and_not_admissiblyRational (hp : 1 ≤ 2) :
    C2Cofinite (tripletWVOA 2 hp) ∧ ¬ AdmissiblyRational (tripletWVOA 2 hp) :=
\end{leancode}
\noindent\textbf{Hypotheses / status.} None\par
\medskip
\begin{btheorem}{18.20}[$W(2)$ is simple and strongly finite]
The lattice $W(2)$ satisfies the $f_2$ certificate; hence it is simple, strongly finite, and neither
admissibly nor Gannon rational (so not strongly rational).
\end{btheorem}
{\raggedright\noindent\textbf{Lean.} \lean{VOA.}\allowbreak\lean{LatGen.}\allowbreak\lean{Wp.}\allowbreak\lean{fTwoZhuCertificateWp},\allowbreak\ \lean{VOA.}\allowbreak\lean{LatGen.}\allowbreak\lean{Wp.}\allowbreak\lean{tripletWVOA_two_isSimple},\allowbreak\ \lean{VOA.}\allowbreak\lean{LatGen.}\allowbreak\lean{Wp.}\allowbreak\lean{tripletWVOA_two_stronglyFinite},\allowbreak\ \lean{VOA.}\allowbreak\lean{LatGen.}\allowbreak\lean{Wp.}\allowbreak\lean{tripletWVOA_two_stronglyFinite_and_irrational},\allowbreak\ \lean{VOA.}\allowbreak\lean{LatGen.}\allowbreak\lean{Wp.}\allowbreak\lean{tripletWVOA_two_stronglyFinite_of_isSimple},\allowbreak\ \lean{VOA.}\allowbreak\lean{LatGen.}\allowbreak\lean{Wp.}\allowbreak\lean{not_stronglyRational_tripletWVOA_two}\par}
{\raggedright\noindent Files \leanpath{VOA/\allowbreak Examples/\allowbreak TripletWpTwoZhuFinal.lean},
\leanpath{TripletWpStronglyFinite.lean}.\par}
\begin{leancode}
theorem fTwoZhuCertificateWp (hp : 1 ≤ 2) : FTwoZhuCertificateWp hp :=

theorem tripletWVOA_two_isSimple (hp : 1 ≤ 2) : (tripletWVOA 2 hp).IsSimple :=

theorem tripletWVOA_two_stronglyFinite (hp : 1 ≤ 2) :
    StronglyFinite (tripletWVOA 2 hp) :=

theorem tripletWVOA_two_stronglyFinite_and_irrational (hp : 1 ≤ 2) :
    StronglyFinite (tripletWVOA 2 hp) ∧
      ¬ AdmissiblyRational (tripletWVOA 2 hp) ∧
      ¬ GannonRational (tripletWVOA 2 hp) :=
\end{leancode}
\begin{leancode}
theorem tripletWVOA_two_stronglyFinite_of_isSimple (hp : 1 ≤ 2)
    (hsimple : (tripletWVOA 2 hp).IsSimple) : StronglyFinite (tripletWVOA 2 hp) := by

theorem not_stronglyRational_tripletWVOA_two (hp : 1 ≤ 2) :
    ¬ StronglyRational (tripletWVOA 2 hp) := by
\end{leancode}
\noindent\textbf{Hypotheses / status.} None\par
\medskip
\begin{bdefinition}{18.22}[The Abe model of $W(2)$]
$\leighteenSF^+$ is the even part of the symplectic-fermion super VOA \lean{sympFermSVOA}
(Chapter~17, Theorem~17.36): the span of the word states
$\chi^{a_1}_{-k_1-1}\cdots\chi^{a_n}_{-k_n-1}\leighteenvac$ with $n$ even, with the restricted vertex
operators, vacuum and conformal vector $\omega_{SF}=-\chi^0_{-1}\chi^1_{-1}\leighteenvac$.
\end{bdefinition}
{\raggedright\noindent\textbf{Lean.} \lean{VOA.}\allowbreak\lean{SympFerm.}\allowbreak\lean{sympFermSVOA_integralWeights},\allowbreak\ \lean{VOA.}\allowbreak\lean{SympFerm.}\allowbreak\lean{tripletW2},\allowbreak\ \lean{VOA.}\allowbreak\lean{SympFerm.}\allowbreak\lean{tripletW2_centralCharge},\allowbreak\ \lean{VOA.}\allowbreak\lean{SympFerm.}\allowbreak\lean{tripletW2_grading},\allowbreak\ \lean{VOA.}\allowbreak\lean{SympFerm.}\allowbreak\lean{evenPartSF_eq_span},\allowbreak\ \lean{VOA.}\allowbreak\lean{SympFerm.}\allowbreak\lean{tripletW2_grading_one}\par}
{\raggedright\noindent File \leanpath{VOA/\allowbreak Examples/\allowbreak TripletW2.lean}.\par}
\begin{leancode}
theorem sympFermSVOA_integralWeights (h : ℂ) (hh : ¬ ∃ n : ℤ, (n : ℂ) = h) :
    ∀ x ∈ sympFermSVOA.parity 0, x ∈ sympFermSVOA.grading h → x = 0 := by

noncomputable def tripletW2 : VertexOperatorAlgebra (evenOdd QSF 0) :=
  sympFermSVOA.evenPartVOA sympFermSVOA_integralWeights

@[simp] theorem tripletW2_centralCharge : tripletW2.centralCharge = -2 := rfl

@[simp] theorem tripletW2_grading (n : ℤ) :
    tripletW2.grading n = (gradingSF (n : ℂ)).comap (evenOdd QSF 0).subtype := rfl

theorem evenPartSF_eq_span :
    (evenOdd QSF (0 : ZMod 2) : Submodule ℂ FockSF)
      = Submodule.span ℂ (stWordSF '' {l : SFWord | Even l.length}) := by

theorem tripletW2_grading_one : tripletW2.grading 1 = ⊥ := by
\end{leancode}
\noindent\textbf{Hypotheses / status.} None\par
\medskip
\begin{btheorem}{18.25}[Abe's $C_2$ relations in the Abe model]
In $R_{W(2)}$, with $\bar\omega,\bar E,\bar H,\bar F$ the classes of the generators:
\begin{equation*}
\bar E^2=\bar F^2=0,\quad \bar E\bar H=\bar F\bar H=0,\quad 2\bar H^2+\bar E\bar F=0,\quad
2\bar H^2=\tfrac{16}{9}\bar\omega^3,\quad \bar H^4=0,\quad \bar\omega^6=0 .
\end{equation*}
Abe's sharper $\bar\omega^5=0$ follows from his (unproved here) $\bar\omega^2\bar H=0$.
\end{btheorem}
{\raggedright\noindent\textbf{Lean.} \lean{VOA.}\allowbreak\lean{AbeRel.}\allowbreak\lean{wbar},\allowbreak\ \lean{VOA.}\allowbreak\lean{AbeRel.}\allowbreak\lean{Ebar},\allowbreak\ \lean{VOA.}\allowbreak\lean{AbeRel.}\allowbreak\lean{Hbar},\allowbreak\ \lean{VOA.}\allowbreak\lean{AbeRel.}\allowbreak\lean{Fbar},\allowbreak\ \lean{VOA.}\allowbreak\lean{AbeRel.}\allowbreak\lean{Ebar_sq},\allowbreak\ \lean{VOA.}\allowbreak\lean{AbeRel.}\allowbreak\lean{Fbar_sq},\allowbreak\ \lean{VOA.}\allowbreak\lean{AbeRel.}\allowbreak\lean{Ebar_mul_Hbar},\allowbreak\ \lean{VOA.}\allowbreak\lean{AbeRel.}\allowbreak\lean{two_Hbar_sq_add_Ebar_Fbar},\allowbreak\ \lean{VOA.}\allowbreak\lean{AbeRel.}\allowbreak\lean{two_Hbar_sq_eq},\allowbreak\ \lean{VOA.}\allowbreak\lean{AbeRel.}\allowbreak\lean{Hbar_pow_four},\allowbreak\ \lean{VOA.}\allowbreak\lean{AbeRel.}\allowbreak\lean{wbar_pow_six},\allowbreak\ \lean{VOA.}\allowbreak\lean{AbeRel.}\allowbreak\lean{OmegaSqHbarZero},\allowbreak\ \lean{VOA.}\allowbreak\lean{AbeRel.}\allowbreak\lean{OmegaPowFiveZero},\allowbreak\ \lean{VOA.}\allowbreak\lean{AbeRel.}\allowbreak\lean{omegaPowFive_of_omegaSqHbar}\par}
{\raggedright\noindent File \leanpath{VOA/\allowbreak Examples/\allowbreak TripletW2AbeRelations.lean}.\par}
\begin{leancode}
noncomputable def wbar : A₂ tripletW2 := mk₂ tripletW2 abeOmega

noncomputable def Ebar : A₂ tripletW2 := mk₂ tripletW2 abeE

noncomputable def Hbar : A₂ tripletW2 := mk₂ tripletW2 abeH

noncomputable def Fbar : A₂ tripletW2 := mk₂ tripletW2 abeF

theorem Ebar_sq : Ebar * Ebar = 0 := by

theorem Fbar_sq : Fbar * Fbar = 0 := by

theorem Ebar_mul_Hbar : Ebar * Hbar = 0 := by

theorem two_Hbar_sq_add_Ebar_Fbar : (2 : ℂ) • (Hbar * Hbar) + Ebar * Fbar = 0 := by

theorem two_Hbar_sq_eq : (2 : ℂ) • (Hbar * Hbar) = ((16 : ℂ) / 9) • (wbar * (wbar * wbar)) := by

theorem Hbar_pow_four : Hbar ^ 4 = 0 := by

theorem wbar_pow_six : wbar ^ 6 = 0 := by

def OmegaSqHbarZero : Prop := wbar * (wbar * Hbar) = 0

def OmegaPowFiveZero : Prop := wbar ^ 5 = 0

theorem omegaPowFive_of_omegaSqHbar (h : OmegaSqHbarZero) : OmegaPowFiveZero := by
\end{leancode}
\noindent\textbf{Hypotheses / status.} None\par
\medskip
\begin{btheorem}{18.29}[A logarithmic module from a Zhu-algebra Jordan block]
Let $V$ be a VOA with $A(V)$ and all $A_{n,0}(V)$ finite-dimensional. For every $h$ there is a logarithmic module
$M_h=M(U_h)$ (the Dong--Jiang induced module on $U_h=A(V)/A(V)\omega_h^2$) on which $L_0=\rho([\omega])+n$ on level
$n$. If moreover the Dong--Jiang gate $A_{0,0}(V)=A(V)$ holds and \lean{OmegaJordanAt V h}, then
$L_0^{\mathrm{nil}}\ne0$ on $M_h$: there are $x,y$ of generalized weight $h$ with $y\ne0$ and $L_0x=hx+y$.
For the Abe model the two finiteness hypotheses follow from $C_2$-cofiniteness.
\end{btheorem}
{\raggedright\noindent\textbf{Lean.} \lean{VOA.}\allowbreak\lean{djLogVModule_L0nil_ne_zero},\allowbreak\ \lean{VOA.}\allowbreak\lean{omegaJordan_nil_sq},\allowbreak\ \lean{VOA.}\allowbreak\lean{SympFerm.}\allowbreak\lean{tripletW2Log},\allowbreak\ \lean{VOA.}\allowbreak\lean{SympFerm.}\allowbreak\lean{tripletW2Log_jordan_pair},\allowbreak\ \lean{VOA.}\allowbreak\lean{SympFerm.}\allowbreak\lean{tripletW2_zhuFiniteness_of_C2Cofinite}\par}
{\raggedright\noindent Files \leanpath{VOA/\allowbreak Examples/\allowbreak TripletW2Log.lean},
\leanpath{TripletW2C2.lean}.\par}
\begin{leancode}
theorem djLogVModule_L0nil_ne_zero (W : VertexOperatorAlgebra 𝒱) (U : AmMod W 0) (h : ℂ) (k : ℕ)
    (hnil : amModOmegaNil U h ^ k = 0)
    (hfin : ∀ n : ℕ, FiniteDimensional ℂ (DJPiece W U n))
    (hgate : DJAnnEqAn W 0) (hne : amModOmegaNil U h ≠ 0) :
    (djLogVModule W U h k hnil hfin).L0nil ≠ 0 := by

theorem omegaJordan_nil_sq (h : ℂ) :
    amModOmegaNil (omegaJordanAMod W h).toAmModZero h ^ 2 = 0 := by

noncomputable def tripletW2Log (h : ℂ)
    (hAnm : ∀ n : ℕ, FiniteDimensional ℂ (AnmFull tripletW2 n 0))
    (hA : FiniteDimensional ℂ (A tripletW2)) :
    LogVModule tripletW2
      (DJCarrier tripletW2 (omegaJordanAMod tripletW2 h).toAmModZero) :=

theorem tripletW2Log_jordan_pair (h : ℂ)
    (hAnm : ∀ n : ℕ, FiniteDimensional ℂ (AnmFull tripletW2 n 0))
    (hA : FiniteDimensional ℂ (A tripletW2)) (hgate : DJAnnEqAn tripletW2 0)
    (hnd : OmegaJordanAt tripletW2 h) :
    ∃ x y : DJCarrier tripletW2 (omegaJordanAMod tripletW2 h).toAmModZero,
      x ∈ (tripletW2Log h hAnm hA).weight h ∧ y ∈ (tripletW2Log h hAnm hA).weight h ∧
        y ≠ 0 ∧ (tripletW2Log h hAnm hA).L0op x = h • x + y :=
\end{leancode}
\begin{leancode}
theorem tripletW2_zhuFiniteness_of_C2Cofinite (hC2 : C2Cofinite tripletW2) :
    (∀ n : ℕ, FiniteDimensional ℂ (AnmFull tripletW2 n 0))
      ∧ FiniteDimensional ℂ (A tripletW2) :=
\end{leancode}
\noindent\textbf{Hypotheses / status.} \lean{hAnm}, \lean{hA}, \lean{hgate}, \lean{hnd}\par
\medskip
\endgroup
\setcounter{chapter}{18}
\chapter{Graded unitarity and Lee--Yang}
\begingroup
\providecommand{\lnineteenvac}{\mathbf{1}}
\providecommand{\lnineteenV}{\mathcal{V}}
\providecommand{\lnineteenL}{L}
\providecommand{\lnineteenG}{G}
\providecommand{\lnineteenLay}{\mathcal{L}}
\providecommand{\lnineteenh}{\langle\!\langle}
\providecommand{\lnineteenhr}{\rangle\!\rangle}
\providecommand{\lnineteencb}{c_{2,2n+1}}
\providecommand{\lnineteenPi}{\Pi}
\providecommand{\lnineteensl}{\mathfrak{sl}_2}
\providecommand{\lnineteenR}{\mathbb{R}}
\noindent{\small 7 of the chapter's 52 numbered statements: its headline results and every definition they use.}\par\medskip
\begin{btheorem}{19.4}[A bilinear form is never positive on a nonzero complex subspace]
Let $\beta$ be a $\C$-bilinear form on a complex vector space $E$ and $S\le E$ a complex subspace.
If $\beta(x,x)$ is a positive real number for every nonzero $x\in S$, then $S = 0$.
\end{btheorem}
{\raggedright\noindent\textbf{Lean.} \lean{VOA.}\allowbreak\lean{eq_bot_of_bilinear_pos}\par}
{\raggedright\noindent File \leanpath{VOA/\allowbreak Core/\allowbreak GradedUnitarity.lean}.\par}
\begin{leancode}
theorem eq_bot_of_bilinear_pos (β : BilinForm ℂ E) (S : Submodule ℂ E)
    (hpos : ∀ x ∈ S, x ≠ 0 → ∃ r : ℝ, 0 < r ∧ β x x = (r : ℂ)) : S = ⊥ := by
\end{leancode}
\noindent\textbf{Hypotheses / status.} None; unconditional\par
\medskip
\begin{btheorem}{19.12}[The two orthogonal complements agree on $\rho$-stable subspaces]
Let $S\le V_n$ be $\rho$-stable ($\rho(S)\subseteq S$) and suppose $s(n)\ne0$. Then
$\{y : h_n(x,y) = 0\ \forall x\in S\} = \{y : B(x,y) = 0\ \forall x\in S\}$.
\end{btheorem}
{\raggedright\noindent\textbf{Lean.} \lean{VOA.}\allowbreak\lean{RhoStableOn},\allowbreak\ \lean{VOA.}\allowbreak\lean{orthogonal_gradedSesq_eq}\par}
{\raggedright\noindent File \leanpath{VOA/\allowbreak Core/\allowbreak GradedUnitarity.lean}.\par}
\begin{leancode}
def RhoStableOn (R : RealStructure W) {n : ℤ} (S : Submodule ℂ (W.grading n)) : Prop :=
  ∀ x ∈ S, ∃ y ∈ S, (y : 𝒱) = R.rho (x : 𝒱)

theorem orthogonal_gradedSesq_eq (B : InvariantForm W) (R : RealStructure W) {n : ℤ}
    {S : Submodule ℂ (W.grading n)} (hS : RhoStableOn R S) (hp : R.parity n ≠ 0) :
    (gradedSesq B R n).orthogonal S = (LiRadical.gradedForm B n).orthogonal S := by
\end{leancode}
\noindent\textbf{Hypotheses / status.} $\rho$-stable, $s(n)\ne0$\par
\medskip
\begin{btheorem}{19.24}[Canonical real structures on the boundary models]
For every $n\in\N$, $c_{2,2n+1} = 1-12\lambda_n^2$ with the real number
$\lambda_n = \sqrt{(2n-1)^2/(4(2n+1))}$, so $L(\lnineteencb,0)$ carries a canonical real structure.
\end{btheorem}
{\raggedright\noindent\textbf{Lean.} \lean{VOA.}\allowbreak\lean{VirGradedUnitarity.}\allowbreak\lean{boundaryBackgroundChargeSq},\allowbreak\ \lean{VOA.}\allowbreak\lean{VirGradedUnitarity.}\allowbreak\lean{boundaryBackgroundCharge_spec},\allowbreak\ \lean{VOA.}\allowbreak\lean{VirGradedUnitarity.}\allowbreak\lean{boundaryVirasoroRealStructure}\par}
{\raggedright\noindent File \leanpath{VOA/\allowbreak Examples/\allowbreak VirasoroBoundaryRealStructure.lean}.\par}
\begin{leancode}
noncomputable def boundaryBackgroundChargeSq (n : ℕ) : ℝ :=
  (2 * (n : ℝ) - 1) ^ 2 / (4 * (2 * n + 1))

theorem boundaryBackgroundCharge_spec (n : ℕ) :
    1 - 12 * boundaryBackgroundCharge n ^ 2 = boundaryCharge n := by

noncomputable def boundaryVirasoroRealStructure (n : ℕ) :
    RealStructure (virasoroSimple (boundaryCharge n)) :=
  virasoroSimpleRealStructure (boundaryBackgroundCharge_spec n)
    (boundaryBackgroundCharge_fixed n)
\end{leancode}
\noindent\textbf{Hypotheses / status.} $c = 1-12\lambda^2$, $\lambda\in\lnineteenR$; boundary: none\par
\medskip
\begin{bdefinition}{19.51}[Layer signs, step nondegeneracy, and the two analytic inputs]
For a family $R$ of real structures on $V^x$, $x<-2$, write $\mathrm{Sign}_R(x,E)$ for clause (vi) at weight $E$
and $\mathrm{Nd}(x,E)$ for clause (iv) at weight $E$, both for the ordinary flag and the form $B_x$.
\begin{enumerate}
\item The \emph{anchor} for $R$: for every $E$ there is $K\in\lnineteenR$ such that $\mathrm{Sign}_R(x,E)$ holds for
  every $x<-2$ with $x<K$.
\item \emph{Rigidity} for $R$: for every $E$, if $\mathrm{Nd}(x,E)$ holds for every $x<-2$, then
  $\mathrm{Sign}_R(x_0,E)$ for one $x_0<-2$ implies $\mathrm{Sign}_R(x_1,E)$ for every $x_1<-2$.
\end{enumerate}
\end{bdefinition}
{\raggedright\noindent\textbf{Lean.} \lean{VOA.}\allowbreak\lean{AffineSl2BelowCritical.}\allowbreak\lean{LayerSigns},\allowbreak\ \lean{VOA.}\allowbreak\lean{AffineSl2BelowCritical.}\allowbreak\lean{LayerNondeg},\allowbreak\ \lean{VOA.}\allowbreak\lean{AffineSl2BelowCritical.}\allowbreak\lean{LayerSignAnchor},\allowbreak\ \lean{VOA.}\allowbreak\lean{AffineSl2BelowCritical.}\allowbreak\lean{LayerSignRigidity}\par}
{\raggedright\noindent {\raggedright File \leanpath{VOA/\allowbreak Examples/\allowbreak AffineSl2BelowCriticalGradedUnitarity.lean}.\par}\par}
\begin{leancode}
def LayerSigns (R : RealFamily) (x : ℝ) (hx : x < -2) (E : ℤ) : Prop :=
  LayerPositive (ordinaryFlag (x : ℂ) (noncritical_of_lt hx)) E
    (LiRadical.gradedForm (universalForm (x : ℂ) (noncritical_of_lt hx)) E)
    (gradedSesq (universalForm (x : ℂ) (noncritical_of_lt hx)) (R x hx) E)

def LayerNondeg (x : ℝ) (hx : x < -2) (E : ℤ) : Prop :=
  LayerNondegenerate (ordinaryFlag (x : ℂ) (noncritical_of_lt hx)) E
    (LiRadical.gradedForm (universalForm (x : ℂ) (noncritical_of_lt hx)) E)

def LayerSignAnchor (R : RealFamily) : Prop :=
  ∀ E : ℤ, ∃ K : ℝ, ∀ (x : ℝ) (hx : x < -2), x < K → LayerSigns R x hx E

def LayerSignRigidity (R : RealFamily) : Prop :=
  ∀ E : ℤ, (∀ (x : ℝ) (hx : x < -2), LayerNondeg x hx E) →
    ∀ (x₀ : ℝ) (hx₀ : x₀ < -2) (x₁ : ℝ) (hx₁ : x₁ < -2),
      LayerSigns R x₀ hx₀ E → LayerSigns R x₁ hx₁ E
\end{leancode}
\noindent\textbf{Hypotheses / status.} Definitions; both proved for the compact family\par
\medskip
\begin{btheorem}{19.53}[Anchor: layer signs at large negative level]
\textbf{(a) Matrix version.} Let $I$ be a finite set with a length $\ell:I\to\N$, $P$ an $I\times I$ matrix of
complex polynomials, and $d:I\to\lnineteenR_{>0}$, such that
\begin{enumerate}
\item $\deg P_{ij}\le\min(\ell(i),\ell(j))$;
\item if $\ell(i) = \ell(j)$, the coefficient of $X^{\ell(i)}$ in $P_{ij}$ is $\delta_{ij}d_i$;
\item $P(x)$ is Hermitian for all real $x<x_*$ (some $x_*\in\lnineteenR$; \lean{x₀} in the Lean statement).
\end{enumerate}
Then there is $K$ such that for every real $x<K$ and every $r$, $\textsc{LayerSign}_r(P(x))$ holds.

\textbf{(b) Affine version.} For the compact family, the anchor of Definition~19.51 holds: for every
weight $E$ there is $K$ such that the weight-$E$ layer signs hold at every real level $x<\min(K,-2)$.
\end{btheorem}
{\raggedright\noindent\textbf{Lean.} \lean{VOA.}\allowbreak\lean{MatrixRigidity.}\allowbreak\lean{matrixLayerSign_eventually},\allowbreak\ \lean{VOA.}\allowbreak\lean{AffineSl2BelowCritical.}\allowbreak\lean{gramPoly_natDegree_le},\allowbreak\ \lean{VOA.}\allowbreak\lean{AffineSl2BelowCritical.}\allowbreak\lean{gramPoly_coeff_top},\allowbreak\ \lean{VOA.}\allowbreak\lean{AffineSl2BelowCritical.}\allowbreak\lean{layerSignAnchor_compact}\par}
{\raggedright\noindent {\raggedright Files \leanpath{VOA/\allowbreak Examples/\allowbreak AffineSl2BelowCriticalMatrixAnchor.lean} (a) and
\leanpath{VOA/\allowbreak Examples/\allowbreak AffineSl2BelowCriticalAnchor.lean} (b, and the two polynomial facts).\par}\par}
\begin{leancode}
theorem matrixLayerSign_eventually (P : Matrix ι ι ℂ[X]) (d : ι → ℝ) (hd : ∀ i, 0 < d i)
    (hdeg : ∀ i j, (P i j).natDegree ≤ min (ℓ i) (ℓ j))
    (htop : ∀ i j, ℓ i = ℓ j → (P i j).coeff (ℓ i) = if i = j then (d i : ℂ) else 0)
    (x₀ : ℝ) (hherm : ∀ x : ℝ, x < x₀ → (P.map (Polynomial.eval (x : ℂ))).IsHermitian) :
    ∃ K : ℝ, ∀ x : ℝ, x < K → ∀ r, MatrixLayerSign ℓ (P.map (Polynomial.eval (x : ℂ))) r := by

theorem gramPoly_natDegree_le {E : ℤ} (i j : EIdx E) :
    (gramPoly E i j).natDegree ≤ min (eLen E i) (eLen E j) := by
...
theorem gramPoly_coeff_top {E : ℤ} (i j : EIdx E) (hlen : eLen E i = eLen E j) :
    (gramPoly E i j).coeff (eLen E i) = if i = j then ((topWeight i : ℝ) : ℂ) else 0 := by
...
theorem layerSignAnchor_compact : LayerSignAnchor compactRealFamily := by
\end{leancode}
\noindent\textbf{Hypotheses / status.} Degree and top-coefficient bounds (proved for \lean{gramPoly}); unconditional\par
\medskip
\begin{btheorem}{19.54}[Rigidity of layer signs]
\textbf{(a) Matrix version.} Let $X$ be a preconnected topological space and $H:X\to\mathrm{Mat}_I(\C)$ a continuous
family of Hermitian matrices with $\textsc{StepNondeg}(H(t))$ for every $t$. If $\textsc{LayerSign}_r(H(t_0))$ holds for one
$t_0$, it holds for every $t_1$.

\textbf{(b) Affine version.} Rigidity (Definition~19.51) holds for the compact family. Consequently, for
every real $k<-2$, graded unitarity of $V^k(\lnineteensl)$ with the compact structure and the ordinary flag follows from
the anchor alone.
\end{btheorem}
{\raggedright\noindent\textbf{Lean.} \lean{VOA.}\allowbreak\lean{MatrixRigidity.}\allowbreak\lean{matrixLayerSign_rigid},\allowbreak\ \lean{VOA.}\allowbreak\lean{AffineSl2BelowCritical.}\allowbreak\lean{gram_isHermitian},\allowbreak\ \lean{VOA.}\allowbreak\lean{AffineSl2BelowCritical.}\allowbreak\lean{gram_eq_eval},\allowbreak\ \lean{VOA.}\allowbreak\lean{AffineSl2BelowCritical.}\allowbreak\lean{layerSignRigidity_compact},\allowbreak\ \lean{VOA.}\allowbreak\lean{AffineSl2BelowCritical.}\allowbreak\lean{universal_gradedUnitary_belowCritical_of_anchor}\par}
{\raggedright\noindent {\raggedright Files \leanpath{VOA/\allowbreak Examples/\allowbreak AffineSl2BelowCriticalMatrixRigidity.lean} (a) and
\leanpath{VOA/\allowbreak Examples/\allowbreak AffineSl2BelowCriticalRigidity.lean} (b).\par}\par}
\begin{leancode}
theorem matrixLayerSign_rigid {X : Type*} [TopologicalSpace X] [PreconnectedSpace X]
    (H : X → Matrix ι ι ℂ) (hc : Continuous H) (hH : ∀ t, (H t).IsHermitian)
    (hnd : ∀ t, MatrixStepNondeg ℓ (H t)) (r : ℕ) {t₀ t₁ : X}
    (h₀ : MatrixLayerSign ℓ (H t₀) r) : MatrixLayerSign ℓ (H t₁) r := by

theorem gram_isHermitian (E : ℤ) : (gram hk (compactRealStructure hreal hk) E).IsHermitian := by
...
theorem gram_eq_eval (hk0 : k ≠ 0) (E : ℤ) (i j : EIdx E) :
    gram hk (compactRealStructure hreal hk) E i j = (gramPoly E i j).eval k := by
...
theorem layerSignRigidity_compact : LayerSignRigidity compactRealFamily := by
...
theorem universal_gradedUnitary_belowCritical_of_anchor
    (hanchor : LayerSignAnchor compactRealFamily) (k : ℝ) (hk : k < -2) :
    GradedUnitary (universalForm (k : ℂ) (noncritical_of_lt hk)) (compactRealFamily k hk)
      (ordinaryFlag (k : ℂ) (noncritical_of_lt hk)) :=
  universal_gradedUnitary_belowCritical_compact hanchor layerSignRigidity_compact k hk
\end{leancode}
\noindent\textbf{Hypotheses / status.} Continuous Hermitian family, nondegenerate steps; unconditional\par
\medskip
\begin{btheorem}{19.55}[Nondegeneracy of the full form below critical level]
For every real $k<-2$ and every $E$, the form $B_k$ is nondegenerate on $V^k_E$: if $v\in V^k_E$ and $B_k(y,v) = 0$
for all $y\in V^k_E$, then $v = 0$.
\end{btheorem}
{\raggedright\noindent\textbf{Lean.} \lean{VOA.}\allowbreak\lean{AffineSl2BelowCritical.}\allowbreak\lean{universalForm_separating_of_lt}\par}
{\raggedright\noindent {\raggedright File \leanpath{VOA/\allowbreak Examples/\allowbreak AffineSl2BelowCriticalNondegenerate.lean}.\par}\par}
\begin{leancode}
theorem universalForm_separating_of_lt {x : ℝ} (hx : x < -2) (E : ℤ) :
    ∀ v ∈ (universal (x : ℂ) (noncritical_of_lt hx)).grading E,
      (∀ y ∈ (universal (x : ℂ) (noncritical_of_lt hx)).grading E,
        (universalForm (x : ℂ) (noncritical_of_lt hx)).form y v = 0) → v = 0 := by
\end{leancode}
\noindent\textbf{Hypotheses / status.} None; unconditional\par
\medskip
\endgroup
\part{Two recent theorems, formalized}
\setcounter{chapter}{19}
\chapter[Classical freeness at the boundary level: Arakawa--Dai]{Classical freeness at the boundary admissible level: the Arakawa--Dai theorems}
\begingroup
\providecommand{\ltwentyvac}{\mathbf{1}}
\providecommand{\ltwentysl}{\mathfrak{sl}_2}
\providecommand{\ltwentyR}{R}
\providecommand{\ltwentygr}{\operatorname{gr}^F}
\providecommand{\ltwentyJ}{J_\infty}
\providecommand{\ltwentyQ}{Q}
\providecommand{\ltwentyA}{A_m}
\providecommand{\ltwentyB}{B_m}
\providecommand{\ltwentyCm}{C_m}
\providecommand{\ltwentyqq}[1]{(q)_{#1}}
\providecommand{\ltwentyord}{\operatorname{ord}}
\noindent{\small 9 of the chapter's 56 numbered statements: its headline results and every definition they use.}\par\medskip
\begin{bdefinition}{20.8}[The fermionic chain and the Hilbert-bound series]
For $\ell,n\ge0$ define $K_\ell(n)\in\Q[[q]]$ by
\begin{equation*}
K_0(n)=\frac{(q)_n^2}{(q)_{2n+1}},\qquad K_{\ell+1}(n)=\sum_{j=0}^n\frac{q^{j(j+1)}}{(q)_{n-j}}K_\ell(j),
\end{equation*}
and the \emph{fermionic trace} $T_m=\sum_{n\ge0}q^{n(n+1)}K_{m-1}(n)$. Unwinding the recursion,
$T_m=\sum_{n_1\ge\cdots\ge n_m\ge0}q^{\sum n_j(n_j+1)}(q)_{n_m}^2/\bigl((q)_{n_1-n_2}\cdots
(q)_{n_{m-1}-n_m}(q)_{2n_m+1}\bigr)$, the sum in (20.3). The \emph{Hilbert-bound series}
$S_m$ (equation (5.12) of the paper) is the same chain with $K_0$ replaced by the three-species residual sum
\begin{equation*}
\mathrm{res}(n)=\sum_{s\ge0}q^{(n+1)s}\,\frac{1}{(q)_s}\,E(n+s+1)\,E(n+1)\,\frac{1}{(q)_n},\qquad
E(A)=\sum_{j\ge0}\frac{q^{Aj}}{(q)_j}.
\end{equation*}
\end{bdefinition}
{\raggedright\noindent\textbf{Lean.} \lean{VOA.}\allowbreak\lean{QSeries.}\allowbreak\lean{fermionicKernel},\allowbreak\ \lean{VOA.}\allowbreak\lean{QSeries.}\allowbreak\lean{fermionicTrace},\allowbreak\ \lean{VOA.}\allowbreak\lean{QSeries.}\allowbreak\lean{residualSum},\allowbreak\ \lean{VOA.}\allowbreak\lean{QSeries.}\allowbreak\lean{eulerSum},\allowbreak\ \lean{VOA.}\allowbreak\lean{QSeries.}\allowbreak\lean{powerShiftSum},\allowbreak\ \lean{VOA.}\allowbreak\lean{QSeries.}\allowbreak\lean{hilbertBoundKernel},\allowbreak\ \lean{VOA.}\allowbreak\lean{QSeries.}\allowbreak\lean{hilbertBoundSeries}\par}
{\raggedright\noindent File \leanpath{VOA/\allowbreak Commutative/\allowbreak QChainTrace.lean}.\par}
\begin{leancode}
def fermionicKernel : ℕ → ℕ → PowerSeries K
  | 0, n => qFactorial PowerSeries.X n ^ 2 * reciprocalQFactorial (2 * n + 1)
  | ℓ + 1, n => ∑ j ∈ Finset.range (n + 1),
      PowerSeries.X ^ (j * (j + 1)) * reciprocalQFactorial (n - j) * fermionicKernel ℓ j

def fermionicTrace (m : ℕ) : PowerSeries K :=
  VOA.LocallyFinitePowerSeriesSum.sum (fun n =>
    PowerSeries.X ^ (n * (n + 1)) * fermionicKernel (m - 1) n)
\end{leancode}
{\raggedright\noindent Files \leanpath{VOA/\allowbreak Commutative/\allowbreak QBinomialResidual.lean},
\leanpath{VOA/\allowbreak Commutative/\allowbreak QEuler.lean},
\leanpath{VOA/\allowbreak Commutative/\allowbreak QSeriesSum.lean}.\par}
\begin{leancode}
def residualSum (n : ℕ) : PowerSeries R :=
  powerShiftSum (n + 1) fun s =>
    reciprocalQFactorial s * eulerSum (n + s + 1) *
      eulerSum (n + 1) * reciprocalQFactorial n

def eulerSum (A : ℕ) : PowerSeries R :=
  powerShiftSum A reciprocalQFactorial

def powerShiftSum (A : ℕ) (B : ℕ → PowerSeries R) : PowerSeries R :=
  VOA.LocallyFinitePowerSeriesSum.sum (fun j => PowerSeries.X ^ (A * j) * B j)
\end{leancode}
{\raggedright\noindent File \leanpath{VOA/\allowbreak Commutative/\allowbreak QHilbertProduct.lean}.\par}
\begin{leancode}
def hilbertBoundKernel : ℕ → ℕ → PowerSeries K
  | 0, n => residualSum n
  | ℓ + 1, n => ∑ j ∈ Finset.range (n + 1),
      PowerSeries.X ^ (j * (j + 1)) * reciprocalQFactorial (n - j) * hilbertBoundKernel ℓ j

def hilbertBoundSeries (m : ℕ) : PowerSeries K :=
  VOA.LocallyFinitePowerSeriesSum.sum (fun n =>
    PowerSeries.X ^ (n * (n + 1)) * hilbertBoundKernel (m - 1) n)
\end{leancode}
\noindent\textbf{Hypotheses / status.} Definitions; finiteness proved\par
\medskip
\begin{btheorem}{20.11}[The arc universal property]
For every commutative $R$-algebra $A$, restriction along the universal arc (the tautological algebra map
$R[x]/I\to(\ltwentyJ(R[x]/I))[[t]]$, $x_i\mapsto\sum_nx_{i,n}t^n$) gives a bijection
$\operatorname{Hom}_{R\text{-alg}}(\ltwentyJ(R[x]/I),A)\cong\operatorname{Hom}_{R\text{-alg}}(R[x]/I,A[[t]])$.
If $e\colon R[x]/I\cong B$ is an isomorphism of algebras, composing with $e$ gives
$\operatorname{Hom}(\ltwentyJ(R[x]/I),A)\cong\operatorname{Hom}(B,A[[t]])$.
\end{btheorem}
{\raggedright\noindent\textbf{Lean.} \lean{VOA.}\allowbreak\lean{InfiniteJet.}\allowbreak\lean{seriesJetEval},\allowbreak\ \lean{VOA.}\allowbreak\lean{InfiniteJet.}\allowbreak\lean{toArc},\allowbreak\ \lean{VOA.}\allowbreak\lean{InfiniteJet.}\allowbreak\lean{arcEquiv},\allowbreak\ \lean{VOA.}\allowbreak\lean{InfiniteJet.}\allowbreak\lean{existsUnique_arcLift},\allowbreak\ \lean{VOA.}\allowbreak\lean{InfiniteJet.}\allowbreak\lean{arcEquivOfBaseAlgebraEquiv},\allowbreak\ \lean{VOA.}\allowbreak\lean{InfiniteJet.}\allowbreak\lean{baseToJet_injective}\par}
{\raggedright\noindent File \leanpath{VOA/\allowbreak Commutative/\allowbreak InfiniteJetArc.lean}.\par}
\begin{leancode}
def seriesJetEval (x : σ → PowerSeries A) : JetPoly R σ →ₐ[R] A :=
  MvPolynomial.aeval fun v => PowerSeries.coeff v.2 (x v.1)

def toArc (I : Ideal (MvPolynomial σ R)) (f : PresentedJet I →ₐ[R] A) :
    (MvPolynomial σ R ⧸ I) →ₐ[R] PowerSeries A :=
  Ideal.Quotient.liftₐ I (MvPolynomial.aeval (pointSeries I f))
    (pointSeries_relations I f)

def arcEquiv (I : Ideal (MvPolynomial σ R)) :
    (PresentedJet I →ₐ[R] A) ≃ ((MvPolynomial σ R ⧸ I) →ₐ[R] PowerSeries A) where
  toFun := toArc I
  invFun := ofArc I
  left_inv := ofArc_toArc I
  right_inv := toArc_ofArc I

theorem existsUnique_arcLift (I : Ideal (MvPolynomial σ R))
    (u : (MvPolynomial σ R ⧸ I) →ₐ[R] PowerSeries A) :
    ∃! f : PresentedJet I →ₐ[R] A,
      (PowerSeries.mapAlgHom f).comp (universalArc I) = u := by
\end{leancode}
{\raggedright\noindent File \leanpath{VOA/\allowbreak Commutative/\allowbreak InfiniteJetBaseEmbedding.lean}.\par}
\begin{leancode}
def arcEquivOfBaseAlgebraEquiv {A B : Type*} [CommRing A] [Algebra R A]
    [CommRing B] [Algebra R B] (I : Ideal (MvPolynomial σ R))
    (e : (MvPolynomial σ R ⧸ I) ≃ₐ[R] B) :
    (PresentedJet I →ₐ[R] A) ≃ (B →ₐ[R] PowerSeries A) where
  toFun f := (toArc I f).comp e.symm.toAlgHom
  invFun u := ofArc I (u.comp e.toAlgHom)
  left_inv f := by simp [AlgHom.comp_assoc]
  right_inv u := by simp [AlgHom.comp_assoc]

theorem baseToJet_injective (I : Ideal (MvPolynomial σ R)) : Function.Injective (baseToJet I) := by
\end{leancode}
\noindent\textbf{Hypotheses / status.} Definitions; arc property proved\par
\medskip
\begin{bdefinition}{20.12}[Poisson vertex algebra, in mode form]
A Poisson vertex algebra over $\C$ is a commutative $\C$-algebra $A$ with a derivation $\partial$ and bilinear
products $a_{(n)}b$, $n\in\N$, such that for all $a,b,c$:
\begin{enumerate}
\item (locality) $a_{(n)}b=0$ for $n\gg0$;
\item (sesquilinearity) $(\partial a)_{(0)}b=0$, $(\partial a)_{(n+1)}b=-(n+1)\,a_{(n)}b$, and
  $\partial(a_{(n)}b)=(\partial a)_{(n)}b+a_{(n)}(\partial b)$;
\item (Leibniz) $a_{(n)}(bc)=(a_{(n)}b)c+b(a_{(n)}c)$;
\item (skew-symmetry) $a_{(n)}b=\sum_{j\ge0}\frac{(-1)^{n+j+1}}{j!}\partial^j(b_{(n+j)}a)$;
\item (Jacobi) $a_{(m)}(b_{(n)}c)-b_{(n)}(a_{(m)}c)=\sum_{j=0}^m\binom mj(a_{(j)}b)_{(m+n-j)}c$.
\end{enumerate}
Equivalently, $\{a_\lambda b\}=\sum_n\frac{\lambda^n}{n!}a_{(n)}b$ is a $\lambda$-bracket satisfying the usual
PVA axioms. A homomorphism is an algebra map commuting with $\partial$ and every $(n)$-product; an
equivalence is a bijective homomorphism (whose inverse is then automatically a homomorphism).
\end{bdefinition}
{\raggedright\noindent\textbf{Lean.} \lean{VOA.}\allowbreak\lean{PoissonVertexAlgebra},\allowbreak\ \lean{VOA.}\allowbreak\lean{PoissonVertexAlgebra.}\allowbreak\lean{Hom},\allowbreak\ \lean{VOA.}\allowbreak\lean{PoissonVertexAlgebra.}\allowbreak\lean{Equiv},\allowbreak\ \lean{VOA.}\allowbreak\lean{PoissonVertexAlgebra.}\allowbreak\lean{Equiv.}\allowbreak\lean{ofBijective},\allowbreak\ \lean{VOA.}\allowbreak\lean{PoissonVertexAlgebra.}\allowbreak\lean{Equiv.}\allowbreak\lean{symm}\par}
{\raggedright\noindent File \leanpath{VOA/\allowbreak Commutative/\allowbreak PoissonVertexAlgebra.lean}.\par}
\begin{leancode}
structure PoissonVertexAlgebra (A : Type*) [CommRing A] [Algebra ℂ A] where
  translation : Derivation ℂ A A
  mode : ℕ → A →ₗ[ℂ] A →ₗ[ℂ] A
  locality : ∀ x y, ∃ B : ℕ, ∀ n, B ≤ n → mode n x y = 0
  translation_left_zero : ∀ x y, mode 0 (translation x) y = 0
  translation_left_succ : ∀ n x y,
    mode (n + 1) (translation x) y = (-((n + 1 : ℕ) : ℂ)) • mode n x y
  translation_mode : ∀ n x y,
    translation (mode n x y) = mode n (translation x) y + mode n x (translation y)
  leibniz : ∀ n x y z, mode n x (y * z) = mode n x y * z + y * mode n x z
  skew : ∀ n x y, mode n x y = ∑ᶠ j : ℕ,
    (j.factorial : ℂ)⁻¹ • (-1 : ℂ) ^ (n + j + 1) •
      ((translation.toLinearMap ^ j) (mode (n + j) y x))
  jacobi : ∀ m n x y z,
    mode m x (mode n y z) - mode n y (mode m x z) =
      ∑ j : Fin (m + 1), (m.choose j.val : ℂ) • mode (m + n - j.val) (mode j.val x y) z
\end{leancode}
{\raggedright\noindent Files \leanpath{VOA/\allowbreak Commutative/\allowbreak PoissonVertexHom.lean},
\leanpath{VOA/\allowbreak Commutative/\allowbreak PoissonVertexEquiv.lean}.\par}
\begin{leancode}
structure Hom (P : VOA.PoissonVertexAlgebra A) (Q : VOA.PoissonVertexAlgebra B) where
  toAlgHom : A →ₐ[ℂ] B
  map_translation : ∀ x, toAlgHom (P.translation x) = Q.translation (toAlgHom x)
  map_mode : ∀ n x y, toAlgHom (P.mode n x y) = Q.mode n (toAlgHom x) (toAlgHom y)

structure Equiv (P : VOA.PoissonVertexAlgebra A) (Q : VOA.PoissonVertexAlgebra B) where
  toAlgEquiv : A ≃ₐ[ℂ] B
  map_translation : ∀ x, toAlgEquiv (P.translation x) = Q.translation (toAlgEquiv x)
  map_mode : ∀ n x y, toAlgEquiv (P.mode n x y) = Q.mode n (toAlgEquiv x) (toAlgEquiv y)

def Equiv.ofBijective {P : VOA.PoissonVertexAlgebra A} {Q : VOA.PoissonVertexAlgebra B}
    (f : Hom P Q) (hf : Function.Bijective f.toAlgHom) : Equiv P Q where
  toAlgEquiv := AlgEquiv.ofBijective f.toAlgHom hf
  map_translation := f.map_translation
  map_mode := f.map_mode

def Equiv.symm {P : VOA.PoissonVertexAlgebra A} {Q : VOA.PoissonVertexAlgebra B}
    (e : Equiv P Q) : Equiv Q P where
\end{leancode}
\noindent\textbf{Hypotheses / status.} $1\le m$; unconditional\par
\medskip
\begin{btheorem}{20.18}[The vacuum character of $L_{k_m}(\ltwentysl)$]
For every $m\ge1$, $p=2m+1$:
\begin{equation*}
\sum_{N\ge0}\dim L_{k_m}(\ltwentysl)_N\,q^N=\frac{(q^p;q^p)_\infty^3}{(q)_\infty^3}=S_m(q)
=\frac{1}{(q)_\infty^2}\,T_m(q),
\end{equation*}
with $S_m$, $T_m$ as in Definition~20.8; the last expression is the fermionic chain of
(20.3).
\end{btheorem}
{\raggedright\noindent\textbf{Lean.} \lean{VOA.}\allowbreak\lean{AffineSl2BoundaryCharacter.}\allowbreak\lean{coeff_product_le_character},\allowbreak\ \lean{VOA.}\allowbreak\lean{AffineSl2BoundaryCharacter.}\allowbreak\lean{character_eq_product},\allowbreak\ \lean{VOA.}\allowbreak\lean{AffineSl2BoundaryCharacter.}\allowbreak\lean{character_eq_chain},\allowbreak\ \lean{VOA.}\allowbreak\lean{AffineSl2BoundaryCharacter.}\allowbreak\lean{character_eq_fermionic_chain}\par}
{\raggedright\noindent File \leanpath{VOA/\allowbreak Examples/\allowbreak AffineSl2BoundaryMainTheorems.lean}.\par}
\begin{leancode}
theorem coeff_product_le_character (m : ℕ) (hm : 1 ≤ m) (N : ℕ) :
    PowerSeries.coeff N (product m) ≤ PowerSeries.coeff N (character m) := by

theorem character_eq_product (m : ℕ) (hm : 1 ≤ m) :
    character m = product m := by

theorem character_eq_chain (m : ℕ) (hm : 1 ≤ m) :
    character m = QSeries.hilbertBoundSeries m :=
  character_eq_chain_of_product m hm (character_eq_product m hm)

theorem character_eq_fermionic_chain (m : ℕ) (hm : 1 ≤ m) :
    character m = (QSeries.reciprocalEuler : PowerSeries ℚ)^2 * QSeries.fermionicTrace m := by
\end{leancode}
\noindent\textbf{Hypotheses / status.} $1\le m$; unconditional, \lean{pbw}-free\par
\medskip
\begin{btheorem}{20.20}[Classical freeness of $L_{k_m}(\ltwentysl)$]
For every $m\ge1$ the canonical map $\pi\colon\ltwentyJ\ltwentyA\to\ltwentygr L_{k_m}(\ltwentysl)$ is bijective, and is an
isomorphism of Poisson vertex algebras which, in both directions, maps the conformal weight $N$ piece onto the
conformal weight $N$ piece.
\end{btheorem}
{\raggedright\noindent\textbf{Lean.} \lean{VOA.}\allowbreak\lean{AffineSl2BoundaryClassicalFreeness.}\allowbreak\lean{canonicalMap_bijective},\allowbreak\ \lean{VOA.}\allowbreak\lean{AffineSl2BoundaryClassicalFreeness.}\allowbreak\lean{boundary_classically_free},\allowbreak\ \lean{VOA.}\allowbreak\lean{AffineSl2BoundaryClassicalFreeness.}\allowbreak\lean{arcPoissonEquiv},\allowbreak\ \lean{VOA.}\allowbreak\lean{AffineSl2BoundaryClassicalFreeness.}\allowbreak\lean{arcPoissonEquiv_conformal_iff},\allowbreak\ \lean{VOA.}\allowbreak\lean{AffineSl2BoundaryClassicalFreeness.}\allowbreak\lean{arcPoissonEquivOfCharacter}\par}
{\raggedright\noindent File \leanpath{VOA/\allowbreak Examples/\allowbreak AffineSl2BoundaryMainTheorems.lean}.\par}
\begin{leancode}
theorem canonicalMap_bijective (m : ℕ) (hm : 1 ≤ m) :
    Function.Bijective (canonicalMap m hm) :=
  canonicalMap_bijective_of_character m hm (character_eq_product m hm)

theorem boundary_classically_free (m : ℕ) (hm : 1 ≤ m) :
    Function.Bijective (uniformCandidate m hm).arcLiPoissonHom.toAlgHom :=
  ⟨arcLiPoissonHom_injective_of_character m hm (character_eq_product m hm),
    (uniformCandidate m hm).arcLiPoissonHom_surjective⟩

def arcPoissonEquiv (m : ℕ) (hm : 1 ≤ m) :
    PoissonVertexAlgebra.Equiv (arcPVA m)
      (liPoissonVertexAlgebra (AffineSl2BoundaryCharacter.vacuum m)) :=
  arcPoissonEquivOfCharacter m hm (character_eq_product m hm)

theorem arcPoissonEquiv_conformal_iff (m : ℕ) (hm : 1 ≤ m) (N : ℕ)
    (x : InfiniteJet.PresentedJet (baseIdeal (R := ℂ) m)) :
    (arcPoissonEquiv m hm).toAlgEquiv x ∈
      BoundaryArcConformalSqueeze.liConformalPiece (AffineSl2BoundaryCharacter.vacuum m) N ↔
      x ∈ arcConformalPiece m N :=
  arcPoissonEquivOfCharacter_conformal_iff m hm (character_eq_product m hm) N x
\end{leancode}
{\raggedright\noindent File \leanpath{VOA/\allowbreak Examples/\allowbreak AffineSl2BoundaryClassicalFreeness.lean}.\par}
\begin{leancode}
def arcPoissonEquivOfCharacter (m : ℕ) (hm : 1 ≤ m)
    (hcharacter : character m = product m) :
    PoissonVertexAlgebra.Equiv (arcPVA m) (liPoissonVertexAlgebra (AffineSl2BoundaryCharacter.vacuum m)) :=
  PoissonVertexAlgebra.Equiv.ofBijective (uniformCandidate m hm).arcLiPoissonHom
    ⟨arcLiPoissonHom_injective_of_character m hm hcharacter,
      (uniformCandidate m hm).arcLiPoissonHom_surjective⟩
\end{leancode}
\noindent\textbf{Hypotheses / status.} $1\le m$; unconditional\par
\medskip
\begin{btheorem}{20.21}[The $C_2$-algebra and the arc universal property]
For $m\ge1$: (1) the map $\C[e,h,f]\to R_{L_{k_m}}$, $x_i\mapsto[\overline{x_i}]$, induces an isomorphism
$\ltwentyA=\C[e,h,f]/(e\ltwentyQ^m,h\ltwentyQ^m,f\ltwentyQ^m)\cong R_{L_{k_m}(\ltwentysl)}$; (2) for every commutative
$\C$-algebra $A$, algebra maps $\ltwentyJ\ltwentyA\to A$ correspond bijectively to algebra maps
$R_{L_{k_m}}\to A[[t]]$, via the specific bijection obtained from Theorem~20.11 and (1); (3) this bijection is natural in $A$. So
$\ltwentyJ\ltwentyA$ represents the arc functor $A\mapsto\operatorname{Hom}(R_{L_{k_m}},A[[t]])$, i.e.\ it is the arc algebra
$\ltwentyJ R_{L_{k_m}}$. (Naturality is not part of the type of \lean{actualC2ArcEquiv}; it is the separate theorem
\lean{actualC2ArcEquiv_natural}; see the warning below.)
\end{btheorem}
{\raggedright\noindent\textbf{Lean.} \lean{VOA.}\allowbreak\lean{AffineSl2BoundaryClassicalFreeness.}\allowbreak\lean{baseC2AlgEquiv},\allowbreak\ \lean{VOA.}\allowbreak\lean{AffineSl2BoundaryClassicalFreeness.}\allowbreak\lean{actualC2ArcEquiv},\allowbreak\ \lean{VOA.}\allowbreak\lean{AffineSl2BoundaryClassicalFreeness.}\allowbreak\lean{actualC2ArcEquiv_natural},\allowbreak\ \lean{VOA.}\allowbreak\lean{AffineSl2BoundarySingularFamily.}\allowbreak\lean{Candidate.}\allowbreak\lean{baseC2AlgEquiv},\allowbreak\ \lean{VOA.}\allowbreak\lean{AffineSl2BoundarySingularFamily.}\allowbreak\lean{Candidate.}\allowbreak\lean{actualC2ArcEquiv},\allowbreak\ \lean{VOA.}\allowbreak\lean{AffineSl2BoundarySingularFamily.}\allowbreak\lean{Candidate.}\allowbreak\lean{actualC2ArcEquiv_natural},\allowbreak\ \lean{VOA.}\allowbreak\lean{AffineSl2BoundarySingularFamily.}\allowbreak\lean{Candidate.}\allowbreak\lean{baseIdeal_eq_simpleSymbol_kernel},\allowbreak\ \lean{VOA.}\allowbreak\lean{baseC2Map}\par}
{\raggedright\noindent File \leanpath{VOA/\allowbreak Examples/\allowbreak AffineSl2BoundaryMainTheorems.lean}.\par}
\begin{leancode}
def baseC2AlgEquiv (m : ℕ) (hm : 1 ≤ m) :
    (BasePoly ℂ ⧸ baseIdeal (R := ℂ) m) ≃ₐ[ℂ] A₂ (AffineSl2BoundaryCharacter.vacuum m) :=
  baseC2AlgEquivOfCharacter m hm (character_eq_product m hm)

def actualC2ArcEquiv {A : Type*} [CommRing A] [Algebra ℂ A]
    (m : ℕ) (hm : 1 ≤ m) :
    (InfiniteJet.PresentedJet (baseIdeal (R := ℂ) m) →ₐ[ℂ] A) ≃
      (A₂ (AffineSl2BoundaryCharacter.vacuum m) →ₐ[ℂ] PowerSeries A) :=
  actualC2ArcEquivOfCharacter m hm (character_eq_product m hm)

theorem actualC2ArcEquiv_natural {A B : Type*} [CommRing A] [Algebra ℂ A]
    [CommRing B] [Algebra ℂ B] (m : ℕ) (hm : 1 ≤ m)
    (f : InfiniteJet.PresentedJet (baseIdeal (R := ℂ) m) →ₐ[ℂ] A) (g : A →ₐ[ℂ] B) :
    actualC2ArcEquiv m hm (g.comp f) = (PowerSeries.mapAlgHom g).comp (actualC2ArcEquiv m hm f) :=
\end{leancode}
{\raggedright\noindent File \leanpath{VOA/\allowbreak Filtered/\allowbreak BoundarySimpleArcPoisson.lean}.\par}
\begin{leancode}
def Candidate.baseC2AlgEquiv
    (hinj : Function.Injective C.arcLiPoissonHom.toAlgHom) :
    (BoundaryAuxiliary.BasePoly ℂ ⧸ BoundaryAuxiliary.baseIdeal (R := ℂ) m) ≃ₐ[ℂ]
      A₂ (BoundaryArcLiMap.simpleVOA C) :=
  baseC2AlgEquiv_of_arc_injective (BoundaryArcLiMap.simpleVOA C) (fun _ => 1)
    (fun i => ⟨BoundaryArcLiMap.gen C i,BoundaryArcLiMap.gen_mem C i⟩)
    (BoundaryAuxiliary.baseIdeal m) C.baseC2Relations (BoundaryArcLiMap.baseIdeal_le_kernel C)
    hinj C.baseC2Eval_surjective

def Candidate.actualC2ArcEquiv {A : Type*} [CommRing A] [Algebra ℂ A]
    (hinj : Function.Injective C.arcLiPoissonHom.toAlgHom) :
    (InfiniteJet.PresentedJet (BoundaryAuxiliary.baseIdeal (R := ℂ) m) →ₐ[ℂ] A) ≃
      (A₂ (BoundaryArcLiMap.simpleVOA C) →ₐ[ℂ] PowerSeries A) :=
  InfiniteJet.arcEquivOfBaseAlgebraEquiv (BoundaryAuxiliary.baseIdeal m) (C.baseC2AlgEquiv hinj)

theorem Candidate.actualC2ArcEquiv_natural {A B : Type*} [CommRing A] [Algebra ℂ A]
    [CommRing B] [Algebra ℂ B] (hinj : Function.Injective C.arcLiPoissonHom.toAlgHom)
    (f : InfiniteJet.PresentedJet (BoundaryAuxiliary.baseIdeal (R := ℂ) m) →ₐ[ℂ] A)
    (g : A →ₐ[ℂ] B) :
    C.actualC2ArcEquiv hinj (g.comp f) =
      (PowerSeries.mapAlgHom g).comp (C.actualC2ArcEquiv hinj f) :=

theorem Candidate.baseIdeal_eq_simpleSymbol_kernel
    (hinj : Function.Injective C.arcLiPoissonHom.toAlgHom) :
    BoundaryAuxiliary.baseIdeal (R := ℂ) m = RingHom.ker (simpleSymbol C) := by
\end{leancode}
{\raggedright\noindent File \leanpath{VOA/\allowbreak Filtered/\allowbreak ArcBaseExactness.lean}.\par}
\begin{leancode}
def baseC2Map : (MvPolynomial σ ℂ ⧸ I) →ₐ[ℂ] A₂ W :=
  Ideal.Quotient.liftₐ I (MvPolynomial.aeval (fun i => mk₂ W (gen i).1)) (fun _ hf => hI hf)
\end{leancode}
\noindent\textbf{Hypotheses / status.} $1\le m$; unconditional\par
\medskip
\begin{bdefinition}{20.22}[The canonical map, the universal arc of $R_V$, and the zero-jet embedding]
Let $\bar e,\bar h,\bar f\in V_1$ be the images in $V$ of the three current states $e_{(-1)}\ltwentyvac$, $h_{(-1)}\ltwentyvac$,
$f_{(-1)}\ltwentyvac$, and write $(x_0,x_1,x_2)=(e,h,f)$.
\begin{enumerate}
\item The \emph{canonical map} $\kappa\colon\ltwentyA\to R_V$ is the algebra map with $\kappa(x_i)=[\bar x_i]$, the $C_2$-class
  of the $i$-th current. (It is well defined because the boundary relations $x_i\ltwentyQ^m$ vanish in $R_V$,
  Section~20.4.)
\item The \emph{universal arc of $R_V$} is the algebra map
  $\gamma\colon R_V\to\ltwentyB[[t]]$ obtained by composing $\kappa^{-1}\colon R_V\to\ltwentyA$ (Theorem~20.23
  below shows $\kappa$ is bijective) with the tautological arc $\ltwentyA\to\ltwentyB[[t]]$,
  $x_i\mapsto\sum_{n\ge0}x_{i,n}t^n$ of Theorem~20.11.
\item The \emph{zero-jet embedding} $\iota_0\colon R_V\to\ltwentyB$ is $\kappa^{-1}$ followed by the constant-arc inclusion
  $\ltwentyA\to\ltwentyB$, $x_i\mapsto x_{i,0}$; equivalently, the constant coefficient of $\gamma$.
\end{enumerate}
\end{bdefinition}
{\raggedright\noindent\textbf{Lean.} \lean{VOA.}\allowbreak\lean{AffineSl2BoundaryClassicalFreeness.}\allowbreak\lean{canonicalC2Map},\allowbreak\ \lean{VOA.}\allowbreak\lean{AffineSl2BoundaryClassicalFreeness.}\allowbreak\lean{c2UniversalArc},\allowbreak\ \lean{VOA.}\allowbreak\lean{AffineSl2BoundaryClassicalFreeness.}\allowbreak\lean{c2ToArc},\allowbreak\ \lean{VOA.}\allowbreak\lean{AffineSl2BoundaryClassicalFreeness.}\allowbreak\lean{c2ToArc_eq_constantCoeff}\par}
{\raggedright\noindent File \leanpath{VOA/\allowbreak Examples/\allowbreak AffineSl2BoundaryTheorem11.lean}.\par}
\begin{leancode}
def canonicalC2Map (m : ℕ) (hm : 1 ≤ m) :
    (BasePoly ℂ ⧸ baseIdeal (R := ℂ) m) →ₐ[ℂ] A₂ (AffineSl2BoundaryCharacter.vacuum m) :=
  baseC2Map (BoundaryArcLiMap.simpleVOA (uniformCandidate m hm)) (fun _ => 1)
    (fun i => ⟨BoundaryArcLiMap.gen (uniformCandidate m hm) i,
      BoundaryArcLiMap.gen_mem (uniformCandidate m hm) i⟩)
    (baseIdeal m) (uniformCandidate m hm).baseC2Relations

def c2UniversalArc (m : ℕ) (hm : 1 ≤ m) :
    A₂ (AffineSl2BoundaryCharacter.vacuum m) →ₐ[ℂ]
      PowerSeries (InfiniteJet.PresentedJet (baseIdeal (R := ℂ) m)) :=
  (InfiniteJet.universalArc (baseIdeal (R := ℂ) m)).comp (baseC2AlgEquiv m hm).symm.toAlgHom

def c2ToArc (m : ℕ) (hm : 1 ≤ m) :
    A₂ (AffineSl2BoundaryCharacter.vacuum m) →ₐ[ℂ]
      InfiniteJet.PresentedJet (baseIdeal (R := ℂ) m) :=
  (InfiniteJet.baseToJet (baseIdeal (R := ℂ) m)).comp (baseC2AlgEquiv m hm).symm.toAlgHom

theorem c2ToArc_eq_constantCoeff (m : ℕ) (hm : 1 ≤ m)
    (a : A₂ (AffineSl2BoundaryCharacter.vacuum m)) :
    c2ToArc m hm a = PowerSeries.constantCoeff (c2UniversalArc m hm a) := rfl
\end{leancode}
\noindent\textbf{Hypotheses / status.} $1\le m$; unconditional\par
\medskip
\begin{btheorem}{20.26}[Arakawa--Dai, Theorem~1.1, in the paper's form]
Let $m\ge1$, $k=-2+2/(2m+1)$, $V=L_k(\ltwentysl)$ (the quotient of $V^k(\ltwentysl)$ by its maximal proper graded ideal),
$R_V=V/C_2(V)$, $\ltwentyQ=h^2+4ef$. Realise $\ltwentyJ R_V$ as $\ltwentyB=\ltwentyJ\bigl(\C[e,h,f]/(e\ltwentyQ^m,h\ltwentyQ^m,
f\ltwentyQ^m)\bigr)$ with the arc $\gamma\colon R_V\to\ltwentyB[[t]]$ of Definition~20.22. Then:
\begin{enumerate}
\item the canonical map $\kappa\colon\C[e,h,f]/(e\ltwentyQ^m,h\ltwentyQ^m,f\ltwentyQ^m)\to R_V$ is bijective;
\item $\kappa$ sends $e,h,f$ to the $C_2$-classes of the currents $e_{(-1)}\ltwentyvac$, $h_{(-1)}\ltwentyvac$, $f_{(-1)}\ltwentyvac$;
\item $(\ltwentyB,\gamma)$ has the universal property of the arc algebra of $R_V$: for every commutative $\C$-algebra $A$ and
  every algebra map $u\colon R_V\to A[[t]]$ there is a unique $f\colon\ltwentyB\to A$ with $f_*\circ\gamma=u$;
\item the Poisson vertex algebra isomorphism $\pi\colon\ltwentyB\to\ltwentygr V$ restricts on $R_V$ (via the zero-jet
  embedding $\iota_0$) to the canonical inclusion $c\colon R_V\to\ltwentygr V$;
\item $\pi$ is the unique algebra map $\ltwentyB\to\ltwentygr V$ that commutes with the translations and restricts to $c$; that is,
  $\pi$ is the canonical map $\ltwentyJ R_V\to\ltwentygr V$;
\item for every $N$, $\pi(x)$ has conformal weight $N$ if and only if $x$ has conformal weight $N$.
\end{enumerate}
Consequently, the canonical map $\ltwentyJ R_V\to\ltwentygr V$ is an isomorphism of conformally graded Poisson vertex algebras.
\end{btheorem}
{\raggedright\noindent\textbf{Lean.} \lean{VOA.}\allowbreak\lean{AffineSl2BoundaryClassicalFreeness.}\allowbreak\lean{arakawaDai_theorem_1_1}\par}
{\raggedright\noindent File \leanpath{VOA/\allowbreak Examples/\allowbreak AffineSl2BoundaryTheorem11.lean}.\par}
\begin{leancode}
theorem arakawaDai_theorem_1_1 (m : ℕ) (hm : 1 ≤ m) :
    Function.Bijective (canonicalC2Map m hm) ∧
    (∀ i : Fin 3,
      canonicalC2Map m hm (Ideal.Quotient.mk (baseIdeal (R := ℂ) m) (MvPolynomial.X i)) =
        mk₂ (AffineSl2BoundaryCharacter.vacuum m)
          ((QuotVOA.maxGradedIdeal (AffineVOA.affineSl2VOA_closed (level m)
            (level_noncritical m))).mkQ
            (AffineVOA.J (level m) (AffineSl2UniversalC2PBW.symbolToCurrentColor i) (-1)
              (AffineVOA.vacuumA (level m))))) ∧
    (∀ (A : Type) [CommRing A] [Algebra ℂ A]
        (u : A₂ (AffineSl2BoundaryCharacter.vacuum m) →ₐ[ℂ] PowerSeries A),
      ∃! f : InfiniteJet.PresentedJet (baseIdeal (R := ℂ) m) →ₐ[ℂ] A,
        (PowerSeries.mapAlgHom f).comp (c2UniversalArc m hm) = u) ∧
    (∀ a, (arcPoissonEquiv m hm).toAlgEquiv (c2ToArc m hm a) =
      c2ToLiAlgebra (AffineSl2BoundaryCharacter.vacuum m) a) ∧
    (∀ F : InfiniteJet.PresentedJet (baseIdeal (R := ℂ) m) →ₐ[ℂ]
        liAssociatedGraded (AffineSl2BoundaryCharacter.vacuum m),
      (∀ x, F ((arcPVA m).translation x) =
        (liPoissonVertexAlgebra (AffineSl2BoundaryCharacter.vacuum m)).translation (F x)) →
      (∀ a, F (c2ToArc m hm a) = c2ToLiAlgebra (AffineSl2BoundaryCharacter.vacuum m) a) →
      F = (arcPoissonEquiv m hm).toAlgEquiv.toAlgHom) ∧
    (∀ (N : ℕ) (x : InfiniteJet.PresentedJet (baseIdeal (R := ℂ) m)),
      (arcPoissonEquiv m hm).toAlgEquiv x ∈
        BoundaryArcConformalSqueeze.liConformalPiece (AffineSl2BoundaryCharacter.vacuum m) N ↔
        x ∈ arcConformalPiece m N) :=
  ⟨canonicalC2Map_bijective m hm, canonicalC2Map_mk_X m hm,
    fun _ _ _ u => c2UniversalArc_universal m hm u, arcPoissonEquiv_c2ToArc m hm,
    eq_arcPoissonEquiv_of_c2ToArc m hm, arcPoissonEquiv_conformal_iff m hm⟩
\end{leancode}
\noindent\textbf{Hypotheses / status.} $1\le m$; unconditional; $\ltwentyJ R_V$ realised as presented arc algebra with universal arc\par
\medskip
\begin{btheorem}{20.47}[Li's filtration of $V(c,0)$ and of $L(c,0)$ is the word filtration]
For every $c\in\C$ and every $p\in\N$: $F^pV(c,0)=C_p$, and $F^pL(c,0)$ is the image of $C_p$ under the quotient map
$V(c,0)\to L(c,0)$.
\end{btheorem}
{\raggedright\noindent\textbf{Lean.} \lean{VOA.}\allowbreak\lean{VirasoroLi.}\allowbreak\lean{liFiltration_eq_virWordPiece},\allowbreak\ \lean{VOA.}\allowbreak\lean{VirasoroLi.}\allowbreak\lean{simpleLiFiltration_eq_virWordPiece},\allowbreak\ \lean{VOA.}\allowbreak\lean{VirasoroLi.}\allowbreak\lean{virWordNegativeModeClosed}\par}
{\raggedright\noindent File \leanpath{VOA/\allowbreak Filtered/\allowbreak VirasoroWordLi.lean}.\par}
\begin{leancode}
theorem liFiltration_eq_virWordPiece (p : ℕ) :
    liFiltration (virasoroVOA c) p = virWordPiece c p :=

theorem simpleLiFiltration_eq_virWordPiece (p : ℕ) :
    liFiltration (VirSimple.virasoroSimple c) p =
      (virWordPiece c p).map (VirSimple.maxIdealVir c).mkQ := by

theorem virWordNegativeModeClosed (q n : ℕ) (u : VirUniv c) {v : VirUniv c}
    (hv : v ∈ virWordPiece c q) :
    (virasoroVOA c).smul u (-(n : ℤ) - 1) v ∈ virWordPiece c (q + n) := by
\end{leancode}
\noindent\textbf{Hypotheses / status.} Every $c$; unconditional\par
\medskip
\endgroup
\setcounter{chapter}{20}
\chapter{Huang's modular tensor category and the congruence property}
\begingroup
\providecommand{\ltwooneSL}{\mathrm{SL}_2}
\providecommand{\ltwooneRep}{\operatorname{Rep}}
\noindent{\small 9 of the chapter's 30 numbered statements: its headline results and every definition they use.}\par\medskip
\begin{bdefinition}{21.1}[Nice VOA; Chapter~22, Definitions~22.2 and~22.3]
A VOA $V$ is \emph{pseudo-unitary} if every irreducible ordinary module $M\not\cong V$ has all its $L_0$-eigenvalues of
strictly positive real part. $V$ is \emph{nice} if it is strongly rational (simple, of CFT type, self-dual,
$C_2$-cofinite and rational; Chapter~12, Definition~12.43) and pseudo-unitary.
\end{bdefinition}
{\raggedright\noindent\textbf{Lean.} \lean{VOA.}\allowbreak\lean{Classification.}\allowbreak\lean{COne.}\allowbreak\lean{Nice}\par}
{\raggedright\noindent File \leanpath{VOA/\allowbreak Classification/\allowbreak COne/\allowbreak Statement.lean}.\par}
\begin{leancode}
def Nice (V : VertexOperatorAlgebra 𝒱) : Prop :=
  StronglyRational V ∧ PseudoUnitary V
\end{leancode}
\medskip
\begin{btheorem}{21.5}[Huang--Lepowsky: the tensor category]
For a strongly rational $V$ and any choice $F$ of fusion products, the category of ordinary $V$-modules carries a monoidal
structure with tensor product $\boxtimes_F$, unit $V$, the analytically defined associator and the unitors.
\end{btheorem}
{\raggedright\noindent\textbf{Lean.} \lean{VOA.}\allowbreak\lean{Huang.}\allowbreak\lean{TensorCat.}\allowbreak\lean{monoidalCategoryOfStronglyRational_pbw},\allowbreak\ \lean{VOA.}\allowbreak\lean{Huang.}\allowbreak\lean{TensorCat.}\allowbreak\lean{Ordinary},\allowbreak\ \lean{VOA.}\allowbreak\lean{Huang.}\allowbreak\lean{FusionProduct.}\allowbreak\lean{IsUniversal},\allowbreak\ \lean{VOA.}\allowbreak\lean{Huang.}\allowbreak\lean{FusionProduct.}\allowbreak\lean{Product},\allowbreak\ \lean{VOA.}\allowbreak\lean{Huang.}\allowbreak\lean{TensorCat.}\allowbreak\lean{FusionChoice},\allowbreak\ \lean{VOA.}\allowbreak\lean{Huang.}\allowbreak\lean{TensorCat.}\allowbreak\lean{tensorObj}\par}
{\raggedright\noindent File \leanpath{VOA/\allowbreak Huang/\allowbreak PentagonCloseMain.lean}.\par}
\begin{leancode}
abbrev monoidalCategoryOfStronglyRational_pbw (hV : StronglyRational W) :
    MonoidalCategory (Ordinary W) :=
  monoidalCategoryOfStronglyRationalSewing F hV (fivePointSewingInput_stronglyRational_pbw F hV)
    (vacuumCreationCompatibility_pbw F hV)

\end{leancode}
{\raggedright\noindent File \leanpath{VOA/\allowbreak Huang/\allowbreak TensorCatOrdinary.lean}.\par}
\begin{leancode}
structure Ordinary where
  carrier : Type u
  [addCommGroup : AddCommGroup carrier]
  [module : Module ℂ carrier]
  isMod : OrdVModule W carrier
\end{leancode}
{\raggedright\noindent File \leanpath{VOA/\allowbreak Huang/\allowbreak FusionProductBasic.lean}.\par}
\begin{leancode}
def IsUniversal {A₁ : OrdVModule W M₁} {A₂ : OrdVModule W M₂}
    {AT : OrdVModule W T} (Y : IntertwiningOperator A₁ A₂ AT) : Prop :=
  ∀ (N : Type u) [AddCommGroup N] [Module ℂ N] (AN : OrdVModule W N), UniversalFor Y AN
\end{leancode}
{\raggedright\noindent File \leanpath{VOA/\allowbreak Huang/\allowbreak FusionProductBasic.lean}.\par}
\begin{leancode}
structure Product (A₁ : OrdVModule W M₁) (A₂ : OrdVModule W M₂) where
  carrier : Type u
  [addCommGroup : AddCommGroup carrier]
  [module : Module ℂ carrier]
  isMod : OrdVModule W carrier
  intertwiner : IntertwiningOperator A₁ A₂ isMod
  universal : IsUniversal intertwiner
\end{leancode}
{\raggedright\noindent File \leanpath{VOA/\allowbreak Huang/\allowbreak TensorCatBifunctor.lean}.\par}
\begin{leancode}
abbrev FusionChoice (W : VertexOperatorAlgebra V) :=
  ∀ X Y : Ordinary W, Product X.isMod Y.isMod
\end{leancode}
{\raggedright\noindent File \leanpath{VOA/\allowbreak Huang/\allowbreak TensorCatBifunctor.lean}.\par}
\begin{leancode}
def tensorObj (X Y : Ordinary W) : Ordinary W := Ordinary.ofProduct (F X Y)
\end{leancode}
\noindent\textbf{Hypotheses / status.} Construction; \lean{VOA.pbw}\par
\medskip
\begin{btheorem}{21.6}[The braiding]
The braiding given by the skew-symmetry of intertwining operators makes this monoidal category braided.
\end{btheorem}
{\raggedright\noindent\textbf{Lean.} \lean{VOA.}\allowbreak\lean{Huang.}\allowbreak\lean{TensorCat.}\allowbreak\lean{HexagonClose.}\allowbreak\lean{braidedCategory_stronglyRational_pbw},\allowbreak\ \lean{VOA.}\allowbreak\lean{Huang.}\allowbreak\lean{TensorCat.}\allowbreak\lean{BraidRibbon.}\allowbreak\lean{braidedCategoryOfStronglyRational_pbw_braiding},\allowbreak\ \lean{VOA.}\allowbreak\lean{Huang.}\allowbreak\lean{TensorCat.}\allowbreak\lean{BraidRibbon.}\allowbreak\lean{braiding},\allowbreak\ \lean{VOA.}\allowbreak\lean{Huang.}\allowbreak\lean{TensorCat.}\allowbreak\lean{BraidRibbon.}\allowbreak\lean{braidingP_intertwiner}\par}
{\raggedright\noindent File \leanpath{VOA/\allowbreak Huang/\allowbreak HexagonCloseResidual.lean}.\par}
\begin{leancode}
abbrev braidedCategory_stronglyRational_pbw :
    @BraidedCategory (Ordinary W) _ (monoidalCategoryOfStronglyRational_pbw F hV) :=
  braidedCategoryOfStronglyRational_pbw F hV (hexagonInput_pbw F hV)
\end{leancode}
{\raggedright\noindent File \leanpath{VOA/\allowbreak Huang/\allowbreak BraidRibbonHexagon.lean}.\par}
\begin{leancode}
theorem braidedCategoryOfStronglyRational_pbw_braiding (hH : HexagonInput F) (X Y : Ordinary W) :
    @BraidedCategory.braiding (Ordinary W) _ (monoidalCategoryOfStronglyRational_pbw F hV)
      (braidedCategoryOfStronglyRational_pbw F hV hH) X Y = braiding F X Y := rfl
\end{leancode}
{\raggedright\noindent File \leanpath{VOA/\allowbreak Huang/\allowbreak BraidRibbonBraiding.lean}.\par}
\begin{leancode}
abbrev braiding (X Y : Ordinary W) : tensorObj F X Y ≅ tensorObj F Y X :=
  braidingP Phase.upper F X Y
\end{leancode}
{\raggedright\noindent File \leanpath{VOA/\allowbreak Huang/\allowbreak BraidRibbonBraiding.lean}.\par}
\begin{leancode}
@[simp] theorem braidingP_intertwiner (X Y : Ordinary W) (n : ℂ) (x : X.carrier)
    (y : Y.carrier) :
    (braidingP p F X Y).hom.val ((F X Y).intertwiner.val n x y) =
      (skewP p Y X (tensorObj F Y X) (F Y X).intertwiner).val n x y :=
\end{leancode}
\noindent\textbf{Hypotheses / status.} Construction; \lean{VOA.pbw}\par
\medskip
\begin{btheorem}{21.7}[Huang's rigidity theorem, right-rigid form, for nice $V$]
For a nice (strongly rational and pseudo-unitary) $V$ the monoidal category of ordinary modules is right rigid: every
object has a right dual.
\end{btheorem}
{\raggedright\noindent\textbf{Lean.} \lean{VOA.}\allowbreak\lean{Huang.}\allowbreak\lean{HuangRigid.}\allowbreak\lean{rightRigid_ordinary}\par}
{\raggedright\noindent File \leanpath{VOA/\allowbreak Huang/\allowbreak HuangRigidOrdinary.lean}.\par}
\begin{leancode}
theorem rightRigid_ordinary (F : FusionChoice W) (hNice : Nice W) :
    letI := monoidalCategoryOfStronglyRational_pbw F hNice.1
    Nonempty (RightRigidCategory (Ordinary W)) := by
\end{leancode}
\noindent\textbf{Hypotheses / status.} Proved for \emph{nice} $V$; \lean{VOA.pbw}\par
\medskip
\begin{btheorem}{21.8}[Huang's rigidity theorem, two-sided, for nice $V$]
For a nice (strongly rational and pseudo-unitary) $V$ and any choice $F$ of fusion products, the monoidal category of ordinary
$V$-modules of Theorem~21.5 is rigid: every object has both a right dual and a left dual.
\end{btheorem}
{\raggedright\noindent\textbf{Lean.} \lean{VOA.}\allowbreak\lean{Huang.}\allowbreak\lean{HuangRigid.}\allowbreak\lean{rigid_ordinary},\allowbreak\ \lean{VOA.}\allowbreak\lean{Huang.}\allowbreak\lean{HuangRigid.}\allowbreak\lean{rigidCategoryOrdinary}\par}
{\raggedright\noindent File \leanpath{VOA/\allowbreak Huang/\allowbreak RigidOrdinary.lean}.\par}
\begin{leancode}
theorem rigid_ordinary (F : FusionChoice W) (hNice : Nice W) :
    letI := monoidalCategoryOfStronglyRational_pbw F hNice.1
    Nonempty (RigidCategory (Ordinary W)) := by

def rigidCategoryOrdinary (F : FusionChoice W) (hNice : Nice W) :
    letI := monoidalCategoryOfStronglyRational_pbw F hNice.1
    RigidCategory (Ordinary W) :=
  Classical.choice (rigid_ordinary F hNice)
\end{leancode}
\noindent\textbf{Hypotheses / status.} Proved for \emph{nice} $V$; \lean{VOA.pbw}\par
\medskip
\begin{btheorem}{21.10}[$V'\cong V$ and the vacuum-fusion chain, for strongly rational $V$]
Let $V$ be strongly rational. (i) If $A$ is a $V$-module with $A\cong V$, then $A'\cong V$; no pseudo-unitarity is used.
(ii) For a modular-data family $(M_a)$ with $M_{v_0}\cong V$: $N_{ab}^{v_0}\le[M_b\cong M_a']$; $N_{ab}^{v_0}=1$ (in
particular $\neq0$, and $\ge1$) when $M_a'\cong M_b$; and $\dim I\binom{V}{M_b\,M_a}\le1$. So $N_{ab}^{0}=\delta_{b,a'}$
holds without pseudo-unitarity.
\end{btheorem}
{\raggedright\noindent\textbf{Lean.} \lean{VOA.}\allowbreak\lean{contra_vacuum_equiv_of_stronglyRational},\allowbreak\ \lean{VOA.}\allowbreak\lean{Huang.}\allowbreak\lean{VerlindeVacuum.}\allowbreak\lean{famFusion_vacuum_le_of_stronglyRational}\par}
{\raggedright\noindent File \leanpath{VOA/\allowbreak Modules/\allowbreak SelfDualContragredient.lean}.\par}
\begin{leancode}
theorem contra_vacuum_equiv_of_stronglyRational (hSR : StronglyRational W) (A : OrdVModule W M)
    (_hirr : A.IsIrreducible) (e : ModuleEquiv W A (adjointModule W)) :
    Nonempty (ModuleEquiv W (contragredientModule A) (adjointModule W)) :=
  contragredient_equiv_of_selfDual hSR.toStronglyFinite.selfDual A e
\end{leancode}
{\raggedright\noindent File \leanpath{VOA/\allowbreak Huang/\allowbreak VerlindeVacuumFusionSR.lean} (one of five copies).\par}
\begin{leancode}
theorem famFusion_vacuum_le_of_stronglyRational (hSR : StronglyRational W) (v0 : ι)
    (hv0 : Nonempty (ModuleEquiv W (D.Mfam v0).isMod (adjointModule W))) (a b : ι) :
    famFusion D a b v0
      ≤ if Nonempty (ModuleEquiv W (D.Mfam b).isMod (contragredientModule (D.Mfam a).isMod))
        then 1 else 0 := by
\end{leancode}
\medskip
\begin{bdefinition}{21.11}[Huang's Verlinde theorem, fusion form, for nice $V$]
For a nice (strongly rational and pseudo-unitary) $V$ and a complete irredundant family of irreducibles with Zhu's modular matrix $S$, the vacuum row is nonzero,
$S_{0a}=S_{a0}$, and $S_{ax}S_{bx}=S_{0x}\sum_cN_{ab}^cS_{cx}$ with the true fusion rules $N_{ab}^c$.
\end{bdefinition}
{\raggedright\noindent\textbf{Lean.} \lean{VOA.}\allowbreak\lean{Huang.}\allowbreak\lean{VerlindeFusion.}\allowbreak\lean{HuangVerlindeFusionInput},\allowbreak\ \lean{VOA.}\allowbreak\lean{Classification.}\allowbreak\lean{COne.}\allowbreak\lean{huangVerlindeFusionInput_holds}\par}
{\raggedright\noindent File \leanpath{VOA/\allowbreak Huang/\allowbreak VerlindeFusion.lean}.\par}
\begin{leancode}
def HuangVerlindeFusionInput : Prop :=
  ∀ (𝒱 : Type u) [AddCommGroup 𝒱] [Module ℂ 𝒱] (W : VertexOperatorAlgebra 𝒱), Nice W →
    ∀ (ι : Type w) [Fintype ι] [DecidableEq ι] (D : VOA.ModularData W ι),
      ZhuModularityAnalytic D → Irredundant D.Mfam → Complete D.Mfam →
      ∀ v0 : ι, Nonempty (ModuleEquiv W (D.Mfam v0).isMod (adjointModule W)) →
        (∀ x, D.repS v0 x ≠ 0) ∧ (∀ a, D.repS v0 a = D.repS a v0) ∧
        ∀ a b x, D.repS a x * D.repS b x
          = D.repS v0 x * ∑ c, (famFusion D a b c : ℂ) * D.repS c x
\end{leancode}
{\raggedright\noindent File \leanpath{VOA/\allowbreak Classification/\allowbreak COne/\allowbreak MainVerlindeProved.lean}.\par}
\begin{leancode}
theorem huangVerlindeFusionInput_holds : VOA.Huang.VerlindeFusion.HuangVerlindeFusionInput.{u, u + 1} :=
  VOA.Huang.VerlindeVacuum.huangVerlindeFusionInput_of_eigen_rho
    VOA.Huang.GenusOneG1f.fusionEigenInput_holds VOA.Huang.VerlindeRho.rhoNegOneInput_pbw
\end{leancode}
\noindent\textbf{Hypotheses / status.} Proved for \emph{nice} $V$; \lean{VOA.pbw}\par
\medskip
\begin{btheorem}{21.16}[Ng--Schauenburg, kernel half: kernels of modular representations are congruence]
For every modular category there is $N>0$ such that $\Gamma(N)$ acts by scalars in every lift of the categorical projective
$\ltwooneSL(\Z)$-representation.
\end{btheorem}
{\raggedright\noindent\textbf{Lean.} \lean{VOA.}\allowbreak\lean{Huang.}\allowbreak\lean{DLNBridge.}\allowbreak\lean{NSProjectiveKernelAll},\allowbreak\ \lean{VOA.}\allowbreak\lean{Huang.}\allowbreak\lean{DLNBridge.}\allowbreak\lean{nsProjectiveKernelAll_holds},\allowbreak\ \lean{VOA.}\allowbreak\lean{Huang.}\allowbreak\lean{NgSchauenburg.}\allowbreak\lean{IsModularRepresentation},\allowbreak\ \lean{VOA.}\allowbreak\lean{Huang.}\allowbreak\lean{NgSchauenburg.}\allowbreak\lean{InPrincipalCongruence},\allowbreak\ \lean{VOA.}\allowbreak\lean{Huang.}\allowbreak\lean{NgSchauenburg.}\allowbreak\lean{ProjectiveKernelContainsGamma}\par}
{\raggedright\noindent File \leanpath{VOA/\allowbreak Huang/\allowbreak NSCongruenceBridge.lean}.\par}
\begin{leancode}
def NSProjectiveKernelAll : Prop :=
  ∀ (C : Type w) [Category.{v} C] [MonoidalCategory C] [RightRigidCategory C]
    [Preadditive C] [Linear ℂ C] [MonoidalPreadditive C] [MonoidalLinear ℂ C]
    [HasKernels C] [HasFiniteBiproducts C] [BraidedCategory C]
    (F : SphericalFusion C) (M : Modular F),
    ∃ N : ℕ, 0 < N ∧ ∀ ρ : SL(2, ℤ) →* ((F.Label → ℂ) ≃ₗ[ℂ] (F.Label → ℂ)),
      IsModularRepresentation F M ρ → ProjectiveKernelContainsGamma ρ N
\end{leancode}
{\raggedright\noindent File \leanpath{VOA/\allowbreak Huang/\allowbreak NSSMoveC4Assembly.lean}.\par}
\begin{leancode}
theorem nsProjectiveKernelAll_holds : NSProjectiveKernelAll.{v, w} :=
\end{leancode}
{\raggedright\noindent File \leanpath{VOA/\allowbreak Huang/\allowbreak NgSchauenburgCongruence.lean}.\par}
\begin{leancode}
def IsModularRepresentation (F : SphericalFusion C) (M : Modular F)
    (ρ : SL(2, ℤ) →* ((F.Label → ℂ) ≃ₗ[ℂ] (F.Label → ℂ))) : Prop :=
  ∃ a b : ℂˣ,
    (∀ z i, ρ ModularGroup.S z i = (a : ℂ) * ∑ j, categoricalS F i j * z j) ∧
    (∀ z i, ρ ModularGroup.T z i = (b : ℂ) * M.twistScalar i * z i)
\end{leancode}
{\raggedright\noindent File \leanpath{VOA/\allowbreak Huang/\allowbreak NgSchauenburgCongruence.lean}.\par}
\begin{leancode}
def InPrincipalCongruence (N : ℕ) (γ : SL(2, ℤ)) : Prop :=
  ∀ i j : Fin 2, (γ i j : ZMod N) = (1 : Matrix (Fin 2) (Fin 2) (ZMod N)) i j
\end{leancode}
{\raggedright\noindent File \leanpath{VOA/\allowbreak Huang/\allowbreak NgSchauenburgCongruence.lean}.\par}
\begin{leancode}
def ProjectiveKernelContainsGamma (ρ : SL(2, ℤ) →* (W ≃ₗ[ℂ] W)) (N : ℕ) : Prop :=
  ∀ γ, InPrincipalCongruence N γ → ∃ c : ℂˣ, ∀ w, ρ γ w = (c : ℂ) • w
\end{leancode}
\noindent\textbf{Hypotheses / status.} Proved\par
\medskip
\begin{btheorem}{21.20}[The Dong--Lin--Ng congruence property, named and hypothesis-free, for nice $V$]
Let $V$ be a nice VOA, $D$ its modular data, $\rho_V$ Zhu's representation on the labels of $D$ (any Zhu datum $D'$ on them).
\begin{enumerate}
\item (\emph{Character form.}) For every level $N$ of $D$, every $\gamma\in\Gamma(N)$, every label $a$ and every $\tau$ with
  $\operatorname{Im}\tau>0$: $\operatorname{ch}_a(\gamma\tau)=\operatorname{ch}_a(\tau)$.
\item (\emph{Galois symmetry, nontrivial case.}) For every level $N$, every $\ell$ invertible modulo $N$ with
  $\ell\not\equiv\pm1$, and every $\gamma\equiv\operatorname{diag}(\ell,\ell^{-1})\pmod N$, $\gamma$ acts on the characters by a
  signed permutation.
\item (\emph{Kernel form.}) There is $L>0$ with $\Gamma(L)\subseteq\ker\rho_V$.
\item (\emph{Kernel form at every level.}) For every level $N$ of $D$, $\Gamma(N)\subseteq\ker\rho_V$.
\end{enumerate}
\end{btheorem}
{\raggedright\noindent\textbf{Lean.} \lean{VOA.}\allowbreak\lean{Huang.}\allowbreak\lean{CongruenceNamed.}\allowbreak\lean{dlnCongruence_and_galois_holds},\allowbreak\ \lean{VOA.}\allowbreak\lean{Huang.}\allowbreak\lean{CongruenceNamed.}\allowbreak\lean{dlnCongruenceInput_holds},\allowbreak\ \lean{VOA.}\allowbreak\lean{Huang.}\allowbreak\lean{CongruenceNamed.}\allowbreak\lean{galoisSymmetryNontrivialInput_holds},\allowbreak\ \lean{VOA.}\allowbreak\lean{Huang.}\allowbreak\lean{CongruenceNamed.}\allowbreak\lean{zhuCongruenceInput_holds},\allowbreak\ \lean{VOA.}\allowbreak\lean{Huang.}\allowbreak\lean{CongruenceNamed.}\allowbreak\lean{zhu_kernelContainsGamma_at_level}\par}
{\raggedright\noindent File \leanpath{VOA/\allowbreak Huang/\allowbreak CongruenceNamed.lean}.\par}
\begin{leancode}
theorem dlnCongruence_and_galois_holds :
    CongruenceGaloisSplit.DLNCongruenceInput.{u} ∧
      CongruenceGaloisSplit.GaloisSymmetryNontrivialInput.{u} :=
  NSGaloisCore.dln_inputs_of_beyl
    (DLNSchur.oddPrimePowerSchur_iff_largePrimeGaussLowerRelations.mpr
      DLNSchur.largePrimeGaussLowerRelations)
    DLNSchur.twoPowerInflationFromTwo_holds

theorem dlnCongruenceInput_holds : CongruenceGaloisSplit.DLNCongruenceInput.{u} :=
  dlnCongruence_and_galois_holds.1

theorem galoisSymmetryNontrivialInput_holds :
    CongruenceGaloisSplit.GaloisSymmetryNontrivialInput.{u} :=
  dlnCongruence_and_galois_holds.2

theorem zhuCongruenceInput_holds : DLNBridge.ZhuCongruenceInput.{u} :=
  DLNBeyl.zhuCongruence_of_beyl
    (DLNSchur.beylInflationSL2_of_large DLNSchur.largePrimeGaussLowerRelations)
    (DLNBridge.zhuProjectiveCongruence_of_NSKernel DLNBridge.nsProjectiveKernelAll_holds
      ZhuRep.zhuModularRepresentationInput_holds)

theorem zhu_kernelContainsGamma_at_level {𝒱 : Type u} [AddCommGroup 𝒱] [Module ℂ 𝒱]
    {V : VertexOperatorAlgebra 𝒱} (hV : VOA.Classification.COne.Nice V)
    (D : VOA.Classification.COne.ModularData V) (hD : D.IsModularData)
    (D' : VOA.ModularData V D.ι) (hZD : DLNBridge.IsZhuDatum D hD D') {N : ℕ}
    (hN : D.IsLevel N) : DLN.KernelContainsGamma D'.rep N :=
  DLNBridge.kernel_at_level zhuCongruenceInput_holds hV D hD D' hZD hN
\end{leancode}
\noindent\textbf{Hypotheses / status.} \textbf{proved, hypothesis-free}; kernel half of NS~6.8 only; \lean{VOA.pbw}\par
\medskip
\endgroup
\setcounter{chapter}{21}
\chapter{The classification of nice central-charge-one vertex operator algebras (Gannon--Rayhaun)}
\begingroup
\providecommand{\ltwotwoS}{S_{00}}
\providecommand{\ltwotwoch}{\operatorname{ch}}
\providecommand{\ltwotwoH}{\mathbb{H}}
\providecommand{\ltwotwoSL}{\mathrm{SL}_2(\Z)}
\noindent{\small 9 of the chapter's 57 numbered statements: its headline results and every definition they use.}\par\medskip
\begin{bdefinition}{22.2}[Pseudo-unitarity; Gannon--Rayhaun Def.~2.4]
A VOA $V$ is \emph{pseudo-unitary} if for every irreducible ordinary $V$-module $M\not\cong V$ and every
$\alpha\in\C$ with $M_\alpha\neq0$, $\operatorname{Re}\alpha>0$.
\end{bdefinition}
{\raggedright\noindent\textbf{Lean.} \lean{VOA.}\allowbreak\lean{Classification.}\allowbreak\lean{COne.}\allowbreak\lean{PseudoUnitary}\par}
{\raggedright\noindent File \leanpath{VOA/\allowbreak Classification/\allowbreak COne/\allowbreak Statement.lean}.\par}
\begin{leancode}
def PseudoUnitary (V : VertexOperatorAlgebra 𝒱) : Prop :=
  ∀ (M : Type u) [AddCommGroup M] [Module ℂ M] (N : OrdVModule V M),
    N.IsIrreducible → IsEmpty (ModuleEquiv V N (adjointModule V)) →
      ∀ α : ℂ, N.weight α ≠ ⊥ → 0 < α.re
\end{leancode}
\noindent\textbf{Hypotheses / status.} Definitions\par
\medskip
\begin{bdefinition}{22.3}[Nice; Gannon--Rayhaun Def.~2.6]
$V$ is \emph{nice} if it is strongly rational and pseudo-unitary.
\end{bdefinition}
{\raggedright\noindent\textbf{Lean.} \lean{VOA.}\allowbreak\lean{Classification.}\allowbreak\lean{COne.}\allowbreak\lean{Nice}\par}
{\raggedright\noindent File \leanpath{VOA/\allowbreak Classification/\allowbreak COne/\allowbreak Statement.lean}.\par}
\begin{leancode}
def Nice (V : VertexOperatorAlgebra 𝒱) : Prop :=
  StronglyRational V ∧ PseudoUnitary V
\end{leancode}
\noindent\textbf{Hypotheses / status.} Definitions\par
\medskip
\begin{btheorem}{22.20}[GR Lemma~3.11: the theta weights are non-negative]
Suppose a power series $F$ equals $\sum_hA_h\eta\chi_h$ with $A_h\in\N$, $A_0=1$, and also
$F=\sum_{t\in T}\sum_{j\bmod r_t}a_t(j)\,\theta_{j,r_t,t}$ with all $t$ square-free. Then for each $t\in T$,
$w_t:=(r_t\sqrt{2t})^{-1}\sum_ja_t(j)$ is a non-negative real number.
\end{btheorem}
{\raggedright\noindent\textbf{Lean.} \lean{VOA.}\allowbreak\lean{Classification.}\allowbreak\lean{COne.}\allowbreak\lean{theta_weight_nonneg}\par}
{\raggedright\noindent File \leanpath{VOA/\allowbreak Classification/\allowbreak COne/\allowbreak WeightPositivity.lean}.\par}
\begin{leancode}
theorem theta_weight_nonneg (F : PowerSeries ℂ) (A : ℕ → ℕ) (hA0 : A 0 = 1)
    (hVir : F = map (Int.castRingHom ℂ) (etaVirasoroSum A))
    (T : Finset ℕ) (hT : ∀ t ∈ T, Squarefree t) (r : ℕ → ℕ) [∀ t, NeZero (r t)]
    (a : (t : ℕ) → ZMod (r t) → ℂ)
    (hθ : F = ∑ t ∈ T, ∑ j : ZMod (r t), a t j • map (Int.castRingHom ℂ) (thetaSeries (r t) j t))
    {t : ℕ} (ht : t ∈ T) :
    ∃ x : ℝ, 0 ≤ x ∧
      (1 / ((r t : ℂ) * (Real.sqrt (2 * t) : ℂ))) * ∑ j : ZMod (r t), a t j = x :=
\end{leancode}
\noindent\textbf{Hypotheses / status.} None; Lean's 3\par
\medskip
\begin{btheorem}{22.28}[Gannon--Rayhaun Theorem~3.1, forward direction]
Every nice vertex operator algebra of central charge $1$ is isomorphic to some $V_{2m}$ ($m\ge1$), some
$V_{2m}^+$ ($m\ge2$), or a fixed-point algebra $V_2^G$ with $G\cong A_4$, $S_4$ or $A_5$.
\end{btheorem}
{\raggedright\noindent\textbf{Lean.} \lean{VOA.}\allowbreak\lean{Classification.}\allowbreak\lean{COne.}\allowbreak\lean{gannonRayhaun_forward}\par}
{\raggedright\noindent File \leanpath{VOA/\allowbreak Classification/\allowbreak COne/\allowbreak MainUnconditional.lean}.\par}
\begin{leancode}
theorem gannonRayhaun_forward : GannonRayhaunStatement.{u} :=
  gannonRayhaun_beyl_large VOA.Huang.DLNSchur.largePrimeGaussLowerRelations
\end{leancode}
\noindent\textbf{Hypotheses / status.} \textbf{none}; Lean's 3, via thm \lean{VOA.pbw}\par
\medskip
\begin{btheorem}{22.32}[Pairwise distinctness]
A VOA carries at most one label of the list.
\end{btheorem}
{\raggedright\noindent\textbf{Lean.} \lean{VOA.}\allowbreak\lean{Classification.}\allowbreak\lean{COne.}\allowbreak\lean{plusExceptionalResidual},\allowbreak\ \lean{VOA.}\allowbreak\lean{Classification.}\allowbreak\lean{COne.}\allowbreak\lean{exceptionalResidual},\allowbreak\ \lean{VOA.}\allowbreak\lean{Classification.}\allowbreak\lean{COne.}\allowbreak\lean{gannonRayhaunDistinct}\par}
{\raggedright\noindent File \leanpath{VOA/\allowbreak Classification/\allowbreak COne/\allowbreak DistinctResiduals.lean}.\par}
\begin{leancode}
theorem plusExceptionalResidual : PlusExceptionalResidual.{u} := by
\end{leancode}
{\raggedright\noindent File \leanpath{VOA/\allowbreak Classification/\allowbreak COne/\allowbreak DistinctResiduals.lean}.\par}
\begin{leancode}
theorem exceptionalResidual : ExceptionalResidual.{u} := by
\end{leancode}
{\raggedright\noindent File \leanpath{VOA/\allowbreak Classification/\allowbreak COne/\allowbreak DistinctResiduals.lean}.\par}
\begin{leancode}
theorem gannonRayhaunDistinct : GannonRayhaunDistinct.{u} :=
  gannonRayhaunDistinct_of_residuals plusExceptionalResidual exceptionalResidual
\end{leancode}
\noindent\textbf{Hypotheses / status.} \textbf{none}; Lean's 3\par
\medskip
\begin{btheorem}{22.33}[$V_{2m}$ is nice]
For every $m\ge1$, $V_{2m}$ is strongly rational and pseudo-unitary.
\end{btheorem}
{\raggedright\noindent\textbf{Lean.} \lean{VOA.}\allowbreak\lean{Classification.}\allowbreak\lean{COne.}\allowbreak\lean{V2m_stronglyRational},\allowbreak\ \lean{VOA.}\allowbreak\lean{Classification.}\allowbreak\lean{COne.}\allowbreak\lean{V2m_pseudoUnitary},\allowbreak\ \lean{VOA.}\allowbreak\lean{Classification.}\allowbreak\lean{COne.}\allowbreak\lean{V2m_nice}\par}
{\raggedright\noindent File \leanpath{VOA/\allowbreak Classification/\allowbreak COne/\allowbreak LatticeStronglyRational.lean}.\par}
\begin{leancode}
theorem V2m_stronglyRational (m : ℕ) (hm : 1 ≤ m) : StronglyRational (V2m m hm) :=
\end{leancode}
{\raggedright\noindent File \leanpath{VOA/\allowbreak Classification/\allowbreak COne/\allowbreak LatticeNice.lean}.\par}
\begin{leancode}
theorem V2m_pseudoUnitary : PseudoUnitary (V2m m hm) :=
\end{leancode}
{\raggedright\noindent File \leanpath{VOA/\allowbreak Classification/\allowbreak COne/\allowbreak LatticeNice.lean}.\par}
\begin{leancode}
theorem V2m_nice : Nice (V2m m hm) := ⟨V2m_stronglyRational m hm, V2m_pseudoUnitary m hm⟩
\end{leancode}
\noindent\textbf{Hypotheses / status.} None; Lean's 3\par
\medskip
\begin{btheorem}{22.48}[Step 5b: the Jacobi bridge holds]
For every VOA $V$, every subgroup $G\le\Aut V$ (finite or not), every proof of finite generation and of
$C_2$-cofiniteness of $V^G$, the twisted Jacobi bridge holds.
\end{btheorem}
{\raggedright\noindent\textbf{Lean.} \lean{VOA.}\allowbreak\lean{Classification.}\allowbreak\lean{COne.}\allowbreak\lean{McRaeO5b.}\allowbreak\lean{twistedJacobiBridge_holds},\allowbreak\ \lean{VOA.}\allowbreak\lean{Classification.}\allowbreak\lean{COne.}\allowbreak\lean{McRaeO5b.}\allowbreak\lean{residue_of_glue}\par}
{\raggedright\noindent File \leanpath{VOA/\allowbreak Classification/\allowbreak COne/\allowbreak McRaeO5bMain.lean}.\par}
\begin{leancode}
theorem twistedJacobiBridge_holds (G : Subgroup (Aut W)) (hfg : FinGen (orbifold G).adjointRestrict)
    (hC2 : C2Cofinite (orbifoldVOA G)) : TwistedJacobiBridge G hfg hC2 := by
\end{leancode}
{\raggedright\noindent File \leanpath{VOA/\allowbreak Classification/\allowbreak COne/\allowbreak McRaeO5bLaurent.lean}.\par}
\begin{leancode}
theorem residue_of_glue {s t q : ℤ → ℂ} {n₀ n₁ m₁ : ℤ}
    (hs0 : ∀ n, n₀ < n → s n = 0) (ht0 : ∀ n, n < n₁ → t n = 0) (hq0 : ∀ m, m < m₁ → q m = 0)
    {Ψ : ℂ → ℂ} (hΨ : ∀ ζ : ℂ, ζ ≠ 0 → ζ ≠ 1 → DifferentiableAt ℂ Ψ ζ)
    (hS : ∀ ζ : ℂ, 1 < ‖ζ‖ → HasSum (fun n : ℤ => s n * ζ ^ n) (Ψ ζ))
    (hT : ∀ ζ : ℂ, ζ ≠ 0 → ‖ζ‖ < 1 → HasSum (fun n : ℤ => t n * ζ ^ n) (Ψ ζ))
    (hQ : ∀ r : ℂ, r ≠ 0 → ‖r‖ < 1 → HasSum (fun m : ℤ => q m * r ^ m) (Ψ (1 + r)))
    (a ℓ : ℤ) : resInf a ℓ s - resZero a ℓ t = resOne a ℓ q := by
\end{leancode}
\noindent\textbf{Hypotheses / status.} \textbf{none}; Lean's 3, via thm \lean{VOA.pbw}\par
\medskip
\begin{btheorem}{22.54}[Rationality of the orbifolds of $V_2$: McRae's Corollary 4.23]
For every finite subgroup $G\le\Aut V_2$ with $V_2^G$ $C_2$-cofinite, $V_2^G$ is admissibly rational.
\end{btheorem}
{\raggedright\noindent\textbf{Lean.} \lean{VOA.}\allowbreak\lean{Classification.}\allowbreak\lean{COne.}\allowbreak\lean{twistedJacobiBridgeAll_holds},\allowbreak\ \lean{VOA.}\allowbreak\lean{Classification.}\allowbreak\lean{COne.}\allowbreak\lean{finiteV2OrbifoldRationalOfC2}\par}
{\raggedright\noindent File \leanpath{VOA/\allowbreak Classification/\allowbreak COne/\allowbreak MainConverseFinal.lean}.\par}
\begin{leancode}
theorem twistedJacobiBridgeAll_holds : McRaeO7.TwistedJacobiBridgeAll := by
  intro G _ hC2
  exact McRaeO5b.twistedJacobiBridge_holds G _ hC2
\end{leancode}
{\raggedright\noindent File \leanpath{VOA/\allowbreak Classification/\allowbreak COne/\allowbreak MainConverseFinal.lean}.\par}
\begin{leancode}
theorem finiteV2OrbifoldRationalOfC2 : FiniteV2OrbifoldRationalOfC2Fact :=
  McRaeO7.finiteV2OrbifoldRationalOfC2_of_bridge twistedJacobiBridgeAll_holds
\end{leancode}
\noindent\textbf{Hypotheses / status.} \textbf{none}; Lean's 3, via thm \lean{VOA.pbw}\par
\medskip
\begin{btheorem}{22.59}[Gannon--Rayhaun Theorem~3.1 and Proposition~3.2, with the library's rationality convention]
A vertex operator algebra (with carrier in universe $0$) is nice (``rational'' read as \lean{AdmissiblyRational},
Remark~22.4) and has central charge $1$ if and only if it is
isomorphic to one of $V_{2m}$ ($m\ge1$), $V_{2m}^+$ ($m\ge2$), $V_2^{A_4}$, $V_2^{S_4}$, $V_2^{A_5}$; and it is
isomorphic to exactly one of them.
\end{btheorem}
{\raggedright\noindent\textbf{Lean.} \lean{VOA.}\allowbreak\lean{Classification.}\allowbreak\lean{COne.}\allowbreak\lean{gannonRayhaunConverse_holds},\allowbreak\ \lean{VOA.}\allowbreak\lean{Classification.}\allowbreak\lean{COne.}\allowbreak\lean{gannonRayhaun_classification}\par}
{\raggedright\noindent File \leanpath{VOA/\allowbreak Classification/\allowbreak COne/\allowbreak MainConverseFinal.lean}.\par}
\begin{leancode}
theorem gannonRayhaunConverse_holds : GannonRayhaunConverse.{0} :=
  McRaeO7.gannonRayhaunConverse_of_bridge twistedJacobiBridgeAll_holds
\end{leancode}
{\raggedright\noindent File \leanpath{VOA/\allowbreak Classification/\allowbreak COne/\allowbreak MainConverseFinal.lean}.\par}
\begin{leancode}
theorem gannonRayhaun_classification :
    GannonRayhaunStatement.{0} ∧ GannonRayhaunDistinct.{0} ∧ GannonRayhaunConverse.{0} :=
  ⟨gannonRayhaun_forward, gannonRayhaunDistinct, gannonRayhaunConverse_holds⟩
\end{leancode}
\noindent\textbf{Hypotheses / status.} \textbf{none}; Lean's 3, via thm \lean{VOA.pbw}; universe 0\par
\medskip
\endgroup
\setcounter{chapter}{22}
\chapter[Affine VOAs for every simple Lie algebra; the exceptional Lie algebras]{Affine vertex algebras for every simple Lie algebra, and the exceptional Lie algebras}
\begingroup
\providecommand{\lxxiiig}{\mathfrak g}
\providecommand{\lxxiiivac}{|0\rangle}
\providecommand{\lxxiiiV}{V^k(\mathfrak g)}
\providecommand{\lxxiiiL}{L_k(\mathfrak g)}
\providecommand{\lxxiiihv}{h^\vee}
\providecommand{\lxxiiiad}{\operatorname{ad}}
\providecommand{\lxxiiiKil}{\operatorname{Kil}}
\providecommand{\lxxiiiCas}{C}
\noindent{\small 7 of the chapter's 37 numbered statements: its headline results and every definition they use.}\par\medskip
\begin{bproposition}{23.3}[Finite identities]
For Lie data as above:
\begin{enumerate}
\item (Casimir in coordinates) $\sum_{a,d}\big(\sum_b\kappa^{ab}f^{bc}{}_d\big)f^{ad}{}_e=2h\,\delta_{ce}$;
\item (trace of $\lxxiiiad$ vanishes) $\sum_a f^{ac}{}_a=0$ for every $c$ (note $\tr\lxxiiiad(b_c)=\sum_af^{ca}{}_a=-\sum_af^{ac}{}_a$);
\item (central charge numerator) $\sum_{a,b}\kappa^{ab}\kappa_{ab}=|\iota|=\dim\lxxiiig$.
\end{enumerate}
\end{bproposition}
{\raggedright\noindent\textbf{Lean.} \lean{VOA.}\allowbreak\lean{AffineGeneral.}\allowbreak\lean{LieData.}\allowbreak\lean{struct_casimir_contract},\allowbreak\ \lean{VOA.}\allowbreak\lean{AffineGeneral.}\allowbreak\lean{LieData.}\allowbreak\lean{structureConstant_diag_sum},\allowbreak\ \lean{VOA.}\allowbreak\lean{AffineGeneral.}\allowbreak\lean{LieData.}\allowbreak\lean{trace_kappaInv_formTable}\par}
{\raggedright\noindent File \leanpath{VOA/\allowbreak AffineGeneral/\allowbreak LieData.lean}.\par}
\begin{leancode}
theorem struct_casimir_contract (c e : ι) :
    (∑ a, ∑ d, (∑ b, D.kappaInv a b * D.structureConstant b c d) * D.structureConstant a d e) =
      if c = e then 2 * D.h else 0 := by

theorem structureConstant_diag_sum (c : ι) : (∑ a, D.structureConstant a c a) = 0 := by

theorem trace_kappaInv_formTable :
    (∑ a, ∑ b, D.kappaInv a b * D.formTable a b) = Fintype.card ι := by
\end{leancode}
\noindent\textbf{Hypotheses / status.} Unconditional\par
\medskip
\begin{btheorem}{23.10}[Universal affine VOA of a simple Lie algebra]
Let $\lxxiiig$ be a finite-dimensional simple complex Lie algebra.
\begin{enumerate}
\item For every nondegenerate symmetric invariant $\kappa$ there is Lie data indexed by $\{0,\dots,\dim\lxxiiig-1\}$, with
  $h=\frac12\mu$, $\mu$ the Casimir eigenvalue; for every $k$ with $k+h\ne0$ the universal affine VOA exists, with central
  charge $\dim\lxxiiig\cdot k/(k+h)$.
\item The Killing form $\lxxiiiKil(x,y)=\tr(\lxxiiiad x\,\lxxiiiad y)$ is symmetric, invariant and nondegenerate, and in its
  normalization $h=\tfrac12$.
\item For every complex number $\nu\neq0$ the form $\kappa=\lxxiiiKil/(2\nu)$ gives Lie data with Casimir number $h=\nu$, and for
  $k+\nu\neq0$ the VOA $V^k(\lxxiiig)$ in this normalization has central charge $\dim\lxxiiig\cdot k/(k+\nu)$. (Here $\nu$ is a free
  parameter, Lean's \lean{hv}; only for $\nu=\lxxiiihv(\lxxiiig)$ is $\lxxiiiKil/(2\nu)$ the standard normalization, which must be proved
  separately for each type.)
\end{enumerate}
\end{btheorem}
{\raggedright\noindent\textbf{Lean.} \lean{VOA.}\allowbreak\lean{AffineGeneral.}\allowbreak\lean{LieData.}\allowbreak\lean{ofSimple},\allowbreak\ \lean{VOA.}\allowbreak\lean{AffineGeneral.}\allowbreak\lean{simpleData},\allowbreak\ \lean{VOA.}\allowbreak\lean{AffineGeneral.}\allowbreak\lean{universalAffineVOA},\allowbreak\ \lean{VOA.}\allowbreak\lean{AffineGeneral.}\allowbreak\lean{universalAffineVOA_centralCharge},\allowbreak\ \lean{VOA.}\allowbreak\lean{AffineGeneral.}\allowbreak\lean{killingInvariant},\allowbreak\ \lean{VOA.}\allowbreak\lean{AffineGeneral.}\allowbreak\lean{killingAffineVOA},\allowbreak\ \lean{VOA.}\allowbreak\lean{AffineGeneral.}\allowbreak\lean{killing_h},\allowbreak\ \lean{VOA.}\allowbreak\lean{AffineGeneral.}\allowbreak\lean{standardData},\allowbreak\ \lean{VOA.}\allowbreak\lean{AffineGeneral.}\allowbreak\lean{standardData_h},\allowbreak\ \lean{VOA.}\allowbreak\lean{AffineGeneral.}\allowbreak\lean{standardAffineVOA},\allowbreak\ \lean{VOA.}\allowbreak\lean{AffineGeneral.}\allowbreak\lean{standardAffineVOA_centralCharge}\par}
{\raggedright\noindent File \leanpath{VOA/\allowbreak AffineGeneral/\allowbreak Simple.lean} (and
\leanpath{VOA/\allowbreak AffineGeneral/\allowbreak Casimir.lean} for \lean{LieData.ofSimple}).\par}
\begin{leancode}
def ofSimple [FiniteDimensional ℂ 𝔤] [LieAlgebra.IsSimple ℂ 𝔤] (hinv : κ.IsInvariant)
    (hnd : LinearMap.BilinForm.Nondegenerate κ) (b : Module.Basis ι ℂ 𝔤) : LieData κ ι :=

def simpleData (κ : LinearMap.BilinForm ℂ 𝔤) (hinv : κ.IsInvariant)
    (hnd : LinearMap.BilinForm.Nondegenerate κ) : LieData κ (Fin (Module.finrank ℂ 𝔤)) :=
  LieData.ofSimple hinv hnd (Module.finBasis ℂ 𝔤)

def universalAffineVOA (κ : LinearMap.BilinForm ℂ 𝔤) (hinv : κ.IsInvariant)
    (hnd : LinearMap.BilinForm.Nondegenerate κ) (k : ℂ) (hk : k + (simpleData 𝔤 κ hinv hnd).h ≠ 0) :
    VertexOperatorAlgebra (AffineLie.vacuumSpace κ hinv k) :=
  affineVOA (simpleData 𝔤 κ hinv hnd) k hk

theorem universalAffineVOA_centralCharge (κ : LinearMap.BilinForm ℂ 𝔤) (hinv : κ.IsInvariant)
    (hnd : LinearMap.BilinForm.Nondegenerate κ) (k : ℂ) (hk : k + (simpleData 𝔤 κ hinv hnd).h ≠ 0) :
    (universalAffineVOA 𝔤 κ hinv hnd k hk).centralCharge =
      (Module.finrank ℂ 𝔤 : ℂ) * k / (k + (simpleData 𝔤 κ hinv hnd).h) := by

theorem killingInvariant : (killingForm ℂ 𝔤).IsInvariant where

def killingAffineVOA (k : ℂ)
    (hk : k + (simpleData 𝔤 (killingForm ℂ 𝔤) (killingInvariant 𝔤)
      (LieAlgebra.IsKilling.killingForm_nondegenerate ℂ 𝔤)).h ≠ 0) :
    VertexOperatorAlgebra (AffineLie.vacuumSpace (killingForm ℂ 𝔤) (killingInvariant 𝔤) k) :=

theorem killing_h :
    (simpleData 𝔤 (killingForm ℂ 𝔤) (killingInvariant 𝔤)
      (LieAlgebra.IsKilling.killingForm_nondegenerate ℂ 𝔤)).h = 1 / 2 := by

def standardData (hv : ℂ) (hhv : hv ≠ 0) :
    LieData ((2 * hv)⁻¹ • killingForm ℂ 𝔤) (Fin (Module.finrank ℂ 𝔤)) :=

theorem standardData_h (hv : ℂ) (hhv : hv ≠ 0) :
    (standardData 𝔤 hv hhv).h = hv := by

def standardAffineVOA (hv : ℂ) (hhv : hv ≠ 0) (k : ℂ) (hk : k + hv ≠ 0) :
    VertexOperatorAlgebra (Vacuum (standardData 𝔤 hv hhv) k) :=
  affineVOA (standardData 𝔤 hv hhv) k (by rw [standardData_h]; exact hk)

theorem standardAffineVOA_centralCharge (hv : ℂ) (hhv : hv ≠ 0) (k : ℂ)
    (hk : k + hv ≠ 0) :
    (standardAffineVOA 𝔤 hv hhv k hk).centralCharge = (Module.finrank ℂ 𝔤 : ℂ) * k / (k + hv) := by
\end{leancode}
\noindent\textbf{Hypotheses / status.} Simple; $h=\lxxiiihv$ only per type (\lean{*_standard})\par
\medskip
\begin{btheorem}{23.19}[Singular vectors die in $L_k(\mathfrak g)$]
A homogeneous positive-mode singular vector of positive weight lies in the maximal ideal. In particular a quadratic state
$t=\sum_i c_i J^{a_i}_{-1}J^{b_i}_{-1}\lxxiiivac$ lies in it as soon as, for all $x,g$,
\begin{equation*}
\sum_i c_i\Big(\delta_{a_i g}\,k\kappa_{x b_i}+\delta_{b_i g}\,k\kappa_{x a_i}+\sum_c f^{x a_i}{}_c f^{c b_i}{}_g\Big)=0,
\qquad
\sum_i c_i\,k\sum_c f^{x a_i}{}_c\kappa_{c b_i}=0 .
\end{equation*}
\end{btheorem}
{\raggedright\noindent\textbf{Lean.} \lean{VOA.}\allowbreak\lean{AffineGeneral.}\allowbreak\lean{mem_maxIdeal_of_SingP},\allowbreak\ \lean{VOA.}\allowbreak\lean{AffineGeneral.}\allowbreak\lean{quadState},\allowbreak\ \lean{VOA.}\allowbreak\lean{AffineGeneral.}\allowbreak\lean{quadCoef1},\allowbreak\ \lean{VOA.}\allowbreak\lean{AffineGeneral.}\allowbreak\lean{quadCoef2},\allowbreak\ \lean{VOA.}\allowbreak\lean{AffineGeneral.}\allowbreak\lean{quadState_mem_maxIdeal}\par}
{\raggedright\noindent Files \leanpath{VOA/\allowbreak AffineGeneral/\allowbreak Singular.lean},
\leanpath{VOA/\allowbreak AffineGeneral/\allowbreak SingularQuad.lean}.\par}
\begin{leancode}
theorem mem_maxIdeal_of_SingP {t : Vacuum D k} (ht : SingP D k t)
    {w : ℤ} (hw : 0 < w) (hmem : t ∈ gradingZA D k w) : t ∈ maxIdeal D k hk := by

def quadState (terms : List (ℂ × ι × ι)) : Vacuum D k :=
  (terms.map fun t => t.1 • st2 D k t.2.1 t.2.2).sum

def quadCoef1 (terms : List (ℂ × ι × ι)) (x g : ι) : ℂ :=
  (terms.map fun t => t.1 * ((if t.2.1 = g then k * formTable D x t.2.2 else 0)
    + (if t.2.2 = g then k * formTable D x t.2.1 else 0)
    + ∑ c, structureConstant D x t.2.1 c * structureConstant D c t.2.2 g)).sum

def quadCoef2 (terms : List (ℂ × ι × ι)) (x : ι) : ℂ :=
  (terms.map fun t => t.1 * (k * ∑ c, structureConstant D x t.2.1 c * formTable D c t.2.2)).sum

theorem quadState_mem_maxIdeal (hk : k + D.h ≠ 0) (terms : List (ℂ × ι × ι))
    (h1 : ∀ x g, quadCoef1 D k terms x g = 0) (h2 : ∀ x, quadCoef2 D k terms x = 0) :
    quadState D k terms ∈ maxIdeal D k hk :=
  mem_maxIdeal_of_SingP D k hk (singP_quad D k terms h1 h2) (by norm_num) (quadState_mem D k terms)
\end{leancode}
\noindent\textbf{Hypotheses / status.} Weight $>0$\par
\medskip
\begin{bproposition}{23.21}[Kernel facts for current coordinates]
Let $\mathrm{cls}:\C[X_1,\dots,X_n]\to\operatorname{gr}^F L_k(\lxxiiig)$ send $X_i$ to the Li-degree-zero class of the current
$[J^i_{-1}\lxxiiivac]$. Then (i) the kernel is stable under the coadjoint derivations $\mathrm{ad}^*(b)$, $X_i\mapsto\sum f^{bi}{}_cX_c$,
and (ii) if a quadratic state $\sum c_iJ^{a_i}_{-1}J^{b_i}_{-1}\lxxiiivac$ lies in $I_{\max}$, then
$\sum c_iX_{a_i}X_{b_i}\in\ker\mathrm{cls}$.
\end{bproposition}
{\raggedright\noindent\textbf{Lean.} \lean{VOA.}\allowbreak\lean{AffineGeneral.}\allowbreak\lean{clsG},\allowbreak\ \lean{VOA.}\allowbreak\lean{AffineGeneral.}\allowbreak\lean{clsG_adPoly_eq_zero},\allowbreak\ \lean{VOA.}\allowbreak\lean{AffineGeneral.}\allowbreak\lean{clsG_quad_eq_zero}\par}
{\raggedright\noindent File \leanpath{VOA/\allowbreak AffineGeneral/\allowbreak LiPresentation.lean}.\par}
\begin{leancode}
abbrev clsG : MvPolynomial (Fin n) ℂ →ₐ[ℂ] liAssociatedGraded W' :=
  aeval fun i => liDegreeZeroClass W' 1 ⟨g i, hg i⟩

theorem clsG_adPoly_eq_zero (b : Fin n) {p : MvPolynomial (Fin n) ℂ}
    (hp : clsG (simpleVOA D k hk) (genS D k hk) (genS_weight_one D k hk) p = 0) :
    clsG (simpleVOA D k hk) (genS D k hk) (genS_weight_one D k hk)
      (adPolyT (structureConstant D) b p) = 0 :=

theorem clsG_quad_eq_zero (terms : List (ℂ × Fin n × Fin n)) (h : quadState D k terms ∈ maxIdeal D k hk) :
    clsG (simpleVOA D k hk) (genS D k hk) (genS_weight_one D k hk) (quadPoly terms) = 0 := by
\end{leancode}
\noindent\textbf{Hypotheses / status.} Unconditional\par
\medskip
\begin{btheorem}{23.23}[Root-space criterion]
Let $L$ be a complex Lie algebra, $R$ a finite nonempty set, $e,h:R\to L$, $p:R\times R\to\C$, and $-:R\to R$, such that
\begin{enumerate}
\item the $e_\beta$ and $h_\alpha$ together span $L$, and $e_\beta\ne0$;
\item $[h_\alpha,e_\beta]=p(\alpha,\beta)e_\beta$ and $[h_\alpha,h_\beta]=0$;
\item $p(\alpha,\alpha)\ne0$, and the weights separate: for $\alpha\ne\beta$ some $\gamma$ has $p(\gamma,\alpha)\ne p(\gamma,\beta)$;
\item $[e_\alpha,e_{-\alpha}]=c_\alpha h_\alpha$ with $c_\alpha\neq0$;
\item no nonzero element of $\operatorname{span}(h)$ commutes with all $e_\beta$;
\item the graph on $R$ with an edge $\alpha\to\beta$ when $p(\alpha,\beta)\ne0$ is connected (every $\beta$ is reachable from every
  $\alpha$).
\end{enumerate}
Then $L$ is simple.
\end{btheorem}
{\raggedright\noindent\textbf{Lean.} \lean{VOA.}\allowbreak\lean{AffineGeneral.}\allowbreak\lean{RootSpace.}\allowbreak\lean{isSimple_of_roots}\par}
{\raggedright\noindent File \leanpath{VOA/\allowbreak AffineGeneral/\allowbreak RootSpaceSimple.lean}.\par}
\begin{leancode}
theorem isSimple_of_roots [Fintype R] [DecidableEq R] [Nonempty R]
    (hspan : span ℂ (Set.range e ∪ Set.range hb) = ⊤)
    (he0 : ∀ β, e β ≠ 0)
    (hbe : ∀ α β, ⁅hb α, e β⁆ = pair α β • e β)
    (hhb : ∀ α β, ⁅hb α, hb β⁆ = 0)
    (hdiag : ∀ α, pair α α ≠ 0)
    (sep : ∀ α β, α ≠ β → ∃ γ, pair γ α ≠ pair γ β)
    (neg : R → R) (hneg : ∀ α, ∃ c : ℂ, c ≠ 0 ∧ ⁅e α, e (neg α)⁆ = c • hb α)
    (center : ∀ y ∈ span ℂ (Set.range hb), (∀ β, ⁅y, e β⁆ = 0) → y = 0)
    (conn : ∀ α β, Relation.ReflTransGen (fun a b => pair a b ≠ 0) α β) :
    LieAlgebra.IsSimple ℂ L where
\end{leancode}
\noindent\textbf{Hypotheses / status.} Listed\par
\medskip
\begin{btheorem}{23.27}[Frenkel--Kac brackets, spanning, dimension and simplicity]
Let $L$ be even and positive definite.
\begin{enumerate}
\item $[h_\alpha,e_\beta]=\langle\alpha,\beta\rangle e_\beta$, $[h_\alpha,h_\beta]=0$, $[e_\alpha,e_{-\alpha}]=\varepsilon(\alpha,-\alpha)h_\alpha$, $e_\beta\ne0$.
\item If the roots generate $L$, the $e_\beta$ and $h_\beta$ span $(V_L)_1$.
\item $\dim(V_L)_1=\operatorname{rank}L+\#\{\beta:\langle\beta,\beta\rangle=2\}$.
\item If a family $s$ of roots generates $L$ and is connected (the graph $\langle s_i,s_j\rangle\ne0$ is connected), $(V_L)_1$ is simple.
\end{enumerate}
\end{btheorem}
{\raggedright\noindent\textbf{Lean.} \lean{VOA.}\allowbreak\lean{AffineGeneral.}\allowbreak\lean{LatticeLie.}\allowbreak\lean{hV_eV},\allowbreak\ \lean{VOA.}\allowbreak\lean{AffineGeneral.}\allowbreak\lean{LatticeLie.}\allowbreak\lean{hV_hV},\allowbreak\ \lean{VOA.}\allowbreak\lean{AffineGeneral.}\allowbreak\lean{LatticeLie.}\allowbreak\lean{eV_eV_neg},\allowbreak\ \lean{VOA.}\allowbreak\lean{AffineGeneral.}\allowbreak\lean{LatticeLie.}\allowbreak\lean{eV_ne_zero},\allowbreak\ \lean{VOA.}\allowbreak\lean{AffineGeneral.}\allowbreak\lean{LatticeLie.}\allowbreak\lean{span_eV_hV},\allowbreak\ \lean{VOA.}\allowbreak\lean{AffineGeneral.}\allowbreak\lean{LatticeLie.}\allowbreak\lean{finrank_V1},\allowbreak\ \lean{VOA.}\allowbreak\lean{AffineGeneral.}\allowbreak\lean{LatticeLie.}\allowbreak\lean{isSimple_of_simpleRoots}\par}
{\raggedright\noindent File \leanpath{VOA/\allowbreak AffineGeneral/\allowbreak LatticeV1Simple.lean}.\par}
\begin{leancode}
theorem hV_eV (α β : Root L hPos) :
    ⁅hV L hEven hPos α, eV L hEven hPos β⁆ = ((L.form α.1 β.1 : ℤ) : ℂ) • eV L hEven hPos β := by

theorem hV_hV (α β : Root L hPos) : ⁅hV L hEven hPos α, hV L hEven hPos β⁆ = 0 := by

theorem eV_eV_neg (α : Root L hPos) :
    ⁅eV L hEven hPos α, eV L hEven hPos (Root.neg L hPos α)⁆
      = latticeCocycle L α.1 (-α.1) • hV L hEven hPos α := by

theorem eV_ne_zero (β : Root L hPos) : eV L hEven hPos β ≠ 0 := by

theorem span_eV_hV (hgen : RootsGenerate L) :
    span ℂ (Set.range (eV L hEven hPos) ∪ Set.range (hV L hEven hPos)) = ⊤ := by

theorem finrank_V1 : Module.finrank ℂ (V1 L hEven hPos) = rank L + (normEq L hPos 2).card := by

theorem isSimple_of_simpleRoots {ι : Type*} [Nonempty ι] (s : ι → Root L hPos)
    (hsgen : AddSubgroup.closure (Set.range fun i => (s i).1) = ⊤)
    (hsconn : ∀ i j, Relation.ReflTransGen (fun i j => L.form (s i).1 (s j).1 ≠ 0) i j) :
    LieAlgebra.IsSimple ℂ (V1 L hEven hPos) := by
\end{leancode}
\noindent\textbf{Hypotheses / status.} Connected generating roots\par
\medskip
\begin{btheorem}{23.32}[$\lxxiiiKil/(2\lxxiiihv)$ is the standard normalization]
Let $L$ be generated by $n=\operatorname{rank}L$ simple roots $s$. Then the root vectors together with the simple coroots form a basis of
$(V_L)_1$, and $\lxxiiiKil(h_\alpha,h_\alpha)=\sum_{\beta}\langle\alpha,\beta\rangle^2$ (sum over roots). For $A_2,D_4,E_6,E_7,E_8$ this sum is
$12,24,48,72,120=4\lxxiiihv$ for every root $\alpha$, so $\big(\lxxiiiKil/(2\lxxiiihv)\big)(h_\alpha,h_\alpha)=2$.
\end{btheorem}
{\raggedright\noindent\textbf{Lean.} \lean{VOA.}\allowbreak\lean{AffineGeneral.}\allowbreak\lean{LatticeLie.}\allowbreak\lean{v1Basis},\allowbreak\ \lean{VOA.}\allowbreak\lean{AffineGeneral.}\allowbreak\lean{LatticeLie.}\allowbreak\lean{killing_hV_hV},\allowbreak\ \lean{VOA.}\allowbreak\lean{AffineGeneral.}\allowbreak\lean{Normalization.}\allowbreak\lean{e8_killing},\allowbreak\ \lean{VOA.}\allowbreak\lean{AffineGeneral.}\allowbreak\lean{Normalization.}\allowbreak\lean{e8_standard}\par}
{\raggedright\noindent Files \leanpath{VOA/\allowbreak AffineGeneral/\allowbreak LatticeKilling.lean},
\leanpath{VOA/\allowbreak AffineGeneral/\allowbreak LatticeNormalization.lean}.\par}
\begin{leancode}
def v1Basis (hrank : rank L = n) : Module.Basis (Root L hPos ⊕ Fin n) ℂ (V1 L hEven hPos) :=
  basisOfTopLeSpanOfCardEqFinrank (Sum.elim (eV L hEven hPos) fun i => hV L hEven hPos (s i))
    (span_eV_hV_simple hEven s hsgen).ge
    (by rw [finrank_V1, Fintype.card_sum, Fintype.card_fin, hrank, Fintype.card_coe, add_comm])

theorem killing_hV_hV (hrank : rank L = n) (α : Root L hPos) :
    killingForm ℂ (V1 L hEven hPos) (hV L hEven hPos α) (hV L hEven hPos α)
      = ∑ β : Root L hPos, ((L.form α.1 β.1 : ℤ) : ℂ) ^ 2 := by

theorem e8_killing (β : Root e8 e8_posDef) :
    killingForm ℂ e8Lie (hV e8 e8_even e8_posDef β) (hV e8 e8_even e8_posDef β) = 120 := by

theorem e8_standard (β : Root e8 e8_posDef) :
    ((2 * 30 : ℂ)⁻¹ • killingForm ℂ e8Lie) (hV e8 e8_even e8_posDef β)
      (hV e8 e8_even e8_posDef β) = 2 := by
\end{leancode}
\noindent\textbf{Hypotheses / status.} Unconditional\par
\medskip
\endgroup
\setcounter{chapter}{23}
\chapter[The Deligne series I: presentations, carriers, divisor comparison]{The Deligne series I: presentations, carriers and the divisor comparison}
\begingroup
\providecommand{\ltwentyfourg}{\mathfrak g}
\providecommand{\ltwentyfourJ}{J_\infty}
\providecommand{\ltwentyfourgr}{\operatorname{gr}^{\mathrm{Li}}}
\providecommand{\ltwentyfourIns}{I_2^{\mathrm{ns}}}
\providecommand{\ltwentyfourRt}{\widetilde R}
\providecommand{\ltwentyfourHilb}{\operatorname{Hilb}}
\providecommand{\ltwentyfourch}{\operatorname{ch}}
\providecommand{\ltwentyfourvac}{\mathbf 1}
\noindent{\small 14 of the chapter's 38 numbered statements: its headline results and every definition they use.}\par\medskip
\begin{btheorem}{24.6}[Exact presentation from classical freeness]
Let $I\subseteq\ker(\operatorname{aeval}x)$. If $(x,I)$ is classically free, then $\ker(\operatorname{aeval}x)=I$.
\end{btheorem}
{\raggedright\noindent\textbf{Lean.} \lean{VOA.}\allowbreak\lean{presentedJetOrderZero},\allowbreak\ \lean{VOA.}\allowbreak\lean{presentedJetRetraction},\allowbreak\ \lean{VOA.}\allowbreak\lean{ker_eq_of_liPresentationClassicallyFree}\par}
{\raggedright\noindent File \leanpath{VOA/\allowbreak DeligneAudit/\allowbreak C2PresentationOfFreeness.lean}.\par}
\begin{leancode}
noncomputable def presentedJetOrderZero (I : Ideal (MvPolynomial σ ℂ)) :
    MvPolynomial σ ℂ →ₐ[ℂ] PresentedJet (R := ℂ) I :=
  (presentedJetMk (R := ℂ) I).comp (MvPolynomial.rename fun i : σ => (i, 0))

noncomputable def presentedJetRetraction (I : Ideal (MvPolynomial σ ℂ)) :
    PresentedJet (R := ℂ) I →ₐ[ℂ] (MvPolynomial σ ℂ ⧸ I) :=

theorem ker_eq_of_liPresentationClassicallyFree (W : VertexOperatorAlgebra V)
    (x : σ → liAssociatedGraded W) (I : Ideal (MvPolynomial σ ℂ))
    (hI : I ≤ RingHom.ker (MvPolynomial.aeval x).toRingHom)
    (hfree : LiPresentationClassicallyFree W x I hI) :
    RingHom.ker (MvPolynomial.aeval x).toRingHom = I := by
\end{leancode}
\noindent\textbf{Hypotheses / status.} Unconditional\par
\medskip
\begin{bdefinition}{24.8}[The $A_2$ relation ideal]
In the coordinates $X_0,\dots,X_7$ of $e_1,e_2,e_{12},f_1,f_2,f_{12},h_1,h_2\in\mathfrak{sl}_3$, put
$Z=3\sum_aX_ag^a$ (with $g^a$ the trace-dual basis) and let $P_a$ be one third of the entry of $Z^2$ at the
position of the $a$-th generator (and of the two diagonal differences for $a=6,7$). Then
$\ltwentyfourIns=\operatorname{span}(P_0,\dots,P_7)\cong V(\theta)$, and $I_2^{\mathrm{ns}}$ (the ideal)
is generated by these eight quadrics.
\end{bdefinition}
{\raggedright\noindent\textbf{Lean.} \lean{VOA.}\allowbreak\lean{DeligneAudit.}\allowbreak\lean{A2.}\allowbreak\lean{q},\allowbreak\ \lean{VOA.}\allowbreak\lean{DeligneAudit.}\allowbreak\lean{A2.}\allowbreak\lean{relationPoly},\allowbreak\ \lean{VOA.}\allowbreak\lean{DeligneAudit.}\allowbreak\lean{A2.}\allowbreak\lean{I2ns}\par}
{\raggedright\noindent File \leanpath{VOA/\allowbreak DeligneAudit/\allowbreak A2Relations.lean}.\par}
\begin{leancode}
abbrev q (a b : Fin 8) : MvPolynomial (Fin 8) ℂ := X a * X b

def relationPoly : Fin 8 → MvPolynomial (Fin 8) ℂ :=
  ![ -- weight α₁:  e₁ (h₁ + 2h₂) + 3 e₁₂ f₂
     q 0 6 + 2 * q 0 7 + 3 * q 2 4,
     -- weight α₂:  -e₂ (2h₁ + h₂) + 3 e₁₂ f₁
     -(2 * q 1 6) - q 1 7 + 3 * q 2 3,
     -- weight θ (highest):  e₁₂ (h₁ - h₂) + 3 e₁ e₂
     q 2 6 - q 2 7 + 3 * q 0 1,
     -- weight -α₁:  f₁ (h₁ + 2h₂) + 3 e₂ f₁₂
     q 3 6 + 2 * q 3 7 + 3 * q 1 5,
     -- weight -α₂:  -f₂ (2h₁ + h₂) + 3 e₁ f₁₂
     -(2 * q 4 6) - q 4 7 + 3 * q 0 5,
     -- weight -θ:  f₁₂ (h₁ - h₂) + 3 f₁ f₂
     q 5 6 - q 5 7 + 3 * q 3 4,
     -- weight 0:  h₁² + 2 h₁h₂ + 3 e₁₂ f₁₂ - 3 e₂ f₂
     q 6 6 + 2 * q 6 7 + 3 * q 2 5 - 3 * q 1 4,
     -- weight 0:  -h₂² - 2 h₁h₂ - 3 e₁₂ f₁₂ + 3 e₁ f₁
     -q 7 7 - 2 * q 6 7 - 3 * q 2 5 + 3 * q 0 3 ]

def I2ns : Ideal (MvPolynomial (Fin 8) ℂ) := Ideal.span (Set.range relationPoly)
\end{leancode}
\medskip
\begin{btheorem}{24.12}[The explicit presentations for all seven types]
For each $X\in\{A_2,G_2,D_4,F_4,E_6,E_7,E_8\}$ and the simple affine VOA $W_X=L_k(\ltwentyfourg_X)$ at the Deligne
level, the ideal generated by the explicit relations of $\ltwentyfourIns$ is contained in the kernel of the
current-coordinate map, so $(\mathrm{currents},(\ltwentyfourIns))$ is a current $C_2$ presentation.
\end{btheorem}
{\raggedright\noindent\textbf{Lean.} \lean{VOA.}\allowbreak\lean{DeligneAudit.}\allowbreak\lean{A2.}\allowbreak\lean{I2ns_le_ker},\allowbreak\ \lean{VOA.}\allowbreak\lean{DeligneAudit.}\allowbreak\lean{A2.}\allowbreak\lean{explicitPresentation},\allowbreak\ \lean{VOA.}\allowbreak\lean{DeligneAudit.}\allowbreak\lean{D4Lie.}\allowbreak\lean{I2ns_le_ker4},\allowbreak\ \lean{VOA.}\allowbreak\lean{DeligneAudit.}\allowbreak\lean{D4Lie.}\allowbreak\lean{explicitPresentation4},\allowbreak\ \lean{VOA.}\allowbreak\lean{DeligneAudit.}\allowbreak\lean{E7Lie.}\allowbreak\lean{I2ns_le_ker7},\allowbreak\ \lean{VOA.}\allowbreak\lean{DeligneAudit.}\allowbreak\lean{E7Lie.}\allowbreak\lean{explicitPresentation7},\allowbreak\ \lean{VOA.}\allowbreak\lean{DeligneAudit.}\allowbreak\lean{E8Lie.}\allowbreak\lean{I2ns_le_ker8},\allowbreak\ \lean{VOA.}\allowbreak\lean{DeligneAudit.}\allowbreak\lean{E8Lie.}\allowbreak\lean{explicitPresentation8}\par}
{\raggedright\noindent File \leanpath{VOA/\allowbreak DeligneAudit/\allowbreak A2RelationsPresentation.lean}.\par}
\begin{leancode}
theorem I2ns_le_ker : I2ns ≤ RingHom.ker
    (MvPolynomial.aeval fun i =>
      liDegreeZeroClass W 1 ⟨gen i, gen_weight_one i⟩).toRingHom :=
  I2ns_le_ker_of_two hT_A2 clsG_relationPoly_two

def explicitPresentation : CurrentC2Presentation .A2 W where
  gen := gen
  gen_weight_one := gen_weight_one
  ideal := I2ns
  relations := I2ns_le_ker
\end{leancode}
{\raggedright\noindent File \leanpath{VOA/\allowbreak DeligneAudit/\allowbreak G2LiePresentation.lean}.\par}
\begin{leancode}
theorem I2ns_le_ker : G2data.I ≤ RingHom.ker
    (MvPolynomial.aeval fun i =>
      liDegreeZeroClass g2W 1 ⟨genS g2CertData (-5/3) g2_level_ne i,
        genS_weight_one g2CertData (-5/3) g2_level_ne i⟩).toRingHom := by

def explicitPresentation : CurrentC2Presentation .G2 g2W where
  gen := genS g2CertData (-5/3) g2_level_ne
  gen_weight_one := genS_weight_one g2CertData (-5/3) g2_level_ne
  ideal := G2data.I
  relations := I2ns_le_ker
\end{leancode}
{\raggedright\noindent File \leanpath{VOA/\allowbreak DeligneAudit/\allowbreak D4LiePresentation.lean}.\par}
\begin{leancode}
theorem I2ns_le_ker4 : DeligneCarrier.D4.D4data.I ≤ RingHom.ker
    (MvPolynomial.aeval fun i => liDegreeZeroClass W4 1 ⟨gen4 i, gen4_weight_one i⟩).toRingHom := by

def explicitPresentation4 : CurrentC2Presentation .D4 W4 where
  gen := gen4
  gen_weight_one := gen4_weight_one
  ideal := DeligneCarrier.D4.D4data.I
  relations := I2ns_le_ker4
\end{leancode}
{\raggedright\noindent File \leanpath{VOA/\allowbreak DeligneAudit/\allowbreak F4LiePresentation.lean}.\par}
\begin{leancode}
theorem I2ns_le_ker : F4data.I ≤ RingHom.ker
    (MvPolynomial.aeval fun i =>
      liDegreeZeroClass f4W 1 ⟨genS f4CertData (-5/2) f4_level_ne i,
        genS_weight_one f4CertData (-5/2) f4_level_ne i⟩).toRingHom := by

def explicitPresentation : CurrentC2Presentation .F4 f4W where
  gen := genS f4CertData (-5/2) f4_level_ne
  gen_weight_one := genS_weight_one f4CertData (-5/2) f4_level_ne
  ideal := F4data.I
  relations := I2ns_le_ker
\end{leancode}
{\raggedright\noindent File \leanpath{VOA/\allowbreak DeligneAudit/\allowbreak E6Presentation.lean}.\par}
\begin{leancode}
theorem I2ns_le_ker : E6data.I ≤ RingHom.ker
    (MvPolynomial.aeval fun i =>
      liDegreeZeroClass W6 1 ⟨genS e6LieData k6 hk6 i, genS_weight_one e6LieData k6 hk6 i⟩).toRingHom := by

def explicitPresentation : CurrentC2Presentation .E6 W6 where
  gen := genS e6LieData k6 hk6
  gen_weight_one := genS_weight_one e6LieData k6 hk6
  ideal := E6data.I
  relations := I2ns_le_ker
\end{leancode}
{\raggedright\noindent File \leanpath{VOA/\allowbreak DeligneAudit/\allowbreak E7LiePresentation.lean}.\par}
\begin{leancode}
theorem I2ns_le_ker7 : DeligneCarrier.E7.E7data.I ≤ RingHom.ker
    (MvPolynomial.aeval fun i => liDegreeZeroClass W7 1 ⟨gen7 i, gen7_weight_one i⟩).toRingHom := by

def explicitPresentation7 : CurrentC2Presentation .E7 W7 where
  gen := gen7
  gen_weight_one := gen7_weight_one
  ideal := DeligneCarrier.E7.E7data.I
  relations := I2ns_le_ker7
\end{leancode}
{\raggedright\noindent File \leanpath{VOA/\allowbreak DeligneAudit/\allowbreak E8Presentation.lean}.\par}
\begin{leancode}
theorem I2ns_le_ker8 : E8data.I ≤ RingHom.ker
    (MvPolynomial.aeval fun i => liDegreeZeroClass W8 1 ⟨gen8 i, gen8_weight_one i⟩).toRingHom := by

def explicitPresentation8 : CurrentC2Presentation .E8 W8 where
  gen := gen8
  gen_weight_one := gen8_weight_one
  ideal := E8data.I
  relations := I2ns_le_ker8
\end{leancode}
\noindent\textbf{Hypotheses / status.} \textbf{unconditional}\par
\medskip
\begin{btheorem}{24.14}[Current-jet surjectivity for every Deligne type]
(i) For any case, VOA and presentation, if the currents span Li's filtration then
\lean{CurrentJetSurjectivity} holds. (ii) For every Lie data $D$ and non-critical $k$, the currents of
$L_k(\ltwentyfourg)$ span Li's filtration. Hence current-jet surjectivity holds for every explicit
presentation of Theorem~24.12.
\end{btheorem}
{\raggedright\noindent\textbf{Lean.} \lean{VOA.}\allowbreak\lean{DeligneAudit.}\allowbreak\lean{currentJetSurjectivity_of_genWordSpan},\allowbreak\ \lean{VOA.}\allowbreak\lean{DeligneAudit.}\allowbreak\lean{A2.}\allowbreak\lean{genWordsSpan},\allowbreak\ \lean{VOA.}\allowbreak\lean{GeneratorModeTracked},\allowbreak\ \lean{VOA.}\allowbreak\lean{generatorWordsSpanLiFiltration_of_tracked},\allowbreak\ \lean{VOA.}\allowbreak\lean{generatorModeTracked_of_current_comm},\allowbreak\ \lean{VOA.}\allowbreak\lean{generatorWordsSpanLiFiltration_quotient}\par}
{\raggedright\noindent File \leanpath{VOA/\allowbreak DeligneAudit/\allowbreak A2CurrentJetSurjectivity.lean}.\par}
\begin{leancode}
theorem currentJetSurjectivity_of_genWordSpan {V : Type*} [AddCommGroup V] [Module ℂ V]
    {case : HigherRankDeligneCase} {W : VertexOperatorAlgebra V}
    (P : CurrentC2Presentation case W)
    (h : GeneratorWordsSpanLiFiltration W P.gen (fun _ => 1)) :
    CurrentJetSurjectivity case W P :=
\end{leancode}
{\raggedright\noindent File \leanpath{VOA/\allowbreak AffineGeneral/\allowbreak DeligneInstance.lean}.\par}
\begin{leancode}
theorem genWordsSpan :
    GeneratorWordsSpanLiFiltration (simpleVOA D k hk) (presentation case D k hk).gen (fun _ => 1) :=
\end{leancode}
{\raggedright\noindent File \leanpath{VOA/\allowbreak Filtered/\allowbreak GeneratorCurrentLiSpanning.lean}.\par}
\begin{leancode}
def GeneratorModeTracked : Prop :=
  ∀ (i : Fin g) (m : ℤ) (q : ℕ) (v : V), v ∈ genWordPiece W gen wt q →
    W.smul (gen i) m v ∈ genWordPiece W gen wt (currentModeTarget q m)

theorem generatorWordsSpanLiFiltration_of_tracked (htrack : GeneratorModeTracked W gen wt)
    (hspan : Submodule.span ℂ (Set.range (genWordState W gen wt)) = ⊤) :
    GeneratorWordsSpanLiFiltration W gen wt := by

theorem generatorModeTracked_of_current_comm (hJ : ∀ i m v, W.smul (gen i) m v = J i m v)
    (hcomm : ∀ i j m p, J i m * J j p - J j p * J i m =
      (∑ c, f i j c • J c (m + p)) +
        (if m + p = 0 then cen i j m • (1 : Module.End ℂ V) else 0))
    (hvac : ∀ i m, 0 ≤ m → J i m W.vacuum = 0) :
    GeneratorModeTracked W gen wt := by

theorem generatorWordsSpanLiFiltration_quotient (h : GeneratorWordsSpanLiFiltration W gen wt) :
    GeneratorWordsSpanLiFiltration (QuotVOA.quotientVOA W I hI hgr)
      (fun i => I.mkQ (gen i)) wt := by
\end{leancode}
\noindent\textbf{Hypotheses / status.} Unconditional\par
\medskip
\begin{bdefinition}{24.16}[The current carrier $W_m$ and the auxiliary carrier]
Fix a basis $\sigma$ of $\ltwentyfourg^*$ and quadrics $Q$. A \emph{tensor polynomial} with $m$ current and $d$
auxiliary particles is a function $t\mapsto p_t$ from labellings $t:\{1..m\}\to\sigma$ to polynomials in
$z_1,\dots,z_{m+d}$. It is \emph{symmetric} if simultaneous permutation of roots and labels fixes it. For a
pair $i<j$ of current slots, a relation $q$ and spectator labels $t$, the \emph{collision readback} is
$\sum_{x,y}q(x,y)\,p_{t[i\mapsto x,j\mapsto y]}|_{z_j=z_i}$. The \emph{current carrier}
$\mathrm{currentCarrier}(Q,m,d)$ consists of the symmetric tensor polynomials all of whose readbacks vanish.
The \emph{auxiliary carrier} additionally requires, for every pair involving an auxiliary particle, that the
polynomial vanish on that collision. $\mathrm{degreePiece}(C,r)$ is the subspace of $C$ of polynomials
homogeneous of degree $r$. $W_m:=\mathrm{currentCarrier}(I_2,m,0)$.
\end{bdefinition}
{\raggedright\noindent\textbf{Lean.} \lean{VOA.}\allowbreak\lean{QuadraticCurrentFunctionalCarrier.}\allowbreak\lean{TensorPolynomial},\allowbreak\ \lean{VOA.}\allowbreak\lean{QuadraticCurrentFunctionalCarrier.}\allowbreak\lean{Symmetric},\allowbreak\ \lean{VOA.}\allowbreak\lean{QuadraticCurrentFunctionalCarrier.}\allowbreak\lean{currentReadback},\allowbreak\ \lean{VOA.}\allowbreak\lean{QuadraticCurrentFunctionalCarrier.}\allowbreak\lean{CurrentEquations},\allowbreak\ \lean{VOA.}\allowbreak\lean{QuadraticCurrentFunctionalCarrier.}\allowbreak\lean{currentCarrier},\allowbreak\ \lean{VOA.}\allowbreak\lean{QuadraticCurrentFunctionalCarrier.}\allowbreak\lean{auxiliaryCarrier},\allowbreak\ \lean{VOA.}\allowbreak\lean{DeligneAuxiliaryCollision.}\allowbreak\lean{Particle},\allowbreak\ \lean{VOA.}\allowbreak\lean{DeligneAuxiliaryCollision.}\allowbreak\lean{conductor},\allowbreak\ \lean{VOA.}\allowbreak\lean{DeligneAuxiliaryCollision.}\allowbreak\lean{conductorDegree},\allowbreak\ \lean{VOA.}\allowbreak\lean{QuadraticCurrentFunctionalCarrier.}\allowbreak\lean{degreePiece}\par}
{\raggedright\noindent File \leanpath{VOA/\allowbreak Commutative/\allowbreak QuadraticCurrentFunctionalCarrier.lean}.\par}
\begin{leancode}
abbrev TensorPolynomial (m d : ℕ) :=
  (Fin m → σ) → MvPolynomial (Particle m d) K

def Symmetric {m d : ℕ} (p : TensorPolynomial (K := K) (σ := σ) m d) : Prop :=
  ∀ (e : Equiv.Perm (Particle m d)) (g : Equiv.Perm (Fin m)),
    ExtendsCurrentPermutation e g → ∀ t,
      MvPolynomial.rename e (p t) = p (t ∘ g.symm)

def currentReadback {m d : ℕ} (q : σ × σ → K)
    (ij : IncreasingPair (Fin m)) (t : Fin m → σ)
    (p : TensorPolynomial (K := K) (σ := σ) m d) :=
  projectedDiagonal (ij.1.1.castAdd d) (ij.1.2.castAdd d)
    ((Fin.castAdd_injective m d).ne (ne_of_gt ij.2)) q
    (fun xy => p (insertPair ij t xy))

def CurrentEquations {m d : ℕ} (Q : ρ → σ × σ → K)
    (p : TensorPolynomial (K := K) (σ := σ) m d) : Prop :=
  ∀ r ij t, currentReadback (Q r) ij t p = 0

def currentCarrier (Q : ρ → σ × σ → K) (m d : ℕ) :
    Submodule K (TensorPolynomial (K := K) (σ := σ) m d) where
  carrier := {p | Symmetric p ∧ CurrentEquations Q p}

def auxiliaryCarrier (Q : ρ → σ × σ → K) (m d : ℕ) :
    Submodule K (TensorPolynomial (K := K) (σ := σ) m d) where
  carrier := {p | p ∈ currentCarrier Q m d ∧
    ∀ t (ij : Pair m d), m ≤ ij.1.2.val →
      collisionMap ij.1.2 ij.1.1 (ne_of_lt ij.2) (p t) = 0}
\end{leancode}
{\raggedright\noindent File \leanpath{VOA/\allowbreak Commutative/\allowbreak DeligneAuxiliaryCollision.lean}.\par}
\begin{leancode}
abbrev Particle (m d : ℕ) := Fin (m + d)

def conductor (m d : ℕ) : MvPolynomial (Particle m d) K :=
  ∏ ij : Pair m d, factor (K := K) ij ^ order ij

def conductorDegree (m d : ℕ) : ℕ := ∑ ij : Pair m d, order ij
\end{leancode}
{\raggedright\noindent File \leanpath{VOA/\allowbreak Commutative/\allowbreak QuadraticCurrentGradedCarrier.lean}.\par}
\begin{leancode}
def degreePiece {m d : ℕ}
    (C : Submodule K (TensorPolynomial (K := K) (σ := σ) m d)) (r : ℕ) :
    Submodule K C where
  carrier := {p | ∀ t, (p.1 t).IsWeightedHomogeneous (fun _ => 1) r}
\end{leancode}
\noindent\textbf{Hypotheses / status.} Definitions\par
\medskip
\begin{btheorem}{24.17}[The restricted dual of a jet piece is a carrier piece (B4)]
For any field $K$, quadrics $Q$, and $r\le N$: the dual of the piece of the order-$N$ presented jet algebra of
$\mathrm{quadIdeal}(Q)$ with $m$ particles and $r$ derivatives is linearly isomorphic to
$\mathrm{degreePiece}(\mathrm{currentCarrier}(Q,m,0),r)$.
\end{btheorem}
{\raggedright\noindent\textbf{Lean.} \lean{VOA.}\allowbreak\lean{QuadraticJetRestrictedDual.}\allowbreak\lean{restrictedDualEquivCurrentCarrier}\par}
{\raggedright\noindent File \leanpath{VOA/\allowbreak Commutative/\allowbreak QuadraticJetRestrictedDual.lean}.\par}
\begin{leancode}
def restrictedDualEquivCurrentCarrier (Q : ρ → σ × σ → K) (N m r : ℕ) (hr : r ≤ N) :
    (quotientPiece Q N m r →ₗ[K] K) ≃ₗ[K] degreePiece (currentCarrier Q m 0) r :=
  (pieceDualEquivCoefficients Q N m r).trans (fixedQuadraticReindexingEquiv Q N m r hr)
\end{leancode}
\noindent\textbf{Hypotheses / status.} Unconditional\par
\medskip
\begin{btheorem}{24.18}[Geometric carrier embedding for every quadratic presentation]
If the ideal of a current $C_2$ presentation is $\mathrm{quadIdeal}(Q)$, then
$\mathrm{Carrier}_{p,N}:=\mathrm{degreePiece}(\mathrm{currentCarrier}(Q,N-p,0),p)$ receives an injective (for
$p\le N$ bijective) restricted-dual map in every bidegree; hence a \lean{GeometricCarrierEmbedding}.
\end{btheorem}
{\raggedright\noindent\textbf{Lean.} \lean{VOA.}\allowbreak\lean{DeligneAudit.}\allowbreak\lean{CarrierEmbedding.}\allowbreak\lean{carrierPiece},\allowbreak\ \lean{VOA.}\allowbreak\lean{DeligneAudit.}\allowbreak\lean{CarrierEmbedding.}\allowbreak\lean{Carrier},\allowbreak\ \lean{VOA.}\allowbreak\lean{DeligneAudit.}\allowbreak\lean{CarrierEmbedding.}\allowbreak\lean{geometricCarrierEmbedding_of_quadraticIdeal}\par}
{\raggedright\noindent File \leanpath{VOA/\allowbreak DeligneAudit/\allowbreak A2CarrierEmbedding.lean}.\par}
\begin{leancode}
abbrev carrierPiece (Q : ρ → σ × σ → K) (p N : ℕ) : Submodule K (currentCarrier Q (N - p) 0) :=

abbrev Carrier (Q : ρ → σ × σ → K) (p N : ℕ) : Type _ := carrierPiece Q p N

def geometricCarrierEmbedding_of_quadraticIdeal {ρ : Type}
    (P : CurrentC2Presentation case W)
    (Q : ρ → Fin case.currentDimension × Fin case.currentDimension → ℂ)
    (hQ : P.ideal = quadraticIdeal Q) :
    GeometricCarrierEmbedding case W P (Carrier Q) where
  toCarrier := (carrierMaps Q P.ideal hQ).1
  injective := (carrierMaps Q P.ideal hQ).2
\end{leancode}
\noindent\textbf{Hypotheses / status.} Unconditional (weak alone)\par
\medskip
\begin{btheorem}{24.19}[The character package is one inequality]
Given current-jet surjectivity and a geometric carrier embedding, $\dim W_N\le\sum_{p\le N}\dim\mathrm{Carrier}_{p,N}$
holds automatically; hence \lean{OrdinaryCarrierCharacter} follows from the reverse inequality. For a
quadratic presentation, classical freeness and $\ker=\mathrm{quadIdeal}(Q)$ follow from current-jet surjectivity
and the named residual $\sum_p\dim\mathrm{Carrier}(Q)_{p,N}\le\dim W_N$.
\end{btheorem}
{\raggedright\noindent\textbf{Lean.} \lean{VOA.}\allowbreak\lean{DeligneAudit.}\allowbreak\lean{CarrierEmbedding.}\allowbreak\lean{QuadraticCarrierUpperBound},\allowbreak\ \lean{VOA.}\allowbreak\lean{DeligneAudit.}\allowbreak\lean{CarrierEmbedding.}\allowbreak\lean{caseClassicallyFree_of_quadraticCarrierUpperBound}\par}
{\raggedright\noindent File \leanpath{VOA/\allowbreak DeligneAudit/\allowbreak A2CarrierUpperBound.lean}.\par}
\begin{leancode}
structure QuadraticCarrierUpperBound {ρ : Type} {n : ℕ} (Q : ρ → Fin n × Fin n → ℂ)
    (W : VertexOperatorAlgebra V) : Prop where
  le : ∀ N, (∑ p ∈ Finset.range (N + 1), Module.finrank ℂ (Carrier Q p N)) ≤
    Module.finrank ℂ (W.grading N)

theorem caseClassicallyFree_of_quadraticCarrierUpperBound {ρ : Type}
    (P : CurrentC2Presentation case W)
    (Q : ρ → Fin case.currentDimension × Fin case.currentDimension → ℂ)
    (hQ : P.ideal = quadraticIdeal Q)
    (hsurj : CurrentJetSurjectivity case W P)
    (hbound : QuadraticCarrierUpperBound Q W) :
    CaseClassicallyFree case W P ∧
      RingHom.ker
          (MvPolynomial.aeval fun i =>
            liDegreeZeroClass W 1 ⟨P.gen i, P.gen_weight_one i⟩).toRingHom =
        P.ideal :=
\end{leancode}
\noindent\textbf{Hypotheses / status.} \lean{QuadraticCarrierUpperBound}\par
\medskip
\begin{btheorem}{24.25}[Carrier side of the chain, any type]
For carrier data $\mathscr C$, a table $\mathrm{hilb}$ of Demazure Hilbert series and a family $F$: if
\lean{SectorFusionBound D F} and $\mathrm{hilbSeries}(\mathrm{hilb})_n=M_{n,0}$ (U6), then
$\mathrm{carrierDim}(N)\le\sum_{n\le N}[q^N]q^n\,\mathrm{hilb}_n/(q)_n$ for every $N$.
\end{btheorem}
{\raggedright\noindent\textbf{Lean.} \lean{VOA.}\allowbreak\lean{DeligneAudit.}\allowbreak\lean{DeligneCarrier.}\allowbreak\lean{hilbSeries},\allowbreak\ \lean{VOA.}\allowbreak\lean{DeligneAudit.}\allowbreak\lean{DeligneCarrier.}\allowbreak\lean{carrierDim_le_demazureSum}\par}
{\raggedright\noindent File \leanpath{VOA/\allowbreak DeligneAudit/\allowbreak DeligneCarrierDemazure.lean}.\par}
\begin{leancode}
def hilbSeries (hilb : ℕ → ℕ → ℕ) (n : ℕ) : PowerSeries ℚ :=
  PowerSeries.mk fun N => (hilb n N : ℚ)

theorem carrierDim_le_demazureSum (D : DeligneCarrierData) (hilb : ℕ → ℕ → ℕ)
    (F : ℕ → PowerSeries ℚ) (hF : SectorFusionBound D F)
    (hU6 : MixedFusionRecursion F (hilbSeries hilb)) (N : ℕ) :
    carrierDim D N ≤ demazureSumCoeff hilb N :=
\end{leancode}
\noindent\textbf{Hypotheses / status.} Sector bound, U6\par
\medskip
\begin{bdefinition}{24.27}[Loop operators, $M_m$, $B_m$, the discriminant]
For structure constants $f$ of $\ltwentyfourg$ in the chosen basis, the loop operator $x\otimes t^r$ acts on tensor
polynomials by the coadjoint action in each slot, weighted by $z_i^r$:
$(x\otimes t^r\cdot p)_t=-\sum_iz_i^r\sum_cf(x,t_i,c)\,p_{t[i\mapsto c]}$. Let $\delta_{c_0}^{\otimes m}$ be the
tensor polynomial equal to $1$ at the labelling $t\equiv c_0$ and $0$ elsewhere, $M_m$ the span of all loop words
applied to it, $B_m$ the symmetric polynomials and $\Delta=\prod_{i\neq j}(z_i-z_j)\in B_m$.
\end{bdefinition}
{\raggedright\noindent\textbf{Lean.} \lean{VOA.}\allowbreak\lean{DeligneAudit.}\allowbreak\lean{SectorLoop.}\allowbreak\lean{loopOp},\allowbreak\ \lean{VOA.}\allowbreak\lean{DeligneAudit.}\allowbreak\lean{SectorLoop.}\allowbreak\lean{deltaVec},\allowbreak\ \lean{VOA.}\allowbreak\lean{DeligneAudit.}\allowbreak\lean{SectorLoop.}\allowbreak\lean{Mspace},\allowbreak\ \lean{VOA.}\allowbreak\lean{DeligneAudit.}\allowbreak\lean{SectorFiber.}\allowbreak\lean{SymAlg},\allowbreak\ \lean{VOA.}\allowbreak\lean{DeligneAudit.}\allowbreak\lean{A2.}\allowbreak\lean{discS}\par}
{\raggedright\noindent File \leanpath{VOA/\allowbreak DeligneAudit/\allowbreak A2SectorLoop.lean}.\par}
\begin{leancode}
def loopOp (m : ℕ) (x : σ) (r : ℕ) :
    TensorPolynomial (K := K) (σ := σ) m 0 →ₗ[K] TensorPolynomial (K := K) (σ := σ) m 0 where
  toFun p t := -∑ i, zr i ^ r * ∑ c, C (f x (t i) c) * p (Function.update t i c)

def deltaVec (m : ℕ) (c₀ : σ) : TensorPolynomial (K := K) (σ := σ) m 0 :=
  fun t => if t = fun _ => c₀ then 1 else 0

def Mspace (m : ℕ) (c₀ : σ) : Submodule K (TensorPolynomial (K := K) (σ := σ) m 0) :=
  Submodule.span K (Set.range fun l => applyWord f m l (deltaVec m c₀))
\end{leancode}
{\raggedright\noindent File \leanpath{VOA/\allowbreak DeligneAudit/\allowbreak A2SectorFiber.lean}.\par}
\begin{leancode}
abbrev SymAlg := MvPolynomial.symmetricSubalgebra (Fin (m + 0)) K
\end{leancode}
{\raggedright\noindent File \leanpath{VOA/\allowbreak DeligneAudit/\allowbreak A2SectorB8.lean}.\par}
\begin{leancode}
def discS (m : ℕ) : SymAlg (K := ℂ) m :=
  ⟨∏ i : Fin (m + 0), ∏ j : Fin (m + 0), if i = j then 1 else (X i - X j), by
\end{leancode}
\medskip
\begin{bdefinition}{24.28}[Sector type data]
Sector data extend carrier data by: structure constants $f$ satisfying the Jacobi identity; ad-stability of
$I_2=\ltwentyfourIns\oplus\C Q$; an index $c_0$ with no relation containing $X_{c_0}^2$; the local Weyl data for
$W_{\mathrm{loc}}(m\theta)$ (raising vectors killing $\delta_{c_0}$, Cartan vectors acting by $\theta(h)$, simple
lowering vectors with exponents $\theta(h_i)$, a weight function compatible with $f$ and a combinatorial
no-weight fact, a fusion set); and a \emph{depth} $L$ such that every basis vector of $\ltwentyfourg^*$ is a
combination of coadjoint words of length $\le L$ applied to $\delta_{c_0}$.
\end{bdefinition}
{\raggedright\noindent\textbf{Lean.} \lean{VOA.}\allowbreak\lean{DeligneAudit.}\allowbreak\lean{DeligneSector.}\allowbreak\lean{slotA},\allowbreak\ \lean{VOA.}\allowbreak\lean{DeligneAudit.}\allowbreak\lean{DeligneSector.}\allowbreak\lean{slotWords},\allowbreak\ \lean{VOA.}\allowbreak\lean{DeligneAudit.}\allowbreak\lean{DeligneSector.}\allowbreak\lean{DeligneSectorData},\allowbreak\ \lean{VOA.}\allowbreak\lean{DeligneAudit.}\allowbreak\lean{SectorLoop.}\allowbreak\lean{JacobiCoord},\allowbreak\ \lean{VOA.}\allowbreak\lean{DeligneAudit.}\allowbreak\lean{SectorLoop.}\allowbreak\lean{RelationsStable}\par}
{\raggedright\noindent File \leanpath{VOA/\allowbreak DeligneAudit/\allowbreak DeligneSectorData.lean}.\par}
\begin{leancode}
def slotA {n : ℕ} (f : Fin n → Fin n → Fin n → ℂ) (x : Fin n) (v : Fin n → ℂ) : Fin n → ℂ :=
  fun y => -∑ d, f x y d * v d

def slotWords {n : ℕ} (f : Fin n → Fin n → Fin n → ℂ) (c₀ : Fin n) (L : ℕ) :
    Submodule ℂ (Fin n → ℂ) :=
  Submodule.span ℂ {v | ∃ l : List (Fin n), l.length ≤ L ∧ v = wordA f l (Pi.single c₀ 1)}

structure DeligneSectorData extends DeligneCarrierData where
  ρFin : Fintype ρ
  f : Fin n → Fin n → Fin n → ℂ
  jacobi : JacobiCoord f
  stable : RelationsStable f toDeligneCarrierData.fullMatrix
  c₀ : Fin n
  diag : ∀ r, toDeligneCarrierData.fullMatrix r (c₀, c₀) = 0
  raising : Fin n → Prop
  raising_table : ∀ e, raising e → ∀ y, f e y c₀ = 0
  cartan : Fin n → Prop
  thetaH : Fin n → ℕ
  cartan_table : ∀ h, cartan h → ∀ y, f h y c₀ = if y = c₀ then -(thetaH h : ℂ) else 0
  hθ : Fin n
  hθ_cartan : cartan hθ
  hθ_ne : thetaH hθ ≠ 0
  lowering : Fin n → Prop
  lowExp : Fin n → ℕ
  R : ℕ
  wt : Fin n → Fin R → ℤ
  table_wt : ∀ x y c, wt c ≠ wt x + wt y → f x y c = 0
  lowering_noweight : ∀ i, lowering i → ∀ (m : ℕ) (t : Fin m → Fin n),
    ∑ j, wt (t j) ≠ m • wt c₀ - (m * lowExp i + 1) • wt i
  fusion : Fin n → Prop
  depth : ℕ
  cyclic : ∀ c, (Pi.single c 1 : Fin n → ℂ) ∈ slotWords f c₀ depth
  cartanN : Fin n → Prop := fun _ => False
  thetaN : Fin n → ℕ := fun _ => 0
  cartanN_table : ∀ h, cartanN h → ∀ y, f h y c₀ = if y = c₀ then (thetaN h : ℂ) else 0 := by
    intro h hh; exact False.elim hh
\end{leancode}
{\raggedright\noindent File \leanpath{VOA/\allowbreak DeligneAudit/\allowbreak A2SectorRep.lean}.\par}
\begin{leancode}
def JacobiCoord : Prop :=
  ∀ a b x c, (∑ d, f a x d * f b d c) - (∑ d, f b x d * f a d c) = -∑ e, f a b e * f e x c
\end{leancode}
{\raggedright\noindent File \leanpath{VOA/\allowbreak DeligneAudit/\allowbreak A2SectorLoop.lean}.\par}
\begin{leancode}
def RelationsStable (Q : ρ → σ × σ → K) : Prop :=
  ∀ x r, quadraticRelation (adMatrix f x (Q r)) ∈
    Submodule.span K (Set.range fun s => quadraticRelation (Q s))
\end{leancode}
\noindent\textbf{Hypotheses / status.} Definitions; U5 named input\par
\medskip
\begin{bdefinition}{24.29}[Current modules, local Weyl relations, input U5]
A \emph{current module} is a vector space with operators $\rho(x\otimes t^r)$, $r\ge0$, satisfying
$[x\otimes t^r,y\otimes t^s]=\sum_ef(x,y,e)\,e\otimes t^{r+s}$. For a vector $w$, $\mathrm{cycDeg}(w,b)$ is the span of
loop words of total $t$-degree $b$ applied to $w$. The \emph{local Weyl and fusion relations} for $m\theta$ are:
raising vectors $\otimes t^r$ kill $w$; Cartan $h\otimes t^r$ kills $w$ for $r\ge1$ and acts by $m\theta(h)$ for
$r=0$; $(f_i\otimes1)^{m\theta(h_i)+1}w=0$; $(x\otimes t^m)w=0$ for $x$ in the fusion set. \emph{Input U5} for a
family $F$: every such module has $\dim\mathrm{cycDeg}(w,b)\le[q^b]F_m$, and
$\sum_{b<B}\dim\mathrm{cycDeg}(w,b)\le(\dim\ltwentyfourg)^m$ for every $B$.
The cyclic module generated by $w$ subject only to the local Weyl relations (the first three) is the \emph{local Weyl
module} $W_{\mathrm{loc}}(m\theta)$ of Chari and Pressley, the largest graded cyclic $\ltwentyfourg[t]$-module of highest weight
$m\theta$ that is finite-dimensional; adding the fusion relations gives the \emph{fusion module} $\mathcal F_m$. Every
module satisfying the relations, with $w$ cyclic, is a quotient of $\mathcal F_m$, so U5 is a statement about $\mathcal F_m$.
\end{bdefinition}
{\raggedright\noindent\textbf{Lean.} \lean{VOA.}\allowbreak\lean{DeligneAudit.}\allowbreak\lean{SectorWeyl.}\allowbreak\lean{CurrentModule},\allowbreak\ \lean{VOA.}\allowbreak\lean{DeligneAudit.}\allowbreak\lean{SectorWeyl.}\allowbreak\lean{cycDeg},\allowbreak\ \lean{VOA.}\allowbreak\lean{DeligneAudit.}\allowbreak\lean{DeligneSector.}\allowbreak\lean{LocalWeylRelG},\allowbreak\ \lean{VOA.}\allowbreak\lean{DeligneAudit.}\allowbreak\lean{DeligneSector.}\allowbreak\lean{FusionModuleBoundG}\par}
{\raggedright\noindent File \leanpath{VOA/\allowbreak DeligneAudit/\allowbreak A2SectorWeyl.lean}.\par}
\begin{leancode}
structure CurrentModule [Fintype σ] (V : Type*) [AddCommGroup V] [Module K V] where
  op : σ → ℕ → V →ₗ[K] V
  comm : ∀ a b r s v, op a r (op b s v) - op b s (op a r v) = ∑ e, f a b e • op e (r + s) v

def cycDeg (ρ : CurrentModule f V) (w : V) (b : ℕ) : Submodule K V :=
  Submodule.span K {v | ∃ l, SectorGraded.wordDeg l = b ∧ v = wordOp ρ l w}
\end{leancode}
{\raggedright\noindent File \leanpath{VOA/\allowbreak DeligneAudit/\allowbreak DeligneSectorWeyl.lean}.\par}
\begin{leancode}
structure LocalWeylRelG {V : Type*} [AddCommGroup V] [Module ℂ V]
    (ρ : CurrentModule S.f V) (w : V) (m : ℕ) : Prop where
  raising : ∀ e, S.raising e → ∀ r, ρ.op e r w = 0
  cartan_pos : ∀ h, S.cartan h → ∀ r, 1 ≤ r → ρ.op h r w = 0
  cartan_zero : ∀ h, S.cartan h → ρ.op h 0 w = ((m * S.thetaH h : ℕ) : ℂ) • w
  lowering : ∀ i, S.lowering i → (ρ.op i 0)^[m * S.lowExp i + 1] w = 0
  fusion : ∀ x, S.fusion x → ρ.op x m w = 0
  cartanN_pos : ∀ h, S.cartanN h → ∀ r, 1 ≤ r → ρ.op h r w = 0
  cartanN_zero : ∀ h, S.cartanN h → ρ.op h 0 w = -((m * S.thetaN h : ℕ) : ℂ) • w

structure FusionModuleBoundG (F : ℕ → PowerSeries ℚ) : Prop where
  le : ∀ (m : ℕ) (V : Type) [AddCommGroup V] [Module ℂ V] (ρ : CurrentModule S.f V) (w : V),
    LocalWeylRelG S ρ w m → ∀ b, (Module.finrank ℂ (cycDeg ρ w b) : ℚ) ≤ PowerSeries.coeff b (F m)
  total : ∀ (m : ℕ) (V : Type) [AddCommGroup V] [Module ℂ V] (ρ : CurrentModule S.f V) (w : V),
    LocalWeylRelG S ρ w m → ∀ B : ℕ,
      ∑ b ∈ Finset.range B, Module.finrank ℂ (cycDeg ρ w b) ≤ S.n ^ m
\end{leancode}
\noindent\textbf{Hypotheses / status.} Definitions; U5 named input\par
\medskip
\begin{btheorem}{24.33}[B6: freeness from U5]
For sector data $S$ and a family $F$ satisfying U5, $M_m$ is a finite free $B_m$-module for every $m$.
\end{btheorem}
{\raggedright\noindent\textbf{Lean.} \lean{VOA.}\allowbreak\lean{DeligneAudit.}\allowbreak\lean{DeligneSector.}\allowbreak\lean{MFreenessG},\allowbreak\ \lean{VOA.}\allowbreak\lean{DeligneAudit.}\allowbreak\lean{DeligneSector.}\allowbreak\lean{mFreeness}\par}
{\raggedright\noindent File \leanpath{VOA/\allowbreak DeligneAudit/\allowbreak DeligneSectorB6.lean}.\par}
\begin{leancode}
structure MFreenessG : Prop where
  free : ∀ m, Module.Free (SymAlg (K := ℂ) m) (MspaceR S m)
  finite : ∀ m, Module.Finite (SymAlg (K := ℂ) m) (MspaceR S m)

theorem mFreeness (hU5 : FusionModuleBoundG S F) : MFreenessG S := by
\end{leancode}
\noindent\textbf{Hypotheses / status.} U5\par
\medskip
\begin{btheorem}{24.41}[The divisor comparison and the sector fusion bound, all seven types]
For each $X\in\{A_2,G_2,D_4,F_4,E_6,E_7,E_8\}$ the divisor comparison holds for the type's sector data, and therefore,
for every family $F$, U5 for $F$ implies the sector fusion bound \lean{SectorFusionBound Xdata F}. For $A_2$, U5 is
taken in the $A_2$-specific form \lean{SectorWeyl.FusionModuleBound F} below, and the Lean theorem takes the divisor
comparison as a second hypothesis, discharged by the proved theorem \lean{A2.divisorComparison}.
\end{btheorem}
{\raggedright\noindent\textbf{Lean.} \lean{VOA.}\allowbreak\lean{DeligneAudit.}\allowbreak\lean{A2.}\allowbreak\lean{divisorComparison},\allowbreak\ \lean{VOA.}\allowbreak\lean{DeligneAudit.}\allowbreak\lean{DeligneSector.}\allowbreak\lean{A2_sectorFusionBound_generic},\allowbreak\ \lean{VOA.}\allowbreak\lean{DeligneAudit.}\allowbreak\lean{SectorWeyl.}\allowbreak\lean{LocalWeylRel},\allowbreak\ \lean{VOA.}\allowbreak\lean{DeligneAudit.}\allowbreak\lean{SectorWeyl.}\allowbreak\lean{FusionModuleBound},\allowbreak\ \lean{VOA.}\allowbreak\lean{DeligneAudit.}\allowbreak\lean{A2.}\allowbreak\lean{DivisorComparison},\allowbreak\ \lean{VOA.}\allowbreak\lean{DeligneAudit.}\allowbreak\lean{DeligneSector.}\allowbreak\lean{fusionModuleBoundG_of},\allowbreak\ \lean{VOA.}\allowbreak\lean{DeligneAudit.}\allowbreak\lean{DeligneSector.}\allowbreak\lean{divisorComparisonG_of},\allowbreak\ \lean{VOA.}\allowbreak\lean{DeligneAudit.}\allowbreak\lean{DeligneDivisor.}\allowbreak\lean{G2.}\allowbreak\lean{g2_pairComparison},\allowbreak\ \lean{VOA.}\allowbreak\lean{DeligneAudit.}\allowbreak\lean{DeligneDivisor.}\allowbreak\lean{G2.}\allowbreak\lean{g2_divisorComparison},\allowbreak\ \lean{VOA.}\allowbreak\lean{DeligneAudit.}\allowbreak\lean{DeligneDivisor.}\allowbreak\lean{G2.}\allowbreak\lean{G2_sectorFusionBound_of_U5},\allowbreak\ \lean{VOA.}\allowbreak\lean{DeligneAudit.}\allowbreak\lean{DeligneDivisor.}\allowbreak\lean{D4.}\allowbreak\lean{d4_pairComparison},\allowbreak\ \lean{VOA.}\allowbreak\lean{DeligneAudit.}\allowbreak\lean{DeligneDivisor.}\allowbreak\lean{D4.}\allowbreak\lean{d4_divisorComparison},\allowbreak\ \lean{VOA.}\allowbreak\lean{DeligneAudit.}\allowbreak\lean{DeligneDivisor.}\allowbreak\lean{D4.}\allowbreak\lean{D4_sectorFusionBound_of_U5},\allowbreak\ \lean{VOA.}\allowbreak\lean{DeligneAudit.}\allowbreak\lean{DeligneDivisor.}\allowbreak\lean{D4.}\allowbreak\lean{d4_carrierDim_le_demazureSum_of_U5_U6},\allowbreak\ \lean{VOA.}\allowbreak\lean{DeligneAudit.}\allowbreak\lean{DeligneDivisor.}\allowbreak\lean{F4.}\allowbreak\lean{f4_pairComparison},\allowbreak\ \lean{VOA.}\allowbreak\lean{DeligneAudit.}\allowbreak\lean{DeligneDivisor.}\allowbreak\lean{F4.}\allowbreak\lean{f4_divisorComparison},\allowbreak\ \lean{VOA.}\allowbreak\lean{DeligneAudit.}\allowbreak\lean{DeligneDivisor.}\allowbreak\lean{F4.}\allowbreak\lean{F4_sectorFusionBound_of_U5},\allowbreak\ \lean{VOA.}\allowbreak\lean{DeligneAudit.}\allowbreak\lean{DeligneDivisor.}\allowbreak\lean{E6.}\allowbreak\lean{e6_pairComparison},\allowbreak\ \lean{VOA.}\allowbreak\lean{DeligneAudit.}\allowbreak\lean{DeligneDivisor.}\allowbreak\lean{E6.}\allowbreak\lean{e6_divisorComparison},\allowbreak\ \lean{VOA.}\allowbreak\lean{DeligneAudit.}\allowbreak\lean{DeligneDivisor.}\allowbreak\lean{E6.}\allowbreak\lean{E6_sectorFusionBound_of_U5},\allowbreak\ \lean{VOA.}\allowbreak\lean{DeligneAudit.}\allowbreak\lean{DeligneDivisor.}\allowbreak\lean{E7.}\allowbreak\lean{e7_pairComparison},\allowbreak\ \lean{VOA.}\allowbreak\lean{DeligneAudit.}\allowbreak\lean{DeligneDivisor.}\allowbreak\lean{E7.}\allowbreak\lean{e7_divisorComparison},\allowbreak\ \lean{VOA.}\allowbreak\lean{DeligneAudit.}\allowbreak\lean{DeligneDivisor.}\allowbreak\lean{E7.}\allowbreak\lean{E7_sectorFusionBound_of_U5},\allowbreak\ \lean{VOA.}\allowbreak\lean{DeligneAudit.}\allowbreak\lean{DeligneDivisor.}\allowbreak\lean{E8.}\allowbreak\lean{e8_pairComparison},\allowbreak\ \lean{VOA.}\allowbreak\lean{DeligneAudit.}\allowbreak\lean{DeligneDivisor.}\allowbreak\lean{E8.}\allowbreak\lean{e8_divisorComparison},\allowbreak\ \lean{VOA.}\allowbreak\lean{DeligneAudit.}\allowbreak\lean{DeligneDivisor.}\allowbreak\lean{E8.}\allowbreak\lean{E8_sectorFusionBound_of_U5}\par}
{\raggedright\noindent File \leanpath{VOA/\allowbreak DeligneAudit/\allowbreak A2DivisorFinal.lean}.\par}
\begin{leancode}
theorem divisorComparison : DivisorComparison := divisorComparison_of_pair pairComparison
\end{leancode}
{\raggedright\noindent File \leanpath{VOA/\allowbreak DeligneAudit/\allowbreak DeligneSectorA2.lean}.\par}
\begin{leancode}
theorem A2_sectorFusionBound_generic (F : ℕ → PowerSeries ℚ) (hU5 : FusionModuleBound F)
    (hD : A2.DivisorComparison) : SectorFusionBound A2data F :=
  sectorFusionBoundG A2sector F (fusionModuleBoundG_of F hU5) (divisorComparisonG_of hD)
\end{leancode}
{\raggedright\noindent The two $A_2$-specific hypotheses. File \leanpath{VOA/\allowbreak DeligneAudit/\allowbreak A2SectorWeyl.lean}.\par}
\begin{leancode}
structure LocalWeylRel {V : Type*} [AddCommGroup V] [Module ℂ V]
    (ρ : CurrentModule (bracketTable) V) (w : V) (m : ℕ) : Prop where
  raising : ∀ (e : Fin 8), e.val < 3 → ∀ r, ρ.op e r w = 0
  cartan_pos : ∀ (h : Fin 8), 6 ≤ h.val → ∀ r, 1 ≤ r → ρ.op h r w = 0
  cartan_zero : ∀ (h : Fin 8), 6 ≤ h.val → ρ.op h 0 w = (m : ℂ) • w
  lowering_f1 : (ρ.op 3 0)^[m + 1] w = 0
  lowering_f2 : (ρ.op 4 0)^[m + 1] w = 0
  fusion : ∀ (fα : Fin 8), 3 ≤ fα.val → fα.val < 6 → ρ.op fα m w = 0

structure FusionModuleBound (F : ℕ → PowerSeries ℚ) : Prop where
  le : ∀ (m : ℕ) (V : Type) [AddCommGroup V] [Module ℂ V] (ρ : CurrentModule (bracketTable) V)
    (w : V), LocalWeylRel ρ w m → ∀ b, (Module.finrank ℂ (cycDeg ρ w b) : ℚ) ≤ coeff b (F m)
  /-- The total-dimension form of U5 (Ravinder Prop. 1): `dim W_loc(mθ)/⟨(f_α ⊗ t^m) w⟩ ≤ 8^m`,
  read on the `t`-graded pieces of any module satisfying the relations. -/
  total : ∀ (m : ℕ) (V : Type) [AddCommGroup V] [Module ℂ V] (ρ : CurrentModule (bracketTable) V)
    (w : V), LocalWeylRel ρ w m → ∀ B : ℕ,
      ∑ b ∈ Finset.range B, Module.finrank ℂ (cycDeg ρ w b) ≤ 8 ^ m
\end{leancode}
{\raggedright\noindent File \leanpath{VOA/\allowbreak DeligneAudit/\allowbreak A2SectorB8.lean}.\par}
\begin{leancode}
structure DivisorComparison : Prop where
  le : ∀ m, ∀ x ∈ currentCarrier fullMatrix m 0, ∀ p : HeightOnePrime (SymAlg (K := ℂ) m),
    discS m ∈ p.asIdeal → ∃ s : SymAlg (K := ℂ) m, s ∉ p.asIdeal ∧
      s • x ∈ Mspace bracketTable m (5 : Fin 8)
\end{leancode}
{\raggedright\noindent File \leanpath{VOA/\allowbreak DeligneAudit/\allowbreak DeligneSectorA2.lean}.\par}
\begin{leancode}
theorem fusionModuleBoundG_of (F : ℕ → PowerSeries ℚ) (hU5 : FusionModuleBound F) :
    FusionModuleBoundG A2sector F where

theorem divisorComparisonG_of (hD : A2.DivisorComparison) : DivisorComparisonG A2sector where
\end{leancode}
{\raggedright\noindent File \leanpath{VOA/\allowbreak DeligneAudit/\allowbreak G2Divisor.lean}.\par}
\begin{leancode}
theorem g2_pairComparison : PairComparisonG G2sector :=
  pairComparisonG_of_cert (S := G2sector) qI hqD 4 18 48384 (by norm_num) (by norm_num) UD hAD hUD hVD

theorem g2_divisorComparison : DivisorComparisonG G2sector :=
  divisorComparisonG_of_pair g2_pairComparison

theorem G2_sectorFusionBound_of_U5 (F : ℕ → PowerSeries ℚ) (hU5 : FusionModuleBoundG G2sector F) :
    SectorFusionBound G2data F :=
  sectorFusionBoundG G2sector F hU5 g2_divisorComparison
\end{leancode}
{\raggedright\noindent File \leanpath{VOA/\allowbreak DeligneAudit/\allowbreak DeligneDivisorD4.lean}.\par}
\begin{leancode}
theorem d4_pairComparison : PairComparisonG D4sector :=
  pairComparisonG_of_cert (S := D4sector) qI hq4 4 8 384 (by norm_num) (by norm_num) U4 hA4 hU4 hV4

theorem d4_divisorComparison : DivisorComparisonG D4sector :=
  divisorComparisonG_of_pair d4_pairComparison

theorem D4_sectorFusionBound_of_U5 (F : ℕ → PowerSeries ℚ) (hU5 : FusionModuleBoundG D4sector F) :
    SectorFusionBound D4data F :=
  D4_sectorFusionBound F hU5 d4_divisorComparison

theorem d4_carrierDim_le_demazureSum_of_U5_U6 (F : ℕ → PowerSeries ℚ)
    (hU5 : FusionModuleBoundG D4sector F)
    (hU6 : MixedFusionRecursion F (hilbSeries d4DemazureHilb)) (N : ℕ) :
    carrierDim D4data N ≤ demazureSumCoeff d4DemazureHilb N :=
  d4_carrierDim_le_demazureSum_of_sector F hU5 d4_divisorComparison hU6 N
\end{leancode}
{\raggedright\noindent File \leanpath{VOA/\allowbreak DeligneAudit/\allowbreak F4Divisor.lean}.\par}
\begin{leancode}
theorem f4_pairComparison : PairComparisonG F4sector :=
  pairComparisonG_of_cert (S := F4sector) qI hqD 4 48 559104 (by norm_num) (by norm_num) UD hAD hUD hVD

theorem f4_divisorComparison : DivisorComparisonG F4sector :=
  divisorComparisonG_of_pair f4_pairComparison

theorem F4_sectorFusionBound_of_U5 (F : ℕ → PowerSeries ℚ) (hU5 : FusionModuleBoundG F4sector F) :
    SectorFusionBound F4data F :=
  sectorFusionBoundG F4sector F hU5 f4_divisorComparison
\end{leancode}
{\raggedright\noindent File \leanpath{VOA/\allowbreak DeligneAudit/\allowbreak E6Divisor.lean}.\par}
\begin{leancode}
theorem e6_pairComparison : PairComparisonG E6sector :=
  pairComparisonG_of_cert (S := E6sector) qI hq6 4 12 216 (by norm_num) (by norm_num) U6 hA6 hU6 hV6

theorem e6_divisorComparison : DivisorComparisonG E6sector :=
  divisorComparisonG_of_pair e6_pairComparison

theorem E6_sectorFusionBound_of_U5 (F : ℕ → PowerSeries ℚ) (hU5 : FusionModuleBoundG E6sector F) :
    SectorFusionBound E6data F :=
  sectorFusionBoundG E6sector F hU5 e6_divisorComparison
\end{leancode}
{\raggedright\noindent File \leanpath{VOA/\allowbreak DeligneAudit/\allowbreak E7Divisor.lean}.\par}
\begin{leancode}
theorem e7_pairComparison : PairComparisonG E7sector :=
  pairComparisonG_of_cert (S := E7sector) qI hq7 4 24 672 (by norm_num) (by norm_num) U7 hA7 hU7 hV7

theorem e7_divisorComparison : DivisorComparisonG E7sector :=
  divisorComparisonG_of_pair e7_pairComparison

theorem E7_sectorFusionBound_of_U5 (F : ℕ → PowerSeries ℚ) (hU5 : FusionModuleBoundG E7sector F) :
    SectorFusionBound E7data F :=
  sectorFusionBoundG E7sector F hU5 e7_divisorComparison
\end{leancode}
{\raggedright\noindent File \leanpath{VOA/\allowbreak DeligneAudit/\allowbreak E8Divisor.lean}.\par}
\begin{leancode}
theorem e8_pairComparison : PairComparisonG E8sector :=
  pairComparisonG_of_cert (S := E8sector) qI hq8 4 120 96 (by norm_num) (by norm_num) U8 hA8 hU8 hV8

theorem e8_divisorComparison : DivisorComparisonG E8sector :=
  divisorComparisonG_of_pair e8_pairComparison

theorem E8_sectorFusionBound_of_U5 (F : ℕ → PowerSeries ℚ) (hU5 : FusionModuleBoundG E8sector F) :
    SectorFusionBound E8data F :=
  sectorFusionBoundG E8sector F hU5 e8_divisorComparison
\end{leancode}
\noindent\textbf{Hypotheses / status.} \textbf{unconditional}\par
\medskip
\endgroup
\setcounter{chapter}{24}
\chapter{The Deligne series II: characters, the Arakawa--Kawasetsu MLDE and \texorpdfstring{$q$}{q}-series identities}
\begingroup
\providecommand{\ltwentyfiveg}{\mathfrak g}
\providecommand{\ltwentyfivevac}{|0\rangle}
\providecommand{\ltwentyfiveD}{\mathcal D}
\providecommand{\ltwentyfivehv}{h^{\vee}}
\providecommand{\ltwentyfiveEu}{\delta_q}
\noindent{\small 17 of the chapter's 37 numbered statements: its headline results and every definition they use.}\par\medskip
\begin{btheorem}{25.2}[The lower bound suffices]
Given a presentation $P$, current-jet surjectivity and an injective geometric carrier embedding, the carrier
character lower bound implies the ordinary carrier character identity.
\end{btheorem}
{\raggedright\noindent\textbf{Lean.} \lean{VOA.}\allowbreak\lean{DeligneAudit.}\allowbreak\lean{ordinaryCarrierCharacter_of_le}\par}
{\raggedright\noindent File \leanpath{VOA/\allowbreak DeligneAudit/\allowbreak A2Character.lean}.\par}
\begin{leancode}
theorem ordinaryCarrierCharacter_of_le
    (P : CurrentC2Presentation case W)
    (Carrier : ℕ → ℕ → Type uC)
    [∀ p N, AddCommGroup (Carrier p N)]
    [∀ p N, Module ℂ (Carrier p N)]
    [∀ p N, FiniteDimensional ℂ (Carrier p N)]
    (hsurj : CurrentJetSurjectivity case W P)
    (hgeom : GeometricCarrierEmbedding case W P Carrier)
    (hlow : CarrierCharacterLowerBound W Carrier) :
    OrdinaryCarrierCharacter W Carrier
\end{leancode}
\noindent\textbf{Hypotheses / status.} Packages; unconditional\ impl.\par
\medskip
\begin{bdefinition}{25.3}[Coefficient arithmetic; the Demazure sum]
For coefficient tables $f,g:\N\to\N$ let $(f\ast g)(N)=\sum_{i=0}^Nf(i)g(N-i)$ (the Cauchy product), and let
$p_{\le n}(N)$ be the number of partitions of $N$ into parts of size at most $n$, i.e.\ $[q^N]1/(q;q)_n$. For a
table $\mathrm{hilb}(n,N)=[q^N]\mathrm{Hilb}_qD(1,n\theta)$, the Demazure sum is
\begin{equation*}
D_N=\sum_{n=0}^{N}\big(\mathrm{hilb}(n,\cdot)\ast p_{\le n}\big)(N-n),
\end{equation*}
which is the $q^N$ coefficient of (25.2).
\end{bdefinition}
{\raggedright\noindent\textbf{Lean.} \lean{VOA.}\allowbreak\lean{DeligneAudit.}\allowbreak\lean{A2.}\allowbreak\lean{conv},\allowbreak\ \lean{VOA.}\allowbreak\lean{DeligneAudit.}\allowbreak\lean{A2.}\allowbreak\lean{partitionsLe},\allowbreak\ \lean{VOA.}\allowbreak\lean{DeligneAudit.}\allowbreak\lean{DeligneCharacter.}\allowbreak\lean{demazureSumCoeff}\par}
{\raggedright\noindent Files \leanpath{VOA/\allowbreak DeligneAudit/\allowbreak A2CharacterInterface.lean} and
\leanpath{VOA/\allowbreak DeligneAudit/\allowbreak DeligneCharacter.lean}.\par}
\begin{leancode}
def conv (f g : ℕ → ℕ) (N : ℕ) : ℕ :=
  ∑ i ∈ Finset.range (N + 1), f i * g (N - i)

def partitionsLe : ℕ → ℕ → ℕ
  | 0, N => if N = 0 then 1 else 0
  | n + 1, N => ∑ j ∈ Finset.range (N / (n + 1) + 1), partitionsLe n (N - j * (n + 1))

def demazureSumCoeff (hilb : ℕ → ℕ → ℕ) (N : ℕ) : ℕ :=
  ∑ n ∈ Finset.range (N + 1), conv (hilb n) (partitionsLe n) (N - n)
\end{leancode}
\noindent\textbf{Hypotheses / status.} Definitions\par
\medskip
\begin{bdefinition}{25.4}[Named inputs of the character side, integer form]
For a table \lean{hilb}, a target sequence $t:\N\to\N$ and a VOA $W$:
\begin{itemize}
\item the \emph{star identity} $D_N=t_N$ for all $N$;
\item the \emph{target lower bound} $t_N\le\dim W_N$ for all $N$;
\item the \emph{Demazure bound} $D_N\le\dim W_N$ for all $N$.
\end{itemize}
\end{bdefinition}
{\raggedright\noindent\textbf{Lean.} \lean{VOA.}\allowbreak\lean{DeligneAudit.}\allowbreak\lean{DeligneCharacter.}\allowbreak\lean{StarIdentityFor},\allowbreak\ \lean{VOA.}\allowbreak\lean{DeligneAudit.}\allowbreak\lean{DeligneCharacter.}\allowbreak\lean{TargetLowerBound},\allowbreak\ \lean{VOA.}\allowbreak\lean{DeligneAudit.}\allowbreak\lean{DeligneCharacter.}\allowbreak\lean{DemazureBoundFor},\allowbreak\ \lean{VOA.}\allowbreak\lean{DeligneAudit.}\allowbreak\lean{DeligneCharacter.}\allowbreak\lean{demazureBoundFor_of_inputs}\par}
{\raggedright\noindent File \leanpath{VOA/\allowbreak DeligneAudit/\allowbreak DeligneCharacter.lean}.\par}
\begin{leancode}
structure StarIdentityFor (hilb : ℕ → ℕ → ℕ) (target : ℕ → ℕ) : Prop where
  eq : ∀ N, demazureSumCoeff hilb N = target N

structure TargetLowerBound (W : VertexOperatorAlgebra V) (target : ℕ → ℕ) : Prop where
  le : ∀ N : ℕ, target N ≤ Module.finrank ℂ (W.grading N)

structure DemazureBoundFor (W : VertexOperatorAlgebra V) (hilb : ℕ → ℕ → ℕ) : Prop where
  le : ∀ N : ℕ, demazureSumCoeff hilb N ≤ Module.finrank ℂ (W.grading N)

theorem demazureBoundFor_of_inputs {W : VertexOperatorAlgebra V} {hilb : ℕ → ℕ → ℕ}
    {target : ℕ → ℕ} (hstar : StarIdentityFor hilb target) (hlow : TargetLowerBound W target) :
    DemazureBoundFor W hilb
\end{leancode}
\medskip
\begin{bdefinition}{25.7}[The AK coefficients]
$\sigma_k(n)=\sum_{d\mid n}d^k$. For $h\in\N$, $c_N=\mathtt{mldeCoeff}\ h\ N\in\Q$ is defined by $c_0=1$
and (25.4) solved for $c_N$.
\end{bdefinition}
{\raggedright\noindent\textbf{Lean.} \lean{VOA.}\allowbreak\lean{DeligneAudit.}\allowbreak\lean{DeligneCharacter.}\allowbreak\lean{sigmaK},\allowbreak\ \lean{VOA.}\allowbreak\lean{DeligneAudit.}\allowbreak\lean{DeligneCharacter.}\allowbreak\lean{mldeList},\allowbreak\ \lean{VOA.}\allowbreak\lean{DeligneAudit.}\allowbreak\lean{DeligneCharacter.}\allowbreak\lean{mldeCoeff}\par}
{\raggedright\noindent File \leanpath{VOA/\allowbreak DeligneAudit/\allowbreak DeligneCharacterTypes.lean}.\par}
\begin{leancode}
def sigmaK (k n : ℕ) : ℕ := ∑ d ∈ n.divisors, d ^ k

def mldeList (h : ℕ) : ℕ → List ℚ
  | 0 => [1]
  | n + 1 =>
      let l := mldeList h n
      let N := n + 1
      l ++ [(∑ r ∈ Finset.range N,
          ((10 * ((h : ℚ) ^ 2 - 1) * sigmaK 3 (r + 1) -
            (24 * ((N : ℚ) - (r + 1)) + 2 * ((h : ℚ) + 1)) * sigmaK 1 (r + 1)) *
              l.getD (N - (r + 1)) 0)) / ((N : ℚ) * (6 * N + h))]

def mldeCoeff (h N : ℕ) : ℚ := (mldeList h N).getD N 0
\end{leancode}
\noindent\textbf{Hypotheses / status.} Definition\par
\medskip
\begin{btheorem}{25.8}[The recurrence, and uniqueness]
Call $d:\N\to K$ ($K$ a field of characteristic $0$) a \emph{normalised AK solution} with parameter $h$ if $d_0=1$
and (25.4) holds for all $N\ge1$. Then $c=\mathtt{mldeCoeff}\ h$ is one, and every normalised AK
solution equals $c$ (cast into $K$).
\end{btheorem}
{\raggedright\noindent\textbf{Lean.} \lean{VOA.}\allowbreak\lean{DeligneAudit.}\allowbreak\lean{DeligneMLDE.}\allowbreak\lean{akKernel},\allowbreak\ \lean{VOA.}\allowbreak\lean{DeligneAudit.}\allowbreak\lean{DeligneMLDE.}\allowbreak\lean{IsAKSolution},\allowbreak\ \lean{VOA.}\allowbreak\lean{DeligneAudit.}\allowbreak\lean{DeligneMLDE.}\allowbreak\lean{mldeCoeff_isAKSolution},\allowbreak\ \lean{VOA.}\allowbreak\lean{DeligneAudit.}\allowbreak\lean{DeligneMLDE.}\allowbreak\lean{eq_mldeCoeff_of_isAKSolution}\par}
{\raggedright\noindent File \leanpath{VOA/\allowbreak DeligneAudit/\allowbreak DeligneMLDE.lean}.\par}
\begin{leancode}
def akKernel (h N r : ℕ) : ℚ :=
  10 * ((h : ℚ) ^ 2 - 1) * sigmaK 3 r - (24 * ((N : ℚ) - r) + 2 * ((h : ℚ) + 1)) * sigmaK 1 r

structure IsAKSolution {K : Type*} [Field K] [CharZero K] (h : ℕ) (d : ℕ → K) : Prop where
  zero : d 0 = 1
  step : ∀ N : ℕ, 1 ≤ N →
    (N : K) * (6 * N + h) * d N
      = ∑ r ∈ Finset.range N, (akKernel h N (r + 1) : K) * d (N - (r + 1))

theorem mldeCoeff_isAKSolution (h : ℕ) : IsAKSolution (K := ℚ) h (mldeCoeff h) where

theorem eq_mldeCoeff_of_isAKSolution {K : Type*} [Field K] [CharZero K] {h : ℕ} {d : ℕ → K}
    (hd : IsAKSolution h d) : ∀ N, d N = (mldeCoeff h N : K)
\end{leancode}
\noindent\textbf{Hypotheses / status.} Unconditional\par
\medskip
\begin{bdefinition}{25.13}[Vacuum quartic relation; formal MLDE]
For a VOA $W$ and $\kappa\in\C$, the \emph{vacuum quartic relation} with constant $\kappa$ is the trace identity
$S_W(L[-2]^2\ltwentyfivevac)=\kappa\,G_4\star f$ on the adjoint module. The \emph{formal MLDE} with constant $\kappa$ is
$\ltwentyfiveEu^2f+G_2\star(2\ltwentyfiveEu f)+G_4\star(\frac c2f)=\kappa\,G_4\star f$.
\end{bdefinition}
{\raggedright\noindent\textbf{Lean.} \lean{VOA.}\allowbreak\lean{DeligneAudit.}\allowbreak\lean{DeligneMLDE.}\allowbreak\lean{VacuumQuarticRelation},\allowbreak\ \lean{VOA.}\allowbreak\lean{DeligneAudit.}\allowbreak\lean{DeligneMLDE.}\allowbreak\lean{FormalMLDE}\par}
{\raggedright\noindent File \leanpath{VOA/\allowbreak DeligneAudit/\allowbreak DeligneMLDEBlock.lean}.\par}
\begin{leancode}
structure VacuumQuarticRelation (W : VertexOperatorAlgebra 𝒱) (κ : ℂ) : Prop where
  eq : torusOnePoint (adjointModule W) (Lsq W (-2) (Lsq W (-2) W.vacuum))
    = κ • qConv (zhuG 2) (torusPartition (adjointModule W))

def FormalMLDE (κ : ℂ) : Prop :=
  qDeriv (qDeriv (torusPartition (adjointModule W)))
      + (qConv (zhuG 1) ((2 : ℂ) • qDeriv (torusPartition (adjointModule W)))
        + qConv (zhuG 2) ((W.centralCharge / 2) • torusPartition (adjointModule W)))
    = κ • qConv (zhuG 2) (torusPartition (adjointModule W))
\end{leancode}
\noindent\textbf{Hypotheses / status.} Definitions\par
\medskip
\begin{btheorem}{25.14}[The constant is forced; the AK recurrence]
Let $W$ have central charge $c=-2(h+1)$ and $\ltwentyfivevac\ne0$. If the formal MLDE (or the vacuum
quartic relation, which implies it) holds with \emph{some} constant $\kappa$, then $\kappa=(h+1)(5h-6)$, and the graded dimensions
$N\mapsto\dim W_N$ form a normalised AK solution with parameter $h$. Hence $\dim W_N=c_N$ and
\lean{TargetLowerBoundQ W (mldeCoeff h)}.
\end{btheorem}
{\raggedright\noindent\textbf{Lean.} \lean{VOA.}\allowbreak\lean{DeligneAudit.}\allowbreak\lean{DeligneMLDE.}\allowbreak\lean{kappa_eq_of_formalMLDE},\allowbreak\ \lean{VOA.}\allowbreak\lean{DeligneAudit.}\allowbreak\lean{DeligneMLDE.}\allowbreak\lean{isAKSolution_of_vacuumQuarticRelation},\allowbreak\ \lean{VOA.}\allowbreak\lean{DeligneAudit.}\allowbreak\lean{DeligneMLDE.}\allowbreak\lean{targetLowerBoundQ_of_vacuumQuarticRelation}\par}
{\raggedright\noindent File \leanpath{VOA/\allowbreak DeligneAudit/\allowbreak DeligneMLDEBlock.lean}.\par}
\begin{leancode}
theorem kappa_eq_of_formalMLDE (h : ℕ) (hc : W.centralCharge = -2 * ((h : ℂ) + 1))
    (hvac : W.vacuum ≠ 0) {κ : ℂ} (hE : FormalMLDE W κ) :
    κ = ((h : ℂ) + 1) * (5 * h - 6)

theorem isAKSolution_of_vacuumQuarticRelation (h : ℕ)
    (hc : W.centralCharge = -2 * ((h : ℂ) + 1)) (hvac : W.vacuum ≠ 0) {κ : ℂ}
    (hrel : VacuumQuarticRelation W κ) :
    IsAKSolution h (fun N : ℕ => (Module.finrank ℂ (W.grading N) : ℂ))

theorem targetLowerBoundQ_of_vacuumQuarticRelation (W : VertexOperatorAlgebra 𝒱) (h : ℕ)
    (hc : W.centralCharge = -2 * ((h : ℂ) + 1)) (hvac : W.vacuum ≠ 0)
    {κ : ℂ} (hrel : VacuumQuarticRelation W κ) :
    DeligneCharacter.TargetLowerBoundQ W (DeligneCharacter.mldeCoeff h)
\end{leancode}
\noindent\textbf{Hypotheses / status.} $c$, $\ltwentyfivevac\ne0$, relation\par
\medskip
\begin{btheorem}{25.16}[Square-bracket current algebra]
Let $W$ be any VOA. For $a,b\in W_1$ with $a_{(1)}b=\beta\ltwentyfivevac$, and all $m,n\in\Z$,
\begin{equation*}
[a[m],b[n]]=(a_{(0)}b)[m+n]+m\beta\,\delta_{m+n,0}.
\end{equation*}
If moreover $a$ is primary ($L_1a=0$), then $[L[m],a[n]]=-n\,a[m+n]$. For $a\in W_{P+1}$,
$(L_{-1}a)[n]=-n\,a[n-1]-(P+1)\,a[n]$ (the translation rule).
\end{btheorem}
{\raggedright\noindent\textbf{Lean.} \lean{VOA.}\allowbreak\lean{ZhuCoordinate.}\allowbreak\lean{sqCurrent_comm},\allowbreak\ \lean{VOA.}\allowbreak\lean{ZhuCoordinate.}\allowbreak\lean{Lsq_sqCurrent_comm},\allowbreak\ \lean{VOA.}\allowbreak\lean{ZhuCoordinate.}\allowbreak\lean{vSqZ_L_neg_one}\par}
{\raggedright\noindent File \leanpath{VOA/\allowbreak Modularity/\allowbreak ZhuCoordinateCurrent.lean}.\par}
\begin{leancode}
theorem sqCurrent_comm {a b : 𝒱} (ha : a ∈ W.grading 1) (hb : b ∈ W.grading 1)
    {β : ℂ} (hp : W.smul a 1 b = β • W.vacuum) (m n : ℤ) (x : 𝒱) :
    vSqZ W a 1 m (vSqZ W b 1 n x) - vSqZ W b 1 n (vSqZ W a 1 m x)
      = vSqZ W (W.smul a 0 b) 1 (m + n) x + (if m + n = 0 then ((m : ℂ) * β) • x else 0)

theorem Lsq_sqCurrent_comm {a : 𝒱} (ha : a ∈ W.grading 1) (ha1 : W.L 1 a = 0) (m n : ℤ)
    (x : 𝒱) :
    Lsq W m (vSqZ W a 1 n x) - vSqZ W a 1 n (Lsq W m x) = (-(n : ℂ)) • vSqZ W a 1 (m + n) x

theorem vSqZ_L_neg_one {a : 𝒱} (P : ℕ) (ha : a ∈ W.grading ((P + 1 : ℕ) : ℤ)) (n : ℤ) :
    vSqZ W (W.L (-1) a) (P + 2) n
      = (-(n : ℂ)) • vSqZ W a (P + 1) (n - 1) - ((P : ℂ) + 1) • vSqZ W a (P + 1) n
\end{leancode}
\noindent\textbf{Hypotheses / status.} Unconditional\par
\medskip
\begin{btheorem}{25.21}[From carrier data to the target bound]
Suppose given the following data and hypotheses.
\begin{enumerate}
\item \emph{Carrier data} $CD$ (Chapter~24): a number $n$ of currents (for the intended use,
  $n=\dim\ltwentyfiveg$); quadrics $\mathrm{rel}_r\in\operatorname{Sym}^2$ spanning $I_2^{\rm ns}$; a Casimir quadric $Q$;
  and a certificate of U3, i.e.\ $X_aQ\in(\mathrm{rel}_r)_r$ for every $a$.
\item A Lie algebra record $D$ (a \lean{LieData}, Chapter~23) on the index set $\{0,\dots,n-1\}$, with dual
  Coxeter number $D.h$ in the normalisation of its invariant form.
\item A level $k\in\C$ with $k+D.h\ne0$, so that the Sugawara construction and $L_k=$\lean{simpleVOA D k hk} exist.
\item \lean{hcas}: a number $c_0\in\C$ with $\mathrm{Cas}=c_0\,Q$. Here $\mathrm{Cas}=\sum K_{ij}X_iX_j$ is the Sugawara
  quadric of $D$, and $Q$ is the certificate's Casimir.
\item \lean{hrel}: every relation $\mathrm{rel}_r$ has vanishing symbol in $L_k$ (Theorem~25.20).
\end{enumerate}
Then the vacuum quartic relation holds for $L_k$ with $\kappa=\frac54c^2+\frac{11}2c$. If moreover
\begin{enumerate}\setcounter{enumi}{5}
\item \lean{hc}: $h\in\N$ satisfies $nk/(k+D.h)=-2(h+1)$, i.e.\ the central charge is $-2(h+1)$,
\end{enumerate}
then $\dim L_k(\ltwentyfiveg)_N\ge c_N=\mathtt{mldeCoeff}\ h\ N$ for all $N$.
\end{btheorem}
{\raggedright\noindent\textbf{Lean.} \lean{VOA.}\allowbreak\lean{DeligneAudit.}\allowbreak\lean{VacuumRelation.}\allowbreak\lean{vacuumQuarticRelation_of_carrierData},\allowbreak\ \lean{VOA.}\allowbreak\lean{DeligneAudit.}\allowbreak\lean{VacuumRelation.}\allowbreak\lean{targetLowerBound_of_carrierData}\par}
{\raggedright\noindent File \leanpath{VOA/\allowbreak DeligneAudit/\allowbreak DeligneVacuumRelationCarrier.lean}.\par}
\begin{leancode}
theorem vacuumQuarticRelation_of_carrierData (CD : DeligneCarrierData) (D : LieData κ (Fin CD.n))
    (k : ℂ) (hk : k + D.h ≠ 0) (c₀ : ℂ)
    (hcas : (affineCurrentData D k hk).casPoly = MvPolynomial.C c₀ * CD.Q)
    (hrel : ∀ r, (affineCurrentData D k hk).sym (quadraticRelation (CD.rel r)) = 0) :
    VacuumQuarticRelation (simpleVOA D k hk)
      (5 / 4 * (simpleVOA D k hk).centralCharge ^ 2 + 11 / 2 * (simpleVOA D k hk).centralCharge) :=

theorem targetLowerBound_of_carrierData (CD : DeligneCarrierData) (D : LieData κ (Fin CD.n))
    (k : ℂ) (hk : k + D.h ≠ 0) (h : ℕ)
    (hc : (CD.n : ℂ) * k / (k + D.h) = -2 * ((h : ℂ) + 1)) (c₀ : ℂ)
    (hcas : (affineCurrentData D k hk).casPoly = MvPolynomial.C c₀ * CD.Q)
    (hrel : ∀ r, (affineCurrentData D k hk).sym (quadraticRelation (CD.rel r)) = 0) :
    DeligneCharacter.TargetLowerBoundQ (simpleVOA D k hk) (DeligneCharacter.mldeCoeff h)
\end{leancode}
\noindent\textbf{Hypotheses / status.} \lean{hc}, \lean{hcas}, \lean{hrel}\par
\medskip
\begin{btheorem}{25.22}[The character bound for every Deligne type]
For each row of the table, the simple affine VOA $W=L_k(\ltwentyfiveg)$ satisfies, with no hypothesis,
$\dim W_N\ge c_N=\mathtt{mldeCoeff}\ h^\vee\ N$ for all $N$ (\lean{TargetLowerBoundQ W (mldeCoeff h∨)}). The last column
says what else is stated as a Lean theorem.
\begin{center}\small
\begin{tabular}{p{0.05\textwidth}p{0.06\textwidth}p{0.04\textwidth}p{0.05\textwidth}p{0.27\textwidth}p{0.38\textwidth}}
\toprule
type & $k$ & $h^\vee$ & $\kappa$ & bound $\dim W_N\ge c_N$ & also stated\\
\midrule
$A_2$ & $-3/2$ & 3 & 36 & via \lean{finrank_eq_mldeCoeff'} & the equality $\dim W_N=c_N$; \lean{vacuumQuarticRelation} ($\kappa=36$)\\
$G_2$ & $-5/3$ & 4 & 70 & \lean{g2_targetLowerBound} & \lean{g2_vacuumQuarticRelation} (as $\frac54c^2+\frac{11}2c$); $c$ and $\ltwentyfivevac\ne0$\\
$D_4$ & $-2$ & 6 & 168 & \lean{targetLowerBoundQ4} & \lean{vacuumQuarticRelation4} ($\kappa=168$)\\
$F_4$ & $-5/2$ & 9 & 390 & \lean{f4_targetLowerBound} & \lean{f4_vacuumQuarticRelation} (as $\frac54c^2+\frac{11}2c$); $c$ and $\ltwentyfivevac\ne0$\\
$E_6$ & $-3$ & 12 & 702 & \lean{e6_targetLowerBound} & only the inputs: \lean{e6_hc}, \lean{e6_casimir_match}, \lean{e6_hrel}\\
$E_7$ & $-4$ & 18 & 1596 & \lean{e7_targetLowerBound} & only the inputs: \lean{e7_hc}, \lean{e7_casimir_match}, \lean{e7_hrel}\\
$E_8$ & $-6$ & 30 & 4464 & \lean{targetLowerBoundQ8} & \lean{vacuumQuarticRelation8} ($\kappa=4464$)\\
\bottomrule
\end{tabular}
\end{center}
In every row $\kappa=(h^\vee+1)(5h^\vee-6)$, as Theorem~25.14 forces.
\end{btheorem}
{\raggedright\noindent\textbf{Lean.} \lean{VOA.}\allowbreak\lean{DeligneAudit.}\allowbreak\lean{A2.}\allowbreak\lean{vacuumQuarticRelation},\allowbreak\ \lean{VOA.}\allowbreak\lean{DeligneAudit.}\allowbreak\lean{A2.}\allowbreak\lean{finrank_eq_mldeCoeff'},\allowbreak\ \lean{VOA.}\allowbreak\lean{DeligneAudit.}\allowbreak\lean{D4Lie.}\allowbreak\lean{vacuumQuarticRelation4},\allowbreak\ \lean{VOA.}\allowbreak\lean{DeligneAudit.}\allowbreak\lean{D4Lie.}\allowbreak\lean{targetLowerBoundQ4},\allowbreak\ \lean{VOA.}\allowbreak\lean{DeligneAudit.}\allowbreak\lean{DeligneCarrier.}\allowbreak\lean{E6.}\allowbreak\lean{e6_hc},\allowbreak\ \lean{VOA.}\allowbreak\lean{DeligneAudit.}\allowbreak\lean{DeligneCarrier.}\allowbreak\lean{E6.}\allowbreak\lean{e6_casimir_match},\allowbreak\ \lean{VOA.}\allowbreak\lean{DeligneAudit.}\allowbreak\lean{DeligneCarrier.}\allowbreak\lean{E6.}\allowbreak\lean{e6_hrel},\allowbreak\ \lean{VOA.}\allowbreak\lean{DeligneAudit.}\allowbreak\lean{DeligneCarrier.}\allowbreak\lean{E6.}\allowbreak\lean{e6_targetLowerBound},\allowbreak\ \lean{VOA.}\allowbreak\lean{DeligneAudit.}\allowbreak\lean{E7Lie.}\allowbreak\lean{e7_hc},\allowbreak\ \lean{VOA.}\allowbreak\lean{DeligneAudit.}\allowbreak\lean{E7Lie.}\allowbreak\lean{e7_casimir_match},\allowbreak\ \lean{VOA.}\allowbreak\lean{DeligneAudit.}\allowbreak\lean{E7Lie.}\allowbreak\lean{e7_hrel},\allowbreak\ \lean{VOA.}\allowbreak\lean{DeligneAudit.}\allowbreak\lean{E7Lie.}\allowbreak\lean{e7_targetLowerBound},\allowbreak\ \lean{VOA.}\allowbreak\lean{DeligneAudit.}\allowbreak\lean{E8Lie.}\allowbreak\lean{targetLowerBoundQ8},\allowbreak\ \lean{VOA.}\allowbreak\lean{DeligneAudit.}\allowbreak\lean{E8Lie.}\allowbreak\lean{vacuumQuarticRelation8},\allowbreak\ \lean{VOA.}\allowbreak\lean{DeligneAudit.}\allowbreak\lean{G2Lie.}\allowbreak\lean{g2_vacuumQuarticRelation},\allowbreak\ \lean{VOA.}\allowbreak\lean{DeligneAudit.}\allowbreak\lean{G2Lie.}\allowbreak\lean{g2W_centralCharge},\allowbreak\ \lean{VOA.}\allowbreak\lean{DeligneAudit.}\allowbreak\lean{G2Lie.}\allowbreak\lean{g2W_vacuum_ne_zero},\allowbreak\ \lean{VOA.}\allowbreak\lean{DeligneAudit.}\allowbreak\lean{G2Lie.}\allowbreak\lean{g2_targetLowerBound},\allowbreak\ \lean{VOA.}\allowbreak\lean{DeligneAudit.}\allowbreak\lean{F4Lie.}\allowbreak\lean{f4_vacuumQuarticRelation},\allowbreak\ \lean{VOA.}\allowbreak\lean{DeligneAudit.}\allowbreak\lean{F4Lie.}\allowbreak\lean{f4W_centralCharge},\allowbreak\ \lean{VOA.}\allowbreak\lean{DeligneAudit.}\allowbreak\lean{F4Lie.}\allowbreak\lean{f4W_vacuum_ne_zero},\allowbreak\ \lean{VOA.}\allowbreak\lean{DeligneAudit.}\allowbreak\lean{F4Lie.}\allowbreak\lean{f4_targetLowerBound}\par}
{\raggedright\noindent Files \leanpath{VOA/\allowbreak DeligneAudit/\allowbreak DeligneVacuumRelationA2.lean},
\leanpath{DeligneVacuumRelationD4.lean}, \leanpath{E6Vacuum.lean}, \leanpath{E7LieVacuum.lean}, \leanpath{E7LieZhu.lean},
\leanpath{DeligneVacuumRelationE8.lean}, \leanpath{G2LieVacuum.lean}, \leanpath{F4LieVacuum.lean} (all in
\leanpath{VOA/\allowbreak DeligneAudit/\allowbreak }). Proof bodies are omitted: each block ends at \lean{:=} or \lean{:= by}. The
line above each block names the type and its \lean{W}.\par}
\begin{leancode}
theorem vacuumQuarticRelation : VacuumQuarticRelation W 36 := by

theorem finrank_eq_mldeCoeff' (N : ℕ) :
    (Module.finrank ℂ (W.grading N) : ℚ) = DeligneCharacter.mldeCoeff 3 N :=
\end{leancode}
\begin{leancode}
theorem vacuumQuarticRelation4 : VacuumQuarticRelation W4 168 := by

theorem targetLowerBoundQ4 : DeligneCharacter.TargetLowerBoundQ W4 (DeligneCharacter.mldeCoeff 6) :=
\end{leancode}
\begin{leancode}
theorem e6_hc : ((E6data.n : ℕ) : ℂ) * k6 / (k6 + e6LieData.h) = -2 * (((12 : ℕ) : ℂ) + 1) := by

theorem e6_casimir_match : E6cd.casPoly = MvPolynomial.C (2 / 3 : ℂ) * E6data.Q :=

theorem e6_hrel (r : Fin 650) : E6cd.sym (quadraticRelation (E6data.rel r)) = 0 := by

theorem e6_targetLowerBound : DeligneCharacter.TargetLowerBoundQ W6 (DeligneCharacter.mldeCoeff 12) :=
\end{leancode}
\begin{leancode}
theorem e7_hc : ((DeligneCarrier.E7.E7data.n : ℕ) : ℂ) * (-4) / (-4 + e7CertData.h) =
    -2 * (((18 : ℕ) : ℂ) + 1) := by

theorem e7_casimir_match :
    (VacuumRelation.affineCurrentData (ι := Fin DeligneCarrier.E7.E7data.n) e7CertData (-4)
      hk7).casPoly = MvPolynomial.C (1 / 2 : ℂ) * DeligneCarrier.E7.E7data.Q :=

theorem e7_hrel (r' : Fin 1539) :
    (affineCurrentData (ι := Fin DeligneCarrier.E7.E7data.n) e7CertData (-4) hk7).sym
      (quadraticRelation (DeligneCarrier.E7.E7data.rel r')) = 0 :=

theorem e7_targetLowerBound : DeligneCharacter.TargetLowerBoundQ W7 (DeligneCharacter.mldeCoeff 18) :=
\end{leancode}
\begin{leancode}
theorem targetLowerBoundQ8 :
    DeligneCharacter.TargetLowerBoundQ W8 (DeligneCharacter.mldeCoeff 30) :=

theorem vacuumQuarticRelation8 : VacuumQuarticRelation W8 4464 := by
\end{leancode}
\begin{leancode}
theorem g2_vacuumQuarticRelation :
    DeligneMLDE.VacuumQuarticRelation g2W
      (5 / 4 * g2W.centralCharge ^ 2 + 11 / 2 * g2W.centralCharge) :=

theorem g2W_centralCharge : g2W.centralCharge = -2 * (((4 : ℕ) : ℂ) + 1) := by

theorem g2W_vacuum_ne_zero : g2W.vacuum ≠ 0 := simple_vacuum_ne_zero' _ _ _

theorem g2_targetLowerBound :
    DeligneCharacter.TargetLowerBoundQ g2W (DeligneCharacter.mldeCoeff 4) :=
\end{leancode}
\begin{leancode}
theorem f4_vacuumQuarticRelation :
    DeligneMLDE.VacuumQuarticRelation f4W
      (5 / 4 * f4W.centralCharge ^ 2 + 11 / 2 * f4W.centralCharge) :=

theorem f4W_centralCharge : f4W.centralCharge = -2 * (((9 : ℕ) : ℂ) + 1) := by

theorem f4W_vacuum_ne_zero : f4W.vacuum ≠ 0 := simple_vacuum_ne_zero' _ _ _

theorem f4_targetLowerBound :
    DeligneCharacter.TargetLowerBoundQ f4W (DeligneCharacter.mldeCoeff 9) :=
\end{leancode}
\noindent\textbf{Hypotheses / status.} \textbf{unconditional}\par
\medskip
\begin{btheorem}{25.26}[Euler's pentagonal theorem; the sum side of the quintuple product]
$(q;q)_\infty=\sum_{r\in\Z}(-1)^rq^{r(3r-1)/2}$. With $J_1(u)=\sum(-1)^nu^np^{n(n-1)/2}$ and $J_2(u)=\sum u^mp^{m^2}$,
$J_1(u)J_2(u)=(p;p)_\infty\sum_j\big(u^{3j}p^{j(3j-1)}-u^{3j+1}p^{j(3j+1)}\big)$ (the sum side of the quintuple
product identity).
\end{btheorem}
{\raggedright\noindent\textbf{Lean.} \lean{VOA.}\allowbreak\lean{QSeries.}\allowbreak\lean{JTP.}\allowbreak\lean{pentSeries},\allowbreak\ \lean{VOA.}\allowbreak\lean{QSeries.}\allowbreak\lean{JTP.}\allowbreak\lean{euler_eq_pentagonal},\allowbreak\ \lean{VOA.}\allowbreak\lean{QSeries.}\allowbreak\lean{JTP.}\allowbreak\lean{quintuple_sum}\par}
{\raggedright\noindent Files \leanpath{VOA/\allowbreak Commutative/\allowbreak JacobiTriplePentagonal.lean},
\leanpath{VOA/\allowbreak Commutative/\allowbreak JacobiTripleQuintuple.lean}.\par}
\begin{leancode}
def pentSeries : PowerSeries R := zSeries (fun r => negSign r) pentagonal

theorem euler_eq_pentagonal : (euler : PowerSeries R) = pentSeries := by

theorem quintuple_sum (u : Rˣ) :
    fsr (fun n : ℤ => negSign n * ((u ^ n : Rˣ) : R)) pentExp *
        fsr (fun m : ℤ => ((u ^ m : Rˣ) : R)) sqExp =
      euler * (fsr (fun j : ℤ => ((u ^ (3 * j) : Rˣ) : R)) (fun j => 2 * pentagonal j) -
        fsr (fun j : ℤ => ((u ^ (3 * j + 1) : Rˣ) : R)) (fun j => 2 * pentagonal (-j)))
\end{leancode}
\medskip
\begin{btheorem}{25.27}[Legendre's four-triangular-numbers theorem]
In $\Z[[q]]$: $\displaystyle\sum_{v\ge0}\frac{(2v+1)q^v}{1-q^{2v+1}}=(q;q)_\infty^4\prod_{k\ge1}(1+q^k)^8$.
\end{btheorem}
{\raggedright\noindent\textbf{Lean.} \lean{VOA.}\allowbreak\lean{QSeries.}\allowbreak\lean{JTP.}\allowbreak\lean{lambertS},\allowbreak\ \lean{VOA.}\allowbreak\lean{QSeries.}\allowbreak\lean{JTP.}\allowbreak\lean{lambertS_eq}\par}
{\raggedright\noindent File \leanpath{VOA/\allowbreak Commutative/\allowbreak FourTriangular.lean}.\par}
\begin{leancode}
def lambertS : PowerSeries R :=
  powerShiftSum 1 fun v => ((2 * v + 1 : ℕ) : PowerSeries R) * geomInvSeries (2 * v + 1)

theorem lambertS_eq :
    (lambertS : PowerSeries ℤ) = euler ^ 4 * infProd fun k => (1 + X ^ (k + 1)) ^ 8
\end{leancode}
\noindent\textbf{Hypotheses / status.} Unconditional\par
\medskip
\begin{btheorem}{25.28}[Ramanujan's three identities]
With $E_2=1-24\sum\sigma_1q^n$, $E_4=1+240\sum\sigma_3q^n$, $E_6=1-504\sum\sigma_5q^n$ in $\Q[[q]]$:
\begin{equation*}
12\ltwentyfiveEu E_2=E_2^2-E_4,\qquad 3\ltwentyfiveEu E_4=E_2E_4-E_6,\qquad 2\ltwentyfiveEu E_6=E_2E_6-E_4^2 .
\end{equation*}
\end{btheorem}
{\raggedright\noindent\textbf{Lean.} \lean{VOA.}\allowbreak\lean{DeligneAudit.}\allowbreak\lean{DeligneCharacter.}\allowbreak\lean{D4MLDE.}\allowbreak\lean{sigS},\allowbreak\ \lean{VOA.}\allowbreak\lean{DeligneAudit.}\allowbreak\lean{DeligneCharacter.}\allowbreak\lean{D4MLDE.}\allowbreak\lean{E2s},\allowbreak\ \lean{VOA.}\allowbreak\lean{DeligneAudit.}\allowbreak\lean{DeligneCharacter.}\allowbreak\lean{D4MLDE.}\allowbreak\lean{E4s},\allowbreak\ \lean{VOA.}\allowbreak\lean{DeligneAudit.}\allowbreak\lean{DeligneCharacter.}\allowbreak\lean{D4MLDE.}\allowbreak\lean{E6s},\allowbreak\ \lean{VOA.}\allowbreak\lean{DeligneAudit.}\allowbreak\lean{DeligneCharacter.}\allowbreak\lean{D4MLDE.}\allowbreak\lean{RamanujanIdentities},\allowbreak\ \lean{VOA.}\allowbreak\lean{DeligneAudit.}\allowbreak\lean{DeligneCharacter.}\allowbreak\lean{D4MLDE.}\allowbreak\lean{ramanujanIdentities},\allowbreak\ \lean{VOA.}\allowbreak\lean{DeligneAudit.}\allowbreak\lean{DeligneMLDE.}\allowbreak\lean{Heat.}\allowbreak\lean{ramanujan},\allowbreak\ \lean{VOA.}\allowbreak\lean{DeligneAudit.}\allowbreak\lean{DeligneMLDE.}\allowbreak\lean{EightSq.}\allowbreak\lean{eight_sq}\par}
{\raggedright\noindent Files \leanpath{VOA/\allowbreak DeligneAudit/\allowbreak DeligneCharacterD4MLDE.lean},
\leanpath{DeligneMLDERamanujan.lean}, \leanpath{DeligneMLDERamanujanFull.lean},
\leanpath{DeligneMLDERamanujanAdapters.lean}, \leanpath{DeligneMLDEEightSq.lean}.\par}
\begin{leancode}
def sigS (k : ℕ) : PowerSeries ℚ := PowerSeries.mk fun n => if n = 0 then 0 else (sigmaK k n : ℚ)

def E2s : PowerSeries ℚ := 1 - 24 * sigS 1
def E4s : PowerSeries ℚ := 1 + 240 * sigS 3
def E6s : PowerSeries ℚ := 1 - 504 * sigS 5

structure RamanujanIdentities : Prop where
  e2 : 12 * eulerOp E2s = E2s ^ 2 - E4s
  e4 : 3 * eulerOp E4s = E2s * E4s - E6s
  e6 : 2 * eulerOp E6s = E2s * E6s - E4s ^ 2

theorem ramanujanIdentities : RamanujanIdentities := by
\end{leancode}
\begin{leancode}
theorem ramanujan : C 12 * (P1 * P1) = C 5 * P3 + P1 - C 6 * eulerOp P1 := by

theorem eight_sq (n : ℕ) :
    ∑ p ∈ Finset.antidiagonal n, e4 p.1 * e4 p.2 = e8 n := by
\end{leancode}
\noindent\textbf{Hypotheses / status.} Unconditional (Mathlib Sturm)\par
\medskip
\begin{btheorem}{25.30}[Lorenz's formula; Borwein's cubic identity; the $G_2$ target]
Let $a=\theta_{A_2}=\sum_{x,y\in\Z}q^{x^2-xy+y^2}$, $b^3=(q;q)^9_\infty/(q^3;q^3)^3_\infty$, $c^3=27q(q^3;q^3)^9_\infty/(q;q)^3_\infty$.
Then $a=1+6\sum_{n\ge1}\big(\sum_{d\mid n}\chi_{-3}(d)\big)q^n$ (Lorenz), $a^3=b^3+c^3$ (Borwein), $2a^2=3E_2(q^3)-E_2$,
$10a^4=E_4+9E_4(q^3)$. Consequently the Kac--Wakimoto form $T=\theta_{A_2}(q)(q^3;q^3)^6_\infty/(q;q)^8_\infty$ has
$[q^N]T=\mathtt{mldeCoeff}\ 4\ N$.
\end{btheorem}
{\raggedright\noindent\textbf{Lean.} \lean{VOA.}\allowbreak\lean{DeligneAudit.}\allowbreak\lean{G2Star.}\allowbreak\lean{lambertA2},\allowbreak\ \lean{VOA.}\allowbreak\lean{DeligneAudit.}\allowbreak\lean{G2Star.}\allowbreak\lean{LorenzA2},\allowbreak\ \lean{VOA.}\allowbreak\lean{DeligneAudit.}\allowbreak\lean{G2Star.}\allowbreak\lean{lorenzA2},\allowbreak\ \lean{VOA.}\allowbreak\lean{DeligneAudit.}\allowbreak\lean{G2Star.}\allowbreak\lean{G2Level3Core},\allowbreak\ \lean{VOA.}\allowbreak\lean{DeligneAudit.}\allowbreak\lean{G2Star.}\allowbreak\lean{g2Level3Core},\allowbreak\ \lean{VOA.}\allowbreak\lean{DeligneAudit.}\allowbreak\lean{G2Star.}\allowbreak\lean{g2KW},\allowbreak\ \lean{VOA.}\allowbreak\lean{DeligneAudit.}\allowbreak\lean{G2Star.}\allowbreak\lean{G2KWEqMLDE},\allowbreak\ \lean{VOA.}\allowbreak\lean{DeligneAudit.}\allowbreak\lean{G2Star.}\allowbreak\lean{g2KWEqMLDE}\par}
{\raggedright\noindent Files \leanpath{VOA/\allowbreak DeligneAudit/\allowbreak G2StarRiccati.lean},
\leanpath{G2StarMLDE.lean}, \leanpath{G2StarLorenz.lean}, \leanpath{G2StarReduction.lean},
\leanpath{G2StarUnconditional.lean}.\par}
\begin{leancode}
def lambertA2 : PowerSeries ℚ :=
  PowerSeries.mk fun n => if n = 0 then 1 else 6 * ∑ d ∈ n.divisors, chi3 d

def LorenzA2 : Prop := AQ = lambertA2

theorem lorenzA2 : LorenzA2 := aq_eq_one_add_six_L0.trans lambertA2_eq_L0.symm

structure G2Level3Core : Prop where
  l1 : AQ ^ 3 = BQ + CQ
  l2 : 2 * AQ ^ 2 = 3 * E2q3 - E2q
  l3 : 10 * AQ ^ 4 = E4q + 9 * E4q3

theorem g2Level3Core : G2Level3Core := g2Level3Core_of_eisLambert g2Level3EisLambert

def g2KW : PowerSeries ℤ := thetaA2 * euler3 ^ 6 * reciprocalEuler ^ 8

structure G2KWEqMLDE : Prop where
  eq : ∀ N, ((coeff N g2KW : ℤ) : ℚ) = mldeCoeff 4 N

theorem g2KWEqMLDE : G2KWEqMLDE := g2KWEqMLDE_of_eisLambert g2Level3EisLambert
\end{leancode}
\noindent\textbf{Hypotheses / status.} Unconditional\par
\medskip
\begin{btheorem}{25.32}[$D_4$, $E_6$, $E_7$, $E_8$ closed forms]
With $H_{E_6}=7E_6^2+5E_4^3-12E_2E_4E_6$, $H_{E_7}=211E_2E_4E_6^2+77E_2E_4^4-91E_6^3-197E_4^3E_6$,
$H_{E_8}=318910E_2E_4^4E_6^2+145105E_2E_4E_6^4+33649E_2E_4^7-43225E_6^5-139129E_4^6E_6-315310E_4^3E_6^3$:
\begin{align*}
\mathtt{mldeCoeff}\ 12\ N&=[q^{N+2}]\,H_{E_6}(q)_\infty^{-22}/3991680,\\
\mathtt{mldeCoeff}\ 18\ N&=[q^{N+3}]\,H_{E_7}(q)_\infty^{-34}/42343741440,\\
\mathtt{mldeCoeff}\ 30\ N&=[q^{N+5}]\,H_{E_8}(q)_\infty^{-58}/13909248338175590400,\\
\mathtt{mldeCoeff}\ 6\ N&=[q^{N+1}]\,\tfrac1{240}\big(\ltwentyfiveEu E_4\big)(q)_\infty^{-10}.
\end{align*}
The last is the coefficientwise form of $[q^N]\,q^{-7/12}\,\ltwentyfiveEu E_4/(240\,\eta^{10})$: since
$\eta^{10}=q^{10/24}(q)_\infty^{10}$, the prefactor is $q^{-7/12-5/12}=q^{-1}$. Check: $\frac1{240}\ltwentyfiveEu E_4=\sum
n\sigma_3(n)q^n=q+18q^2+\dots$, so $c_0=1$ and $c_1=18+10=28$ (Example~25.11).
\end{btheorem}
{\raggedright\noindent\textbf{Lean.} \lean{VOA.}\allowbreak\lean{DeligneAudit.}\allowbreak\lean{DeligneCharacter.}\allowbreak\lean{E6MLDE.}\allowbreak\lean{H6},\allowbreak\ \lean{VOA.}\allowbreak\lean{DeligneAudit.}\allowbreak\lean{DeligneCharacter.}\allowbreak\lean{E6MLDE.}\allowbreak\lean{e6Target},\allowbreak\ \lean{VOA.}\allowbreak\lean{DeligneAudit.}\allowbreak\lean{DeligneCharacter.}\allowbreak\lean{E6MLDE.}\allowbreak\lean{e6Target_eq_mldeCoeff},\allowbreak\ \lean{VOA.}\allowbreak\lean{DeligneAudit.}\allowbreak\lean{DeligneCharacter.}\allowbreak\lean{MLDEClosed.}\allowbreak\lean{e7Target_eq_mldeCoeff},\allowbreak\ \lean{VOA.}\allowbreak\lean{DeligneAudit.}\allowbreak\lean{DeligneCharacter.}\allowbreak\lean{MLDEClosed.}\allowbreak\lean{e8Target_eq_mldeCoeff},\allowbreak\ \lean{VOA.}\allowbreak\lean{DeligneAudit.}\allowbreak\lean{DeligneCharacter.}\allowbreak\lean{D4MLDE.}\allowbreak\lean{d4E4Target},\allowbreak\ \lean{VOA.}\allowbreak\lean{DeligneAudit.}\allowbreak\lean{DeligneCharacter.}\allowbreak\lean{D4MLDE.}\allowbreak\lean{d4E4Target_eq_mldeCoeff}\par}
{\raggedright\noindent Files \leanpath{VOA/\allowbreak DeligneAudit/\allowbreak DeligneCharacterE6MLDE.lean},
\leanpath{DeligneCharacterE78MLDE.lean}, \leanpath{DeligneCharacterD4MLDE.lean}.\par}
\begin{leancode}
def H6 : PowerSeries ℚ := 7 * E6s ^ 2 + 5 * E4s ^ 3 - 12 * E2s * E4s * E6s

def e6Target (N : ℕ) : ℚ := coeff N g6 / 3991680

theorem e6Target_eq_mldeCoeff (N : ℕ) : e6Target N = mldeCoeff 12 N := by
\end{leancode}
\begin{leancode}
theorem e7Target_eq_mldeCoeff (N : ℕ) :
    coeff (N + 3) (HE7 * Qk 34) / 42343741440 = mldeCoeff 18 N := by

theorem e8Target_eq_mldeCoeff (N : ℕ) :
    coeff (N + 5) (HE8 * Qk 58) / 13909248338175590400 = mldeCoeff 30 N := by
\end{leancode}
\begin{leancode}
def d4E4Target (N : ℕ) : ℚ := coeff N gser / 240

theorem d4E4Target_eq_mldeCoeff (hR : RamanujanIdentities) (N : ℕ) :
    d4E4Target N = mldeCoeff 6 N := by
\end{leancode}
\noindent\textbf{Hypotheses / status.} Unconditional\par
\medskip
\begin{btheorem}{25.36}[The $A_2$ star identity]
With $\mathrm{Hilb}_qD(1,n\theta)=P_{(2n,n,0)}(1,1,1;q,0)$ (the $q$-Whittaker polynomial, by Gelfand--Tsetlin
branching at $t=0$),
\begin{equation*}
\sum_{n\ge0}\frac{q^nP_{(2n,n,0)}(1,1,1;q,0)}{(q;q)_n}=\prod_{r\ge1}(1+q^r)^8 .
\end{equation*}
\end{btheorem}
{\raggedright\noindent\textbf{Lean.} \lean{VOA.}\allowbreak\lean{DeligneAudit.}\allowbreak\lean{A2.}\allowbreak\lean{demazureHilb},\allowbreak\ \lean{VOA.}\allowbreak\lean{DeligneAudit.}\allowbreak\lean{A2.}\allowbreak\lean{StarIdentity},\allowbreak\ \lean{VOA.}\allowbreak\lean{DeligneAudit.}\allowbreak\lean{A2.}\allowbreak\lean{StarReduction},\allowbreak\ \lean{VOA.}\allowbreak\lean{DeligneAudit.}\allowbreak\lean{A2.}\allowbreak\lean{ShiftedMehler},\allowbreak\ \lean{VOA.}\allowbreak\lean{DeligneAudit.}\allowbreak\lean{A2.}\allowbreak\lean{FourTriangular},\allowbreak\ \lean{VOA.}\allowbreak\lean{DeligneAudit.}\allowbreak\lean{A2.}\allowbreak\lean{starIdentity}\par}
{\raggedright\noindent Files \leanpath{VOA/\allowbreak DeligneAudit/\allowbreak A2CharacterInterface.lean},
\leanpath{A2CharacterStarChain.lean}, \leanpath{A2Headline.lean}.\par}
\begin{leancode}
def demazureHilb (n N : ℕ) : ℕ := gl3Hilb (2 * n) n 0 N

structure StarIdentity : Prop where
  eq : ∀ N, demazureCoeff N = productCoeff N

structure StarReduction : Prop where
  eq : (euler : PowerSeries ℤ) * starF = powerShiftSum 1 fun v => shiftedU (2 * v)

structure ShiftedMehler : Prop where
  eq : ∀ m : ℕ, shiftedU m * (euler : PowerSeries ℤ) ^ 3 * (1 - q ^ (m + 1)) = (m + 1 : ℕ)

structure FourTriangular : Prop where
  eq : (powerShiftSum 1 fun v => ((2 * v + 1 : ℕ) : PowerSeries ℤ) * geomInv (2 * v + 1)) =
    (euler : PowerSeries ℤ) ^ 4 * fermionInf

theorem starIdentity : StarIdentity := starIdentity_of_fourTriangular fourTriangular
\end{leancode}
\noindent\textbf{Hypotheses / status.} Unconditional ($q$-series)\par
\medskip
\begin{btheorem}{25.37}[$D_4$: the series \lean{d4SchurQ} $=$ MLDE solution; Demazure $=$ \lean{d4SchurQ} under positivity]
Work over $\Bbbk=\Q((r))$, $r=q^{1/2}$. Let
\begin{equation*}
\Theta(x)=\sum_{k\in\Z}(-1)^k\big(q^{\binom k2}+q^{\binom{k+1}2}\big)x^{2k},\qquad
M(s,x,z)=\frac{(s^2;q)_\infty(s^2x^2;q)_\infty(s^2x^{-2};q)_\infty}
{\prod_{\epsilon=\pm1}\prod_{i=1}^4\prod_{\sigma=\pm1}(s\,x^\epsilon z_i^\sigma;q)_\infty},
\end{equation*}
a power series in $s$ with coefficients Laurent polynomials in $x,z_1,\dots,z_4$ over $\Bbbk$; let
$I(s,z)=\mathrm{CT}_x\,\Theta(x)M(s,x,z)$ (constant term in $x$), and let
\begin{equation*}
\mathtt{d4SchurQ}_N=\sum_{m=0}^{N}\Big[r^{2(N-m)}\Big]\ \tfrac12\,[s^{2m}]\,I(s,z)\big|_{z=1}.
\end{equation*}
Then, unconditionally, $\mathtt{d4SchurQ}_N=[q^{N+1}]\,\frac1{240}(\ltwentyfiveEu E_4)\,(q)_\infty^{-10}=c_N$ for $h=6$; in
modular notation this is $[q^N]\,q^{-7/12}\,\ltwentyfiveEu E_4/(240\eta^{10})$ (Theorem~25.32). If moreover
\lean{D4DemazurePositivity} holds, the $D_4$ Demazure sum equals $\mathtt{d4SchurQ}_N$.
\end{btheorem}
{\raggedright\noindent\textbf{Lean.} \lean{VOA.}\allowbreak\lean{DeligneAudit.}\allowbreak\lean{DeligneCharacter.}\allowbreak\lean{D4MLDE.}\allowbreak\lean{D4SchurEqE4},\allowbreak\ \lean{VOA.}\allowbreak\lean{DeligneAudit.}\allowbreak\lean{DeligneCharacter.}\allowbreak\lean{D4Wp.}\allowbreak\lean{d4SchurEqE4},\allowbreak\ \lean{VOA.}\allowbreak\lean{DeligneAudit.}\allowbreak\lean{DeligneCharacter.}\allowbreak\lean{D4Wp.}\allowbreak\lean{d4SchurQ_eq_mldeCoeff_holds},\allowbreak\ \lean{VOA.}\allowbreak\lean{DeligneAudit.}\allowbreak\lean{DeligneCharacter.}\allowbreak\lean{D4K.}\allowbreak\lean{Kq},\allowbreak\ \lean{VOA.}\allowbreak\lean{DeligneAudit.}\allowbreak\lean{DeligneCharacter.}\allowbreak\lean{D4K.}\allowbreak\lean{R},\allowbreak\ \lean{VOA.}\allowbreak\lean{DeligneAudit.}\allowbreak\lean{DeligneCharacter.}\allowbreak\lean{D4K.}\allowbreak\lean{S},\allowbreak\ \lean{VOA.}\allowbreak\lean{DeligneAudit.}\allowbreak\lean{DeligneCharacter.}\allowbreak\lean{D4K.}\allowbreak\lean{tFactors},\allowbreak\ \lean{VOA.}\allowbreak\lean{DeligneAudit.}\allowbreak\lean{DeligneCharacter.}\allowbreak\lean{D4K.}\allowbreak\lean{hyperFactors},\allowbreak\ \lean{VOA.}\allowbreak\lean{DeligneAudit.}\allowbreak\lean{DeligneCharacter.}\allowbreak\lean{D4K.}\allowbreak\lean{kernelM},\allowbreak\ \lean{VOA.}\allowbreak\lean{DeligneAudit.}\allowbreak\lean{DeligneCharacter.}\allowbreak\lean{D4K.}\allowbreak\lean{thK},\allowbreak\ \lean{VOA.}\allowbreak\lean{DeligneAudit.}\allowbreak\lean{DeligneCharacter.}\allowbreak\lean{D4K.}\allowbreak\lean{ctTheta},\allowbreak\ \lean{VOA.}\allowbreak\lean{DeligneAudit.}\allowbreak\lean{DeligneCharacter.}\allowbreak\lean{D4K.}\allowbreak\lean{Ikernel},\allowbreak\ \lean{VOA.}\allowbreak\lean{DeligneAudit.}\allowbreak\lean{DeligneCharacter.}\allowbreak\lean{D4K.}\allowbreak\lean{ev1},\allowbreak\ \lean{VOA.}\allowbreak\lean{DeligneAudit.}\allowbreak\lean{DeligneCharacter.}\allowbreak\lean{D4K.}\allowbreak\lean{d4SchurQ},\allowbreak\ \lean{VOA.}\allowbreak\lean{DeligneAudit.}\allowbreak\lean{DeligneCharacter.}\allowbreak\lean{D4K.}\allowbreak\lean{D4DemazurePositivity},\allowbreak\ \lean{VOA.}\allowbreak\lean{DeligneAudit.}\allowbreak\lean{DeligneCharacter.}\allowbreak\lean{D4DemazureEqSchur},\allowbreak\ \lean{VOA.}\allowbreak\lean{DeligneAudit.}\allowbreak\lean{DeligneCharacter.}\allowbreak\lean{d4DemazureEqSchur_of_positivity},\allowbreak\ \lean{VOA.}\allowbreak\lean{DeligneAudit.}\allowbreak\lean{DeligneCharacter.}\allowbreak\lean{D4Wp.}\allowbreak\lean{d4DemazurePositivity_of_U5Z}\par}
{\raggedright\noindent Files \leanpath{VOA/\allowbreak DeligneAudit/\allowbreak DeligneCharacterD4MLDE.lean},
\leanpath{DeligneCharacterD4Wp12.lean}, \leanpath{DeligneCharacterD4Schur.lean}, \leanpath{DeligneCharacterD4Final.lean},
\leanpath{D4CharacterTop.lean}.\par}
\begin{leancode}
def D4SchurEqE4 : Prop := ∀ N, D4K.d4SchurQ N = d4E4Target N

theorem d4SchurEqE4 : D4MLDE.D4SchurEqE4 := by

theorem d4SchurQ_eq_mldeCoeff_holds (N : ℕ) : D4K.d4SchurQ N = mldeCoeff 6 N :=
\end{leancode}
\begin{leancode}
abbrev Kq := LaurentSeries ℚ

abbrev R := AddMonoidAlgebra Kq V

abbrev S := PowerSeries R

def tFactors : S :=
  EPos 1 0 2 two_ne_zero * EPos 1 (2 • xv) 2 two_ne_zero * EPos 1 (-2 • xv) 2 two_ne_zero

def hyperFactors : S := (hyperVecs.map fun v => EInv 1 v 1 one_ne_zero).prod

def kernelM : S := tFactors * hyperFactors

def thK (i : ℤ) : Kq := if Even i then sgn (i / 2) * (qC (i / 2) + qC (i / 2 + 1)) else 0

abbrev ctTheta : R →ₗ[Kq] R := ctFun thK

def Ikernel : S := smap ctTheta kernelM

def ev1 : R →ₗ[Kq] Kq := Finsupp.linearCombination Kq fun _ => (1 : Kq)

def d4SchurQ (N : ℕ) : ℚ :=
  ∑ m ∈ Finset.range (N + 1), ((2 : Kq)⁻¹ * ev1 (coeff (2 * m) Ikernel)).coeff (2 * ((N : ℤ) - m))
\end{leancode}
{\raggedright\noindent Files \leanpath{DeligneCharacterD4Schur.lean}, \leanpath{DeligneCharacterD4Final.lean},
\leanpath{DeligneCharacterD4Positivity.lean}.\par}
\begin{leancode}
structure D4DemazurePositivity : Prop where
  energy : ∀ m : ℕ, ∀ t ∈ demazureChar cartanD4 thetaD4 wordD4 m, 0 ≤ t.1.2 + (m : ℤ) * m
  coeff : ∀ (m : ℕ) (N : ℕ), 0 ≤ hilbInt m N

def D4DemazureEqSchur : Prop := ∀ N, (demazureSumCoeff d4DemazureHilb N : ℚ) = d4SchurQ N

theorem d4DemazureEqSchur_of_positivity (hpos : D4DemazurePositivity) : D4DemazureEqSchur :=

theorem d4DemazurePositivity_of_U5Z
    (hU5 : DeligneSector.FusionModuleBoundG DeligneSector.D4.D4sector d4FusionFZ) :
    D4K.D4DemazurePositivity :=
\end{leancode}
\noindent\textbf{Hypotheses / status.} Unconditional; \lean{D4DemazurePositivity} (\lean{coeff} only from open U5)\par
\medskip
\endgroup
\setcounter{chapter}{25}
\chapter[The Deligne series III: fusion modules and Ravinder's kernel]{The Deligne series III: fusion modules, Demazure modules and Ravinder's kernel}
\begingroup
\providecommand{\ltwentysixg}{\mathfrak g}
\providecommand{\ltwentysixgt}{\mathfrak g[t]}
\providecommand{\ltwentysixWloc}{W_{\mathrm{loc}}}
\providecommand{\ltwentysixD}{\mathcal D}
\providecommand{\ltwentysixF}{\mathcal F}
\providecommand{\ltwentysixU}{U}
\providecommand{\ltwentysixcyc}{\operatorname{cyc}}
\providecommand{\ltwentysixHilb}{\operatorname{Hilb}}
\providecommand{\ltwentysixqb}[2]{\genfrac{[}{]}{0pt}{}{#1}{#2}_q}
\noindent{\small 19 of the chapter's 35 numbered statements: its headline results and every definition they use.}\par\medskip
\begin{bdefinition}{26.2}[Evaluation modules; fusion products]
(a) For a $\ltwentysixg$-module $V$ and $z\in\C$ the \textbf{evaluation module} $V(z)$ is $V$ with $x(r)$ acting by $z^rx$
(with $0^0=1$). At $z=0$ every $x(r)$, $r\ge1$, acts by $0$, and $V(0)$ is a graded module concentrated in degree $0$.
We write $V(\lambda)$ for the irreducible $\ltwentysixg$-module of highest weight $\lambda$; $V(\theta)=\ltwentysixg$.

(b) More generally, for a graded cyclic $\ltwentysixgt$-module $V$ and $z\in\C$, the \textbf{translate} $V^{z}$ is $V$ with
$x(r)$ acting as $x\otimes(t+z)^r=\sum_s\binom rs z^{r-s}x(s)$. For $V=V(\lambda)(0)$ this is $V(\lambda)(z)$.

(c) Let $V_1,\dots,V_m$ be graded cyclic $\ltwentysixgt$-modules with cyclic vectors $v_j$, and $z_1,\dots,z_m$ distinct.
Put $T=V_1^{z_1}\otimes\dots\otimes V_m^{z_m}$ (with $x(r)$ acting by the Leibniz rule) and $v=v_1\otimes\dots\otimes v_m$.
Feigin and Loktev show that $T$ is cyclic on $v$. Its \textbf{degree filtration} is
\begin{equation*}
\Phi_{\le d}=\operatorname{span}\{x_1(r_1)\cdots x_\ell(r_\ell)v:\ r_1+\dots+r_\ell\le d\},\qquad \Phi_{<d}=\Phi_{\le d-1}.
\end{equation*}
Then $\Phi_{\le d}\subseteq\Phi_{\le d+1}$, $x(r)\Phi_{\le d}\subseteq\Phi_{\le d+r}$, and $\bigcup_d\Phi_{\le d}=T$. The
\textbf{fusion product} is the associated graded space
\begin{equation*}
V_1*\dots*V_m=\bigoplus_{d\ge0}\Phi_{\le d}/\Phi_{<d},
\end{equation*}
on which $x(r)$ acts by the map $\Phi_{\le d}/\Phi_{<d}\to\Phi_{\le d+r}/\Phi_{<d+r}$ that it induces. We write $V^{*m}$ for
$m$ factors equal to $V$.
\end{bdefinition}
\noindent\emph{(The chapter quotes no Lean declaration for this statement; see book p.~1390.)}\par
\medskip
\begin{bdefinition}{26.6}[The $A_2$ relations and two named statements]
Let $\rho$ be a current module and $w\in V$.
\begin{enumerate}
\item[(a)] $w$ satisfies the \emph{relations of $M_{a,b}$} if, with $k=a+b$:
$e_c(r)w=0$ for $c=e_1,e_2,e_\theta$ and all $r$; $h(r)w=0$ for $r\ge1$ and $h(0)w=k\,w$ for $h=h_1,h_2$;
$f_1(0)^{k+1}w=f_2(0)^{k+1}w=0$; $f_1(k)w=f_2(k)w=0$; and $f_\theta(2a+b)w=0$.
\item[(b)] \textbf{(TB)}, the \emph{total bound} \lean{FusionTotalBound}: for every $m$ and every module with the local Weyl and
fusion relations for $m\theta$, every $\ltwentysixcyc(w)_d$ is finite-dimensional and $\sum_{d<B}\dim\ltwentysixcyc(w)_d\le8^m$
for every $B$.
\item[(c)] \textbf{(KL)}, the \emph{kernel relation} \lean{RavinderKernelRel}: for $a\ge1$, if $w$ satisfies the
relations of $M_{a,b}$, then $f_\theta(2a+b-1)w$ satisfies those of $M_{a-1,b}$.
\end{enumerate}
\end{bdefinition}
{\raggedright\noindent\textbf{Lean.} \lean{VOA.}\allowbreak\lean{DeligneAudit.}\allowbreak\lean{SectorWeyl.}\allowbreak\lean{MabRel},\allowbreak\ \lean{VOA.}\allowbreak\lean{DeligneAudit.}\allowbreak\lean{SectorWeyl.}\allowbreak\lean{FusionTotalBound},\allowbreak\ \lean{VOA.}\allowbreak\lean{DeligneAudit.}\allowbreak\lean{SectorWeyl.}\allowbreak\lean{RavinderKernelRel}\par}
{\raggedright\noindent File \leanpath{VOA/\allowbreak DeligneAudit/\allowbreak DeligneFusionRel.lean}.\par}
\begin{leancode}
structure MabRel {V : Type} [AddCommGroup V] [Module ℂ V]
    (ρ : CurrentModule (bracketTable) V) (w : V) (a b : ℕ) : Prop where
  raising : ∀ (e : Fin 8), e.val < 3 → ∀ r, ρ.op e r w = 0
  cartan_pos : ∀ (h : Fin 8), 6 ≤ h.val → ∀ r, 1 ≤ r → ρ.op h r w = 0
  cartan_zero : ∀ (h : Fin 8), 6 ≤ h.val → ρ.op h 0 w = ((a + b : ℕ) : ℂ) • w
  lowering_f1 : (ρ.op 3 0)^[a + b + 1] w = 0
  lowering_f2 : (ρ.op 4 0)^[a + b + 1] w = 0
  trunc_f1 : ρ.op 3 (a + b) w = 0
  trunc_f2 : ρ.op 4 (a + b) w = 0
  trunc_theta : ρ.op 5 (2 * a + b) w = 0

def FusionTotalBound : Prop :=
  ∀ (m : ℕ) (V : Type) [AddCommGroup V] [Module ℂ V] (ρ : CurrentModule (bracketTable) V) (w : V),
    LocalWeylRel ρ w m → (∀ b, FiniteDimensional ℂ (cycDeg ρ w b)) ∧
      ∀ B : ℕ, ∑ b ∈ Finset.range B, Module.finrank ℂ (cycDeg ρ w b) ≤ 8 ^ m

def RavinderKernelRel : Prop :=
  ∀ (a b : ℕ) (V : Type) [AddCommGroup V] [Module ℂ V] (ρ : CurrentModule (bracketTable) V) (w : V),
    1 ≤ a → MabRel ρ w a b → MabRel ρ (ρ.op 5 (2 * a + b - 1) w) (a - 1) b
\end{leancode}
\noindent\textbf{Hypotheses / status.} Definition\par
\medskip
\begin{btheorem}{26.7}[Finiteness clause for $\mathfrak{sl}_3$]
If $w$ satisfies the local Weyl and fusion relations for $m\theta$, then every $\ltwentysixcyc(w)_b$ is
finite-dimensional. Consequently \lean{FusionTotalBound} is equivalent to its dimension part \lean{FusionTotalDim}: for
every module with these relations and every $B$, $\sum_{b<B}\dim\ltwentysixcyc(w)_b\le8^m$.
\end{btheorem}
{\raggedright\noindent\textbf{Lean.} \lean{VOA.}\allowbreak\lean{DeligneAudit.}\allowbreak\lean{SectorWeyl.}\allowbreak\lean{cycDeg_finiteDimensional},\allowbreak\ \lean{VOA.}\allowbreak\lean{DeligneAudit.}\allowbreak\lean{SectorWeyl.}\allowbreak\lean{FusionTotalDim},\allowbreak\ \lean{VOA.}\allowbreak\lean{DeligneAudit.}\allowbreak\lean{SectorWeyl.}\allowbreak\lean{fusionTotalBound_of_dim}\par}
{\raggedright\noindent File \leanpath{VOA/\allowbreak DeligneAudit/\allowbreak DeligneFusionFinite.lean}.\par}
\begin{leancode}
theorem cycDeg_finiteDimensional {w : V} {m : ℕ} (hW : LocalWeylRel ρ w m) (b : ℕ) :
    FiniteDimensional ℂ (cycDeg ρ w b) := by

def FusionTotalDim : Prop :=
  ∀ (m : ℕ) (V : Type) [AddCommGroup V] [Module ℂ V] (ρ : CurrentModule (bracketTable) V) (w : V),
    LocalWeylRel ρ w m → ∀ B : ℕ, ∑ b ∈ Finset.range B, Module.finrank ℂ (cycDeg ρ w b) ≤ 8 ^ m

theorem fusionTotalBound_of_dim (h : FusionTotalDim) : FusionTotalBound :=
\end{leancode}
\noindent\textbf{Hypotheses / status.} None; proved\par
\medskip
\begin{bdefinition}{26.8}[The fusion characters as the inverse of the mixed recursion]
For any sequence of power series $\mathrm{Dem}_n$ define recursively
\begin{equation*}
\mathrm{fusionOf}(\mathrm{Dem})_n=\mathrm{Dem}_n-\sum_{d=1}^{n}q^{dn}\ltwentysixqb nd\,\mathrm{fusionOf}(\mathrm{Dem})_{n-d}.
\end{equation*}
For $A_2$, $\mathrm{fusionF}=\mathrm{fusionOf}(\ltwentysixD)$ with $\ltwentysixD_n$ the Gelfand--Tsetlin polynomial
\lean{demazureHilbSeries n}.
\end{bdefinition}
{\raggedright\noindent\textbf{Lean.} \lean{VOA.}\allowbreak\lean{DeligneAudit.}\allowbreak\lean{DeligneCarrier.}\allowbreak\lean{fusionOf},\allowbreak\ \lean{VOA.}\allowbreak\lean{DeligneAudit.}\allowbreak\lean{DeligneCarrier.}\allowbreak\lean{mixedFusionRecursion_fusionOf},\allowbreak\ \lean{VOA.}\allowbreak\lean{DeligneAudit.}\allowbreak\lean{DeligneCarrier.}\allowbreak\lean{fusionOf_unique},\allowbreak\ \lean{VOA.}\allowbreak\lean{DeligneAudit.}\allowbreak\lean{A2.}\allowbreak\lean{fusionF},\allowbreak\ \lean{VOA.}\allowbreak\lean{DeligneAudit.}\allowbreak\lean{A2.}\allowbreak\lean{mixedFusionRecursion_fusionF}\par}
{\raggedright\noindent File \leanpath{VOA/\allowbreak DeligneAudit/\allowbreak D4FusionF.lean}.\par}
\begin{leancode}
def fusionOf (Dem : ℕ → PowerSeries ℚ) : ℕ → PowerSeries ℚ
  | n => Dem n -
      ∑ d : Fin n, X ^ ((d.1 + 1) * n) * qbinom (X : PowerSeries ℚ) n (d.1 + 1) *
        fusionOf Dem (n - (d.1 + 1))
decreasing_by omega

theorem mixedFusionRecursion_fusionOf (Dem : ℕ → PowerSeries ℚ) :
    MixedFusionRecursion (fusionOf Dem) Dem := by

theorem fusionOf_unique {F Dem : ℕ → PowerSeries ℚ} (h : MixedFusionRecursion F Dem) (n : ℕ) :
    F n = fusionOf Dem n := by
\end{leancode}
{\raggedright\noindent File \leanpath{VOA/\allowbreak DeligneAudit/\allowbreak DeligneFusionInverse.lean}.\par}
\begin{leancode}
def fusionF : ℕ → PowerSeries ℚ
  | n => demazureHilbSeries n -
      ∑ d : Fin n, X ^ ((d.1 + 1) * n) * qbinom (X : PowerSeries ℚ) n (d.1 + 1) * fusionF (n - (d.1 + 1))
decreasing_by omega

theorem mixedFusionRecursion_fusionF : MixedFusionRecursion fusionF demazureHilbSeries := by
\end{leancode}
\noindent\textbf{Hypotheses / status.} None; proved\par
\medskip
\begin{bdefinition}{26.11}[The squeeze inputs and U7]
\begin{enumerate}
\item[(a)] For $f,g\in\Q[[q]]$ write $f\le g$ (\emph{coefficientwise order}) if $[q^N]f\le[q^N]g$ for all $N$.
\item[(b)] A family $M_{a,b}\in\Q[[q]]$ satisfies the \emph{squeeze inputs} if (nonneg) its coefficients are $\ge0$;
(R) $M_{a+1,b}\le M_{a,b+1}+q^{2a+b+1}M_{a,b}$; (D) $\ltwentysixD_n\le M_{n,0}$; (T) $\sum_{N<B}[q^N]M_{0,m}\le8^m$ for all
$m,B$. These are the three inequalities of Section~26.2.5.
\item[(c)] It satisfies the \emph{universal bound} if every module whose generator satisfies the relations of $M_{a,b}$ has
$\dim\ltwentysixcyc(w)_N\le[q^N]M_{a,b}$.
\item[(d)] \textbf{U7 at $n$} (\lean{DemazureRealization n}): there is a current module and a vector $w$ satisfying the
relations of $M_{n,0}$ with $\dim\ltwentysixcyc(w)_N\ge[q^N]\ltwentysixD_n$ for every $N$.
\end{enumerate}
\end{bdefinition}
{\raggedright\noindent\textbf{Lean.} \lean{VOA.}\allowbreak\lean{DeligneAudit.}\allowbreak\lean{A2.}\allowbreak\lean{CoeffLe},\allowbreak\ \lean{VOA.}\allowbreak\lean{DeligneAudit.}\allowbreak\lean{A2.}\allowbreak\lean{FusionSqueezeInputs},\allowbreak\ \lean{VOA.}\allowbreak\lean{DeligneAudit.}\allowbreak\lean{A2.}\allowbreak\lean{MabUniversalBound},\allowbreak\ \lean{VOA.}\allowbreak\lean{DeligneAudit.}\allowbreak\lean{A2.}\allowbreak\lean{DemazureRealization}\par}
{\raggedright\noindent File \leanpath{VOA/\allowbreak DeligneAudit/\allowbreak DeligneFusionSqueeze.lean}.\par}
\begin{leancode}
def CoeffLe (f g : PowerSeries ℚ) : Prop := ∀ N, coeff N f ≤ coeff N g

structure FusionSqueezeInputs (M : ℕ → ℕ → PowerSeries ℚ) : Prop where
  /-- Graded characters have nonnegative coefficients. -/
  nonneg : ∀ a b N, 0 ≤ coeff N (M a b)
  /-- Spanning half of Ravinder's exact sequences. -/
  rec_le : ∀ a b, CoeffLe (M (a + 1) b) (M a (b + 1) + X ^ (2 * a + b + 1) * M a b)
  /-- `Hilb_q D(1,nθ) ≤ ch W_loc(nθ)`, with U7 (Demazure = Gelfand–Tsetlin). -/
  dem_le : ∀ n, CoeffLe (demazureHilbSeries n) (M n 0)
  /-- F4: `dim M_{0,m} ≤ 8^m`. -/
  total_le : ∀ m B, ∑ N ∈ Finset.range B, coeff N (M 0 m) ≤ 8 ^ m
\end{leancode}
{\raggedright\noindent File \leanpath{VOA/\allowbreak DeligneAudit/\allowbreak DeligneFusionDemazure.lean}.\par}
\begin{leancode}
def MabUniversalBound (M : ℕ → ℕ → PowerSeries ℚ) : Prop :=
  ∀ (a b : ℕ) (V : Type) [AddCommGroup V] [Module ℂ V] (ρ : CurrentModule (bracketTable) V)
    (w : V), MabRel ρ w a b → ∀ N, (Module.finrank ℂ (cycDeg ρ w N) : ℚ) ≤ coeff N (M a b)

def DemazureRealization (n : ℕ) : Prop :=
  ∃ (V : Type) (_ : AddCommGroup V) (_ : Module ℂ V) (ρ : CurrentModule (bracketTable) V) (w : V),
    MabRel ρ w n 0 ∧ ∀ N, demazureHilb n N ≤ Module.finrank ℂ (cycDeg ρ w N)
\end{leancode}
\noindent\textbf{Hypotheses / status.} Proved for all $n$\par
\medskip
\begin{btheorem}{26.12}[Ravinder's squeeze; U5 from the squeeze]
If $M$ satisfies the squeeze inputs, then $M_{0,m}=\mathrm{fusionF}_m$ for every $m$, and $\mathrm{fusionF}$ has nonnegative
coefficients. If moreover every module with the local Weyl and fusion relations for $m\theta$ is bounded degreewise by
$M_{0,m}$, then U5 holds for $\mathrm{fusionF}$. Finally U5 for $\mathrm{fusionF}$ follows from the universal bound, U7 for
all $n$, (nonneg), (R) and (T).
\end{btheorem}
{\raggedright\noindent\textbf{Lean.} \lean{VOA.}\allowbreak\lean{DeligneAudit.}\allowbreak\lean{A2.}\allowbreak\lean{squeeze},\allowbreak\ \lean{VOA.}\allowbreak\lean{DeligneAudit.}\allowbreak\lean{A2.}\allowbreak\lean{fusionF_coeff_nonneg},\allowbreak\ \lean{VOA.}\allowbreak\lean{DeligneAudit.}\allowbreak\lean{A2.}\allowbreak\lean{fusionModuleBound_of_squeeze},\allowbreak\ \lean{VOA.}\allowbreak\lean{DeligneAudit.}\allowbreak\lean{A2.}\allowbreak\lean{dem_le_of_realization},\allowbreak\ \lean{VOA.}\allowbreak\lean{DeligneAudit.}\allowbreak\lean{A2.}\allowbreak\lean{fusionModuleBound_of_realizations}\par}
{\raggedright\noindent File \leanpath{VOA/\allowbreak DeligneAudit/\allowbreak DeligneFusionSqueeze.lean}.\par}
\begin{leancode}
theorem squeeze (h : FusionSqueezeInputs M) (m : ℕ) : M 0 m = fusionF m := by

theorem fusionF_coeff_nonneg (h : FusionSqueezeInputs M) (m N : ℕ) :
    0 ≤ coeff N (fusionF m) := by

theorem fusionModuleBound_of_squeeze (h : FusionSqueezeInputs M)
    (huniv : ∀ (m : ℕ) (V : Type) [AddCommGroup V] [Module ℂ V]
      (ρ : CurrentModule (bracketTable) V) (w : V), LocalWeylRel ρ w m →
        ∀ b, (Module.finrank ℂ (cycDeg ρ w b) : ℚ) ≤ coeff b (M 0 m)) :
    FusionModuleBound fusionF := by
\end{leancode}
{\raggedright\noindent File \leanpath{VOA/\allowbreak DeligneAudit/\allowbreak DeligneFusionDemazure.lean}.\par}
\begin{leancode}
theorem dem_le_of_realization (hU : MabUniversalBound M) {n : ℕ} (hD : DemazureRealization n) :
    CoeffLe (demazureHilbSeries n) (M n 0) := by

theorem fusionModuleBound_of_realizations (hU : MabUniversalBound M)
    (hD : ∀ n, DemazureRealization n) (hnonneg : ∀ a b N, 0 ≤ coeff N (M a b))
    (hrec : ∀ a b, CoeffLe (M (a + 1) b) (M a (b + 1) + X ^ (2 * a + b + 1) * M a b))
    (htot : ∀ m B, ∑ N ∈ Finset.range B, coeff N (M 0 m) ≤ 8 ^ m) :
    FusionModuleBound fusionF :=
\end{leancode}
\noindent\textbf{Hypotheses / status.} Squeeze inputs; proved\par
\medskip
\begin{bdefinition}{26.13}[The universal modules $U_{a,b}$ and their characters]
Let $\mathrm{FreeW}$ be the vector space with basis the loop words (lists of letters $(c,r)\in\{0,\dots,7\}\times\N$).
Let $J_{a,b}\subseteq\mathrm{FreeW}$ be spanned by: all \emph{commutation instances}
$p\,xy\,s-p\,yx\,s-\sum_ef(x,y,e)\,p\,(e,r_x+r_y)\,s$ (for prefixes $p$, suffixes $s$); all words $p\,t$ with $t$ a
relation tail of $M_{a,b}$ (a single raising letter, a positive-grade Cartan letter, $f_i(0)^{a+b+1}$, $f_1(a+b)$,
$f_2(a+b)$, $f_\theta(2a+b)$); and all $p\,h(0)-(a+b)\,p$. Then $U_{a,b}=\mathrm{FreeW}/J_{a,b}$, with $x_c(r)$
acting by prepending the letter, and generator $w_U=[\,\text{empty word}\,]$. Its character is
$M_{a,b}^{\mathrm{univ}}(q)=\sum_N\dim\ltwentysixcyc(w_U)_N\,q^N$ (\lean{Mchar a b}).
\end{bdefinition}
{\raggedright\noindent\textbf{Lean.} \lean{VOA.}\allowbreak\lean{DeligneAudit.}\allowbreak\lean{SectorWeyl.}\allowbreak\lean{Ltr},\allowbreak\ \lean{VOA.}\allowbreak\lean{DeligneAudit.}\allowbreak\lean{SectorWeyl.}\allowbreak\lean{FreeW},\allowbreak\ \lean{VOA.}\allowbreak\lean{DeligneAudit.}\allowbreak\lean{SectorWeyl.}\allowbreak\lean{commInst},\allowbreak\ \lean{VOA.}\allowbreak\lean{DeligneAudit.}\allowbreak\lean{SectorWeyl.}\allowbreak\lean{relTail},\allowbreak\ \lean{VOA.}\allowbreak\lean{DeligneAudit.}\allowbreak\lean{SectorWeyl.}\allowbreak\lean{relGen},\allowbreak\ \lean{VOA.}\allowbreak\lean{DeligneAudit.}\allowbreak\lean{SectorWeyl.}\allowbreak\lean{Jrel},\allowbreak\ \lean{VOA.}\allowbreak\lean{DeligneAudit.}\allowbreak\lean{SectorWeyl.}\allowbreak\lean{UMod},\allowbreak\ \lean{VOA.}\allowbreak\lean{DeligneAudit.}\allowbreak\lean{SectorWeyl.}\allowbreak\lean{univRep},\allowbreak\ \lean{VOA.}\allowbreak\lean{DeligneAudit.}\allowbreak\lean{SectorWeyl.}\allowbreak\lean{wU},\allowbreak\ \lean{VOA.}\allowbreak\lean{DeligneAudit.}\allowbreak\lean{SectorWeyl.}\allowbreak\lean{Mchar},\allowbreak\ \lean{VOA.}\allowbreak\lean{DeligneAudit.}\allowbreak\lean{SectorWeyl.}\allowbreak\lean{univ_mabRel}\par}
{\raggedright\noindent File \leanpath{VOA/\allowbreak DeligneAudit/\allowbreak DeligneFusionUniv.lean}.\par}
\begin{leancode}
abbrev Ltr := Fin 8 × ℕ

abbrev FreeW := List Ltr →₀ ℂ

def commInst (p s : List Ltr) (x y : Ltr) : FreeW :=
  sw (p ++ x :: y :: s) - sw (p ++ y :: x :: s) -
    ∑ e, bracketTable x.1 y.1 e • sw (p ++ (e, x.2 + y.2) :: s)

def relTail (a b : ℕ) : Set (List Ltr) :=
  {t | (∃ e r, e.val < 3 ∧ t = [(e, r)]) ∨ (∃ h r, 6 ≤ h.val ∧ 1 ≤ r ∧ t = [(h, r)]) ∨
    t = List.replicate (a + b + 1) (3, 0) ∨ t = List.replicate (a + b + 1) (4, 0) ∨
    t = [(3, a + b)] ∨ t = [(4, a + b)] ∨ t = [(5, 2 * a + b)]}

def relGen (a b : ℕ) : Set FreeW :=
  {g | ∃ p : List Ltr, (∃ s x y, g = commInst p s x y) ∨ (∃ t ∈ relTail a b, g = sw (p ++ t)) ∨
    (∃ h : Fin 8, 6 ≤ h.val ∧ g = sw (p ++ [(h, 0)]) - ((a + b : ℕ) : ℂ) • sw p)}

def Jrel (a b : ℕ) : Submodule ℂ FreeW := Submodule.span ℂ (relGen a b)

abbrev UMod (a b : ℕ) := FreeW ⧸ Jrel a b

def univRep (a b : ℕ) : CurrentModule (bracketTable) (UMod a b) where
  op c r := opU a b (c, r)
  comm c d r s v := by

def wU (a b : ℕ) : UMod a b := Submodule.Quotient.mk (sw [])

def Mchar (a b : ℕ) : PowerSeries ℚ :=
  PowerSeries.mk fun N => (Module.finrank ℂ (cycDeg (univRep a b) (wU a b) N) : ℚ)
\end{leancode}
\begin{leancode}
theorem univ_mabRel (a b : ℕ) : MabRel (univRep a b) (wU a b) a b where
\end{leancode}
\noindent\textbf{Hypotheses / status.} None; proved\par
\medskip
\begin{btheorem}{26.14}[The exact-sequence step: quotient and rank--nullity]
(i) If $w$ satisfies the relations of $M_{a+1,b}$ and $S$ is an invariant subspace containing $f_\theta(2a+b+1)w$, then
the image of $w$ in $V/S$ satisfies the relations of $M_{a,b+1}$. (ii) Assume the kernel relation (KL). Then for every
degree $N$, with $s=2a+b+1$,
\begin{equation*}
\dim\ltwentysixcyc_{U_{a+1,b}}(w_U)_N\ \le\ \dim\ltwentysixcyc_{U_{a,b+1}}(w_U)_N+[s\le N]\,\dim\ltwentysixcyc_{U_{a,b}}(w_U)_{N-s},
\end{equation*}
provided the two right-hand pieces are finite-dimensional, and then the left piece is finite-dimensional too.
\end{btheorem}
{\raggedright\noindent\textbf{Lean.} \lean{VOA.}\allowbreak\lean{DeligneAudit.}\allowbreak\lean{SectorWeyl.}\allowbreak\lean{mabRel_quot},\allowbreak\ \lean{VOA.}\allowbreak\lean{DeligneAudit.}\allowbreak\lean{SectorWeyl.}\allowbreak\lean{vK},\allowbreak\ \lean{VOA.}\allowbreak\lean{DeligneAudit.}\allowbreak\lean{SectorWeyl.}\allowbreak\lean{fd_step},\allowbreak\ \lean{VOA.}\allowbreak\lean{DeligneAudit.}\allowbreak\lean{SectorWeyl.}\allowbreak\lean{finrank_step}\par}
{\raggedright\noindent File \leanpath{VOA/\allowbreak DeligneAudit/\allowbreak DeligneFusionGrade.lean}.\par}
\begin{leancode}
theorem mabRel_quot {w : V} {a b : ℕ} (h : MabRel ρ w (a + 1) b)
    (hv : ρ.op 5 (2 * a + b + 1) w ∈ S) : MabRel (quotRep ρ S hS) (S.mkQ w) a (b + 1) where
\end{leancode}
\begin{leancode}
def vK : UMod (a + 1) b := (univRep (a + 1) b).op 5 (2 * a + b + 1) (wU (a + 1) b)

theorem fd_step (hR : RavinderKernelRel) (N : ℕ)
    [FiniteDimensional ℂ (cycDeg (univRep a (b + 1)) (wU a (b + 1)) N)]
    [FiniteDimensional ℂ (cycDeg (univRep a b) (wU a b) (N - (2 * a + b + 1)))] :
    FiniteDimensional ℂ (cycDeg (univRep (a + 1) b) (wU (a + 1) b) N) := by

theorem finrank_step (hR : RavinderKernelRel) (N : ℕ)
    [FiniteDimensional ℂ (cycDeg (univRep a (b + 1)) (wU a (b + 1)) N)]
    [FiniteDimensional ℂ (cycDeg (univRep a b) (wU a b) (N - (2 * a + b + 1)))] :
    Module.finrank ℂ (cycDeg (univRep (a + 1) b) (wU (a + 1) b) N) ≤
      Module.finrank ℂ (cycDeg (univRep a (b + 1)) (wU a (b + 1)) N) +
        if 2 * a + b + 1 ≤ N then
          Module.finrank ℂ (cycDeg (univRep a b) (wU a b) (N - (2 * a + b + 1))) else 0 := by
\end{leancode}
\noindent\textbf{Hypotheses / status.} (KL), which is proved\par
\medskip
\begin{btheorem}{26.15}[The $A_2$ assembly]
(1) \emph{Structural statements.} Assume (TB) and (KL). Then all degree pieces of all $U_{a,b}$ are finite-dimensional; if
moreover U7 holds for every $n$, the universal characters \lean{Mchar} satisfy the squeeze inputs.
(2) \emph{Corollary.} If \lean{FusionTotalDim} holds and U7 holds for every $n$, then $L_{-3/2}(\mathfrak{sl}_3)$ is classically
free with kernel $(I_2^{\mathrm{ns}})$. ((KL) is a theorem, Section~26.7, and (TB) follows from
\lean{FusionTotalDim} by Theorem~26.7, so neither appears as a hypothesis of (2).)
\end{btheorem}
{\raggedright\noindent\textbf{Lean.} \lean{VOA.}\allowbreak\lean{DeligneAudit.}\allowbreak\lean{A2.}\allowbreak\lean{univ_fd},\allowbreak\ \lean{VOA.}\allowbreak\lean{DeligneAudit.}\allowbreak\lean{A2.}\allowbreak\lean{fusionSqueezeInputs_Mchar},\allowbreak\ \lean{VOA.}\allowbreak\lean{DeligneAudit.}\allowbreak\lean{A2.}\allowbreak\lean{classicallyFree_of_totalBound},\allowbreak\ \lean{VOA.}\allowbreak\lean{DeligneAudit.}\allowbreak\lean{A2.}\allowbreak\lean{classicallyFree_of_totalDim}\par}
{\raggedright\noindent File \leanpath{VOA/\allowbreak DeligneAudit/\allowbreak DeligneFusionAssembly.lean}.\par}
\begin{leancode}
theorem univ_fd (hF : FusionTotalBound) (hR : RavinderKernelRel) :
    ∀ a b N, FiniteDimensional ℂ (cycDeg (univRep a b) (SectorWeyl.wU a b) N)

theorem fusionSqueezeInputs_Mchar (hF : FusionTotalBound) (hR : RavinderKernelRel)
    (hD : ∀ n, DemazureRealization n) : FusionSqueezeInputs Mchar where

theorem classicallyFree_of_totalBound (hF : FusionTotalBound) (hD : ∀ n, DemazureRealization n) :
    CaseClassicallyFree .A2 W explicitPresentation ∧
      RingHom.ker (MvPolynomial.aeval fun i =>
        liDegreeZeroClass W 1 ⟨explicitPresentation.gen i,
          explicitPresentation.gen_weight_one i⟩).toRingHom = I2ns :=

theorem classicallyFree_of_totalDim (hF : FusionTotalDim) (hD : ∀ n, DemazureRealization n) :
    CaseClassicallyFree .A2 W explicitPresentation ∧
      RingHom.ker (MvPolynomial.aeval fun i =>
        liDegreeZeroClass W 1 ⟨explicitPresentation.gen i,
          explicitPresentation.gen_weight_one i⟩).toRingHom = I2ns :=
\end{leancode}
\noindent\textbf{Hypotheses / status.} \lean{FusionTotalDim}, U7\par
\medskip
\begin{btheorem}{26.16}[Heisenberg nilpotency; Garland's identity with $P(u)=1$]
(1) Let $X,Y,Z$ be operators with $[X,Y]=Z$ and $Z$ commuting with $X$ and $Y$. If $X^{p+1}w=Y^{p+1}w=0$ then
$Z^{2p+1}w=0$. (2) Let $E_1$, $H_k$ ($k\ge1$), $F_k$ ($k\ge0$) be operators with $[E_1,F_k]=H_{k+1}$,
$[H_{m+1},F_k]=-2F_{m+1+k}$ and $[F_j,F_k]=0$, and let $E_1w=0$, $H_{m+1}w=0$ for all $m$. If $F_0^Nw=0$, then for
$r\ge1$ and $r+s\ge N$,
\begin{equation*}
[u^s]\,F(u)^r\,w=\sum_{k_1+\dots+k_r=s}F_{k_1}\cdots F_{k_r}\,w=0 .
\end{equation*}
\end{btheorem}
{\raggedright\noindent\textbf{Lean.} \lean{VOA.}\allowbreak\lean{DeligneAudit.}\allowbreak\lean{Garland.}\allowbreak\lean{heis_nilpotent},\allowbreak\ \lean{VOA.}\allowbreak\lean{DeligneAudit.}\allowbreak\lean{Garland.}\allowbreak\lean{Fs},\allowbreak\ \lean{VOA.}\allowbreak\lean{DeligneAudit.}\allowbreak\lean{Garland.}\allowbreak\lean{garland}\par}
{\raggedright\noindent File \leanpath{VOA/\allowbreak DeligneAudit/\allowbreak DeligneFusionGarland.lean}.\par}
\begin{leancode}
variable {X Y Z : Module.End ℂ V} (hXY : X * Y - Y * X = Z) (hXZ : X * Z = Z * X)
  (hYZ : Y * Z = Z * Y)

theorem heis_nilpotent (hXZ : X * Z = Z * X) {p : ℕ} {w : V} (hX : (X ^ (p + 1)) w = 0)
    (hY : (Y ^ (p + 1)) w = 0) : (Z ^ (2 * p + 1)) w = 0 := by

def Fs : PowerSeries (Module.End ℂ V) := PowerSeries.mk F


variable {E1 H F} (hEF : ∀ k, E1 * F k - F k * E1 = H (k + 1))
  (hHF : ∀ m k, H (m + 1) * F k - F k * H (m + 1) = (-2 : ℂ) • F (m + 1 + k))

theorem garland (r : ℕ) (hr : 1 ≤ r) {N : ℕ} (hN : (F 0 ^ N) w = 0) (s : ℕ) (hs : N ≤ r + s) :
    (coeff s (Fs F ^ r) : Module.End ℂ V) w = 0 := by
\end{leancode}
\noindent\textbf{Hypotheses / status.} Five \lean{include}d hyps\par
\medskip
\begin{btheorem}{26.17}[Ravinder's kernel lemma for $\mathfrak{sl}_3$]
Let $w$ satisfy the relations of $M_{a,b}$. Then $f_\theta(0)^{2(a+b)+1}w=0$ and
$\sum_{i+j=s}f_\theta(i)f_\theta(j)w=0$ for $s+1\ge2(a+b)$. If $a\ge1$, the kernel vector $v=f_\theta(2a+b-1)w$ satisfies
$f_1(a-1+b)v=f_2(a-1+b)v=0$ and $f_\theta(2(a-1)+b)v=0$, and therefore all relations of $M_{a-1,b}$.
\end{btheorem}
{\raggedright\noindent\textbf{Lean.} \lean{VOA.}\allowbreak\lean{DeligneAudit.}\allowbreak\lean{SectorWeyl.}\allowbreak\lean{RavinderKernelTrunc},\allowbreak\ \lean{VOA.}\allowbreak\lean{DeligneAudit.}\allowbreak\lean{SectorWeyl.}\allowbreak\lean{ravinderKernelRel_of_trunc},\allowbreak\ \lean{VOA.}\allowbreak\lean{DeligneAudit.}\allowbreak\lean{SectorWeyl.}\allowbreak\lean{ftheta_pow},\allowbreak\ \lean{VOA.}\allowbreak\lean{DeligneAudit.}\allowbreak\lean{SectorWeyl.}\allowbreak\lean{theta_garland},\allowbreak\ \lean{VOA.}\allowbreak\lean{DeligneAudit.}\allowbreak\lean{SectorWeyl.}\allowbreak\lean{ravinderKernelTrunc},\allowbreak\ \lean{VOA.}\allowbreak\lean{DeligneAudit.}\allowbreak\lean{SectorWeyl.}\allowbreak\lean{ravinderKernelRel}\par}
{\raggedright\noindent File \leanpath{VOA/\allowbreak DeligneAudit/\allowbreak DeligneFusionKernel.lean}.\par}
\begin{leancode}
def RavinderKernelTrunc : Prop :=
  ∀ (a b : ℕ) (V : Type) [AddCommGroup V] [Module ℂ V] (ρ : CurrentModule (bracketTable) V) (w : V),
    1 ≤ a → MabRel ρ w a b →
      ρ.op 3 (a - 1 + b) (ρ.op 5 (2 * a + b - 1) w) = 0 ∧
      ρ.op 4 (a - 1 + b) (ρ.op 5 (2 * a + b - 1) w) = 0 ∧
      ρ.op 5 (2 * (a - 1) + b) (ρ.op 5 (2 * a + b - 1) w) = 0

theorem ravinderKernelRel_of_trunc (hT : RavinderKernelTrunc) : RavinderKernelRel := by
\end{leancode}
{\raggedright\noindent File \leanpath{VOA/\allowbreak DeligneAudit/\allowbreak DeligneFusionGarlandSl3.lean}.\par}
\begin{leancode}
theorem ftheta_pow (h : MabRel ρ w a b) : ((ρ.op 5 0) ^ (2 * (a + b) + 1)) w = 0 := by

theorem theta_garland (h : MabRel ρ w a b) (s : ℕ) (hs : 2 * (a + b) ≤ s + 1) :
    ∑ x ∈ Finset.antidiagonal s, ρ.op 5 x.1 (ρ.op 5 x.2 w) = 0 := by

theorem ravinderKernelTrunc : RavinderKernelTrunc := by

theorem ravinderKernelRel : RavinderKernelRel := ravinderKernelRel_of_trunc ravinderKernelTrunc
\end{leancode}
\noindent\textbf{Hypotheses / status.} None; proved\par
\medskip
\begin{btheorem}{26.18}[Chari--Pressley, Prop.~6.1]
In $R=\C[z_0,z_1,\dots]$ let $Z(u)=\sum_iz_iu^i$, $x(r,s)=[u^s]Z(u)^r$ and $J_m=(x(r,s):r\ge1,\ r+s\ge m+1)$. Then every
$p\in R$ is a combination of \emph{box monomials} $z_{i_1}\cdots z_{i_\ell}$ with $i_j+\ell\le m$, plus an element of $J_m$.
Consequently, in every $\mathfrak{sl}_2[t]$-module, if $w$ is killed by $e\otimes t$ and by $h(k)$, $k\ge1$, and
$f(0)^{m+1}w=0$, then every polynomial in the $f(k)$ applied to $w$ is a combination of box monomials applied to $w$.
\end{btheorem}
{\raggedright\noindent\textbf{Lean.} \lean{VOA.}\allowbreak\lean{DeligneAudit.}\allowbreak\lean{Sl2Weyl.}\allowbreak\lean{Zs},\allowbreak\ \lean{VOA.}\allowbreak\lean{DeligneAudit.}\allowbreak\lean{Sl2Weyl.}\allowbreak\lean{xx},\allowbreak\ \lean{VOA.}\allowbreak\lean{DeligneAudit.}\allowbreak\lean{Sl2Weyl.}\allowbreak\lean{J},\allowbreak\ \lean{VOA.}\allowbreak\lean{DeligneAudit.}\allowbreak\lean{Sl2Weyl.}\allowbreak\lean{BoxP},\allowbreak\ \lean{VOA.}\allowbreak\lean{DeligneAudit.}\allowbreak\lean{Sl2Weyl.}\allowbreak\lean{Bk},\allowbreak\ \lean{VOA.}\allowbreak\lean{DeligneAudit.}\allowbreak\lean{Sl2Weyl.}\allowbreak\lean{T},\allowbreak\ \lean{VOA.}\allowbreak\lean{DeligneAudit.}\allowbreak\lean{Sl2Weyl.}\allowbreak\lean{box_span},\allowbreak\ \lean{VOA.}\allowbreak\lean{DeligneAudit.}\allowbreak\lean{Sl2Weyl.}\allowbreak\lean{weyl_span_box}\par}
{\raggedright\noindent File \leanpath{VOA/\allowbreak DeligneAudit/\allowbreak DeligneFusionSl2.lean}.\par}
\begin{leancode}
def Zs : PowerSeries R := PowerSeries.mk fun i => z i

def xx (r s : ℕ) : R := PowerSeries.coeff s (Zs ^ r)

def J (m : ℕ) : Ideal R := Ideal.span {p | ∃ r s, 1 ≤ r ∧ m + 1 ≤ r + s ∧ p = xx r s}

def BoxP (m : ℕ) (f : ℕ →₀ ℕ) : Prop := ∀ i ∈ f.support, i + Finsupp.degree f ≤ m

def Bk (m k : ℕ) : Submodule ℂ R :=
  Submodule.span ℂ {p | ∃ f, BoxP m f ∧ f 0 ≤ k ∧ p = MvPolynomial.monomial f 1}

def T (m k : ℕ) : Submodule ℂ R := Bk m k ⊔ (J m).restrictScalars ℂ

theorem box_span : ∀ m (p : R), p ∈ T m m := by
\end{leancode}
{\raggedright\noindent File \leanpath{VOA/\allowbreak DeligneAudit/\allowbreak DeligneFusionSl2Module.lean}.\par}
\begin{leancode}
theorem weyl_span_box {E1 : Module.End ℂ V} {H : ℕ → Module.End ℂ V}
    (hEF : ∀ k, E1 * F k - F k * E1 = H (k + 1))
    (hHF : ∀ m k, H (m + 1) * F k - F k * H (m + 1) = (-2 : ℂ) • F (m + 1 + k))
    {y : V} (hE : E1 y = 0) (hH : ∀ m, H (m + 1) y = 0) (m : ℕ) (hN : (F 0 ^ (m + 1)) y = 0)
    (p : R) :
    evAlg F hFF p y ∈
      Submodule.span ℂ {v | ∃ f, BoxP m f ∧ v = evAlg F hFF (MvPolynomial.monomial f 1) y} :=
\end{leancode}
\noindent\textbf{Hypotheses / status.} None; proved\par
\medskip
\begin{btheorem}{26.20}[Chari--Loktev spanning and count for $\mathfrak{sl}_3$]
For every $N$, \lean{CLSpanning N} and \lean{CLCount N} hold. Consequently, in every module with the $M_{N,0}$ relations,
$\dim\ltwentysixcyc(w)_d\le[q^d]\ltwentysixD_N$ for every $d$.
\end{btheorem}
{\raggedright\noindent\textbf{Lean.} \lean{VOA.}\allowbreak\lean{DeligneAudit.}\allowbreak\lean{CLAssembly.}\allowbreak\lean{clSpanning},\allowbreak\ \lean{VOA.}\allowbreak\lean{DeligneAudit.}\allowbreak\lean{CLAssembly.}\allowbreak\lean{finrank_cycDeg_le},\allowbreak\ \lean{VOA.}\allowbreak\lean{DeligneAudit.}\allowbreak\lean{CL.}\allowbreak\lean{clCount}\par}
{\raggedright\noindent File \leanpath{VOA/\allowbreak DeligneAudit/\allowbreak DeligneFusionCLAssembly.lean}.\par}
\begin{leancode}
theorem clSpanning (N : ℕ) : CLSpanning N := by

theorem finrank_cycDeg_le {V : Type} [AddCommGroup V] [Module ℂ V] {ρ : CurrentModule (bracketTable) V}
    {w : V} {N : ℕ} (hR : MabRel ρ w N 0) (d : ℕ) :
    Module.finrank ℂ (cycDeg ρ w d) ≤ A2.demazureHilb N d := by
\end{leancode}
{\raggedright\noindent File \leanpath{VOA/\allowbreak DeligneAudit/\allowbreak DeligneFusionCLCount.lean}.\par}
\begin{leancode}
theorem clCount (N : ℕ) : CLCount N := by
\end{leancode}
\noindent\textbf{Hypotheses / status.} None; proved\par
\medskip
\begin{bdefinition}{26.23}[Filtered relations; the evaluation model]
Let $T$ be a current module (no grading) and $w\in T$; let $\Phi_{<d}$ be the span of words of degree $<d$ applied to $w$
and $\Phi_{\le d}=\Phi_{<d+1}$. Say $w$ satisfies the \emph{filtered $M_{N,0}$ relations} if the raising and lowering
relations hold exactly, $h(r)w\in\C w$, $h(0)w=Nw$, and the truncations hold modulo lower filtration:
$f_i(N)w\in\Phi_{<N}$, $f_\theta(2N)w\in\Phi_{<2N}$. The \emph{evaluation model} is $T_N=\C^{\{1,2,3\}^{2N}}$, the tensor
product of $N$ copies of the standard representation $\C^3$ and $N$ copies of its dual, at the distinct points
$z=0,1,\dots,2N-1$, with $x(k)$ acting by $\sum_iz_i^k\,x^{(i)}$ ($x^{(i)}$ the action on factor $i$).
\end{bdefinition}
{\raggedright\noindent\textbf{Lean.} \lean{VOA.}\allowbreak\lean{DeligneAudit.}\allowbreak\lean{GrFilt.}\allowbreak\lean{Fbelow},\allowbreak\ \lean{VOA.}\allowbreak\lean{DeligneAudit.}\allowbreak\lean{GrFilt.}\allowbreak\lean{Fle},\allowbreak\ \lean{VOA.}\allowbreak\lean{DeligneAudit.}\allowbreak\lean{GrFilt.}\allowbreak\lean{FiltRel},\allowbreak\ \lean{VOA.}\allowbreak\lean{DeligneAudit.}\allowbreak\lean{EvalModel.}\allowbreak\lean{matT},\allowbreak\ \lean{VOA.}\allowbreak\lean{DeligneAudit.}\allowbreak\lean{EvalModel.}\allowbreak\lean{Site},\allowbreak\ \lean{VOA.}\allowbreak\lean{DeligneAudit.}\allowbreak\lean{EvalModel.}\allowbreak\lean{Cfg},\allowbreak\ \lean{VOA.}\allowbreak\lean{DeligneAudit.}\allowbreak\lean{EvalModel.}\allowbreak\lean{TT},\allowbreak\ \lean{VOA.}\allowbreak\lean{DeligneAudit.}\allowbreak\lean{EvalModel.}\allowbreak\lean{R},\allowbreak\ \lean{VOA.}\allowbreak\lean{DeligneAudit.}\allowbreak\lean{EvalModel.}\allowbreak\lean{z},\allowbreak\ \lean{VOA.}\allowbreak\lean{DeligneAudit.}\allowbreak\lean{EvalModel.}\allowbreak\lean{loc},\allowbreak\ \lean{VOA.}\allowbreak\lean{DeligneAudit.}\allowbreak\lean{EvalModel.}\allowbreak\lean{op}\par}
{\raggedright\noindent File \leanpath{VOA/\allowbreak DeligneAudit/\allowbreak DeligneFusionGrFilt.lean}.\par}
\begin{leancode}
def Fbelow (d : ℕ) : Submodule ℂ T :=
  Submodule.span ℂ {v | ∃ l : List (Fin 8 × ℕ), SectorGraded.wordDeg l < d ∧ v = wordOp ρ l w}

abbrev Fle (d : ℕ) : Submodule ℂ T := Fbelow ρ w (d + 1)

structure FiltRel (N : ℕ) : Prop where
  raising : ∀ (e : Fin 8), e.val < 3 → ∀ r, ρ.op e r w = 0
  cartan : ∀ (h : Fin 8), 6 ≤ h.val → ∀ r, ρ.op h r w ∈ ℂ ∙ w
  cartan_zero : ∀ (h : Fin 8), 6 ≤ h.val → ρ.op h 0 w = (N : ℂ) • w
  lowering_f1 : (ρ.op 3 0)^[N + 1] w = 0
  lowering_f2 : (ρ.op 4 0)^[N + 1] w = 0
  trunc_f1 : ρ.op 3 N w ∈ Fbelow ρ w N
  trunc_f2 : ρ.op 4 N w ∈ Fbelow ρ w N
  trunc_theta : ρ.op 5 (2 * N) w ∈ Fbelow ρ w (2 * N)
\end{leancode}
{\raggedright\noindent File \leanpath{VOA/\allowbreak DeligneAudit/\allowbreak DeligneFusionEvalModel.lean}.\par}
\begin{leancode}
def matT : Fin 8 → Matrix (Fin 3) (Fin 3) ℂ :=
  ![Matrix.single 0 1 1, Matrix.single 1 2 1, Matrix.single 0 2 1, Matrix.single 1 0 1,
    Matrix.single 2 1 1, Matrix.single 2 0 1, Matrix.single 0 0 1 - Matrix.single 1 1 1,
    Matrix.single 1 1 1 - Matrix.single 2 2 1]

abbrev Site := Fin N ⊕ Fin N
abbrev Cfg := Site N → Fin 3
abbrev TT := Cfg N → ℂ

def R : Site N → Fin 8 → Matrix (Fin 3) (Fin 3) ℂ
  | .inl _, c => matT c
  | .inr _, c => -(Matrix.transpose (matT c))

def z : Site N → ℂ
  | .inl i => ((i : ℕ) : ℂ)
  | .inr i => ((N + i : ℕ) : ℂ)

def loc (i : Site N) (X : Matrix (Fin 3) (Fin 3) ℂ) : Module.End ℂ (TT N) where
  toFun f γ := ∑ k, X (γ i) k * f (Function.update γ i k)
  map_add' f g := by ext γ; simp [mul_add, Finset.sum_add_distrib]
  map_smul' c f := by ext γ; simp [Finset.mul_sum, mul_left_comm]

def op (c : Fin 8) (k : ℕ) : Module.End ℂ (TT N) := ∑ i, (z N i) ^ k • loc i (R N i c)
\end{leancode}
\medskip
\begin{btheorem}{26.24}[U7 for $A_2$]
(1) Let $T$ be a \emph{finite-dimensional} current module. If $w\in T$ satisfies the filtered $M_{N,0}$ relations and
$\dim\Phi_{\le D}\ge9^N$ for some $D$, then \lean{DemazureRealization N} holds. (2) The evaluation model satisfies the hypotheses of (1). Hence
\lean{DemazureRealization N} holds for every $N$.
\end{btheorem}
{\raggedright\noindent\textbf{Lean.} \lean{VOA.}\allowbreak\lean{DeligneAudit.}\allowbreak\lean{GrFilt.}\allowbreak\lean{demazureRealization_of_filt},\allowbreak\ \lean{VOA.}\allowbreak\lean{DeligneAudit.}\allowbreak\lean{EvalModel.}\allowbreak\lean{demazureRealization_all}\par}
{\raggedright\noindent File \leanpath{VOA/\allowbreak DeligneAudit/\allowbreak DeligneFusionGrFilt.lean}.\par}
\begin{leancode}
theorem demazureRealization_of_filt {N : ℕ} (h : FiltRel ρ w N) (D : ℕ)
    (hD : 9 ^ N ≤ Module.finrank ℂ (Fle ρ w D)) : A2.DemazureRealization N := by
\end{leancode}
{\raggedright\noindent File \leanpath{VOA/\allowbreak DeligneAudit/\allowbreak DeligneFusionEvalModel.lean}.\par}
\begin{leancode}
theorem demazureRealization_all (N : ℕ) : A2.DemazureRealization N := by
\end{leancode}
\noindent\textbf{Hypotheses / status.} None; proved\par
\medskip
\begin{btheorem}{26.27}[Each refinement implies the total bound]
(1) \lean{RavinderKernelLower} implies $\sum_{N<B}\dim\ltwentysixcyc_{U_{a,b}}(w_U)_N\le9^a8^b$ for all $a,b,B$, hence
\lean{FusionTotalDim}. (2) In a module with the relations of $M_{0,N}$, every mode $x(k)$ with $k\ge N$ kills the whole
cyclic module. (3) \lean{IdealLower} implies \lean{FusionTotalDim}.
\end{btheorem}
{\raggedright\noindent\textbf{Lean.} \lean{VOA.}\allowbreak\lean{DeligneAudit.}\allowbreak\lean{A2.}\allowbreak\lean{uTot_le},\allowbreak\ \lean{VOA.}\allowbreak\lean{DeligneAudit.}\allowbreak\lean{A2.}\allowbreak\lean{fusionTotalDim_of_kernelLower},\allowbreak\ \lean{VOA.}\allowbreak\lean{DeligneAudit.}\allowbreak\lean{A2.}\allowbreak\lean{trunc_all},\allowbreak\ \lean{VOA.}\allowbreak\lean{DeligneAudit.}\allowbreak\lean{A2.}\allowbreak\lean{fusionTotalDim_of_idealLower}\par}
{\raggedright\noindent File \leanpath{VOA/\allowbreak DeligneAudit/\allowbreak DeligneFusionKernelLower.lean}.\par}
\begin{leancode}
theorem uTot_le (hK : RavinderKernelLower) : ∀ b a B, uTot a b B ≤ 9 ^ a * 8 ^ b

theorem fusionTotalDim_of_kernelLower (hK : RavinderKernelLower) : FusionTotalDim := by
\end{leancode}
{\raggedright\noindent File \leanpath{VOA/\allowbreak DeligneAudit/\allowbreak DeligneFusionIdeal.lean}.\par}
\begin{leancode}
theorem trunc_all {V : Type} [AddCommGroup V] [Module ℂ V] {ρ : CurrentModule (bracketTable) V}
    {w : V} {N : ℕ} (h : MabRel ρ w 0 N) :
    ∀ (l : List (Fin 8 × ℕ)) (c : Fin 8) (k : ℕ), N ≤ k → ρ.op c k (wordOp ρ l w) = 0

theorem fusionTotalDim_of_idealLower (hI : IdealLower) : SectorWeyl.FusionTotalDim := by
\end{leancode}
\noindent\textbf{Hypotheses / status.} The named \lean{Prop}; proved\par
\medskip
\begin{btheorem}{26.28}[The $A_2$ case: final forms in this chapter]
$L_{-3/2}(\mathfrak{sl}_3)$ is classically free for its explicit presentation, with kernel $(I_2^{\mathrm{ns}})$, assuming
either \lean{RavinderKernelLower} or \lean{FusionTotalDim}.
\end{btheorem}
{\raggedright\noindent\textbf{Lean.} \lean{VOA.}\allowbreak\lean{DeligneAudit.}\allowbreak\lean{A2.}\allowbreak\lean{classicallyFree_of_kernelLower'},\allowbreak\ \lean{VOA.}\allowbreak\lean{DeligneAudit.}\allowbreak\lean{A2.}\allowbreak\lean{classicallyFree_of_totalDim'}\par}
{\raggedright\noindent File \leanpath{VOA/\allowbreak DeligneAudit/\allowbreak DeligneFusionU7.lean}.\par}
\begin{leancode}
theorem classicallyFree_of_kernelLower' (hK : RavinderKernelLower) :
    CaseClassicallyFree .A2 W explicitPresentation ∧
      RingHom.ker (MvPolynomial.aeval fun i =>
        liDegreeZeroClass W 1 ⟨explicitPresentation.gen i,
          explicitPresentation.gen_weight_one i⟩).toRingHom = I2ns :=

theorem classicallyFree_of_totalDim' (hF : SectorWeyl.FusionTotalDim) :
    CaseClassicallyFree .A2 W explicitPresentation ∧
      RingHom.ker (MvPolynomial.aeval fun i =>
        liDegreeZeroClass W 1 ⟨explicitPresentation.gen i,
          explicitPresentation.gen_weight_one i⟩).toRingHom = I2ns :=
\end{leancode}
\noindent\textbf{Hypotheses / status.} \textbf{\lean{FusionTotalDim}} (proved, so unconditional in Ch.~27), resp.\ \textbf{\lean{RavinderKernelLower}}\par
\medskip
\begin{btheorem}{26.30}[The $A_2$ total bound]
\leavevmode
\begin{enumerate}
\item The two base families hold: \lean{CVUniv.sharpRectangularBases} and \lean{CVUniv.residualTheta3Bound}.
\item If both base families hold, then \lean{FusionTotalDim}: for every $m$, every current module and every vector $w$
  satisfying the local Weyl and fusion relations for $m\theta$, $\sum_{b<B}\dim\ltwentysixcyc(w)_b\le8^m$ for all $B$.
\item Hence \lean{FusionTotalDim} holds, with no hypothesis.
\end{enumerate}
\end{btheorem}
{\raggedright\noindent\textbf{Lean.} \lean{VOA.}\allowbreak\lean{DeligneAudit.}\allowbreak\lean{A2.}\allowbreak\lean{CVUniv.}\allowbreak\lean{sharpRectangularBases},\allowbreak\ \lean{VOA.}\allowbreak\lean{DeligneAudit.}\allowbreak\lean{A2.}\allowbreak\lean{CVUniv.}\allowbreak\lean{residualTheta3Bound},\allowbreak\ \lean{VOA.}\allowbreak\lean{DeligneAudit.}\allowbreak\lean{A2.}\allowbreak\lean{CVUniv.}\allowbreak\lean{fusionTotalDim_of_cvBases},\allowbreak\ \lean{VOA.}\allowbreak\lean{DeligneAudit.}\allowbreak\lean{A2.}\allowbreak\lean{CVUniv.}\allowbreak\lean{fusionTotalDim}\par}
{\raggedright\noindent Files \leanpath{VOA/\allowbreak DeligneAudit/\allowbreak DeligneFusionCVRectFilteredRectangular.lean} (1, first
family), \leanpath{VOA/\allowbreak DeligneAudit/\allowbreak DeligneFusionCVResidualThetaBoundFinal.lean} (1, second family),
\leanpath{VOA/\allowbreak DeligneAudit/\allowbreak DeligneFusionCVTotal.lean} (2) and
\leanpath{VOA/\allowbreak DeligneAudit/\allowbreak A2Unconditional.lean} (3).\par}
\begin{leancode}
theorem sharpRectangularBases : SharpRectangularBases := by

theorem residualTheta3Bound : ResidualTheta3Bound := by

theorem fusionTotalDim_of_cvBases (hR : SharpRectangularBases)
    (hΘ : ResidualTheta3Bound) : FusionTotalDim :=

theorem CVUniv.fusionTotalDim : SectorWeyl.FusionTotalDim :=
  CVUniv.fusionTotalDim_of_cvBases CVUniv.sharpRectangularBases CVUniv.residualTheta3Bound
\end{leancode}
\noindent\textbf{Hypotheses / status.} \textbf{none}; proved (record)\par
\medskip
\begin{btheorem}{26.41}[$e^{N\theta}$ satisfies Ravinder's $M_{N,0}$ relations; lattice U7]
(1) For lattice realization data with $\langle\theta,\theta\rangle=2$ and every $N$, $e^{N\theta}$ satisfies \lean{ravFam S N 0}.
(2) Hence \lean{LatDim Dem N} implies \lean{DemazureRealizationG S (ravFam S) Dem N}.
(3) Realization data exist for $D_4$, $E_6$, $E_7$, $E_8$ (\lean{d4LatData}, \lean{e6LatData}, \lean{e7LatData}, \lean{e8LatData},
with $\langle\theta,\theta\rangle=2$ proved as \lean{d4_hθ}, \dots, \lean{e8_hθ}); for these types the relations half of U7 is
unconditional.
(4) (Corollary for $D_4$.) If \lean{FusionTotalBoundG D4sector} holds and \lean{D4LatticeDim N} holds for every $N$, then $D_4$ is
classically free for its explicit presentation, with the expected kernel (\lean{d4_classicallyFree_of_totalBound_latDim}). The
$E_6$, $E_7$ analogues take additional ``star'' inputs (Chapter~25).
There are no lattice data for $G_2$, $F_4$: the Frenkel--Kac construction realizes level-one currents of \emph{simply laced}
algebras only, since it needs all roots to have norm $2$ in an even lattice.
\end{btheorem}
{\raggedright\noindent\textbf{Lean.} \lean{VOA.}\allowbreak\lean{DeligneAudit.}\allowbreak\lean{LatticeRealization.}\allowbreak\lean{LatData.}\allowbreak\lean{sat},\allowbreak\ \lean{VOA.}\allowbreak\lean{DeligneAudit.}\allowbreak\lean{LatticeRealization.}\allowbreak\lean{LatData.}\allowbreak\lean{demazureRealizationG},\allowbreak\ \lean{VOA.}\allowbreak\lean{DeligneAudit.}\allowbreak\lean{LatticeRealization.}\allowbreak\lean{D4.}\allowbreak\lean{d4_lattice_sat},\allowbreak\ \lean{VOA.}\allowbreak\lean{DeligneAudit.}\allowbreak\lean{LatticeRealization.}\allowbreak\lean{D4.}\allowbreak\lean{D4LatticeDim},\allowbreak\ \lean{VOA.}\allowbreak\lean{DeligneAudit.}\allowbreak\lean{LatticeRealization.}\allowbreak\lean{D4.}\allowbreak\lean{d4_classicallyFree_of_totalBound_latDim},\allowbreak\ \lean{VOA.}\allowbreak\lean{DeligneAudit.}\allowbreak\lean{LatticeRealization.}\allowbreak\lean{E8.}\allowbreak\lean{e8_lattice_sat},\allowbreak\ \lean{VOA.}\allowbreak\lean{DeligneAudit.}\allowbreak\lean{LatticeRealization.}\allowbreak\lean{E8.}\allowbreak\lean{E8LatticeDim},\allowbreak\ \lean{VOA.}\allowbreak\lean{DeligneAudit.}\allowbreak\lean{LatticeRealization.}\allowbreak\lean{E8.}\allowbreak\lean{e8_demazureRealization}\par}
{\raggedright\noindent File \leanpath{VOA/\allowbreak DeligneAudit/\allowbreak LatticeRealizationGeneric.lean}.\par}
\begin{leancode}
theorem sat (hθ : L.form D.θ D.θ = 2) (N : ℕ) : Sat S.f D.ρ (D.vN N) (ravFam S N 0) := by

theorem demazureRealizationG (hθ : L.form D.θ D.θ = 2) {Dem : ℕ → PowerSeries ℚ} {N : ℕ}
    (hdim : D.LatDim Dem N) : DemazureRealizationG S (ravFam S) Dem N :=
\end{leancode}
{\raggedright\noindent File \leanpath{VOA/\allowbreak DeligneAudit/\allowbreak LatticeRealizationD4.lean}.\par}
\begin{leancode}
theorem d4_lattice_sat (N : ℕ) : Sat D4sector.f d4LatData.ρ (d4LatData.vN N) (ravFam D4sector N 0) :=

def D4LatticeDim (N : ℕ) : Prop :=
  d4LatData.LatDim DeligneCharacter.D4Wp.hilbSeriesZ N

theorem d4_classicallyFree_of_totalBound_latDim (hF : FusionTotalBoundG D4sector)
    (hdim : ∀ N, D4LatticeDim N) :
    VOA.DeligneHigherRankConditionalCapstone.CaseClassicallyFree .D4 D4Lie.W4 D4Lie.explicitPresentation4 ∧
      RingHom.ker (MvPolynomial.aeval fun i =>
        liDegreeZeroClass D4Lie.W4 1 ⟨D4Lie.explicitPresentation4.gen i,
          D4Lie.explicitPresentation4.gen_weight_one i⟩).toRingHom =
        D4Lie.explicitPresentation4.ideal :=
\end{leancode}
{\raggedright\noindent File \leanpath{VOA/\allowbreak DeligneAudit/\allowbreak LatticeRealizationE8.lean}.\par}
\begin{leancode}
theorem e8_lattice_sat (N : ℕ) : Sat E8sector.f e8LatData.ρ (e8LatData.vN N) (ravFam E8sector N 0) :=

def E8LatticeDim (N : ℕ) : Prop :=
  e8LatData.LatDim (DeligneCarrier.hilbSeries DeligneCharacter.e8DemazureHilb) N

theorem e8_demazureRealization {N : ℕ} (h : E8LatticeDim N) :
    DemazureRealizationG E8sector (ravFam E8sector) (DeligneCarrier.hilbSeries DeligneCharacter.e8DemazureHilb) N :=
\end{leancode}
\noindent\textbf{Hypotheses / status.} \textbf{\lean{FusionTotalBoundG}}, \textbf{\lean{D4LatticeDim}}\par
\medskip
\endgroup
\setcounter{chapter}{26}
\chapter[The Deligne series IV: final theorems and named inputs]{The Deligne series IV: the final theorems and the named inputs}
\begingroup
\providecommand{\ltwentyseveng}{\mathfrak g}
\providecommand{\ltwentysevenIns}{I_2^{\mathrm{ns}}}
\providecommand{\ltwentysevenHilb}{\operatorname{Hilb}}
\providecommand{\ltwentysevenvac}{\mathbf 1}
\providecommand{\ltwentysevenJ}{J_\infty}
\providecommand{\ltwentysevenF}{\mathcal F}
\providecommand{\ltwentysevengr}{\operatorname{gr}}
\providecommand{\ltwentysevendg}{d_{\mathfrak g}}
\noindent{\small 7 of the chapter's 21 numbered statements: its headline results and every definition they use.}\par\medskip
\begin{blemma}{27.6}[The fusion module is a truncated Weyl module]
Let $m\ge0$ and let $w$ be a vector of a $\ltwentyseveng[t]$-module satisfying the local Weyl relations of
$W_{\mathrm{loc}}(m\theta)$, together with $x(m)w=0$ for every $x$ in a set $\mathcal S$ of negative root vectors with
$\{f_\alpha:\langle\alpha,\theta^\vee\rangle=1\}\cup\{f_\theta\}\subseteq\mathcal S$. Then $(\ltwentyseveng\otimes t^m\C[t])$
annihilates the cyclic module $U(\ltwentyseveng[t])w$. Consequently the universal module with these relations,
$\ltwentysevenF_m=M_{0,m}$, is the \emph{truncated Weyl module} $W(m\theta,m)$ of Kus--Littelmann (Definition~4.1): the
quotient of $W_{\mathrm{loc}}(m\theta)$ on which $\ltwentyseveng\otimes t^m\C[t]$ acts by zero. In particular $M_{0,m}$ does not
depend on the choice of $\mathcal S$ within these bounds.
\end{blemma}
\noindent\emph{(The chapter quotes no Lean declaration for this statement; see book p.~1460.)}\par
\noindent\textbf{Hypotheses / status.} On paper; no Lean statement\par
\medskip
\begin{bdefinition}{27.7}[Input R]
Let $S$ be sector data, $\mathrm{spec}$ a family of relation specifications indexed by $\N\times\N$, and
$\mathrm{Dem}:\N\to\Q[[q]]$ a family of power series. For $j\in\N$, a \emph{Demazure realization at $j$} is a current module
$V$ and a vector $w\in V$ satisfying $\mathrm{spec}(j,0)$ such that, for every degree $N$,
\begin{equation*}
[q^N]\,\mathrm{Dem}_j\ \le\ \dim\bigl(\text{degree-}N\text{ part of the cyclic span of }w\bigr).
\end{equation*}
In the final theorems $\mathrm{spec}=$ \lean{ravFam S}, so $\mathrm{spec}(j,0)$ are the relations of $M_{j,0}$, and $\mathrm{Dem}_j$
is the type's Demazure table from Chapter~25. Input R is ``a Demazure realization exists at every $j$''.
\end{bdefinition}
{\raggedright\noindent\textbf{Lean.} \lean{VOA.}\allowbreak\lean{DeligneAudit.}\allowbreak\lean{UnivG.}\allowbreak\lean{DemazureRealizationG},\allowbreak\ \lean{VOA.}\allowbreak\lean{DeligneAudit.}\allowbreak\lean{DeligneCarrier.}\allowbreak\lean{hilbSeries},\allowbreak\ \lean{VOA.}\allowbreak\lean{DeligneAudit.}\allowbreak\lean{DeligneCarrier.}\allowbreak\lean{g2HilbInt},\allowbreak\ \lean{VOA.}\allowbreak\lean{DeligneAudit.}\allowbreak\lean{DeligneCarrier.}\allowbreak\lean{g2HilbSeriesZ}\par}
{\raggedright\noindent File \leanpath{VOA/\allowbreak DeligneAudit/\allowbreak DeligneFusionUnivGSector.lean}.\par}
\begin{leancode}
def DemazureRealizationG (spec : ℕ → ℕ → RelSpec S.n) (Dem : ℕ → PowerSeries ℚ) (n : ℕ) : Prop :=
  ∃ (V : Type) (_ : AddCommGroup V) (_ : Module ℂ V) (ρ : CurrentModule S.f V) (w : V),
    Sat S.f ρ w (spec n 0) ∧ ∀ N, coeff N (Dem n) ≤ (Module.finrank ℂ (cycDeg ρ w N) : ℚ)
\end{leancode}
{\raggedright\noindent File \leanpath{VOA/\allowbreak DeligneAudit/\allowbreak DeligneCarrierDemazure.lean}.\par}
\begin{leancode}
def hilbSeries (hilb : ℕ → ℕ → ℕ) (n : ℕ) : PowerSeries ℚ :=
  PowerSeries.mk fun N => (hilb n N : ℚ)
\end{leancode}
{\raggedright\noindent File \leanpath{VOA/\allowbreak DeligneAudit/\allowbreak DeligneCharacterG2FusionZ.lean}.\par}
\begin{leancode}
def g2HilbInt (m : ℕ) (N : ℤ) : ℤ := hilbIntG cartanG2 thetaG2 wordG2 m N
\end{leancode}
\begin{leancode}
def g2HilbSeriesZ (n : ℕ) : PowerSeries ℚ := PowerSeries.mk fun N => (g2HilbInt n N : ℚ)
\end{leancode}
\noindent\textbf{Hypotheses / status.} $A_2$, $D_4$: proved; others: \textbf{hypothesis}\par
\medskip
\begin{btheorem}{27.8}[Input R for $D_4$]
For every $n$ there is a current module for the $D_4$ structure constants with a vector satisfying the relations
$\mathrm{rav}(n,2n)$ whose degree-$N$ piece has dimension at least $[q^N]$ of the unclamped $D_4$ Demazure series, for every $N$.
\end{btheorem}
{\raggedright\noindent\textbf{Lean.} \lean{VOA.}\allowbreak\lean{DeligneAudit.}\allowbreak\lean{D4Char.}\allowbreak\lean{d4_demazureRealization}\par}
{\raggedright\noindent File \leanpath{VOA/\allowbreak DeligneAudit/\allowbreak D4CharDemazureRealization.lean}.\par}
\begin{leancode}
theorem d4_demazureRealization :
    ∀ n, DemazureRealizationG D4sector (ravFam D4sector) D4Wp.hilbSeriesZ n :=
\end{leancode}
\noindent\textbf{Hypotheses / status.} Unconditional; aggregate build green 2026-11-03; axioms standard\par
\medskip
\begin{btheorem}{27.13}[$A_2$: classical freeness of $L_{-3/2}(\mathfrak{sl}_3)$]
Let $W=L_{-3/2}(\mathfrak{sl}_3)$ and let $P$ be its explicit current $C_2$ presentation, with generators the eight currents and
ideal $(\ltwentysevenIns)$ generated by the eight quadrics $P_0,\dots,P_7$ of Definition~24.8. Then $W$ is classically
free for $P$, and $\ker(\operatorname{Sym}\mathfrak{sl}_3\to R_W)=(\ltwentysevenIns)$; that is, $R_W=\operatorname{Sym}\mathfrak{sl}_3/(\ltwentysevenIns)$
and $\ltwentysevenJ R_W\to\ltwentysevengr^{\mathrm{Li}}W$ is an isomorphism. The same conclusion is also stated under the (unproved)
hypothesis \lean{RavinderKernelLower}.
\end{btheorem}
{\raggedright\noindent\textbf{Lean.} \lean{VOA.}\allowbreak\lean{DeligneAudit.}\allowbreak\lean{A2.}\allowbreak\lean{classicallyFree},\allowbreak\ \lean{VOA.}\allowbreak\lean{DeligneAudit.}\allowbreak\lean{DeligneSeries.}\allowbreak\lean{A2_classicallyFree_final},\allowbreak\ \lean{VOA.}\allowbreak\lean{DeligneAudit.}\allowbreak\lean{DeligneSeries.}\allowbreak\lean{A2_classicallyFree_final_S2}\par}
{\raggedright\noindent File \leanpath{VOA/\allowbreak DeligneAudit/\allowbreak A2Unconditional.lean}.\par}
\begin{leancode}
theorem classicallyFree :
    CaseClassicallyFree .A2 W explicitPresentation ∧
      RingHom.ker (MvPolynomial.aeval fun i =>
        liDegreeZeroClass W 1 ⟨explicitPresentation.gen i,
          explicitPresentation.gen_weight_one i⟩).toRingHom = I2ns :=
  classicallyFree_of_totalDim' CVUniv.fusionTotalDim
\end{leancode}
{\raggedright\noindent File \leanpath{VOA/\allowbreak DeligneAudit/\allowbreak DeligneSeriesA2.lean}.\par}
\begin{leancode}
theorem A2_classicallyFree_final :
    CaseClassicallyFree .A2 VOA.DeligneAudit.A2.W VOA.DeligneAudit.A2.explicitPresentation ∧
      RingHom.ker (MvPolynomial.aeval fun i =>
        liDegreeZeroClass VOA.DeligneAudit.A2.W 1 ⟨VOA.DeligneAudit.A2.explicitPresentation.gen i,
          VOA.DeligneAudit.A2.explicitPresentation.gen_weight_one i⟩).toRingHom =
        VOA.DeligneAudit.A2.I2ns :=
  VOA.DeligneAudit.A2.classicallyFree

/-- The older A₂ S2 interface, retained for reuse. -/
theorem A2_classicallyFree_final_S2 (hK : VOA.DeligneAudit.A2.RavinderKernelLower) :
    CaseClassicallyFree .A2 VOA.DeligneAudit.A2.W VOA.DeligneAudit.A2.explicitPresentation ∧
      RingHom.ker (MvPolynomial.aeval fun i =>
        liDegreeZeroClass VOA.DeligneAudit.A2.W 1 ⟨VOA.DeligneAudit.A2.explicitPresentation.gen i,
          VOA.DeligneAudit.A2.explicitPresentation.gen_weight_one i⟩).toRingHom =
        VOA.DeligneAudit.A2.I2ns :=
  VOA.DeligneAudit.A2.classicallyFree_of_kernelLower' hK
\end{leancode}
\medskip
\begin{btheorem}{27.14}[$D_4$: the final theorem]
Let $W=L_{-2}(\mathfrak{so}_8)$, built from the certificate-basis Lie data, with its explicit presentation by $105$ quadrics. Suppose
Input T ($\dim\ltwentysevenF_m\le28^m$) and Input R for the unclamped $D_4$ Demazure series. Then $W$ is classically free for its
presentation and $\ker$ is the ideal generated by the $105$ quadrics.
\end{btheorem}
{\raggedright\noindent\textbf{Lean.} \lean{VOA.}\allowbreak\lean{DeligneAudit.}\allowbreak\lean{DeligneSeries.}\allowbreak\lean{D4_classicallyFree_final}\par}
{\raggedright\noindent File \leanpath{VOA/\allowbreak DeligneAudit/\allowbreak DeligneSeries.lean}.\par}
\begin{leancode}
theorem D4_classicallyFree_final (hF : VOA.DeligneAudit.FiniteG.FusionTotalDimG D4.D4sector)
    (hD : ∀ n, DemazureRealizationG D4.D4sector (ravFam D4.D4sector)
      VOA.DeligneAudit.DeligneCharacter.D4Wp.hilbSeriesZ n) :
    CaseClassicallyFree .D4 VOA.DeligneAudit.D4Lie.W4 VOA.DeligneAudit.D4Lie.explicitPresentation4 ∧
      RingHom.ker (MvPolynomial.aeval fun i =>
        liDegreeZeroClass VOA.DeligneAudit.D4Lie.W4 1
          ⟨VOA.DeligneAudit.D4Lie.explicitPresentation4.gen i,
            VOA.DeligneAudit.D4Lie.explicitPresentation4.gen_weight_one i⟩).toRingHom =
        VOA.DeligneAudit.D4Lie.explicitPresentation4.ideal :=
  d4_classicallyFree_of_totalBound (VOA.DeligneAudit.FiniteG.D4.d4_fusionTotalBoundG_of_dim hF) hD
\end{leancode}
\noindent\textbf{Hypotheses / status.} T, R\par
\medskip
\begin{btheorem}{27.15}[$D_4$: classical freeness from Input T alone]
Let $W=L_{-2}(\mathfrak{so}_8)$, built from the certificate-basis Lie data, with its explicit presentation by $105$ quadrics. Suppose
Input T ($\dim\ltwentysevenF_m\le28^m$). Then $W$ is classically free for its presentation and $\ker$ is the ideal generated by the
$105$ quadrics.
\end{btheorem}
{\raggedright\noindent\textbf{Lean.} \lean{VOA.}\allowbreak\lean{DeligneAudit.}\allowbreak\lean{DeligneSeries.}\allowbreak\lean{D4_classicallyFree_of_inputT}\par}
{\raggedright\noindent File \leanpath{VOA/\allowbreak DeligneAudit/\allowbreak DeligneSeriesD4.lean}.\par}
\begin{leancode}
theorem D4_classicallyFree_of_inputT (hF : VOA.DeligneAudit.FiniteG.FusionTotalDimG D4.D4sector) :
    CaseClassicallyFree .D4 VOA.DeligneAudit.D4Lie.W4 VOA.DeligneAudit.D4Lie.explicitPresentation4 ∧
      RingHom.ker (MvPolynomial.aeval fun i =>
        liDegreeZeroClass VOA.DeligneAudit.D4Lie.W4 1
          ⟨VOA.DeligneAudit.D4Lie.explicitPresentation4.gen i,
            VOA.DeligneAudit.D4Lie.explicitPresentation4.gen_weight_one i⟩).toRingHom =
        VOA.DeligneAudit.D4Lie.explicitPresentation4.ideal :=
\end{leancode}
\noindent\textbf{Hypotheses / status.} T only; aggregate build green 2026-11-03; axioms standard\par
\medskip
\begin{btheorem}{27.23}[$G_2$ and $F_4$ concrete headlines]
For $W=L_{-5/3}(\mathfrak g_2)$, classical freeness and $\ker=$ the ideal of the listed quadrics follow from U5 for the clamped fusion series
and Input S; or from U5 and the two halves of the star identity. For $W=L_{-5/2}(\mathfrak f_4)$ they follow from U5 and Input S.
\end{btheorem}
{\raggedright\noindent\textbf{Lean.} \lean{VOA.}\allowbreak\lean{DeligneAudit.}\allowbreak\lean{G2Lie.}\allowbreak\lean{g2_classicallyFree_concrete},\allowbreak\ \lean{VOA.}\allowbreak\lean{DeligneAudit.}\allowbreak\lean{G2Lie.}\allowbreak\lean{g2_classicallyFree_concrete_of_fermionic},\allowbreak\ \lean{VOA.}\allowbreak\lean{DeligneAudit.}\allowbreak\lean{F4Lie.}\allowbreak\lean{f4_classicallyFree_concrete'},\allowbreak\ \lean{VOA.}\allowbreak\lean{DeligneAudit.}\allowbreak\lean{G2Star.}\allowbreak\lean{G2DemazureFermionic},\allowbreak\ \lean{VOA.}\allowbreak\lean{DeligneAudit.}\allowbreak\lean{G2Star.}\allowbreak\lean{G2FermionicKW}\par}
{\raggedright\noindent File \leanpath{VOA/\allowbreak DeligneAudit/\allowbreak G2LieHeadline.lean}.\par}
\begin{leancode}
theorem g2_classicallyFree_concrete (hU5 : FusionModuleBoundG G2sector g2FusionF)
    (hB : G2StarBound) :
    CaseClassicallyFree .G2 g2W explicitPresentation ∧
      RingHom.ker (MvPolynomial.aeval fun i =>
        liDegreeZeroClass g2W 1 ⟨explicitPresentation.gen i,
          explicitPresentation.gen_weight_one i⟩).toRingHom = explicitPresentation.ideal :=
\end{leancode}
\begin{leancode}
theorem g2_classicallyFree_concrete_of_fermionic (hU5 : FusionModuleBoundG G2sector g2FusionF)
    (hF1 : G2Star.G2DemazureFermionic) (hF2 : G2Star.G2FermionicKW) :
    CaseClassicallyFree .G2 g2W explicitPresentation ∧
      RingHom.ker (MvPolynomial.aeval fun i =>
        liDegreeZeroClass g2W 1 ⟨explicitPresentation.gen i,
          explicitPresentation.gen_weight_one i⟩).toRingHom = explicitPresentation.ideal :=
\end{leancode}
{\raggedright\noindent File \leanpath{VOA/\allowbreak DeligneAudit/\allowbreak F4LieHeadline.lean}.\par}
\begin{leancode}
theorem f4_classicallyFree_concrete' (hU5 : FusionModuleBoundG F4sector f4FusionF)
    (hB : F4StarBound) :
    CaseClassicallyFree .F4 f4W explicitPresentation ∧
      RingHom.ker (MvPolynomial.aeval fun i =>
        liDegreeZeroClass f4W 1 ⟨explicitPresentation.gen i,
          explicitPresentation.gen_weight_one i⟩).toRingHom = explicitPresentation.ideal :=
\end{leancode}
{\raggedright\noindent File \leanpath{VOA/\allowbreak DeligneAudit/\allowbreak G2StarFermionic.lean}.\par}
\begin{leancode}
structure G2DemazureFermionic : Prop where
  eq : ∀ n N, g2DemazureHilb n N = g2FermHilb n N

/-- **Named (one `q`-series identity)**: the fermionic Demazure sum is the Kac–Wakimoto form. -/
structure G2FermionicKW : Prop where
  eq : ∀ N, ((demazureSumCoeff g2FermHilb N : ℕ) : ℤ) = coeff N g2KW
\end{leancode}
\noindent\textbf{Hypotheses / status.} U5, star input\par
\medskip
\endgroup
\setcounter{chapter}{27}
\chapter[Graded unitarity through the BRST coset]{Graded unitarity through the BRST coset: from affine $\mathfrak{sl}_2$ at level $-4/3$ to Lee--Yang}
\begingroup
\providecommand{\ltwentyeightsl}{\mathfrak{sl}_2}
\providecommand{\ltwentyeightvac}{\mathbf 1}
\providecommand{\ltwentyeightgr}{\operatorname{gr}}
\providecommand{\ltwentyeightGU}{\mathrm{GU}}
\providecommand{\ltwentyeightQ}{\mathbf Q}
\providecommand{\ltwentyeightLY}{L(-\tfrac{22}{5},0)}
\providecommand{\ltwentyeightA}{L_{-4/3}}
\providecommand{\ltwentyeightinv}{\operatorname{inv}}
\noindent{\small 17 of the chapter's 47 numbered statements: its headline results and every definition they use.}\par\medskip
\begin{bdefinition}{28.4}[The corrected filtration]
A \emph{corrected word} is a finite list of letters $(g,m)$, where $g$ is a current $e,f,h$ or the Sugawara vector $\omega$,
and $m\in\N$ is a mode depth. The letter $(g,m)$ acts by $g_{(-m-1)}$: for a current this is $J^g_{-m-1}$, for $\omega$ it is
$L_{-m-2}$. The \emph{corrected filtration} $F_pV^k$ is the span of the states $g^1_{(-m_1-1)}\cdots g^\ell_{(-m_\ell-1)}\ltwentyeightvac$
with $\ell\le p$ letters, intersected with each conformal weight space. On a quotient $V^k/I$ by a graded ideal the flag is the
image. The flag on $A$ is the image of the corrected flag of $V^{-4/3}$.
\end{bdefinition}
{\raggedright\noindent\textbf{Lean.} \lean{VOA.}\allowbreak\lean{AffineVOA.}\allowbreak\lean{CorrectedR.}\allowbreak\lean{Generator},\allowbreak\ \lean{VOA.}\allowbreak\lean{AffineVOA.}\allowbreak\lean{CorrectedR.}\allowbreak\lean{generatorState},\allowbreak\ \lean{VOA.}\allowbreak\lean{AffineVOA.}\allowbreak\lean{CorrectedR.}\allowbreak\lean{correctedWordState},\allowbreak\ \lean{VOA.}\allowbreak\lean{AffineVOA.}\allowbreak\lean{CorrectedR.}\allowbreak\lean{correctedAffineRPiece},\allowbreak\ \lean{VOA.}\allowbreak\lean{AffineVOA.}\allowbreak\lean{CorrectedR.}\allowbreak\lean{correctedAffineRQuotientFlag},\allowbreak\ \lean{VOA.}\allowbreak\lean{LeeYangTransferAssembly.}\allowbreak\lean{affineFlag}\par}
{\raggedright\noindent File \leanpath{VOA/\allowbreak Examples/\allowbreak AffineSl2GoodFiltration.lean} (a file of Chapter~9).\par}
\begin{leancode}
inductive Generator
  | current (a : Fin 3)
  | sugawara
  deriving DecidableEq

def generatorState (k : ℂ) (hk : k + 2 ≠ 0) : Generator → WakModule k
  | .current a => currentState k a
  | .sugawara => (affineSl2VOA_closed k hk).conformalVector

def correctedWordState (k : ℂ) (hk : k + 2 ≠ 0) :
    List (Generator × ℕ) → WakModule k
  | [] => (affineSl2VOA_closed k hk).vacuum
  | gr :: rest =>
      (affineSl2VOA_closed k hk).smul (generatorState k hk gr.1)
        (-(gr.2 : ℤ) - 1) (correctedWordState k hk rest)

def correctedAffineRPiece (k : ℂ) (hk : k + 2 ≠ 0) (p : ℕ) :
    Submodule ℂ (WakModule k) :=
  Submodule.span ℂ (Set.range (fun w : {w : List (Generator × ℕ) // w.length ≤ p} =>
    correctedWordState k hk w.1))

noncomputable def correctedAffineRQuotientFlag (k : ℂ) (hk : k + 2 ≠ 0)
    (I : Submodule ℂ (WakModule k))
    (hI : (affineSl2VOA_closed k hk).IsIdeal I)
    (hgr : QuotVOA.IsGradedIdeal (affineSl2VOA_closed k hk) I) :
    GradedFlag (QuotVOA.quotientVOA (affineSl2VOA_closed k hk) I hI hgr) :=
  (correctedAffineRFlag k hk).quotient I hI hgr

\end{leancode}
{\raggedright\noindent File \leanpath{VOA/\allowbreak Examples/\allowbreak LeeYangTransferAssemblyAffine.lean}.\par}
\begin{leancode}
def affineFlag : GradedFlag AVOA :=
  CorrectedR.correctedAffineRQuotientFlag ALevel aNoncritical
    (QuotVOA.maxGradedIdeal (universal ALevel aNoncritical))
    (QuotVOA.maxGradedIdeal_isIdeal _)
    (QuotVOA.maxGradedIdeal_isGradedIdeal _)
\end{leancode}
\noindent\textbf{Hypotheses / status.} Definitions\par
\medskip
\begin{bdefinition}{28.5}[Compact conjugation]
For real $k\ne-2$, the \emph{compact conjugation} of $V^k(\ltwentyeightsl)$ is the unique anti-linear map $\rho$ that fixes
$\ltwentyeightvac$, commutes with vertex operators in the sense $\rho(u_{(n)}v)=\rho(u)_{(n)}\rho(v)$, and acts on the currents by
$\rho(e)=-f$, $\rho(f)=-e$, $\rho(h)=-h$; that is, $J^a\mapsto-J^{a^*}$ with $a^*$ the ``adjoint colour''. It is an involution,
and it descends to every quotient by an ideal that it preserves; on $A$ we use the descended map.
\end{bdefinition}
{\raggedright\noindent\textbf{Lean.} \lean{VOA.}\allowbreak\lean{AffineSl2BelowCritical.}\allowbreak\lean{IsCompactRealStructure},\allowbreak\ \lean{VOA.}\allowbreak\lean{AffineSl2BelowCritical.}\allowbreak\lean{compactRealStructure},\allowbreak\ \lean{VOA.}\allowbreak\lean{LeeYangTransferAssembly.}\allowbreak\lean{affineRealStructure},\allowbreak\ \lean{VOA.}\allowbreak\lean{LeeYangTransferAssembly.}\allowbreak\lean{affineRealStructure_current}\par}
{\raggedright\noindent File \leanpath{VOA/\allowbreak Examples/\allowbreak AffineSl2BelowCriticalCompact.lean}.\par}
\begin{leancode}
structure IsCompactRealStructure (R : RealStructure (universal k hk)) : Prop where
  current : ∀ a : Fin 3, R.rho (currentState k a) = -currentState k (adjointColour a)
  parity_ne_zero : ∀ n : ℤ, R.parity n ≠ 0
\end{leancode}
{\raggedright\noindent File \leanpath{VOA/\allowbreak Examples/\allowbreak AffineSl2BelowCriticalRealStructure.lean}.\par}
\begin{leancode}
def compactRealStructure : RealStructure (universal k hk) where
\end{leancode}
{\raggedright\noindent File \leanpath{VOA/\allowbreak Examples/\allowbreak LeeYangTransferAssemblyAffine.lean}.\par}
\begin{leancode}
def affineRealStructure : RealStructure AVOA where

theorem affineRealStructure_current (a : Fin 3) :
    affineRealStructure.rho (simpleCurrent ALevel aNoncritical a) =
      -simpleCurrent ALevel aNoncritical (adjointColour a) := by
\end{leancode}
\noindent\textbf{Hypotheses / status.} Constructed\par
\medskip
\begin{bdefinition}{28.7}[The Lee--Yang datum]
On the Lee--Yang side the form is the canonical normalized form \lean{Bsimp} of $L(c,0)$, the flag is the length flag, and the real
structure is the partition-basis conjugation of Theorem~19.23, at $c=c_{2,5}=1-12\lambda^2$ with the real background
charge $\lambda=\lambda_2$ of Theorem~19.24.
\end{bdefinition}
{\raggedright\noindent\textbf{Lean.} \lean{VOA.}\allowbreak\lean{LeeYangTransferAssembly.}\allowbreak\lean{affineForm},\allowbreak\ \lean{VOA.}\allowbreak\lean{LeeYangTransferAssembly.}\allowbreak\lean{leeYangRealStructure}\par}
{\raggedright\noindent File \leanpath{VOA/\allowbreak Examples/\allowbreak LeeYangTransferAssemblyAffine.lean}.\par}
\begin{leancode}
abbrev affineForm : InvariantForm AVOA := simpleForm ALevel aNoncritical

def leeYangRealStructure : RealStructure (VirSimple.virasoroSimple LeeYang.leeYangC) :=
  virasoroSimpleRealStructure
    ((VirGradedUnitarity.boundaryBackgroundCharge_spec 2).trans
      VirGradedUnitarity.boundaryCharge_two_eq_leeYangC)
    (VirGradedUnitarity.boundaryBackgroundCharge_fixed 2)
\end{leancode}
\noindent\textbf{Hypotheses / status.} Constructed\par
\medskip
\begin{btheorem}{28.8}[Affine graded unitarity implies Lee--Yang graded unitarity]
Suppose that $A=L_{-4/3}(\ltwentyeightsl)$, with its normalized canonical invariant form, the descended compact conjugation
(Definition~28.5) and the corrected flag (Definition~28.4), is graded unitary. Then the
Lee--Yang VOA $\ltwentyeightLY$, with its normalized canonical form, its partition-basis conjugation and its length flag
(Definition~28.7), is graded unitary.
\end{btheorem}
{\raggedright\noindent\textbf{Lean.} \lean{VOA.}\allowbreak\lean{AdamovicSymbol.}\allowbreak\lean{leeYang_gradedUnitary_of_affine}\par}
{\raggedright\noindent File \leanpath{VOA/\allowbreak Examples/\allowbreak AdamovicDescCyclic.lean}.\par}
\begin{leancode}
theorem leeYang_gradedUnitary_of_affine
    (hA : GradedUnitary affineForm affineRealStructure affineFlag) :
    VirGradedUnitarity.VirGradedUnitary LeeYang.leeYangC leeYangRealStructure :=
  leeYang_gradedUnitary_of_affine_wordSymbols_jet hA fun r => Desc.wordSymbolDescends (r + 2)
\end{leancode}
\noindent\textbf{Hypotheses / status.} \textbf{conditional on \lean{hA}} (open)\par
\medskip
\begin{btheorem}{28.11}[Ghost level and nilpotency]
(1) The ghost currents form an affine $\ltwentyeightsl$ action of level $+4=2h^\vee$ on the ghost system.
(2) For a translated affine action of level $k$ on a matter super vertex algebra $V$ with nonzero vacuum (and the parity and grading
conditions needed to form the tensor product), $\ltwentyeightQ^2=0$ if and only if $k=-4$.
\end{btheorem}
{\raggedright\noindent\textbf{Lean.} \lean{VOA.}\allowbreak\lean{Sl2SemiInfinite.}\allowbreak\lean{ghostAffineAction},\allowbreak\ \lean{VOA.}\allowbreak\lean{Sl2SemiInfinite.}\allowbreak\lean{semiInfiniteQ_sq_zero_iff}\par}
{\raggedright\noindent Files \leanpath{VOA/\allowbreak Examples/\allowbreak Sl2SemiInfiniteGhostAffine.lean}, \leanpath{Sl2SemiInfiniteAnomalyCriterion.lean}.\par}
\begin{leancode}
noncomputable def ghostAffineAction : AffineSl2Action sl2GhostSVA 4 :=

theorem semiInfiniteQ_sq_zero_iff (A : AffineSl2Action V k)
    (H : TranslationCovariant V) (hV : TensorVA.ParityCompatible V)
    (hg : TensorVA.TensorGradingHyp V sl2GhostSVA) (hvac : V.vacuum ≠ 0) :
    semiInfiniteQ A * semiInfiniteQ A = 0 ↔ k = -4 := by
\end{leancode}
\noindent\textbf{Hypotheses / status.} Structural hypotheses; proved\par
\medskip
\begin{btheorem}{28.16}[Coset states are relative cocycles; the map to $H^0$]
For a commutant state $x$, the state $x\otimes\ltwentyeightvac_U\otimes\ltwentyeightvac_{\mathrm{gh}}$ is relative and $\ltwentyeightQ$-closed, and
$x\mapsto[x\otimes\ltwentyeightvac\otimes\ltwentyeightvac]$ is a linear map into the ghost-number-zero cohomology which intertwines the vertex
products.
\end{btheorem}
{\raggedright\noindent\textbf{Lean.} \lean{VOA.}\allowbreak\lean{Sl2CosetConcrete.}\allowbreak\lean{cosetToComplex_apply},\allowbreak\ \lean{VOA.}\allowbreak\lean{Sl2CosetConcrete.}\allowbreak\lean{cosetToComplex_closed},\allowbreak\ \lean{VOA.}\allowbreak\lean{Sl2CosetConcrete.}\allowbreak\lean{cosetToComplex_relative},\allowbreak\ \lean{VOA.}\allowbreak\lean{Sl2CosetConcrete.}\allowbreak\lean{cosetToComplex_mode},\allowbreak\ \lean{VOA.}\allowbreak\lean{Sl2CosetConcrete.}\allowbreak\lean{leeYangCosetToH0}\par}
{\raggedright\noindent File \leanpath{VOA/\allowbreak Examples/\allowbreak Sl2CosetConcreteCohomology.lean}.\par}
\begin{leancode}
theorem cosetToComplex_apply (x : CosetSpace) :
    cosetToComplex x = (x.val ⊗ₜ[ℂ] USVA.vacuum) ⊗ₜ[ℂ] sl2GhostSVA.vacuum := rfl

theorem cosetToComplex_closed (x : CosetSpace) : leeYangBRSTComplex.Q (cosetToComplex x) = 0 :=

theorem cosetToComplex_relative (x : CosetSpace) :
    cosetToComplex x ∈ relativeCarrier leeYangMatterAction.toAffineSl2Action :=

theorem cosetToComplex_mode (x y : CosetSpace) (n : ℤ) :
    cosetToComplex (cosetMode diagonalAction.toAffineSl2Action diagonalAction.translation x n y) =
      mode (brstSVA.Y (cosetToComplex x)) n (cosetToComplex y) :=

def leeYangCosetToH0 : CosetSpace →ₗ[ℂ] leeYangBRSTComplex.H0 :=
\end{leancode}
\noindent\textbf{Hypotheses / status.} Unconditional\par
\medskip
\begin{btheorem}{28.17}[The Virasoro image surjects onto Lee--Yang]
Let $\phi:V(-\tfrac{22}{5},0)\to H^0$ be the composite of the VOA map $V(c,0)\to(\text{coset})$ determined by $\omega\mapsto\omega_{\mathrm{coset}}$
with the map of Theorem~28.16. Then $\ker\phi$ is contained in the maximal ideal of $V(c,0)$.
\end{btheorem}
{\raggedright\noindent\textbf{Lean.} \lean{VOA.}\allowbreak\lean{LeeYangShortcut.}\allowbreak\lean{conformalEval_hom},\allowbreak\ \lean{VOA.}\allowbreak\lean{LeeYangShortcut.}\allowbreak\lean{virasoroToCoset},\allowbreak\ \lean{VOA.}\allowbreak\lean{LeeYangShortcut.}\allowbreak\lean{virasoroToH0},\allowbreak\ \lean{virasoroToH0_ker_le_maxIdeal}\par}
{\raggedright\noindent Files \leanpath{VOA/\allowbreak Examples/\allowbreak LeeYangShortcutVirasoro.lean},
\leanpath{LeeYangShortcutImage.lean}.\par}
\begin{leancode}
theorem conformalEval_hom (W : VertexOperatorAlgebra E) {c : ℂ}
    (hc : W.centralCharge = c) :
    IsVOAHom (virasoroVOA c) W (conformalEval W hc) := by

def virasoroToCoset : VirUniv LeeYang.leeYangC →ₗ[ℂ] Sl2CosetConcrete.CosetSpace :=
  conformalEval Sl2CosetConcrete.cosetVOA rfl

def virasoroToH0 : VirUniv LeeYang.leeYangC →ₗ[ℂ] leeYangBRSTComplex.H0 :=
  leeYangCosetToH0.comp virasoroToCoset

theorem virasoroToH0_ker_le_maxIdeal :
    LinearMap.ker virasoroToH0 ≤ maxIdealVir LeeYang.leeYangC :=
\end{leancode}
\noindent\textbf{Hypotheses / status.} Unconditional\par
\medskip
\begin{btheorem}{28.21}[The harmonic bound]
Let $H(Q^-)_{(j,g)}$ be the cohomology of the leading piece $Q^-$ alone in bidegree $(j,g)$. Harmonic vectors of bidegree $(j,g)$ inject into
$H(Q^-)_{(j,g)}$. Consequently, for a finite set $S$ of values of $2R$, the classes of $\ker D/\operatorname{im}D$ with a representative of
ghost number $g$ and $2R\in S$ span a space of dimension at most $\sum_{j\in S}\dim H(Q^-)_{(j,g)}$.
\end{btheorem}
{\raggedright\noindent\textbf{Lean.} \lean{VOA.}\allowbreak\lean{GradedUnitaryCohomology.}\allowbreak\lean{KahlerComplex.}\allowbreak\lean{leadCycles},\allowbreak\ \lean{VOA.}\allowbreak\lean{GradedUnitaryCohomology.}\allowbreak\lean{KahlerComplex.}\allowbreak\lean{leadBoundaries},\allowbreak\ \lean{VOA.}\allowbreak\lean{GradedUnitaryCohomology.}\allowbreak\lean{KahlerComplex.}\allowbreak\lean{leadCohom},\allowbreak\ \lean{VOA.}\allowbreak\lean{GradedUnitaryCohomology.}\allowbreak\lean{KahlerComplex.}\allowbreak\lean{harmDeg},\allowbreak\ \lean{VOA.}\allowbreak\lean{GradedUnitaryCohomology.}\allowbreak\lean{KahlerComplex.}\allowbreak\lean{harmToLead_injective},\allowbreak\ \lean{VOA.}\allowbreak\lean{GradedUnitaryCohomology.}\allowbreak\lean{KahlerComplex.}\allowbreak\lean{cohomFilt},\allowbreak\ \lean{VOA.}\allowbreak\lean{GradedUnitaryCohomology.}\allowbreak\lean{KahlerComplex.}\allowbreak\lean{rowSet},\allowbreak\ \lean{VOA.}\allowbreak\lean{GradedUnitaryCohomology.}\allowbreak\lean{KahlerComplex.}\allowbreak\lean{finrank_cohomFilt_row_le_lead}\par}
{\raggedright\noindent Files \leanpath{VOA/\allowbreak Examples/\allowbreak GradedUnitaryCohomologyKahler.lean},
\leanpath{GradedUnitaryCohomologyFiltered.lean}.\par}
\begin{leancode}
def leadCycles (a : ℤ × ℤ) : Submodule ℂ K := LinearMap.ker C.Qm ⊓ C.deg a

def leadBoundaries (a : ℤ × ℤ) : Submodule ℂ (C.leadCycles a) :=
  (Submodule.map C.Qm (C.deg (a - shQm))).comap (C.leadCycles a).subtype

abbrev leadCohom (a : ℤ × ℤ) := C.leadCycles a ⧸ C.leadBoundaries a

def harmDeg (a : ℤ × ℤ) : Submodule ℂ K := C.harm ⊓ C.deg a

theorem harmToLead_injective (a : ℤ × ℤ) : Function.Injective (C.harmToLead a) := by

def cohomFilt (A : Set (ℤ × ℤ)) : Submodule ℂ (cycleQuotient C.D) :=
  Submodule.map (boundariesInCycles C.D).mkQ
    (Submodule.comap (LinearMap.ker C.D).subtype (C.degOn A))

def rowSet (S : Finset ℤ) (g : ℤ) : Set (ℤ × ℤ) := {a | a.1 ∈ S ∧ a.2 = g}

theorem finrank_cohomFilt_row_le_lead (S : Finset ℤ) (g : ℤ) :
    finrank ℂ (C.cohomFilt (rowSet S g)) ≤ ∑ j ∈ S, finrank ℂ (C.leadCohom (j, g)) := by
\end{leancode}
\noindent\textbf{Hypotheses / status.} Unconditional\par
\medskip
\begin{btheorem}{28.27}[Graded unitarity at one weight from a family of cohomology classes]
Let $C$ be a Kähler complex with a GU datum, $W$ a VOA with a flag $F$ at weight $n$ and forms $\beta$, $h$. Let $f:U\to\ker D/\operatorname{im}D$
and $g:U\to W_n$ be linear, $g$ surjective, $\ker f\subseteq\ker g$, and $P_p\le U$ an exhaustive filtration with $g(P_p)=F_pW_n$ and
$f(P_p)\subseteq F_{p}(\ker D/\operatorname{im}D)$ (classes with representatives of $R\le p$ at ghost number zero). Suppose
$\sum_{j=0}^{2p}\dim H(Q^-)_{(j,0)}=\dim F_pW_n$ for every $p$, that $\beta$ and $h$ pulled back along $g$ agree with $B$ and $h$ on the harmonic
lifts of $f$, and that harmonic vectors of odd $2R$ at ghost number zero vanish. Then $W$ is graded unitary at weight $n$.
\end{btheorem}
{\raggedright\noindent\textbf{Lean.} \lean{VOA.}\allowbreak\lean{LeeYangShortcut.}\allowbreak\lean{gradedUnitaryAt_of_class_family},\allowbreak\ \lean{VOA.}\allowbreak\lean{LeeYangTransferAssembly.}\allowbreak\lean{IntegralHarmonicSupport}\par}
{\raggedright\noindent File \leanpath{VOA/\allowbreak Examples/\allowbreak LeeYangShortcutHodge.lean}.\par}
\begin{leancode}
theorem gradedUnitaryAt_of_class_family {U : Type*} [AddCommGroup U] [Module ℂ U]
    (C : KahlerComplex K) (G : GUData C) (F : GradedFlag W) (n : ℤ)
    (β : LinearMap.BilinForm ℂ (W.grading n)) (h : SesqForm (W.grading n))
    (f : U →ₗ[ℂ] cycleQuotient C.D) (g : U →ₗ[ℂ] W.grading n)
    (hgs : Function.Surjective g) (hker : LinearMap.ker f ≤ LinearMap.ker g)
    (P : ℕ → Submodule ℂ U)
    (hquotient : ∀ p, (P p).map g = F.pieceIn p n)
    (hexhaustive : ∀ x : U, ∃ p, x ∈ P p)
    (hfiltered : ∀ p, (P p).map f ≤
      C.cohomFilt (KahlerComplex.rowSet (Finset.Icc 0 (2 * (p : ℤ))) 0))
    (hcharacter : ∀ p : ℕ,
      (∑ j ∈ Finset.Icc 0 (2 * (p : ℤ)), Module.finrank ℂ (C.leadCohom (j, 0))) =
        Module.finrank ℂ (F.pieceIn p n))
    (hB : ∀ x y, β (g x) (g y) = G.B (harmonicClassLift C f x) (harmonicClassLift C f y))
    (hh : ∀ x y, h (g x) (g y) = G.h (harmonicClassLift C f x) (harmonicClassLift C f y))
    (hintegral : IntegralHarmonicSupport C) : GradedUnitaryAt F n β h := by
\end{leancode}
{\raggedright\noindent File \leanpath{VOA/\allowbreak Examples/\allowbreak LeeYangTransferAssemblyHodge.lean}.\par}
\begin{leancode}
def IntegralHarmonicSupport (C : KahlerComplex K) : Prop :=
  ∀ p : ℕ, C.harmDeg (2 * (p : ℤ) + 1, 0) = ⊥
\end{leancode}
\noindent\textbf{Hypotheses / status.} Explicit hypotheses\par
\medskip
\begin{btheorem}{28.29}[$L_1(\ltwentyeightsl)$ is unitary]
With its normalized canonical form and the descended compact conjugation, $B=L_1(\ltwentyeightsl)$ is unitary. Moreover its conformal vector
is $\omega=\tfrac14 h_{-1}h_{-1}\ltwentyeightvac$.
\end{btheorem}
{\raggedright\noindent\textbf{Lean.} \lean{VOA.}\allowbreak\lean{ConcreteHodge.}\allowbreak\lean{bRealStructure},\allowbreak\ \lean{VOA.}\allowbreak\lean{ConcreteHodge.}\allowbreak\lean{bForm},\allowbreak\ \lean{VOA.}\allowbreak\lean{LevelOneUnitary.}\allowbreak\lean{bUnitary},\allowbreak\ \lean{VOA.}\allowbreak\lean{LevelOneHeis.}\allowbreak\lean{hB},\allowbreak\ \lean{VOA.}\allowbreak\lean{LevelOneHeis.}\allowbreak\lean{bConformal_eq}\par}
{\raggedright\noindent Files \leanpath{VOA/\allowbreak Examples/\allowbreak BGPackageBFactor.lean}, \leanpath{LevelOneUnitaryPositivity.lean},
\leanpath{LevelOneHeisSingular.lean}.\par}
\begin{leancode}
def bRealStructure : RealStructure BVOA where

abbrev bForm : InvariantForm BVOA := simpleForm 1 bNoncritical

theorem bUnitary : ConcreteHodge.Unitary bForm bRealStructure := by

abbrev hB : BSpace := bAction.current 2

theorem bConformal_eq : BVOA.conformalVector = (1 / 4 : ℂ) • BVOA.smul hB (-1) hB := by
\end{leancode}
\noindent\textbf{Hypotheses / status.} \textbf{unconditional}\par
\medskip
\begin{bdefinition}{28.31}[Goodness of the actual differential]
Give each factor its $R$-grading: on $A$ the orthogonal layers of the corrected flag ($R=p$ on layer $p$), on $B$ everything in $R=0$, on $U$ the
layers of the ordinary flag, and on the ghosts a positive grading by weights. The differential is \emph{good} on relative states if it maps
bidegree $(2R,\mathrm{gh})$ into $(2R\pm1,\mathrm{gh}+1)$, with no $R$-preserving component.
\end{bdefinition}
{\raggedright\noindent\textbf{Lean.} \lean{VOA.}\allowbreak\lean{ConcreteHodge.}\allowbreak\lean{aDeg},\allowbreak\ \lean{VOA.}\allowbreak\lean{ConcreteHodge.}\allowbreak\lean{bDeg},\allowbreak\ \lean{VOA.}\allowbreak\lean{ConcreteHodge.}\allowbreak\lean{uDeg},\allowbreak\ \lean{VOA.}\allowbreak\lean{ConcreteHodge.}\allowbreak\lean{matterDeg},\allowbreak\ \lean{VOA.}\allowbreak\lean{ConcreteHodge.}\allowbreak\lean{brstDeg},\allowbreak\ \lean{VOA.}\allowbreak\lean{ConcreteHodge.}\allowbreak\lean{RelativeDifferentialGood}\par}
{\raggedright\noindent File \leanpath{VOA/\allowbreak Examples/\allowbreak ConcreteGoodInterpreted.lean}.\par}
\begin{leancode}
def aDeg : (ℤ × ℤ) × ℂ → Submodule ℂ ASpace := layerDeg affineForm affineFlag

def bDeg : (ℤ × ℤ) × ℂ → Submodule ℂ BSpace := trivialDeg BVOA

def uDeg : (ℤ × ℤ) × ℂ → Submodule ℂ USpace := layerDeg auxiliaryForm auxiliaryFlag

def matterDeg : (ℤ × ℤ) × ℂ → Submodule ℂ MatterSpace := tensorDeg (tensorDeg aDeg bDeg) uDeg

def brstDeg {PG : SesqForm Sl2GhostSpace} (Gh : GhostRGrading PG) :
    (ℤ × ℤ) × ℂ → Submodule ℂ BRSTSpace :=
  tensorDeg matterDeg Gh.deg

def RelativeDifferentialGood {PG : SesqForm Sl2GhostSpace} (Gh : GhostRGrading PG) : Prop :=
  ∀ (n j g : ℤ) (x : BRSTSpace), x ∈ relativeWeight n → x ∈ brstDeg Gh ((j, g), (n : ℂ)) →
    leeYangBRSTComplex.Q x ∈
      brstDeg Gh ((j + 1, g + 1), (n : ℂ)) ⊔ brstDeg Gh ((j - 1, g + 1), (n : ℂ))

\end{leancode}
\noindent\textbf{Hypotheses / status.} \lean{hA}, \lean{hB}\par
\medskip
\begin{btheorem}{28.36}[The Virasoro classes are filtered]
Assume $\ltwentyeightGU(A)$. On the actual complex, at every conformal weight $n$ and for every $p$, the cohomology class of a universal
Virasoro word of length $\le p$ (through the coset stress tensor, Theorem~28.16) has a representative of $2R\in[0,2p]$ at
ghost number zero. This is step (b) of Remark~28.1.
\end{btheorem}
{\raggedright\noindent\textbf{Lean.} \lean{VOA.}\allowbreak\lean{LeeYangShortcut.}\allowbreak\lean{UniversalClassFiltered},\allowbreak\ \lean{VOA.}\allowbreak\lean{ConcreteHodge.}\allowbreak\lean{endpointKahler},\allowbreak\ \lean{VOA.}\allowbreak\lean{ConcreteHodge.}\allowbreak\lean{endpointKahler_filtered}\par}
{\raggedright\noindent Files \leanpath{VOA/\allowbreak Examples/\allowbreak LeeYangShortcutAssembly.lean}, \leanpath{InvariantSymbolBridge.lean}.\par}
\begin{leancode}
def UniversalClassFiltered (S : ConcreteKahlerGeometry) : Prop :=
  ∀ n (hp : IsPosOn (S.geometry n).P ⊤) p,
    (universalPrefix p n).map (modelUniversalClass S n hp) ≤
      (S.model n hp).cohomFilt (KahlerComplex.rowSet (Finset.Icc 0 (2 * (p : ℤ))) 0)

abbrev endpointKahler (hA : GradedUnitary affineForm affineRealStructure affineFlag) :
    ConcreteKahlerGeometry :=
  sp2Kahler hA LevelOneUnitary.bUnitary bgPG_pos bgGh bgGh_isStd
    ((bgGhostRemainderI hA ghostConjRhoD_proved ghostConjBE_proved).pkg
      LevelOneUnitary.bUnitary).sp2

theorem endpointKahler_filtered (hA : GradedUnitary affineForm affineRealStructure affineFlag) :
    UniversalClassFiltered (endpointKahler hA) :=
\end{leancode}
\noindent\textbf{Hypotheses / status.} \lean{hA}\par
\medskip
\begin{bdefinition}{28.39}[The gap-two coefficient and the invariant-symbol statements]
$G_r(n)$ is the coefficient of $z^n$ in $z^{r(r+1)}/\prod_{j=1}^r(1-z^j)$, i.e.\ the number of partitions of $n$ into exactly $r$ parts, all $\ge2$,
with consecutive parts differing by at least $2$. In the Kähler coordinate $j=2R$ put $\mathrm{lc}(n,j)=G_{j/2}(n)$ for even $j\ge0$ and $0$
otherwise. The \emph{invariant-symbol bound} for a concrete Kähler geometry says $\dim H(Q^-)_{(j,0)}\le\mathrm{lc}(n,j)$ at every weight $n$; the
\emph{dimension} statement says equality. The \emph{Lee--Yang PBW character} says $\dim F_pL_n=\sum_{j=0}^{2p}\mathrm{lc}(n,j)$.
\end{bdefinition}
{\raggedright\noindent\textbf{Lean.} \lean{VOA.}\allowbreak\lean{LeeYangTransferAssembly.}\allowbreak\lean{gapTwoCoefficient},\allowbreak\ \lean{VOA.}\allowbreak\lean{LeeYangTransferAssembly.}\allowbreak\lean{leadingCoefficient},\allowbreak\ \lean{VOA.}\allowbreak\lean{LeeYangTransferAssembly.}\allowbreak\lean{LeeYangPBWCharacter},\allowbreak\ \lean{VOA.}\allowbreak\lean{LeeYangTransferAssembly.}\allowbreak\lean{LeadingInvariantSymbolDimension},\allowbreak\ \lean{VOA.}\allowbreak\lean{ConcreteHodge.}\allowbreak\lean{LeadingInvariantSymbolBound}\par}
{\raggedright\noindent File \leanpath{VOA/\allowbreak Examples/\allowbreak LeeYangTransferAssemblyCharacter.lean}.\par}
\begin{leancode}
def gapTwoCoefficient (n : ℤ) (r : ℕ) : ℕ := by
  classical
  exact Fintype.card {a : Fin r → Fin (n.toNat + 1) //
    ((r * (r + 1) + ∑ j : Fin r, (j.val + 1) * (a j).val : ℕ) : ℤ) = n}

def leadingCoefficient (n j : ℤ) : ℕ :=
  if 0 ≤ j ∧ j % 2 = 0 then gapTwoCoefficient n (j.toNat / 2) else 0

def LeeYangPBWCharacter : Prop :=
  ∀ (n : ℤ) (p : ℕ),
    Module.finrank ℂ ((virGradedFlag leeYangC).pieceIn p n) =
      ∑ j ∈ Finset.Icc 0 (2 * (p : ℤ)), leadingCoefficient n j
\end{leancode}
{\raggedright\noindent Files \leanpath{VOA/\allowbreak Examples/\allowbreak LeeYangTransferAssemblyInputs.lean}, \leanpath{InvariantSymbolBridge.lean}.\par}
\begin{leancode}
def LeadingInvariantSymbolDimension (S : ConcreteKahlerGeometry) : Prop :=
  ∀ n (hp : IsPosOn (S.geometry n).P ⊤) j,
    Module.finrank ℂ ((S.model n hp).leadCohom (j, 0)) = leadingCoefficient n j

def LeadingInvariantSymbolBound (S : ConcreteKahlerGeometry) : Prop :=
  ∀ n (hp : IsPosOn (S.geometry n).P ⊤) j,
    Module.finrank ℂ ((S.model n hp).leadCohom (j, 0)) ≤ leadingCoefficient n j
\end{leancode}
\noindent\textbf{Hypotheses / status.} Definitions\par
\medskip
\begin{btheorem}{28.48}[The fermionic jet identity, and strictness]
(1) $\sum_{r\ge0}q^r\sum_{k=0}^{2r}\left[\begin{smallmatrix}2r\\k\end{smallmatrix}\right]_q/(q)_r=\bigl((q^3;q^3)_\infty/(q;q)_\infty\bigr)^3$.
(2) $[q^N]$ of the fermionic series is at most $\sum_{r\le N}\dim A_{r,N}$; hence the jet character bound holds. (3) Given leading-symbol maps in every degree
$r\ge1$ and the jet character bound, strictness holds in every degree $r\ge1$; with Theorem~28.45 this gives strictness outright.
\end{btheorem}
{\raggedright\noindent\textbf{Lean.} \lean{VOA.}\allowbreak\lean{JetFermionic.}\allowbreak\lean{jetFermionicIdentity},\allowbreak\ \lean{VOA.}\allowbreak\lean{JetProduct.}\allowbreak\lean{coeff_jetSeries_le},\allowbreak\ \lean{VOA.}\allowbreak\lean{SymbolModel.}\allowbreak\lean{grJetStrict_of_leadingSymbol_product}\par}
{\raggedright\noindent Files \leanpath{VOA/\allowbreak Examples/\allowbreak JetFermionicIdentity.lean}, \leanpath{JetProductCount.lean},
\leanpath{JetStrictEndpoint.lean}.\par}
\begin{leancode}
theorem jetFermionicIdentity : JetFermionicIdentity := by

theorem coeff_jetSeries_le (N : ℕ) :
    PowerSeries.coeff N jetSeries ≤
      ∑ r ∈ Finset.range (N + 1), (Module.finrank ℂ (jetDeg r (N : ℤ)) : ℚ) := by

def grJetStrict_of_leadingSymbol_product (sym : ∀ r, JetLeadingSymbol (r + 1))
    (chi : JetProductBound) (r : ℕ) : GrJetStrict (r + 1) :=
  grJetStrict_succ_of_leadingSymbol sym (jetCharacterBound_of_product chi) r
\end{leancode}
\noindent\textbf{Hypotheses / status.} \textbf{unconditional}; (3) from inputs\par
\medskip
\begin{btheorem}{28.50}[Gap-two spanning and the upper bound]
In $L$, the length-$\le p$ step of the PBW flag at weight $N$ is spanned by the gap-two words of length $\le p$ (words $L_{-n_1}\cdots L_{-n_q}\ltwentyeightvac$ with
$n_1\ge n_2+2\ge\dots$ and $n_q\ge2$). Hence $\dim F_pL_N\le\sum_{j=0}^{2p}\mathrm{lc}(N,j)$.
\end{btheorem}
{\raggedright\noindent\textbf{Lean.} \lean{VOA.}\allowbreak\lean{LeeYangCharacter.}\allowbreak\lean{GapTwo},\allowbreak\ \lean{VOA.}\allowbreak\lean{LeeYangCharacter.}\allowbreak\lean{gapSpan},\allowbreak\ \lean{VOA.}\allowbreak\lean{LeeYangCharacter.}\allowbreak\lean{goodFil_leeYang_eq_gapSpan},\allowbreak\ \lean{VOA.}\allowbreak\lean{LeeYangCharacter.}\allowbreak\lean{finrank_pieceIn_le}\par}
{\raggedright\noindent Files \leanpath{VOA/\allowbreak Examples/\allowbreak LeeYangCharacterSpanning.lean}, \leanpath{LeeYangCharacterReduction.lean}.\par}
\begin{leancode}
def GapTwo (w : List ℕ) : Prop := w.IsChain (fun a b => b + 2 ≤ a)

noncomputable def gapSpan (ℓ N : ℕ) : Submodule ℂ (VirUniv leeYangC) :=
  Submodule.span ℂ
    {x | ∃ u : List ℕ, GapTwo u ∧ energyVir u = N ∧ u.length ≤ ℓ ∧ negWordVir leeYangC u = x}

theorem goodFil_leeYang_eq_gapSpan (p N : ℕ) :
    goodFil leeYangC p (N : ℤ) = (gapSpan p N).map (maxIdealVir leeYangC).mkQ := by

theorem finrank_pieceIn_le (n : ℤ) (p : ℕ) :
    Module.finrank ℂ ((virGradedFlag leeYangC).pieceIn p n) ≤
      ∑ j ∈ Finset.Icc 0 (2 * (p : ℤ)), leadingCoefficient n j := by
\end{leancode}
\noindent\textbf{Hypotheses / status.} Unconditional\par
\medskip
\begin{btheorem}{28.53}[The Lee--Yang vacuum character]
$\mathrm{rc}(N)\le\operatorname{rank}G_N\le\dim L_N$ for every $N$; hence $\dim L_N=\mathrm{rc}(N)=\sum_rG_r(N)$, and the Rogers--Ramanujan lower bound holds.
\end{btheorem}
{\raggedright\noindent\textbf{Lean.} \lean{VOA.}\allowbreak\lean{RRLower.}\allowbreak\lean{gramRank_le_finrank},\allowbreak\ \lean{VOA.}\allowbreak\lean{RRLower.}\allowbreak\lean{leeYangGramRankBound_holds},\allowbreak\ \lean{VOA.}\allowbreak\lean{RRLower.}\allowbreak\lean{leeYangRochaCaridiCharacter_holds},\allowbreak\ \lean{VOA.}\allowbreak\lean{RRLower.}\allowbreak\lean{leeYangRRLowerBound_holds}\par}
{\raggedright\noindent Files \leanpath{VOA/\allowbreak Examples/\allowbreak RRLowerGram.lean}, \leanpath{LYGramRank.lean}.\par}
{\raggedright\noindent File \leanpath{VOA/\allowbreak Examples/\allowbreak LYGramRankCharacter.lean}.\par}
\begin{leancode}
theorem gramRank_le_finrank (N : ℕ) :
    (vermaGramMatrix leeYangC 0 N).rank ≤
      Module.finrank ℂ ((virasoroSimple leeYangC).grading (N : ℤ)) := by

theorem leeYangGramRankBound_holds : LeeYangGramRankBound := by

theorem leeYangRochaCaridiCharacter_holds : LeeYangRochaCaridiCharacter :=

theorem leeYangRRLowerBound_holds : LeeYangCharacter.LeeYangRRLowerBound :=
\end{leancode}
\noindent\textbf{Hypotheses / status.} \textbf{unconditional}\par
\medskip
\begin{btheorem}{28.54}[The Lee--Yang PBW character, filtration level by filtration level]
For every weight $N$ and every $p$, $\dim F_pL_N=\sum_{s\le p}G_s(N)$; that is, the statement \lean{LeeYangPBWCharacter} of
Definition~28.39 holds. It is equivalent to the Rogers--Ramanujan lower bound.
\end{btheorem}
{\raggedright\noindent\textbf{Lean.} \lean{VOA.}\allowbreak\lean{LeeYangCharacter.}\allowbreak\lean{leeYangPBWCharacter_iff_lowerBound}\par}
{\raggedright\noindent File \leanpath{VOA/\allowbreak Examples/\allowbreak LeeYangCharacterReduction.lean}.\par}
\begin{leancode}
theorem leeYangPBWCharacter_iff_lowerBound : LeeYangPBWCharacter ↔ LeeYangRRLowerBound :=
\end{leancode}
\noindent\textbf{Hypotheses / status.} \textbf{unconditional}\par
\medskip
\endgroup
\setcounter{chapter}{28}
\chapter[The (2,2n+3) series: coset transfer at every boundary level]{The \texorpdfstring{$(2,2n+3)$}{(2,2n+3)} series: coset transfer of graded unitarity at every boundary level}
\begingroup
\providecommand{\ltwentyninesl}{\mathfrak{sl}_2}
\providecommand{\ltwentyninegr}{\operatorname{gr}}
\providecommand{\ltwentyninevac}{|0\rangle}
\providecommand{\ltwentynineGU}{\mathrm{GU}}
\providecommand{\ltwentynineJ}{J_\infty}
\providecommand{\ltwentynineinv}{\mathcal I}
\providecommand{\ltwentyninegc}{g}
\providecommand{\ltwentynineRn}{R_n}
\providecommand{\ltwentynineS}{S_{n,r}}
\providecommand{\ltwentynineR}{\mathbb{R}}
\providecommand{\ltwentynineln}[1]{\nolinkurl{#1}}
\noindent{\small 20 of the chapter's 55 numbered statements: its headline results and every definition they use.}\par\medskip
\begin{bdefinition}{29.3}[Coset level data]
A \emph{coset level datum} consists of complex numbers $k_A$ (the level of the boundary factor), $k_{\mathrm{act}}$
(the index of its current action), $k_U$ (the auxiliary level), a real number $k_U^{\ltwentynineR}$, a target central charge
$c$, and a natural number $m$, such that: $k_{\mathrm{act}}=k_A$; $k_A+2\ne0$; $\overline{k_A}=k_A$;
$k_U=k_U^{\ltwentynineR}$ and $k_U^{\ltwentynineR}<-2$; $k_U+2\ne0$; $k_{\mathrm{act}}+3\ne0$; $k_{\mathrm{act}}+1+k_U=-4$;
$\frac{3k_A}{k_A+2}+1-\frac{3(k_{\mathrm{act}}+1)}{k_{\mathrm{act}}+3}=c$; and $c_{2,2m+1}=c$.
\end{bdefinition}
{\raggedright\noindent\textbf{Lean.} \lean{VOA.}\allowbreak\lean{TwoSeries.}\allowbreak\lean{CosetLevels}\par}
{\raggedright\noindent File \leanpath{VOA/\allowbreak Examples/\allowbreak TwoSeriesLevels.lean}.\par}
\begin{leancode}
structure CosetLevels where
  /-- Level of the boundary affine carrier. -/
  kA : ℂ
  /-- Index of the boundary current action (equal to `kA`). -/
  kAct : ℂ
  /-- Level of the auxiliary universal factor. -/
  kU : ℂ
  /-- The auxiliary level as a real number. -/
  kUr : ℝ
  /-- Target central charge. -/
  c : ℂ
  /-- Index of `c` in the boundary family `c_{2,2m+1}`. -/
  m : ℕ
  kAct_eq : kAct = kA
  aNoncritical : kA + 2 ≠ 0
  aReal : (starRingEnd ℂ) kA = kA
  kU_eq : kU = (kUr : ℂ)
  kUr_lt : kUr < -2
  uNoncritical : kU + 2 ≠ 0
  diagNoncritical : kAct + 1 + 2 ≠ 0
  /-- The total matter level cancels the ghost level `-4`. -/
  sum : kAct + 1 + kU = -4
  /-- Coset central charge `c(k) + c(1) - c(k+1)`. -/
  charge : 3 * kA / (kA + 2) + 1 - 3 * (kAct + 1) / (kAct + 1 + 2) = c
  /-- The target is the `m`th boundary minimal model. -/
  boundary : boundaryCharge m = c
\end{leancode}
\medskip
\begin{btheorem}{29.4}[The coset central charge]
For every $n\in\N$, with $k=k_n=-2+\frac{2}{2n+1}$,
\begin{equation*}
\frac{3k}{k+2}+1-\frac{3(k+1)}{k+3}\;=\;c_{2,2n+3}\;=\;1-\frac{3(2n+1)^2}{2n+3}.
\end{equation*}
\end{btheorem}
{\raggedright\noindent\textbf{Lean.} \lean{VOA.}\allowbreak\lean{TwoSeries.}\allowbreak\lean{twoSeries_charge}\par}
{\raggedright\noindent File \leanpath{VOA/\allowbreak Examples/\allowbreak TwoSeriesLevels.lean}.\par}
\begin{leancode}
theorem twoSeries_charge (n : ℕ) :
    3 * level n / (level n + 2) + 1 - 3 * (level n + 1) / (level n + 1 + 2) =
      boundaryCharge (n + 1) := by
\end{leancode}
\noindent\textbf{Hypotheses / status.} None\par
\medskip
\begin{bdefinition}{29.6}[The data on the boundary factor]
For a level datum $L$, let $A=L_{k_A}(\ltwentyninesl)$ be the simple quotient of $V^{k_A}(\ltwentyninesl)$ by its
maximal graded ideal. Its data are: the canonical invariant form normalised by $(\ltwentyninevac|\ltwentyninevac)=1$;
the \emph{corrected current-and-stress flag} (the PBW-type filtration of Chapter~28 in which a current mode
$J^a_{-j}$ and a Sugawara mode $L_{-j}$ each raise the degree by one, descended to the quotient); and the compact
conjugation (the antilinear involution fixing the compact real form $\mathfrak{su}(2)$), descended from
$V^{k_A}$ to $A$.
\end{bdefinition}
{\raggedright\noindent\textbf{Lean.} \lean{VOA.}\allowbreak\lean{TwoSeries.}\allowbreak\lean{affineForm},\allowbreak\ \lean{VOA.}\allowbreak\lean{TwoSeries.}\allowbreak\lean{affineFlag},\allowbreak\ \lean{VOA.}\allowbreak\lean{TwoSeries.}\allowbreak\lean{affineConj}\par}
{\raggedright\noindent File \leanpath{VOA/\allowbreak Examples/\allowbreak TwoSeriesAffine.lean}.\par}
\begin{leancode}
abbrev affineForm : InvariantForm (AVOA L) := simpleForm L.kA L.aNoncritical

def affineFlag : GradedFlag (AVOA L) :=
  CorrectedR.correctedAffineRQuotientFlag L.kA L.aNoncritical
    (QuotVOA.maxGradedIdeal (universal L.kA L.aNoncritical))
    (QuotVOA.maxGradedIdeal_isIdeal _)
    (QuotVOA.maxGradedIdeal_isGradedIdeal _)

def affineConj : ASpace L →ₛₗ[starRingEnd ℂ] ASpace L :=
  (QuotVOA.maxGradedIdeal (universal L.kA L.aNoncritical)).mapQ
    (QuotVOA.maxGradedIdeal (universal L.kA L.aNoncritical))
    (affineUniversalConj L) (affineUniversalConj_maxIdeal_stable L)
\end{leancode}
\medskip
\begin{btheorem}{29.11}[Coset transfer for any level datum]
Let $L$ be a level datum. If $A=L_{k_A}(\ltwentyninesl)$ is graded unitary, and a Hodge package and a squeeze input
are given, then $L(c,0)$ is graded unitary for its partition-basis real structure. In particular, for the $n$-th
member of the series: $\ltwentynineGU(L_{k_n})\wedge\mathrm{Hodge}\wedge\mathrm{Squeeze}\Rightarrow
\ltwentynineGU(M(2,2n+3))$.
\end{btheorem}
{\raggedright\noindent\textbf{Lean.} \lean{VOA.}\allowbreak\lean{TwoSeries.}\allowbreak\lean{cosetLevels_gradedUnitary_of_affine},\allowbreak\ \lean{VOA.}\allowbreak\lean{TwoSeries.}\allowbreak\lean{twoSeries_gradedUnitary_of_affine}\par}
{\raggedright\noindent File \leanpath{VOA/\allowbreak Examples/\allowbreak TwoSeriesAssembly.lean}.\par}
\begin{leancode}
theorem cosetLevels_gradedUnitary_of_affine
    (hA : GradedUnitary (affineForm L) (affineRealStructure L) (affineFlag L))
    (hodge : BRSTHodgeInputs L) (hsq : SqueezeInput L hodge.kahler) :
    VirGradedUnitary L.c (targetRealStructure L) := by

theorem twoSeries_gradedUnitary_of_affine (n : ℕ)
    (hA : GradedUnitary (affineForm (twoSeriesLevels n)) (affineRealStructure (twoSeriesLevels n))
      (affineFlag (twoSeriesLevels n)))
    (hodge : BRSTHodgeInputs (twoSeriesLevels n))
    (hsq : SqueezeInput (twoSeriesLevels n) hodge.kahler) :
    VirGradedUnitary (twoSeriesLevels n).c (targetRealStructure (twoSeriesLevels n)) :=
  cosetLevels_gradedUnitary_of_affine (twoSeriesLevels n) hA hodge hsq
\end{leancode}
\noindent\textbf{Hypotheses / status.} \lean{hA}, Hodge, squeeze\par
\medskip
\begin{btheorem}{29.12}[The Hodge package exists at every level]
For every level datum $L$ and every proof $h_A$ of $\ltwentynineGU(A)$ there is a BRST/Hodge package. Consequently
$\ltwentynineGU(A)\wedge\mathrm{Squeeze}\Rightarrow\ltwentynineGU(L(c,0))$.
\end{btheorem}
{\raggedright\noindent\textbf{Lean.} \lean{VOA.}\allowbreak\lean{TwoSeries.}\allowbreak\lean{brstHodgeInputs},\allowbreak\ \lean{VOA.}\allowbreak\lean{TwoSeries.}\allowbreak\lean{twoSeries_gradedUnitary_of_affine_hodgeFree}\par}
{\raggedright\noindent File \leanpath{VOA/\allowbreak Examples/\allowbreak TwoSeriesHodgeInputs.lean}.\par}
\begin{leancode}
def brstHodgeInputs
    (hA : GradedUnitary (affineForm L) (affineRealStructure L) (affineFlag L)) :
    BRSTHodgeInputs L where

theorem twoSeries_gradedUnitary_of_affine_hodgeFree (n : ℕ)
    (hA : GradedUnitary (affineForm (twoSeriesLevels n)) (affineRealStructure (twoSeriesLevels n))
      (affineFlag (twoSeriesLevels n)))
    (hsq : SqueezeInput (twoSeriesLevels n) (brstHodgeInputs (twoSeriesLevels n) hA).kahler) :
    VirGradedUnitary (twoSeriesLevels n).c (targetRealStructure (twoSeriesLevels n)) :=
  cosetLevels_gradedUnitary_of_affine_hodgeFree (twoSeriesLevels n) hA hsq
\end{leancode}
\noindent\textbf{Hypotheses / status.} \lean{hA}\par
\medskip
\begin{bdefinition}{29.13}[The affine invariant spaces]
For $r\in\N$ and $N\in\Z$, $\ltwentynineinv_{r,N}$ is the space of tensors in $\ltwentyninegr_rA\otimes B$ of total
conformal weight $N$ annihilated by every diagonal non-negative current mode $J^a_m\otimes1+1\otimes J^a_m$
($a\in\{e,f,h\}$, $m\ge0$). Here $\ltwentyninegr_rA$ is the $r$-th layer of the corrected flag of $A$ with its induced
mode action, and $B=L_1(\ltwentyninesl)$.
\end{bdefinition}
{\raggedright\noindent\textbf{Lean.} \lean{VOA.}\allowbreak\lean{TwoSeries.}\allowbreak\lean{affineInvariants}\par}
{\raggedright\noindent File \leanpath{VOA/\allowbreak Examples/\allowbreak TwoSeriesKoszulInvariants.lean}.\par}
\begin{leancode}
def affineInvariants (r : ℕ) (N : ℤ) : Submodule ℂ (GrA L r ⊗[ℂ] BSpace) :=
  VOA.CyclicEval.tensorWeight (grDeg L r) VOA.CyclicEval.bWeight N ⊓ jointKer (grDiag L r)
\end{leancode}
\medskip
\begin{btheorem}{29.15}[Transfer with an invariant count]
For every $n$: if $\ltwentynineGU(L_{k_n})$ holds and a function $\mathrm{coeff}$ satisfies
$\dim\ltwentynineinv_{r,N}\le\mathrm{coeff}(N,r)$ for all $N,r$ and
$\sum_{r\le p}\mathrm{coeff}(N,r)\le\dim F_pT_N$ for all $N,p$, then $M(2,2n+3)$ is graded unitary.
\end{btheorem}
{\raggedright\noindent\textbf{Lean.} \lean{VOA.}\allowbreak\lean{TwoSeries.}\allowbreak\lean{twoSeries_gradedUnitary_of_affine_invariants}\par}
{\raggedright\noindent File \leanpath{VOA/\allowbreak Examples/\allowbreak TwoSeriesKoszulInvariants.lean}.\par}
\begin{leancode}
theorem twoSeries_gradedUnitary_of_affine_invariants (n : ℕ)
    (hA : GradedUnitary (affineForm (twoSeriesLevels n)) (affineRealStructure (twoSeriesLevels n))
      (affineFlag (twoSeriesLevels n)))
    (coeff : ℤ → ℕ → ℕ)
    (hinv : ∀ (N : ℤ) (r : ℕ),
      Module.rank ℂ (affineInvariants (twoSeriesLevels n) r N) ≤ coeff N r)
    (hpbw : ∀ (N : ℤ) (p : ℕ), ∑ r ∈ Finset.range (p + 1), coeff N r ≤
      Module.finrank ℂ ((virGradedFlag (twoSeriesLevels n).c).pieceIn p N)) :
    VirGradedUnitary (twoSeriesLevels n).c (targetRealStructure (twoSeriesLevels n)) :=
\end{leancode}
\noindent\textbf{Hypotheses / status.} \lean{hA}, \lean{hinv}, \lean{hpbw}\par
\medskip
\begin{bdefinition}{29.16}[Gordon-good words; the Andrews--Gordon count]
A list $w=(w_0,\dots,w_{\ell-1})$ is \emph{Gordon-good of order $n$} if $w_i\ge w_{i+n}+2$ whenever both entries
exist. The \emph{Andrews--Gordon count} $\ltwentyninegc_n(N,r)$ is the number of sorted Gordon-good words of length
exactly $r$ and energy $N$ (and $0$ for $N<0$). Equivalently, it is the number of partitions of $N$ into exactly $r$
parts $\ge2$ in which, writing $f_j$ for the multiplicity of the part $j$, $f_j+f_{j+1}\le n$ for every $j$.
\end{bdefinition}
{\raggedright\noindent\textbf{Lean.} \lean{VOA.}\allowbreak\lean{TwoSeriesCharacter.}\allowbreak\lean{GordonGood},\allowbreak\ \lean{VOA.}\allowbreak\lean{TwoSeriesCharacter.}\allowbreak\lean{tupleList},\allowbreak\ \lean{VOA.}\allowbreak\lean{TwoSeriesCharacter.}\allowbreak\lean{gordonWords},\allowbreak\ \lean{VOA.}\allowbreak\lean{TwoSeriesCharacter.}\allowbreak\lean{gordonCoeff}\par}
{\raggedright\noindent Files \leanpath{VOA/\allowbreak Examples/\allowbreak TwoSeriesCharacterSpanning.lean} (Gordon-good) and
\leanpath{VOA/\allowbreak Examples/\allowbreak TwoSeriesCharacterReduction.lean} (the count).\par}
\begin{leancode}
def GordonGood (n : ℕ) (w : List ℕ) : Prop :=
  ∀ i < w.length - n, w.getD (i + n) 0 + 2 ≤ w.getD i 0

def tupleList {N r : ℕ} (a : Fin r → Fin (N + 1)) : List ℕ := List.ofFn fun j => (a j : ℕ)

def gordonWords (n N r : ℕ) : Finset (List ℕ) :=
  ((Finset.univ : Finset (Fin r → Fin (N + 1))).image tupleList).filter
    (fun u => SortedWord u ∧ GordonGood n u ∧ energyVir u = N)

def gordonCoeff (n : ℕ) (N : ℤ) (r : ℕ) : ℕ :=
  if 0 ≤ N then (gordonWords n N.toNat r).card else 0
\end{leancode}
\medskip
\begin{btheorem}{29.19}[Jet relations of the series]
For every $n\in\N$, $V(c_{2,2n+3},0)$ has jet relations of order $n+1$.
\end{btheorem}
{\raggedright\noindent\textbf{Lean.} \lean{VOA.}\allowbreak\lean{TwoSeries.}\allowbreak\lean{twoSeries_hasJetRelations_all}\par}
{\raggedright\noindent File \leanpath{VOA/\allowbreak Examples/\allowbreak TwoSeriesCharacterEndpoint.lean}.\par}
\begin{leancode}
theorem twoSeries_hasJetRelations_all (n : ℕ) :
    VOA.TwoSeriesCharacter.HasJetRelations (twoSeriesLevels n).c n := by
\end{leancode}
\noindent\textbf{Hypotheses / status.} None\par
\medskip
\begin{btheorem}{29.26}[The vacuum character of $M(2,2n+3)$]
For all $n,N\in\N$,
\begin{equation*}
\dim L(c_{2,2n+3},0)_N\;=\;\sum_{r\le N}\ltwentyninegc_n(N,r)\;=\;\mathrm{rcq}_n(N).
\end{equation*}
In particular the total lower bound $\mathrm{GordonLowerBound}(c_{2,2n+3},n)$ holds.
\end{btheorem}
{\raggedright\noindent\textbf{Lean.} \lean{VOA.}\allowbreak\lean{TwoSeries.}\allowbreak\lean{gordonLowerBound_twoSeries},\allowbreak\ \lean{VOA.}\allowbreak\lean{TwoSeries.}\allowbreak\lean{twoSeries_character}\par}
{\raggedright\noindent File \leanpath{VOA/\allowbreak Examples/\allowbreak TwoSeriesCharacterAGIdentity.lean}.\par}
\begin{leancode}
theorem gordonLowerBound_twoSeries (n : ℕ) :
    VOA.TwoSeriesCharacter.GordonLowerBound (twoSeriesLevels n).c n := by

theorem twoSeries_character (n N : ℕ) :
    Module.finrank ℂ ((VOA.VirSimple.virasoroSimple (twoSeriesLevels n).c).grading (N : ℤ)) =
      ∑ r ∈ Finset.range (N + 1), VOA.TwoSeriesCharacter.gordonCoeff n (N : ℤ) r ∧
    ((∑ r ∈ Finset.range (N + 1), VOA.TwoSeriesCharacter.gordonCoeff n (N : ℤ) r : ℕ) : ℤ) =
      VOA.TwoSeriesCharacter.KacGram.rcq n N :=
\end{leancode}
\noindent\textbf{Hypotheses / status.} None; \textbf{unconditional}\par
\medskip
\begin{btheorem}{29.27}[The refined PBW character of $M(2,2n+3)$]
For every $n\in\N$, $N\in\Z$, $p\in\N$: $\sum_{r\le p}\ltwentyninegc_n(N,r)\le\dim F_pL(c_{2,2n+3},0)_N$. Combined with
Theorem~29.20, the inequality is an equality.
\end{btheorem}
{\raggedright\noindent\textbf{Lean.} \lean{VOA.}\allowbreak\lean{TwoSeries.}\allowbreak\lean{twoSeries_hpbw}\par}
{\raggedright\noindent File \leanpath{VOA/\allowbreak Examples/\allowbreak TwoSeriesCharacterAGIdentity.lean}.\par}
\begin{leancode}
theorem twoSeries_hpbw (n : ℕ) :
    ∀ (N : ℤ) (p : ℕ), ∑ r ∈ Finset.range (p + 1), VOA.TwoSeriesCharacter.gordonCoeff n N r ≤
      Module.finrank ℂ ((virGradedFlag (twoSeriesLevels n).c).pieceIn p N) :=
\end{leancode}
\noindent\textbf{Hypotheses / status.} None\par
\medskip
\begin{btheorem}{29.31}[The socle-word count is the Andrews--Gordon count]
For all $n,N,r\in\N$, $\mathrm{swc}_n(N,r)=\ltwentyninegc_n(N,r)$.
\end{btheorem}
{\raggedright\noindent\textbf{Lean.} \lean{VOA.}\allowbreak\lean{SocleCount.}\allowbreak\lean{socleWordCount_eq}\par}
{\raggedright\noindent File \leanpath{VOA/\allowbreak Examples/\allowbreak TwoSeriesSocleCountGordon.lean}.\par}
\begin{leancode}
theorem socleWordCount_eq (n N r : ℕ) :
    socleWordCount n (N : ℤ) r = gordonCoeff n (N : ℤ) r := by
\end{leancode}
\noindent\textbf{Hypotheses / status.} None\par
\medskip
\begin{btheorem}{29.33}[Transfer from the socle data of $\ltwentyninegr_rA$]
For every $n$: if $L_{k_n}(\ltwentyninesl)$ is graded unitary and $\ltwentyninegr_rA$ (for every $r$, with the socle
words of size $r$ as labels and the Cartan image $I_N\cdot\ltwentyninegr_rA$) has socle data, then $M(2,2n+3)$ is
graded unitary.
\end{btheorem}
{\raggedright\noindent\textbf{Lean.} \lean{VOA.}\allowbreak\lean{TwoSeries.}\allowbreak\lean{TwoSeriesSocleSpectrum},\allowbreak\ \lean{VOA.}\allowbreak\lean{TwoSeries.}\allowbreak\lean{twoSeries_gradedUnitary_of_socleSpectrum}\par}
{\raggedright\noindent File \leanpath{VOA/\allowbreak Examples/\allowbreak TwoSeriesSocleDual.lean}.\par}
\begin{leancode}
structure TwoSeriesSocleSpectrum (n : ℕ) : Prop where
  rank_le : ∀ (r : ℕ) (d : ℤ) (c : ℂ),
    finrank ℂ (sing (grRho (twoSeriesLevels n) r) (grDeg (twoSeriesLevels n) r) d c) ≤
      Nat.card {i : SocleIndex n r //
        (socleDegree i.2.1 : ℤ) = d ∧ c = -(2 * (((i.1 : ℕ) : ℕ) : ℂ))}
  transverse : ∀ (r : ℕ) (d : ℤ) (N : ℕ),
    Disjoint (sing (grRho (twoSeriesLevels n) r) (grDeg (twoSeriesLevels n) r) d
        (-(2 * (N : ℂ))))
      (avoid (grRho (twoSeriesLevels n) r) (grDeg (twoSeriesLevels n) r)
        (grCart (twoSeriesLevels n) r) d N)

theorem twoSeries_gradedUnitary_of_socleSpectrum (n : ℕ)
    (hA : GradedUnitary (affineForm (twoSeriesLevels n)) (affineRealStructure (twoSeriesLevels n))
      (affineFlag (twoSeriesLevels n)))
    (S : TwoSeriesSocleSpectrum n) :
    VirGradedUnitary (twoSeriesLevels n).c (targetRealStructure (twoSeriesLevels n)) :=
  twoSeries_gradedUnitary_of_socleInput n hA S.toSocleInput
\end{leancode}
\noindent\textbf{Hypotheses / status.} \lean{hA}, socle data\par
\medskip
\begin{bdefinition}{29.34}[The coefficient ring $\ltwentynineRn$]
Let $\C[T]=\C[T_0,T_1,\dots]$ and $T(s)=\sum_jT_js^j$. The \emph{jet ideal} of $T^n$ is the ideal generated by the
coefficients of $T(s)^n$, and $\ltwentynineRn=\C[T]/(\text{jet ideal of }T^n)$, with $t_j$ the image of $T_j$.
\end{bdefinition}
{\raggedright\noindent\textbf{Lean.} \lean{VOA.}\allowbreak\lean{TwoSeries.}\allowbreak\lean{StrictModel.}\allowbreak\lean{TPoly},\allowbreak\ \lean{VOA.}\allowbreak\lean{TwoSeries.}\allowbreak\lean{StrictModel.}\allowbreak\lean{tSer},\allowbreak\ \lean{VOA.}\allowbreak\lean{TwoSeries.}\allowbreak\lean{StrictModel.}\allowbreak\lean{jetIdealT},\allowbreak\ \lean{VOA.}\allowbreak\lean{TwoSeries.}\allowbreak\lean{StrictModel.}\allowbreak\lean{Rn},\allowbreak\ \lean{VOA.}\allowbreak\lean{TwoSeries.}\allowbreak\lean{StrictModel.}\allowbreak\lean{tn}\par}
{\raggedright\noindent File \leanpath{VOA/\allowbreak Examples/\allowbreak TwoSeriesStrictEndpoint.lean}.\par}
{\raggedright\noindent File \leanpath{VOA/\allowbreak Examples/\allowbreak TwoSeriesJetRing.lean}.\par}
\begin{leancode}
abbrev TPoly : Type := MvPolynomial ℕ ℂ

def tSer : PowerSeries TPoly := PowerSeries.mk X

def jetIdealT (n : ℕ) : Ideal TPoly :=
  Ideal.span (Set.range fun j => PowerSeries.coeff j (tSer ^ n))

abbrev Rn (n : ℕ) : Type := TPoly ⧸ jetIdealT n

def tn (n : ℕ) (j : ℕ) : Rn n := Ideal.Quotient.mk (jetIdealT n) (X j)
\end{leancode}
\medskip
\begin{bdefinition}{29.36}[The model $\ltwentynineS$]
Let $\ltwentynineRn[u_j,v_j]$ be the polynomial ring in two sequences of jet variables over $\ltwentynineRn$. Each
generator letter of the corrected PBW basis of $V^k(\ltwentyninesl)$ gets a \emph{model symbol}: a current letter
$(J^a,j)$ gets the $j$-th jet of $u^2$, $uv$ or $v^2$ (the Adamovi\'c symbols of Chapter~28); the Sugawara
letter $(L,j)$ gets the $j$-th jet of $-(u'v-uv')$ plus $t_j$. A word's symbol is the product of its letters' symbols.
The model $\ltwentynineS$ is the span of the symbols of words of length $r$; its degree-$d$ piece is the span of symbols
of words of length $r$ and weight $d$.
\end{bdefinition}
{\raggedright\noindent\textbf{Lean.} \lean{VOA.}\allowbreak\lean{TwoSeries.}\allowbreak\lean{StrictModel.}\allowbreak\lean{MPoly},\allowbreak\ \lean{VOA.}\allowbreak\lean{TwoSeries.}\allowbreak\lean{StrictModel.}\allowbreak\lean{mletter},\allowbreak\ \lean{VOA.}\allowbreak\lean{TwoSeries.}\allowbreak\lean{StrictModel.}\allowbreak\lean{mword},\allowbreak\ \lean{VOA.}\allowbreak\lean{TwoSeries.}\allowbreak\lean{StrictModel.}\allowbreak\lean{modelImage},\allowbreak\ \lean{VOA.}\allowbreak\lean{TwoSeries.}\allowbreak\lean{StrictModel.}\allowbreak\lean{modelDeg}\par}
{\raggedright\noindent Files \leanpath{VOA/\allowbreak Examples/\allowbreak TwoSeriesStrictRing.lean} and
\leanpath{VOA/\allowbreak Examples/\allowbreak TwoSeriesStrictWords.lean}.\par}
\begin{leancode}
abbrev MPoly : Type := MvPolynomial (Fin 2 × ℕ) R

def mletter : Generator × ℕ → MPoly R
  | (.current a, n) => ι R (letterSym (.current a, n))
  | (.sugawara, n) => ι R (letterSym (.sugawara, n)) + MvPolynomial.C (t n)

def mword (w : List (Generator × ℕ)) : MPoly R := (w.map (mletter R t)).prod

def modelImage (r : ℕ) : Submodule ℂ (MPoly R) :=
  Submodule.span ℂ {Q | ∃ w : List (Generator × ℕ), w.length = r ∧ mword R t w = Q}

def modelDeg (r : ℕ) (d : ℤ) : Submodule ℂ (MPoly R) :=
  Submodule.span ℂ
    {Q | ∃ w : List (Generator × ℕ), w.length = r ∧ wordWt w = d ∧ mword R t w = Q}
\end{leancode}
\medskip
\begin{btheorem}{29.39}[Transfer from the strict model]
For every $n\ge1$: graded unitarity of $L_{k_n}$, detection, the model character bound and model socle data for every
$r$ imply graded unitarity of $M(2,2n+3)$.
\end{btheorem}
{\raggedright\noindent\textbf{Lean.} \lean{VOA.}\allowbreak\lean{TwoSeries.}\allowbreak\lean{StrictModel.}\allowbreak\lean{twoSeries_gradedUnitary_of_strictModel}\par}
{\raggedright\noindent File \leanpath{VOA/\allowbreak Examples/\allowbreak TwoSeriesStrictEndpoint.lean}.\par}
\begin{leancode}
theorem twoSeries_gradedUnitary_of_strictModel (n : ℕ) (hn : 1 ≤ n)
    (hA : GradedUnitary (affineForm (twoSeriesLevels n)) (affineRealStructure (twoSeriesLevels n))
      (affineFlag (twoSeriesLevels n)))
    (det : TwoSeriesModelDetects n) (chi : TwoSeriesModelCharacter n)
    (spec : ∀ r, ModelSocleSpectrum n r) :
    VirGradedUnitary (twoSeriesLevels n).c (targetRealStructure (twoSeriesLevels n)) :=
  twoSeries_gradedUnitary_of_socleModels n hA
    fun r => twoSeriesSocleModel_of_strict hn det chi r (spec r)
\end{leancode}
\noindent\textbf{Hypotheses / status.} \lean{hA}, det, chi, spec\par
\medskip
\begin{btheorem}{29.41}[The fermionic series is the product; the model character bound]
For every $n\in\N$, $\mathrm{ferm}_n=\bigl((q^{2n+1};q^{2n+1})_\infty/(q;q)_\infty\bigr)^3$. For every $n\ge1$ the
model character bound holds.
\end{btheorem}
{\raggedright\noindent\textbf{Lean.} \lean{VOA.}\allowbreak\lean{TwoSeries.}\allowbreak\lean{ModelChar.}\allowbreak\lean{fermSeries_eq_product},\allowbreak\ \lean{VOA.}\allowbreak\lean{TwoSeries.}\allowbreak\lean{ModelChar.}\allowbreak\lean{twoSeriesModelCharacter_all}\par}
{\raggedright\noindent Files \leanpath{VOA/\allowbreak Examples/\allowbreak TwoSeriesModelCharFermionic.lean} and
\leanpath{VOA/\allowbreak Examples/\allowbreak TwoSeriesModelCharEndpoint.lean}.\par}
\begin{leancode}
theorem fermSeries_eq_product (n : ℕ) :
    fermSeries n = AffineSl2BoundaryCharacter.product n := by

theorem twoSeriesModelCharacter_all (n : ℕ) (hn : 1 ≤ n) : TwoSeriesModelCharacter n := by
  intro N
  rw [← fermSeries_eq_product n]
  exact fermSeries_le_model hn N
\end{leancode}
\noindent\textbf{Hypotheses / status.} None ($n\ge1$)\par
\medskip
\begin{btheorem}{29.51}[The slice form of Arakawa's theorem, unconditionally]
For every level datum $L$ (Definition~29.3) whose affine level $k=k_A$ is not a natural number, the slice map
at $c=c_{\mathrm{DS}}(k)$ sends the maximal ideal of $V^k(\ltwentyninesl)$ into the maximal ideal of $V(c_{\mathrm{DS}}(k),0)$, in every
Rees degree. For an arbitrary non-critical $k\notin\N$ (no level datum) the same holds for the ungraded map $W_1$. In particular, for every $m$, the slice map of
the member $n=m+1$ at $c=c_{2,2m+3}$ has the slice Arakawa property.
\end{btheorem}
{\raggedright\noindent\textbf{Lean.} \lean{VOA.}\allowbreak\lean{TwoSeries.}\allowbreak\lean{SliceMap.}\allowbreak\lean{W1_maxIdeal},\allowbreak\ \lean{VOA.}\allowbreak\lean{TwoSeries.}\allowbreak\lean{SliceMap.}\allowbreak\lean{sliceArakawa_unconditional},\allowbreak\ \lean{VOA.}\allowbreak\lean{TwoSeries.}\allowbreak\lean{SliceMap.}\allowbreak\lean{twoSeries_kA_not_nat},\allowbreak\ \lean{VOA.}\allowbreak\lean{TwoSeries.}\allowbreak\lean{SliceMap.}\allowbreak\lean{twoSeries_sliceArakawa}\par}
{\raggedright\noindent File \leanpath{VOA/\allowbreak Examples/\allowbreak TwoSeriesArakawaUnconditional.lean}.\par}
\begin{leancode}
theorem W1_maxIdeal (hkN : ∀ i : ℕ, k ≠ i) {y : WakModule k} (hy : y ∈ Imax k hk) :
    W1 (cDS k) k hk y ∈ maxIdealVir (cDS k) := by

theorem sliceArakawa_unconditional (L : CosetLevels) (hkN : ∀ i : ℕ, L.kA ≠ i) :
    SliceArakawa (dsSliceMap (cDS L.kA) L) := by

theorem twoSeries_kA_not_nat (m : ℕ) (i : ℕ) : (twoSeriesLevels (m + 1)).kA ≠ i := by

theorem twoSeries_sliceArakawa (m : ℕ) :
    SliceArakawa (dsSliceMap (twoSeriesLevels m).c (twoSeriesLevels (m + 1))) := by
\end{leancode}
\noindent\textbf{Hypotheses / status.} None\par
\medskip
\begin{btheorem}{29.58}[The generation input for every $n\ge2$; the arrangement inclusions for every $n\ge1$]
For every $n\ge2$, $\mathrm{GI}_n$ holds; hence $\mathrm{AI}_n$ holds. For every $n\ge1$, $\mathrm{SAI}_n$ holds.
\end{btheorem}
{\raggedright\noindent\textbf{Lean.} \lean{VOA.}\allowbreak\lean{TwoSeries.}\allowbreak\lean{GenIncl.}\allowbreak\lean{genInclusion},\allowbreak\ \lean{VOA.}\allowbreak\lean{TwoSeries.}\allowbreak\lean{GenIncl.}\allowbreak\lean{arrangementInclusion},\allowbreak\ \lean{VOA.}\allowbreak\lean{TwoSeries.}\allowbreak\lean{GenIncl.}\allowbreak\lean{symArrangementInclusion}\par}
{\raggedright\noindent Files \leanpath{VOA/\allowbreak Examples/\allowbreak TwoSeriesGenIncl.lean} (and
\leanpath{VOA/\allowbreak Examples/\allowbreak TwoSeriesArrangementGen.lean} for the passage $\mathrm{GI}\Rightarrow\mathrm{AI}$).\par}
\begin{leancode}
theorem genInclusion {n : ℕ} (hn : 2 ≤ n) : GenInclusion n := by

theorem arrangementInclusion {n : ℕ} (hn : 2 ≤ n) : ArrangementInclusion n :=
  arrangementInclusion_of_genInclusion hn (genInclusion hn)

theorem symArrangementInclusion {n : ℕ} (hn : 1 ≤ n) :
    ArrangementSym.SymArrangementInclusion n := by
\end{leancode}
\noindent\textbf{Hypotheses / status.} $2\le n$ (resp.\ $1\le n$)\par
\medskip
\begin{btheorem}{29.59}[Graded unitarity transfers along the coset, at every boundary level]
For every $m\in\N$: if $L_{-2+2/(2m+3)}(\ltwentyninesl)$, with its canonical form, compact conjugation and corrected
current-and-stress flag, is graded unitary, then the minimal model $M(2,2m+5)=L(c_{2,2m+5},0)$, with its canonical form,
partition-basis real structure and length flag, is graded unitary.
\end{btheorem}
{\raggedright\noindent\textbf{Lean.} \lean{VOA.}\allowbreak\lean{TwoSeries.}\allowbreak\lean{JetPoincareTwo.}\allowbreak\lean{twoSeries_gradedUnitary_of_affine_final}\par}
{\raggedright\noindent File \leanpath{VOA/\allowbreak Examples/\allowbreak TwoSeriesGenIncl.lean}.\par}
\begin{leancode}
theorem twoSeries_gradedUnitary_of_affine_final (m : ℕ)
    (hA : GradedUnitary (affineForm (twoSeriesLevels (m + 1)))
      (affineRealStructure (twoSeriesLevels (m + 1))) (affineFlag (twoSeriesLevels (m + 1)))) :
    VirGradedUnitary (twoSeriesLevels (m + 1)).c (targetRealStructure (twoSeriesLevels (m + 1))) :=
  twoSeries_gradedUnitary_of_affine_symArrangement m hA
    (GenIncl.symArrangementInclusion (by omega))
\end{leancode}
\noindent\textbf{Hypotheses / status.} \lean{hA} only; conditional\par
\medskip
\endgroup
\part{The Moonshine module}
\chapter{Lattices, codes and the Leech lattice}
\begingroup
\providecommand{\chtwoG}{\mathcal{G}}
\providecommand{\chtwoF}{\mathbb{F}_2}
\providecommand{\chtwowt}{\operatorname{wt}}
\providecommand{\chtworank}{\operatorname{rank}}
\noindent{\small 6 of the chapter's 42 numbered statements: its headline results and every definition they use.}\par\medskip
\begin{btheorem}{30.17}[Index--determinant relation]
Let $f:M\to L$ be an injective isometric map of integral lattices whose image has the same rank
as $L$, and let $k=[L:f(M)]$.  Then $\det M=k^2\det L$.
\end{btheorem}
{\raggedright\noindent\textbf{Lean.} \lean{IntegralLattice.}\allowbreak\lean{det_of_injective_index}\par}
\begin{leancode}
theorem det_of_injective_index (f : Γ →ₗ[ℤ] Λ) (hf : Function.Injective f)
    (hform : ∀ x y : Γ, Λ.bilin (f x) (f y) = Γ.bilin x y)
    (hrank : finrank ℤ (LinearMap.range f) = Λ.rank)
    {k : ℕ} (hk : Nat.card (Λ ⧸ LinearMap.range f) = k) :
    Γ.det = (k : ℤ) ^ 2 * Λ.det
\end{leancode}
\noindent File \leanpath{VOA/\allowbreak LatticesAndCodes/\allowbreak Lattice/\allowbreak Index.lean}.
\noindent\textbf{Hypotheses / status.} Isometric, full rank\par
\medskip
\begin{btheorem}{30.21}[MacWilliams identity]
For a linear code $C\subseteq(\Z/q)^n$,
$|C|\,W_{C^\perp}(X,Y)=W_C\bigl(X+(q-1)Y,\;X-Y\bigr)$.
\end{btheorem}
{\raggedright\noindent\textbf{Lean.} \lean{Code.}\allowbreak\lean{macWilliams_eval},\allowbreak\ \lean{Code.}\allowbreak\lean{macWilliams}\par}
\begin{leancode}
theorem macWilliams_eval (C : Code q n) (A B : ℂ) :
    (C.toFinset.card : ℂ) * eval ![A, B] C.dual.hammingEnumerator
      = eval ![A + ((q : ℂ) - 1) * B, A - B] C.hammingEnumerator

theorem macWilliams (C : Code q n) :
    (C.toFinset.card : ℂ) • C.dual.hammingEnumerator
      = bind₁ ![X 0 + MvPolynomial.C ((q : ℂ) - 1) * X 1, X 0 - X 1] C.hammingEnumerator
\end{leancode}
\noindent File \leanpath{VOA/\allowbreak LatticesAndCodes/\allowbreak Codes/\allowbreak MacWilliams.lean}.
\noindent\textbf{Hypotheses / status.} $q\ge1$\par
\medskip
\begin{btheorem}{30.24}[{$\chtwoG$ is a doubly even self-dual $[24,12,8]$ code}]
$\dim\chtwoG=12$, $\chtwoG$ is self-dual and doubly even, and its minimum distance is $8$.
\end{btheorem}
{\raggedright\noindent\textbf{Lean.} \lean{Code.}\allowbreak\lean{golay_rank},\allowbreak\ \lean{Code.}\allowbreak\lean{golay_isSelfDual},\allowbreak\ \lean{Code.}\allowbreak\lean{golay_isDoublyEven},\allowbreak\ \lean{Code.}\allowbreak\lean{golay_minDist}\par}
\begin{leancode}
theorem golay_rank : golay.rank = 12

theorem golay_isSelfDual : golay.IsSelfDual

theorem golay_isDoublyEven : golay.IsDoublyEven

theorem golay_minDist : golay.minDist = 8
\end{leancode}
\noindent File \leanpath{VOA/\allowbreak LatticesAndCodes/\allowbreak Codes/\allowbreak Golay.lean}.
\noindent\textbf{Hypotheses / status.} None\par
\medskip
\begin{btheorem}{30.25}[Weight enumerator of $\chtwoG$]
$W_{\chtwoG}(X,Y)=X^{24}+759X^{16}Y^8+2576X^{12}Y^{12}+759X^8Y^{16}+Y^{24}$.  In particular
$\chtwoG$ has $759$ codewords of weight $8$, $2576$ of weight $12$ and $759$ of weight $16$.
\end{btheorem}
{\raggedright\noindent\textbf{Lean.} \lean{Code.}\allowbreak\lean{golay_weightCount_eight},\allowbreak\ \lean{Code.}\allowbreak\lean{golay_weightCount_twelve},\allowbreak\ \lean{Code.}\allowbreak\lean{golay_weightCount_sixteen_eq},\allowbreak\ \lean{Code.}\allowbreak\lean{golay_hammingEnumerator}\par}
\begin{leancode}
theorem golay_weightCount_eight : golay.weightCount 8 = 759

theorem golay_weightCount_twelve : golay.weightCount 12 = 2576

theorem golay_weightCount_sixteen_eq : golay.weightCount 16 = 759

theorem golay_hammingEnumerator :
    golay.hammingEnumerator
      = X 0 ^ 24 + MvPolynomial.C 759 * (X 0 ^ 16 * X 1 ^ 8)
        + MvPolynomial.C 2576 * (X 0 ^ 12 * X 1 ^ 12)
        + MvPolynomial.C 759 * (X 0 ^ 8 * X 1 ^ 16) + X 1 ^ 24
\end{leancode}
\noindent File \leanpath{VOA/\allowbreak LatticesAndCodes/\allowbreak Codes/\allowbreak GolayEnumerator.lean}.
\noindent\textbf{Hypotheses / status.} None\par
\medskip
\begin{btheorem}{30.29}[Properties of $\Gamma_C$]
Let $C\subseteq\chtwoF^n$ be self-orthogonal of dimension $k$.  Then $\Gamma_C$ is positive
definite of rank $n$; $\det\Gamma_C=2^{n-2k}$; $\Gamma_C$ is even iff $C$ is doubly even; and
$\Gamma_C$ is unimodular iff $C$ is self-dual.
\end{btheorem}
{\raggedright\noindent\textbf{Lean.} \lean{latticeOfCode_isPosDef},\allowbreak\ \lean{latticeOfCode_rank},\allowbreak\ \lean{latticeOfCode_isEven},\allowbreak\ \lean{latticeOfCode_det},\allowbreak\ \lean{latticeOfCode_isUnimodular}\par}
\begin{leancode}
theorem latticeOfCode_isPosDef (C : Code 2 n) (hC : C ≤ C.dual) :
    (latticeOfCode C hC).IsPosDef

@[simp] theorem latticeOfCode_rank (C : Code 2 n) (hC : C ≤ C.dual) :
    (latticeOfCode C hC).rank = n

theorem latticeOfCode_isEven (C : Code 2 n) (hC : C ≤ C.dual) :
    (latticeOfCode C hC).IsEven ↔ C.IsDoublyEven

theorem latticeOfCode_det (C : Code 2 n) (hC : C ≤ C.dual) :
    (latticeOfCode C hC).det = 2 ^ (n - 2 * C.rank)

theorem latticeOfCode_isUnimodular (C : Code 2 n) (hC : C ≤ C.dual) :
    (latticeOfCode C hC).IsUnimodular ↔ C.IsSelfDual
\end{leancode}
\noindent File \leanpath{VOA/\allowbreak LatticesAndCodes/\allowbreak Construction/\allowbreak CodeToLattice.lean}.
\noindent\textbf{Hypotheses / status.} $C\subseteq C^\perp$\par
\medskip
\begin{btheorem}{30.42}[Rank divisible by $8$; Ebeling, Thm.~2.1]
If $L$ is even, unimodular and positive definite, then $8\mid\operatorname{rank}L$, and
$\vartheta_L$ is a modular form of weight $\operatorname{rank}L/2$ for $SL_2(\Z)$.
\end{btheorem}
{\raggedright\noindent\textbf{Lean.} \lean{IntegralLattice.}\allowbreak\lean{eight_dvd_rank},\allowbreak\ \lean{IntegralLattice.}\allowbreak\lean{thetaModularForm},\allowbreak\ \lean{IntegralLattice.}\allowbreak\lean{theta_isModularForm}\par}
{\raggedright\noindent File \leanpath{VOA/\allowbreak LatticesAndCodes/\allowbreak Theta/\allowbreak ModularForm.lean}.\par}
\begin{leancode}
theorem eight_dvd_rank (h : Γ.IsPosDef) (he : Γ.IsEven) (hu : Γ.IsUnimodular) : 8 ∣ Γ.rank

def thetaModularForm (Γ : IntegralLattice) (h : Γ.IsPosDef) (he : Γ.IsEven)
    (hu : Γ.IsUnimodular) : ModularForm 𝒮ℒ ((Γ.rank / 2 : ℕ) : ℤ) where
  toFun := Γ.thetaFun

theorem theta_isModularForm (h : Γ.IsPosDef) (he : Γ.IsEven) (hu : Γ.IsUnimodular) :
    8 ∣ Γ.rank ∧
      ∃ f : ModularForm 𝒮ℒ ((Γ.rank / 2 : ℕ) : ℤ), ∀ z : ℍ, f z = Γ.theta (z : ℂ)
\end{leancode}
\noindent\textbf{Hypotheses / status.} Even, unimod., pos.\ definition\par
\medskip
\endgroup
\chapter{Heisenberg and lattice vertex algebras}
\begingroup
\providecommand{\chthreeV}{V}
\providecommand{\chthreeF}{\mathcal{F}}
\providecommand{\chthreeeps}{\varepsilon}
\providecommand{\chthreecoc}{\operatorname{coc}}
\providecommand{\chthreeR}{\mathbb{R}}
\noindent{\small 7 of the chapter's 35 numbered statements: its headline results and every definition they use.}\par\medskip
\begin{btheorem}{31.3}[Heisenberg relations]
For all colours $a,b$ and all $m,n\in\Z$,
$[h_a(m),h_b(n)] = m\,k\,\delta_{ab}\,\delta_{m+n,0}\cdot\mathrm{id}$ on $\mathcal F_r$.
\end{btheorem}
{\raggedright\noindent\textbf{Lean.} \lean{VOA.}\allowbreak\lean{HeisenbergN.}\allowbreak\lean{oscN_comm}\par}
{\raggedright\noindent File \leanpath{VOA/\allowbreak Examples/\allowbreak HeisenbergNFock.lean}.\par}
\begin{leancode}
theorem oscN_comm (r : ℕ) (k : ℂ) (a b : Fin r) (m n : ℤ) :
    oscN r k a m * oscN r k b n - oscN r k b n * oscN r k a m
      = (if a = b ∧ m + n = 0 then (m : ℂ) * k else 0) • LinearMap.id
\end{leancode}
\noindent\textbf{Hypotheses / status.} Unconditional\par
\medskip
\begin{btheorem}{31.11}[Properties of $\varepsilon_{\mathcal B}$]
For every basis $\mathcal B$ and all $\alpha,\alpha',\beta,\beta',\gamma\in L$:
\begin{enumerate}[label=(\roman*)]
\item $\varepsilon(\alpha+\alpha',\beta)=\varepsilon(\alpha,\beta)\varepsilon(\alpha',\beta)$
  and $\varepsilon(\alpha,\beta+\beta')=\varepsilon(\alpha,\beta)\varepsilon(\alpha,\beta')$;
\item $\varepsilon(0,\beta)=\varepsilon(\alpha,0)=1$ and $\varepsilon(\alpha,\beta)^2=1$;
\item the 2-cocycle law (31.8);
\item $\chthreecoc(\alpha,\beta)+\chthreecoc(\beta,\alpha)=\langle\alpha,\beta\rangle
  -\sum_i\langle B_i,B_i\rangle\alpha_i\beta_i$;
\item if $L$ is even, the commutator property (31.9).
\end{enumerate}
\end{btheorem}
{\raggedright\noindent\textbf{Lean.} \lean{VOA.}\allowbreak\lean{IntegralLattice.}\allowbreak\lean{cocycle_add_left},\allowbreak\ \lean{VOA.}\allowbreak\lean{IntegralLattice.}\allowbreak\lean{cocycle_add_right},\allowbreak\ \lean{VOA.}\allowbreak\lean{IntegralLattice.}\allowbreak\lean{cocycle_zero_left},\allowbreak\ \lean{VOA.}\allowbreak\lean{IntegralLattice.}\allowbreak\lean{cocycle_zero_right},\allowbreak\ \lean{VOA.}\allowbreak\lean{IntegralLattice.}\allowbreak\lean{cocycle_sq},\allowbreak\ \lean{VOA.}\allowbreak\lean{IntegralLattice.}\allowbreak\lean{cocycle_cocycleLaw},\allowbreak\ \lean{VOA.}\allowbreak\lean{IntegralLattice.}\allowbreak\lean{cocExp_add_swap},\allowbreak\ \lean{VOA.}\allowbreak\lean{IntegralLattice.}\allowbreak\lean{cocycle_comm}\par}
\begin{leancode}
theorem cocycle_add_left (α α' β : L.carrier) :
    cocycle L B (α + α') β = cocycle L B α β * cocycle L B α' β
theorem cocycle_add_right (α β β' : L.carrier) :
    cocycle L B α (β + β') = cocycle L B α β * cocycle L B α β'
theorem cocycle_zero_left (β : L.carrier) : cocycle L B 0 β = 1
theorem cocycle_zero_right (α : L.carrier) : cocycle L B α 0 = 1
theorem cocycle_sq (α β : L.carrier) : cocycle L B α β * cocycle L B α β = 1
theorem cocycle_cocycleLaw (α β γ : L.carrier) :
    cocycle L B α β * cocycle L B (α + β) γ = cocycle L B β γ * cocycle L B α (β + γ)
theorem cocExp_add_swap (α β : L.carrier) :
    cocExp L B α β + cocExp L B β α
      = L.form α β - ∑ i : Fin n, L.form (B i) (B i) * B.repr α i * B.repr β i
theorem cocycle_comm (hEven : L.Even) (α β : L.carrier) :
    cocycle L B α β * cocycle L B β α = (-1 : ℂ) ^ (L.form α β)
\end{leancode}
\noindent File \leanpath{VOA/\allowbreak Foundations/\allowbreak LatticeCocycle.lean}.
\noindent\textbf{Hypotheses / status.} Unconditional\par
\medskip
\begin{bdefinition}{31.18}[Grading of $V_L$]
For $\alpha\in L$ let $|\alpha|^2/2:=\langle\alpha,\alpha\rangle/2$ (an integer when $L$ is even).
For $(w,\alpha)\in\Z\times L$ let $V_{(w,\alpha)}=(\mathcal F)_w\otimes e^\alpha$ (zero if $w<0$),
and for $n\in\Z$ let $(V_L)_n=\sum_\alpha V_{(n-|\alpha|^2/2,\alpha)}$.
\end{bdefinition}
{\raggedright\noindent\textbf{Lean.} \lean{VOA.}\allowbreak\lean{LatGen.}\allowbreak\lean{halfNorm},\allowbreak\ \lean{VOA.}\allowbreak\lean{LatGen.}\allowbreak\lean{bigradα},\allowbreak\ \lean{VOA.}\allowbreak\lean{LatGen.}\allowbreak\lean{gradingLatGen}\par}
\begin{leancode}
noncomputable def halfNorm (α : L.carrier) : ℤ := L.form α α / 2

noncomputable def bigradα (wα : ℤ × L.carrier) : Submodule ℂ (LatModuleGen L) :=
  (HeisenbergN.gradingZN (rank L) wα.1).map (mkEα L wα.2)

noncomputable def gradingLatGen (k : ℤ) : Submodule ℂ (LatModuleGen L) :=
  ⨆ α : L.carrier, bigradα L (k - halfNorm L α, α)
\end{leancode}
\noindent File \leanpath{VOA/\allowbreak Examples/\allowbreak LatticeGenGrading.lean}.
\medskip
\begin{btheorem}{31.26}[The reconstructed vertex map]
Let $L$ be even and positive-definite. There is a linear map $Y\colon V_L\to\End(V_L)[[z^{\pm1}]]$
which equals (31.13) on every word state and is a local, $T$-covariant field map
with the creation property.
\end{btheorem}
{\raggedright\noindent\textbf{Lean.} \lean{VOA.}\allowbreak\lean{LatGen.}\allowbreak\lean{latYfullGen_of_locality},\allowbreak\ \lean{VOA.}\allowbreak\lean{LatGen.}\allowbreak\lean{YfullGen}\par}
\begin{leancode}
theorem latYfullGen_of_locality (hEven : L.Even) (hPos : L.PosDef) :
    ∃ Yfull : LatModuleGen L →ₗ[ℂ] FormalDist (End ℂ (LatModuleGen L)),
      (∀ (lα : List (ℕ × Fin (rank L)) × L.carrier),
        Yfull (Function.uncurry (wordStateLatGen L) lα)
          = Function.uncurry (fun l α => YwordLatGen L l α) lα) ∧
      (∀ u : LatModuleGen L, IsField (Yfull u)) ∧
      (∀ u v : LatModuleGen L, Local (Yfull u) (Yfull v)) ∧
      (∀ u : LatModuleGen L, (Yfull u) 0 (latVacuumGen L) = u) ∧
      (∀ u : LatModuleGen L, ∀ q : ℤ, q < 0 → (Yfull u) q (latVacuumGen L) = 0) ∧
      (∀ u : LatModuleGen L, Reconstruction.DCovariant (TrecGen L) (Yfull u))

noncomputable def YfullGen (hEven : L.Even) (hPos : L.PosDef) :
    LatModuleGen L →ₗ[ℂ] FormalDist (End ℂ (LatModuleGen L)) :=
  (latYfullGen_of_locality L hEven hPos).choose
\end{leancode}
\noindent Files \leanpath{VOA/\allowbreak Examples/\allowbreak LatticeGenWordFields.lean}, \leanpath{VOA/\allowbreak Examples/\allowbreak LatticeGenVOAInstance.lean}.
\noindent\textbf{Hypotheses / status.} E, P\par
\medskip
\begin{btheorem}{31.27}[The lattice vertex operator algebra]
For every even positive-definite lattice $L$ of rank $r$, $V_L$ with the vertex map of
Theorem~31.26, vacuum $1\otimes e^0$, conformal vector
$\omega=\frac12\sum_ah_a(-1)^2\mathbb 1$ and the grading of Definition~31.18
is a vertex operator algebra of CFT type with central charge $r$.
\end{btheorem}
{\raggedright\noindent\textbf{Lean.} \lean{VOA.}\allowbreak\lean{LatGen.}\allowbreak\lean{latticeVOAGen},\allowbreak\ \lean{VOA.}\allowbreak\lean{LatGen.}\allowbreak\lean{latticeVOAGen_centralCharge}\par}
\begin{leancode}
noncomputable def latticeVOAGen (hEven : L.Even) (hPos : L.PosDef) :
    VertexOperatorAlgebra (LatModuleGen L) where
  grading := gradingLatGen L
  isInternal := gradingLatGen_isInternal_of_posDef L hEven hPos
  finrank := gradingLatGen_finrank_of_posDef L hEven hPos
  Y := YfullGen L hEven hPos
  vacuum := latVacuumGen L
  vacuum_mem := by
    rw [gradingLatGen_zero_eq_span_vacuum L hEven hPos]
    exact Submodule.mem_span_singleton_self (latVacuumGen L)
  grading_mode := fun u hu v hv j => YfullGen_grading_mode L hEven hPos u hu v hv j
  vacuum_Y := YfullGen_vacuum L hEven hPos
  creation := YfullGen_creation L hEven hPos
  locality := YfullGen_local L hEven hPos
  regularity := YfullGen_regularity L hEven hPos
  conformalVector := omegaLatGen L
  conformal_mem := omegaLatGen_mem L
  centralCharge := (rank L : ℂ)
  virasoro := YfullGen_virasoro L hEven hPos
  L0_grading := fun v hv => YfullGen_L0_grading L hEven hPos v hv
  translation := YfullGen_translation_omega L hEven hPos
  vacuum_grade := gradingLatGen_zero_eq_span_vacuum L hEven hPos
  positivity := fun n hn => gradingLatGen_positivity L hEven hPos n hn

theorem latticeVOAGen_centralCharge (hEven : L.Even) (hPos : L.PosDef) :
    (latticeVOAGen L hEven hPos).centralCharge = (rank L : ℂ) := rfl
\end{leancode}
\noindent File \leanpath{VOA/\allowbreak Examples/\allowbreak LatticeGenInstance.lean}.
\noindent\textbf{Hypotheses / status.} E, P\par
\medskip
\begin{btheorem}{31.39}[Functoriality]
Let $K,L$ be even positive-definite lattices and $\iota\colon K\to L$ an isometry. Then
\begin{enumerate}[label=(\roman*)]
\item $\Phi_\iota$ is a homomorphism of vertex algebras $V_K\to V_L$ (preserves $\mathbb 1$ and all
  products $u_{(n)}v$);
\item if $\operatorname{rank}K=\operatorname{rank}L$, $\Phi_\iota(\omega_K)=\omega_L$ and $\Phi_\iota$
  is injective, with image the set of $x\in V_L$ all of whose momenta lie in $\iota(K)$;
\item $\Phi_\iota\theta_K=\theta_L\Phi_\iota$;
\item $\Phi_\iota g_{\chi\circ\iota}=g_\chi\Phi_\iota$ for every $\chi\colon L\to\Z/2$.
\end{enumerate}
\end{btheorem}
{\raggedright\noindent\textbf{Lean.} \lean{VOA.}\allowbreak\lean{LatGen.}\allowbreak\lean{latMorph_isVOAHom},\allowbreak\ \lean{VOA.}\allowbreak\lean{LatGen.}\allowbreak\lean{latMorph_mode},\allowbreak\ \lean{VOA.}\allowbreak\lean{LatGen.}\allowbreak\lean{latMorph_omega},\allowbreak\ \lean{VOA.}\allowbreak\lean{LatGen.}\allowbreak\lean{latMorph_injective},\allowbreak\ \lean{VOA.}\allowbreak\lean{LatGen.}\allowbreak\lean{mem_range_latMorph_iff},\allowbreak\ \lean{VOA.}\allowbreak\lean{LatGen.}\allowbreak\lean{latMorph_thetaLat},\allowbreak\ \lean{VOA.}\allowbreak\lean{LatGen.}\allowbreak\lean{latMorph_chiLat}\par}
\begin{leancode}
theorem latMorph_isVOAHom :
    IsVOAHom (latticeVOAGen K hKe hK) (latticeVOAGen L hLe hL) (latMorph K L ι)
theorem latMorph_mode (u : LatModuleGen K) (n : ℤ) (v : LatModuleGen K) :
    latMorph K L ι (mode (YfullGen K hKe hK u) n v)
      = mode (YfullGen L hLe hL (latMorph K L ι u)) n (latMorph K L ι v)
theorem latMorph_omega : latMorph K L ι (omegaLatGen K) = omegaLatGen L
theorem latMorph_injective : Function.Injective (latMorph K L ι)
theorem mem_range_latMorph_iff (x : LatModuleGen L) :
    x ∈ LinearMap.range (latMorph K L ι) ↔ ∀ γ, γ ∉ Set.range ι → coordEα L γ x = 0
theorem latMorph_thetaLat (x : LatModuleGen K) :
    latMorph K L ι (thetaLat K x) = thetaLat L (latMorph K L ι x)
theorem latMorph_chiLat (χ : L.carrier →+ ZMod 2) (x : LatModuleGen K) :
    latMorph K L ι (chiLat K (χ.comp ι.toAddMonoidHom) x) = chiLat L χ (latMorph K L ι x)
\end{leancode}
\noindent File \leanpath{VOA/\allowbreak Moonshine/\allowbreak LatticeMorph.lean}.
\noindent\textbf{Hypotheses / status.} E, P (both), I\par
\medskip
\begin{btheorem}{31.42}[Rigidity]
Let $L$ be even positive-definite and $\psi\in\Aut(V_L)$ such that $\psi(\alpha(-1)\mathbb 1)=
-\alpha(-1)\mathbb 1$ for all $\alpha\in L$ (equivalently $\psi(h_a(-1)\mathbb 1)=-h_a(-1)\mathbb1$
for all $a$) and $\psi(\mathcal F\otimes e^\gamma)\subseteq\mathcal F\otimes e^{-\gamma}$ for all
$\gamma$. Then $\psi=g_d\theta$ for a character $d\colon L\to\C^\times$, and $\psi=g_e\theta
g_e^{-1}$ for a character $e$; in particular $\psi$ is conjugate to $\theta$ in $\Aut(V_L)$.
\end{btheorem}
{\raggedright\noindent\textbf{Lean.} \lean{VOA.}\allowbreak\lean{LatGen.}\allowbreak\lean{exists_char_of_lift},\allowbreak\ \lean{VOA.}\allowbreak\lean{LatGen.}\allowbreak\lean{lift_eq_conj_theta},\allowbreak\ \lean{VOA.}\allowbreak\lean{LatGen.}\allowbreak\lean{lift_eq_conj_theta'},\allowbreak\ \lean{VOA.}\allowbreak\lean{Aut.}\allowbreak\lean{pullback},\allowbreak\ \lean{VOA.}\allowbreak\lean{Aut.}\allowbreak\lean{pullback_apply}\par}
\begin{leancode}
theorem exists_char_of_lift (ψ : Aut (latticeVOAGen L hEven hPos))
    (hh : ∀ a, ψ (heisSt L a) = -heisSt L a)
    (hsec : ∀ (γ : L.carrier) (x : LatModuleGen L), x ∈ chargeSec L γ → ψ x ∈ chargeSec L (-γ)) :
    ∃ d : L.carrier →+ Additive ℂˣ, ∀ x, ψ x = charLat L d (thetaLat L x)

theorem lift_eq_conj_theta (ψ : Aut (latticeVOAGen L hEven hPos))
    (hh : ∀ a, ψ (heisSt L a) = -heisSt L a)
    (hsec : ∀ (γ : L.carrier) (x : LatModuleGen L), x ∈ chargeSec L γ → ψ x ∈ chargeSec L (-γ)) :
    ∃ e : L.carrier →+ Additive ℂˣ, ∀ x, ψ x = charLat L e (thetaLat L (charLat L (-e) x))

theorem lift_eq_conj_theta' (ψ : Aut (latticeVOAGen L hEven hPos))
    (hh : ∀ α, ψ (hSt L α) = -hSt L α)
    (hsec : ∀ (γ : L.carrier) (x : LatModuleGen L), x ∈ chargeSec L γ → ψ x ∈ chargeSec L (-γ)) :
    ∃ e : L.carrier →+ Additive ℂˣ, ∀ x, ψ x = charLat L e (thetaLat L (charLat L (-e) x))
\end{leancode}
\begin{leancode}
noncomputable def Aut.pullback (hΦ : IsVOAHom V W Φ)
    (hω : Φ V.conformalVector = W.conformalVector) : Aut V
theorem Aut.pullback_apply (hΦ : IsVOAHom V W Φ)
    (hω : Φ V.conformalVector = W.conformalVector) (x : 𝒱) :
    Φ (Aut.pullback Φ hinj g hg hΦ hω x) = g (Φ x)
\end{leancode}
\noindent File \leanpath{VOA/\allowbreak Moonshine/\allowbreak LatticeMorphRigid.lean}.
\noindent\textbf{Hypotheses / status.} E, P, two hyps on $\psi$\par
\medskip
\endgroup
\chapter{The involution \texorpdfstring{$\theta$}{theta} and the fixed-point algebra \texorpdfstring{$V_\Lambda^+$}{V+}}
\begingroup
\providecommand{\chfourone}{\mathbf{1}}
\noindent{\small 6 of the chapter's 37 numbered statements: its headline results and every definition they use.}\par\medskip
\begin{bproposition}{32.9}
For every integral lattice $L$ and all $\alpha,\beta\in L$:
$\varepsilon(-\alpha,\beta)=\varepsilon(\alpha,\beta)=\varepsilon(\alpha,-\beta)$; hence
$\varepsilon(-\alpha,-\beta)=\varepsilon(\alpha,\beta)$.
\end{bproposition}
{\raggedright\noindent\textbf{Lean.} \lean{VOA.}\allowbreak\lean{LatGen.}\allowbreak\lean{latticeCocycle_neg_left},\allowbreak\ \lean{VOA.}\allowbreak\lean{LatGen.}\allowbreak\lean{latticeCocycle_neg_right},\allowbreak\ \lean{VOA.}\allowbreak\lean{LatGen.}\allowbreak\lean{latticeCocycle_neg_neg}\par}
\begin{leancode}
-- VOA/Moonshine/LatticeInvolution.lean
theorem latticeCocycle_neg_left (α β : L.carrier) :
    latticeCocycle L (-α) β = latticeCocycle L α β
theorem latticeCocycle_neg_right (α β : L.carrier) :
    latticeCocycle L α (-β) = latticeCocycle L α β
theorem latticeCocycle_neg_neg (α β : L.carrier) :
    latticeCocycle L (-α) (-β) = latticeCocycle L α β
\end{leancode}
\noindent File \leanpath{VOA/\allowbreak Moonshine/\allowbreak LatticeInvolution.lean}.
\noindent\textbf{Hypotheses / status.} None\par
\medskip
\begin{btheorem}{32.13}[$\theta$ intertwines the vertex operators]
Let $L$ be even and positive-definite. Then for all $u,v\in V_L$ and $n\in\Z$,
\begin{equation*}
  \theta\big(u_{(n)}v\big)=(\theta u)_{(n)}(\theta v),
\end{equation*}
equivalently (32.2).
\end{btheorem}
{\raggedright\noindent\textbf{Lean.} \lean{VOA.}\allowbreak\lean{LatGen.}\allowbreak\lean{thetaLat_YfullGen},\allowbreak\ \lean{VOA.}\allowbreak\lean{LatGen.}\allowbreak\lean{thetaLat_mode}\par}
\begin{leancode}
-- VOA/Moonshine/LatticeInvolution.lean
theorem thetaLat_YfullGen (u : LatModuleGen L) (t : ℤ) (x : LatModuleGen L) :
    thetaLat L (YfullGen L hEven hPos u t x)
      = YfullGen L hEven hPos (thetaLat L u) t (thetaLat L x)

theorem thetaLat_mode (u : LatModuleGen L) (n : ℤ) (v : LatModuleGen L) :
    thetaLat L (mode (YfullGen L hEven hPos u) n v)
      = mode (YfullGen L hEven hPos (thetaLat L u)) n (thetaLat L v)
\end{leancode}
\noindent File \leanpath{VOA/\allowbreak Moonshine/\allowbreak LatticeInvolution.lean}.
\noindent\textbf{Hypotheses / status.} $L$ even, pos.-definition\par
\medskip
\begin{btheorem}{32.29}[Fock generating functions]
In the ring $\C[[q]]$ of formal power series,
\begin{equation*}
  \sum_{w\ge0}P(r,w)q^w=\prod_{n\ge1}(1-q^n)^{-r},
  \qquad
  \sum_{w\ge0}\big(P^+(r,w)-P^-(r,w)\big)q^w=\prod_{n\ge1}(1+q^n)^{-r},
\end{equation*}
the infinite products converging in the $q$-adic (coefficientwise) topology.
\end{btheorem}
{\raggedright\noindent\textbf{Lean.} \lean{VOA.}\allowbreak\lean{LatGen.}\allowbreak\lean{fockGenFun},\allowbreak\ \lean{VOA.}\allowbreak\lean{LatGen.}\allowbreak\lean{fockSignedGenFun},\allowbreak\ \lean{VOA.}\allowbreak\lean{LatGen.}\allowbreak\lean{hasProd_fockGenFun},\allowbreak\ \lean{VOA.}\allowbreak\lean{LatGen.}\allowbreak\lean{hasProd_fockSignedGenFun}\par}
\begin{leancode}
-- VOA/Moonshine/LatticeGradedDim.lean
noncomputable def fockGenFun (r : ℕ) : PowerSeries ℂ := PowerSeries.mk fun w => (fockCount r w : ℂ)
noncomputable def fockSignedGenFun (r : ℕ) : PowerSeries ℂ :=
  PowerSeries.mk fun w => (fockSigned r w : ℂ)

open scoped PowerSeries.WithPiTopology in
theorem hasProd_fockGenFun (r : ℕ) :
    HasProd (fun i : ℕ => ((1 - (PowerSeries.X : ℂ⟦X⟧) ^ (i + 1))⁻¹) ^ r) (fockGenFun r)

open scoped PowerSeries.WithPiTopology in
theorem hasProd_fockSignedGenFun (r : ℕ) :
    HasProd (fun i : ℕ => ((1 + (PowerSeries.X : ℂ⟦X⟧) ^ (i + 1))⁻¹) ^ r)
      (fockSignedGenFun r)
\end{leancode}
\noindent File \leanpath{VOA/\allowbreak Moonshine/\allowbreak LatticeGradedDim.lean}.
\noindent\textbf{Hypotheses / status.} None\par
\medskip
\begin{btheorem}{32.31}[Graded dimension of $V_L$]
Let $L$ be even and positive-definite of rank $r$. For every $k\in\N$,
\begin{equation*}
  \dim(V_L)_k=\sum_{n=0}^{k}\#\{\alpha\in L:\langle\alpha,\alpha\rangle=2n\}\cdot P(r,k-n),
\end{equation*}
the coefficient of $q^k$ in $\Theta_L(q)\prod_{n\ge1}(1-q^n)^{-r}$, where
$\Theta_L(q)=\sum_{\alpha\in L}q^{\langle\alpha,\alpha\rangle/2}$ is the theta series of $L$.
\end{btheorem}
{\raggedright\noindent\textbf{Lean.} \lean{VOA.}\allowbreak\lean{LatGen.}\allowbreak\lean{normEq},\allowbreak\ \lean{VOA.}\allowbreak\lean{LatGen.}\allowbreak\lean{mem_normEq},\allowbreak\ \lean{VOA.}\allowbreak\lean{LatGen.}\allowbreak\lean{finrank_gradingLatGen_eq_sum_theta}\par}
\begin{leancode}
-- VOA/Moonshine/LatticeGradedDim.lean
noncomputable def normEq (m : ℤ) : Finset L.carrier := (normEq_finite L hPos m).toFinset
theorem mem_normEq {m : ℤ} {α : L.carrier} : α ∈ normEq L hPos m ↔ L.form α α = m

theorem finrank_gradingLatGen_eq_sum_theta (hEven : L.Even) (k : ℕ) :
    finrank ℂ (gradingLatGen L k)
      = ∑ n ∈ Finset.range (k + 1), (normEq L hPos (2 * n)).card * fockCount (rank L) (k - n)
\end{leancode}
\noindent File \leanpath{VOA/\allowbreak Moonshine/\allowbreak LatticeGradedDim.lean}.
\noindent\textbf{Hypotheses / status.} Even, pos.-definition\par
\medskip
\begin{btheorem}{32.36}[Generation]
Let $S\subseteq V_\Leech^+$ be a linear subspace such that $\chfourone\in S$, $u_{(n)}v\in S$ for all
$u,v\in S$ and $n\in\Z$, and $x_\beta\in S$ for every $\beta\in\Leech$. Then $S=V_\Leech^+$.
\end{btheorem}
{\raggedright\noindent\textbf{Lean.} \lean{VOA.}\allowbreak\lean{Moonshine.}\allowbreak\lean{leechPlus_generated}\par}
\begin{leancode}
-- VOA/Moonshine/Generation.lean
theorem leechPlus_generated : ∀ S : Submodule ℂ VPlus, leechVOAPlus.vacuum ∈ S →
    (∀ u ∈ S, ∀ v ∈ S, ∀ n : ℤ, leechVOAPlus.smul u n v ∈ S) →
    (∀ β : leech.carrier, frameX0P β ∈ S) → S = ⊤
\end{leancode}
\noindent File \leanpath{VOA/\allowbreak Moonshine/\allowbreak Generation.lean}.
\noindent\textbf{Hypotheses / status.} None\par
\medskip
\begin{btheorem}{32.45}[Properties of the lattice form]
Let $L$ be even. Then:
\begin{enumerate}[label=(\alph*)]
\item (invariance at $e^\gamma$) for all $\gamma\in L$, $n\in\Z$ and $x,y\in V_L$, with
  $N=\tfrac12\langle\gamma,\gamma\rangle$,
  \begin{equation*}B(e^\gamma_{(n)}x,y)=(-1)^NB(x,e^\gamma_{(2N-n-2)}y);\end{equation*}
\item (nondegeneracy) if $x\neq0$ there is $y$ with $B(x,y)\neq0$;
\item ($\theta$-invariance) $B(\theta x,\theta y)=B(x,y)$;
\item (orthogonality) distinct weight spaces of $V_L$ are $B$-orthogonal.
\end{enumerate}
\end{btheorem}
{\raggedright\noindent\textbf{Lean.} \lean{VOA.}\allowbreak\lean{Moonshine.}\allowbreak\lean{Nondeg.}\allowbreak\lean{latForm_Yeαvec},\allowbreak\ \lean{VOA.}\allowbreak\lean{Moonshine.}\allowbreak\lean{Nondeg.}\allowbreak\lean{latForm_nondeg},\allowbreak\ \lean{VOA.}\allowbreak\lean{Moonshine.}\allowbreak\lean{Nondeg.}\allowbreak\lean{latForm_thetaLat},\allowbreak\ \lean{VOA.}\allowbreak\lean{Moonshine.}\allowbreak\lean{Nondeg.}\allowbreak\lean{latForm_orth_grading}\par}
\begin{leancode}
-- VOA/Moonshine/NondegLattice.lean
theorem latForm_Yeαvec (hEven : L.Even) (γ : L.carrier) (n : ℤ) (x y : LatModuleGen L) :
    latForm L (mode (Yeαvec L γ) n x) y
      = (-1 : ℂ) ^ (LatGen.halfNorm L γ) *
          latForm L x (mode (Yeαvec L γ) (2 * LatGen.halfNorm L γ - n - 2) y)
theorem latForm_nondeg (x : LatModuleGen L) (hx : x ≠ 0) : ∃ y, latForm L x y ≠ 0
theorem latForm_thetaLat (x y : LatModuleGen L) :
    latForm L (thetaLat L x) (thetaLat L y) = latForm L x y
theorem latForm_orth_grading {m n : ℤ} (hmn : m ≠ n) {x y : LatModuleGen L}
    (hx : x ∈ gradingLatGen L m) (hy : y ∈ gradingLatGen L n) : latForm L x y = 0
\end{leancode}
\noindent File \leanpath{VOA/\allowbreak Moonshine/\allowbreak NondegLattice.lean}.
\noindent\textbf{Hypotheses / status.} $L$ even\par
\medskip
\endgroup
\chapter{Graded dimensions and the modular function \texorpdfstring{$J$}{J}}
\begingroup
\providecommand{\chfiveH}{\mathfrak{H}}
\providecommand{\chfiveSL}{\mathrm{SL}_2(\Z)}
\noindent{\small 8 of the chapter's 44 numbered statements: its headline results and every definition they use.}\par\medskip
\begin{btheorem}{33.5}[$j$ is a Hauptmodul; argued on paper only]
Every $\chfiveSL$-invariant holomorphic function on $\chfiveH$ with at most a pole of order $m$ at the
cusp is a polynomial of degree $\le m$ in $j$. In particular $J$ is the unique invariant function
of the form $q^{-1} + O(q)$.
\end{btheorem}
\noindent\emph{(The chapter quotes no Lean declaration for this statement; see book p.~1740.)}\par
\medskip
\begin{bdefinition}{33.9}[Formal $\Delta$]
$\Delta = (E_4^3 - E_6^2)/1728\in\C[[q]]$.
\end{bdefinition}
{\raggedright\noindent\textbf{Lean.} \lean{VOA.}\allowbreak\lean{deltaSeries}\par}
\begin{leancode}
noncomputable def deltaSeries : PowerSeries ℂ :=
  (1728 : ℂ)⁻¹ • (E4Series ^ 3 - E6Series ^ 2)
\end{leancode}
\noindent File \leanpath{VOA/\allowbreak Modularity/\allowbreak JFunction.lean}.
\medskip
\begin{btheorem}{33.18}[Generating function of the twisted Fock space]
With $x = q^{1/2}$, in $\C[[x]]$,
\begin{equation*}
  \sum_{d\ge0}P_T(r,d)\,x^d = \prod_{i\ge0}\bigl(1 - x^{2i+1}\bigr)^{-r}
  \quad\Bigl(= \prod_{n\ge1}\bigl(1-q^{n-\frac12}\bigr)^{-r}\Bigr).
\end{equation*}
\end{btheorem}
{\raggedright\noindent\textbf{Lean.} \lean{VOA.}\allowbreak\lean{TwistedSector.}\allowbreak\lean{twGenFun},\allowbreak\ \lean{VOA.}\allowbreak\lean{TwistedSector.}\allowbreak\lean{hasProd_twGenFun}\par}
\begin{leancode}
noncomputable def twGenFun (r : ℕ) : PowerSeries ℂ := PowerSeries.mk fun d => (twCount r d : ℂ)

theorem hasProd_twGenFun (r : ℕ) :
    HasProd (fun i : ℕ => ((1 - (PowerSeries.X : ℂ⟦X⟧) ^ (2 * i + 1))⁻¹) ^ r) (twGenFun r)
\end{leancode}
\noindent File \leanpath{VOA/\allowbreak Moonshine/\allowbreak TwistedSectorDim.lean}.
\noindent\textbf{Hypotheses / status.} None\par
\medskip
\begin{btheorem}{33.21}[Rank $24$: fixed $=$ integral weight, nothing below $2$]
For $r = 24$ and any $T$: the fixed part is exactly the sum of the integer-weight eigenspaces; there
are no fixed vectors of weight $<2$; and $\dim(S[\theta]\otimes T)^+_{n+2} = P_T(24,2n+1)\dim T$ for all
$n\ge0$.
\end{btheorem}
{\raggedright\noindent\textbf{Lean.} \lean{VOA.}\allowbreak\lean{TwistedSector.}\allowbreak\lean{twFixed_eq_iSup_int},\allowbreak\ \lean{VOA.}\allowbreak\lean{TwistedSector.}\allowbreak\lean{twFixed_inf_twWeight_of_lt_two},\allowbreak\ \lean{VOA.}\allowbreak\lean{TwistedSector.}\allowbreak\lean{finrank_twFixed_inf_twWeight_24}\par}
{\raggedright\noindent File \leanpath{VOA/\allowbreak Moonshine/\allowbreak TwistedSectorDim.lean}.\par}
\begin{leancode}
theorem twFixed_eq_iSup_int : twFixed 24 T = ⨆ n : ℤ, twWeight 24 T n

theorem twFixed_inf_twWeight_of_lt_two {q : ℚ} (hq : q < 2) :
    twFixed 24 T ⊓ twWeight 24 T q = ⊥

theorem finrank_twFixed_inf_twWeight_24 (n : ℕ) :
    finrank ℂ (twFixed 24 T ⊓ twWeight 24 T (((n : ℤ) + 2 : ℤ) : ℚ) :
        Submodule ℂ (FockT 24 ⊗[ℂ] T))
      = twCount 24 (2 * n + 1) * finrank ℂ T
\end{leancode}
\noindent\textbf{Hypotheses / status.} None\par
\medskip
\begin{btheorem}{33.23}[Twisted dimensions in weights $\le3$]
$(V_\Leech^T)^+_0 = (V_\Leech^T)^+_1 = 0$, $\dim(V_\Leech^T)^+_2 = 24\cdot4096 = 98304$,
$\dim(V_\Leech^T)^+_3 = 2624\cdot4096 = 10747904$. Consequently
\begin{gather*}
  \dim(V_\Leech^+)_2 + \dim(V_\Leech^T)^+_2 = 98580 + 98304 = 196884,\\
  \dim(V_\Leech^+)_3 + \dim(V_\Leech^T)^+_3 = 10745856 + 10747904 = 21493760 .
\end{gather*}
\end{btheorem}
{\raggedright\noindent\textbf{Lean.} \lean{VOA.}\allowbreak\lean{Moonshine.}\allowbreak\lean{leechTwisted_finrank_fixed_two},\allowbreak\ \lean{VOA.}\allowbreak\lean{Moonshine.}\allowbreak\lean{leechTwisted_finrank_fixed_three},\allowbreak\ \lean{VOA.}\allowbreak\lean{Moonshine.}\allowbreak\lean{leechTwisted_fixed_weight_zero_one},\allowbreak\ \lean{VOA.}\allowbreak\lean{Moonshine.}\allowbreak\lean{finrank_sectors_two},\allowbreak\ \lean{VOA.}\allowbreak\lean{Moonshine.}\allowbreak\lean{finrank_sectors_three}\par}
\begin{leancode}
theorem leechTwisted_finrank_fixed_two :
    finrank ℂ (twFixed 24 leechT ⊓ twWeight 24 leechT 2 : Submodule ℂ leechTwistedSpace)
      = 98304

theorem leechTwisted_finrank_fixed_three :
    finrank ℂ (twFixed 24 leechT ⊓ twWeight 24 leechT 3 : Submodule ℂ leechTwistedSpace)
      = 10747904

theorem leechTwisted_fixed_weight_zero_one :
    twFixed 24 leechT ⊓ twWeight 24 leechT 0 = ⊥ ∧ twFixed 24 leechT ⊓ twWeight 24 leechT 1 = ⊥

theorem finrank_sectors_two :
    finrank ℂ (leechVOAPlus.grading 2)
      + finrank ℂ (twFixed 24 leechT ⊓ twWeight 24 leechT 2 : Submodule ℂ leechTwistedSpace)
      = 196884

theorem finrank_sectors_three :
    finrank ℂ (leechVOAPlus.grading 3)
      + finrank ℂ (twFixed 24 leechT ⊓ twWeight 24 leechT 3 : Submodule ℂ leechTwistedSpace)
      = 21493760
\end{leancode}
\noindent File \leanpath{VOA/\allowbreak Moonshine/\allowbreak TwistedSectorDim.lean}.
\noindent\textbf{Hypotheses / status.} None\par
\medskip
\begin{btheorem}{33.35}[Jacobi's product, formally]
$\Delta = q\,D$ in $\C[[q]]$, for the formal $\Delta = (E_4^3-E_6^2)/1728$; equivalently $\Delta/q = D$, and
(A): $(\Delta/q)\cdot\prod(1-q^n)^{-24} = 1$.
\end{btheorem}
{\raggedright\noindent\textbf{Lean.} \lean{VOA.}\allowbreak\lean{Moonshine.}\allowbreak\lean{deltaSeries_eq_X_mul_deltaProd},\allowbreak\ \lean{VOA.}\allowbreak\lean{Moonshine.}\allowbreak\lean{deltaOverQ_eq_deltaProd},\allowbreak\ \lean{VOA.}\allowbreak\lean{Moonshine.}\allowbreak\lean{deltaProductIdentity}\par}
\begin{leancode}
theorem deltaSeries_eq_X_mul_deltaProd : deltaSeries = X * deltaProd

theorem deltaOverQ_eq_deltaProd : deltaOverQ = deltaProd

theorem deltaProductIdentity : DeltaProductIdentity
\end{leancode}
\noindent File \leanpath{VOA/\allowbreak Moonshine/\allowbreak CharacterDelta.lean}.
\noindent\textbf{Hypotheses / status.} None\par
\medskip
\begin{btheorem}{33.36}[Weight-$12$ forms are determined by two coefficients]
For every $f\in M_{12}$ with $q$-expansion $\sum a_nq^n$:
$\ \sum a_nq^n = a_0E_4^3 + \frac{a_1 - 720a_0}{1728}(E_4^3 - E_6^2)$ in $\C[[q]]$.
\end{btheorem}
{\raggedright\noindent\textbf{Lean.} \lean{VOA.}\allowbreak\lean{Moonshine.}\allowbreak\lean{qExpansion_weight_twelve}\par}
\begin{leancode}
theorem qExpansion_weight_twelve (f : ModularForm 𝒮ℒ ((12 : ℕ) : ℤ)) :
    qExpansion 1 (f : ℍ → ℂ)
      = (qExpansion 1 (f : ℍ → ℂ)).coeff 0 • E4Series ^ 3
        + (((qExpansion 1 (f : ℍ → ℂ)).coeff 1 - 720 * (qExpansion 1 (f : ℍ → ℂ)).coeff 0) / 1728)
          • (E4Series ^ 3 - E6Series ^ 2)
\end{leancode}
\noindent File \leanpath{VOA/\allowbreak Moonshine/\allowbreak CharacterDelta.lean}.
\noindent\textbf{Hypotheses / status.} $f\in M_{12}$\par
\medskip
\begin{btheorem}{33.47}[$\chi = qJ$]
In $\C[[q]]$,
\begin{equation*}
  \sum_{n\ge0}\Bigl(\dim(V_\Leech^+)_n + \dim(V_\Leech^T)^+_n\Bigr)q^n = q\,J(q),
\end{equation*}
equivalently, for every $n\in\N$, $\dim(V_\Leech^+)_n + \dim(V_\Leech^T)^+_n = [q^n]\,qJ$.
\end{btheorem}
{\raggedright\noindent\textbf{Lean.} \lean{VOA.}\allowbreak\lean{Moonshine.}\allowbreak\lean{sectorSeries_eq_JShift},\allowbreak\ \lean{VOA.}\allowbreak\lean{Moonshine.}\allowbreak\lean{finrank_sectors_eq_coeff_JShift}\par}
\begin{leancode}
theorem sectorSeries_eq_JShift : sectorSeries = JShift :=
  sectorSeries_eq_JShift_of_levelTwo levelTwoIdentity

theorem finrank_sectors_eq_coeff_JShift (n : ℕ) :
    ((finrank ℂ (leechVOAPlus.grading (n : ℤ))
      + finrank ℂ (twFixed 24 leechT ⊓ twWeight 24 leechT ((n : ℤ) : ℚ) :
          Submodule ℂ leechTwistedSpace) : ℕ) : ℂ)
      = coeff n JShift
\end{leancode}
\noindent File \leanpath{VOA/\allowbreak Moonshine/\allowbreak Character.lean}.
\noindent\textbf{Hypotheses / status.} None\par
\medskip
\endgroup
\chapter{Twisted modules, the twisted Fock space and the module \texorpdfstring{$T$}{T}}
\begingroup
\providecommand{\chsixS}{S[\theta]}
\providecommand{\chsixeps}{\varepsilon}
\noindent{\small 6 of the chapter's 27 numbered statements: its headline results and every definition they use.}\par\medskip
\begin{bproposition}{34.6}
Every ordinary $V$-module is a $1$-twisted module with $N=1$, with $u_{(k/1)}=u_{(k)}$ and the same
grading.
\end{bproposition}
{\raggedright\noindent\textbf{Lean.} \lean{VOA.}\allowbreak\lean{OrdVModule.}\allowbreak\lean{toTwisted}\par}
{\raggedright\noindent File \leanpath{VOA/\allowbreak Modularity/\allowbreak Twisted.lean}.\par}
\begin{leancode}
noncomputable def OrdVModule.toTwisted {M : Type u} [AddCommGroup M] [Module ℂ M]
    (Nm : OrdVModule W M) : TwistedModule (1 : Aut W) 1 M where
\end{leancode}
\noindent\textbf{Hypotheses / status.} Unconditional\par
\medskip
\begin{btheorem}{34.8}[Dead-end certificate]
Let $N\ge2$, and let $Y_M:V\to(\Z\to\End M)$ be linear with $\mathbf 1_{(k/N)}=\delta_{k,-N}$.  If
$\mathrm{JacobiAllIndices}(g,N,Y_M)$ holds, then $M=0$: every $x\in M$ is zero.
\end{btheorem}
{\raggedright\noindent\textbf{Lean.} \lean{VOA.}\allowbreak\lean{jacobiAllIndices_eq_zero}\par}
{\raggedright\noindent File \leanpath{VOA/\allowbreak Moonshine/\allowbreak HeisenbergTwistedModuleDefect.lean}.\par}
\begin{leancode}
theorem jacobiAllIndices_eq_zero {g : Aut W} {N : ℕ} (hN : 2 ≤ N) {M : Type*}
    [AddCommGroup M] [Module ℂ M] (YM : 𝒱 →ₗ[ℂ] (ℤ → End ℂ M))
    (hvac : ∀ k : ℤ, YM W.vacuum k = if k = -(N : ℤ) then 1 else 0)
    (hJ : JacobiAllIndices g N YM) (x : M) : x = 0
\end{leancode}
\noindent\textbf{Hypotheses / status.} $N\ge2$, vacuum\par
\medskip
\begin{btheorem}{34.23}[Dimension of the simple module]
Under the same hypothesis, $\C^c[V]\cong M_k(\C)$ with $k^2=2^n$; hence $n$ is even and
$k=2^{n/2}$.  Put $T=\C^k$ with $\C^c[V]$ acting through this isomorphism.  Then $T$ is a simple
module of dimension $2^{n/2}$, and every simple $\C^c[V]$-module is isomorphic to $T$.
\end{btheorem}
{\raggedright\noindent\textbf{Lean.} \lean{VOA.}\allowbreak\lean{Moonshine.}\allowbreak\lean{Extraspecial.}\allowbreak\lean{repDim},\allowbreak\ \lean{VOA.}\allowbreak\lean{Moonshine.}\allowbreak\lean{Extraspecial.}\allowbreak\lean{repDim_eq},\allowbreak\ \lean{VOA.}\allowbreak\lean{Moonshine.}\allowbreak\lean{Extraspecial.}\allowbreak\lean{Rep},\allowbreak\ \lean{VOA.}\allowbreak\lean{Moonshine.}\allowbreak\lean{Extraspecial.}\allowbreak\lean{repEquiv},\allowbreak\ \lean{VOA.}\allowbreak\lean{Moonshine.}\allowbreak\lean{Extraspecial.}\allowbreak\lean{finrank_rep},\allowbreak\ \lean{VOA.}\allowbreak\lean{Moonshine.}\allowbreak\lean{Extraspecial.}\allowbreak\lean{isSimpleModule_rep},\allowbreak\ \lean{VOA.}\allowbreak\lean{Moonshine.}\allowbreak\lean{Extraspecial.}\allowbreak\lean{nonempty_linearEquiv_rep},\allowbreak\ \lean{VOA.}\allowbreak\lean{Moonshine.}\allowbreak\lean{Extraspecial.}\allowbreak\lean{even_of_isUnit}\par}
\begin{leancode}
noncomputable def repDim (hN : IsUnit (M + Mᵀ).det) : ℕ := (exists_matrix M hN).choose
\end{leancode}
{\raggedright\noindent File \leanpath{VOA/\allowbreak Moonshine/\allowbreak ExtraspecialModule.lean}.\par}
\begin{leancode}
theorem repDim_eq (hN : IsUnit (M + Mᵀ).det) : repDim M hN = 2 ^ (n / 2)
\end{leancode}
{\raggedright\noindent File \leanpath{VOA/\allowbreak Moonshine/\allowbreak ExtraspecialModule.lean}.\par}
\begin{leancode}
def Rep (hN : IsUnit (M + Mᵀ).det) : Type := Fin (repDim M hN) → ℂ
\end{leancode}
{\raggedright\noindent File \leanpath{VOA/\allowbreak Moonshine/\allowbreak ExtraspecialModule.lean}.\par}
\begin{leancode}
noncomputable def repEquiv (hN : IsUnit (M + Mᵀ).det) : alg M ≃ₐ[ℂ] Module.End ℂ (Rep M hN) :=
  (exists_matrix M hN).choose_spec.2.some.trans Matrix.toLinAlgEquiv'
\end{leancode}
{\raggedright\noindent File \leanpath{VOA/\allowbreak Moonshine/\allowbreak ExtraspecialModule.lean}.\par}
\begin{leancode}
theorem finrank_rep (hN : IsUnit (M + Mᵀ).det) :
    Module.finrank ℂ (Rep M hN) = 2 ^ (n / 2)
\end{leancode}
{\raggedright\noindent File \leanpath{VOA/\allowbreak Moonshine/\allowbreak ExtraspecialModule.lean}.\par}
\begin{leancode}
instance isSimpleModule_rep (hN : IsUnit (M + Mᵀ).det) : IsSimpleModule (alg M) (Rep M hN)
\end{leancode}
{\raggedright\noindent File \leanpath{VOA/\allowbreak Moonshine/\allowbreak ExtraspecialModule.lean}.\par}
\begin{leancode}
theorem nonempty_linearEquiv_rep (hN : IsUnit (M + Mᵀ).det) (S : Type*) [AddCommGroup S]
    [Module (alg M) S] [IsSimpleModule (alg M) S] :
    Nonempty (S ≃ₗ[alg M] Rep M hN)
\end{leancode}
{\raggedright\noindent File \leanpath{VOA/\allowbreak Moonshine/\allowbreak ExtraspecialModule.lean}.\par}
\begin{leancode}
theorem even_of_isUnit (hN : IsUnit (M + Mᵀ).det) : Even n
\end{leancode}
\noindent\textbf{Hypotheses / status.} N\par
\medskip
\begin{btheorem}{34.24}
(a) $\mathrm{red}:L\to\mathbb F_2^n$ is surjective with kernel $2L$, so it identifies $L/2L$ with
$\mathbb F_2^n$.  (b) $\chsixeps(\alpha,\beta)=c(\bar\alpha,\bar\beta)$ with $c$ built from
\lean{cocMat}; in particular $\chsixeps$ depends only on classes mod $2L$.  (c) If $L$ is even and
unimodular, the commutator form on $L/2L$ is nondegenerate.
\end{btheorem}
{\raggedright\noindent\textbf{Lean.} \lean{VOA.}\allowbreak\lean{IntegralLattice.}\allowbreak\lean{red_surjective},\allowbreak\ \lean{VOA.}\allowbreak\lean{IntegralLattice.}\allowbreak\lean{red_eq_zero_iff},\allowbreak\ \lean{VOA.}\allowbreak\lean{IntegralLattice.}\allowbreak\lean{latticeCocycle_eq_coc},\allowbreak\ \lean{VOA.}\allowbreak\lean{IntegralLattice.}\allowbreak\lean{latticeCocycle_add_two_nsmul},\allowbreak\ \lean{VOA.}\allowbreak\lean{IntegralLattice.}\allowbreak\lean{isUnit_det_cocMat}\par}
{\raggedright\noindent File \leanpath{VOA/\allowbreak Moonshine/\allowbreak ExtraspecialModule.lean}.\par}
\begin{leancode}
theorem red_surjective : Function.Surjective (red L)
\end{leancode}
{\raggedright\noindent File \leanpath{VOA/\allowbreak Moonshine/\allowbreak ExtraspecialModule.lean}.\par}
\begin{leancode}
theorem red_eq_zero_iff (α : L.carrier) : red L α = 0 ↔ ∃ γ, α = 2 • γ
\end{leancode}
{\raggedright\noindent File \leanpath{VOA/\allowbreak Moonshine/\allowbreak ExtraspecialModule.lean}.\par}
\begin{leancode}
theorem latticeCocycle_eq_coc (α β : L.carrier) :
    latticeCocycle L α β = coc (cocMat L) (red L α) (red L β)
\end{leancode}
{\raggedright\noindent File \leanpath{VOA/\allowbreak Moonshine/\allowbreak ExtraspecialModule.lean}.\par}
\begin{leancode}
theorem latticeCocycle_add_two_nsmul (α β γ : L.carrier) :
    latticeCocycle L (α + 2 • γ) β = latticeCocycle L α β ∧
      latticeCocycle L α (β + 2 • γ) = latticeCocycle L α β
\end{leancode}
{\raggedright\noindent File \leanpath{VOA/\allowbreak Moonshine/\allowbreak ExtraspecialModule.lean}.\par}
\begin{leancode}
theorem isUnit_det_cocMat (hEven : L.Even) (hU : IsUnit (gramMat L).det) :
    IsUnit (cocMat L + (cocMat L)ᵀ).det
\end{leancode}
\noindent\textbf{Hypotheses / status.} (E, U)\par
\medskip
\begin{btheorem}{34.25}
Let $L$ be even unimodular and $T_L$ the simple $\C^\chsixeps[L/2L]$-module.  Then for all
$\alpha,\beta,\gamma\in L$:
$e_\alpha e_\beta=\chsixeps(\alpha,\beta)e_{\alpha+\beta}$, $e_0=1$, $e_{-\alpha}=e_\alpha$,
$e_{\alpha+2\gamma}=e_\alpha$ on $T_L$; and $\dim T_L=2^{\operatorname{rank}L/2}$.  For the Leech
lattice, $\dim T=4096$.
\end{btheorem}
{\raggedright\noindent\textbf{Lean.} \lean{VOA.}\allowbreak\lean{IntegralLattice.}\allowbreak\lean{eT_mul},\allowbreak\ \lean{VOA.}\allowbreak\lean{IntegralLattice.}\allowbreak\lean{eT_zero},\allowbreak\ \lean{VOA.}\allowbreak\lean{IntegralLattice.}\allowbreak\lean{eT_neg},\allowbreak\ \lean{VOA.}\allowbreak\lean{IntegralLattice.}\allowbreak\lean{eT_add_two_nsmul},\allowbreak\ \lean{VOA.}\allowbreak\lean{IntegralLattice.}\allowbreak\lean{finrank_extraT},\allowbreak\ \lean{VOA.}\allowbreak\lean{Moonshine.}\allowbreak\lean{leechT},\allowbreak\ \lean{VOA.}\allowbreak\lean{Moonshine.}\allowbreak\lean{finrank_leechT}\par}
{\raggedright\noindent File \leanpath{VOA/\allowbreak Moonshine/\allowbreak ExtraspecialModule.lean}.\par}
\begin{leancode}
theorem eT_mul (α β : L.carrier) :
    eT L hEven hU α * eT L hEven hU β = latticeCocycle L α β • eT L hEven hU (α + β)
\end{leancode}
{\raggedright\noindent File \leanpath{VOA/\allowbreak Moonshine/\allowbreak ExtraspecialModule.lean}.\par}
\begin{leancode}
theorem eT_zero : eT L hEven hU 0 = 1
\end{leancode}
{\raggedright\noindent File \leanpath{VOA/\allowbreak Moonshine/\allowbreak ExtraspecialModule.lean}.\par}
\begin{leancode}
theorem eT_neg (α : L.carrier) : eT L hEven hU (-α) = eT L hEven hU α
\end{leancode}
{\raggedright\noindent File \leanpath{VOA/\allowbreak Moonshine/\allowbreak ExtraspecialModule.lean}.\par}
\begin{leancode}
theorem eT_add_two_nsmul (α γ : L.carrier) : eT L hEven hU (α + 2 • γ) = eT L hEven hU α
\end{leancode}
{\raggedright\noindent File \leanpath{VOA/\allowbreak Moonshine/\allowbreak ExtraspecialModule.lean}.\par}
\begin{leancode}
theorem finrank_extraT : Module.finrank ℂ (ExtraT L hEven hU) = 2 ^ (L.rank / 2)
\end{leancode}
{\raggedright\noindent File \leanpath{VOA/\allowbreak Moonshine/\allowbreak ExtraspecialModule.lean}.\par}
\begin{leancode}
noncomputable abbrev leechT : Type := ExtraT leech leech_even leech_isUnit_det_gramMat
\end{leancode}
{\raggedright\noindent File \leanpath{VOA/\allowbreak Moonshine/\allowbreak ExtraspecialModule.lean}.\par}
\begin{leancode}
theorem finrank_leechT : Module.finrank ℂ leechT = 4096
\end{leancode}
\noindent\textbf{Hypotheses / status.} E, U\par
\medskip
\begin{btheorem}{34.30}
For a quadratic sign $s$ and $u\in T^*$: $B_u$ is symmetric; $B_u(Xt,t')=B_u(t,\iota(X)t')$ for all
$X\in\End(T)$; and if $B_u\ne0$ then $B_u$ is nondegenerate.  Consequently, if there is an
idempotent $P\in\End(T)$ with $\iota(P)=P$ and $\tr P=1$, then $T$ carries a nondegenerate symmetric
bilinear form $B$ with $B(Xt,t')=B(t,\iota(X)t')$ for all $X$.
\end{btheorem}
{\raggedright\noindent\textbf{Lean.} \lean{VOA.}\allowbreak\lean{Moonshine.}\allowbreak\lean{Extraspecial.}\allowbreak\lean{formOf_symm},\allowbreak\ \lean{VOA.}\allowbreak\lean{Moonshine.}\allowbreak\lean{Extraspecial.}\allowbreak\lean{formOf_invariant},\allowbreak\ \lean{VOA.}\allowbreak\lean{Moonshine.}\allowbreak\lean{Extraspecial.}\allowbreak\lean{formOf_nondegenerate},\allowbreak\ \lean{VOA.}\allowbreak\lean{Moonshine.}\allowbreak\lean{Extraspecial.}\allowbreak\lean{exists_form_of_idempotent}\par}
{\raggedright\noindent File \leanpath{VOA/\allowbreak Moonshine/\allowbreak TFormAlg.lean}.\par}
\begin{leancode}
theorem formOf_symm {s : (Fin n → ZMod 2) → ℂ} (hs : IsQuadSign M s)
    (u : Module.Dual ℂ (Rep M hN)) (t t' : Rep M hN) :
    formOf M hN s u t t' = formOf M hN s u t' t
\end{leancode}
{\raggedright\noindent File \leanpath{VOA/\allowbreak Moonshine/\allowbreak TFormAlg.lean}.\par}
\begin{leancode}
theorem formOf_invariant {s : (Fin n → ZMod 2) → ℂ} (hs : IsQuadSign M s)
    (u : Module.Dual ℂ (Rep M hN)) (X : Module.End ℂ (Rep M hN)) (t t' : Rep M hN) :
    formOf M hN s u (X t) t' = formOf M hN s u t (iotaT M hN s X t')
\end{leancode}
{\raggedright\noindent File \leanpath{VOA/\allowbreak Moonshine/\allowbreak TFormAlg.lean}.\par}
\begin{leancode}
theorem formOf_nondegenerate {s : (Fin n → ZMod 2) → ℂ} (hs : IsQuadSign M s)
    (u : Module.Dual ℂ (Rep M hN)) {v w : Rep M hN} (hvw : formOf M hN s u v w ≠ 0) :
    (formOf M hN s u).Nondegenerate
\end{leancode}
{\raggedright\noindent File \leanpath{VOA/\allowbreak Moonshine/\allowbreak TFormAlg.lean}.\par}
\begin{leancode}
theorem exists_form_of_idempotent {s : (Fin n → ZMod 2) → ℂ} (hs : IsQuadSign M s)
    {P : Module.End ℂ (Rep M hN)} (hP : IsIdempotentElem P) (hιP : iotaT M hN s P = P)
    (htr : trace ℂ (Rep M hN) P = 1) :
    ∃ B : LinearMap.BilinForm ℂ (Rep M hN), (∀ t t', B t t' = B t' t) ∧ B.Nondegenerate ∧
      ∀ X t t', B (X t) t' = B t (iotaT M hN s X t')
\end{leancode}
\noindent\textbf{Hypotheses / status.} N, Q, $P$\par
\medskip
\endgroup
\chapter{The twisted Heisenberg module}
\begingroup
\providecommand{\chsevenF}{\widetilde{F}}
\providecommand{\chsevenst}{\mathbin{\star}}
\noindent{\small 9 of the chapter's 41 numbered statements: its headline results and every definition they use.}\par\medskip
\begin{bdefinition}{35.4}[Twisted eigenspaces; twisted module]
Let $W$ be a VOA on $\mathcal V$ and $g\in\Aut(W)$. For $N\ge1$ and $j\in\N$ put
$\mathcal V^j = \{v : g v = e^{-2\pi i j/N}v\}$. A $g$-twisted module of order $N$ on $M$ is a
linear map $u\mapsto (u_{(k/N)})_{k\in\Z}$, $\mathcal V\to (\Z\to\End M)$, with: (i)
$\mathbf 1_{(k/N)} = \delta_{k,-N}\,\mathrm{id}$; (ii) lower truncation; (iii) for $u\in\mathcal
V^j$, $u_{(k/N)} = 0$ unless $k\equiv -j \pmod N$; (iv) for $u\in\mathcal V^j$, $v\in\mathcal
V$, $x\in M$, $\ell, m\in \Z$ and $k\equiv -j \pmod N$,
\begin{multline}
  \sum_{i\in\Z}(-1)^i\tbinom{\ell}{i}\Bigl(u_{(\ell+\frac kN-i)}\,v_{(\frac mN+i)}x - (-1)^\ell
  v_{(\ell+\frac mN-i)}\,u_{(\frac kN+i)}x\Bigr)\\
  = \sum_{i\ge 0}\tbinom{k/N}{i}\,(u_{(\ell+i)}v)_{(\frac{k+m}{N}-i)}x;
\end{multline}
(v) a direct-sum decomposition $M = \bigoplus_{\alpha\in\C}M_\alpha$ with $\omega_{(1)} = L_0$
acting on $M_\alpha$ by $\alpha$, each $M_\alpha$ finite-dimensional, and $\mathrm{Re}\,\alpha$
bounded below on the nonzero $M_\alpha$.
\end{bdefinition}
{\raggedright\noindent\textbf{Lean.} \lean{VOA.}\allowbreak\lean{Aut.}\allowbreak\lean{twistedEigenspace},\allowbreak\ \lean{VOA.}\allowbreak\lean{TwistedModule},\allowbreak\ \lean{VOA.}\allowbreak\lean{TwistedJacobi}\par}
{\raggedright\noindent File \leanpath{VOA/\allowbreak Modularity/\allowbreak Twisted.lean}.\par}
\begin{leancode}
noncomputable def Aut.twistedEigenspace (g : Aut W) (N : ℕ) (j : ℕ) : Submodule ℂ 𝒱 :=
  Module.End.eigenspace (g.toLinearEquiv : Module.End ℂ 𝒱)
    (Complex.exp (-(2 * (Real.pi : ℂ) * Complex.I * j) / N))

structure TwistedModule (g : Aut W) (N : ℕ) (M : Type u) [AddCommGroup M] [Module ℂ M] where
  YM : 𝒱 →ₗ[ℂ] (ℤ → Module.End ℂ M)
  vacuum_YM : ∀ k : ℤ, YM W.vacuum k = if k = -(N : ℤ) then 1 else 0
  truncation : ∀ (u : 𝒱) (x : M), ∃ K : ℤ, ∀ k : ℤ, K ≤ k → YM u k x = 0
  support : ∀ (j : ℕ) (u : 𝒱), u ∈ g.twistedEigenspace N j →
    ∀ k : ℤ, (k + j) % (N : ℤ) ≠ 0 → YM u k = 0
  jacobi : ∀ (j : ℕ) (u : 𝒱), u ∈ g.twistedEigenspace N j → ∀ (v : 𝒱) (x : M) (ℓ k m : ℤ),
    (k + j) % (N : ℤ) = 0 →
    (∑ᶠ i : ℤ, (negOnePow i * cbz ℓ i) •
        (YM u (ℓ * N + k - i * N) (YM v (m + i * N) x)
          - negOnePow ℓ • YM v (ℓ * N + m - i * N) (YM u (k + i * N) x)))
      = ∑ᶠ i : ℕ, cbinom ((k : ℂ) / N) i • YM (W.smul u (ℓ + i) v) (k + m - i * N) x
  weight : ℂ → Submodule ℂ M
  isInternal : DirectSum.IsInternal weight
  L0_weight : ∀ (α : ℂ) (x : M), x ∈ weight α → YM W.conformalVector N x = α • x
  finrank : ∀ α, FiniteDimensional ℂ (weight α)
  boundedBelow : ∃ B : ℝ, ∀ α : ℂ, weight α ≠ ⊥ → B ≤ α.re
\end{leancode}
\begin{leancode}
def TwistedJacobi (g : Aut W) (N : ℕ) {M : Type*} [AddCommGroup M] [Module ℂ M]
    (YM : 𝒱 →ₗ[ℂ] (ℤ → End ℂ M)) : Prop :=
  ∀ (j : ℕ) (u : 𝒱), u ∈ g.twistedEigenspace N j → ∀ (v : 𝒱) (x : M) (ℓ k m : ℤ),
    (k + j) % (N : ℤ) = 0 →
    (∑ᶠ i : ℤ, (negOnePow i * cbz ℓ i) •
        (YM u (ℓ * N + k - i * N) (YM v (m + i * N) x)
          - negOnePow ℓ • YM v (ℓ * N + m - i * N) (YM u (k + i * N) x)))
      = ∑ᶠ i : ℕ, cbinom ((k : ℂ) / N) i • YM (W.smul u (ℓ + i) v) (k + m - i * N) x
\end{leancode}
\medskip
\begin{btheorem}{35.5}
For every $r\in\N$ and level $k\ne0$, $\sigma$ is an automorphism of the Heisenberg VOA
$\lean{heisenbergNVOA r k hk}$: it is a linear bijection with $\sigma(u_{(n)}v) =
(\sigma u)_{(n)}(\sigma v)$ for all $u,v,n$, $\sigma\mathbf 1 = \mathbf 1$ and $\sigma\omega
=\omega$. Moreover $\sigma^2 = 1$.
\end{btheorem}
{\raggedright\noindent\textbf{Lean.} \lean{VOA.}\allowbreak\lean{Aut},\allowbreak\ \lean{VOA.}\allowbreak\lean{HeisenbergN.}\allowbreak\lean{sigmaHeisAut},\allowbreak\ \lean{VOA.}\allowbreak\lean{HeisenbergN.}\allowbreak\lean{sigmaHeisAut_sq},\allowbreak\ \lean{VOA.}\allowbreak\lean{HeisenbergN.}\allowbreak\lean{heisVOA1},\allowbreak\ \lean{VOA.}\allowbreak\lean{HeisenbergN.}\allowbreak\lean{sigma1}\par}
{\raggedright\noindent Files \leanpath{VOA/\allowbreak Core/\allowbreak Automorphism.lean}, \leanpath{VOA/\allowbreak Moonshine/\allowbreak HeisenbergTwistedModuleSigma.lean}.\par}
\begin{leancode}
structure Aut (W : VertexOperatorAlgebra 𝒱) where
  toLinearEquiv : 𝒱 ≃ₗ[ℂ] 𝒱
  map_smul' : ∀ (u : 𝒱) (n : ℤ) (v : 𝒱),
    toLinearEquiv (W.smul u n v) = W.smul (toLinearEquiv u) n (toLinearEquiv v)
  map_vacuum : toLinearEquiv W.vacuum = W.vacuum
  map_omega : toLinearEquiv W.conformalVector = W.conformalVector

noncomputable def sigmaHeisAut (k : ℂ) (hk : k ≠ 0) : Aut (heisenbergNVOA r k hk) where
  toLinearEquiv := sigmaN r
  -- (the three proof fields map_smul', map_vacuum, map_omega are elided)
theorem sigmaHeisAut_sq (k : ℂ) (hk : k ≠ 0) : sigmaHeisAut r k hk * sigmaHeisAut r k hk = 1

noncomputable abbrev heisVOA1 : VertexOperatorAlgebra (FockN r) :=
  heisenbergNVOA r 1 one_ne_zero
noncomputable abbrev sigma1 : Aut (heisVOA1 r) := sigmaHeisAut r 1 one_ne_zero
\end{leancode}
\medskip
\begin{blemma}{35.8}[Complex Chu--Vandermonde]
For $x,y\in\C$ and $s\in\N$, $\binom{x+y}{s} = \sum_{i+t=s}\binom xi\binom yt$.
\end{blemma}
{\raggedright\noindent\textbf{Lean.} \lean{VOA.}\allowbreak\lean{TwistedReduction.}\allowbreak\lean{cbinom_add}\par}
\begin{leancode}
theorem cbinom_add (x y : ℂ) (s : ℕ) :
    cbinom (x + y) s = ∑ ij ∈ Finset.antidiagonal s, cbinom x ij.1 * cbinom y ij.2
\end{leancode}
\noindent File \leanpath{VOA/\allowbreak Moonshine/\allowbreak HeisenbergTwistedModuleReduction.lean}.
\medskip
\begin{btheorem}{35.10}[Twisted Jacobi from locality and the product formula]
Let $Y_M$ be lower truncated, $u,v\in\mathcal V$, $e,e'\in\Z$. Suppose $\chsevenF_e(u)$ and
$\chsevenF_{e'}(v)$ are mutually local and satisfy the twisted product formula (35.7)
at every $\ell'\ge\ell$ (resp.\ every $\ell'$). Then for all $x$ and all $m,n\in\Z$ the twisted
Jacobi identity (35.4) holds for $(u,v)$ at $\ell$ (resp.\ every $\ell$), with
$k=Nm+e$ and the mode index of $v$ equal to $Nn+e'$.
\end{btheorem}
{\raggedright\noindent\textbf{Lean.} \lean{VOA.}\allowbreak\lean{TwistedReduction.}\allowbreak\lean{twisted_jacobi_of_local_pf_ge},\allowbreak\ \lean{VOA.}\allowbreak\lean{TwistedReduction.}\allowbreak\lean{twisted_jacobi_of_local_pf}\par}
\begin{leancode}
theorem twisted_jacobi_of_local_pf_ge (W : VertexOperatorAlgebra 𝒱)
    (YM : 𝒱 →ₗ[ℂ] (ℤ → End ℂ M))
    {N : ℕ} (hN : 0 < N) (htr : ∀ (u : 𝒱) (x : M), ∃ K : ℤ, ∀ k : ℤ, K ≤ k → YM u k x = 0)
    (u v : 𝒱) (e e' : ℤ)
    (hloc : Local (resc (YM u) N e) (resc (YM v) N e')) (ℓ : ℤ)
    (hpf : ∀ (ℓ' s : ℤ) (x : M), ℓ ≤ ℓ' →
      mode (Fld.prod ℓ' (rescFld hN (htr u) e) (rescFld hN (htr v) e')).1 s x
        = ∑ᶠ t : ℕ, cbinom ((e : ℂ) / N) t •
            YM (W.smul u (ℓ' + t) v) (N * (s - t) + e + e') x)
    (x : M) (m n : ℤ) :
    (∑ᶠ i : ℤ, (negOnePow i * cbz ℓ i) •
        (YM u (ℓ * N + (N * m + e) - i * N) (YM v ((N * n + e') + i * N) x)
          - negOnePow ℓ • YM v (ℓ * N + (N * n + e') - i * N) (YM u ((N * m + e) + i * N) x)))
      = ∑ᶠ i : ℕ, cbinom (((N * m + e : ℤ) : ℂ) / N) i •
          YM (W.smul u (ℓ + i) v) ((N * m + e) + (N * n + e') - i * N) x

theorem twisted_jacobi_of_local_pf (W : VertexOperatorAlgebra 𝒱) (YM : 𝒱 →ₗ[ℂ] (ℤ → End ℂ M))
    {N : ℕ} (hN : 0 < N) (htr : ∀ (u : 𝒱) (x : M), ∃ K : ℤ, ∀ k : ℤ, K ≤ k → YM u k x = 0)
    (u v : 𝒱) (e e' : ℤ)
    (hloc : Local (resc (YM u) N e) (resc (YM v) N e'))
    (hpf : ∀ (ℓ s : ℤ) (x : M),
      mode (Fld.prod ℓ (rescFld hN (htr u) e) (rescFld hN (htr v) e')).1 s x
        = ∑ᶠ t : ℕ, cbinom ((e : ℂ) / N) t •
            YM (W.smul u (ℓ + t) v) (N * (s - t) + e + e') x)
    (x : M) (ℓ m n : ℤ) :
    (∑ᶠ i : ℤ, (negOnePow i * cbz ℓ i) •
        (YM u (ℓ * N + (N * m + e) - i * N) (YM v ((N * n + e') + i * N) x)
          - negOnePow ℓ • YM v (ℓ * N + (N * n + e') - i * N) (YM u ((N * m + e) + i * N) x)))
      = ∑ᶠ i : ℕ, cbinom (((N * m + e : ℤ) : ℂ) / N) i •
          YM (W.smul u (ℓ + i) v) ((N * m + e) + (N * n + e') - i * N) x
\end{leancode}
\noindent File \leanpath{VOA/\allowbreak Moonshine/\allowbreak HeisenbergTwistedModuleReduction.lean}.
\medskip
\begin{btheorem}{35.15}[The generator commutator]
\begin{equation*}
  [F_a(z_1),F_b(z_2)] = \delta_{ab}\Bigl(z_2\,\partial_{z_2}\delta(z_1,z_2) +
  \tfrac12\,\delta(z_1,z_2)\Bigr).
\end{equation*}
Consequently $F_a$, $F_b$ are mutually local (of order $2$), and
$F_{a(0)}F_b = \tfrac12\delta_{ab}\,1$, $F_{a(1)}F_b = \delta_{ab}\,z$, $F_{a(\ell)}F_b = 0$ for
$\ell\ge2$.
\end{btheorem}
{\raggedright\noindent\textbf{Lean.} \lean{VOA.}\allowbreak\lean{Moonshine.}\allowbreak\lean{opeOne},\allowbreak\ \lean{VOA.}\allowbreak\lean{Moonshine.}\allowbreak\lean{opeZero},\allowbreak\ \lean{VOA.}\allowbreak\lean{Moonshine.}\allowbreak\lean{bracketDist_Fgen},\allowbreak\ \lean{VOA.}\allowbreak\lean{Moonshine.}\allowbreak\lean{Fgen_local},\allowbreak\ \lean{VOA.}\allowbreak\lean{Moonshine.}\allowbreak\lean{prod_Fgen_zero},\allowbreak\ \lean{VOA.}\allowbreak\lean{Moonshine.}\allowbreak\lean{prod_Fgen_one},\allowbreak\ \lean{VOA.}\allowbreak\lean{Moonshine.}\allowbreak\lean{prod_Fgen_of_two_le}\par}
\begin{leancode}
noncomputable def opeOne (a b : Fin r) : FormalDist (End ℂ (FockT r)) :=
  fun t => if t = 1 then (if a = b then (1 : ℂ) else 0) • LinearMap.id else 0
noncomputable def opeZero (a b : Fin r) : FormalDist (End ℂ (FockT r)) :=
  fun t => if t = 0 then (if a = b then (1 / 2 : ℂ) else 0) • LinearMap.id else 0

theorem bracketDist_Fgen (a b : Fin r) :
    bracketDist (Fgen r a) (Fgen r b)
      = deltaDeriv 1 (opeOne r a b) + deltaDeriv 0 (opeZero r a b)
theorem Fgen_local (a b : Fin r) : Local (Fgen r a) (Fgen r b)
theorem prod_Fgen_zero (a b : Fin r) :
    (Fld.prod 0 (FgenFld r a) (FgenFld r b)).1 = opeZero r a b
theorem prod_Fgen_one (a b : Fin r) :
    (Fld.prod 1 (FgenFld r a) (FgenFld r b)).1 = opeOne r a b
theorem prod_Fgen_of_two_le (a b : Fin r) {ℓ : ℤ} (hℓ : 2 ≤ ℓ) :
    (Fld.prod ℓ (FgenFld r a) (FgenFld r b)).1 = 0
\end{leancode}
\noindent File \leanpath{VOA/\allowbreak Moonshine/\allowbreak HeisenbergTwistedModuleGenerators.lean}.
\medskip
\begin{btheorem}{35.20}[Locality is inherited]
If the fields of $S$ are pairwise mutually local, then so are all fields of $\langle S\rangle$.
In particular, for $S = \{F_a : a<r\}$, every two fields in $\langle S\rangle$ are mutually local.
\end{btheorem}
{\raggedright\noindent\textbf{Lean.} \lean{VOA.}\allowbreak\lean{TwistedReduction.}\allowbreak\lean{tgen_local},\allowbreak\ \lean{VOA.}\allowbreak\lean{TwistedReduction.}\allowbreak\lean{tgenerated_local},\allowbreak\ \lean{VOA.}\allowbreak\lean{Moonshine.}\allowbreak\lean{genSet},\allowbreak\ \lean{VOA.}\allowbreak\lean{Moonshine.}\allowbreak\lean{twistedSystem_local}\par}
\begin{leancode}
theorem tgen_local (S : Set (Fld M)) (hS : ∀ a ∈ S, ∀ b ∈ S, Local a.1 b.1) :
    ∀ (N i j : ℕ) (a b : Fld M), i + j ≤ N → TGen S i a → TGen S j b → Local a.1 b.1
theorem tgenerated_local {S : Set (Fld M)} (hS : ∀ a ∈ S, ∀ b ∈ S, Local a.1 b.1)
    {a b : Fld M} (ha : a ∈ TGenerated S) (hb : b ∈ TGenerated S) : Local a.1 b.1

def genSet : Set (Fld (FockT r)) := Set.range (FgenFld r)
theorem twistedSystem_local {a b : Fld (FockT r)} (ha : a ∈ TGenerated (genSet r))
    (hb : b ∈ TGenerated (genSet r)) : Local a.1 b.1
\end{leancode}
\noindent File \leanpath{VOA/\allowbreak Moonshine/\allowbreak HeisenbergTwistedModuleLocality.lean}.
\medskip
\begin{bdefinition}{35.22}[The twisted vertex operator]
For a monomial $x^m$ of degree $d$ with canonical word $l_m$,
$Y^T(x^m)_{(k/2)} = \widetilde Y(l_m)_{((k-d)/2)}$ if $k\equiv d\pmod2$, and $0$ otherwise.
$Y^T$ is extended linearly to $\lean{FockN r}$.
\end{bdefinition}
{\raggedright\noindent\textbf{Lean.} \lean{VOA.}\allowbreak\lean{HeisenbergN.}\allowbreak\lean{monoWordN},\allowbreak\ \lean{VOA.}\allowbreak\lean{Moonshine.}\allowbreak\lean{YTmon},\allowbreak\ \lean{VOA.}\allowbreak\lean{Moonshine.}\allowbreak\lean{YT}\par}
{\raggedright\noindent Files \leanpath{VOA/\allowbreak Examples/\allowbreak HeisenbergNVertex.lean}, \leanpath{VOA/\allowbreak Moonshine/\allowbreak HeisenbergTwistedModuleVertex.lean}.\par}
\begin{leancode}
noncomputable def monoWordN (a : (ℕ × Fin r) →₀ ℕ) : List (ℕ × Fin r) :=
  a.support.toList.flatMap (fun ia => List.replicate (a ia) ia)

noncomputable def YTmon (m : (ℕ × Fin r) →₀ ℕ) : ℤ → End ℂ (FockT r) := fun k =>
  if (k - (monoWordN m).length) % 2 = 0 then
    mode (Ywt r (monoWordN m)).1 ((k - (monoWordN m).length) / 2)
  else 0

noncomputable def YT : FockN r →ₗ[ℂ] (ℤ → End ℂ (FockT r)) :=
  Finsupp.linearCombination ℂ (YTmon r)
\end{leancode}
\medskip
\begin{btheorem}{35.23}
(i) $Y^T(\mathbf 1)_{(k/2)} = \delta_{k,-2}\,\mathrm{id}$; (ii) $Y^T(h_a(-1)\mathbf 1)_{(k/2)} =
h_a(k/2)$; (iii) lower truncation; (iv) the support condition: for $u\in\mathcal V^j$,
$Y^T(u)_{(k/2)} = 0$ unless $k\equiv -j\pmod2$.
\end{btheorem}
{\raggedright\noindent\textbf{Lean.} \lean{VOA.}\allowbreak\lean{Moonshine.}\allowbreak\lean{YT_vacuum},\allowbreak\ \lean{VOA.}\allowbreak\lean{Moonshine.}\allowbreak\lean{YT_X},\allowbreak\ \lean{VOA.}\allowbreak\lean{Moonshine.}\allowbreak\lean{YT_truncation},\allowbreak\ \lean{VOA.}\allowbreak\lean{Moonshine.}\allowbreak\lean{YT_support}\par}
\begin{leancode}
theorem YT_vacuum (k : ℤ) : YT r (vacuumN r) k = if k = -2 then 1 else 0
theorem YT_X (a : Fin r) : YT r (X (0, a)) = modeT r a
theorem YT_truncation (u : FockN r) (x : FockT r) :
    ∃ K : ℤ, ∀ k : ℤ, K ≤ k → YT r u k x = 0
theorem YT_support (j : ℕ) (u : FockN r) (hu : u ∈ (sigma1 r).twistedEigenspace 2 j)
    (k : ℤ) (hk : (k + j) % ((2 : ℕ) : ℤ) ≠ 0) : YT r u k = 0
\end{leancode}
\noindent File \leanpath{VOA/\allowbreak Moonshine/\allowbreak HeisenbergTwistedModuleVertex.lean}.
\medskip
\begin{btheorem}{35.43}[The twisted Heisenberg module]
For every $r\in\N$, the twisted Fock space $S[\theta]$, with the twisted vertex operator $Y^T$ of
Definition~35.22 and the grading $M_\alpha$, is a $\sigma$-twisted module of order $2$
for the level-one rank-$r$ Heisenberg VOA, in the sense of Definition~35.4. Its
generators act by the twisted bosons, and its $L_0$ is $\sum h_a(-n-\frac12)h_a(n+\frac12) +
\frac r{16}$.
\end{btheorem}
{\raggedright\noindent\textbf{Lean.} \lean{VOA.}\allowbreak\lean{Moonshine.}\allowbreak\lean{heisTwistedModule},\allowbreak\ \lean{VOA.}\allowbreak\lean{Moonshine.}\allowbreak\lean{heisTwistedModule_YM}\par}
\begin{leancode}
noncomputable def heisTwistedModule : TwistedModule (sigma1 r) 2 (FockT r) where
  YM := YT r
  vacuum_YM k := by
    change YT r (vacuumN r) k = _
    rw [YT_vacuum]
    push_cast
    rfl
  truncation := YT_truncation r
  support := fun j u hu k hk => YT_support r j u hu k hk
  jacobi := twistedJacobi_YT r
  weight := weightC r
  isInternal := weightC_isInternal r
  L0_weight := fun α x hx => YT_omega_two_weightC r hx
  finrank := weightC_finiteDimensional r
  boundedBelow := weightC_boundedBelow r

@[simp] theorem heisTwistedModule_YM : (heisTwistedModule r).YM = YT r := rfl
\end{leancode}
\noindent File \leanpath{VOA/\allowbreak Moonshine/\allowbreak HeisenbergTwistedModuleInstance.lean}.
\medskip
\endgroup
\chapter{The twisted Leech lattice module and \texorpdfstring{$(V_\Lambda^T)^+$}{the twisted plus part}}
\begingroup
\providecommand{\cheightS}{S[\theta]}
\providecommand{\cheightT}{\mathcal{T}}
\providecommand{\cheightE}[1]{E^{#1}}
\providecommand{\cheightpair}[2]{\langle #1,#2\rangle}
\noindent{\small 6 of the chapter's 46 numbered statements: its headline results and every definition they use.}\par\medskip
\begin{blemma}{36.6}[Characterisation and the exponential law]
Let $\theta$ be the Euler operator on $R[[X]]$ and $D(d)=\sum_{k\ge1}d_kX^k$.
\begin{enumerate}
\item If $F\in R[[X]]$ satisfies $\theta F=D(d)F$ and $F_0=1$, then $F=\sum_p\mathrm{nexp}(d)_pX^p$.
\item $\mathrm{nexp}(d+e)=\mathrm{nexp}(d)\cdot\mathrm{nexp}(e)$ as power series.
\item If $\varphi:R\to S$ is a $\C$-algebra map, $\varphi(\mathrm{nexp}(d)_p)=\mathrm{nexp}(\varphi\circ d)_p$.
\item $\mathrm{nexp}(k\mapsto t^kd_k)_p=t^p\,\mathrm{nexp}(d)_p$.
\item If $\delta$ is a $\C$-derivation of $R$ with $\delta(d_k)=\delta_{k,K}\,K\,c$ for some
  $K\ge1$, then $\delta(E_p)=c\,E_{p-K}$ for $p\ge K$ and $\delta(E_p)=0$ for $p<K$.
\end{enumerate}
\end{blemma}
{\raggedright\noindent\textbf{Lean.} \lean{VOA.}\allowbreak\lean{Moonshine.}\allowbreak\lean{LatticeTwisted.}\allowbreak\lean{eq_nexpPS},\allowbreak\ \lean{VOA.}\allowbreak\lean{Moonshine.}\allowbreak\lean{LatticeTwisted.}\allowbreak\lean{nexpPS_add},\allowbreak\ \lean{VOA.}\allowbreak\lean{Moonshine.}\allowbreak\lean{LatticeTwisted.}\allowbreak\lean{nexp_map},\allowbreak\ \lean{VOA.}\allowbreak\lean{Moonshine.}\allowbreak\lean{LatticeTwisted.}\allowbreak\lean{nexp_scale},\allowbreak\ \lean{VOA.}\allowbreak\lean{Moonshine.}\allowbreak\lean{LatticeTwisted.}\allowbreak\lean{nexp_derivation}\par}
\begin{leancode}
theorem eq_nexpPS (d : ℕ → R) (F : R⟦X⟧) (h : theta F = Dser d * F) (h0 : coeff 0 F = 1) :
    F = nexpPS d
theorem nexpPS_add (d e : ℕ → R) : nexpPS (d + e) = nexpPS d * nexpPS e
theorem nexp_map {S : Type*} [CommRing S] [Algebra ℂ S] (φ : R →ₐ[ℂ] S) (d : ℕ → R) (p : ℕ) :
    φ (nexp d p) = nexp (fun k => φ (d k)) p
theorem nexp_scale (t : R) (d : ℕ → R) (p : ℕ) :
    nexp (fun k => t ^ k * d k) p = t ^ p * nexp d p
theorem nexp_derivation (δ : Derivation ℂ R R) (d : ℕ → R) (K : ℕ) (hK : 0 < K) (c : R)
    (hd : ∀ k, δ (d k) = if k = K then (K : ℂ) • c else 0) (p : ℕ) :
    δ (nexp d p) = if K ≤ p then c * nexp d (p - K) else 0
\end{leancode}
\noindent File \leanpath{VOA/\allowbreak Moonshine/\allowbreak LatticeTwistedExp.lean}.
\noindent\textbf{Hypotheses / status.} None\par
\medskip
\begin{btheorem}{36.17}[Locality of the dressings]
If $\cheightpair vu=n\in\Z$ and $N\ge|n|$, then for all $A,B\in\Z$ and $x\in\cheightS$
\begin{equation*}
  \sum_{r=0}^{N}(-1)^{N-r}\binom Nr\Bigl(\mathrm{DD}_v(A-2r)\mathrm{DD}_u(B-2(N-r))
  -(-1)^{|n|}\mathrm{DD}_u(B-2(N-r))\mathrm{DD}_v(A-2r)\Bigr)x=0,
\end{equation*}
that is, $(w_1^2-w_2^2)^N\bigl(Y_v(w_1)Y_u(w_2)-(-1)^nY_u(w_2)Y_v(w_1)\bigr)=0$ with
$Y_v(w)=\cheightE{-}\cheightE{+}(-v,w)$.
\end{btheorem}
{\raggedright\noindent\textbf{Lean.} \lean{VOA.}\allowbreak\lean{Moonshine.}\allowbreak\lean{LatticeTwisted.}\allowbreak\lean{DD_bracket_local}\par}
\begin{leancode}
theorem DD_bracket_local (v u : Fin r → ℂ) (n : ℤ) (hn : pair v u = n) (N : ℕ)
    (h1 : -(N : ℤ) ≤ n) (h2 : n ≤ N) (A B : ℤ) (x : FockT r) :
    ∑ r ∈ Finset.range (N + 1), ((-1 : ℂ) ^ (N - r) * (N.choose r : ℂ)) •
      (DD v (A - 2 * r) (DD u (B - 2 * (N - r : ℕ)) x)
        - (-1 : ℂ) ^ n.natAbs • DD u (B - 2 * (N - r : ℕ)) (DD v (A - 2 * r) x)) = 0
\end{leancode}
\noindent File \leanpath{VOA/\allowbreak Moonshine/\allowbreak LatticeTwistedTwoPoint.lean}.
\noindent\textbf{Hypotheses / status.} $\langle v,u\rangle\in\Z$, $|n|\le N$\par
\medskip
\begin{btheorem}{36.21}[Properties of $Y^T(e^\alpha)$]
Let $L$ be even, positive definite, unimodular.
\begin{enumerate}
\item $e_\beta e_\alpha=(-1)^{\cheightpair\alpha\beta}e_\alpha e_\beta$ on $\cheightT_L$.
\item For all $\alpha,\beta\in L$ and residues $e,e'$, the rescaled fields of $Y^T(e^\alpha)$ and
  $Y^T(e^\beta)$ are mutually local.
\item $Y^T(e^{-\alpha})_{(k/2)}=(-1)^kY^T(e^\alpha)_{(k/2)}$.  Consequently
  $e^\alpha+e^{-\alpha}$ has no odd modes and $e^\alpha-e^{-\alpha}$ no even modes.
\item $Y^T(e^0)_{(k/2)}=\delta_{k,-2}$.
\end{enumerate}
\end{btheorem}
{\raggedright\noindent\textbf{Lean.} \lean{VOA.}\allowbreak\lean{Moonshine.}\allowbreak\lean{LatticeTwisted.}\allowbreak\lean{eT_swap},\allowbreak\ \lean{VOA.}\allowbreak\lean{Moonshine.}\allowbreak\lean{LatticeTwisted.}\allowbreak\lean{YTL_local},\allowbreak\ \lean{VOA.}\allowbreak\lean{Moonshine.}\allowbreak\lean{LatticeTwisted.}\allowbreak\lean{YTL_neg},\allowbreak\ \lean{VOA.}\allowbreak\lean{Moonshine.}\allowbreak\lean{LatticeTwisted.}\allowbreak\lean{YTL_zero}\par}
{\raggedright\noindent File \leanpath{VOA/\allowbreak Moonshine/\allowbreak LatticeTwistedVertex.lean}.\par}
\begin{leancode}
theorem eT_swap (α β : L.carrier) :
    eT L hEven hU β * eT L hEven hU α
      = (-1 : ℂ) ^ (L.form α β).natAbs • (eT L hEven hU α * eT L hEven hU β)

theorem YTL_local (hPos : L.PosDef) (α β : L.carrier) (e e' : ℤ) :
    Local (resc (YTL L hEven hU α) 2 e) (resc (YTL L hEven hU β) 2 e') :=
  YTg_local (chargeCoord L α) (chargeCoord L β) (L.form α β) (sum_chargeCoord L hPos α β) _ _ _ _
    _ _ (eT_swap L hEven hU α β) e e'

theorem YTL_neg (α : L.carrier) (k : ℤ) :
    YTL L hEven hU (-α) k = (-1 : ℂ) ^ k • YTL L hEven hU α k

theorem YTL_zero (k : ℤ) : YTL L hEven hU 0 k = if k = -2 then 1 else 0
\end{leancode}
\noindent\textbf{Hypotheses / status.} Lattice hypothesis\par
\medskip
\begin{btheorem}{36.23}[Locality of the twisted lattice system]
For $L$ even, positive definite, unimodular, any two fields in the generated twisted system
$\mathrm{TGenerated}(\text{latGenSet})$ are mutually local.
\end{btheorem}
{\raggedright\noindent\textbf{Lean.} \lean{VOA.}\allowbreak\lean{Moonshine.}\allowbreak\lean{LatticeTwisted.}\allowbreak\lean{latticeTwistedSystem_local}\par}
{\raggedright\noindent File \leanpath{VOA/\allowbreak Moonshine/\allowbreak LatticeTwistedVertex.lean}.\par}
\begin{leancode}
theorem latticeTwistedSystem_local (hPos : L.PosDef)
    {a b : Fld (FockT (rank L) ⊗[ℂ] Rep (cocMat L) (isUnit_det_cocMat L hEven hU))}
    (ha : a ∈ TGenerated (latGenSet L hEven hU)) (hb : b ∈ TGenerated (latGenSet L hEven hU)) :
    Local a.1 b.1 :=
  tgenerated_local (latGenSet_local L hEven hU hPos) ha hb
\end{leancode}
\noindent\textbf{Hypotheses / status.} Lattice hypothesis\par
\medskip
\begin{btheorem}{36.27}[Generators and the reduction to exponentials]
Let $L$ be even, positive definite, unimodular.
\begin{enumerate}
\item For every Heisenberg generator, every state $w$ and every $\ell$:
  $\mathrm{star}_{1/2}(F_b)_\ell F_e(w)=F_{e+1}(h_{b\,(\ell)}w)$.
\item If $\mathrm{PF}(e^\alpha,e^\beta)$ holds for all $\alpha,\beta\in L$, then $\mathrm{PF}(u,v)$
  holds for all $u,v\in V_L$.
\end{enumerate}
\end{btheorem}
{\raggedright\noindent\textbf{Lean.} \lean{VOA.}\allowbreak\lean{Moonshine.}\allowbreak\lean{LatticeTwisted.}\allowbreak\lean{genPF},\allowbreak\ \lean{VOA.}\allowbreak\lean{Moonshine.}\allowbreak\lean{LatticeTwisted.}\allowbreak\lean{PF_of_base}\par}
{\raggedright\noindent File \leanpath{VOA/\allowbreak Moonshine/\allowbreak LatticeTwistedProduct.lean}.\par}
\begin{leancode}
theorem genPF (b : Fin (rank L)) (e' ℓ : ℤ) (w : LatModuleGen L) :
    star (1 / 2) (FgT b) ℓ (FL L hEven hU e' w)
      = FL L hEven hU (e' + 1) ((latticeVOAGen L hEven hPos).smul (hgen L b) ℓ w)

theorem PF_of_base (hbase : ∀ α β : L.carrier, PF L hEven hU hPos (echg L α) (echg L β))
    (u v : LatModuleGen L) : PF L hEven hU hPos u v :=
\end{leancode}
\noindent\textbf{Hypotheses / status.} Lattice hypothesis (+ \lean{hbase})\par
\medskip
\begin{btheorem}{36.48}[Weights and dimensions]
\begin{enumerate}
\item For $q\in\Q$: $\mathrm{TP}_q=\mathrm{twFixed}\cap\mathrm{twWeight}_q$, the graded piece used in
  Chapter~33.
\item If $\mathrm{TP}_\alpha\ne0$ then $\alpha=n\in\Z$ with $n\ge2$.
\item $\dim\mathrm{TP}_2=98304=24\cdot4096$.
\item For every $n\in\N$: $\dim(V_\Lambda^+)_n+\dim\mathrm{TP}_n=[q^n]\bigl(q\,J(q)\bigr)$, $J=j-744$.
\end{enumerate}
\end{btheorem}
{\raggedright\noindent\textbf{Lean.} \lean{VOA.}\allowbreak\lean{Moonshine.}\allowbreak\lean{map_plusModule_weight},\allowbreak\ \lean{VOA.}\allowbreak\lean{Moonshine.}\allowbreak\lean{plusModule_weight_ne_bot},\allowbreak\ \lean{VOA.}\allowbreak\lean{Moonshine.}\allowbreak\lean{finrank_plusModule_weight_two},\allowbreak\ \lean{VOA.}\allowbreak\lean{Moonshine.}\allowbreak\lean{finrank_plusModule_eq_coeff_JShift}\par}
\begin{leancode}
theorem map_plusModule_weight (q : ℚ) :
    (plusModule.weight (q : ℂ)).map TP.subtype = twFixed 24 leechT ⊓ twWeight 24 leechT q

theorem plusModule_weight_ne_bot {α : ℂ} (h : plusModule.weight α ≠ ⊥) :
    ∃ n : ℤ, α = n ∧ 2 ≤ n

theorem finrank_plusModule_weight_two : finrank ℂ (plusModule.weight 2) = 98304

theorem finrank_plusModule_eq_coeff_JShift (n : ℕ) :
    ((finrank ℂ (leechVOAPlus.grading (n : ℤ))
      + finrank ℂ (plusModule.weight ((n : ℤ) : ℂ)) : ℕ) : ℂ)
      = PowerSeries.coeff n JShift
\end{leancode}
\noindent File \leanpath{VOA/\allowbreak Moonshine/\allowbreak TwistedPlus.lean}.
\noindent\textbf{Hypotheses / status.} None\par
\medskip
\endgroup
\chapter{Intertwiners and invariant forms: the operators \texorpdfstring{$W$}{W} and \texorpdfstring{$\overline{W}$}{Wbar}}
\begingroup
\providecommand{\chnineVP}{V_{\Lambda}^{+}}
\providecommand{\chnineTP}{(V_{\Lambda}^{T})^{+}}
\providecommand{\chnineWb}{\overline{W}}
\providecommand{\chnineWd}{W^{\dagger}}
\providecommand{\chnineV}{\mathcal{V}}
\noindent{\small 7 of the chapter's 52 numbered statements: its headline results and every definition they use.}\par\medskip
\begin{btheorem}{37.13}[$L_{-1}$-derivative property]
$W(L_{-1}t)_{(n)}=-n\,W(t)_{(n-1)}$ for all $n$; equivalently $W(L_{-1}t,z)=\partial_zW(t,z)$.
\end{btheorem}
{\raggedright\noindent\textbf{Lean.} \lean{VOA.}\allowbreak\lean{SkewIntertwiner.}\allowbreak\lean{wMode_Lm1},\allowbreak\ \lean{VOA.}\allowbreak\lean{SkewIntertwiner.}\allowbreak\lean{Wsk_deriv}\par}
\begin{leancode}
theorem wMode_Lm1 (A : WeakModule W M) (t : M) (n : ℤ) (φ : 𝒱) :
    wMode A (Lm1 A t) n φ = (-(n : ℂ)) • wMode A t (n - 1) φ

theorem Wsk_deriv (A : WeakModule W M) (t : M) :
    Wsk A (Lm1 A t) = deriv (Wsk A t)
\end{leancode}
\noindent File \leanpath{VOA/\allowbreak Moonshine/\allowbreak SkewIntertwiner.lean}.
\noindent\textbf{Hypotheses / status.} None\par
\medskip
\begin{btheorem}{37.15}[Intertwining commutator formula for $W$]
Let $M$ be a weak $\chnineV$-module, $u\in\chnineV$, $t\in M$, and let $J\in\N$ be such that
$u_{(j)}t=0$ for all $j\ge J$. Then for \emph{all} $m,n\in\Z$ and $\varphi\in\chnineV$,
\begin{equation}
u^M_{(m)}\,W(t)_{(n)}\varphi-W(t)_{(n)}\big(u_{(m)}\varphi\big)
=\sum_{r=0}^{J-1}\binom{m}{r}\,W\big(u^M_{(r)}t\big)_{(m+n-r)}\varphi .
\end{equation}
\end{btheorem}
{\raggedright\noindent\textbf{Lean.} \lean{VOA.}\allowbreak\lean{SkewIntertwiner.}\allowbreak\lean{wMode_comm}\par}
{\raggedright\noindent File \leanpath{VOA/\allowbreak Moonshine/\allowbreak SkewIntertwiner.lean}.\par}
\begin{leancode}
theorem wMode_comm (A : WeakModule W M) (u : 𝒱) (m : ℤ) (t : M) (n : ℤ) (φ : 𝒱) (J : ℕ)
    (hJ : ∀ j : ℕ, J ≤ j → A.act u j t = 0) :
    A.act u m (wMode A t n φ) - wMode A t n (W.smul u m φ)
      = ∑ r ∈ Finset.range J, gBinom m r • wMode A (A.act u r t) (m + n - r) φ
\end{leancode}
\noindent\textbf{Hypotheses / status.} $J$ past truncation\par
\medskip
\begin{btheorem}{37.17}[$W$ is local with every module vertex operator]
For $u\in\chnineV$ and $t\in M$, let $J$ be past the truncation of $u$ on $t$. Then
\begin{equation*}
(z_1-z_2)^J\Big(Y_M(u,z_1)\,W(t,z_2)-W(t,z_2)\,Y(u,z_1)\Big)=0
\end{equation*}
as a two-variable distribution with values in $\operatorname{Hom}(\chnineV,M)$. In particular such
an exponent exists.
\end{btheorem}
{\raggedright\noindent\textbf{Lean.} \lean{VOA.}\allowbreak\lean{SkewIntertwiner.}\allowbreak\lean{wBracket},\allowbreak\ \lean{VOA.}\allowbreak\lean{SkewIntertwiner.}\allowbreak\lean{Wsk_local_of_bound},\allowbreak\ \lean{VOA.}\allowbreak\lean{SkewIntertwiner.}\allowbreak\lean{Wsk_local}\par}
{\raggedright\noindent File \leanpath{VOA/\allowbreak Moonshine/\allowbreak SkewIntertwiner.lean}.\par}
\begin{leancode}
noncomputable def wBracket (A : WeakModule W M) (u : 𝒱) (t : M) : FormalDist₂ (𝒱 →ₗ[ℂ] M) :=
  fun a b => (A.YM u a).comp (Wsk A t b) - (Wsk A t b).comp (W.Y u a)

theorem Wsk_local_of_bound (A : WeakModule W M) (u : 𝒱) (t : M) (J : ℕ)
    (hJ : ∀ j : ℕ, J ≤ j → A.act u j t = 0) :
    mulBySubPow J (wBracket A u t) = 0

theorem Wsk_local (A : WeakModule W M) (u : 𝒱) (t : M) :
    ∃ N : ℕ, mulBySubPow N (wBracket A u t) = 0
\end{leancode}
\noindent\textbf{Hypotheses / status.} None\par
\medskip
\begin{btheorem}{37.51}[Intertwining commutator formula for $\chnineWb$]
Assume: $B_V$ graded orthogonal, graded nondegenerate and invariant; $B_M$ invariant at every
homogeneous state; and $L_1\chnineV_1=0$. Then for \emph{every} $u\in\chnineV$, $m,n\in\Z$, $t,a\in M$
and $J$ past the truncation of $u$ on $t$,
\begin{equation}
u_{(m)}\chnineWb(t)_{(n)}a-\chnineWb(t)_{(n)}\big(u_{(m)}a\big)=\sum_{r<J}\binom mr\,\chnineWb\big(u_{(r)}t\big)_{(m+n-r)}a .
\end{equation}
\end{btheorem}
{\raggedright\noindent\textbf{Lean.} \lean{VOA.}\allowbreak\lean{Moonshine.}\allowbreak\lean{SkewAdjoint.}\allowbreak\lean{WbarMode_comm}\par}
{\raggedright\noindent File \leanpath{VOA/\allowbreak Moonshine/\allowbreak WbarComm.lean}.\par}
\begin{leancode}
theorem WbarMode_comm (hVo : GradedOrth W BV) (hVnd : GradedNondeg W BV) (hVinv : VInv W BV)
    (hMinv : ∀ (wt : ℤ) (u : 𝒱), u ∈ W.grading wt → FHLInv N BM wt u)
    (hqpV : ∀ v, v ∈ W.grading 1 → W.L 1 v = 0) (u : 𝒱) (m : ℤ) (t : M) (n : ℤ) (a : M)
    (J : ℕ) (hJ : ∀ j : ℕ, J ≤ j → N.act u j t = 0) :
    W.smul u m (WbarMode BV N BM hO n t a) - WbarMode BV N BM hO n t (N.act u m a)
      = ∑ r ∈ Finset.range J, gBinom m r • WbarMode BV N BM hO (m + n - r) (N.act u r t) a
\end{leancode}
\noindent\textbf{Hypotheses / status.} \lean{hVo,hVnd,hVinv,hMinv,hqpV}\par
\medskip
\begin{btheorem}{37.52}[Locality of $Y(u,z_1)$ and $\chnineWb(t,z_2)$]
Under the hypotheses of Theorem~37.51, for every $u\in\chnineV$, $t\in M$ there is $N$
with $(z_1-z_2)^N\big(Y(u,z_1)\chnineWb(t,z_2)-\chnineWb(t,z_2)Y_M(u,z_1)\big)=0$ in
$\operatorname{Hom}(M,\chnineV)[[z_1^{\pm1},z_2^{\pm1}]]$. Moreover
$[L_{-1},\chnineWb(t,z)]=\partial_z\chnineWb(t,z)$, in coefficients
$L_{-1}\circ\chnineWb(t)_k-\chnineWb(t)_k\circ L_{-1}^M=(k+1)\chnineWb(t)_{k+1}$.
\end{btheorem}
{\raggedright\noindent\textbf{Lean.} \lean{VOA.}\allowbreak\lean{Moonshine.}\allowbreak\lean{SkewAdjoint.}\allowbreak\lean{Wbar_local},\allowbreak\ \lean{VOA.}\allowbreak\lean{Moonshine.}\allowbreak\lean{SkewAdjoint.}\allowbreak\lean{Wbar_homCov}\par}
{\raggedright\noindent File \leanpath{VOA/\allowbreak Moonshine/\allowbreak WbarComm.lean}.\par}
\begin{leancode}
theorem Wbar_local (hVo : GradedOrth W BV) (hVnd : GradedNondeg W BV)
    (hVinv : VInv W BV) (hMinv : ∀ (wt : ℤ) (u : 𝒱), u ∈ W.grading wt → FHLInv N BM wt u)
    (hqpV : ∀ v, v ∈ W.grading 1 → W.L 1 v = 0) (u : 𝒱) (t : M) :
    ∃ N₀ : ℕ, mulBySubPow N₀
      (fun a b => (W.Y u a).comp (Wbar BV N BM hO t b) - (Wbar BV N BM hO t b).comp (N.YM u a) :
        FormalDist₂ (M →ₗ[ℂ] 𝒱)) = 0

theorem Wbar_homCov (hVo : GradedOrth W BV) (hVnd : GradedNondeg W BV)
    (hVinv : VInv W BV) (hMinv : ∀ (wt : ℤ) (u : 𝒱), u ∈ W.grading wt → FHLInv N BM wt u)
    (hqpV : ∀ v, v ∈ W.grading 1 → W.L 1 v = 0) (t : M) (k : ℤ) :
    W.L (-1) ∘ₗ Wbar BV N BM hO t k - Wbar BV N BM hO t k ∘ₗ Lm1 N.toWeakModule
      = (k + 1 : ℤ) • Wbar BV N BM hO t (k + 1)
\end{leancode}
\noindent\textbf{Hypotheses / status.} Same\par
\medskip
\begin{btheorem}{37.57}[The vacuum pairing]
\textup{(i)} For $t,t'\in(\chnineTP)_2$ (the lowest weight space),
$\chnineWb(t)_{(3)}t'=B_T(t',t)\,\mathbf1$.
\textup{(ii)} For every nonzero $t\in(\chnineTP)_2$ there are $t'\in(\chnineTP)_2$ and $c\neq0$ with
$\chnineWb(t)_{(3)}t'=c\,\mathbf1$.
\textup{(iii)} $\chnineWb\neq0$.
\end{btheorem}
{\raggedright\noindent\textbf{Lean.} \lean{VOA.}\allowbreak\lean{Moonshine.}\allowbreak\lean{leechWbarMode_three_weight_two},\allowbreak\ \lean{VOA.}\allowbreak\lean{Moonshine.}\allowbreak\lean{leechWbar_vacuum_pairing_ne_zero},\allowbreak\ \lean{VOA.}\allowbreak\lean{Moonshine.}\allowbreak\lean{leechWbar_ne_zero}\par}
\begin{leancode}
theorem leechWbarMode_three_weight_two {t t' : TP} (ht : t ∈ plusModule.weight ((2 : ℤ) : ℂ))
    (ht' : t' ∈ plusModule.weight ((2 : ℤ) : ℂ)) :
    leechWbarMode 3 t t' = TPForm t' t • leechVOAPlus.vacuum

theorem leechWbar_vacuum_pairing_ne_zero {t : TP} (ht : t ∈ plusModule.weight ((2 : ℤ) : ℂ))
    (ht0 : t ≠ 0) :
    ∃ t' ∈ plusModule.weight ((2 : ℤ) : ℂ), ∃ c : ℂ, c ≠ 0 ∧
      leechWbarMode 3 t t' = c • leechVOAPlus.vacuum

theorem leechWbar_ne_zero : leechWbar ≠ 0
\end{leancode}
\noindent File \leanpath{VOA/\allowbreak Moonshine/\allowbreak VnatGenuine.lean}.
\noindent\textbf{Hypotheses / status.} $t,t'$ of weight 2\par
\medskip
\begin{btheorem}{37.58}[Skew-symmetry of $\chnineWb$]
For homogeneous $t_1\in(\chnineTP)_a$, $t_2\in(\chnineTP)_b$ and every $n\in\Z$,
\begin{equation*}
\chnineWb(t_1)_{(n)}t_2=\sum_{i\ge0}\frac{(-1)^{n+i+1}}{i!}\,L_{-1}^{\,i}\Big(\chnineWb(t_2)_{(n+i)}t_1\Big).
\end{equation*}
\end{btheorem}
{\raggedright\noindent\textbf{Lean.} \lean{VOA.}\allowbreak\lean{Moonshine.}\allowbreak\lean{M8h0.}\allowbreak\lean{TPForm_Lm1},\allowbreak\ \lean{VOA.}\allowbreak\lean{Moonshine.}\allowbreak\lean{M8h0.}\allowbreak\lean{leechBV_Lm1_pow},\allowbreak\ \lean{VOA.}\allowbreak\lean{Moonshine.}\allowbreak\lean{M8h0.}\allowbreak\lean{leechWbar_skew}\par}
\begin{leancode}
theorem TPForm_Lm1 (x y : TP) : TPForm (leechLTP x) y = TPForm x (Lp1 plusModule y)

theorem leechBV_Lm1_pow (i : ℕ) (x y : SkewVP) :
    leechBV (((leechVOAPlus.L (-1)) ^ i) x) y = leechBV x (((leechVOAPlus.L 1) ^ i) y)

theorem leechWbar_skew {a b : ℤ} {t₁ t₂ : TP} (ht₁ : t₁ ∈ plusModule.weight (a : ℂ))
    (ht₂ : t₂ ∈ plusModule.weight (b : ℂ)) (n : ℤ) :
    leechWbarMode n t₁ t₂ = ∑ᶠ i : ℕ, (((i.factorial : ℂ)⁻¹) * negOnePow (n + i + 1)) •
      ((leechVOAPlus.L (-1)) ^ i) (leechWbarMode (n + i) t₂ t₁)
\end{leancode}
\noindent File \leanpath{VOA/\allowbreak Moonshine/\allowbreak M8h0Skew.lean}.
\noindent\textbf{Hypotheses / status.} $t_1,t_2$ homogeneous\par
\medskip
\endgroup
\chapter{The vertex algebra on \texorpdfstring{$V^\natural$}{V natural} modulo one relation}
\begingroup
\providecommand{\chtenWb}{\overline{W}}
\providecommand{\chtenTP}{T^{+}}
\providecommand{\chtenVP}{V_{\Lambda}^{+}}
\providecommand{\chtenB}{B^{\natural}}
\noindent{\small 8 of the chapter's 65 numbered statements: its headline results and every definition they use.}\par\medskip
\begin{btheorem}{38.16}[Diagonal--off-diagonal locality]
$\operatorname{diag}(a(z_1),b(z_1))$ and $\operatorname{offdiag}(c'(z_2),d'(z_2))$ are mutually
local provided that, for some $N$,
\begin{equation*}
  (z_1-z_2)^N\bigl(a(z_1)c'(z_2)-c'(z_2)b(z_1)\bigr)=0
  \quad\text{and}\quad
  (z_1-z_2)^N\bigl(b(z_1)d'(z_2)-d'(z_2)a(z_1)\bigr)=0 .
\end{equation*}
\end{btheorem}
{\raggedright\noindent\textbf{Lean.} \lean{VOA.}\allowbreak\lean{Moonshine.}\allowbreak\lean{VnatBlock.}\allowbreak\lean{local_diag_off}\par}
{\raggedright\noindent File \leanpath{VOA/\allowbreak Moonshine/\allowbreak VnatBlock.lean}.\par}
\begin{leancode}
theorem local_diag_off {a : FormalDist (End ℂ V)} {b : FormalDist (End ℂ M)}
    {c' : FormalDist (M →ₗ[ℂ] V)} {d' : FormalDist (V →ₗ[ℂ] M)}
    (hc : ∃ N : ℕ, mulBySubPow N
      ((fun m n => a m ∘ₗ c' n - c' n ∘ₗ b m) : FormalDist₂ (M →ₗ[ℂ] V)) = 0)
    (hd : ∃ N : ℕ, mulBySubPow N
      ((fun m n => b m ∘ₗ d' n - d' n ∘ₗ a m) : FormalDist₂ (V →ₗ[ℂ] M)) = 0) :
    Local (diagF a b) (offF c' d') := by
\end{leancode}
\noindent\textbf{Hypotheses / status.} Two block localities\par
\medskip
\begin{btheorem}{38.19}[Locality of module fields]
For every weak module $M$ of $W$ and all $u,v\in W$, the fields $Y_M(u,z)$ and $Y_M(v,z)$ are
mutually local.
\end{btheorem}
{\raggedright\noindent\textbf{Lean.} \lean{VOA.}\allowbreak\lean{Moonshine.}\allowbreak\lean{VnatBlock.}\allowbreak\lean{local_YM}\par}
{\raggedright\noindent File \leanpath{VOA/\allowbreak Moonshine/\allowbreak VnatBlock.lean}.\par}
\begin{leancode}
theorem local_YM (A : WeakModule W M) (u v : 𝒱) : Local (A.YM u) (A.YM v) := by
\end{leancode}
\noindent\textbf{Hypotheses / status.} None\par
\medskip
\begin{btheorem}{38.41}[The VOA $\Vnat$, conditional on twisted--twisted locality]
Assume $Y^\natural_T(t)\sim Y^\natural_T(t')$ for all $t,t'\in\chtenTP$.  Then
$(\Vnat,\ (\Vnat_n)_n,\ Y^\natural,\ \mathbf1^\natural,\ \omega^\natural,\ c=24)$ is a vertex
operator algebra in the sense of the library.
\end{btheorem}
{\raggedright\noindent\textbf{Lean.} \lean{VOA.}\allowbreak\lean{Moonshine.}\allowbreak\lean{VnatAxioms.}\allowbreak\lean{vnatVOA},\allowbreak\ \lean{VOA.}\allowbreak\lean{Moonshine.}\allowbreak\lean{VnatAxioms.}\allowbreak\lean{vnat_VOA_of_htt}\par}
{\raggedright\noindent File \leanpath{VOA/\allowbreak Moonshine/\allowbreak VnatAxioms.lean}.\par}
\begin{leancode}
noncomputable def vnatVOA (htt : ∀ t t' : TP, Local (YnatT leechWbar t) (YnatT leechWbar t')) :
    VertexOperatorAlgebra Vnat where
  grading := gradingNat
  isInternal := gradingNat_isInternal
  finrank := finiteDimensional_gradingNat
  Y := YN
  vacuum := vacNat
  vacuum_mem := vacNat_mem
  grading_mode := fun x hx y hy n => YN_grading_mode x hx y hy n
  vacuum_Y := YN_vacuum
  creation := YN_creation
  locality := YN_local_of_htt htt
  regularity := YN_regularity
  conformalVector := omegaNat
  conformal_mem := omegaNat_mem
  centralCharge := 24
  virasoro := YN_virasoro
  L0_grading := fun v hv => YN_L0_grading v hv
  translation := YN_translation
  vacuum_grade := gradingNat_zero
  positivity := fun _ hn => gradingNat_neg hn

theorem vnat_VOA_of_htt (htt : ∀ t t' : TP, Local (YnatT leechWbar t) (YnatT leechWbar t')) :
    ∃ V : VertexOperatorAlgebra Vnat,
      V.Y = Ynat wbarL ∧ V.vacuum = vacNat ∧ V.conformalVector = omegaNat ∧
        V.centralCharge = 24 ∧ V.grading = gradingNat :=
  ⟨vnatVOA htt, rfl, rfl, rfl, rfl, rfl⟩
\end{leancode}
\noindent\textbf{Hypotheses / status.} Htt\par
\medskip
\begin{btheorem}{38.44}
$Y^\natural_T[\overline W'](t_1)\sim Y^\natural_T[\overline W'](t_2)$ if and only if
$\mathrm{VP}(t_1,t_2)$ and $\mathrm{TP}(t_1,t_2)$ hold.
\end{btheorem}
{\raggedright\noindent\textbf{Lean.} \lean{VOA.}\allowbreak\lean{Moonshine.}\allowbreak\lean{YnatT_local_iff}\par}
\begin{leancode}
theorem YnatT_local_iff (Wbar : TP → FormalDist (TP →ₗ[ℂ] SkewVP)) (t₁ t₂ : TP) :
    Local (YnatT Wbar t₁) (YnatT Wbar t₂) ↔
      VPBlockLocal Wbar t₁ t₂ ∧ TPBlockLocal Wbar t₁ t₂ :=
  local_off_off_iff _ _ _ _
\end{leancode}
\noindent File \leanpath{VOA/\allowbreak Moonshine/\allowbreak TTBlock.lean}.
\noindent\textbf{Hypotheses / status.} None\par
\medskip
\begin{btheorem}{38.45}[Two-variable Goddard lemma]
Let $V$ be a vector space, $\mathrm{vac}\in V$, $A(w)\in\End(V)[[w^\pm]]$ with $A_p\,\mathrm{vac}=0$ for
$p<0$ (raw indexing: no negative powers of $w$ on the vacuum), and
$B(z_1,z_2)\in\End(V)[[z_1^\pm,z_2^\pm]]$ with $B_{m,n}\mathrm{vac}=0$ for all $m,n$.  If
$(w-z_1)^{M_1}(w-z_2)^{M_2}[A(w),B(z_1,z_2)]=0$, then $B_{m,n}(A_p\,\mathrm{vac})=0$ for all $p,m,n$.
\end{btheorem}
{\raggedright\noindent\textbf{Lean.} \lean{VOA.}\allowbreak\lean{Moonshine.}\allowbreak\lean{M8h0.}\allowbreak\lean{goddard₂}\par}
\begin{leancode}
theorem goddard₂ (vac : V) (A : FormalDist (End ℂ V)) (B : FormalDist₂ (End ℂ V)) (M₁ M₂ : ℕ)
    (hloc : mulBySub₁₂Pow M₁ (mulBySub₁₃Pow M₂ (brk3 A B)) = 0)
    (hneg : ∀ p : ℤ, p < 0 → A p vac = 0) (hvac : ∀ m n, B m n vac = 0) (p m n : ℤ) :
    B m n (A p vac) = 0 := by
\end{leancode}
\noindent File \leanpath{VOA/\allowbreak Moonshine/\allowbreak M8h0Goddard.lean}.
\noindent\textbf{Hypotheses / status.} Displayed\par
\medskip
\begin{btheorem}{38.55}[Four explicit pairs]
There exist $x_0\in T^+_0$, $x_1\in T^+_1$ (the two frame sectors of $\chtenTP$) such that (4.23)
on the four pairs $(x_i,x_j)$ implies the conclusion of Theorem~38.41.
\end{btheorem}
{\raggedright\noindent\textbf{Lean.} \lean{VOA.}\allowbreak\lean{Moonshine.}\allowbreak\lean{M8h0.}\allowbreak\lean{vnat_VOA_of_TPBlockLocal_frame}\par}
{\raggedright\noindent File \leanpath{VOA/\allowbreak Moonshine/\allowbreak M8h0Gap.lean}.\par}
\begin{leancode}
theorem vnat_VOA_of_TPBlockLocal_frame :
    ∃ x : ZMod 2 → TP, (∀ s, x s ∈ TPframe s) ∧
      ((∀ i j, TPBlockLocal leechWbar (x i) (x j)) →
        ∃ V : VertexOperatorAlgebra Vnat,
          V.Y = Ynat VnatAxioms.wbarL ∧ V.vacuum = vacNat ∧ V.conformalVector = omegaNat ∧
            V.centralCharge = 24 ∧ V.grading = gradingNat) := by
\end{leancode}
\noindent\textbf{Hypotheses / status.} (4.23) on 4 chosen pairs\par
\medskip
\begin{btheorem}{38.66}[$\Vnat$ is simple]
Assume htt.  Then $\mathtt{vnatVOA}(\mathrm{htt})$ is simple, and the canonical invariant form of the
library on it coincides with $\chtenB$ and is nondegenerate on every weight space.
\end{btheorem}
{\raggedright\noindent\textbf{Lean.} \lean{VOA.}\allowbreak\lean{Moonshine.}\allowbreak\lean{SimpleNat.}\allowbreak\lean{vnatVOA_isSimple},\allowbreak\ \lean{VOA.}\allowbreak\lean{Moonshine.}\allowbreak\lean{SimpleNat.}\allowbreak\lean{vnatVOA_invariantForm_nondegenerate},\allowbreak\ \lean{VOA.}\allowbreak\lean{Moonshine.}\allowbreak\lean{SimpleNat.}\allowbreak\lean{vnatVOA_canonicalForm_eq}\par}
\begin{leancode}
theorem vnatVOA_isSimple : (vnatVOA htt).IsSimple :=
  (LiRadical.isSimple_iff_leftRadical_eq_bot (natInvForm_isFullyInvariant htt)
    (natInvForm_normalized htt)).2 (natInvForm_leftRadical htt)

theorem vnatVOA_invariantForm_nondegenerate :
    (natInvForm htt).IsFullyInvariant ∧ (natInvForm htt).Normalized ∧
      InvariantForm.leftRadical (natInvForm htt) = ⊥ ∧
        GradedNondeg (vnatVOA htt) (natInvForm htt).form :=

theorem vnatVOA_canonicalForm_eq :
    (canonicalInvariantForm (vnatVOA_vacuum_ne_zero htt) (vnatVOA_L1_one htt)).form = natForm :=
\end{leancode}
\noindent File \leanpath{VOA/\allowbreak Moonshine/\allowbreak SimpleVnat.lean}.
\noindent\textbf{Hypotheses / status.} Htt\par
\medskip
\begin{btheorem}{38.68}[$\Vnat$ modulo (4.23) on four pairs]
There exist $x_0\in T^+_0$, $x_1\in T^+_1$ such that, if DGM (4.23) holds on the four pairs
$(x_i,x_j)$, then there is a simple vertex operator algebra $V$ on $\Vnat$ whose vertex map is the
block map $Y^\natural[\chtenWb]$ of (38.1), with vacuum $(\mathbf 1,0)$, conformal
vector $(\omega,0)$, central charge $24$, $V_1=0$, $\dim V_2=196884$, and $\dim V_n=[q^n]\,qJ(q)$
for all $n\ge0$.
\end{btheorem}
{\raggedright\noindent\textbf{Lean.} \lean{VOA.}\allowbreak\lean{Moonshine.}\allowbreak\lean{moonshine_of_frameLocality}\par}
\begin{leancode}
theorem moonshine_of_frameLocality :
    ∃ x : ZMod 2 → TP, (∀ s, x s ∈ TPframe s) ∧
      ((∀ i j, TPBlockLocal leechWbar (x i) (x j)) →
        ∃ V : VertexOperatorAlgebra Vnat,
          V.Y = Ynat wbarL ∧ V.vacuum = vacNat ∧ V.conformalVector = omegaNat ∧
          V.centralCharge = 24 ∧
          V.grading 1 = ⊥ ∧
          finrank ℂ (V.grading 2) = 196884 ∧
          (∀ n : ℕ, (finrank ℂ (V.grading (n : ℤ)) : ℂ) = PowerSeries.coeff n JShift) ∧
          V.IsSimple) := by
\end{leancode}
\noindent File \leanpath{VOA/\allowbreak Moonshine/\allowbreak Main.lean}.
\noindent\textbf{Hypotheses / status.} (4.23) on 4 chosen pairs\par
\medskip
\endgroup
\chapter{The Golay frame and the untwisted triality}
\begingroup
\providecommand{\chelLz}{\Lambda_0}
\providecommand{\chelVz}{V_0^{+}}
\providecommand{\chelVt}{V_3^{+}}
\providecommand{\chelN}{N}
\noindent{\small 6 of the chapter's 42 numbered statements: its headline results and every definition they use.}\par\medskip
\begin{btheorem}{39.6}[$V_\Leech^+=\chelVz\oplus\chelVt$, and the $\Z/2$-grading of products]
(i) $V_\Leech^+=V_{(0)}\oplus V_{(1)}$ (internal direct sum).
(ii) For $u\in V_{(i)}$, $v\in V_{(j)}$ and all $n\in\Z$: $u_{(n)}v\in V_{(i+j)}$.
(iii) Each $V_{(i)}$ is stable under the projections onto the $L_0$-weight spaces.
\end{btheorem}
{\raggedright\noindent\textbf{Lean.} \lean{VOA.}\allowbreak\lean{Moonshine.}\allowbreak\lean{frameSector_isInternal},\allowbreak\ \lean{VOA.}\allowbreak\lean{Moonshine.}\allowbreak\lean{frameSector_smul_mem},\allowbreak\ \lean{VOA.}\allowbreak\lean{Moonshine.}\allowbreak\lean{frameSector_graded}\par}
{\raggedright\noindent File \leanpath{VOA/\allowbreak Moonshine/\allowbreak FrameSectors.lean}.\par}
\begin{leancode}
/-- **`V_Λ⁺ = V₀⁺ ⊕ V₃⁺`** (internal direct sum). -/
theorem frameSector_isInternal : DirectSum.IsInternal frameSector

/-- **The `ℤ/2`-grading of products**: `u ∈ V_i`, `v ∈ V_j ⇒ u₍ₙ₎v ∈ V_{i+j}`. -/
theorem frameSector_smul_mem {i j : ZMod 2} {u v : VPlus} (hu : u ∈ frameSector i)
    (hv : v ∈ frameSector j) (n : ℤ) : leechVOAPlus.smul u n v ∈ frameSector (i + j)

/-- The sectors are graded: each is stable under the weight projections of `V_Λ⁺`. -/
theorem frameSector_graded (i : ZMod 2) (n : ℤ) {v : VPlus} (hv : v ∈ frameSector i) :
    piHom leechVOAPlus n v ∈ frameSector i
\end{leancode}
\medskip
\begin{btheorem}{39.10}[$98304$ norm-four vectors outside $\Lambda_0$]
$\#\{\alpha\in\Leech\setminus\Lambda_0:\langle\alpha,\alpha\rangle=4\}=24\cdot2^{12}=98304$.
\end{btheorem}
{\raggedright\noindent\textbf{Lean.} \lean{VOA.}\allowbreak\lean{Moonshine.}\allowbreak\lean{oddFour},\allowbreak\ \lean{VOA.}\allowbreak\lean{Moonshine.}\allowbreak\lean{card_normFour_frameChar_one}\par}
\begin{leancode}
/-- The norm-`4` vectors outside `Λ₀`. -/
noncomputable def oddFour : Finset leech.carrier :=
  (normFour leech leech_posDef).filter fun α => frameChar α = 1

/-- **`#{α ∈ Λ \ Λ₀ | ⟨α,α⟩ = 4} = 98304`.** -/
theorem card_normFour_frameChar_one : oddFour.card = 98304
\end{leancode}
\noindent File \leanpath{VOA/\allowbreak Moonshine/\allowbreak FrameCount.lean}.
\medskip
\begin{btheorem}{39.15}[The Leech frame split of $T$]
For $\Leech$, $w=8e_0$: $T=T_+\oplus T_-$ with $\dim T_+=\dim T_-=2048$; $e_\alpha$ preserves $T_\pm$
for $\alpha\in\Lambda_0$ and swaps them for $\alpha\notin\Lambda_0$; and on the twisted space
$S[\theta]\otimes T$ the operator $\mathrm{id}\otimes J$ gives
$S[\theta]\otimes T=(S[\theta]\otimes T_+)\oplus(S[\theta]\otimes T_-)$.
\end{btheorem}
{\raggedright\noindent\textbf{Lean.} \lean{VOA.}\allowbreak\lean{Moonshine.}\allowbreak\lean{frameW},\allowbreak\ \lean{VOA.}\allowbreak\lean{Moonshine.}\allowbreak\lean{frameW_sign},\allowbreak\ \lean{VOA.}\allowbreak\lean{Moonshine.}\allowbreak\lean{finrank_leechFrameT_plus},\allowbreak\ \lean{VOA.}\allowbreak\lean{Moonshine.}\allowbreak\lean{finrank_leechFrameT_minus},\allowbreak\ \lean{VOA.}\allowbreak\lean{Moonshine.}\allowbreak\lean{eT_mem_leechFrameT},\allowbreak\ \lean{VOA.}\allowbreak\lean{Moonshine.}\allowbreak\lean{leechTw_frame_isCompl}\par}
\begin{leancode}
/-- The frame vector `w = 8e₀ ∈ Λ₀`. -/
noncomputable def frameW : leech.carrier := (⟨_, eightE0_mem⟩ : ↥Leech.carrier)

theorem frameW_sign (α : leech.carrier) :
    (-1 : ℂ) ^ (leech.form frameW α) = chiSign leech frameChar α

/-- **`dim T₊ = 2048`.** -/
theorem finrank_leechFrameT_plus : finrank ℂ (leechFrameT 0) = 2048

/-- **`dim T₋ = 2048`.** -/
theorem finrank_leechFrameT_minus : finrank ℂ (leechFrameT 1) = 2048

/-- `e_α` preserves `T_±` for `α ∈ Λ₀` and swaps them for `α ∉ Λ₀`. -/
theorem eT_mem_leechFrameT (α : leech.carrier) {s : ZMod 2}
    {x : TT leech leech_even leech_isUnit_det_gramMat} (hx : x ∈ leechFrameT s) :
    eT leech leech_even leech_isUnit_det_gramMat α x ∈ leechFrameT (s + frameChar α)

/-- The split of the twisted carrier `S[θ] ⊗ T = (S[θ] ⊗ T₊) ⊕ (S[θ] ⊗ T₋)`. -/
theorem leechTw_frame_isCompl :
    IsCompl (twFrame leech leech_even leech_isUnit_det_gramMat (TwistedHeisenberg.FockT 24) frameW 0)
      (twFrame leech leech_even leech_isUnit_det_gramMat (TwistedHeisenberg.FockT 24) frameW 1)
\end{leancode}
\noindent File \leanpath{VOA/\allowbreak Moonshine/\allowbreak FrameTwisted.lean}.
\medskip
\begin{btheorem}{39.19}[$\chelVz$ is cyclic over itself]
$\chelVz$ is cyclic with cyclic vector $\mathbf 1$.
\end{btheorem}
{\raggedright\noindent\textbf{Lean.} \lean{VOA.}\allowbreak\lean{Moonshine.}\allowbreak\lean{cyclic_V0}\par}
{\raggedright\noindent File \leanpath{VOA/\allowbreak Moonshine/\allowbreak CyclicFrameSectors.lean}.\par}
\begin{leancode}
/-- **`V₀⁺` is cyclic over itself**, generated by `𝟙`. -/
theorem cyclic_V0 : ∃ x₀ ∈ frameSector 0,
    Submodule.span ℂ {y | ∃ u ∈ frameSector 0, ∃ n : ℤ, y = leechVOAPlus.smul u n x₀}
      = frameSector 0
\end{leancode}
\medskip
\begin{btheorem}{39.28}[Each twisted frame sector is cyclic]
Let $0\ne t_0\in T_{(s)}$ and $x_0=h_0(-\tfrac12)\mathbf 1\otimes t_0$ (weight $\tfrac32+\tfrac12=2$).
Then $x_0\in T_{(s)}^+$ and
$\operatorname{span}\{u_{(n)}x_0 : u\in\chelVz,\ n\in\Z\}=T_{(s)}^+$, where $u_{(n)}$ is the action of
$V_\Leech^+$ on $(V_\Leech^T)^+$.  In particular each $T_{(s)}^+$ is a cyclic $\chelVz$-module.
\end{btheorem}
{\raggedright\noindent\textbf{Lean.} \lean{VOA.}\allowbreak\lean{Moonshine.}\allowbreak\lean{frameX0T},\allowbreak\ \lean{VOA.}\allowbreak\lean{Moonshine.}\allowbreak\lean{span_smul_frameX0T},\allowbreak\ \lean{VOA.}\allowbreak\lean{Moonshine.}\allowbreak\lean{cyclic_TPframe}\par}
\begin{leancode}
/-- The cyclic vector `h₀(−½)𝟙 ⊗ t₀` (weight `2`). -/
noncomputable def frameX0T (t₀ : TLeechT) : Tw := (X (0, 0) : FockT 24) ⊗ₜ t₀

/-- **The explicit form**: for `0 ≠ t₀ ∈ T_s`, `x₀ = h₀(−½)𝟙 ⊗ t₀` generates `T_s⁺` over `V₀⁺`. -/
theorem span_smul_frameX0T {s : ZMod 2} {t₀ : TLeechT} (ht₀ : t₀ ∈ leechFrameT s) (h0 : t₀ ≠ 0) :
    Submodule.span ℂ {y | ∃ u ∈ frameSector 0, ∃ n : ℤ,
        y = mode (Moonshine.plusModule.YM u) n (⟨frameX0T t₀, frameX0T_mem_TP t₀⟩ : TP)}
      = TPframe s

/-- **M8f, twisted half: each twisted frame sector `T_s⁺` is a cyclic `V₀⁺`-module.** -/
theorem cyclic_TPframe (s : ZMod 2) : ∃ x₀ ∈ TPframe s,
    Submodule.span ℂ {y | ∃ u ∈ frameSector 0, ∃ n : ℤ, y = mode (Moonshine.plusModule.YM u) n x₀}
      = TPframe s
\end{leancode}
\noindent File \leanpath{VOA/\allowbreak Moonshine/\allowbreak CyclicTwisted.lean}.
\medskip
\begin{btheorem}{39.41}[$V_{\Lambda_0}^+\cong\chelVz$]
\lean{bridgeΛ} induces a linear bijection $V_{\Lambda_0}^+=V_{\Lambda_0}^{\langle\theta\rangle}\to\chelVz$
which is a VOA homomorphism preserving $\omega$.
\end{btheorem}
{\raggedright\noindent\textbf{Lean.} \lean{VOA.}\allowbreak\lean{Moonshine.}\allowbreak\lean{lam0PlusEquiv},\allowbreak\ \lean{VOA.}\allowbreak\lean{Moonshine.}\allowbreak\lean{lam0PlusEquiv_isVOAHom},\allowbreak\ \lean{VOA.}\allowbreak\lean{Moonshine.}\allowbreak\lean{lam0PlusEquiv_omega}\par}
\begin{leancode}
/-- **`V_{Λ₀}⁺ ≃ V₀⁺`** as vector spaces (a VOA isomorphism by `lam0PlusEquiv_isVOAHom`). -/
noncomputable def lam0PlusEquiv : Lam0PlusCarrier ≃ₗ[ℂ] V0Sub.carrier :=
  LinearEquiv.ofBijective lam0PlusToV0 ⟨lam0PlusToV0_injective, lam0PlusToV0_surjective⟩

/-- `lam0PlusEquiv` is a homomorphism of VOAs `V_{Λ₀}⁺ → V₀⁺`. -/
theorem lam0PlusEquiv_isVOAHom : IsVOAHom lam0Plus V0P lam0PlusEquiv.toLinearMap

theorem lam0PlusEquiv_omega : lam0PlusEquiv lam0Plus.conformalVector = V0P.conformalVector
\end{leancode}
\noindent File \leanpath{VOA/\allowbreak Moonshine/\allowbreak FrameBridge.lean}.
\medskip
\endgroup
\chapter{Existence of the twisted triality I: reduction to Zhu generation}
\begingroup
\providecommand{\chtwelveO}{\mathcal{O}}
\providecommand{\chtwelveG}{\mathcal{G}}
\noindent{\small 6 of the chapter's 61 numbered statements: its headline results and every definition they use.}\par\medskip
\begin{blemma}{40.3}[What $\chtwelveO(V)$ contains]
Let $a\in V_p$ with $p\in\N$, $b\in V$, and $y\in V_w$.
\begin{enumerate}[label=(\alph*)]
\item $L_{-1}y+L_0y\in\chtwelveO(V)$; hence $L_{-1}y+w\,y\in\chtwelveO(V)$.
\item (Frenkel--Zhu higher residues) For every $k\in\N$,
  \begin{equation*}
    \mathrm{Pres}(a,p,b,k)\;:=\;\sum_{i=0}^{p+k}\binom{p+k}{i}\,a_{(i-2-k)}b
    \;=\;\Res_zY(a,z)b\,\frac{(1+z)^{p+k}}{z^{2+k}}\;\in\;\chtwelveO(V).
  \end{equation*}
\end{enumerate}
\end{blemma}
{\raggedright\noindent\textbf{Lean.} \lean{VOA.}\allowbreak\lean{Moonshine.}\allowbreak\lean{ZhuGen.}\allowbreak\lean{L_neg_one_add_mem_O},\allowbreak\ \lean{VOA.}\allowbreak\lean{Moonshine.}\allowbreak\lean{ZhuGen.}\allowbreak\lean{L_neg_one_add_wt_mem_O},\allowbreak\ \lean{VOA.}\allowbreak\lean{Pres},\allowbreak\ \lean{VOA.}\allowbreak\lean{Pres_mem_O}\par}
{\raggedright\noindent File \leanpath{VOA/\allowbreak Moonshine/\allowbreak ZhuGenFiltration.lean}.\par}
\begin{leancode}
theorem L_neg_one_add_mem_O (y : 𝒱) : W.L (-1) y + W.L 0 y ∈ O W

theorem L_neg_one_add_wt_mem_O {w : ℤ} {y : 𝒱} (hy : y ∈ W.grading w) :
    W.L (-1) y + (w : ℂ) • y ∈ O W
\end{leancode}
{\raggedright\noindent File \leanpath{VOA/\allowbreak Modules/\allowbreak ZhuModule.lean}.\par}
\begin{leancode}
noncomputable def Pres (W : VertexOperatorAlgebra 𝒱) (M : ℕ) (a b : 𝒱) (n : ℕ) : 𝒱 :=
  ∑ i ∈ Finset.range (M + n + 1), cbz ((M : ℤ) + n) (i : ℤ) • W.smul a ((i : ℤ) - 2 - n) b

theorem Pres_mem_O (W : VertexOperatorAlgebra 𝒱) {M : ℕ} {a : 𝒱}
    (ha : a ∈ W.grading (M : ℤ)) (b : 𝒱) (n : ℕ) : Pres W M a b n ∈ O W
\end{leancode}
\noindent\textbf{Hypotheses / status.} $a$ homogeneous\par
\medskip
\begin{btheorem}{40.5}[Zhu's easy direction]
Let $N$ have lowest weight $h$ and $x\in N_h$.  For all $u,v\in V$: $o(u\ast v)x=o(u)o(v)x$;
$o(a)x=0$ for $a\in\chtwelveO(V)$; $o(u)x\in N_h$; and $o(\mathbf 1)=\mathrm{id}$.
\end{btheorem}
{\raggedright\noindent\textbf{Lean.} \lean{VOA.}\allowbreak\lean{oMod_zhuMul_lowest_all},\allowbreak\ \lean{VOA.}\allowbreak\lean{O_le_ker_rhoLW}\par}
{\raggedright\noindent File \leanpath{VOA/\allowbreak Modules/\allowbreak ZhuModule.lean}.\par}
\begin{leancode}
theorem oMod_zhuMul_lowest_all (N : OrdVModule W Mod) {h : ℂ} (hcw : IsLowestWeight N h)
    (u v : 𝒱) {x : Mod} (hx : x ∈ N.weight h) :
    oMod N (zhuMul W u v) x = oMod N u (oMod N v x)

theorem O_le_ker_rhoLW (N : OrdVModule W Mod) {h : ℂ} (hcw : IsLowestWeight N h) :
    O W ≤ LinearMap.ker (rhoLW N h)
\end{leancode}
\noindent\textbf{Hypotheses / status.} Lowest weight $h$\par
\medskip
\begin{btheorem}{40.19}[Existence criterion on the sectors]
For $\beta\notin\Leech_0$ and $y_0\in(V_\Leech^T)^+$: a $\sigma$-twisted $\Phi$ on $V_3^+$ with
$\Phi(x_\beta)=y_0$ exists iff $(x_\beta,y_0)$ is annihilator-compatible for $(\mathrm{act}_V,\mathrm{act}_T,\sigma)$
with all letters.  Likewise on $T_s^+$ with $x_{t_0}$.
\end{btheorem}
{\raggedright\noindent\textbf{Lean.} \lean{VOA.}\allowbreak\lean{Moonshine.}\allowbreak\lean{TrialityTw.}\allowbreak\lean{sigmaV3T_exists_iff},\allowbreak\ \lean{VOA.}\allowbreak\lean{Moonshine.}\allowbreak\lean{TrialityTw.}\allowbreak\lean{sigmaTT_exists_iff}\par}
{\raggedright\noindent File \leanpath{VOA/\allowbreak Moonshine/\allowbreak TrialityTwSigma.lean}.\par}
\begin{leancode}
theorem sigmaV3T_exists_iff {β : leech.carrier} (hβ : β ∉ lambda0) (y₀ : TP) :
    (∃ Φ : VPlus →ₗ[ℂ] TP, IsSigmaV3T Φ ∧ Φ (frameX0P β) = y₀)
      ↔ AnnCompat actV actT sig Set.univ (frameX0P β) y₀

theorem sigmaTT_exists_iff {s : ZMod 2} {t₀ : TLeechT} (ht₀ : t₀ ∈ leechFrameT s) (h0 : t₀ ≠ 0)
    (y₀ : TP) :
    (∃ Φ : TP →ₗ[ℂ] TP, IsSigmaTT s Φ ∧ Φ (x0T t₀) = y₀)
      ↔ AnnCompat actT actT sig Set.univ (x0T t₀) y₀
\end{leancode}
\noindent\textbf{Hypotheses / status.} Cyclic vector\par
\medskip
\begin{btheorem}{40.25}[Annihilator compatibility from the top level]
Let $A$, $B$ be weak $W$-modules on $M$, $N$; $\beta$ a bilinear form on $N$; $\Omega\le M$;
$\varphi:M\to N$ and $P:M\to M$ linear; $G\subseteq W$.  Assume:
\begin{enumerate}[label=(\roman*)]
\item (top-level compatibility) for every $k$, every $a\in W_k$, every $q\in\Z$ and all
  $x,x'\in\Omega$: $\beta(a_{(q)}\varphi x,\varphi x') = \beta(\varphi(P(a_{(q)}x)),\varphi x')$;
\item (invariance on $G$) for every $c\in G$ and $r\in\Z$ there are finitely many $a_i\in W$, $q_i\in\Z$,
  $t_i\in\C$ with $\beta(y,c_{(r)}y')=\beta(\sum_i t_i\,a_{i\,(q_i)}y,\,y')$ for all $y,y'\in N$;
\item (nondegeneracy) if $y$ lies in the submodule of $B$ generated by $\varphi(\Omega)$ and
  $\beta(y,z)=0$ for all $z\in Y_G$, then $y=0$.
\end{enumerate}
Then for every $x_0\in\Omega$, $(x_0,\varphi x_0)$ is annihilator-compatible for the untwisted
actions of all of $W$ on $A$ and on $B$.
\end{btheorem}
{\raggedright\noindent\textbf{Lean.} \lean{VOA.}\allowbreak\lean{Moonshine.}\allowbreak\lean{SigmaExist.}\allowbreak\lean{annCompat_of_top}\par}
{\raggedright\noindent File \leanpath{VOA/\allowbreak Moonshine/\allowbreak SigmaExistGeneric.lean}.\par}
\begin{leancode}
theorem annCompat_of_top (A : WeakModule W M) (B : WeakModule W N)
    (β : LinearMap.BilinForm ℂ N) (Ω : Submodule ℂ M) (φ : M →ₗ[ℂ] N) (P : M →ₗ[ℂ] M)
    (G : Set 𝒱)
    (hZO : ∀ (k : ℤ) (a : 𝒱), a ∈ W.grading k → ∀ (q : ℤ), ∀ x ∈ Ω, ∀ x' ∈ Ω,
      β (B.act a q (φ x)) (φ x') = β (φ (P (A.act a q x))) (φ x'))
    (hadj : ∀ c ∈ G, ∀ r : ℤ, ∃ (n : ℕ) (a : ℕ → 𝒱) (q : ℕ → ℤ) (t : ℕ → ℂ), ∀ y y' : N,
      β y (B.act c r y') = β (∑ i ∈ Finset.range n, t i • B.act (a i) (q i) y) y')
    (hnd : ∀ y ∈ modeSpan B (φ '' Ω), (∀ z ∈ testSpan B G Ω φ, β y z = 0) → y = 0)
    {x₀ : M} (hx₀ : x₀ ∈ Ω) :
    AnnCompat (fun a q => mode (A.YM a) q) (fun a q => mode (B.YM a) q) id Set.univ x₀ (φ x₀)
\end{leancode}
\noindent\textbf{Hypotheses / status.} (i)--(iii)\par
\medskip
\begin{btheorem}{40.54}[Leading term, commutators and vacuum values]
On the charge-$\gamma$ sector:
\begin{enumerate}[label=(\alph*)]
\item (leading term) if $p$ has degree $\le d$ and weight $w$, then
  $\mathrm{quadF}(0)p - X(i,a)X(j,b)p$ has degree $\le d+1$ and weight $w+i+j+2$;
\item (commutator, weight-one legs) for $s\ge-1$,
  $[\mathrm{quadF}(a,0;b,0;s),X(t,c)]=(t+1)\big(\delta_{cb}X(t+s+2,a)+\delta_{ca}X(t+s+2,b)\big)$;
\item $(h_ah_b)_{(0)}e^\gamma = \gamma_b\,h_a(-1)e^\gamma+\gamma_a\,h_b(-1)e^\gamma$;
\item $(h_a(-1)h_c(-j-1)\mathbf 1)_{(-1)}e^\gamma = \big(X(0,a)X(j,c)+(-1)^j\gamma_c X(j+1,a)+(j+1)\gamma_aX(j+1,c)\big)e^\gamma$.
\end{enumerate}
\end{btheorem}
{\raggedright\noindent\textbf{Lean.} \lean{VOA.}\allowbreak\lean{Moonshine.}\allowbreak\lean{FockSpan.}\allowbreak\lean{Fdw},\allowbreak\ \lean{VOA.}\allowbreak\lean{Moonshine.}\allowbreak\lean{FockSpan.}\allowbreak\lean{quadF_lead},\allowbreak\ \lean{VOA.}\allowbreak\lean{Moonshine.}\allowbreak\lean{FockSpan.}\allowbreak\lean{quadF_X_mul},\allowbreak\ \lean{VOA.}\allowbreak\lean{Moonshine.}\allowbreak\lean{FockSpan.}\allowbreak\lean{quadF_D_one},\allowbreak\ \lean{VOA.}\allowbreak\lean{Moonshine.}\allowbreak\lean{FockSpan.}\allowbreak\lean{quadF_lead_one}\par}
{\raggedright\noindent File \leanpath{VOA/\allowbreak Moonshine/\allowbreak FockSpanOps.lean}.\par}
\begin{leancode}
noncomputable def Fdw (d : ℕ) (w : ℤ) : Submodule ℂ Fk :=
  restrictTotalDegree (ℕ × Fin (rank leech)) ℂ d ⊓ gradingZN (rank leech) w

theorem quadF_lead (γ : leech.carrier) (a : Fin (rank leech)) (i : ℕ) (b : Fin (rank leech))
    (j : ℕ) {d : ℕ} {w : ℤ} {p : Fk} (hp : p ∈ Fdw d w) :
    quadF γ a i b j 0 p - X (i, a) * X (j, b) * p ∈ Fdw (d + 1) (w + i + j + 2)

theorem quadF_X_mul (γ : leech.carrier) (a b : Fin (rank leech)) {s : ℤ} (hs : -1 ≤ s) (t : ℕ)
    (c : Fin (rank leech)) (p : Fk) :
    quadF γ a 0 b 0 s (X (t, c) * p) = X (t, c) * quadF γ a 0 b 0 s p
      + (if c = b then ((t : ℂ) + 1) • (X (((t : ℤ) + s + 2).toNat, a) * p) else 0)
      + (if c = a then ((t : ℂ) + 1) • (X (((t : ℤ) + s + 2).toNat, b) * p) else 0)

theorem quadF_D_one (γ : leech.carrier) (a b : Fin (rank leech)) :
    quadF γ a 0 b 0 (-1) 1
      = chargeCoord leech γ b • (X (0, a) : Fk) + chargeCoord leech γ a • (X (0, b) : Fk)

theorem quadF_lead_one (γ : leech.carrier) (a c : Fin (rank leech)) (j : ℕ) :
    quadF γ a 0 c j 0 1
      = X (0, a) * X (j, c) + ((-1 : ℂ) ^ j * chargeCoord leech γ c) • (X (j + 1, a) : Fk)
        + (((j : ℂ) + 1) * chargeCoord leech γ a) • (X (j + 1, c) : Fk)
\end{leancode}
\noindent\textbf{Hypotheses / status.} None\par
\medskip
\begin{btheorem}{40.61}[Golay facts]
\begin{enumerate}[label=(\alph*)]
\item Every codeword has weight $0,8,12,16$ or $24$.
\item Two distinct octads meet in at most $4$ points.
\item \textbf{(Steiner system $S(5,8,24)$, existence)} Every $5$-subset of the coordinates lies in an
  octad.
\item A word of even weight not in $\chtwelveG$ has odd intersection with some octad.
\item A codeword of weight $12$ or $16$ meets some octad in exactly $6$ points.
\item Every $4$-set meets some octad in exactly $3$ points.
\end{enumerate}
\end{btheorem}
{\raggedright\noindent\textbf{Lean.} \lean{VOA.}\allowbreak\lean{Moonshine.}\allowbreak\lean{GroundLat.}\allowbreak\lean{golay_weight_cases},\allowbreak\ \lean{VOA.}\allowbreak\lean{Moonshine.}\allowbreak\lean{GroundLat.}\allowbreak\lean{octad_inter_le},\allowbreak\ \lean{VOA.}\allowbreak\lean{Moonshine.}\allowbreak\lean{GroundLat.}\allowbreak\lean{steiner},\allowbreak\ \lean{VOA.}\allowbreak\lean{Moonshine.}\allowbreak\lean{GroundLat.}\allowbreak\lean{exists_octad_dot_ne_zero},\allowbreak\ \lean{VOA.}\allowbreak\lean{Moonshine.}\allowbreak\lean{GroundLat.}\allowbreak\lean{exists_octad_inter_six},\allowbreak\ \lean{VOA.}\allowbreak\lean{Moonshine.}\allowbreak\lean{GroundLat.}\allowbreak\lean{exists_octad_inter_three}\par}
{\raggedright\noindent File \leanpath{VOA/\allowbreak Moonshine/\allowbreak GroundLatticeGolay.lean}.\par}
\begin{leancode}
theorem golay_weight_cases {c : Fin 24 → ZMod 2} (hc : c ∈ golay) :
    weight c = 0 ∨ weight c = 8 ∨ weight c = 12 ∨ weight c = 16 ∨ weight c = 24

theorem octad_inter_le {O₁ O₂ : Finset (Fin 24)} (h₁ : IsOct O₁) (h₂ : IsOct O₂) (hne : O₁ ≠ O₂) :
    (O₁ ∩ O₂).card ≤ 4

theorem steiner {T : Finset (Fin 24)} (hT : T.card = 5) : ∃ O, IsOct O ∧ T ⊆ O

theorem exists_octad_dot_ne_zero {x : Fin 24 → ZMod 2} (hx : x ∉ golay) (hw : Even (weight x)) :
    ∃ O, IsOct O ∧ ind O ⬝ᵥ x ≠ 0

theorem exists_octad_inter_six {C : Finset (Fin 24)} (hC : ind C ∈ golay)
    (h : C.card = 12 ∨ C.card = 16) : ∃ O, IsOct O ∧ (O ∩ C).card = 6

theorem exists_octad_inter_three {S : Finset (Fin 24)} (hS : S.card = 4) :
    ∃ O, IsOct O ∧ (O ∩ S).card = 3
\end{leancode}
\noindent\textbf{Hypotheses / status.} As stated\par
\medskip
\endgroup
\chapter{Existence of the twisted triality II: the explicit construction}
\begingroup
\providecommand{\chthirteenT}{\mathcal{T}}
\providecommand{\chthirteenV}{V_0^{+}}
\providecommand{\chthirteenVt}{V_3^{+}}
\providecommand{\chthirteenTs}[1]{T_{#1}^{+}}
\providecommand{\chthirteeno}{\operatorname{o}}
\noindent{\small 6 of the chapter's 54 numbered statements: its headline results and every definition they use.}\par\medskip
\begin{bproposition}{41.3}[source zero modes]
Let $n\in\Leech$.
\begin{enumerate}[label=(\alph*)]
\item $\chthirteeno(q_{ab})\,x_n=n_an_b\,x_n$.
\item If $\gamma\in\Leech_0$, $n\notin\Leech_0$ and both have norm $4$, then $|\langle\gamma,n\rangle|\le2$ and
\begin{equation*}
  \chthirteeno(u_\gamma)\,x_n=[\langle\gamma,n\rangle=-2]\,\varepsilon_\Leech(\gamma,n)\,x_{\gamma+n}
  +[\langle\gamma,n\rangle=2]\,\varepsilon_\Leech(\gamma,-n)\,x_{\gamma-n}.
\end{equation*}
\end{enumerate}
\end{bproposition}
{\raggedright\noindent\textbf{Lean.} \lean{VOA.}\allowbreak\lean{Moonshine.}\allowbreak\lean{SigmaTop.}\allowbreak\lean{actV_qV},\allowbreak\ \lean{VOA.}\allowbreak\lean{Moonshine.}\allowbreak\lean{SigmaTop.}\allowbreak\lean{form_bound},\allowbreak\ \lean{VOA.}\allowbreak\lean{Moonshine.}\allowbreak\lean{SigmaTop.}\allowbreak\lean{actV_lV}\par}
\begin{leancode}
theorem actV_qV (a b : Fin (rank leech)) (n : leech.carrier) :
    actV (qV a b) 1 (frameX0P n)
      = (chargeCoord leech n a * chargeCoord leech n b) • frameX0P n

theorem form_bound {γ n : leech.carrier} (hγ : γ ∈ lambda0) (hn : n ∉ lambda0)
    (hγ4 : leech.form γ γ = 4) (hn4 : leech.form n n = 4) :
    -2 ≤ leech.form γ n ∧ leech.form γ n ≤ 2

theorem actV_lV {γ n : leech.carrier} (hγ : γ ∈ lambda0) (hn : n ∉ lambda0)
    (hγ4 : leech.form γ γ = 4) (hn4 : leech.form n n = 4) :
    actV (lV γ hγ) 1 (frameX0P n)
      = (if leech.form γ n = -2 then latticeCocycle leech γ n • frameX0P (γ + n) else 0)
        + (if leech.form γ n = 2 then latticeCocycle leech γ (-n) • frameX0P (γ - n) else 0)
\end{leancode}
\noindent File \leanpath{VOA/\allowbreak Moonshine/\allowbreak SigmaTopSource.lean}.
\noindent\textbf{Hypotheses / status.} Norms, $\gamma\in\Leech_0$, $n\notin\Leech_0$\par
\medskip
\begin{bproposition}{41.7}[twisted zero modes]
On $h_c(-\tfrac12)\otimes t$, $t\in T$:
\begin{enumerate}[label=(\alph*)]
\item $\chthirteeno(h_a(-1)h_b(-1)\mathbf 1)\,(h_c\otimes t)=\bigl(\tfrac18\delta_{ab}h_c+\tfrac12\delta_{bc}h_a+\tfrac12\delta_{ac}h_b\bigr)\otimes t$;
\item for $\|\alpha\|^2=4$: $\chthirteeno(e^\alpha)\,(h_c\otimes t)=2^{-4}\bigl(h_c-2\alpha_c\,\alpha(-\tfrac12)\bigr)\otimes e_\alpha t$;
\item $\chthirteeno(e^{-\alpha})=\chthirteeno(e^{\alpha})$ on the twisted module (more generally on all even modes).
\end{enumerate}
\end{bproposition}
{\raggedright\noindent\textbf{Lean.} \lean{VOA.}\allowbreak\lean{Moonshine.}\allowbreak\lean{SigmaTop.}\allowbreak\lean{YML_quadSt},\allowbreak\ \lean{VOA.}\allowbreak\lean{Moonshine.}\allowbreak\lean{SigmaTop.}\allowbreak\lean{YML_echg_X},\allowbreak\ \lean{VOA.}\allowbreak\lean{Moonshine.}\allowbreak\lean{SigmaTop.}\allowbreak\lean{YML_echg_neg_two}\par}
{\raggedright\noindent File \leanpath{VOA/\allowbreak Moonshine/\allowbreak SigmaTopTarget.lean}.\par}
\begin{leancode}
theorem YML_quadSt (a b c : Fin (rank L)) (t : Rep (cocMat L) (isUnit_det_cocMat L hEven hU)) :
    YML L hEven hU (quadSt L hEven hPos a b) 2 ((X (0, c) : FockT (rank L)) ⊗ₜ t)
      = ((if a = b then (1 / 8 : ℂ) else 0) • (X (0, c) : FockT (rank L))
          + (if b = c then (1 / 2 : ℂ) else 0) • X (0, a)
          + (if a = c then (1 / 2 : ℂ) else 0) • X (0, b)) ⊗ₜ t

theorem YML_echg_X (α : L.carrier) (hα : L.form α α = 4) (c : Fin (rank L))
    (t : Rep (cocMat L) (isUnit_det_cocMat L hEven hU)) :
    YML L hEven hU (echg L α) 2 ((X (0, c) : FockT (rank L)) ⊗ₜ t)
      = (2 : ℂ) ^ (-(4 : ℤ)) • ((X (0, c) - (2 * chargeCoord L α c) • vVarT (chargeCoord L α) 0)
          ⊗ₜ eT L hEven hU α t)

theorem YML_echg_neg_two (α : L.carrier) :
    YML L hEven hU (echg L (-α)) 2 = YML L hEven hU (echg L α) 2
\end{leancode}
\noindent\textbf{Hypotheses / status.} Norm 4 (for $e^\alpha$)\par
\medskip
\begin{bproposition}{41.11}[zero modes in the frame]
For $v\in\C^{24}$, $t\in T$, $\alpha,\beta\in\Leech$, $\gamma\in\Leech_0$ of norm $4$:
\begin{enumerate}[label=(\alph*)]
\item $\chthirteeno(d(\alpha,\beta))\,x_n=\langle\alpha,n\rangle\langle\beta,n\rangle\,x_n$;
\item $\chthirteeno(d(\alpha,\beta))\,(h_v\otimes t)=h_{\,\frac{\langle\alpha,\beta\rangle}8v+\frac12\langle v,\beta\rangle\alpha+\frac12\langle v,\alpha\rangle\beta}\otimes t$;
\item $\chthirteeno(u_\gamma)\,(h_v\otimes t)=2^{-3}\,h_{v-2\langle v,\gamma\rangle\gamma}\otimes e_\gamma t$;
\item (diagonal letters) $\chthirteeno(d(w_j,w_j))$ acts by $9$ on block $j$ and by $1$ on every other
  block, on \emph{both} $(\chthirteenVt)_2$ and $(T^+)_2$:
  $\chthirteeno(d(w_j,w_j))x_n=(9\text{ if }j=j(n)\text{ else }1)\,x_n$ and
  $\chthirteeno(d(w_j,w_j))(h_{c_m}\otimes t)=(9\text{ if }j=m\text{ else }1)\,h_{c_m}\otimes t$.
\end{enumerate}
\end{bproposition}
{\raggedright\noindent\textbf{Lean.} \lean{VOA.}\allowbreak\lean{Moonshine.}\allowbreak\lean{SigmaTop.}\allowbreak\lean{actV_dV},\allowbreak\ \lean{VOA.}\allowbreak\lean{Moonshine.}\allowbreak\lean{SigmaTop.}\allowbreak\lean{actT_dV_hP},\allowbreak\ \lean{VOA.}\allowbreak\lean{Moonshine.}\allowbreak\lean{SigmaTop.}\allowbreak\lean{actT_lV_hP},\allowbreak\ \lean{VOA.}\allowbreak\lean{Moonshine.}\allowbreak\lean{SigmaTop.}\allowbreak\lean{actT_dV_fw_hP},\allowbreak\ \lean{VOA.}\allowbreak\lean{Moonshine.}\allowbreak\lean{SigmaTop.}\allowbreak\lean{actV_dV_fw}\par}
\begin{leancode}
theorem actV_dV (α β n : leech.carrier) :
    actV (dV α β) 1 (frameX0P n)
      = (((leech.form α n : ℤ) : ℂ) * ((leech.form β n : ℤ) : ℂ)) • frameX0P n

theorem actT_dV_hP (α β : leech.carrier) (v : Fin (rank leech) → ℂ) (t : TLeechT) :
    actT (dV α β) 1 (hP v t) = hP (dL α β v) t

theorem actT_lV_hP {γ : leech.carrier} (hγ : γ ∈ lambda0) (hγ4 : leech.form γ γ = 4)
    (v : Fin (rank leech) → ℂ) (t : TLeechT) :
    actT (lV γ hγ) 1 (hP v t)
      = ((2 : ℂ) ^ (-(3 : ℤ))) • hP (latL γ v) (eT leech leech_even leech_isUnit_det_gramMat γ t)

theorem actT_dV_fw_hP (j m : Fin 24) (t : TLeechT) :
    actT (dV (fw j) (fw j)) 1 (hP (cw m) t) = (if j = m then (9 : ℂ) else 1) • hP (cw m) t

theorem actV_dV_fw {n : leech.carrier} (hn : n ∉ lambda0) (h4 : leech.form n n = 4) (j : Fin 24) :
    actV (dV (fw j) (fw j)) 1 (frameX0P n) = (if j = blk n then (9 : ℂ) else 1) • frameX0P n
\end{leancode}
\noindent File \leanpath{VOA/\allowbreak Moonshine/\allowbreak SigmaTopFrame.lean}.
\noindent\textbf{Hypotheses / status.} As stated\par
\medskip
\begin{bproposition}{41.23}
For $j\ne k$:
\begin{enumerate}[label=(\alph*)]
\item ($\theta$-invariance) the coefficients of $e^{\gamma_+}$ and $e^{-\gamma_+}$ in $\sigma(d(w_j,w_k))$ are equal,
  and so are those of $e^{\pm\gamma_-}$; consequently $\kappa_+=\tfrac12\varepsilon_N(\alpha_j,\alpha_k)\,t_c(a_+)$,
  $\kappa_-=\tfrac12\varepsilon_k\varepsilon_N(\alpha_j,\alpha_k)\,t_c(a_-)$, and $g_e(a_\pm)^2=\varepsilon_j\varepsilon_k$;
\item $\kappa_\pm^2\,\varepsilon_\Leech(\gamma_\pm,\gamma_\pm)=\tfrac14$;
\item $A_\pm^2=\tfrac14\cdot\mathrm{id}_T$. In particular $A_\pm$ is diagonalisable with eigenvalues in
  $\{\pm\tfrac12\}$.
\end{enumerate}
\end{bproposition}
{\raggedright\noindent\textbf{Lean.} \lean{VOA.}\allowbreak\lean{Moonshine.}\allowbreak\lean{SigmaTop.}\allowbreak\lean{dQ_symm},\allowbreak\ \lean{VOA.}\allowbreak\lean{Moonshine.}\allowbreak\lean{SigmaTop.}\allowbreak\lean{kP_eq},\allowbreak\ \lean{VOA.}\allowbreak\lean{Moonshine.}\allowbreak\lean{SigmaTop.}\allowbreak\lean{eV_aP_sq},\allowbreak\ \lean{VOA.}\allowbreak\lean{Moonshine.}\allowbreak\lean{SigmaTop.}\allowbreak\lean{kP_sq},\allowbreak\ \lean{VOA.}\allowbreak\lean{Moonshine.}\allowbreak\lean{SigmaTop.}\allowbreak\lean{AP_mul_self},\allowbreak\ \lean{VOA.}\allowbreak\lean{Moonshine.}\allowbreak\lean{SigmaTop.}\allowbreak\lean{AM_mul_self}\par}
{\raggedright\noindent File \leanpath{VOA/\allowbreak Moonshine/\allowbreak SigmaTop2Eigen.lean}.\par}
\begin{leancode}
theorem dQ_symm {j k : Fin 24} (hjk : j ≠ k) : dQ j k 0 = dQ j k 3 ∧ dQ j k 1 = dQ j k 2

theorem kP_eq {j k : Fin 24} (hjk : j ≠ k) : kP j k = (1 / 2 : ℂ) * cQ j k 0 * tc (aP j k)

theorem eV_aP_sq {j k : Fin 24} (hjk : j ≠ k) : eV (aP j k) * eV (aP j k) = golayEps j * golayEps k

theorem kP_sq {j k : Fin 24} (hjk : j ≠ k) : kP j k * kP j k * εΛ (gP j k) (gP j k) = 1 / 4

theorem AP_mul_self {j k : Fin 24} (hjk : j ≠ k) : AP j k * AP j k = (1 / 4 : ℂ) • 1

theorem AM_mul_self {j k : Fin 24} (hjk : j ≠ k) : AM j k * AM j k = (1 / 4 : ℂ) • 1
\end{leancode}
\noindent\textbf{Hypotheses / status.} $j\ne k$\par
\medskip
\begin{btheorem}{41.37}[the seed exists]
There is a nonzero $t\in T$ with $A_+(j,k)t=-\tfrac12t$ and $A_-(j,k)t=\tfrac12t$ for all $j\ne k$.
Moreover one can take $t\in T_s$ for some $s$.
\end{btheorem}
{\raggedright\noindent\textbf{Lean.} \lean{VOA.}\allowbreak\lean{Moonshine.}\allowbreak\lean{SigmaTop.}\allowbreak\lean{seedExists},\allowbreak\ \lean{VOA.}\allowbreak\lean{Moonshine.}\allowbreak\lean{SigmaTop.}\allowbreak\lean{seed_in_frame}\par}
{\raggedright\noindent File \leanpath{VOA/\allowbreak Moonshine/\allowbreak SigmaTop2Eigen.lean}.\par}
\begin{leancode}
theorem seedExists : SeedExists

theorem seed_in_frame (h : SeedExists) :
    ∃ s : ZMod 2, ∃ t : TLeechT, t ∈ leechFrameT s ∧ t ≠ 0 ∧ IsSeed t
\end{leancode}
\noindent\textbf{Hypotheses / status.} \textbf{none}\par
\medskip
\begin{btheorem}{41.47}[\lean{classify_S0}]
Every $\gamma\in S_0$ (norm $4$, in $\Leech_0$) is either $\pm\gamma_+(j,k)$ or $\pm\gamma_-(j,k)$ for some $j\ne k$ (shape
$(\pm4,\pm4,0^{22})$), or $\gamma=\gamma_S$ for an octad $D$ and an even $S\subseteq D$ (shape $(\pm2^8,0^{16})$).
\end{btheorem}
{\raggedright\noindent\textbf{Lean.} \lean{VOA.}\allowbreak\lean{Moonshine.}\allowbreak\lean{SigmaTop.}\allowbreak\lean{classify_S0}\par}
\begin{leancode}
theorem classify_S0 {γ : leech.carrier} (hγ : γ ∈ halfFourChi 0) :
    (∃ j k : Fin 24, j ≠ k ∧ (γ = gP j k ∨ γ = -gP j k ∨ γ = gM j k ∨ γ = -gM j k)) ∨
    (∃ (D : Finset (Fin 24)) (hDA : ind D ∈ Leech.subA), D.card = 8 ∧
      ∃ S : Finset (Fin 24), _root_.Even S.card ∧ S ⊆ D ∧ γ = inclΛ (aSA hDA S))
\end{leancode}
\noindent File \leanpath{VOA/\allowbreak Moonshine/\allowbreak SigmaTop2Class.lean}.
\noindent\textbf{Hypotheses / status.} $\gamma\in S_0$\par
\medskip
\endgroup
\chapter{The key identity and the completion of the proof}
\begingroup
\providecommand{\chXIVsig}{\sigma^{\natural}}
\providecommand{\chXIVA}{V_0^{+}}
\providecommand{\chXIVB}{V_3^{+}}
\noindent{\small 7 of the chapter's 56 numbered statements: its headline results and every definition they use.}\par\medskip
\begin{btheorem}{42.3}[One vector generates $T^+$]
For every nonzero $t_0\in T$, the single vector $x_0(t_0)=h_0(-\tfrac12)\mathbb 1\otimes t_0$
mode-spans $(V_\Leech^T)^+$.
\end{btheorem}
{\raggedright\noindent\textbf{Lean.} \lean{VOA.}\allowbreak\lean{Moonshine.}\allowbreak\lean{M8hKey.}\allowbreak\lean{modeSpan_x0T}\par}
\begin{leancode}
theorem modeSpan_x0T {t₀ : TLeechT} (h0 : t₀ ≠ 0) : ModeSpan (fun _ : Unit => x0T t₀)
\end{leancode}
\noindent File \leanpath{VOA/\allowbreak Moonshine/\allowbreak M8hKeySpan.lean}.
\noindent\textbf{Hypotheses / status.} $t_0\ne0$\par
\medskip
\begin{btheorem}{42.9}[$\chXIVsig$ is an involution compatible with $\chXIVA$]
For triality data $D$:
\begin{enumerate}[label=(\alph*)]
\item $(\chXIVsig)^2=1$;
\item for every $u\in\chXIVA$ and $k\in\Z$: $\chXIVsig\,Y^\natural(u)_k=Y^\natural(\sigma u)_k\,\chXIVsig$,
  where $Y^\natural(u)$ is the diagonal field $\operatorname{diag}(Y(u,z),Y^T(u,z))$;
  consequently $\chXIVsig Y^\natural(u)\chXIVsig=Y^\natural(\sigma u)$;
\item $\chXIVsig$ commutes with the translation operator $L_{-1}$ of $\Vnat$, and fixes the vacuum.
\end{enumerate}
\end{btheorem}
{\raggedright\noindent\textbf{Lean.} \lean{VOA.}\allowbreak\lean{Moonshine.}\allowbreak\lean{M8hKey.}\allowbreak\lean{sigNat_sq},\allowbreak\ \lean{VOA.}\allowbreak\lean{Moonshine.}\allowbreak\lean{M8hKey.}\allowbreak\lean{sigNat_YnatU_V0},\allowbreak\ \lean{VOA.}\allowbreak\lean{Moonshine.}\allowbreak\lean{M8hKey.}\allowbreak\lean{conjF_YnatU_V0},\allowbreak\ \lean{VOA.}\allowbreak\lean{Moonshine.}\allowbreak\lean{M8hKey.}\allowbreak\lean{sigNat_Tnat},\allowbreak\ \lean{VOA.}\allowbreak\lean{Moonshine.}\allowbreak\lean{M8hKey.}\allowbreak\lean{sigNat_vac}\par}
\begin{leancode}
theorem sigNat_sq : sigNat D * sigNat D = 1

theorem sigNat_YnatU_V0 (u : V0Sub.carrier) (k : ℤ) :
    sigNat D * YnatU (u : VPlus) k = YnatU ((sig u : V0Sub.carrier) : VPlus) k * sigNat D

theorem conjF_YnatU_V0 (u : V0Sub.carrier) :
    conjF (sigNat D) (YnatU (u : VPlus)) = YnatU ((sig u : V0Sub.carrier) : VPlus)

theorem sigNat_Tnat : sigNat D * Tnat = Tnat * sigNat D

theorem sigNat_vac : sigNat D vacNat = vacNat
\end{leancode}
\noindent File \leanpath{VOA/\allowbreak Moonshine/\allowbreak M8hKeySigma.lean}.
\noindent\textbf{Hypotheses / status.} $D$\par
\medskip
\begin{btheorem}{42.10}[Circularity-free uniqueness]
Let $F$ be a formal distribution with coefficients in $\End\Vnat$ that is local with
$Y^\natural(u)$ for every $u\in\chXIVA$.  Fix $\beta\in\Leech\setminus\Leech_0$ and nonzero
$t_0\in T_0$, $t_1\in T_1$.  If every coefficient of $F$ kills the four vectors
\begin{equation*}
  \mathbb 1,\qquad (x_\beta,0),\qquad (0,x_0(t_0)),\qquad (0,x_0(t_1)),
\end{equation*}
then $F=0$.
\end{btheorem}
{\raggedright\noindent\textbf{Lean.} \lean{VOA.}\allowbreak\lean{Moonshine.}\allowbreak\lean{M8hKey.}\allowbreak\lean{eq_zero_of_V0local}\par}
\begin{leancode}
theorem eq_zero_of_V0local (F : FormalDist (End ℂ Vnat))
    (hloc : ∀ u : V0Sub.carrier, Local F (YnatU (u : VPlus)))
    {β : leech.carrier} (hβ : β ∉ lambda0) (t : ZMod 2 → TLeechT)
    (ht : ∀ s, t s ∈ leechFrameT s) (ht0 : ∀ s, t s ≠ 0)
    (hvac : ∀ p, F p vacNat = 0) (h3 : ∀ p, F p (frameX0P β, 0) = 0)
    (hT : ∀ s p, F p (0, x0T (t s)) = 0) : F = 0
\end{leancode}
\noindent File \leanpath{VOA/\allowbreak Moonshine/\allowbreak M8hKeySigma.lean}.
\noindent\textbf{Hypotheses / status.} Cyclic data\par
\medskip
\begin{btheorem}{42.13}[(4.23) from the key identities]
If \lean{SigmaKey D β t} holds, then (i) $\chXIVsig Y^\natural(\Psi x)\chXIVsig=Y^\natural((0,x))$ for
$x=x_0(t_s)$; (ii) (4.23) holds for \emph{all} pairs $t_1,t_2\in T^+$; (iii) $\Vnat$ carries all the
structure listed in the final theorem.
\end{btheorem}
{\raggedright\noindent\textbf{Lean.} \lean{VOA.}\allowbreak\lean{Moonshine.}\allowbreak\lean{M8hKey.}\allowbreak\lean{SigmaKey.}\allowbreak\lean{conj_eq},\allowbreak\ \lean{VOA.}\allowbreak\lean{Moonshine.}\allowbreak\lean{M8hKey.}\allowbreak\lean{tpBlockLocal_of_sigma},\allowbreak\ \lean{VOA.}\allowbreak\lean{Moonshine.}\allowbreak\lean{M8hKey.}\allowbreak\lean{moonshine_of_sigma}\par}
\begin{leancode}
theorem SigmaKey.conj_eq {β : leech.carrier} {t : ZMod 2 → TLeechT} (h : SigmaKey D β t) :
    conjF (sigNat D) (YnatU (D.Ψ (x0T (t D.s)))) = YnatT leechWbar (x0T (t D.s))

theorem tpBlockLocal_of_sigma {β : leech.carrier} {t : ZMod 2 → TLeechT} (h : SigmaKey D β t)
    (t₁ t₂ : TP) : TPBlockLocal leechWbar t₁ t₂

theorem moonshine_of_sigma {β : leech.carrier} {t : ZMod 2 → TLeechT} (h : SigmaKey D β t) :
    ∃ V : VertexOperatorAlgebra Vnat,
      V.Y = Ynat VnatAxioms.wbarL ∧ V.vacuum = vacNat ∧ V.conformalVector = omegaNat ∧
      V.centralCharge = 24 ∧
      V.grading 1 = ⊥ ∧
      Module.finrank ℂ (V.grading 2) = 196884 ∧
      (∀ n : ℕ, (Module.finrank ℂ (V.grading (n : ℤ)) : ℂ) = PowerSeries.coeff n JShift) ∧
      V.IsSimple
\end{leancode}
\noindent File \leanpath{VOA/\allowbreak Moonshine/\allowbreak M8hKey.lean}.
\noindent\textbf{Hypotheses / status.} \lean{SigmaKey}\par
\medskip
\begin{btheorem}{42.14}[Automorphisms are isometries]
Let $W$ be a VOA with $\mathbb 1\ne0$ and $L_1W_1=0$, $B$ a fully invariant normalised form, and
$\alpha\in\Aut W$.  Then $B(\alpha a,\alpha b)=B(a,b)$ for all $a,b$.
\end{btheorem}
{\raggedright\noindent\textbf{Lean.} \lean{VOA.}\allowbreak\lean{Moonshine.}\allowbreak\lean{SigmaKey.}\allowbreak\lean{aut_isometry}\par}
\begin{leancode}
theorem aut_isometry {W : VertexOperatorAlgebra 𝒱} (hne : W.vacuum ≠ 0)
    (hqp : ∀ v, v ∈ W.grading 1 → W.L 1 v = 0) {B : InvariantForm W}
    (hB : B.IsFullyInvariant) (hn : B.Normalized) (α : Aut W) (a b : 𝒱) :
    B.form (α a) (α b) = B.form a b
\end{leancode}
\noindent File \leanpath{VOA/\allowbreak Moonshine/\allowbreak SigmaKeyForm.lean}.
\noindent\textbf{Hypotheses / status.} CFT type, $L_1V_1=0$\par
\medskip
\begin{btheorem}{42.45}[The Frenkel--Zhu reduction and its converse]
\begin{enumerate}[label=(\alph*)]
\item \lean{BimodGenV3} and \lean{TopPairs D} imply \lean{KeyTopAll D}.
\item \lean{KeyTopAll D} implies \lean{TopPairs D} (so \lean{TopPairs} is also necessary).
\item Hence \lean{BimodGenV3} and \lean{TopPairs D} imply \lean{KeySkew D t}, and with \lean{PhiTop},
  \lean{ThetaTop}, all targets.
\end{enumerate}
\end{btheorem}
{\raggedright\noindent\textbf{Lean.} \lean{VOA.}\allowbreak\lean{Moonshine.}\allowbreak\lean{FZKey.}\allowbreak\lean{keyTopAll_of_FZ},\allowbreak\ \lean{VOA.}\allowbreak\lean{Moonshine.}\allowbreak\lean{FZKey.}\allowbreak\lean{topPairs_of_keyTopAll},\allowbreak\ \lean{VOA.}\allowbreak\lean{Moonshine.}\allowbreak\lean{FZKey.}\allowbreak\lean{keySkew_of_FZ},\allowbreak\ \lean{VOA.}\allowbreak\lean{Moonshine.}\allowbreak\lean{FZKey.}\allowbreak\lean{moonshine_of_FZ}\par}
\begin{leancode}
theorem keyTopAll_of_FZ (hgen : BimodGenV3) (htop : TopPairs D) : KeyTopAll D

theorem topPairs_of_keyTopAll (h : KeyTopAll D) : TopPairs D

theorem keySkew_of_FZ {t : ZMod 2 → TLeechT} (ht : ∀ s, t s ∈ leechFrameT s) (ht0 : ∀ s, t s ≠ 0)
    (hgen : BimodGenV3) (htop : TopPairs D) : KeySkew D t

theorem moonshine_of_FZ {β : leech.carrier} (hβ : β ∉ lambda0) {t : ZMod 2 → TLeechT}
    (ht : ∀ s, t s ∈ leechFrameT s) (ht0 : ∀ s, t s ≠ 0)
    (hΦ : PhiTop D t) (hΘ : ThetaTop D t) (hgen : BimodGenV3) (htop : TopPairs D) :
    ∃ V : VertexOperatorAlgebra Vnat,
      V.Y = Ynat VnatAxioms.wbarL ∧ V.vacuum = vacNat ∧ V.conformalVector = omegaNat ∧
      V.centralCharge = 24 ∧
      V.grading 1 = ⊥ ∧
      Module.finrank ℂ (V.grading 2) = 196884 ∧
      (∀ n : ℕ, (Module.finrank ℂ (V.grading (n : ℤ)) : ℂ) = PowerSeries.coeff n JShift) ∧
      V.IsSimple
\end{leancode}
{\raggedright\noindent File \leanpath{VOA/\allowbreak Moonshine/\allowbreak FZBimoduleKey.lean}.\par}
\noindent\textbf{Hypotheses / status.} \lean{BimodGenV3}, \lean{TopPairs}\par
\medskip
\begin{btheorem}{42.58}[The Moonshine module]
There is a vertex operator algebra structure on $\Vnat$ whose vertex operator is the block operator
(42.1), with vacuum $(\mathbb 1,0)$, conformal vector $(\omega,0)$, central charge $24$,
$\Vnat_1=0$, $\dim\Vnat_n=$ the coefficient of $q^n$ in the formal $q$-expansion $q\,J(q)$ for all $n\ge0$
(so $\dim\Vnat_2=196884$), and which is simple.  These data determine the structure uniquely, so the
theorem constructs a specific object.  The Lean statement and its full fidelity check are in
\S42.8.
\end{btheorem}
{\raggedright\noindent\textbf{Lean.} \lean{VOA.}\allowbreak\lean{Moonshine.}\allowbreak\lean{moonshine}\par}
\begin{leancode}
theorem moonshine :
    ∃ V : VertexOperatorAlgebra Vnat,
      V.Y = Ynat VnatAxioms.wbarL ∧ V.vacuum = vacNat ∧ V.conformalVector = omegaNat ∧
      V.centralCharge = 24 ∧
      V.grading 1 = ⊥ ∧
      Module.finrank ℂ (V.grading 2) = 196884 ∧
      (∀ n : ℕ, (Module.finrank ℂ (V.grading (n : ℤ)) : ℂ) = PowerSeries.coeff n JShift) ∧
      V.IsSimple :=
  moonshine_of_topPairs (by norm_num) topPairs_one
\end{leancode}
{\raggedright\noindent File \leanpath{VOA/\allowbreak Moonshine/\allowbreak TopPairs.lean}.\par}
\medskip
\endgroup
\part{Verification}
\chapter{Verification, dead ends, numerics and open problems}
\begingroup
\providecommand{\chXVsq}{\mathbin{\ltimes}}
\noindent{\small 3 of the chapter's 5 numbered statements: its headline results and every definition they use.}\par\medskip
\begin{btheorem}{43.2}[The all-indices twisted Jacobi identity forces $M=0$]
Let $V$ be a vertex operator algebra, $g\in\Aut V$, $N\ge 2$, and let $Y_M$ be a linear map from $V$
to fractional-mode families of endomorphisms of a vector space $M$ such that $Y_M(\mathbf 1,z)=\mathrm{id}$
(only the mode $\mathbf 1_{(-1)}$, label $k=-N$, is nonzero, and it is the identity). If $Y_M$
satisfies the twisted Jacobi identity for all $j,u\in V^j,v,x,\ell,k,m$ without restriction on $k$,
then $M=0$.
\end{btheorem}
{\raggedright\noindent\textbf{Lean.} \lean{VOA.}\allowbreak\lean{jacobiAllIndices_eq_zero}\par}
{\raggedright\noindent File \leanpath{VOA/\allowbreak Moonshine/\allowbreak HeisenbergTwistedModuleDefect.lean}:\par}
\begin{leancode}
theorem jacobiAllIndices_eq_zero {g : Aut W} {N : ℕ} (hN : 2 ≤ N) {M : Type*}
    [AddCommGroup M] [Module ℂ M] (YM : 𝒱 →ₗ[ℂ] (ℤ → End ℂ M))
    (hvac : ∀ k : ℤ, YM W.vacuum k = if k = -(N : ℤ) then 1 else 0)
    (hJ : JacobiAllIndices g N YM) (x : M) : x = 0 := by
\end{leancode}
\noindent\textbf{Hypotheses / status.} Unconditional\par
\medskip
\begin{btheorem}{43.3}[A $\C$-bilinear form is never positive on a nonzero $\C$-subspace]
Let $\beta$ be a $\C$-bilinear form on a complex vector space $E$ and $S\le E$ a subspace. If
$\beta(x,x)$ is a positive real number for every nonzero $x\in S$, then $S=0$.
\end{btheorem}
{\raggedright\noindent\textbf{Lean.} \lean{VOA.}\allowbreak\lean{eq_bot_of_bilinear_pos}\par}
{\raggedright\noindent File \leanpath{VOA/\allowbreak Core/\allowbreak GradedUnitarity.lean}:\par}
\begin{leancode}
theorem eq_bot_of_bilinear_pos (β : BilinForm ℂ E) (S : Submodule ℂ E)
    (hpos : ∀ x ∈ S, x ≠ 0 → ∃ r : ℝ, 0 < r ∧ β x x = (r : ℂ)) : S = ⊥ := by
\end{leancode}
\noindent\textbf{Hypotheses / status.} Unconditional\par
\medskip
\begin{btheorem}{43.6}[golayE is not determined]
There are two characters $e_1,e_2$ of $\Lambda_0$ with values in $\C^\times$, both satisfying the
specification $\psi=g_e\,\theta\,g_e^{-1}$, whose values at the preimage $w_0$ of $8e_0$ differ by a
sign.
\end{btheorem}
{\raggedright\noindent\textbf{Lean.} \lean{VOA.}\allowbreak\lean{Moonshine.}\allowbreak\lean{TrialityTw.}\allowbreak\lean{golayE_not_determined}\par}
{\raggedright\noindent File \leanpath{VOA/\allowbreak Moonshine/\allowbreak TrialityTwSigma.lean}:\par}
\begin{leancode}
theorem golayE_not_determined :
    ∃ e₁ e₂ : lam0.carrier →+ Additive ℂˣ,
      (∀ x, lam0Psi x = charLat lam0 e₁ (thetaLat lam0 (charLat lam0 (-e₁) x))) ∧
      (∀ x, lam0Psi x = charLat lam0 e₂ (thetaLat lam0 (charLat lam0 (-e₂) x))) ∧
      ∃ w₀ : lam0.carrier, inclΛ w₀ = frameW ∧ charVal lam0 e₂ w₀ = -charVal lam0 e₁ w₀ := by
\end{leancode}
\noindent\textbf{Hypotheses / status.} Unconditional\par
\medskip
\endgroup
\end{document}
